\documentclass[11pt]{amsart}
\usepackage[margin=27mm]{geometry}
\usepackage{amsmath,amssymb}
\usepackage{tikz}

% amsart loads amsthm, whose proof environment and \qed clash with pf2's.
% We use pf2's proof environment throughout, as Lamport's format requires.
\let\proof\relax
\let\endproof\relax
\let\qed\relax
\usepackage{pf2}

\usepackage[colorlinks=true,linkcolor=blue,citecolor=blue,urlcolor=blue]{hyperref}

% The author's macros.
%
% some calligraphy
%
\newcommand\cA{{\mathcal A}}
\newcommand\cB{{\mathcal B}}
\newcommand\cC{{\mathcal C}}
\newcommand\cD{{\mathcal D}}
\newcommand\cE{{\mathcal E}}
\newcommand\cF{{\mathcal F}}
\newcommand\cG{{\mathcal G}}
\newcommand\cH{{\mathcal H}}
\newcommand\cI{{\mathcal I}}
\newcommand\cJ{{\mathcal J}}
\newcommand\cK{{\mathcal K}}
\newcommand\cL{{\mathcal L}}
\newcommand\cM{{\mathcal M}}
\newcommand\cN{{\mathcal N}}
\newcommand\cP{{\mathcal P}}
\newcommand\cQ{{\mathcal Q}}
\newcommand\cR{{\mathcal R}}
\newcommand\cS{{\mathcal S}}
\newcommand\cT{{\mathcal T}}
\newcommand\cU{{\mathcal U}}
\newcommand\cV{{\mathcal V}}

%
% fields and rings (and a semigroup)
%
\newcommand\cx{{\mathbb C}}% complexes
\newcommand\fld{{\mathbb F}}
\newcommand\flde{{\mathbb E}}
\newcommand\ints{{\mathbb Z}}
\newcommand\nn{{\mathbb N}}%non-negative integers
\newcommand\re{{\mathbb R}}%reals
\newcommand\rats{{\mathbb Q}}

\newcommand\Zu{{\mathbf u}}
\newcommand\Zv{{\mathbf v}}
\newcommand\Zw{{\mathbf w}}
\newcommand\Zy{{\mathbf y}}
\newcommand\Zz{{\mathbf z}}
\newcommand\Zx{{\mathbf x}}
\newcommand\Zj{{\mathbf j}}


%
% the really useful stuff
%
%\newcommand\opk[1]{\mathop{\hbox{\rm #1}}\nolimits} %% old version
\newcommand\opk[1]{\mathop{\mathrm{#1}}\nolimits}
\newcommand\comp[1]{{\mkern2mu\overline{\mkern-2mu#1}}}
\newcommand\diff{\mathbin{\mkern-1.5mu\setminus\mkern-1.5mu}}% for \setminus
\newcommand\gbin[2]{\genfrac{[}{]}{0pt}{}{#1}{#2}}
\newcommand\inv{^{-1}}
\newcommand\res{\mathbin{\mkern-2.0mu\restriction\mkern-2.0mu}}% ch2
\newcommand\sbs{\subseteq}
\newcommand\sps{\supseteq}
\newcommand\seq[3]{#1_{#2},\ldots,#1_{#3}}
\newcommand\supp{\opk{supp}}
\newcommand\im{\opk{im}}
\newcommand\row{\opk{row}}
\newcommand\pmat[1]{\begin{pmatrix} #1 \end{pmatrix}}
\newcommand\cprod{\mathbin{\square}}
%
% matrix theory
%
\newcommand\ip[2]{\langle#1,#2\rangle}
\newcommand\one{{\bf1}}
\newcommand\rk{\opk{rk}}
\newcommand\tr{\opk{tr}}
\newcommand\col{\opk{col}}
\newcommand\mat[3]{\mathrm{Mat}_{#1\times #2}(#3)}

%
% some group theory
%
\newcommand\aut[1]{{\rm Aut}(#1)}
\newcommand\fx[1]{{\rm fix}(#1)}% ch2
\newcommand\grp[1]{\langle #1\rangle}
\newcommand\nrml{\vartriangleleft}
\newcommand\nrmleq{\trianglelefteq}
\newcommand\alt[1]{\opk{Alt}(#1)}
\DeclareMathOperator{\Sym}{Sym}
% End of the author's macros.
% Macros of the blueprint. The author's own macros are at the start of each document; do not duplicate them here.

% The macros of leanblueprint, as no-ops under pdflatex, so that its dependency graph can be
% generated later: \lean{Decl.name}, \uses{label, label}, \leanok.
\newcommand{\lean}[1]{}
\newcommand{\uses}[1]{}
\newcommand{\leanok}{}

% Results that carry the paper's own number, for instance
%   \begin{cgss}{Lemma}{2.8}\label{lem:2.8} ... \end{cgss}
% and results of ours, numbered within the section, through the usual counters.
% Statements are set upright, since pf2's keywords (Assume, Prove) are small capitals, which have
% no italic shape in Computer Modern.
\theoremstyle{definition}
\makeatletter
\newtheorem*{cgss@env}{\cgss@name}
\newcommand{\cgss@name}{}
% The environment sets an anchor and makes the paper's number the text of \ref, so that
% \ref{lem:2.9} prints 2.9 and links to the lemma.
\newenvironment{cgss}[2]{\renewcommand{\cgss@name}{#1 #2}\begin{cgss@env}%
  \phantomsection\def\@currentlabel{#2}}{\end{cgss@env}}
\makeatother
\newtheorem{lemma}{Lemma}[section]
\newtheorem{proposition}[lemma]{Proposition}
\newtheorem{corollary}[lemma]{Corollary}
\newtheorem{definition}[lemma]{Definition}
\newtheorem*{hypothesis}{Hypothesis}
\theoremstyle{remark}
\newtheorem*{remark}{Remark}

% Citing assumptions as in [L95]: \asm{1} gives <1>, the assumption of the enclosing level-1
% step, and \asm[2]{0} gives <0>:2, the second assumption of the theorem itself. Steps are
% cited with pf2's \stepref{label}.
\newcommand{\asm}[2][]{$\langle#2\rangle$\if\relax\detokenize{#1}\relax\else{:}#1\fi}

% Tags of the structured proofs: a step that departs from the paper, and a step that fills a gap.
\newcommand{\new}{\textbf{[new]}}
\newcommand{\repair}{\textbf{[repair]}}

% The paper's sets of Hypothesis 2.4.
\newcommand{\Ga}{\Gamma_a}
\newcommand{\Gb}{\Gamma_b}
\newcommand{\Gc}{\Gamma_c}

% Vectors that the author's macros do not provide, in its pattern: the vectors a, b, c of Hypothesis 2.4,
% a unit vector e, and roots r, s, t.
\newcommand\Za{{\mathbf a}}
\newcommand\Zb{{\mathbf b}}
\newcommand\Zc{{\mathbf c}}
\newcommand\Ze{{\mathbf e}}
\newcommand\Zr{{\mathbf r}}
\newcommand\Zs{{\mathbf s}}
\newcommand\Zt{{\mathbf t}}

% Additional vector names for the arbitrary adjacent pair in Lemma 2.10.
\newcommand\Zp{{\mathbf p}}
\newcommand\Zq{{\mathbf q}}

% Upright Dynkin type letters, following cgss/notation.md: \typ{A}_\ell, \typ{E}_8.
\newcommand{\typ}[1]{\mathrm{#1}}

% Auxiliary results have descriptive headings, not numbers that could clash with CGSS.
% For example, \begin{namedresult}{Corollary}{of Lemma 2.8} inherits cgss's unique anchor;
% \ref prints the descriptive second argument. Existing result environments are unchanged.
\newenvironment{namedresult}[2]{\begin{cgss}{#1}{#2}}{\end{cgss}}

% The all-positive half root used in the Section 3 spanning and inclusion proofs.
\newcommand\Zh{{\mathbf h}}

% A visible citation with its exact Lean name as metadata and a pdfLaTeX-safe display.
% Unlike \path, the display may render Unicode subscripts and Greek letters as TeX math.
% Example: \leancite{LineSystem.E₈}{\texttt{LineSystem.E\ensuremath{{}_8}}}.
\newcommand{\leancite}[2]{#2}

% Local layout for the merged structured proofs; existing section settings are restored.
\newenvironment{mergedproofsection}{\begingroup\emergencystretch=3em}{\par\endgroup}

% A descriptive library-fact label (the paragraph already supplies its hyperlink anchor).
\makeatletter
\newcommand{\factlabel}[2]{\def\@currentlabel{#1}\label{#2}}
\makeatother

% Section 5 (spectral characterization). A citation of the spectral toolkit
% cgss/proofs/spectraltools/spectraltools.tex, which the blueprint does not input: \ST{Theorem 6.2} prints
% [ST, Theorem 6.2]. Its structured version spectraltools-lamport.tex has the same numbering.
\newcommand{\ST}[1]{[ST,~#1]}
% The spectrum of a matrix or a graph, as a multiset, and the multiplicity of an eigenvalue:
% \Spec(X), \mult_X(\mu). A multiset is written \ms{\mu^{m},\dots}, where \mu^{m} means m copies.
\newcommand{\Spec}{\opk{Spec}}
\newcommand{\mult}{\opk{mult}}
\newcommand{\ms}[1]{\left[\!\left[#1\right]\!\right]}
% The largest eigenvalue of a real symmetric matrix or of a graph, as in the toolkit.
\newcommand{\lmax}{\lambda_{\max}}


\begin{document}

\title{Line graphs, root systems, and elliptic geometry: structured proofs}
\date{2 November 2026}

\begin{abstract}
  We give structured proofs, in the style of Lamport, of every numbered result of Cameron,
  Goethals, Seidel and Shult, \emph{Line graphs, root systems, and elliptic geometry} [CGSS], and
  of the results from other sources that the formalization in Lean needs. Each result carries
  the paper's number, the name of its Lean declaration and the results it uses.
\end{abstract}

\maketitle
\tableofcontents
\clearpage

\section{Conventions}

The purpose of this section is not to give full definitions and background, but rather to
establish notation; we defer to [CGSS] for the mathematics and to \verb|cgss/notation.md| for
the dictionary between the notation of the paper, ours and the Lean names.

\subsection*{Structured proofs}
We follow the format of Lamport [L95]. A result is stated with numbered hypotheses under
\textsc{Assume} and its conclusion under \textsc{Prove}. A \textsc{Proof sketch} may follow; it
is an informal road map. The proof is a list of steps $\langle1\rangle1$, $\langle1\rangle2$,
\dots, ending in a \textsc{Qed} step, and each step at level $k$ has its own proof. The
assumption of the enclosing level-$k$ step is cited as $\langle k\rangle$, its numbered parts as
$\langle k\rangle{:}i$, and those of the result itself as $\langle0\rangle{:}i$. Definitions are
made by \textsc{Let}, witnesses are introduced by \textsc{Pick}, and case splits use
\textsc{Case}, with the proof of the closing \textsc{Qed} explaining why the cases are exhaustive.
We expand a step until each leaf is a single inference from named facts: definitions, earlier
steps, earlier results, Mathlib lemmas, or numbered computations. A step that departs from the
paper's argument is tagged \new, and a step that fills a gap in it is tagged \repair.

Each result records the name of its Lean declaration with \verb|\lean{...}| and the results it
uses with \verb|\uses{...}|; these macros print nothing.

\subsection*{Notation}
We will use functional notation for the entries of matrices and vectors; the $uv$-entry of a
matrix $M$ is written $M(u,v)$, and the entry of a vector $\Zv$ indexed by $u$ is written
$\Zv(u)$. Graphs are $X$, $Y$ and $Z$. A set of lines at $60^\circ$ and $90^\circ$ is recorded by
the set $\Phi$ of vectors of squared norm~2 on its lines, taken with both signs, so that
$\Phi=-\Phi$. We keep the paper's names for the paper's own objects: $\Sigma$, $\Ga$, $\Gb$,
$\Gc$, $\Delta$, and the sets $K$, $L$, $M$, $N$ of the proof of Lemma 2.10. The paper uses three
graph conventions; \verb|cgss/notation.md| gives them three names, which this blueprint uses.

\section{Lines at \texorpdfstring{$60^\circ$ and $90^\circ$}{60 and 90 degrees}}\label{sec:2}
% Long structured steps should leave space at a page foot instead of stretching proof spacing.
\raggedbottom

We follow Section 2 of [CGSS], with the vector encoding of Section~1.
We use the graph names and Lean names of \path{cgss/notation.md}; in Lean,
\texttt{Hyp24} is in the namespace \texttt{LineSystem}. In what follows,
$V$ is a finite-dimensional real inner product space. We write $(\Zx,\Zy)$ for its inner
product, as the paper does, since angle brackets are taken by step numbers.
A \emph{line system} $S$ in $V$ is a finite set $\Phi\subseteq V$ such that
\[
 \Phi=-\Phi,\qquad (\Zr,\Zr)=2\quad(\Zr\in\Phi),\qquad
 (\Zr,\Zs)\in\{-2,-1,0,1,2\}\quad(\Zr,\Zs\in\Phi).
\]
\label{def:line-system}\lean{LineSystem}\uses{}
We count lines as antipodal pairs $\{\Zr,-\Zr\}$, so the number of lines is $|\Phi|/2$.
\lean{LineSystem.numLines}
We do not require $\Phi$ to span $V$ until Theorem 2.13.
We call $S$ \emph{indecomposable} if, for all $P,Q\subseteq V$ with $P\cup Q=\Phi$ and
$(\Zr,\Zs)=0$ for all $\Zr\in P$, $\Zs\in Q$, we have $P=\varnothing$ or $Q=\varnothing$.
\label{def:indecomposable}\lean{LineSystem.IsIndecomposable}\uses{def:line-system}
\new\ This intrinsic definition agrees with the paper's orthogonal-subspace definition
when $\Phi$ spans $V$; without spanning, the paper's literal wording allows the partition
into the whole system in its proper span and an empty part. We also require finiteness
explicitly: in finite dimension this loses no line systems, since their normalized signed
vectors are pairwise at distance at least 1 on the compact unit sphere.
An isomorphism of line systems is a bijection of their vector sets preserving inner products.
\label{def:line-iso}\lean{LineSystem.Iso}\uses{def:line-system}

\paragraph{Elementary identities.}\label{facts:inner}
We will use bilinearity, symmetry, nonnegativity of squared norm, and the implication
$(\Zx,\Zx)=0\Rightarrow\Zx=0$. These are, respectively, Mathlib's\par
\begingroup\raggedright
\texttt{inner\_add\_left}, \texttt{inner\_add\_right},
\texttt{inner\_sub\_left}, \texttt{inner\_sub\_right},
\texttt{inner\_neg\_left}, \texttt{inner\_neg\_right},
\texttt{real\_inner\_smul\_left}, \texttt{real\_inner\_smul\_right},
\texttt{real\_inner\_comm}, \texttt{real\_inner\_self\_nonneg}, and
\texttt{inner\_self\_eq\_zero}, all in
\path{Mathlib/Analysis/InnerProductSpace/Basic.lean}.\par\endgroup
Here \texttt{inner\_self\_eq\_zero} is the equivalence
$(\Zx,\Zx)=0\Longleftrightarrow\Zx=0$; its \texttt{NormedAddCommGroup}
hypothesis holds for our real inner product spaces.
In particular, for vectors of squared norm 2,
\begin{equation}\label{eq:endpoints}
 (\Zr,\Zs)=2\Longleftrightarrow\Zr=\Zs,
 \qquad (\Zr,\Zs)=-2\Longleftrightarrow\Zr=-\Zs.
\end{equation}
Indeed, $(\Zr-\Zs,\Zr-\Zs)=4-2(\Zr,\Zs)$ and
$(\Zr+\Zs,\Zr+\Zs)=4+2(\Zr,\Zs)$ by bilinearity; the forward implications follow
from zero squared norm, and the reverse implications follow by substitution.
Thus distinct antipodal pairs represent distinct lines with angles $60^\circ$ or $90^\circ$.
We will also use the \emph{finite-span lemma}: for a finite family $(\Zr_i)_{i\in I}$
in $V$, every member of its span is $\sum_{i\in I}\alpha(i)\Zr_i$ for some
$\alpha:I\to\re$. This is
\path{Submodule.mem_span_range_iff_exists_fun} in
\path{Mathlib/LinearAlgebra/Finsupp/LinearCombination.lean}.

\begin{cgss}{Definition}{2.1}\label{def:2.1}
 \lean{LineSystem.IsStarClosed}\uses{def:line-system}
 \assume{$S$ is a line system in a finite-dimensional real inner product space $V$.}
 We call $S$ \emph{star-closed} if
 \[
  \forall\Zr,\Zs\in\Phi,\quad (\Zr,\Zs)=-1\ \Longrightarrow\ \Zr+\Zs\in\Phi.
 \]
 The third line of the star on $\Zr$ and $\Zs$ is the line of $\Zr+\Zs$:
 its squared norm is $2+2-2=2$, and its inner product with each summand is 1.
 The three vectors lie in the plane spanned by $\Zr,\Zs$.
\end{cgss}

\begin{cgss}{Definition}{2.2}\label{def:2.2}
 \lean{LineSystem.IsMaximal}\uses{def:line-system}
 \assume{$S$ is a line system in a finite-dimensional real inner product space $V$.}
 We call $S$ \emph{maximal} if, for every $\Zv\in\operatorname{span}(\Phi)$,
 \[
 (\Zv,\Zv)=2\ \text{ and }\
 \bigl(\forall\Zr\in\Phi,\ (\Zv,\Zr)\in\{-2,-1,0,1,2\}\bigr)
 \quad\Longrightarrow\quad \Zv\in\Phi.
 \]
 We include $\pm2$ to allow a line already in $S$; by \eqref{eq:endpoints} these values
 occur only on that same line. \repair\ The paper takes the lines to span $\re^n$ (the first
 paragraph of Section 2) and asks for maximality in $\re^n$; since we do not require $\Phi$ to
 span $V$, we ask for maximality within $\operatorname{span}(\Phi)$. In the whole of $V$ the
 condition would be too strong: $\Ze_1+\Ze_2$ is at $60^\circ$ or $90^\circ$ to every line of
 $\typ{A}_\ell$ in $\re^{\ell+1}$, and Theorem 3.6 would fail (errata, Implicit hypothesis 7). Thus this is
 Definition 2.2 in the signed vector encoding.
\end{cgss}

\begin{cgss}{Lemma}{2.3}\label{lem:2.3}
 \lean{LineSystem.IsMaximal.isStarClosed}\uses{def:2.1,def:2.2,facts:inner}
 \assume{\begin{pfenum}
  \item $S$ is a line system in a finite-dimensional real inner product space $V$.
  \item $S$ is maximal.
 \end{pfenum}}
 \prove{$S$ is star-closed.}
\end{cgss}
\begin{proof}
 \pfsketch\ We will orient the two lines so that their inner product is $-1$ and test their
 sum against every root. This is the angle argument in the paper, including coincident lines.
 \step{23sum}{\assume{$\Zr,\Zs\in\Phi$ and $(\Zr,\Zs)=-1$.}
  \prove{$\Zr+\Zs\in\Phi$.}}
 \begin{proof}
  \step{23norm}{$(\Zr+\Zs,\Zr+\Zs)=2$.}
  \begin{proof}
   \pf\ By bilinearity and \asm{1}, the squared norm is $2-1-1+2=2$.~\qed
  \end{proof}
  \step{23test}{\assume{$\Zt\in\Phi$.}
   \prove{$(\Zr+\Zs,\Zt)\in\{-2,-1,0,1,2\}$.}}
  \begin{proof}
   \step{23same}{\case{$\Zt\in\{\Zr,-\Zr,\Zs,-\Zs\}$.}}
   \begin{proof}
    \pf\ The four values are $1,-1,1,-1$, respectively, by bilinearity and \asm{1}.~\qed
   \end{proof}
   \step{23other}{\case{$\Zt\notin\{\Zr,-\Zr,\Zs,-\Zs\}$.}}
   \begin{proof}
    \step{23small}{$(\Zr,\Zt),(\Zs,\Zt)\in\{-1,0,1\}$.}
    \begin{proof}
     \pf\ The line-system condition and \eqref{eq:endpoints} exclude $\pm2$.~\qed
    \end{proof}
    \qedstep
    \begin{proof}
     \pf\ The sum of two integers between $-1$ and 1 is an integer between $-2$ and 2;
     apply bilinearity to \stepref{23small}.~\qed
    \end{proof}
   \end{proof}
   \qedstep
   \begin{proof}\pf\ The cases \stepref{23same} and \stepref{23other} are exhaustive.~\qed\end{proof}
  \end{proof}
  \step{23span}{$\Zr+\Zs\in\operatorname{span}(\Phi)$.}
  \begin{proof}
   \pf\ Both $\Zr$ and $\Zs$ lie in $\Phi$ by \asm{1}, and a span is closed under
   addition.~\qed
  \end{proof}
  \qedstep
  \begin{proof}\pf\ Apply Definition 2.2 to \stepref{23span}, \stepref{23norm} and
   \stepref{23test}.~\qed\end{proof}
 \end{proof}
 \qedstep
 \begin{proof}\pf\ Step \stepref{23sum} is Definition 2.1.~\qed\end{proof}
\end{proof}

\begin{cgss}{Hypothesis}{2.4}\label{hyp:2.4}
 \lean{LineSystem.Hyp24}\uses{def:line-system,def:indecomposable,def:2.1}
 \assume{\begin{pfenum}
  \item $V$ is a finite-dimensional real inner product space and $S$ is a line system in $V$.
  \item $S$ is indecomposable and star-closed.
  \item $\Za,\Zb\in\Phi$ and $(\Za,\Zb)=-1$.
 \end{pfenum}}
 \pflet{$\Zc\triangleq-\Za-\Zb$, and}
 \begin{align*}
 \Ga&\triangleq\{\Zx\in\Phi:(\Zx,\Za)=0,\ (\Zx,\Zb)=1,\ (\Zx,\Zc)=-1\},\\
 \Gb&\triangleq\{\Zx\in\Phi:(\Zx,\Zb)=0,\ (\Zx,\Zc)=1,\ (\Zx,\Za)=-1\},\\
 \Gc&\triangleq\{\Zx\in\Phi:(\Zx,\Zc)=0,\ (\Zx,\Za)=1,\ (\Zx,\Zb)=-1\},\\
 \Delta&\triangleq\{\Zt\in\Phi:(\Zt,\Za)=(\Zt,\Zb)=(\Zt,\Zc)=0\},\\
 \Sigma&\triangleq\{\Za,\Zb,\Zc\}\cup\Ga\cup\Gb\cup\Gc\cup\Delta.
 \end{align*}
 We will write $X=\opk{Orth}(\Ga)$ for the \emph{orthogonality graph} on $\Ga$: distinct vertices
 $x,y$ are adjacent exactly when $(\Zx,\Zy)=0$. We will write a vertex as $x$ and
 its representing vector as $\Zx$.
 \label{def:gamma-graph}\lean{LineSystem.Hyp24.orthGraph}\uses{hyp:2.4}
 The represented graph of a vector family uses inner product 1; its
 \emph{obtuse graph}, denoted $\opk{Obt}$, uses inner product $-1$.
 We use the orthogonality convention throughout Section 2.
 The \emph{empty graph} $\comp{K_m}$ has $m$ vertices and no edges, with $m=0$ allowed.
\end{cgss}

\noindent\repair\ We require the pair $\Za,\Zb$ explicitly (Implicit hypothesis 4 of
\path{cgss/errata.md}). A one-line system has no such pair;
we will treat it separately in Lemma 2.12 and Theorem 2.13.
\repair\ Also, $\Delta=-\Delta$ here (Error 6), so
$\Sigma$ need not have just one vector on each line. Lemma 2.5 asserts a signed cover.

\begin{namedresult}{Lemma}{(star identities)}\label{lem:star-identities}
 \lean{LineSystem.Hyp24.c_mem}\lean{LineSystem.Hyp24.inner_b_c}
 \lean{LineSystem.Hyp24.inner_c_a}\lean{LineSystem.Hyp24.inner_c_b}
 \uses{hyp:2.4,def:2.1,facts:inner}
 \assume{Hypothesis 2.4 holds in the finite-dimensional real inner product space $V$.}
 \prove{$\Zc\in\Phi$ and
  $(\Za,\Zb)=(\Zb,\Zc)=(\Zc,\Za)=(\Zc,\Zb)=-1$.}
\end{namedresult}
\begin{proof}
 \pfsketch\ \new\ We will state the star identities separately so that later results
 can cite them without citing a step inside the proof of Lemma 2.5.
 \step{25star}{$\Zc\in\Phi$ and all three inner products between $\Za,\Zb,\Zc$ are $-1$.}
 \begin{proof}
  \step{25mem}{$\Za+\Zb\in\Phi$, so $\Zc\in\Phi$.}
  \begin{proof}\pf\ Use star-closure, $(\Za,\Zb)=-1$, and closure under negation.~\qed\end{proof}
  \qedstep
  \begin{proof}
   \pf\ By bilinearity, $(\Za,\Zc)=-2+1=-1$ and $(\Zb,\Zc)=1-2=-1$;
   the remaining products follow from symmetry and the hypothesis $(\Za,\Zb)=-1$.~\qed
  \end{proof}
 \end{proof}
 \qedstep
 \begin{proof}\pf\ Step \stepref{25star} and symmetry give the stated identities.~\qed\end{proof}
\end{proof}

\begin{cgss}{Lemma}{2.5}\label{lem:2.5}
 \lean{LineSystem.Hyp24.Φ_eq_sigmaSet_union_neg}
 \uses{hyp:2.4,lem:star-identities,facts:inner}
 \assume{Hypothesis 2.4 holds in the finite-dimensional real inner product space $V$.}
 \prove{$\Phi=\Sigma\cup(-\Sigma)$; equivalently, $S$ consists exactly of the lines of $\Sigma$.}
\end{cgss}
\begin{proof}
 \pfsketch\ We will exclude equal nonzero inner products with $\Za,\Zb$, then choose the
 sign in each of the remaining seven cases.
 \step{25equal}{\assume{$\Zu\in\Phi$, $\epsilon\in\{-1,1\}$, and
   $(\Zu,\Za)=(\Zu,\Zb)=\epsilon$.}
  \prove{$\Zu=-\epsilon\Zc$.}}
 \begin{proof}
  \step{25zero}{$(\Za+\Zb-\epsilon\Zu,\Za+\Zb-\epsilon\Zu)=0$.}
  \begin{proof}
   \pf\ Bilinearity gives $2+2\epsilon^2-2\epsilon(2\epsilon)=0$, since
   $\epsilon^2=1$ and $(\Za+\Zb,\Za+\Zb)=2$.~\qed
  \end{proof}
  \qedstep
  \begin{proof}
   \pf\ Zero squared norm gives $\Za+\Zb=\epsilon\Zu$; multiply by $\epsilon$ and use
   $\Zc=-\Za-\Zb$ and $\epsilon^2=1$.~\qed
  \end{proof}
 \end{proof}
 \step{25switch}{\assume{$\Zv\in\Phi$.}
  \prove{$\Zv\in\Sigma$ or $-\Zv\in\Sigma$.}}
 \begin{proof}
  \step{25onstar}{\case{$\Zv\in\{\pm\Za,\pm\Zb,\pm\Zc\}$.}}
  \begin{proof}\pf\ Choose the displayed sign of a member of $\{\Za,\Zb,\Zc\}\subseteq\Sigma$.~\qed\end{proof}
  \step{25offstar}{\case{$\Zv\notin\{\pm\Za,\pm\Zb,\pm\Zc\}$.}}
  \begin{proof}
   \step{25poss}{$((\Zv,\Za),(\Zv,\Zb))$ is one of
    $(0,0),(0,1),(0,-1),(1,0),(-1,0),(1,-1),(-1,1)$.}
   \begin{proof}
    \pf\ By \eqref{eq:endpoints}, both entries are in $\{-1,0,1\}$.
    Step \stepref{25equal} excludes the two pairs $(1,1)$ and $(-1,-1)$.~\qed
   \end{proof}
   \step{25table}{We choose representatives by the following table.}
   \begin{proof}
    \pf\ In each row $(\Zv,\Zc)=-(\Zv,\Za)-(\Zv,\Zb)$ by bilinearity;
    substitution gives the defining three inner products of the indicated set.
    \[
    \begin{array}{c|rrrrrrr}
    ((\Zv,\Za),(\Zv,\Zb))&(0,0)&(0,1)&(0,-1)&(1,0)&(-1,0)&(1,-1)&(-1,1)\\\hline
    \text{representative}&\Zv&\Zv&-\Zv&-\Zv&\Zv&\Zv&-\Zv\\
    \text{set}&\Delta&\Ga&\Ga&\Gb&\Gb&\Gc&\Gc
    \end{array}
    \]
    Each representative belongs to $\Phi$ by closure under negation.~\qed
   \end{proof}
   \qedstep
   \begin{proof}\pf\ Apply \stepref{25table} to the exhaustive list \stepref{25poss}.~\qed\end{proof}
  \end{proof}
  \qedstep
  \begin{proof}\pf\ The cases \stepref{25onstar} and \stepref{25offstar} are exhaustive.~\qed\end{proof}
 \end{proof}
 \qedstep
 \begin{proof}
  \pf\ Step \stepref{25switch} gives $\Phi\subseteq\Sigma\cup(-\Sigma)$.
  Conversely, the star identities and the definitions give $\Sigma\subseteq\Phi$,
  and $\Phi=-\Phi$ gives $-\Sigma\subseteq\Phi$.~\qed
 \end{proof}
\end{proof}

\begin{cgss}{Lemma}{2.6}\label{lem:2.6}
 \lean{LineSystem.Hyp24.ΓaEquivΓb}
 \lean{LineSystem.Hyp24.ΓbEquivΓc}\lean{LineSystem.Hyp24.ΓcEquivΓa}\lean{LineSystem.Hyp24.inner_eq_zero_or_one}
 \lean{LineSystem.Hyp24.add_c_mem_Γb}\lean{LineSystem.Hyp24.sub_b_mem_Γc}
 \lean{LineSystem.Hyp24.inner_ΓaEquivΓb}\lean{LineSystem.Hyp24.inner_ΓbEquivΓc}
 \lean{LineSystem.Hyp24.inner_ΓcEquivΓa}
 \lean{LineSystem.Hyp24.inner_eq_zero_or_one_of_mem_Γb}
 \lean{LineSystem.Hyp24.inner_eq_zero_or_one_of_mem_Γc}
 \uses{hyp:2.4,lem:star-identities,facts:inner}
 \assume{Hypothesis 2.4 holds in the finite-dimensional real inner product space $V$.}
 \prove{\begin{pfenum}
  \item The maps $\varphi_c(\Zx)=\Zx+\Zc$, $\varphi_a(\Zy)=\Zy+\Za$, and
   $\varphi_b(\Zz)=\Zz+\Zb$ give inner-product-preserving bijections
   $\Ga\to\Gb\to\Gc\to\Ga$.
  \item $\Gb=\Zc+\Ga$ and $\Gc=-\Zb+\Ga$.
  \item For each $D\in\{\Ga,\Gb,\Gc\}$ and distinct $\Zx,\Zy\in D$,
   $(\Zx,\Zy)\in\{0,1\}$. Diagonal inner products are 2.
 \end{pfenum}}
\end{cgss}
\begin{proof}
 \pfsketch\ We will prove the translation identity for all three cyclic triples at once.
 Their composite is the identity, so each map is bijective. A negative inner product would
 give opposite vectors in two incompatible parts of the partition.
 \step{26cycle}{\assume{$(\Zr,\Zs,\Zt)$ is one of
  $(\Za,\Zb,\Zc),(\Zb,\Zc,\Za),(\Zc,\Za,\Zb)$, and $\Zx\in\Gamma_r$.}
  \prove{$\Zx+\Zt\in\Gamma_s$.}}
 \begin{proof}
  \step{26root}{$\Zx+\Zt\in\Phi$.}
  \begin{proof}
   \pf\ In each of the three listed triples, $(\Zx,\Zt)=-1$ by the definition of
   $\Gamma_r$; apply star-closure and the star identities.~\qed
  \end{proof}
  \step{26values}{$(\Zx+\Zt,\Zr)=-1$, $(\Zx+\Zt,\Zs)=0$, and
   $(\Zx+\Zt,\Zt)=1$.}
  \begin{proof}
   \pf\ By bilinearity, the star identities and the defining products, these values are respectively
   $0-1$, $1-1$, and $-1+2$.~\qed
  \end{proof}
  \qedstep
  \begin{proof}\pf\ Steps \stepref{26root} and \stepref{26values} are the definition of $\Gamma_s$.~\qed\end{proof}
 \end{proof}
 \step{26bijections}{The three translations are bijections and preserve inner products.}
 \begin{proof}
  \step{26inverse}{Each translation has the next two translations as a two-sided inverse.}
  \begin{proof}
   \pf\ By \stepref{26cycle} all composites have the stated domains, and every cyclic
   composite sends $\Zx$ to $\Zx+\Za+\Zb+\Zc=\Zx$.~\qed
  \end{proof}
  \step{26gram}{\assume{$(\Zr,\Zs,\Zt)$ is a listed cyclic triple and
    $\Zx,\Zy\in\Gamma_r$.}
   \prove{$(\Zx+\Zt,\Zy+\Zt)=(\Zx,\Zy)$.}}
  \begin{proof}
   \pf\ By bilinearity the left side is $(\Zx,\Zy)-1-1+2=(\Zx,\Zy)$,
   using the definition of $\Gamma_r$.~\qed
  \end{proof}
  \qedstep
  \begin{proof}\pf\ Combine the two-sided inverses in \stepref{26inverse} with \stepref{26gram}.~\qed\end{proof}
 \end{proof}
 \step{26positive}{\assume{$(\Zr,\Zs,\Zt)$ is a listed cyclic triple and
    $\Zx,\Zy\in\Gamma_r$ are distinct.}
  \prove{$(\Zx,\Zy)\in\{0,1\}$.}}
 \begin{proof}
  \step{26endpoints}{$(\Zx,\Zy)\ne2,-2$.}
  \begin{proof}
   \pf\ The value 2 contradicts distinctness by \eqref{eq:endpoints}.
   The value $-2$ gives $\Zy=-\Zx$, contradicting
   $(\Zx,\Zs)=(\Zy,\Zs)=1$.~\qed
  \end{proof}
  \step{26minus}{$(\Zx,\Zy)\ne-1$.}
  \begin{proof}
   \step{26opposite}{\assume{$(\Zx,\Zy)=-1$.}
    \prove{$\Zx=-(\Zy+\Zt)$.}}
   \begin{proof}
    \pf\ Step \stepref{26cycle} puts $\Zy+\Zt$ in $\Phi$, and bilinearity gives
    $(\Zx,\Zy+\Zt)=-1-1=-2$; apply \eqref{eq:endpoints}.~\qed
   \end{proof}
   \qedstep
   \begin{proof}
    \pf\ The equality in \stepref{26opposite} gives
    $0=(\Zx,\Zr)=-(\Zy+\Zt,\Zr)=1$, by the definition of $\Gamma_r$
    and $(\Zt,\Zr)=-1$ from the star identities, a contradiction.~\qed
   \end{proof}
  \end{proof}
  \qedstep
  \begin{proof}\pf\ The line-system condition, \stepref{26endpoints}, and \stepref{26minus} leave 0 and 1.~\qed\end{proof}
 \end{proof}
 \qedstep
 \begin{proof}
  \pf\ Step \stepref{26bijections} proves part 1. Its first map gives
  $\Gb=\Zc+\Ga$, and its first two maps give
  $\Gc=(\Za+\Zc)+\Ga=-\Zb+\Ga$. Step \stepref{26positive} and the norm condition prove part 3.~\qed
 \end{proof}
\end{proof}

\begin{cgss}{Lemma}{2.7}\label{lem:2.7}
 \lean{LineSystem.Hyp24.Δ_eq_sub}\lean{LineSystem.Hyp24.Φ_eq_of_Γa_eq}
 \lean{LineSystem.Hyp24.Φ_eq_recovery_union_neg}\lean{LineSystem.Hyp24.sigmaSet_eq_recoverySet}
 \uses{hyp:2.4,lem:2.5,lem:2.6,def:indecomposable,facts:inner}
 \assume{Hypothesis 2.4 holds in the finite-dimensional real inner product space $V$.}
 \prove{\[
  \Delta=\{\Zx-\Zy:\Zx,\Zy\in\Ga,\ (\Zx,\Zy)=1\}.
 \]
 Put
 \[
 \cR=\{\Za,\Zb,\Zc\}\cup\Ga\cup(\Zc+\Ga)\cup(-\Zb+\Ga)
       \cup\{\Zx-\Zy:\Zx,\Zy\in\Ga,\ (\Zx,\Zy)=1\}.
 \]
 Then $\Sigma=\cR$ and $\Phi=\cR\cup(-\cR)$.
 In particular, any two instances of Hypothesis 2.4 in $V$ with the same
 $\Za,\Zb$ and $\Ga$ have the same signed root set $\Phi$.}
\end{cgss}
\begin{proof}
 \pfsketch\ We will recover $\Gb,\Gc$ by translation and $\Delta$ by differences.
 A root of $\Delta$ orthogonal to all of $\Ga$ would belong to an orthogonal part of $\Phi$;
 we will give that partition explicitly.
 \step{27forward}{\assume{$\Zx,\Zy\in\Ga$ and $(\Zx,\Zy)=1$.}
  \prove{$\Zx-\Zy\in\Delta$.}}
 \begin{proof}
  \step{27root}{$\Zx-\Zy\in\Phi$.}
  \begin{proof}\pf\ Apply star-closure to $\Zx,-\Zy\in\Phi$, whose inner product is $-1$.~\qed\end{proof}
  \qedstep
  \begin{proof}
   \pf\ The products of $\Zx-\Zy$ with $\Za,\Zb,\Zc$ are
   $0-0,1-1,-1-(-1)$, respectively, by bilinearity. With \stepref{27root}, this defines $\Delta$.~\qed
  \end{proof}
 \end{proof}
 \step{27nonzero}{\assume{$\Zt\in\Delta$ and $(\Zt,\Zx)\ne0$ for some $\Zx\in\Ga$.}
  \prove{$\Zt=\Zu-\Zv$ for some $\Zu,\Zv\in\Ga$ with $(\Zu,\Zv)=1$.}}
 \begin{proof}
  \pflet{Choose $\Zx\in\Ga$ with $(\Zt,\Zx)\ne0$.}
  \step{27sign}{$(\Zt,\Zx)\in\{-1,1\}$.}
  \begin{proof}
   \pf\ If the product were $\pm2$, \eqref{eq:endpoints} would give
   $\Zt=\pm\Zx$, contradicting $(\Zt,\Zb)=0$ and $(\Zx,\Zb)=1$.
   The line-system condition and the nonzero assumption leave $\pm1$.~\qed
  \end{proof}
  \step{27negative}{\case{$(\Zt,\Zx)=-1$.}}
  \begin{proof}
   \step{27add}{$\Zx+\Zt\in\Ga$.}
   \begin{proof}
    \pf\ Star-closure gives membership in $\Phi$; adding $\Zt$, which is orthogonal to
    $\Za,\Zb,\Zc$, preserves the three products defining $\Ga$.~\qed
   \end{proof}
   \qedstep
   \begin{proof}
    \pf\ Choose $\Zu=\Zx+\Zt$, $\Zv=\Zx$. Their difference is $\Zt$, and
    $(\Zu,\Zv)=2-1=1$ by bilinearity and \asm{2}.~\qed
   \end{proof}
  \end{proof}
  \step{27positive}{\case{$(\Zt,\Zx)=1$.}}
  \begin{proof}
   \step{27sub}{$\Zx-\Zt\in\Ga$.}
   \begin{proof}
    \pf\ Star-closure applies to $\Zx,-\Zt$; subtracting $\Zt$ preserves the
    three products defining $\Ga$.~\qed
   \end{proof}
   \qedstep
   \begin{proof}
    \pf\ Choose $\Zu=\Zx$, $\Zv=\Zx-\Zt$. Their difference is $\Zt$, and
    $(\Zu,\Zv)=2-1=1$ by bilinearity and \asm{2}.~\qed
   \end{proof}
  \end{proof}
  \qedstep
  \begin{proof}\pf\ Step \stepref{27sign} makes \stepref{27negative} and \stepref{27positive} exhaustive.~\qed\end{proof}
 \end{proof}
 \step{27partition}{No root of $\Delta$ is orthogonal to all of $\Ga$.}
 \begin{proof}
  \pflet{$P\triangleq\{\Zt\in\Delta:\forall\Zx\in\Ga,\ (\Zt,\Zx)=0\}$ and
   $Q\triangleq\Phi\diff P$.}
  \step{27signstable}{$P=-P$ and $\Za\in Q$.}
  \begin{proof}
   \pf\ The defining products of $P$ are all zero, and thus invariant under negation.
   Also $\Za\in\Phi$ and $(\Za,\Za)=2$, so $\Za\notin\Delta$.~\qed
  \end{proof}
  \step{27sigma}{\assume{$\Zt\in P$ and $\Zu\in\Sigma\diff P$.}
   \prove{$(\Zt,\Zu)=0$.}}
  \begin{proof}
   \step{27core}{\case{$\Zu\in\{\Za,\Zb,\Zc\}\cup\Ga\cup\Gb\cup\Gc$.}}
   \begin{proof}
    \pf\ The definition of $P$ gives orthogonality to $\Za,\Zb,\Zc,\Ga$.
    Lemma 2.6 gives $\Gb=\Zc+\Ga$ and $\Gc=-\Zb+\Ga$; bilinearity gives zero
    on each of these translates.~\qed
   \end{proof}
   \step{27delta}{\case{$\Zu\in\Delta\diff P$.}}
   \begin{proof}
    \step{27witness}{$\Zu=\Zx-\Zy$ for some $\Zx,\Zy\in\Ga$.}
    \begin{proof}
     \pf\ Membership in $\Delta\diff P$ gives a nonzero product with a member of $\Ga$;
     apply \stepref{27nonzero}.~\qed
    \end{proof}
    \qedstep
    \begin{proof}\pf\ By the definition of $P$, $(\Zt,\Zx-\Zy)=0-0=0$.~\qed\end{proof}
   \end{proof}
   \qedstep
   \begin{proof}\pf\ The definition of $\Sigma$ makes \stepref{27core} and \stepref{27delta} exhaustive.~\qed\end{proof}
  \end{proof}
  \step{27orthogonal}{$(\Zt,\Zu)=0$ for all $\Zt\in P$ and $\Zu\in Q$.}
  \begin{proof}
   \pf\ Lemma 2.5 gives $\Zu=\epsilon\Zv$ with $\epsilon\in\{-1,1\}$ and
   $\Zv\in\Sigma$. By \stepref{27signstable}, $\Zv\notin P$.
   Step \stepref{27sigma} and bilinearity give $(\Zt,\Zu)=\epsilon\cdot0=0$.~\qed
  \end{proof}
  \qedstep
  \begin{proof}
   \pf\ We have $P\cup Q=\Phi$. Indecomposability and \stepref{27orthogonal} imply
   $P=\varnothing$ or $Q=\varnothing$; \stepref{27signstable} excludes the second alternative.~\qed
  \end{proof}
 \end{proof}
 \qedstep
 \begin{proof}
  \pf\ Steps \stepref{27forward}, \stepref{27nonzero}, and \stepref{27partition} give both
  inclusions in the formula for $\Delta$. Substitute this equality and Lemma 2.6
  into the definition of $\Sigma$ to obtain $\Sigma=\cR$; Lemma 2.5 gives
  $\Phi=\cR\cup(-\cR)$. Equal $\Za,\Zb,\Ga$ give equal $\Zc=-\Za-\Zb$
  and equal $\cR$, so they give equal root sets.~\qed
 \end{proof}
\end{proof}

\begin{cgss}{Lemma}{2.8}\label{lem:2.8}
  \lean{LineSystem.Hyp24.triangle}\uses{hyp:2.4,lem:2.6,lem:star-identities,facts:inner}
  \assume{\begin{pfenum}
      \item Hypothesis 2.4 holds in the finite-dimensional real inner product space $V$.
      \item $\Zx,\Zy\in\Ga$ and $(\Zx,\Zy)=0$.
    \end{pfenum}}
  \prove{\begin{pfenum}
      \item $\Zz\triangleq\Zb-\Zc-\Zx-\Zy$ lies in $\Ga$.
      \item Each $\Zw\in\Ga\diff\{\Zx,\Zy,\Zz\}$ is orthogonal to exactly one of $\Zx$, $\Zy$,
        $\Zz$.
    \end{pfenum}}
\end{cgss}

\begin{proof}
  \pfsketch\ The vectors $\Zx+\Zc$ and $\Zy-\Zb$ lie in $\Phi$ and have inner product $-1$, so
  star-closure puts their sum, which is $-\Zz$, in $\Phi$. Computing inner products places $\Zz$
  in $\Ga$. For part 2, $(\Zw,\Zx)+(\Zw,\Zy)+(\Zw,\Zz)=2$, and each term is 0 or 1.

  \step{xc}{$\Zx+\Zc\in\Gb$ and $\Zy-\Zb\in\Gc$.}
  \begin{proof}
    \pf\ By \asm[2]{0} and Lemma 2.6.~\qed
  \end{proof}

  \step{ip}{$(\Zx+\Zc,\Zy-\Zb)=-1$.}
  \begin{proof}
    \pf\ We use $(\Zb,\Zc)=-1$ from the star identities.
    \begin{align*}
      (\Zx+\Zc,\Zy-\Zb)&=(\Zx,\Zy)-(\Zx,\Zb)+(\Zc,\Zy)-(\Zc,\Zb)
        &&\text{[bilinearity]}\\
      &=0-1-1+1
        &&\text{[\asm[2]{0}, the definition of $\Ga$, and $(\Zb,\Zc)=-1$]}\\
      &=-1.
    \end{align*}
    \qed
  \end{proof}

  \step{zphi}{$\Zz\in\Phi$.}
  \begin{proof}
    \pf\ By \stepref{xc}, \stepref{ip} and star-closure, $(\Zx+\Zc)+(\Zy-\Zb)=-\Zz$ lies in
    $\Phi$; and $\Phi=-\Phi$.~\qed
  \end{proof}

  \step{zga}{$\Zz\in\Ga$.}
  \begin{proof}
    \pf\ We have $(\Zz,\Za)=0$, $(\Zz,\Zb)=1$ and $(\Zz,\Zc)=-1$, by bilinearity, \asm[2]{0} and
    the definition of $\Ga$; together with \stepref{zphi}, this is the definition of
    $\Ga$.~\qed
  \end{proof}

  \step{w}{\assume{$\Zw\in\Ga\diff\{\Zx,\Zy,\Zz\}$.}
    \prove{$\Zw$ is orthogonal to exactly one of $\Zx$, $\Zy$, $\Zz$.}}
  \begin{proof}
    \step{sum}{$(\Zw,\Zx)+(\Zw,\Zy)+(\Zw,\Zz)=2$.}
    \begin{proof}
      \pf\
      \begin{align*}
        (\Zw,\Zx)+(\Zw,\Zy)+(\Zw,\Zz)&=(\Zw,\Zb-\Zc)&&\text{[the definition of $\Zz$]}\\
        &=1+1&&\text{[\asm{1} and the definition of $\Ga$]}
      \end{align*}
      \qed
    \end{proof}

    \step{zo}{Each of $(\Zw,\Zx)$, $(\Zw,\Zy)$, $(\Zw,\Zz)$ is 0 or 1.}
    \begin{proof}
      \pf\ By \asm{1}, \asm[2]{0}, \stepref{zga} and Lemma 2.6.~\qed
    \end{proof}

    \qedstep
    \begin{proof}
      \pf\ By \stepref{sum} and \stepref{zo}, three numbers in $\{0,1\}$ have sum 2, so exactly
      one of them is 0.~\qed
    \end{proof}
  \end{proof}

  \qedstep
  \begin{proof}
    \pf\ By \stepref{zga} and \stepref{w}.~\qed
  \end{proof}
\end{proof}

\begin{namedresult}{Definition}{(the triangle property)}\label{def:triangle-property}
 \lean{SimpleGraph.HasTriangleProperty}\uses{}
 \assume{$X$ is a simple graph on a set $U$.}
 We say that $X$ has the \emph{triangle property} if every edge $xy$ has exactly one
 common neighbour $z$, and every vertex outside $\{x,y,z\}$ is adjacent to exactly one
 of $x,y,z$. For adjacent $x,y$ we will write $f(x,y)$ for this unique common neighbour;
 the graph will be clear from the argument.
 \new\ We name this predicate once so that we can use the graph argument without a vector
 representation.
\end{namedresult}

\begin{namedresult}{Corollary}{of Lemma 2.8}\label{cor:triangle-property}
 \lean{LineSystem.Hyp24.orthGraph_hasTriangleProperty}\uses{hyp:2.4,lem:2.8,def:triangle-property,facts:inner}
 \assume{Hypothesis 2.4 holds in the finite-dimensional real inner product space $V$, and $X$ is its orthogonality graph on $\Ga$.}
 \prove{$X$ has the triangle property.}
\end{namedresult}
\begin{proof}
 \pfsketch\ \new\ We will record the pairwise orthogonality calculation needed to pass
 from Lemma 2.8 to graph terminology; it does not change that lemma.
 \step{28triangle}{\assume{$x,y$ are adjacent vertices of $X$.}
  \prove{$\Zz\triangleq\Zb-\Zc-\Zx-\Zy$ represents a distinct common neighbour $z$.}}
 \begin{proof}
  \step{28products}{$(\Zx,\Zz)=(\Zy,\Zz)=0$.}
  \begin{proof}
   \pf\ Bilinearity gives $(\Zx,\Zz)=1-(-1)-2-0=0$ and
   $(\Zy,\Zz)=1-(-1)-0-2=0$, using adjacency and the definition of $\Ga$.~\qed
  \end{proof}
  \qedstep
  \begin{proof}
   \pf\ Lemma 2.8 puts $\Zz$ in $\Ga$. Step \stepref{28products} and squared norm 2
   give $z\ne x,y$ and $z\sim x,y$.~\qed
  \end{proof}
 \end{proof}
 \step{28unique}{\assume{$w$ is a common neighbour of adjacent $x,y$.}
  \prove{$\Zw=\Zb-\Zc-\Zx-\Zy$.}}
 \begin{proof}
  \pf\ If $w$ differed from the vertex in \stepref{28triangle}, it would lie outside
  its triangle and be adjacent to two of its vertices, contrary to Lemma 2.8.~\qed
 \end{proof}
 \qedstep
 \begin{proof}
  \pf\ Steps \stepref{28triangle} and \stepref{28unique} give the unique common neighbour;
  Lemma 2.8 gives the assertion about every outside vertex.~\qed
 \end{proof}
\end{proof}

\begin{cgss}{Lemma}{2.9}\label{lem:2.9}
 \lean{LineSystem.Hyp24.orthGraph_iso_friendship_of_forall_adj}\uses{hyp:2.4,def:triangle-property,cor:triangle-property}
 \assume{\begin{pfenum}
  \item Hypothesis 2.4 holds in the finite-dimensional real inner product space $V$.
  \item $X$ is its orthogonality graph and $x\in V(X)$ is adjacent to every other vertex.
 \end{pfenum}}
 \prove{There is an integer $m\geq0$ and a partition of $V(X)\diff\{x\}$ into $m$ pairs,
  each of which makes a triangle with $x$. These are all the edges of $X$; thus $X$ is the friendship graph with $m$ triangles.}
\end{cgss}
\begin{proof}
 \pfsketch\ The unique triangle on each edge through $x$ pairs the other vertices.
 We will use the outside-vertex condition to exclude edges between different pairs.
 \step{29mate}{For $y\ne x$, let $f(x,y)$ be the unique common neighbour of $x,y$.
  Then $f(x,y)\ne y,x$ and $f(x,f(x,y))=y$.}
 \begin{proof}
  \step{29exists}{$f(x,y)$ exists and differs from $x,y$.}
  \begin{proof}\pf\ Apply Corollary~\ref{cor:triangle-property} to the edge $xy$.~\qed\end{proof}
  \qedstep
  \begin{proof}
   \pf\ The vertex $y$ is a common neighbour of $x,f(x,y)$, so uniqueness in
   Corollary~\ref{cor:triangle-property} gives $f(x,f(x,y))=y$.~\qed
  \end{proof}
 \end{proof}
 \step{29pairs}{The sets $\{y,f(x,y)\}$ partition $V(X)\diff\{x\}$ into pairs.}
 \begin{proof}
  \pf\ Step \stepref{29mate} gives a fixed-point-free involution: its two-element
  orbits cover the domain, and intersecting orbits coincide by applying the involution.~\qed
 \end{proof}
 \step{29edges}{There is no edge between two different pairs.}
 \begin{proof}
  \pf\ A vertex in a different pair lies outside $\{x,y,f(x,y)\}$ and is adjacent to $x$.
  Corollary~\ref{cor:triangle-property} excludes adjacency to either $y$ or $f(x,y)$.~\qed
 \end{proof}
 \qedstep
 \begin{proof}
  \pf\ The finite partition in \stepref{29pairs} has some size $m\geq0$.
  Step \stepref{29mate} supplies the triangle edges, and \stepref{29edges} excludes all others.
  When $V(X)=\{x\}$ the partition is empty and $m=0$.~\qed
 \end{proof}
\end{proof}

\begin{namedresult}{Lemma}{(triangle-property normal forms)}\label{lem:triangle-normal-form}
 \lean{SimpleGraph.HasTriangleProperty.normal_form_of_not_forall_adj}
 \uses{def:triangle-property}
 \assume{\begin{pfenum}
  \item $X$ is a finite simple graph with the triangle property.
  \item $X$ has an edge.
  \item No vertex of $X$ is adjacent to every other vertex.
 \end{pfenum}}
 \prove{For every pair of distinct nonadjacent vertices $u,v$, put
 \[
 \begin{array}{ll}
 K=\Gamma_1(u)\cap\Gamma_1(v),&L=\Gamma_1(u)\diff\Gamma_1(v),\\
 M=\Gamma_1(v)\diff\Gamma_1(u),&
 N=V(X)\diff\bigl(\{u,v\}\cup K\cup L\cup M\bigr),
 \end{array}
 \]
 and put $k=|K|$. Then the following hold.
 \begin{pfenum}
  \item $k\in\{2,3,5\}$, $X$ is regular of degree $2k$, $|V(X)|=6k-3$, and
   \repair\ $X$ is strongly regular with parameters $(6k-3,2k,1,k)$
   (Error 12 of \path{cgss/errata.md}).
  \item The sets $\{u\},\{v\},K,L,M,N$ partition $V(X)$, with $|L|=|M|=k$ and
   $|N|=3k-5$. The labels $x'=f(u,x)$ and $x''=f(v,x)$ give bijections
   $K\to L$ and $K\to M$.
  \item $K,L,M$ are independent. For $x,y\in K$, $x\sim y'$ and $x\sim y''$
   hold exactly when $x=y$, and $x'\sim y''$ holds exactly when $x\ne y$.
  \item For $w\in N$, put $A(w)=K\diff\Gamma_1(w)$ and
   $d_K(w)=|K\cap\Gamma_1(w)|$. Then $|A(w)|\in\{0,2\}$,
   $d_K(w)\in\{k-2,k\}$, and $|N\cap\Gamma_1(w)|=d_K(w)$.
   For $x\in K$, $w\sim x'$ and $w\sim x''$ hold exactly when $x\in A(w)$.
  \item The graph on $N$ has $k(k-2)$ edges, each with its third vertex in $K$.
   Moreover,
   \[
    \sum_{w\in N}d_K(w)=k(2k-4),\qquad
    2|\{w\in N:d_K(w)=k\}|=(k-2)(5-k).
   \]
  \item If $k=2$, $N$ consists of one vertex $w$, with $A(w)=K$, and has no edges.
  \item If $k=3$, the graph on $N$ is a claw with centre $w_0$, where $A(w_0)=\varnothing$.
   Its leaves are $w_x$ for $x\in K$, with $A(w_x)=K\diff\{x\}$.
  \item If $k=5$, the map $w\mapsto A(w)$ is a bijection from $N$ to the two-element
   subsets of $K$, and adjacency in $N$ is disjointness of these labels.
 \end{pfenum}
 These rules and the definitions of $K,L,M,N$ specify every adjacency.}
\end{namedresult}
\begin{proof}
 \pfsketch\ \new\ We separate the combinatorial argument from its application to
 line systems. We will keep the paper's partition and prove regularity for arbitrary
 pairs. The invariant $d_K$ gives both the bound on $k$ and the three normal forms.
 All counting identities involving subtraction are read in $\ints$.

 \step{210triangles}{Every edge $xy$ has a unique common neighbour $f(x,y)$.
  Every vertex has positive degree and has a distinct non-neighbour.}
 \begin{proof}
  \step{210third}{The triangle on $xy$ is $\{x,y,f(x,y)\}$ and every outside vertex
   sees exactly one vertex of it.}
  \begin{proof}\pf\ Apply the triangle property in \asm[1]{0}.~\qed\end{proof}
  \step{210positive}{Every vertex has positive degree.}
  \begin{proof}
   \pf\ A triangle exists since $X$ has an edge. Its vertices have degree at least 2,
   and every other vertex has a neighbour in it by \stepref{210third}.~\qed
  \end{proof}
  \step{210nouniversal}{Every vertex has a distinct non-neighbour.}
  \begin{proof}
   \pf\ This is precisely \asm[3]{0}, the absence of a universal vertex.~\qed
  \end{proof}
  \qedstep
  \begin{proof}\pf\ Combine \stepref{210third}, \stepref{210positive}, and \stepref{210nouniversal}.~\qed\end{proof}
 \end{proof}

 \step{210local}{For every pair of distinct nonadjacent vertices $u,v$, define, excluding $u,v$,
  \[
  \begin{array}{ll}
  K=\Gamma_1(u)\cap\Gamma_1(v),&
  L=\Gamma_1(u)\diff\Gamma_1(v),\\
  M=\Gamma_1(v)\diff\Gamma_1(u),&
  N=V(X)\diff\bigl(\{u,v\}\cup K\cup L\cup M\bigr).
  \end{array}
  \]
  These four sets are disjoint. If $k=|K|$, then $|L|=|M|=k>0$ and
  $\deg(u)=\deg(v)=2k$. They admit labels
  $L=\{x':x\in K\}$ and $M=\{x'':x\in K\}$ with the following complete edge rules:
  $K,L,M$ are independent; $x\sim y'$ and $x\sim y''$ exactly when $x=y$;
  $x'\sim y''$ exactly when $x\ne y$.
  For $w\in N$, $w\sim x'$ and $w\sim x''$ hold exactly when $w\not\sim x$.}
 \begin{proof}
  \step{210pairing}{For $p\in\{u,v\}$ and the other vertex $q$, the triangles through $p$
   partition $\Gamma_1(p)$ into pairs; $q$ sees exactly one vertex in each pair.}
  \begin{proof}
   \step{210involution}{On $\Gamma_1(p)$, the map $y\mapsto f(p,y)$ is a
    fixed-point-free involution, and these pairs are all the edges within $\Gamma_1(p)$.}
   \begin{proof}
    \pf\ The proof of \stepref{210triangles} gives existence and uniqueness of the third
    vertex. Thus $f(p,f(p,y))=y$. If two neighbours of $p$ are adjacent, uniqueness of
    the common neighbour on an edge through $p$ makes them one of these pairs.~\qed
   \end{proof}
   \qedstep
   \begin{proof}
    \pf\ Since $p\not\sim q$, the triangle property forces $q$ to see exactly one of
    the two other vertices of each triangle through $p$.~\qed
   \end{proof}
  \end{proof}
  \step{210labels}{For $x\in K$, put $x'=f(u,x)$ and $x''=f(v,x)$.
   These define bijections $K\to L$ and $K\to M$, and give the stated rules inside
   $K\cup L$ and $K\cup M$.}
  \begin{proof}
   \step{210onto}{Each pair through $u$ consists of one point of $K$ and one of $L$;
    each pair through $v$ consists of one point of $K$ and one of $M$.}
   \begin{proof}\pf\ Apply \stepref{210pairing} with $(p,q)=(u,v)$ and $(v,u)$, respectively.~\qed\end{proof}
   \qedstep
   \begin{proof}
    \pf\ Step \stepref{210onto} gives the bijections. Since \stepref{210pairing} lists
    every edge in either neighbourhood, each part is independent and the only edges
    between its two parts are the labelled pairs.~\qed
   \end{proof}
  \end{proof}
  \step{210lm}{$x'\sim y''$ if and only if $x\ne y$.}
  \begin{proof}
   \step{210same}{For $x=y$, $x'\not\sim x''$.}
   \begin{proof}
    \pf\ The vertex $x''$ sees $x$ and does not see $u$. Apply the triangle property to
    $\{u,x,x'\}$ to exclude $x''\sim x'$.~\qed
   \end{proof}
   \step{210different}{For $x\ne y$, $x'\sim y''$.}
   \begin{proof}
    \pf\ The vertex $x'$ sees neither $v$ nor $y$, by its definition and
    \stepref{210labels}. Applied to $\{v,y,y''\}$, the triangle property forces
    $x'\sim y''$.~\qed
   \end{proof}
   \qedstep
   \begin{proof}\pf\ The cases \stepref{210same} and \stepref{210different} are exhaustive.~\qed\end{proof}
  \end{proof}
  \step{210signature}{For $w\in N$ and $x\in K$, the adjacencies to $(x,x',x'')$
   are either $(1,0,0)$ or $(0,1,1)$.}
  \begin{proof}
   \pf\ The vertex $w$ sees neither $u$ nor $v$. It must see one vertex in each
   of $\{u,x,x'\}$ and $\{v,x,x''\}$ by the triangle property.~\qed
  \end{proof}
  \qedstep
  \begin{proof}
   \pf\ The sets partition $V(X)\diff\{u,v\}$ by their two adjacency indicators.
   Step \stepref{210labels} gives the equal cardinalities and degrees; positivity follows
   from \stepref{210triangles}. The remaining assertions are
   \stepref{210labels}, \stepref{210lm}, and \stepref{210signature}.~\qed
  \end{proof}
 \end{proof}

 \step{210regular}{\repair\ $X$ is regular. We fill the paper's missing propagation
  from nonadjacent pairs to arbitrary pairs (Gap 3 of \path{cgss/errata.md}).}
 \begin{proof}
  \step{210nonadjdeg}{Distinct nonadjacent vertices have equal degrees.}
  \begin{proof}\pf\ This is the degree equality in \stepref{210local}, for an arbitrary pair.~\qed\end{proof}
  \step{210adjdeg}{\assume{$p\sim q$.}\prove{$\deg(p)=\deg(q)$.}}
  \begin{proof}
   \pflet{$r=f(p,q)$. Choose $s\ne p$ nonadjacent to $p$, using \stepref{210triangles}.}
   \step{210s}{The vertex $s$ is outside $\{p,q,r\}$ and sees exactly one of $q,r$.}
   \begin{proof}
    \pf\ Both $q,r$ see $p$, so $s\ne q,r$; apply the triangle property and $s\not\sim p$.~\qed
   \end{proof}
   \step{210seesr}{\case{$s\sim r$.}}
   \begin{proof}
    \pf\ Step \stepref{210s} gives $s\not\sim q$. Thus
    $\deg(p)=\deg(s)=\deg(q)$ by \stepref{210nonadjdeg}.~\qed
   \end{proof}
   \step{210seesq}{\case{$s\sim q$.}}
   \begin{proof}
    \step{210prdeg}{$\deg(p)=\deg(r)$.}
    \begin{proof}
     \pf\ Step \stepref{210s} gives $s\not\sim r$;
     apply \stepref{210nonadjdeg} to $p,s$ and $s,r$.~\qed
    \end{proof}
    \pflet{Choose $t\ne q$ nonadjacent to $q$, using \stepref{210triangles}.}
    \step{210t}{The vertex $t$ is outside $\{p,q,r\}$ and sees exactly one of $p,r$.}
    \begin{proof}\pf\ Since $p,r$ see $q$, the vertex $t$ differs from both. The triangle property gives one neighbour in $\{p,q,r\}$; it cannot be $q$.~\qed\end{proof}
    \step{210tdegrees}{$\deg(q)$ equals $\deg(r)$ or $\deg(p)$.}
    \begin{proof}
     \pf\ If $t\sim p$, then $t\not\sim r$, so $\deg(q)=\deg(t)=\deg(r)$.
     If $t\sim r$, then $t\not\sim p$, so $\deg(q)=\deg(t)=\deg(p)$.
     Both equalities use \stepref{210nonadjdeg}; \stepref{210t} exhausts the cases.~\qed
    \end{proof}
    \qedstep
    \begin{proof}\pf\ Use \stepref{210prdeg} in either alternative of \stepref{210tdegrees}.~\qed\end{proof}
   \end{proof}
   \qedstep
   \begin{proof}\pf\ Step \stepref{210s} makes \stepref{210seesr} and \stepref{210seesq} exhaustive.~\qed\end{proof}
  \end{proof}
  \qedstep
  \begin{proof}
   \pf\ Identical vertices have equal degrees; distinct vertices are covered by
   \stepref{210nonadjdeg} and \stepref{210adjdeg}.~\qed
  \end{proof}
 \end{proof}

 \pflet{Fix arbitrary distinct nonadjacent $u,v$ and retain the sets and labels of
 \stepref{210local}. Put $k=|K|$. By \stepref{210regular}, every vertex has degree $2k$.}
 \step{210kge}{$k\geq2$.}
 \begin{proof}
  \step{210kone}{If $k=1$, then $X$ is a single triangle.}
  \begin{proof}
   \pf\ By \stepref{210regular}, every vertex would have degree 2. A triangle exists by
   \stepref{210triangles}; its vertices already have two neighbours inside it, whereas
   every outside vertex would have a neighbour in it. Thus there is no outside vertex.~\qed
  \end{proof}
  \qedstep
  \begin{proof}
   \pf\ Step \stepref{210local} gives $k>0$; \stepref{210kone} contradicts the chosen nonadjacent pair when $k=1$.~\qed
  \end{proof}
 \end{proof}

 \step{210counts}{$|N|=3k-5$, $|V(X)|=6k-3$, and $X$ has strongly regular parameters
  $(6k-3,2k,1,k)$. The graph on $N$ has $k(k-2)$ edges, and the third vertex
  of every such edge lies in $K$.}
 \begin{proof}
  \step{210cross}{There are $4k^2-6k$ edges between $K\cup L\cup M$ and $N$.}
  \begin{proof}
   \step{210degreesintoN}{A vertex in $K$ has $2k-4$ neighbours in $N$;
    a vertex in $L$ or $M$ has $k-1$ neighbours in $N$.}
   \begin{proof}
    \pf\ By \stepref{210local}, $x\in K$ has the four neighbours $u,v,x',x''$
    outside $N$. A vertex $x'\in L$ has $u,x$ and the $k-1$ vertices $y''$ with
    $y\ne x$ outside $N$; a vertex $x''\in M$ has $v,x$ and the $k-1$ vertices
    $y'$ with $y\ne x$. Subtract these counts from degree $2k$.~\qed
   \end{proof}
   \qedstep
   \begin{proof}
    \pf\ Sum \stepref{210degreesintoN} over the three parts, each of size $k$:
    $k(2k-4)+k(k-1)+k(k-1)=4k^2-6k$.~\qed
   \end{proof}
  \end{proof}
  \step{210Nedges}{The subgraph on $N$ has $k(k-2)$ edges, each with its third vertex in $K$.}
  \begin{proof}
   \step{210Nthird}{If $w,z\in N$ are adjacent, then $f(w,z)\in K$.}
   \begin{proof}
    \step{210Noutside}{Both $u$ and $v$ lie outside $\{w,z,f(w,z)\}$.}
    \begin{proof}
     \pf\ They differ from $w,z$ by the definition of $N$.
     The vertex $f(w,z)$ is adjacent to $w$, whereas neither $u$ nor $v$ is;
     thus they also differ from $f(w,z)$.~\qed
    \end{proof}
    \qedstep
    \begin{proof}
     \pf\ By \stepref{210Noutside}, the triangle property applies to $u$ and $v$.
     Each sees neither $w$ nor $z$, so each sees $f(w,z)$. This is membership in $K$.~\qed
    \end{proof}
   \end{proof}
   \step{210Kpairs}{Each $x\in K$ is the third vertex of exactly $k-2$ edges in $N$.}
   \begin{proof}
    \pf\ In the triangle pairing of the $2k$ neighbours of $x$, the pairs
    $\{u,x'\}$ and $\{v,x''\}$ use its four neighbours outside $N$.
    The remaining $2k-4$ neighbours pair within $N$, giving $k-2$ edges.~\qed
   \end{proof}
   \qedstep
   \begin{proof}
    \pf\ Step \stepref{210Nthird} covers every edge in $N$; uniqueness of its third
    vertex prevents double counting. Sum \stepref{210Kpairs} over $|K|=k$.~\qed
   \end{proof}
  \end{proof}
  \step{210cardN}{$|N|=3k-5$ and $|V(X)|=6k-3$.}
  \begin{proof}
   \pf\ Counting the degree sum over $N$, using \stepref{210cross} and \stepref{210Nedges},
   gives
   \begin{align*}
    2k|N|&=2k(k-2)+(4k^2-6k)&&\text{[internal edges count twice]}\\
          &=2k(3k-5)&&\text{[ring arithmetic],}\\
    |N|&=3k-5&&\text{[$2k>0$],}\\
    |V(X)|&=2+3k+(3k-5)=6k-3&&\text{[the partition].}
   \end{align*}
   ~\qed
  \end{proof}
  \qedstep
  \begin{proof}
   \pf\ Step \stepref{210cardN} gives the order and \stepref{210regular} the degree.
   Step \stepref{210triangles} gives one common neighbour on each edge.
   For any nonadjacent pair, \stepref{210local} gives half its degree as the number
   of common neighbours, namely $k$. These are the four strongly regular parameters.~\qed
  \end{proof}
 \end{proof}

 \step{210dK}{\repair\ For $w\in N$, put
  $A(w)=K\diff\Gamma_1(w)$ and $d_K(w)=|K\cap\Gamma_1(w)|$.
  Then $|A(w)|\in\{0,2\}$, $d_K(w)\in\{k-2,k\}$, and
  $|N\cap\Gamma_1(w)|=d_K(w)$. We replace the unsupported pigeonhole argument
  by this invariant (Gap 1 of \path{cgss/errata.md}).}
 \begin{proof}
  \step{210KNpairing}{In the triangle pairing at $w$, its neighbours in $K$ pair
   with its neighbours in $N$. Its neighbours in $L$ pair with its neighbours in $M$.}
  \begin{proof}
   \step{210KtoN}{If $w\sim x\in K$, then $f(w,x)\in N$.}
   \begin{proof}
    \pf\ Each of $u,v$ is adjacent to $x$ and not to $w$. The triangle property
    excludes adjacency to $f(w,x)$. Also $f(w,x)\ne u,v$, since both are nonadjacent
    to $w$. The definition of $N$ applies.~\qed
   \end{proof}
   \step{210NtoK}{If $w\sim z\in N$, then $f(w,z)\in K$.}
   \begin{proof}
    \pf\ Apply the third-vertex assertion of \stepref{210counts} to the edge $wz$.~\qed
   \end{proof}
   \step{210LMclosed}{The partner of a neighbour of $w$ in $L\cup M$ lies in $L\cup M$.}
   \begin{proof}
    \pflet{Let $y\in(L\cup M)\cap\Gamma_1(w)$ and put $z=f(w,y)$.}
    \step{210partnerback}{$f(w,z)=y$.}
    \begin{proof}\pf\ The vertex $y$ is a common neighbour of $w,z$;
     uniqueness in \stepref{210triangles} gives the equality.~\qed\end{proof}
    \step{210partnernotUV}{$z\notin\{u,v\}$.}
    \begin{proof}\pf\ The vertex $z$ is adjacent to $w$, whereas $u,v$ are not.~\qed\end{proof}
    \step{210partnernotK}{$z\notin K$.}
    \begin{proof}\pf\ If $z\in K$, \stepref{210KtoN} gives
     $f(w,z)\in N$, contrary to \stepref{210partnerback} and $y\in L\cup M$.~\qed\end{proof}
    \step{210partnernotN}{$z\notin N$.}
    \begin{proof}\pf\ If $z\in N$, \stepref{210NtoK} gives
     $f(w,z)\in K$, contrary to \stepref{210partnerback} and $y\in L\cup M$.~\qed\end{proof}
    \qedstep
    \begin{proof}\pf\ The partition in \stepref{210local} and
     \stepref{210partnernotUV}, \stepref{210partnernotK}, \stepref{210partnernotN}
     leave only $L\cup M$.~\qed\end{proof}
   \end{proof}
   \qedstep
   \begin{proof}
    \pf\ The involution $y\mapsto f(w,y)$ pairs the neighbours of $w$.
    Steps \stepref{210KtoN} and \stepref{210NtoK} pair all its $K$ and $N$ neighbours.
    Step \stepref{210LMclosed} keeps every remaining pair in $L\cup M$.
    Neither part has internal edges by \stepref{210local}, so each such pair has
    one vertex in each part.~\qed
   \end{proof}
  \end{proof}
  \step{210Asize}{If $A(w)\ne\varnothing$, then $|A(w)|=2$.}
  \begin{proof}
   \pflet{Choose $x\in A(w)$.}
   \step{210common}{The common neighbours of $w,x'$ are exactly
    $\{y'':y\in A(w),\ y\ne x\}$.}
   \begin{proof}
    \pf\ The signature in \stepref{210local} gives $w\sim x'$.
    By \stepref{210KNpairing} their unique common neighbour lies in $M$.
    For $y\in K$, we have $w\sim y''$ exactly when $y\in A(w)$, and
    $x'\sim y''$ exactly when $y\ne x$, by \stepref{210local}.
    Thus the displayed set lists all, and only, common neighbours.~\qed
   \end{proof}
   \qedstep
   \begin{proof}
    \pf\ An edge has exactly one common neighbour by \stepref{210triangles}.
    The bijection $y\mapsto y''$ in \stepref{210local} and \stepref{210common}
    give $|A(w)|-1=1$.~\qed
   \end{proof}
  \end{proof}
  \qedstep
  \begin{proof}
   \pf\ Either $A(w)$ is empty or \stepref{210Asize} applies.
   Since $d_K(w)=k-|A(w)|$, its two possible values follow.
   The $K,N$ pairing in \stepref{210KNpairing} gives the asserted equality of degrees.~\qed
  \end{proof}
 \end{proof}

 \step{210degreeSum}{$\sum_{w\in N}d_K(w)=k(2k-4)$.}
 \begin{proof}
  \pf\ Each $x\in K$ has $2k-4$ neighbours in $N$: its four neighbours outside $N$
  are $u,v,x',x''$ by \stepref{210local}, and its degree is $2k$.
  Count edges between $K,N$ from the two ends.~\qed
 \end{proof}
 \pflet{$h\triangleq|\{w\in N:d_K(w)=k\}|$.}
 \step{210highcount}{$2h=(k-2)(5-k)$.}
 \begin{proof}
  \pf\ By \stepref{210dK} and \stepref{210counts}, the sum in \stepref{210degreeSum} is
  $(|N|-h)(k-2)+hk=(3k-5)(k-2)+2h$.
  Equating the two expressions and subtracting gives the formula.~\qed
 \end{proof}
 \step{210combbound}{\new\ $2\leq k\leq5$. We obtain the upper bound without a
  Gram matrix, so the result applies to an arbitrary graph with the triangle property.}
 \begin{proof}
  \step{210negative}{If $k>5$, then $(k-2)(5-k)<0$.}
  \begin{proof}\pf\ We have $k-2>0$ and $5-k<0$; multiply these inequalities.~\qed\end{proof}
  \qedstep
  \begin{proof}
   \pf\ The cardinality $h$ is nonnegative. Thus \stepref{210highcount} contradicts
   \stepref{210negative} if $k>5$. The lower bound is \stepref{210kge}.~\qed
  \end{proof}
 \end{proof}
 \step{210exclude}{\repair\ We have $k\in\{2,3,5\}$ (Gap 1).}
 \begin{proof}
  \step{210highforces}{If $h>0$, then $k=3$.}
  \begin{proof}
   \pflet{Choose $w_0\in N$ with $d_K(w_0)=k$.}
   \step{210choosez}{There is $z\in N$ adjacent to $w_0$.}
   \begin{proof}\pf\ Steps \stepref{210dK} and \stepref{210kge} give at least two neighbours of $w_0$ in $N$.~\qed\end{proof}
   \pflet{Choose such a $z$ and put $x=f(w_0,z)$.}
   \step{210zone}{$K\cap\Gamma_1(z)=\{x\}$.}
   \begin{proof}
    \pf\ Step \stepref{210counts} puts $x$ in $K$.
    If $y\in K\diff\{x\}$, then $y\sim w_0$ and $y\not\sim x$, since $K$ is independent.
    The triangle property for $\{w_0,z,x\}$ forces $y\not\sim z$.~\qed
   \end{proof}
   \qedstep
   \begin{proof}
    \pf\ Step \stepref{210zone} gives $d_K(z)=1$.
    By \stepref{210dK}, $1=k-2$ or $1=k$; \stepref{210kge} excludes the second,
    so $k=3$.~\qed
   \end{proof}
  \end{proof}
  \qedstep
  \begin{proof}
   \pf\ Step \stepref{210combbound} leaves $k=2,3,4,5$.
   At $k=4$, \stepref{210highcount} gives $h=1$, whereas
   \stepref{210highforces} gives $k=3$. Thus $k=4$ is impossible.~\qed
  \end{proof}
 \end{proof}

 \step{210two}{\repair\ (Gap 2.) If $k=2$, the labelled graph is uniquely determined:
  $N$ consists of one vertex $w$, with $A(w)=K$, and has no edges.}
 \begin{proof}
  \step{210twoN}{$|N|=1$ and $d_K(w)=0$.}
  \begin{proof}
   \pf\ Step \stepref{210counts} gives $|N|=1$.
   A simple graph on one vertex has degree 0, so \stepref{210dK} gives $d_K(w)=0$.~\qed
  \end{proof}
  \qedstep
  \begin{proof}
   \pf\ Step \stepref{210twoN} gives $A(w)=K$.
   The signature in \stepref{210local} makes $w$ adjacent to every vertex of $L\cup M$
   and to none of $K\cup\{u,v\}$. The remaining edges are all listed in
   \stepref{210local}. Thus every edge is determined by labels on the two elements of $K$.~\qed
  \end{proof}
 \end{proof}

 \step{210three}{\repair\ (Gap 2.) If $k=3$, the subgraph on $N$ is a claw. Its centre sees all
  of $K$ and its three leaves correspond bijectively to $K$ by their unique neighbour there.
  These labels determine the whole graph.}
 \begin{proof}
  \step{210centre}{There is exactly one $w_0\in N$ with $d_K(w_0)=3$;
   the other three vertices have $d_K=1$.}
  \begin{proof}
   \pf\ Step \stepref{210highcount} gives $2h=(3-2)(5-3)=2$.
   Thus $h=1$; \stepref{210counts} gives $|N|=4$, and \stepref{210dK}
   gives the values 1 or 3.~\qed
  \end{proof}
  \step{210claw}{The vertex $w_0$ is adjacent to all three other vertices of $N$,
   and they are pairwise nonadjacent.}
  \begin{proof}
   \pf\ By \stepref{210dK}, the degrees inside $N$ are 3 for $w_0$ and 1 for
   the others. Its three edges use the entire degree of each of the other vertices.~\qed
  \end{proof}
  \step{210leaflabels}{For each $x\in K$ there is exactly one leaf $w_x$ with
   $K\cap\Gamma_1(w_x)=\{x\}$.}
  \begin{proof}
   \pf\ By \stepref{210local} and \stepref{210regular}, the vertex $x$ has
   $2k-4=2$ neighbours in $N$, of which one is $w_0$.
   Its other neighbour is a leaf. Conversely, each leaf has exactly one neighbour in $K$
   by \stepref{210centre}. Thus this incidence is a bijection.~\qed
  \end{proof}
  \qedstep
  \begin{proof}
   \pf\ Steps \stepref{210centre}, \stepref{210claw}, and \stepref{210leaflabels}
   determine $N$ and its incidences with $K$.
   The signature in \stepref{210local} determines its incidences with $L,M$;
   the same step already lists all other edges.~\qed
  \end{proof}
 \end{proof}

 \step{210five}{\repair\ (Gap 2.) If $k=5$, the map
  $w\mapsto A(w)=K\diff\Gamma_1(w)$ is a bijection from $N$ to the two-element subsets
  of $K$. Two vertices of $N$ are adjacent exactly when their labels are disjoint.
  These labels determine the whole graph.}
 \begin{proof}
  \step{210pairsizes}{$d_K(w)=3$ and $|A(w)|=2$ for every $w\in N$.}
  \begin{proof}
   \pf\ Step \stepref{210counts} gives $|N|=10$ and \stepref{210degreeSum} gives
   $\sum_{w\in N}d_K(w)=5(10-4)=30$. Each summand is 3 or 5 by \stepref{210dK};
   the minimum total is already 30, so every summand is 3.~\qed
  \end{proof}
  \step{210pairlabel}{For distinct $x,y\in K$, there is exactly one $w\in N$
   with $A(w)=\{x,y\}$.}
  \begin{proof}
   \step{210edgeLM}{$x'$ and $y''$ are adjacent and their third vertex lies in $N$.}
   \begin{proof}
    \step{210noCore}{They have no common neighbour in $\{u,v\}\cup K\cup L\cup M$.}
    \begin{proof}
     \pf\ The vertex $u$ does not see $y''$, and $v$ does not see $x'$.
     Their respective neighbours in $K$ are the distinct vertices $x,y$.
     The set $L$ is independent and contains $x'$; the set $M$ is independent and
     contains $y''$. These observations exclude each of the five parts.~\qed
    \end{proof}
    \qedstep
    \begin{proof}
     \pf\ Step \stepref{210local} gives $x'\sim y''$.
     Their unique common neighbour exists by \stepref{210triangles};
     \stepref{210noCore} and the partition put it in $N$.~\qed
    \end{proof}
   \end{proof}
   \step{210labelvalue}{For $w=f(x',y'')$, $A(w)=\{x,y\}$.}
   \begin{proof}
    \pf\ Since $w$ sees $x',y''$, the signature in \stepref{210local} gives
    $x,y\in A(w)$. Step \stepref{210pairsizes} gives $|A(w)|=2$, so these are all its elements.~\qed
   \end{proof}
   \step{210labelunique}{Any $z\in N$ with $A(z)=\{x,y\}$ equals $w$.}
   \begin{proof}
    \pf\ The signature makes $z$ a common neighbour of $x',y''$;
    uniqueness of the third vertex gives $z=f(x',y'')=w$.~\qed
   \end{proof}
   \qedstep
   \begin{proof}\pf\ Steps \stepref{210labelvalue} and \stepref{210labelunique} give existence and uniqueness.~\qed\end{proof}
  \end{proof}
  \step{210disjoint}{For $w,z\in N$, $w\sim z$ if and only if $A(w)\cap A(z)=\varnothing$.}
  \begin{proof}
   \step{210adjdisjoint}{If $w\sim z$, then $A(w)\cap A(z)=\varnothing$.}
   \begin{proof}
    \pf\ Their third vertex $x$ is in $K$ by \stepref{210counts}.
    Every $y\in K\diff\{x\}$ is
    nonadjacent to $x$, so the triangle property makes it adjacent to exactly one of
    $w,z$; $x$ is adjacent to both. Thus no element of $K$ is missed by both.~\qed
   \end{proof}
   \step{210threechoices}{For fixed $w$, there are exactly three $z\in N$ whose labels
    are disjoint from $A(w)$, and exactly three neighbours of $w$ in $N$.}
   \begin{proof}
    \pf\ The complement of $A(w)$ in $K$ has three elements, and thus three two-element
    subsets. Step \stepref{210pairlabel} gives exactly one vertex for each.
    Steps \stepref{210dK} and \stepref{210pairsizes} give degree 3 within $N$.~\qed
   \end{proof}
   \qedstep
   \begin{proof}
    \pf\ Step \stepref{210adjdisjoint} gives inclusion of the neighbour set in the
    disjoint-label set; \stepref{210threechoices} gives equal finite cardinalities,
    so the two sets coincide.~\qed
   \end{proof}
  \end{proof}
  \qedstep
  \begin{proof}
   \pf\ Steps \stepref{210pairsizes} and \stepref{210pairlabel} give the bijection,
   and \stepref{210disjoint} gives the graph on $N$, the two-subset model of the Petersen
   graph. The labels give all incidences with $K$, their complements give those with
   $L,M$ by \stepref{210local}, and there are no incidences with $u,v$.~\qed
  \end{proof}
 \end{proof}

 \qedstep
 \begin{proof}
  \pf\ Steps \stepref{210local}, \stepref{210regular}, and \stepref{210counts}
  give the partition, edge rules, regularity and parameters.
  Steps \stepref{210dK}, \stepref{210degreeSum}, and \stepref{210highcount}
  give the degree identities. Step \stepref{210exclude} leaves exactly the three cases
  proved in \stepref{210two}, \stepref{210three}, and \stepref{210five}.
  The pair $u,v$ was arbitrary.~\qed
 \end{proof}
\end{proof}

\begin{namedresult}{Lemma}{(isomorphism from the order)}\label{lem:triangle-iso}
 \lean{SimpleGraph.HasTriangleProperty.nonempty_iso_of_card_eq}
 \uses{def:triangle-property,lem:triangle-normal-form}
 \assume{\begin{pfenum}
  \item $X,Y$ are finite simple graphs with the triangle property, and each has an edge.
  \item Neither graph has a vertex adjacent to every other vertex.
  \item $|V(X)|=|V(Y)|$.
 \end{pfenum}}
 \prove{There is a graph isomorphism $X\to Y$.}
\end{namedresult}
\begin{proof}
 \pfsketch\ \repair\ We will apply the normal-form lemma separately to the two graphs
 and match their labels. This completes the uniqueness argument of Gap 2.
 \step{210pairs}{Choose distinct nonadjacent $u,v\in V(X)$ and
  $\widetilde u,\widetilde v\in V(Y)$.}
 \begin{proof}
  \pf\ Each graph has a vertex since it has an edge. By \asm[2]{0}, that vertex
  has a distinct non-neighbour, in each graph.~\qed
 \end{proof}
 \step{210formX}{The normal-form lemma applies to $X,u,v$; write $k_X=|K|$.
  In particular $|V(X)|=6k_X-3$.}
 \begin{proof}\pf\ Apply Lemma~\ref{lem:triangle-normal-form} to the hypotheses on $X$
  and the first pair in \stepref{210pairs}.~\qed\end{proof}
 \step{210formY}{The normal-form lemma applies to $Y,\widetilde u,\widetilde v$;
  write $k_Y=|\widetilde K|$. In particular $|V(Y)|=6k_Y-3$.}
 \begin{proof}\pf\ Apply Lemma~\ref{lem:triangle-normal-form} to the hypotheses on $Y$
  and the second pair in \stepref{210pairs}.~\qed\end{proof}
 \step{210samek}{$k_X=k_Y$.}
 \begin{proof}\pf\ By \asm[3]{0}, \stepref{210formX} and \stepref{210formY},
  $6k_X-3=6k_Y-3$; cancellation gives the equality.~\qed\end{proof}
 \pflet{Write $k$ for this common value, use tildes for the parts of $Y$,
  and choose a bijection $\beta:K\to\widetilde K$, using \stepref{210samek}.}
 \step{210iso}{\repair\ The bijection $\beta$ extends to a graph isomorphism (Gap 2).}
 \begin{proof}
  \step{210extend}{Extend $\beta$ by $u\mapsto\widetilde u$, $v\mapsto\widetilde v$,
   $x'\mapsto\beta(x)'$, $x''\mapsto\beta(x)''$, and the following rule on $N$:
   send the single vertex to the single vertex when $k=2$; send centre to centre and
   $w_x$ to $\widetilde w_{\beta(x)}$ when $k=3$; send label $A$ to label $\beta(A)$
   when $k=5$. The resulting map $\psi$ is a bijection.}
  \begin{proof}
   \step{210extendcore}{The prescribed maps on $\{u\},\{v\},K,L,M$ are bijections
    onto the corresponding parts of $Y$.}
   \begin{proof}\pf\ The label maps from $K$ to $L,M$ are bijections in both
    \stepref{210formX} and \stepref{210formY}; compose these with $\beta$.~\qed\end{proof}
   \step{210extendtwo}{When $k=2$, the map on $N$ is a bijection.}
   \begin{proof}\pf\ Both sets are singletons by part 6 of the normal-form lemma.~\qed\end{proof}
   \step{210extendthree}{When $k=3$, the map on $N$ is a bijection.}
   \begin{proof}\pf\ Part 7 of the normal-form lemma gives one centre and a
    bijective leaf labelling by $K$ in each graph; $\beta$ matches the labels.~\qed\end{proof}
   \step{210extendfive}{When $k=5$, the map on $N$ is a bijection.}
   \begin{proof}\pf\ Part 8 of the normal-form lemma identifies both sets with
    two-element subsets, and $A\mapsto\beta(A)$ is a bijection on these subsets.~\qed\end{proof}
   \qedstep
   \begin{proof}\pf\ The parts partition both vertex sets by the normal-form lemma.
    Step \stepref{210extendcore} and the exhaustive cases
    \stepref{210extendtwo}, \stepref{210extendthree}, \stepref{210extendfive}
    give a bijection on their disjoint unions.~\qed\end{proof}
  \end{proof}
  \step{210preserve}{The map $\psi$ preserves adjacency and nonadjacency.}
  \begin{proof}
   \step{210coreedges}{It preserves all pairs with neither endpoint in $N$.}
   \begin{proof}
    \step{210coreuv}{The pair $u,v$ is a nonedge in each graph.}
    \begin{proof}\pf\ This is the choice in \stepref{210pairs}.~\qed\end{proof}
    \step{210coreuK}{Every pair $u,x$ with $x\in K$ is an edge in each graph.}
    \begin{proof}\pf\ This is the definition of $K$.~\qed\end{proof}
    \step{210coreuL}{Every pair $u,x'$ with $x'\in L$ is an edge in each graph.}
    \begin{proof}\pf\ This is the definition of $L$.~\qed\end{proof}
    \step{210coreuM}{Every pair $u,x''$ with $x''\in M$ is a nonedge in each graph.}
    \begin{proof}\pf\ This is the definition of $M$.~\qed\end{proof}
    \step{210corevK}{Every pair $v,x$ with $x\in K$ is an edge in each graph.}
    \begin{proof}\pf\ This is the definition of $K$.~\qed\end{proof}
    \step{210corevL}{Every pair $v,x'$ with $x'\in L$ is a nonedge in each graph.}
    \begin{proof}\pf\ This is the definition of $L$.~\qed\end{proof}
    \step{210corevM}{Every pair $v,x''$ with $x''\in M$ is an edge in each graph.}
    \begin{proof}\pf\ This is the definition of $M$.~\qed\end{proof}
    \step{210coresamePart}{Pairs within any one part $D\in\{K,L,M\}$ are nonedges in each graph.}
    \begin{proof}\pf\ Each such $D$ is independent by part 3 of the normal-form lemma.~\qed\end{proof}
    \step{210coreKL}{The relation $x\sim y'$ is preserved in both directions.}
    \begin{proof}\pf\ In each graph it is $x=y$, by part 3 of the normal-form lemma; $\beta$ preserves and reflects equality.~\qed\end{proof}
    \step{210coreKM}{The relation $x\sim y''$ is preserved in both directions.}
    \begin{proof}\pf\ In each graph it is $x=y$, by part 3 of the normal-form lemma; $\beta$ preserves and reflects equality.~\qed\end{proof}
    \step{210coreLM}{The relation $x'\sim y''$ is preserved in both directions.}
    \begin{proof}\pf\ In each graph it is $x\ne y$, by part 3 of the normal-form lemma; $\beta$ preserves and reflects inequality.~\qed\end{proof}
    \qedstep
    \begin{proof}\pf\ The listed pair types exhaust the five parts
     $\{u\},\{v\},K,L,M$. Loops are nonedges because the graphs are simple.~\qed\end{proof}
   \end{proof}
   \step{210mixededges}{It preserves all pairs with exactly one endpoint in $N$.}
   \begin{proof}
    \step{210mixeduv}{Vertices of $N$ are adjacent to neither $u$ nor $v$ in either graph.}
    \begin{proof}\pf\ This is the definition of $N$.~\qed\end{proof}
    \step{210mixedtwo}{For $k=2$, $\beta(A(w))=\widetilde A(\psi(w))$.}
    \begin{proof}\pf\ Both missed sets are the whole of $K$ or $\widetilde K$,
     respectively, by part 6 of the normal-form lemma.~\qed\end{proof}
    \step{210mixedthree}{For $k=3$, $\beta(A(w))=\widetilde A(\psi(w))$.}
    \begin{proof}
     \step{210mixedcentre}{The equality holds at the centre.}
     \begin{proof}\pf\ Both missed sets are empty by part 7 of the normal-form lemma.~\qed\end{proof}
     \step{210mixedleaf}{The equality holds at each leaf $w_x$.}
     \begin{proof}\pf\ Its missed set is $K\diff\{x\}$ and its image has missed set
      $\widetilde K\diff\{\beta(x)\}$, by part 7 of the normal-form lemma.~\qed\end{proof}
     \qedstep
     \begin{proof}\pf\ A claw consists of its centre and leaves; apply
      \stepref{210mixedcentre} and \stepref{210mixedleaf}.~\qed\end{proof}
    \end{proof}
    \step{210mixedfive}{For $k=5$, $\beta(A(w))=\widetilde A(\psi(w))$.}
    \begin{proof}\pf\ This is the definition of $\psi$ on the two-subset labels
     in \stepref{210extend}.~\qed\end{proof}
    \step{210mixedK}{Pairs between $N$ and $K$ are preserved in both directions.}
    \begin{proof}\pf\ Adjacency is $x\notin A(w)$ by definition.
     Steps \stepref{210mixedtwo}, \stepref{210mixedthree}, and \stepref{210mixedfive}
     preserve this condition in every possible case.~\qed\end{proof}
    \step{210mixedL}{Pairs between $N$ and $L$ are preserved in both directions.}
    \begin{proof}\pf\ The relation $w\sim x'$ is $x\in A(w)$ by part 4 of the
     normal-form lemma, preserved by the three missed-set equalities above.~\qed\end{proof}
    \step{210mixedM}{Pairs between $N$ and $M$ are preserved in both directions.}
    \begin{proof}\pf\ The relation $w\sim x''$ is $x\in A(w)$ by part 4 of the
     normal-form lemma, preserved by the same equalities.~\qed\end{proof}
    \qedstep
    \begin{proof}\pf\ Steps \stepref{210mixeduv}, \stepref{210mixedK},
     \stepref{210mixedL}, and \stepref{210mixedM} exhaust the other endpoint.~\qed\end{proof}
   \end{proof}
   \step{210internaledges}{It preserves all pairs with both endpoints in $N$.}
   \begin{proof}
    \step{210internaltwo}{When $k=2$, there are no edges in either $N$.}
    \begin{proof}\pf\ Both sets are singletons by part 6 of the normal-form lemma.~\qed\end{proof}
    \step{210internalthree}{When $k=3$, adjacency within $N$ is preserved in both directions.}
    \begin{proof}\pf\ Part 7 of the normal-form lemma says that the edges join
     the centre to a leaf, and \stepref{210extend} preserves these roles.~\qed\end{proof}
    \step{210internalfive}{When $k=5$, adjacency within $N$ is preserved in both directions.}
    \begin{proof}\pf\ Part 8 of the normal-form lemma says that adjacency is disjointness
     of labels; $\beta(A)\cap\beta(B)=\beta(A\cap B)$ since $\beta$ is injective.~\qed\end{proof}
    \qedstep
    \begin{proof}\pf\ The normal-form lemma gives $k\in\{2,3,5\}$, so the three
     preceding cases are exhaustive.~\qed\end{proof}
   \end{proof}
   \qedstep
   \begin{proof}\pf\ Steps \stepref{210coreedges}, \stepref{210mixededges}, and
    \stepref{210internaledges} exhaust all pairs.~\qed\end{proof}
  \end{proof}
  \qedstep
  \begin{proof}\pf\ The bijection \stepref{210extend} and the equivalence of
   adjacency in \stepref{210preserve} define a graph isomorphism.~\qed\end{proof}
 \end{proof}
 \qedstep
 \begin{proof}\pf\ Step \stepref{210iso} gives the required isomorphism.~\qed\end{proof}
\end{proof}

\begin{cgss}{Lemma}{2.10}\label{lem:2.10}
 \lean{LineSystem.Hyp24.card_Γa_mem}
 \lean{LineSystem.Hyp24.orthGraph_iso_of_card_Γa_eq}
 \lean{LineSystem.Hyp24.orthGraph_isSRGWith}
 \uses{hyp:2.4,lem:2.6,lem:2.9,cor:triangle-property,lem:triangle-normal-form,lem:triangle-iso,facts:inner,comp:C1}
 \assume{\begin{pfenum}
  \item Hypothesis 2.4 holds in the finite-dimensional real inner product space $V$,
   and $X$ is its orthogonality graph on $\Ga$.
  \item $X$ has an edge, is noncomplete, and no vertex is adjacent to every other vertex.
 \end{pfenum}}
 \prove{$|\Ga|\in\{9,15,27\}$, and the graph is unique up to isomorphism in each case.
  More explicitly, if another instance of Hypothesis 2.4 in a finite-dimensional real inner
  product space $W$ satisfies the same graph hypotheses and has the same value of $|\Ga|$,
  its orthogonality graph is isomorphic to $X$.
  Moreover, \repair\ $X$ is strongly regular with parameters $(6k-3,2k,1,k)$ for some
  $k\in\{2,3,5\}$ (Error 12).}
\end{cgss}
\begin{proof}
 \pfsketch\ The two graph lemmas give the conclusion. We will also retain the paper's
 Gram-matrix proof of the upper bound on $k$, including computation C1.
 By Lemma 2.9, having a universal vertex is equivalent here to being a friendship graph.
 Thus our precise hypothesis matches the paper's exclusion of that family; we retain its
 noncompleteness hypothesis as well.
 \step{210apply}{For any distinct nonadjacent $u,v$ and $k=|\Gamma_1(u)\cap\Gamma_1(v)|$,
  the normal-form conclusions hold. In particular $|\Ga|\in\{9,15,27\}$ and
  $X$ has the stated strongly regular parameters.}
 \begin{proof}
  \pf\ Corollary~\ref{cor:triangle-property} supplies the triangle property.
  Apply Lemma~\ref{lem:triangle-normal-form} and substitute $k=2,3,5$ in $6k-3$.~\qed
 \end{proof}
 \step{210unique}{Any second instance with the same graph hypotheses and order has
  an orthogonality graph isomorphic to $X$.}
 \begin{proof}
  \raggedright\pf\ Corollary~\ref{cor:triangle-property} gives the triangle property for each
  instance. Apply Lemma~\ref{lem:triangle-iso} to the two graphs.~\qed
 \end{proof}
 \step{210paperBound}{The paper's Gram-matrix argument gives a second proof of $k\leq5$.}
 \begin{proof}
  \pflet{Choose distinct nonadjacent $u,v$ and use the partition and labels of
   Lemma~\ref{lem:triangle-normal-form}. Write $\Zu,\Zv,\Zx,\Zx',\Zx''$ for their
   representing vectors and put $k=|K|\geq2$, by \stepref{210apply}.}
 \step{210gram}{The following are Gram matrices, first for the family
  \[\Zu,\Zv,(\Zx)_{x\in K},(\Zx')_{x\in K},(\Zx'')_{x\in K},\]
  then for its five block sums
  \[\Zu,\Zv,\sum_{x\in K}\Zx,\sum_{x\in K}\Zx',\sum_{x\in K}\Zx''.\]
  \[
  B_0=\begin{pmatrix}
  2&1&0^T&0^T&\Zj^T\\
  1&2&0^T&\Zj^T&0^T\\
  0&0&I+J&J-I&J-I\\
  0&\Zj&J-I&I+J&I\\
  \Zj&0&J-I&I&I+J
  \end{pmatrix},
  \]
  \[
  B=\begin{pmatrix}
  2&1&0&0&k\\
  1&2&0&k&0\\
  0&0&k^2+k&k^2-k&k^2-k\\
  0&k&k^2-k&k^2+k&k\\
  k&0&k^2-k&k&k^2+k
  \end{pmatrix}.
  \]
  Here $I,J$ have size $k$ and $\Zj$ has $k$ entries all equal to 1.
  We number the displayed rows and columns from 1.}
 \begin{proof}
  \step{210entries}{For distinct vertices $x,y$, $(\Zx,\Zy)=0$ on edges and 1 on nonedges.}
  \begin{proof}\pf\ Use the definition of $X$ and Lemma 2.6; diagonal products are 2.~\qed\end{proof}
  \step{210blocks}{The displayed blocks of $B_0$ follow from Lemma~\ref{lem:triangle-normal-form}.}
  \begin{proof}
   \pf\ Independent parts give $I+J$; matched edges give $J-I$;
   all edges except matched pairs give $I$. The products with $\Zu,\Zv$ follow from
   the definitions of $K,L,M$, using \stepref{210entries}.~\qed
  \end{proof}
  \step{210sums}{The matrix $B$ consists of the block sums of $B_0$.}
  \begin{proof}
   \pf\ Bilinearity gives each sum-vector product as the corresponding double sum.
   The sums of the entries of $I+J,J-I,I$ are $k^2+k,k^2-k,k$, respectively;
   those of $\Zj$ and 0 are $k$ and 0.~\qed
  \end{proof}
  \qedstep
  \begin{proof}
   \pf\ Steps \stepref{210blocks} and \stepref{210sums} prove the formulas.
   \repair\ In particular $B(3,1)=0$, correcting the printed 2 (Error 3 of
   \path{cgss/errata.md}): for every $x\in K$, the vector $\Zx$ is orthogonal to $\Zu$.~\qed
  \end{proof}
 \end{proof}

 \step{210bound}{$2\leq k\leq5$.}
 \begin{proof}
  \step{210psd}{The matrix $D$ obtained by deleting row and column 1 of $B$ is positive
   semidefinite, so $\det D\geq0$.}
  \begin{proof}
   \step{210Dgram}{$D$ is the Gram matrix of
    $\Zv,\sum_{x\in K}\Zx,\sum_{x\in K}\Zx',\sum_{x\in K}\Zx''$.}
   \begin{proof}\pf\ Delete the first vector from the block-sum family in
    \stepref{210gram}.~\qed\end{proof}
   \step{210Dpsd}{$D$ is positive semidefinite.}
   \begin{proof}
    \pf\ Apply \texttt{Matrix.posSemidef\_gram} in
    \path{Mathlib/Analysis/InnerProductSpace/GramMatrix.lean}
    directly to the four-vector family in \stepref{210Dgram}.~\qed
   \end{proof}
   \qedstep
   \begin{proof}
    \raggedright\pf\ Apply \path{Matrix.PosSemidef.det_nonneg} in
    \path{Mathlib/Analysis/Matrix/PosDef.lean} to \stepref{210Dpsd}.
    We index the real matrix by \texttt{Fin 4}, a finite type with decidable equality.~\qed
   \end{proof}
  \end{proof}
  \step{210det}{Computation C1: $\det D=2k^5(5-k)$.}\label{comp:C1}
  \lean{Computations.C1.det_blockSumMinor}\uses{}
  \begin{proof}
   \raggedright\pf\ Expanding along the first row $(2,0,k,0)$ gives
   \begin{align*}
    \det D&=2\bigl(k^5(7-k)\bigr)+k\bigl(-4k^4\bigr)
      &&\text{[the two $3\times3$ minors]}\\
     &=2k^5(5-k)&&\text{[ring arithmetic].}
   \end{align*}
   Each $3\times3$ minor is its six-term determinant expansion. Thus C1 is a polynomial
   identity, certified by \path{Matrix.det_succ_row_zero} in
   \path{Mathlib/LinearAlgebra/Matrix/Determinant/Basic.lean}, with $D$ indexed by
   \texttt{Fin (3+1)}. For each minor use \path{Matrix.det_fin_three} in the same file,
   with index type \texttt{Fin 3}, then ring arithmetic.~\qed
  \end{proof}
  \qedstep
  \begin{proof}
   \pf\ Steps \stepref{210psd} and \stepref{210det} give $2k^5(5-k)\geq0$.
   Since $k\geq2$ by \stepref{210apply}, division by $2k^5>0$ gives $k\leq5$.~\qed
  \end{proof}
 \end{proof}

  \qedstep
  \begin{proof}\pf\ Step \stepref{210bound} is the claimed second proof.~\qed\end{proof}
 \end{proof}
 \qedstep
 \begin{proof}\pf\ Steps \stepref{210apply} and \stepref{210unique} give the assertions
  of the lemma; \stepref{210paperBound} records the paper's bound as well.~\qed\end{proof}
\end{proof}

\begin{cgss}{Remark}{2.11}\label{rem:2.11}
 \uses{cor:triangle-property}
 \emph{Side claim, not formalized.} We quote [CGSS]:
 \begin{quote}
 Similar methods may be used to prove the characterization of symplectic and orthogonal
 graphs over $\mathrm{GF}(2)$ by the triangle property, cf. [24, 21].
 \end{quote}
 We attach no Lean declaration to this side claim.
\end{cgss}

\begin{cgss}{Lemma}{2.12}\label{lem:2.12}
 \lean{LineSystem.Iso.ofOrthGraphIso}\lean{LineSystem.Hyp24.orthGraph_cases}
 \lean{LineSystem.IsIndecomposable.eq_pair_or_exists_hyp24}
 \uses{def:line-iso,def:indecomposable,hyp:2.4,lem:star-identities,lem:2.5,lem:2.6,lem:2.7,lem:2.9,lem:2.10,facts:inner}
 \assume{\begin{pfenum}
  \item $S$ is a nonempty indecomposable star-closed line system in a finite-dimensional
   real inner product space $V$.
  \item When a star is chosen, $\Ga$ and $X$ have the meanings of Hypothesis 2.4.
 \end{pfenum}}
 \prove{Either $S$ has one line, or it admits a pair $\Za,\Zb$ as in Hypothesis 2.4
  and, for every such choice, its graph $X$ is one of:
  the empty graph $\comp{K_m}$ on $m\geq0$ vertices; the friendship graph with
  $m\geq1$ triangles; the graphs of Lemma 2.10 with parameters
  $(9,4,1,2)$, $(15,6,1,3)$, or $(27,10,1,5)$.
  For any two instances of Hypothesis 2.4 in finite-dimensional real inner product
  spaces $V,W$, every isomorphism of their graphs on $\Ga$ extends to an isomorphism
  of their line systems, induced by a surjective linear isometry of their spans
  carrying the ordered stars to each other. This holds for every graph type in the list.}
\end{cgss}
\noindent\repair\ Nonemptiness excludes the empty line system, as required by
Implicit hypothesis 4 of \path{cgss/errata.md}.
\begin{proof}
 \pfsketch\ \repair\ We give all four strongly regular parameters in the exceptional
 cases to delimit the paper's uniqueness statement (Error 12). We will recover the signed root set
 from Lemmas 2.5 and 2.7.
 A graph isomorphism preserves the Gram matrix of the roots in
 $\{\Za,\Zb,\Zc\}\cup\Ga$, so it extends linearly to the spans and then carries
 the recovered sets to each other.

 \step{212star}{\repair\ Either $S$ consists of one antipodal pair or it admits Hypothesis 2.4.
  We include the one-line case omitted by the choice of a star in the paper
  (Implicit hypothesis 4).}
 \begin{proof}
  \pflet{Choose $\Zr\in\Phi$, using nonemptiness.}
  \step{212split}{If every root outside $\{\Zr,-\Zr\}$ is orthogonal to $\Zr$,
   then $\Phi=\{\Zr,-\Zr\}$.}
  \begin{proof}
   \pf\ Take $P=\{\Zr,-\Zr\}$ and $Q=\Phi\diff P$.
   The hypothesis and bilinearity make these parts orthogonal.
   Indecomposability and $\Zr\in P$ force $Q=\varnothing$.~\qed
  \end{proof}
  \step{212orient}{Otherwise there are $\Za,\Zb\in\Phi$ with $(\Za,\Zb)=-1$.}
  \begin{proof}
   \pf\ Choose $\Zs\notin\{\Zr,-\Zr\}$ with $(\Zr,\Zs)\ne0$.
   Equation \eqref{eq:endpoints} and the line-system condition give $\pm1$.
   Set $\Za=\Zr$, and set $\Zb=\Zs$ if the value is $-1$, or $\Zb=-\Zs$ if it is 1.
   Closure under negation supplies both choices.~\qed
  \end{proof}
  \qedstep
  \begin{proof}\pf\ The cases \stepref{212split} and \stepref{212orient} are exhaustive.~\qed\end{proof}
 \end{proof}

 \step{212recover}{In any instance of Hypothesis 2.4, put
  $T=\{\Za,\Zb,\Zc\}\cup\Ga$. Then $\Phi=\cR\cup(-\cR)$ for the set $\cR$
  in Lemma 2.7, and $\operatorname{span}(T)=\operatorname{span}(\Phi)$.}
 \begin{proof}
  \step{212formulas}{We have $\Phi=\cR\cup(-\cR)$, where
   \[
   \cR=\{\Za,\Zb,\Zc\}\cup\Ga\cup(\Zc+\Ga)\cup(-\Zb+\Ga)
   \cup\{\Zx-\Zy:\Zx,\Zy\in\Ga,\ (\Zx,\Zy)=1\}.
   \]}
  \begin{proof}\pf\ This is the recovery equality in the statement of Lemma 2.7,
   including both inclusions.~\qed\end{proof}
  \qedstep
  \begin{proof}
   \pf\ Step \stepref{212formulas} recovers the set using only $T$ and its inner products.
   Every displayed vector lies in $\operatorname{span}(T)$.
   Since $T\subseteq\Phi$, the two spans are equal.~\qed
  \end{proof}
 \end{proof}

 \pflet{For the extension assertion, fix any two instances of Hypothesis 2.4 in
  finite-dimensional real inner product spaces $V,W$ and a graph isomorphism
  $\psi:X\to Y$. Use tildes for the second root set and ordered star.
  For $x\in V(X)$, write $\Zs_x$ for the vector at $\psi(x)\in V(Y)$.}
 \step{212metric}{\new\ There is a surjective linear isometry between the two root spans,
  sending $\Za,\Zb,\Zc$ to $\widetilde{\Za},\widetilde{\Zb},\widetilde{\Zc}$
  and $\Zx$ to $\Zs_x$ for every $x\in V(X)$.
  We make the metric meaning of ``determined'' explicit.}
 \begin{proof}
  \pflet{Index two vector families $(\Zr_i),(\Zs_i)$ by the same finite disjoint union
   $I=\{1,2,3\}\sqcup V(X)$.
   The first three values are the ordered stars; the other values are
   $\Zr_x=\Zx$ and the vectors $\Zs_x$ just defined.
   Put $R(\alpha)=\sum_i\alpha(i)\Zr_i$ and $R'(\alpha)=\sum_i\alpha(i)\Zs_i$.
   These are linear maps from $I\to\re$ to $V,W$, respectively, by
   \path{Fintype.linearCombination} in
   \path{Mathlib/LinearAlgebra/Finsupp/LinearCombination.lean}.}
  \step{212gram}{$(\Zr_i,\Zr_j)=(\Zs_i,\Zs_j)$ for all $i,j\in I$.}
  \begin{proof}
   \step{212stargram}{Products within the first three positions are 2 on the diagonal
    and $-1$ off it; products with a vertex of $X$ are $0,1,-1$ in that order.}
   \begin{proof}\pf\ These are the star identities and the defining products of $\Ga$
    in each system.~\qed\end{proof}
   \step{212graphgram}{Products within the graph positions are 2 on the diagonal,
    0 on edges, and 1 on nonedges.}
   \begin{proof}
    \pf\ Use Lemma 2.6 and the orthogonality convention; the bijection $\psi$ preserves
    equality, adjacency, and nonadjacency.~\qed
   \end{proof}
   \qedstep
   \begin{proof}\pf\ Steps \stepref{212stargram} and \stepref{212graphgram}, with symmetry, cover all pairs of positions.~\qed\end{proof}
  \end{proof}
  \step{212relations}{$\ker R\subseteq\ker R'$: for every $\alpha:I\to\re$,
   $R(\alpha)=0$ implies $R'(\alpha)=0$.}
  \begin{proof}
   \step{212norms}{The two sums have the same squared norm.}
   \begin{proof}
    \pf\ Expanding by bilinearity gives
    \[
    \sum_{i,j}\alpha(i)\alpha(j)(\Zr_i,\Zr_j)
     =\sum_{i,j}\alpha(i)\alpha(j)(\Zs_i,\Zs_j)
    \]
    by \stepref{212gram}.~\qed
   \end{proof}
   \qedstep
   \begin{proof}
    \pf\ If the first sum is zero, \stepref{212norms} makes the second sum have squared
    norm zero. The zero-norm implication in the elementary identities makes that sum zero.~\qed
   \end{proof}
  \end{proof}
  \step{212linear}{There is a surjective linear map $F$ of the two spans such that
   \[F\left(\sum_i\alpha(i)\Zr_i\right)=\sum_i\alpha(i)\Zs_i.\]}
  \begin{proof}
   \step{212ranges}{The ranges of $R,R'$ are the spans of $(\Zr_i),(\Zs_i)$, respectively.}
   \begin{proof}
    \raggedright\pf\ Apply \path{Fintype.range_linearCombination} in
    \path{Mathlib/LinearAlgebra/Finsupp/LinearCombination.lean} to each finite family.~\qed
   \end{proof}
   \pflet{Let $R'_0$ be $R'$ with codomain restricted to $\operatorname{range}R'$,
    using \path{LinearMap.rangeRestrict} in
    \path{Mathlib/Algebra/Module/Submodule/Range.lean}.}
   \step{212kernel}{$\ker R\subseteq\ker R'_0$.}
   \begin{proof}
    \pf\ A vector in the subspace $\operatorname{range}R'$ is zero exactly when its
    underlying vector is zero. Thus \stepref{212relations} gives the inclusion.~\qed
   \end{proof}
   \step{212quotient}{There is a linear map
    $\overline R':(I\to\re)/\ker R\to\operatorname{range}R'$
    satisfying $\overline R'([\alpha])=R'_0(\alpha)$.}
   \begin{proof}
    \pf\ Apply \path{Submodule.liftQ} to $R'_0$ and \stepref{212kernel};
    \path{Submodule.liftQ_apply} gives the formula. Both declarations are in
    \path{Mathlib/LinearAlgebra/Quotient/Basic.lean}.~\qed
   \end{proof}
   \step{212rangemap}{There is a linear map
    $F_0:\operatorname{range}R\to\operatorname{range}R'$
    with $F_0(R(\alpha))=R'(\alpha)$ for all $\alpha$.}
   \begin{proof}
    \pflet{Let $E:(I\to\re)/\ker R\to\operatorname{range}R$ be
     \path{LinearMap.quotKerEquivRange} for $R$, in
     \path{Mathlib/LinearAlgebra/Isomorphisms.lean}, and put $F_0=\overline R'\circ E^{-1}$.}
    \step{212compose}{$F_0$ is linear.}
    \begin{proof}\pf\ It is a composite of the linear map in \stepref{212quotient}
     and the inverse of a linear equivalence.~\qed\end{proof}
    \step{212oncombinations}{$F_0(R(\alpha))=R'(\alpha)$.}
    \begin{proof}
     \pf\ The lemma \path{LinearMap.quotKerEquivRange_symm_apply_image} in
     \path{Mathlib/LinearAlgebra/Isomorphisms.lean} gives $E^{-1}(R(\alpha))=[\alpha]$.
     Apply the formula in \stepref{212quotient}.~\qed
    \end{proof}
    \qedstep
    \begin{proof}\pf\ Steps \stepref{212compose} and \stepref{212oncombinations}
     prove both assertions.~\qed\end{proof}
   \end{proof}
   \step{212linearonto}{$F_0$ is surjective.}
   \begin{proof}\pf\ Every member of $\operatorname{range}R'$ has the form $R'(\alpha)$;
    by \stepref{212rangemap}, $R(\alpha)$ is a preimage.~\qed\end{proof}
   \step{212generatorvalues}{$F_0(\Zr_i)=\Zs_i$ for every $i$.}
   \begin{proof}
    \pf\ In \stepref{212rangemap}, take the coefficient function which is 1 at $i$
    and zero elsewhere. The two finite sums are $\Zr_i,\Zs_i$.~\qed
   \end{proof}
   \qedstep
   \begin{proof}\pf\ Identify the ranges with the spans by \stepref{212ranges}.
    The map $F_0$ in \stepref{212rangemap}, with surjectivity from
    \stepref{212linearonto}, is the required $F$ and has the values in
    \stepref{212generatorvalues}.~\qed\end{proof}
  \end{proof}
  \pflet{Choose a surjective linear map $F$ with the formula in \stepref{212linear}.}
  \step{212isometry}{The map $F$ preserves inner products and is injective.}
  \begin{proof}
   \step{212inner}{For vectors $\Zx,\Zy$ in the first span,
    $(F(\Zx),F(\Zy))=(\Zx,\Zy)$.}
   \begin{proof}
    \pf\ By the finite-span lemma, write $\Zx,\Zy$ as combinations with coefficients
    $\alpha,\beta$. The formula in \stepref{212linear} gives their images, and bilinearity gives
    \[
     \sum_{i,j}\alpha(i)\beta(j)(\Zr_i,\Zr_j)
       =\sum_{i,j}\alpha(i)\beta(j)(\Zs_i,\Zs_j)
    \]
    by \stepref{212gram}.~\qed
   \end{proof}
   \step{212injective}{$F$ is injective.}
   \begin{proof}
    \pf\ If $F(\Zx)=F(\Zy)$, linearity gives $F(\Zx-\Zy)=0$.
    Step \stepref{212inner} gives $(\Zx-\Zy,\Zx-\Zy)=0$.
    The zero-norm identity then gives $\Zx=\Zy$.~\qed
   \end{proof}
   \qedstep
   \begin{proof}\pf\ Combine \stepref{212inner} and \stepref{212injective}.~\qed\end{proof}
  \end{proof}
  \qedstep
  \begin{proof}
   \pf\ Steps \stepref{212linear} and \stepref{212isometry} give a linear isometry
   onto the second span. Step \stepref{212recover} identifies these spans with the root spans.
   Taking one coefficient equal to 1 and all others zero in \stepref{212linear}
   gives its prescribed values on the two indexed families.
   This is the equal-Gram argument LA2 from the prerequisites, proved here for the needed families.~\qed
  \end{proof}
 \end{proof}

 \pflet{Choose an isometry $F$ with the values and surjectivity in \stepref{212metric}.}
 \step{212rootiso}{The isometry $F$ in \stepref{212metric} carries the first root set
  bijectively onto the second.}
 \begin{proof}
  \step{212image}{$F(\Phi)\subseteq\widetilde\Phi$.}
  \begin{proof}
   \step{212imagea}{$F(\Za)=\widetilde{\Za}\in\widetilde\Phi$.}
   \begin{proof}\pf\ Use \stepref{212metric} and the second instance of Hypothesis 2.4.~\qed\end{proof}
   \step{212imageb}{$F(\Zb)=\widetilde{\Zb}\in\widetilde\Phi$.}
   \begin{proof}\pf\ Use \stepref{212metric} and the second instance of Hypothesis 2.4.~\qed\end{proof}
   \step{212imagec}{$F(\Zc)=\widetilde{\Zc}\in\widetilde\Phi$.}
   \begin{proof}\pf\ Use \stepref{212metric} and the star identities in the second instance.~\qed\end{proof}
   \step{212imagegamma}{For $x\in V(X)$, $F(\Zx)=\Zs_x\in\widetilde\Ga$.}
   \begin{proof}\pf\ Step \stepref{212metric} gives the value; $\Zs_x$ is the
    vector at $\psi(x)$, a vertex of the second graph.~\qed\end{proof}
   \step{212imageplus}{For $\Zx\in\Ga$, $F(\Zc+\Zx)=\widetilde{\Zc}+\Zs_x\in\widetilde\Phi$.}
   \begin{proof}\pf\ Linearity and \stepref{212metric} give the equality.
    Lemma 2.6 in the second instance, or the reverse inclusion of its recovery
    equality in Lemma 2.7, gives membership.~\qed\end{proof}
   \step{212imageminus}{For $\Zx\in\Ga$, $F(-\Zb+\Zx)=-\widetilde{\Zb}+\Zs_x\in\widetilde\Phi$.}
   \begin{proof}\pf\ Linearity and \stepref{212metric} give the equality.
    Apply the second translation in Lemma 2.6 in the target, equivalently the
    reverse inclusion of its recovery equality in Lemma 2.7.~\qed\end{proof}
   \step{212imagedifference}{If $\Zx,\Zy\in\Ga$ and $(\Zx,\Zy)=1$,
    then $F(\Zx-\Zy)=\Zs_x-\Zs_y\in\widetilde\Phi$.}
   \begin{proof}
    \step{212differenceproduct}{$(\Zs_x,\Zs_y)=1$.}
    \begin{proof}\pf\ The isometry in \stepref{212metric} preserves the given product.~\qed\end{proof}
    \qedstep
    \begin{proof}\pf\ Linearity gives the image formula. The reverse inclusion
     of the second recovery equality in Lemma 2.7 applies to
     \stepref{212differenceproduct}.~\qed\end{proof}
   \end{proof}
   \step{212imagesigns}{If a displayed expression $\Zr$ maps into $\widetilde\Phi$,
    so does $-\Zr$.}
   \begin{proof}\pf\ Linearity gives $F(-\Zr)=-F(\Zr)$ and
    $\widetilde\Phi=-\widetilde\Phi$.~\qed\end{proof}
   \qedstep
   \begin{proof}
    \pf\ The forward inclusion of the first recovery equality in Lemma 2.7
    expresses every root by one of these seven kinds and a sign.
    Steps \stepref{212imagea} through \stepref{212imagedifference}
    handle the expressions, and \stepref{212imagesigns} handles negation.~\qed
   \end{proof}
  \end{proof}
  \step{212onto}{$\widetilde\Phi\subseteq F(\Phi)$.}
  \begin{proof}
   \step{212ontostar}{Each root of the second ordered star has its corresponding
    first-star root as a preimage in $\Phi$.}
   \begin{proof}\pf\ Apply the prescribed star values in \stepref{212metric};
    Hypothesis 2.4 and the star identities put all three preimages in $\Phi$.~\qed\end{proof}
   \step{212ontogamma}{Each $\Zs\in\widetilde\Ga$ has a preimage $\Zx\in\Ga$.}
   \begin{proof}\pf\ Take the inverse image of its vertex under the bijection $\psi$
    and use \stepref{212metric}.~\qed\end{proof}
   \step{212ontoplus}{Each $\widetilde{\Zc}+\Zs$ with $\Zs\in\widetilde\Ga$
    has a preimage in $\Phi$.}
   \begin{proof}\pf\ Choose $\Zx$ by \stepref{212ontogamma}.
    The reverse inclusion of the first recovery equality in Lemma 2.7 puts
    $\Zc+\Zx$ in $\Phi$, and linearity gives its image.~\qed\end{proof}
   \step{212ontominus}{Each $-\widetilde{\Zb}+\Zs$ with $\Zs\in\widetilde\Ga$
    has a preimage in $\Phi$.}
   \begin{proof}\pf\ Choose $\Zx$ by \stepref{212ontogamma}.
    The reverse inclusion of the first recovery equality in Lemma 2.7 puts
    $-\Zb+\Zx$ in $\Phi$, and linearity gives its image.~\qed\end{proof}
   \step{212ontodifference}{Each $\Zs-\Zt$ with $\Zs,\Zt\in\widetilde\Ga$
    and $(\Zs,\Zt)=1$ has a preimage in $\Phi$.}
   \begin{proof}
    \pf\ Choose preimages $\Zx,\Zy$ by \stepref{212ontogamma}.
    The isometry in \stepref{212metric} gives $(\Zx,\Zy)=1$.
    The reverse inclusion of the first recovery equality in Lemma 2.7 puts
    $\Zx-\Zy$ in $\Phi$, and linearity gives its image.~\qed
   \end{proof}
   \step{212ontosigns}{The negative of a target with a preimage in $\Phi$
    also has a preimage in $\Phi$.}
   \begin{proof}\pf\ Negate the preimage, using $\Phi=-\Phi$ and linearity.~\qed\end{proof}
   \qedstep
   \begin{proof}
    \pf\ The forward inclusion of the second recovery equality in Lemma 2.7
    lists exactly the expressions and signs covered by \stepref{212ontostar}
    through \stepref{212ontosigns}.~\qed
   \end{proof}
  \end{proof}
  \qedstep
  \begin{proof}
   \pf\ Injectivity in \stepref{212metric}, \stepref{212image}, and \stepref{212onto}
   make the restriction of $F$ a bijection preserving inner products, as required
   by the definition of a line-system isomorphism.~\qed
  \end{proof}
 \end{proof}

 \step{212graphs}{Every graph arising from a chosen star belongs to the stated list.}
 \begin{proof}
  \step{212empty}{\case{$X$ has no edges.}}
  \begin{proof}\pf\ Its finite vertex set has some cardinality $m\geq0$, giving the first type.~\qed\end{proof}
  \step{212universal}{\case{$X$ has an edge and has a universal vertex.}}
  \begin{proof}
   \pf\ Lemma 2.9 gives a friendship graph with $m$ triangles.
   Since $X$ has an edge, $m\ne0$, so $m\geq1$.~\qed
  \end{proof}
  \step{212remaining}{\case{$X$ has an edge and has no universal vertex.}}
  \begin{proof}
   \pf\ A complete graph with an edge has universal vertices, so this graph is noncomplete.
   Thus Lemma 2.10 applies; its statement gives the displayed strongly regular parameters.~\qed
  \end{proof}
  \qedstep
  \begin{proof}\pf\ These three cases exhaust existence of an edge and of a universal vertex.~\qed\end{proof}
 \end{proof}
 \qedstep
 \begin{proof}
  \pf\ Step \stepref{212star} separates the one-line case; \stepref{212graphs} lists
  every graph in the other case. Steps \stepref{212metric} and \stepref{212rootiso}
  give the claimed uniqueness and its realization by a linear isometry of spans.~\qed
 \end{proof}
\end{proof}

\paragraph{Inputs from Section 3.}\label{input:sec3-models}
\repair\ The line counts and dimensions in Theorem 2.13 require the models of Section 3 (Gap 4 of \path{cgss/errata.md}).
We will use the following facts from Definitions 3.1 to 3.3, Lemma 3.4, and computations C2
and C4. We record them as deferred dependencies, not as conclusions of Section 2.
The labels \path{def:3.1}, \path{def:3.2}, \path{def:3.3}, \path{lem:3.4},
\path{comp:C2} and \path{comp:C4} are those of Section 3, which proves these facts.
Each model is a nonempty indecomposable star-closed line system in its real span. The model
$\typ{A}_1$ consists of one antipodal pair and has rank 1. Each remaining model has a specified
ordered star for which the following data hold. Here rank means dimension of the real span,
and $\delta$ means $|\Delta|/2$, the number of lines in $\Delta$.
\[
\begin{array}{c|c|c|c|c}
 \text{model}&\text{rank}&\text{number of lines}&\text{graph on }\Ga&\delta\\\hline
 \typ{A}_\ell\ (\ell\geq2)&\ell&\ell(\ell+1)/2&\comp{K_{\ell-2}}&\text{not used}\\
 \typ{D}_\ell\ (\ell\geq4)&\ell&\ell(\ell-1)&\text{friendship, }\ell-3\text{ triangles}&\text{not used}\\
 \typ{E}_6&6&36&(9,4,1,2)&6\\
 \typ{E}_7&7&63&(15,6,1,3)&15\\
 \typ{E}_8&8&120&(27,10,1,5)&36
\end{array}
\]
In the exceptional rows the tuple gives strongly regular parameters; the graph is the
orthogonality graph, so it has degree 4, 6, or 10, respectively. We take its uniqueness from
Lemma 2.10, not from the model verification. C2 gives the line counts and C4 verifies
star-closure, indecomposability and the marked-star data; spanning sets in the explicit
models give the displayed ranks. These verifications must not use Theorem 2.13 or Theorem 3.5.
For Definition 3.3 we use fixed representatives of a root and a star in $\typ{E}_8$.
The fixed-root construction of $\typ{E}_7$ and the fixed-star construction of $\typ{E}_6$
suffice here. Prerequisite LN5 is needed to compare a fixed-root construction with the
phrase ``any one of its lines'' in Definition 3.3, but is not used in Theorem 2.13.

\begin{cgss}{Theorem}{2.13}\label{thm:2.13}
 \lean{LineSystem.five_possibilities}\lean{LineSystem.five_possibilities'}
 \lean{LineSystem.five_possibilities_forall_hyp24}
 \lean{LineSystem.Iso.exists_ambient_isometry_of_same_model}
 \lean{LineSystem.Hyp24.six_lines_of_iso_friendship_zero}
 \uses{lem:2.12,lem:2.10,lem:2.9,def:line-iso,facts:inner,input:sec3-models,def:3.1,def:3.2,def:3.3,lem:3.4,comp:C2,comp:C4}
 \assume{\begin{pfenum}
  \item $V$ is a finite-dimensional real inner product space of dimension $\ell$.
  \item $S$ is a nonempty indecomposable star-closed line system in $V$, and
   $\operatorname{span}(\Phi)=V$.
 \end{pfenum}}
 \prove{The possibilities for $S$, up to isomorphism, are the five cases below.
 For every case with more than one line, there exist $\Za,\Zb\in\Phi$ with
 $(\Za,\Zb)=-1$ for which the stated data on $\Ga$ and $\Delta$ hold.
 \repair\ We exclude the empty system and separate the one-line case
 (Implicit hypothesis 4 of \path{cgss/errata.md}).
 \begin{pfenum}
  \item $\ell(\ell+1)/2$ lines, with the empty graph $\comp{K_{\ell-2}}$ on $\Ga$ when $\ell\geq2$;
   for $\ell=1$ there is one line and no chosen star;
  \item $\ell(\ell-1)$ lines for $\ell\geq4$, with the friendship graph of $\ell-3$ triangles on $\Ga$;
  \item 120 lines in dimension 8, with $|\Ga|=27$ and $|\Delta|/2=36$;
  \item 63 lines in dimension 7, with $|\Ga|=15$ and $|\Delta|/2=15$;
  \item 36 lines in dimension 6, with $|\Ga|=9$ and $|\Delta|/2=6$.
 \end{pfenum}
 Each case is unique up to an orthogonal transformation when the ambient space is $\re^\ell$.
 The zero-triangle case of Lemma 2.9 has $\ell=3$ and six lines; it is already in case 1
 and also satisfies the formula $\ell(\ell-1)=6$.}
 The Lean statement gives the data on $\Ga$ and $\Delta$ for some star, while in the paper they
 belong to the star fixed after Hypothesis 2.4, which is arbitrary. Every star gives the same
 data: \path{LineSystem.five_possibilities_forall_hyp24} is the theorem for every star
 $(\Za,\Zb)$ of Hypothesis 2.4 whose system spans $V$, with the graph on $\Ga$ and the counts of
 $\Ga$ and $\Delta$ taken at that star. It leaves out the one-line case, which has no star.
\end{cgss}
\begin{proof}
 \pfsketch\ We will match the graph from Lemma 2.12 with a specified star in one of the
 models. Its graph isomorphism gives a linear isometry of spans, so it transfers both
 rank and number of lines. Thus the numerical conclusions come from Section 3.

 \step{213one}{If $S$ has one line, then $\ell=1$ and a linear isometry of spans
  carries it onto $\typ{A}_1$.}
 \begin{proof}
  \step{213rankone}{The span of $\Phi=\{\Zr,-\Zr\}$ is one-dimensional.}
  \begin{proof}
   \pf\ Squared norm 2 gives $\Zr\ne0$. Thus the singleton $\{\Zr\}$ is a basis
   of this span, which equals $V$ by \asm[2]{0}.~\qed
  \end{proof}
  \step{213onemap}{For either root $\Zs$ of $\typ{A}_1$, the rule
   $F(\alpha\Zr)=\alpha\Zs$ defines a surjective linear isometry of spans.}
  \begin{proof}
   \pf\ Step \stepref{213rankone} gives a unique coefficient $\alpha$ for each vector.
   Addition and scaling of this coefficient make $F$ linear and surjective.
   For any real $\alpha,\beta$, both inner products are $2\alpha\beta$ by bilinearity.
   Thus $F$ preserves inner products.~\qed
  \end{proof}
  \qedstep
  \begin{proof}
   \pf\ The map in \stepref{213onemap} sends $\{\Zr,-\Zr\}$ bijectively onto
   $\{\Zs,-\Zs\}$. The line count is $1=1(1+1)/2$, and
   \stepref{213rankone} gives $\ell=1$.~\qed
  \end{proof}
 \end{proof}

 \step{213match}{If $S$ has more than one line, its chosen-star graph is isomorphic
  to the specified graph of one model in the input table.}
 \begin{proof}
  \pflet{Choose a star using Lemma 2.12 and let $X$ be its graph.}
  \step{213A}{\case{$X=\comp{K_m}$.}}
  \begin{proof}
   \pf\ The specified star of $\typ{A}_{m+2}$ has the empty graph on $(m+2)-2=m$ vertices
   by the Section 3 inputs. A bijection between the finite vertex sets preserves
   adjacency since both graphs have no edges.~\qed
  \end{proof}
  \step{213D}{\case{$X$ is the friendship graph with $m\geq1$ triangles.}}
  \begin{proof}
   \step{213Dsize}{The specified graph in $\typ{D}_{m+3}$ is the friendship graph with $m$ triangles.}
   \begin{proof}\pf\ Use $m+3\geq4$ and the $\typ{D}_\ell$ row of the Section 3 inputs.~\qed\end{proof}
   \qedstep
   \begin{proof}
    \pf\ Send a chosen common vertex to the model's common vertex, choose a bijection
    of the $m$ pairs of remaining vertices, and a bijection within each pair.
    Lemma 2.9 lists exactly the edges preserved by this map, so it is a graph isomorphism.~\qed
   \end{proof}
  \end{proof}
  \step{213E}{\case{$|V(X)|\in\{9,15,27\}$ as in Lemma 2.10.}}
  \begin{proof}
   \step{213Echoose}{Choose $\typ{E}_6,\typ{E}_7,\typ{E}_8$ according as the order is $9,15,27$.
    Its specified graph satisfies the hypotheses of Lemma 2.10.}
   \begin{proof}
    \pf\ The respective orders and degrees in the input table are $(9,4),(15,6),(27,10)$.
    Each degree is positive and smaller than order minus 1, so each graph has an edge, is
    noncomplete, and has no universal vertex. A graph of Lemma 2.9 has a universal
    vertex, so that type is excluded.~\qed
   \end{proof}
   \qedstep
   \begin{proof}
    \pf\ The chosen model is a star-closed indecomposable line system and has the same
    graph order as $X$. Apply the uniqueness assertion of Lemma 2.10.~\qed
   \end{proof}
  \end{proof}
  \qedstep
  \begin{proof}\pf\ Lemma 2.12 makes \stepref{213A}, \stepref{213D}, and \stepref{213E} exhaustive.~\qed\end{proof}
 \end{proof}

 \step{213transfer}{For a model matched in \stepref{213match}, $\ell$ equals its rank and
  $S$ has its number of lines. The marked sets $\Delta$ also have the same number of lines.}
 \begin{proof}
  \step{213spaniso}{The graph isomorphism extends to a surjective linear isometry $F$
   between the spans, carrying the ordered stars to each other and mapping root sets bijectively.}
  \begin{proof}\pf\ Apply the reconstruction and equal-Gram argument of Lemma 2.12.~\qed\end{proof}
  \step{213dimension}{The two spans have equal dimensions.}
  \begin{proof}
   \pf\ A linear isometry is injective; together with surjectivity in
   \stepref{213spaniso}, it sends a basis bijectively to a basis.
   The first span is $V$ by \asm[2]{0}.~\qed
  \end{proof}
  \step{213lines}{The root bijection induces a bijection of antipodal pairs.}
  \begin{proof}
   \pf\ Linearity gives $F(-\Zr)=-F(\Zr)$; injectivity prevents two different
   antipodal pairs from merging, and surjectivity on roots covers every target pair.~\qed
  \end{proof}
  \step{213delta}{The root bijection sends $\Delta$ exactly to the model's marked $\Delta$.}
  \begin{proof}
   \pf\ The definition of $\Delta$ asks for three zero inner products with the ordered
   star. Both $F$ and its inverse preserve these products and the ordered star by
   \stepref{213spaniso}. Thus membership is preserved in both directions.~\qed
  \end{proof}
  \qedstep
  \begin{proof}
   \pf\ Step \stepref{213dimension} transfers rank, \stepref{213lines} transfers line count,
   and \stepref{213delta} with \stepref{213lines} transfers the line count in $\Delta$.~\qed
  \end{proof}
 \end{proof}

 \step{213numbers}{The transferred data give exactly the five numerical cases in the statement.}
 \begin{proof}
  \step{213classical}{The empty-graph case gives $\ell=m+2$ and $\ell(\ell+1)/2$ lines;
   the friendship case with $m\geq1$ triangles gives $\ell=m+3$ and $\ell(\ell-1)$ lines.}
  \begin{proof}
   \pf\ Substitute the ranks of $\typ{A}_{m+2},\typ{D}_{m+3}$ from the input table into
   \stepref{213transfer}, then use their line-count formulas from Lemma 3.4 and C2.~\qed
  \end{proof}
  \step{213exceptional}{The orders $27,15,9$ give respectively the data
   \[(\ell,|\Phi|/2,|\Delta|/2)=(8,120,36),(7,63,15),(6,36,6).\]}
  \begin{proof}
   \pf\ Use the $\typ{E}_8,\typ{E}_7,\typ{E}_6$ rows of the Section 3 inputs in \stepref{213transfer}.
   \repair\ We correct Error 6: we write $|\Delta|/2$ because the signed set $\Delta$ has respectively
   $72,30,12$ vectors. The paper's $|\Delta|=36,15,6$ counts lines.~\qed
  \end{proof}
  \step{213overlap}{\new\ The zero-triangle case of Lemma 2.9 has six lines in dimension 3; we make this overlap explicit.}
  \begin{proof}
   \pf\ Its graph is a single vertex, the empty graph $\comp{K_1}$.
   The first formula in \stepref{213classical} gives $\ell=3$ and $3\cdot4/2=6=3\cdot2$.~\qed
  \end{proof}
  \qedstep
  \begin{proof}
   \pf\ Combine \stepref{213one}, \stepref{213classical}, \stepref{213exceptional},
   and \stepref{213overlap}.~\qed
  \end{proof}
 \end{proof}

 \step{213unique}{Each numerical case, with its stated graph type, is unique up to
  an orthogonal transformation in $\re^\ell$.}
 \begin{proof}
  \step{213samemodel}{Two systems in the same case are isomorphic to the same model.}
  \begin{proof}
   \pf\ In the first family the model rank is $\ell$, so the model is $\typ{A}_\ell$;
   in the second it is $\typ{D}_\ell$; the exceptional rows specify $\typ{E}_8,\typ{E}_7,\typ{E}_6$ individually.
   Use \stepref{213one} for $\ell=1$ and \stepref{213match} otherwise.~\qed
  \end{proof}
  \qedstep
  \begin{proof}
   \pf\ Compose the span isometry to that model with the inverse for the other system.
   Both root spans are the entire $\re^\ell$ by \asm[2]{0}. Thus the composite is a
   surjective linear isometry of $\re^\ell$, an orthogonal transformation, sending one
   line system onto the other.~\qed
  \end{proof}
 \end{proof}
 \qedstep
 \begin{proof}
  \pf\ Step \stepref{213numbers} gives the alternatives and \stepref{213unique} gives
  uniqueness. The asserted dimensions and counts depend on the explicit Section 3 inputs
  recorded above; no classification result of Section 3 was used.~\qed
 \end{proof}
\end{proof}

\section{Root systems}\label{sec:3}
\begingroup
\emergencystretch=3em

We will construct the models before using the classification. All dimensions in this
section are dimensions of spans. We use both signs of every root, with squared norm 2,
as in Section 2. The symbol $\ell$ denotes rank, and $N$ denotes the number of lines.
Coordinates in the text start at 1; Lean's coordinates in \texttt{Fin} start at 0.
An isometric inclusion means a linear isometry of spans carrying the source roots
into the target roots. It is proper when that inclusion of root sets is proper.

\new\ We distinguish a proved blueprint result from an existing Lean declaration.
The tags \verb|\lean| give the Lean names; every one of them names a declaration in
Lean, and \path{dependency-diagram.md}, at the top of the repository, maps the numbered results
to their declarations. In
computational leaves we specify the integer lists, the Boolean test, and the soundness theorem.
The script \path{blueprint/checks/sec3/check_sec3.py} recomputes these tests exactly; it is a
sanity check, not a kernel proof of a new certificate.
No classification of Root Systems from Remark 3.11 is used here.

\begin{namedresult}{Lemma}{(transport on spans)}\label{lem:sec3-transport}
 \lean{LineSystem.Iso.spanIsometry}
 \uses{def:line-iso,def:line-system,def:indecomposable,def:2.1,def:2.2,hyp:2.4,facts:inner}
 \assume{$S,T$ are line systems in finite-dimensional real inner product spaces
 $V,V'$, and $f:\Phi(S)\to\Phi(T)$ is a bijection preserving inner products.}
 \prove{There is a unique surjective linear isometry
 $F:\operatorname{span}\Phi(S)\to\operatorname{span}\Phi(T)$ with $F(\Zr)=f(\Zr)$.
 It preserves line count and rank, star-closure, indecomposability, and maximality
 within the span. Its inverse has the same properties. For $\Za,\Zb\in\Phi(S)$ with
 $(\Za,\Zb)=-1$, the map $f$ carries the sets $\Ga$ and $\Delta$ of
 Hypothesis~\ref{hyp:2.4} for $\Za,\Zb$ bijectively onto those for $f(\Za),f(\Zb)$,
 and preserves orthogonality in $\Ga$.}
\end{namedresult}
\begin{proof}
 \pfsketch\ \new\ We state the transport fact separately so later results do not
 cite a construction hidden inside a proof. We will factor finite linear combinations through their kernels, and
 then transport the defining predicates.
 \step{s3trrelations}{For any real coefficients $a(\Zr)$, the two sums
 $\sum a(\Zr)\Zr$ and $\sum a(\Zr)f(\Zr)$ have equal squared norms.}
 \begin{proof}
  \pf\ Expand both squared norms by bilinearity from Section~\ref{facts:inner}; each term
  agrees by the inner-product hypothesis.~\qed
 \end{proof}
 \step{s3trmap}{The rule $F(\sum a(\Zr)\Zr)=\sum a(\Zr)f(\Zr)$ is a linear map of spans.}
 \begin{proof}
  \step{s3trkernel}{A zero source combination has zero target combination.}
  \begin{proof}\pf\ Apply zero squared norm from Section~\ref{facts:inner} to
   \stepref{s3trrelations}.~\qed\end{proof}
  \step{s3trquotient}{The target combination map factors through the quotient by
   the kernel of the source combination map.}
  \begin{proof}\pf\ Apply \path{Submodule.liftQ} to the kernel inclusion
   in \stepref{s3trkernel}.~\qed\end{proof}
  \qedstep
  \begin{proof}\pf\ Identify the quotient with the source range by
   \path{LinearMap.quotKerEquivRange}, and the range with the span by
   \path{Fintype.range_linearCombination}.~\qed\end{proof}
 \end{proof}
 \step{s3trisometry}{The map $F$ is an isometry and is onto the target span.}
 \begin{proof}
  \step{s3trinner}{It preserves the inner product of two combinations.}
  \begin{proof}\pf\ Bilinearity from Section~\ref{facts:inner} expands that product into
   the preserved products of pairs of roots.~\qed\end{proof}
  \step{s3tronto}{Each target combination has a source combination mapping to it.}
  \begin{proof}\pf\ Reindex its coefficients by the inverse of $f$.~\qed\end{proof}
  \qedstep
  \begin{proof}\pf\ Use \stepref{s3trinner} and \stepref{s3tronto}.~\qed\end{proof}
 \end{proof}
 \step{s3trpredicates}{The asserted invariants are preserved.}
 \begin{proof}
  \step{s3trlines}{Antipodal pairs correspond bijectively.}
  \begin{proof}\pf\ Linearity gives $F(-\Zr)=-F(\Zr)$; use the root bijection.~\qed\end{proof}
  \step{s3trrank}{The ranks agree.}
  \begin{proof}\pf\ A linear bijection carries a basis to a basis.~\qed\end{proof}
  \step{s3trclosed}{Star-closure is preserved in each direction.}
  \begin{proof}\pf\ Transport the sum and inner product in Definition~\ref{def:2.1}
   through $F$ or $F^{-1}$.~\qed\end{proof}
  \step{s3trindecomp}{Indecomposability is preserved in each direction.}
  \begin{proof}\pf\ Pull an orthogonal partition back through the corresponding
   root bijection in the indecomposability definition in Section~\ref{def:indecomposable}.~\qed\end{proof}
  \step{s3trmaximal}{Maximality within the span is preserved in each direction.}
  \begin{proof}\pf\ Transport a candidate vector of Definition~\ref{def:2.2}
   through the surjective span isometry.~\qed\end{proof}
  \step{s3trmarked}{\new\ The map $f$ carries $\Ga$ and $\Delta$ for $\Za,\Zb$ onto
   those for $f(\Za),f(\Zb)$, preserving orthogonality.}
  \begin{proof}\pf\ These sets are defined in Hypothesis~\ref{hyp:2.4} by inner products
   with $\Za$, $\Zb$ and $\Zc=-\Za-\Zb$; $f$ preserves inner products and
   $F(\Zc)=-f(\Za)-f(\Zb)$ by linearity.~\qed\end{proof}
  \qedstep
  \begin{proof}\pf\ The six preceding steps are the six assertions.~\qed\end{proof}
 \end{proof}
 \qedstep
 \begin{proof}\pf\ The finite-span lemma in Section~\ref{facts:inner} forces the formula
  for any linear extension; thus \stepref{s3trmap} is unique. Combine
  \stepref{s3trisometry} and \stepref{s3trpredicates}.~\qed\end{proof}
\end{proof}
\begingroup\raggedright
The quotient and combination declarations are in the Mathlib files
\path{LinearAlgebra/Quotient/Basic.lean}, \path{LinearAlgebra/Isomorphisms.lean},
and \path{LinearAlgebra/Finsupp/LinearCombination.lean}, respectively.\par
\endgroup

\begin{namedresult}{Lemma}{(restriction to the span)}\label{lem:sec3-restrict}
 \lean{LineSystem.restrictSpan}\lean{LineSystem.Iso.restrictSpan}
 \uses{def:line-system,lem:sec3-transport,facts:inner}
 \assume{$S$ is a line system in a finite-dimensional real inner product space $V$,
 and $W=\operatorname{span}\Phi(S)$, with the inner product of $V$.}
 \prove{The vectors of $\Phi(S)$, as vectors of $W$, form the root set of a line system
 $S|_W$ in $W$, and $\Phi(S|_W)$ spans $W$. The identity of $\Phi(S)$ is a bijection
 $\Phi(S|_W)\to\Phi(S)$ preserving inner products. Thus $S|_W$ and $S$ have the same
 number of lines, the same rank, and the same sets $\Ga$ and $\Delta$ with the same
 orthogonalities, and each is nonempty, star-closed or indecomposable exactly when the
 other is. A surjective linear isometry from $W$ that carries $\Phi(S|_W)$ onto a root
 set carries $\Phi(S)$ onto it.}
\end{namedresult}
\begin{proof}
 \pfsketch\ \new\ Theorem 2.13 assumes that the roots span the ambient space. We will
 state the passage to the span once, since Theorems 3.5 and 3.6 use it for $S$ and for
 the models.
 \step{s3reslines}{$S|_W$ is a line system in $W$.}
 \begin{proof}\pf\ The inner product of $W$ is that of $V$ on vectors of $W$, by
  \path{Submodule.coe_inner}; the conditions of the line-system definition in
  Section~\ref{def:line-system} are conditions on these inner products and on negation.~\qed\end{proof}
 \step{s3resspan}{$\Phi(S|_W)$ spans $W$.}
 \begin{proof}\pf\ This is \path{Submodule.span_span_coe_preimage} for the set
  $\Phi(S)$.~\qed\end{proof}
 \step{s3resinner}{The identity of $\Phi(S)$ preserves inner products.}
 \begin{proof}\pf\ Apply \path{Submodule.coe_inner} to two roots.~\qed\end{proof}
 \step{s3resinvariants}{The two systems have the asserted common invariants.}
 \begin{proof}\pf\ Apply Lemma~\ref{lem:sec3-transport} to the bijection of
  \stepref{s3resinner}; nonemptiness is preserved by any bijection.~\qed\end{proof}
 \step{s3resisometry}{The assertion on surjective linear isometries holds.}
 \begin{proof}\pf\ The two root sets are the same vectors of $W$, and
  $W=\operatorname{span}\Phi(S)$ by definition.~\qed\end{proof}
 \qedstep
 \begin{proof}\pf\ Steps \stepref{s3reslines}, \stepref{s3resspan},
  \stepref{s3resinner}, \stepref{s3resinvariants} and \stepref{s3resisometry} are the
  assertions.~\qed\end{proof}
\end{proof}
The two span declarations are in the Mathlib files
\path{Analysis/InnerProductSpace/Subspace.lean} and \path{LinearAlgebra/Span/Defs.lean}.

\begin{namedresult}{Lemma}{(an equal-rank proper inclusion)}\label{lem:sec3-proper-inclusion}
 \lean{LineSystem.not_isMaximal_of_proper_isometric_inclusion}
 \uses{def:line-system,def:2.2,facts:inner}
 \assume{\begin{pfenum}
  \item $S$ and $T$ are line systems in finite-dimensional real inner product spaces.
  \item The spans $W_S=\operatorname{span}\Phi(S)$ and $W_T=\operatorname{span}\Phi(T)$
   have the same dimension.
  \item $G:W_S\to W_T$ is a linear isometry with $G(\Phi(S))\subsetneq\Phi(T)$.
 \end{pfenum}}
 \prove{$S$ is not maximal within $W_S$.}
\end{namedresult}
\begin{proof}
 \pfsketch\ \new\ We will pull back a target root outside the image; it is a candidate
 of Definition 2.2 that is not a root of $S$.
 \step{s3propsurj}{$G$ is onto $W_T$.}
 \begin{proof}\pf\ An isometry is injective, and the dimensions agree by \asm[2]{0};
  apply \path{LinearMap.injective_iff_surjective_of_finrank_eq_finrank}.~\qed\end{proof}
 \step{s3proppick}{\pick\ $\Zt\in\Phi(T)\diff G(\Phi(S))$ and $\Zx\in W_S$ with
  $G(\Zx)=\Zt$.}
 \begin{proof}\pf\ The root $\Zt$ exists by \asm[3]{0}, and $\Zx$ by
  \stepref{s3propsurj}.~\qed\end{proof}
 \step{s3propnorm}{$(\Zx,\Zx)=2$.}
 \begin{proof}\pf\ $G$ preserves inner products and $(\Zt,\Zt)=2$.~\qed\end{proof}
 \step{s3propproducts}{For every $\Zr\in\Phi(S)$, $(\Zx,\Zr)\in\{-2,-1,0,1,2\}$.}
 \begin{proof}\pf\ $(\Zx,\Zr)=(\Zt,G(\Zr))$, and $\Zt,G(\Zr)\in\Phi(T)$ by
  \asm[3]{0}; use the line-system definition in Section~\ref{def:line-system}.~\qed\end{proof}
 \step{s3propnew}{$\Zx\notin\Phi(S)$.}
 \begin{proof}\pf\ Otherwise $\Zt=G(\Zx)$ would lie in $G(\Phi(S))$, contrary to the
  choice of $\Zt$.~\qed\end{proof}
 \qedstep
 \begin{proof}\pf\ Steps \stepref{s3propnorm}, \stepref{s3propproducts} and
  \stepref{s3propnew} give a vector of $W_S$ that Definition~\ref{def:2.2} requires in
  $\Phi(S)$ for a maximal $S$.~\qed\end{proof}
\end{proof}
The rank declaration is in the Mathlib file \path{LinearAlgebra/FiniteDimensional/Lemmas.lean}.

\begin{cgss}{Definition}{3.1}\label{def:3.1}
 \lean{LineSystem.An}\lean{LineSystem.Dn}\uses{def:line-system}
 \assume{$\ell$ is an integer, $\ell\geq1$ for $\typ A_\ell$ and $\ell\geq4$
 for $\typ D_\ell$; $(\Ze_i)$ is the standard orthonormal coordinate basis.}
 \pflet{The signed roots of the two systems are}
 \[
 \begin{split}
 \Phi(\typ A_\ell)&=\{\Ze_i-\Ze_j:1\leq i,j\leq\ell+1,\ i\ne j\},\\
 \Phi(\typ D_\ell)&=\{s\Ze_i+t\Ze_j:1\leq i<j\leq\ell,\ s,t\in\{1,-1\}\}.
 \end{split}
 \]
 These are line systems: each root has squared norm 2; taking negatives stays in
 the displayed set; pairwise products are integers of absolute value at most 2.
 We verify these assertions below with the other models.
 The inclusion $\Phi(\typ A_\ell)\subseteq\Phi(\typ D_{\ell+1})$ uses these coordinates.
 For $\ell<4$, the Lean definition \texttt{Dn} still uses this coordinate formula.
 The classification only uses $\typ D_\ell$ for $\ell\geq4$.
\end{cgss}

\begin{remark}
 \lean{LineSystem.cuboctahedronVertices}\lean{LineSystem.cuboctahedron}
 \lean{LineSystem.cuboctahedronDiagonals}
 \lean{LineSystem.cuboctahedronVertices_eq}\lean{LineSystem.ncard_cuboctahedronVertices}
 \lean{LineSystem.ncard_cubeEdges}\lean{LineSystem.IsCubeEdge.sym2_eq_of_midpoint_eq}
 \lean{LineSystem.cuboctahedron_Φ_inter_span_singleton}
 \lean{LineSystem.numLines_cuboctahedron}\lean{LineSystem.ncard_cuboctahedronDiagonals}
 \lean{LineSystem.cuboctahedron_eq_Dn}\lean{LineSystem.cuboctahedron.isoAn}
 \lean{LineSystem.nonempty_iso_cuboctahedron_An}
 The paper adds that $\typ A_3$ is represented by the 6 diagonals of the cuboctahedron. This is
 a side claim, and we prove it in \path{LineGraphs/SideClaims/Cuboctahedron.lean}. The vertices of
 the cuboctahedron are the midpoints of the twelve edges of the cube with vertices
 $(\pm1,\pm1,\pm1)$, and distinct edges have distinct midpoints. The midpoints are the twelve
 vectors with one coordinate 0 and two coordinates $\pm1$, which is $\Phi(\typ D_3)$ in the
 coordinates above. The vertex opposite to $\Zv$ is $-\Zv$, so each diagonal carries exactly two
 vertices, and the six diagonals are the lines of the line system $\typ D_3$
 (\path{LineSystem.cuboctahedron_eq_Dn}). This system is isomorphic to $\typ A_3$
 (\path{LineSystem.nonempty_iso_cuboctahedron_An}); in Lean the isomorphism is checked on the
 integer roots, by the kernel.
\end{remark}

\begin{cgss}{Definition}{3.2}\label{def:3.2}
 \lean{LineSystem.E₈viaA₈}\lean{LineSystem.E₈viaD₈}
 \lean{LineSystem.E₈viaSteiner}\lean{LineSystem.E₈}
 \lean{LineSystem.steiner348}\uses{def:3.1,def:line-system}
 \assume{Coordinate spaces have their usual real inner products.}
 \pflet{The following three signed sets, with the indicated normalization, define
 the three models of $\typ E_8$. The default is model (ii).}
 \begin{pfenum}
  \item In $\re^9$, take $\Phi(\typ A_8)$ together with both signs of
  $\frac13\sum_{i=1}^9\Ze_i-\sum_{i\in T}\Ze_i$, for $T\in\binom{[9]}3$.
  \item In $\re^8$, take $\Phi(\typ D_8)$ together with
  $\frac12\sum_{i=1}^8\epsilon(i)\Ze_i$, where $\epsilon(i)\in\{1,-1\}$ and
  $\prod_i\epsilon(i)=1$.
  \item In $\re^8$, take $\pm\sqrt2\Ze_i$ and
  $\frac1{\sqrt2}\sum_{i\in B}\epsilon(i)\Ze_i$ for all sign choices and all
  $B\in\mathcal B$, where the blocks of $\mathcal B=S(3,4,8)$ are
  \[
  \begin{array}{rrrrrrr}
  1234&1256&1278&1357&1368&1458&1467\\
  2358&2367&2457&2468&3456&3478&5678.
  \end{array}
  \]
 \end{pfenum}
 \repair\ This fixes the labelling required by Gap 5. We correct the second
 printed ``(i)'' to (ii) (Error 6), and remove the repeated sign condition
 (Error 13). \new\ The factor $1/\sqrt2$ in (iii) changes the printed vectors
 of squared norm 4 to our uniform squared norm 2, without changing their lines.
\end{cgss}

\begin{namedresult}{Lemma}{(integer realizations)}\label{lem:sec3-integer}
 \lean{Computations.C2.model_presentations}
 \uses{def:3.1,def:3.2,def:line-system,facts:inner}
 \assume{The coordinate sets are those of Definitions~\ref{def:3.1} and~\ref{def:3.2}.}
 \prove{They are line systems. Write $\Zr=c\Zv$ with integral $\Zv$ and
 $c^2q=1$. Their presentations are as follows; each list of raw vectors has no repetitions.
 \[
 \begin{array}{c|c|c|l}
 \text{model}&q&c&\text{raw vectors }\Zv\\\hline
 \typ A_\ell&1&1&\Ze_i-\Ze_j\ (i\ne j)\\
 \typ D_\ell&1&1&s\Ze_i+t\Ze_j\ (i<j)\\
 \typ E_8\text{ (i)}&9&1/3&3(\Ze_i-\Ze_j),\quad\pm(\mathbf1-3\mathbf1_T)\\
 \typ E_8\text{ (ii)}&4&1/2&2(s\Ze_i+t\Ze_j),\quad\epsilon\ (\prod\epsilon=1)\\
 \typ E_8\text{ (iii)}&2&1/\sqrt2&\pm2\Ze_i,\quad\epsilon\mathbf1_B
 \end{array}
 \]
 Every raw squared norm is $2q$, and every raw product lies in
 $q\{-2,-1,0,1,2\}$. Each triple of points belongs to exactly one block of $\mathcal B$;
 distinct blocks intersect in 0 or 2 points.}
\end{namedresult}
\begin{proof}
 \pfsketch\ \new\ We make the integer normalization and injective indices explicit
 to match the Lean models. We will check coordinates before using the common integer-model transfer.
 \step{s3intclassical}{The assertions hold for $\typ A_\ell,\typ D_\ell$.}
 \begin{proof}
  \step{s3intclassnorm}{Every such vector has exactly two nonzero coordinates, each a sign.}
  \begin{proof}\pf\ Read the two displays in Definition~\ref{def:3.1}.~\qed\end{proof}
  \step{s3intclassproducts}{Its squared norm is 2 and the products are integers between $-2$ and 2.}
  \begin{proof}\pf\ Sum the two coordinate products from \stepref{s3intclassnorm}.~\qed\end{proof}
  \step{s3intclassunique}{Supports and signs recover the displayed indices uniquely.}
  \begin{proof}\pf\ In $\typ A_\ell$ the positive and negative positions give the ordered
   pair; in $\typ D_\ell$ the sorted support gives the pair and its entries give the signs.~\qed\end{proof}
  \qedstep
  \begin{proof}\pf\ Negating the displayed signs stays in each list;
   use \stepref{s3intclassproducts} and \stepref{s3intclassunique}.~\qed\end{proof}
 \end{proof}
 \step{s3intone}{The assertions hold for model (i).}
 \begin{proof}
  \step{s3intonetriple}{For $t_T=\mathbf1-3\mathbf1_T$,
  $t_T\cdot t_U=9(|T\cap U|-1)$.}
  \begin{proof}\pf\ Expand the indicators and use $|T|=|U|=3$ in a 9-set.~\qed\end{proof}
  \step{s3intonemixed}{$3(\Ze_i-\Ze_j)\cdot t_T=9(\mathbf1_T(j)-\mathbf1_T(i))$.}
  \begin{proof}\pf\ Substitute the $i,j$ entries of $t_T$.~\qed\end{proof}
  \step{s3intoneunique}{The signed triple vectors are distinct and disjoint from the pair vectors.}
  \begin{proof}\pf\ A triple vector has three entries $-2$ and six entries 1;
   its negative has different entry values, while a pair vector has seven zero entries.~\qed\end{proof}
  \qedstep
  \begin{proof}\pf\ Use \stepref{s3intonetriple}, \stepref{s3intonemixed},
   \stepref{s3intoneunique}, and \stepref{s3intclassical} scaled by 3.~\qed\end{proof}
 \end{proof}
 \step{s3inttwo}{The assertions hold for model (ii).}
 \begin{proof}
  \step{s3inttwohalf}{For two parity-even sign vectors $u,v$, their Hamming distance
   $d$ is even, and $u\cdot v=8-2d\in\{-8,-4,0,4,8\}$.}
  \begin{proof}\pf\ The product of the eight coordinatewise products is 1,
   so the number of $-1$ products is even.~\qed\end{proof}
  \step{s3inttwomixed}{$2(s\Ze_i+t\Ze_j)\cdot\epsilon=2(s\epsilon(i)+t\epsilon(j))\in\{-4,0,4\}$.}
  \begin{proof}\pf\ Each summand in parentheses is a sign.~\qed\end{proof}
  \step{s3inttwounique}{The pair and half-root lists are disjoint and individually injective.}
  \begin{proof}\pf\ Pair roots have six zeros; half roots have none and their
   first seven signs determine the last by parity.~\qed\end{proof}
  \qedstep
  \begin{proof}\pf\ Use the last three steps and \stepref{s3intclassical};
   reversing all eight signs preserves parity.~\qed\end{proof}
 \end{proof}
 \step{s3intthree}{The assertions hold for model (iii).}
 \begin{proof}
  \step{s3intblocks}{The stated block sizes, intersections, and unique-triple property hold.}
  \begin{proof}\pf\ Enumerate the 14 displayed blocks and the 56 triples.
   These are the existing kernel results \path{LineSystem.steinerBlock_card},
   \path{LineSystem.steinerBlock_inter_even}, and
   \path{LineSystem.existsUnique_subset_steiner348} in \path{Models.lean}.~\qed\end{proof}
  \step{s3intthreeproducts}{Raw squared norms are 4, and products are even integers of absolute value at most 4.}
  \begin{proof}
   \step{s3intaxes}{Two axes have product 0 or $\pm4$; an axis and a block vector have product 0 or $\pm2$.}
   \begin{proof}\pf\ The axis has its only nonzero entry equal to $\pm2$.~\qed\end{proof}
   \step{s3intblockproducts}{Two block vectors have an even product between $-4$ and 4.}
   \begin{proof}\pf\ By \stepref{s3intblocks}, the product sums 0, 2, or 4 signs.~\qed\end{proof}
   \qedstep
   \begin{proof}\pf\ The two root types are exhausted by \stepref{s3intaxes}
    and \stepref{s3intblockproducts}; each self-product is 4.~\qed\end{proof}
  \end{proof}
  \step{s3intthreeunique}{The raw list is injective and closed under negation.}
  \begin{proof}\pf\ Support size distinguishes axes from blocks; support and entries
   recover every index and sign. All signs are allowed.~\qed\end{proof}
  \qedstep
  \begin{proof}\pf\ Use \stepref{s3intthreeproducts} and \stepref{s3intthreeunique}.~\qed\end{proof}
 \end{proof}
 \qedstep
 \begin{proof}\pf\ Apply \path{LineSystem.ofIntegerRoots} and
  \path{LineSystem.inner_integerVector} to the five checked presentations.
  These declarations are in \path{LineGraphs/LineSystem/IntegerModel.lean}.~\qed\end{proof}
\end{proof}

\begin{cgss}{Definition}{3.3}\label{def:3.3}
 \lean{LineSystem.E₇}\lean{LineSystem.E₆}
 \lean{LineSystem.T₇}\lean{LineSystem.T₆}
 \uses{def:3.2,lem:sec3-integer,def:line-system}
 \assume{$\typ E_8$ is model (ii) in $\re^8$.}
 \pflet{$\Zr_0=\Ze_7-\Ze_8$, $\Zs_0=\Ze_6-\Ze_7$, and}
 \[
 \Phi(\typ E_7)=\{\Zx\in\Phi(\typ E_8):x_7=x_8\},\qquad
 \Phi(\typ E_6)=\{\Zx\in\Phi(\typ E_8):x_6=x_7=x_8\}.
 \]
 Thus $\typ E_7=\typ E_8\cap\Zr_0^\perp$ and
 $\typ E_6=\typ E_8\cap\Zr_0^\perp\cap\Zs_0^\perp$.
 The roots $\Zr_0,\Zs_0,-\Zr_0-\Zs_0$ form a star.
 \repair\ We fix these choices first, and prove independence of both choices below
 (Gap 5); no transitivity hypothesis is part of the definition.
 For the printed counting argument we also name two cuts in model (iii):
 $T_7=\typ E_8^{\rm(iii)}\cap\Ze_8^\perp$ and
 $T_6=T_7\cap(\Ze_5+\Ze_6+\Ze_7+\Ze_8)^\perp$.
 Orthogonal cuts are line systems by restriction of the defining conditions.
\end{cgss}

\begin{namedresult}{Lemma}{(spans of the models)}\label{lem:sec3-spans}
 \lean{LineSystem.finrank_span_An}\lean{LineSystem.finrank_span_Dn}
 \lean{LineSystem.finrank_span_E₈}\lean{LineSystem.finrank_span_E₈viaD₈}
 \lean{LineSystem.finrank_span_E₈viaA₈}\lean{LineSystem.finrank_span_E₈viaSteiner}
 \lean{LineSystem.finrank_span_E₇}\lean{LineSystem.finrank_span_E₆}
 \uses{def:3.1,def:3.2,def:3.3,facts:inner}
 \assume{The models have the parameter ranges in Definitions~\ref{def:3.1}--\ref{def:3.3}.}
 \prove{Their spans and ranks are
 \[
 \begin{array}{c|c|c}
 \typ A_\ell&\{\Zx\in\re^{\ell+1}:\sum x_i=0\}&\ell\\
 \typ D_\ell&\re^\ell&\ell\\
 \typ E_8\text{ (i)}&\{\Zx\in\re^9:\sum x_i=0\}&8\\
 \typ E_8\text{ (ii), (iii)}&\re^8&8\\
 \typ E_7&\{\Zx\in\re^8:x_7=x_8\}&7\\
 \typ E_6&\{\Zx\in\re^8:x_6=x_7=x_8\}&6.
 \end{array}
 \]}
\end{namedresult}
\begin{proof}
 \pfsketch\ \repair\ We supply the dimension inputs deferred by Theorem 2.13
 (Gap 4). We will give spanning vectors with explicit independent coordinates.
 \step{s3spanA}{The $\typ A_\ell$ span is the stated hyperplane of dimension $\ell$.}
 \begin{proof}
  \step{s3spanAzero}{Every root has coordinate sum zero.}
  \begin{proof}\pf\ Its two nonzero coordinates are 1 and $-1$.~\qed\end{proof}
  \step{s3spanAbasis}{The vectors $\Ze_i-\Ze_{\ell+1}$, $1\leq i\leq\ell$, form a basis of this hyperplane.}
  \begin{proof}\pf\ A vector in the hyperplane is
   $\sum_{i=1}^\ell x_i(\Ze_i-\Ze_{\ell+1})$; its first $\ell$ coordinates make
   these coefficients unique.~\qed\end{proof}
  \qedstep
  \begin{proof}\pf\ The basis roots of \stepref{s3spanAbasis} and the containment
   in \stepref{s3spanAzero} give equality.~\qed\end{proof}
 \end{proof}
 \step{s3spanD}{The $\typ D_\ell$ roots span $\re^\ell$.}
 \begin{proof}\pf\ For $j\ne i$,
  $\Ze_i=((\Ze_i+\Ze_j)+(\Ze_i-\Ze_j))/2$ lies in the root span.~\qed\end{proof}
 \step{s3spanEi}{Model (i) spans the zero-sum hyperplane in $\re^9$.}
 \begin{proof}\pf\ Its additional roots have sum $9/3-3=0$, and it contains the
  $\typ A_8$ hyperplane of \stepref{s3spanA}.~\qed\end{proof}
 \step{s3spanEii}{Model (ii) spans $\re^8$.}
 \begin{proof}\pf\ It contains $\typ D_8$, whose span is given by \stepref{s3spanD}.~\qed\end{proof}
 \step{s3spanEiii}{Model (iii) spans $\re^8$.}
 \begin{proof}\pf\ Its eight axis roots are nonzero scalar multiples of a coordinate basis.~\qed\end{proof}
 \step{s3spanSeven}{The $\typ E_7$ span has basis $\Ze_1,\ldots,\Ze_6,\Ze_7+\Ze_8$.}
 \begin{proof}
  \step{s3spanSevenin}{These seven vectors lie in the span.}
  \begin{proof}\pf\ The $\typ D_6$ roots on the first six coordinates give the six axes
   by \stepref{s3spanD}; $\Ze_7+\Ze_8$ is itself an $\typ E_7$ root.~\qed\end{proof}
  \qedstep
  \begin{proof}\pf\ The displayed vectors uniquely expand every vector with
   $x_7=x_8$, which contains all the roots by Definition~\ref{def:3.3}.~\qed\end{proof}
 \end{proof}
 \step{s3spanSix}{The $\typ E_6$ span has basis $\Ze_1,\ldots,\Ze_5,\Ze_6+\Ze_7+\Ze_8$.}
 \begin{proof}
  \step{s3spanSixaxes}{The first five axes belong to the span.}
  \begin{proof}\pf\ Apply \stepref{s3spanD} to the $\typ D_5$ roots on the first five coordinates.~\qed\end{proof}
  \step{s3spanSixtail}{The last basis vector belongs to the span.}
  \begin{proof}\pf\ Subtract $\sum_{i=1}^5\Ze_i$ from twice the root
   $\frac12\mathbf1$, using \stepref{s3spanSixaxes}.~\qed\end{proof}
  \qedstep
  \begin{proof}\pf\ These six vectors uniquely expand every vector with
   $x_6=x_7=x_8$, which contains the roots by Definition~\ref{def:3.3}.~\qed\end{proof}
 \end{proof}
 \qedstep
 \begin{proof}\pf\ The six cases exhaust the models. The cited dimension declarations
  already implement these arguments in \path{LineGraphs/LineSystem/ModelSpans.lean}.~\qed\end{proof}
\end{proof}

\begin{namedresult}{Computation}{C2}\label{comp:C2}
 \lean{Computations.C2.line_counts}
 \lean{LineSystem.numLines_T₇}\lean{LineSystem.numLines_T₆}
 \uses{def:3.1,def:3.2,def:3.3,lem:sec3-integer}
 \assume{The models and the fixed cuts $T_7,T_6$ are as defined above.}
 \prove{Their line counts are $\ell(\ell+1)/2$, $\ell(\ell-1)$,
 120 in each $\typ E_8$ model, 63 for $\typ E_7,T_7$, and 36 for $\typ E_6,T_6$.}
\end{namedresult}
\begin{proof}
 \pfsketch\ We will count injective signed indices and divide by 2.
 \step{s3countA}{$\typ A_\ell$ has $\ell(\ell+1)/2$ lines.}
 \begin{proof}\pf\ There are $(\ell+1)\ell$ ordered distinct pairs in
  Definition~\ref{def:3.1}; injectivity is Lemma~\ref{lem:sec3-integer}.~\qed\end{proof}
 \step{s3countD}{$\typ D_\ell$ has $\ell(\ell-1)$ lines.}
 \begin{proof}\pf\ There are $4\binom\ell2$ signed indices; divide by 2.~\qed\end{proof}
 \step{s3countEone}{Model (i) has $36+84=120$ lines.}
 \begin{proof}\pf\ Its disjoint types contribute $\binom92$ and $\binom93$ antipodal pairs.~\qed\end{proof}
 \step{s3countEtwo}{Model (ii) has $56+64=120$ lines.}
 \begin{proof}\pf\ Its disjoint types contribute $2\binom82$ and $2^7/2$ antipodal pairs.~\qed\end{proof}
 \step{s3countEthree}{Model (iii) has $8+8\cdot14=120$ lines.}
 \begin{proof}\pf\ Each of the 14 supports contributes $2^4/2$ pairs, in addition to the eight axes.~\qed\end{proof}
 \step{s3countSeven}{$\typ E_7$ has $(60+2+64)/2=63$ lines.}
 \begin{proof}
  \step{s3countSevenpairs}{Its pair roots are $\typ D_6$ and $\pm(\Ze_7+\Ze_8)$.}
  \begin{proof}\pf\ The condition $x_7=x_8$ permits a support inside $[6]$, or the
   support $\{7,8\}$ with equal signs.~\qed\end{proof}
  \step{s3countSevenhalf}{It has $2^5\cdot2=64$ half roots.}
  \begin{proof}\pf\ The first six signs have even parity; the last two are equal with two choices.~\qed\end{proof}
  \qedstep
  \begin{proof}\pf\ Add the two disjoint counts and divide by 2.~\qed\end{proof}
 \end{proof}
 \step{s3countSix}{$\typ E_6$ has $(40+32)/2=36$ lines.}
 \begin{proof}
  \step{s3countSixpairs}{Its pair roots are exactly $\typ D_5$ on $[5]$.}
  \begin{proof}\pf\ A two-element support cannot meet the three required equal
   nonzero coordinates $6,7,8$.~\qed\end{proof}
  \step{s3countSixhalf}{There are $2^5=32$ half roots.}
  \begin{proof}\pf\ The first five signs determine the common last sign by the even-parity condition.~\qed\end{proof}
  \qedstep
  \begin{proof}\pf\ Add the disjoint counts and divide by 2.~\qed\end{proof}
 \end{proof}
 \step{s3countTseven}{$T_7$ has $7+8\cdot7=63$ lines.}
 \begin{proof}\pf\ Exactly seven blocks in Definition~\ref{def:3.2} avoid 8;
  each contributes eight lines, in addition to seven axes.~\qed\end{proof}
 \step{s3countTsix}{$T_6$ has $4+8+4\cdot6=36$ lines.}
 \begin{proof}
  \step{s3countTsixaxes}{The four surviving axes have indices $1,2,3,4$.}
  \begin{proof}\pf\ Impose $x_8=0$ and $x_5+x_6+x_7=0$ on an axis.~\qed\end{proof}
  \step{s3countTsixblock}{The block 1234 contributes eight lines.}
  \begin{proof}\pf\ Both linear conditions vanish on every sign choice on this block.~\qed\end{proof}
  \step{s3countTsixother}{Each of the other six blocks avoiding 8 contributes four lines.}
  \begin{proof}\pf\ Each meets $\{5,6,7\}$ twice; those signs must be opposite.
   There are $2\cdot2^2=8$ signed roots per block.~\qed\end{proof}
  \qedstep
  \begin{proof}\pf\ The three disjoint root types give the stated sum.~\qed\end{proof}
 \end{proof}
 \qedstep
 \begin{proof}\pf\ All counts use the injectivity in Lemma~\ref{lem:sec3-integer}.
  The common Lean transfer is \path{LineSystem.IntegerModel.numLines_eq}.
  The finite cardinal certificates \leancite{LineSystem.card_E₇Index}{\texttt{LineSystem.\allowbreak card\_\allowbreak E\ensuremath{{}_7}Index}} and
  \leancite{LineSystem.card_E₆Index}{\texttt{LineSystem.\allowbreak card\_\allowbreak E\ensuremath{{}_6}Index}} already check 126 and 72 in the kernel.~\qed\end{proof}
\end{proof}

\begin{namedresult}{Computation}{C3}\label{comp:C3}
 \lean{LineSystem.E₈viaA₈.isoE₈}\lean{LineSystem.E₈viaSteiner.isoE₈}
 \uses{def:3.2,lem:sec3-integer,lem:sec3-spans,comp:C2,lem:sec3-transport}
 \assume{The three models have the normalization of Definition~\ref{def:3.2}.}
 \prove{Model (i) is isometric to model (ii) by
 \[
 U(\Zx)_i=d_i(x_i+x_9/2)\quad(1\leq i\leq8),\qquad
 d_1=\cdots=d_7=1,\quad d_8=-1,
 \]
 on the zero-sum hyperplane. Model (iii) is isometric to model (ii) by
 $Q/\sqrt2$, where $Q$ is block diagonal with four blocks
 $\left(\begin{smallmatrix}1&1\\1&-1\end{smallmatrix}\right)$
 on the coordinate pairs $12,34,56,78$. Both maps carry the full signed root sets onto one another.}
\end{namedresult}
\begin{proof}
 \pfsketch\ \new\ We give explicit maps, so equivalence of the definitions does not
 use Theorem 3.5. In raw coordinates their formulas are rational, even though the
 normalization of the Steiner map introduces $\sqrt2$.
 \step{s3c3Umetric}{$U$ preserves inner products on the zero-sum hyperplane.}
 \begin{proof}\pf\ For $\sum x_i=\sum y_i=0$, expansion gives
  $\sum_{i=1}^8(x_i+x_9/2)(y_i+y_9/2)=\sum_{i=1}^9x_iy_i$.
  The factors $d_i^2=1$ do not change this sum.~\qed\end{proof}
 \step{s3c3Uroots}{$U$ sends every root of model (i) to a root of model (ii).}
 \begin{proof}
  \step{s3c3Upair}{\case{$\Zx=\Ze_i-\Ze_j$, with $i,j\leq8$.}}
  \begin{proof}\pf\ Its image is $d_i\Ze_i-d_j\Ze_j$, a $\typ D_8$ root.~\qed\end{proof}
  \step{s3c3Unine}{\case{$\Zx=\Ze_i-\Ze_9$, with $i\leq8$.}}
  \begin{proof}\pf\ Before the factor $d$, the image has one entry $1/2$ and
   seven entries $-1/2$. Negating coordinate 8 makes its sign parity even.~\qed\end{proof}
  \step{s3c3Utripleout}{\case{$\Zx=\mathbf1/3-\mathbf1_T$ and $9\notin T$.}}
  \begin{proof}\pf\ Before $d$, the image is $\mathbf1/2-\mathbf1_T$, with
   three negative signs; the coordinate-8 sign change makes their number even.~\qed\end{proof}
  \step{s3c3Utriplein}{\case{$\Zx=\mathbf1/3-\mathbf1_T$ and $T=\{i,j,9\}$.}}
  \begin{proof}\pf\ Its image is $-d_i\Ze_i-d_j\Ze_j$, a $\typ D_8$ root.~\qed\end{proof}
  \step{s3c3Usigns}{Negatives of the listed roots also map to roots.}
  \begin{proof}\pf\ $U$ is linear and model (ii) is closed under negation by
   Lemma~\ref{lem:sec3-integer}.~\qed\end{proof}
  \qedstep
  \begin{proof}\pf\ The pair and triple alternatives in Definition~\ref{def:3.2}
   are exhausted by the five preceding steps.~\qed\end{proof}
 \end{proof}
 \step{s3c3Qmetric}{$Q/\sqrt2$ preserves inner products.}
 \begin{proof}\pf\ Each displayed block has transpose times itself equal to $2I_2$,
  so $Q^TQ=2I_8$.~\qed\end{proof}
 \step{s3c3Qroots}{$Q/\sqrt2$ sends the roots of model (iii) into model (ii).}
 \begin{proof}
  \step{s3c3Qaxis}{An axis root $\pm\sqrt2\Ze_i$ maps to a $\typ D_8$ root.}
  \begin{proof}\pf\ The $i$th column of $Q$ has two entries in $\{1,-1\}$.~\qed\end{proof}
  \step{s3c3Qunion}{\case{$B$ is the union of two pairs among $12,34,56,78$.}}
  \begin{proof}\pf\ On each occupied pair $Q$ sends two signs to one entry $\pm2$
   and one zero; division by 2 gives a $\typ D_8$ root.~\qed\end{proof}
  \step{s3c3Qtransversal}{\case{$B$ meets each of the four pairs once.}}
  \begin{proof}
   \step{s3c3Qeven}{The number of second positions selected by $B$ is even.}
   \begin{proof}\pf\ Read the eight transversal blocks in Definition~\ref{def:3.2}:
    \[\begin{array}{rrrr}1357&1368&1458&1467\\2358&2367&2457&2468.\end{array}\]
    \qed\end{proof}
   \qedstep
   \begin{proof}\pf\ Selecting a first position gives a pair of signs with product 1;
    selecting a second gives product $-1$. Step \stepref{s3c3Qeven} makes the total
    product 1, so $Q\Zv/2$ is a half root of model (ii).~\qed\end{proof}
  \end{proof}
  \qedstep
  \begin{proof}\pf\ The other six blocks are $1234,1256,1278,3456,3478,5678$,
   the unions in \stepref{s3c3Qunion}; together with axes these exhaust the raw types.~\qed\end{proof}
 \end{proof}
 \step{s3c3onto}{Both root maps are onto.}
 \begin{proof}\pf\ Each is injective by its isometry property and maps between two
  finite sets of 240 roots by C2, so each is surjective.~\qed\end{proof}
 \qedstep
 \begin{proof}\pf\ Use \stepref{s3c3onto} and the span descriptions in
  Lemma~\ref{lem:sec3-spans}. Equivalently, the induced index bijections satisfy
  the hypothesis of the existing constructor \path{LineSystem.Iso.ofIntegerModels}.~\qed\end{proof}
\end{proof}
For a kernel certificate the raw maps are
$\Zv\mapsto(d_i(2v_i+v_9)/3)_{i=1}^8$ and $\Zv\mapsto Q\Zv$.
The 240 images in each case are integral. In Lean (\path{Computations/C3.lean}) the kernel
checks that the lists of images are permutations of the roots of model (ii), and checks the
scaled Gram equalities; the script \path{blueprint/checks/sec3/check_sec3.py} checks the images,
their bijections, and the scaled Gram equalities in exact arithmetic.

\begin{namedresult}{Lemma}{(a connected spanning set)}\label{lem:sec3-connected-span}
 \lean{LineSystem.isIndecomposable_of_connected_spanning}
 \uses{def:line-system,def:indecomposable,facts:inner}
 \assume{$S$ is a line system in a finite-dimensional real inner product space;
 $B\subseteq\Phi(S)$ is nonempty, spans $\operatorname{span}\Phi(S)$, and its
 graph joining distinct roots of nonzero inner product is connected.}
 \prove{$S$ is indecomposable.}
\end{namedresult}
\begin{proof}
 \pfsketch\ \new\ We isolate the spanning criterion for the model verifications.
 An orthogonal partition cannot separate an edge of $B$, or a root
 from every vector of a spanning set.
 \step{s3connpartition}{\assume{$P\cup Q=\Phi(S)$ and $P\perp Q$.}
  \prove{$P=\varnothing$ or $Q=\varnothing$.}}
 \begin{proof}
  \step{s3connunique}{No root belongs to both parts.}
  \begin{proof}\pf\ Such a root would have self-product both 0 and 2.~\qed\end{proof}
  \step{s3connedge}{The endpoints of every edge of $B$ lie in the same part.}
  \begin{proof}\pf\ Different parts have inner product zero by \asm{1}.~\qed\end{proof}
  \step{s3connpath}{All of $B$ lies in one part.}
  \begin{proof}\pf\ Choose a vertex of the nonempty $B$ and propagate its part
   along a path to each vertex, using \stepref{s3connedge}.~\qed\end{proof}
  \step{s3connother}{The other part is empty.}
  \begin{proof}
   \step{s3connzero}{Every vector of that part is orthogonal to the whole span.}
   \begin{proof}\pf\ It is orthogonal to $B$ by \asm{1} and
    \stepref{s3connpath}; extend by the finite-span lemma in Section~\ref{facts:inner}.~\qed\end{proof}
   \qedstep
   \begin{proof}\pf\ A root in that part is in the span, so
    \stepref{s3connzero} would give its self-product as 0 instead of 2.~\qed\end{proof}
  \end{proof}
  \qedstep
  \begin{proof}\pf\ Use \stepref{s3connother}, whichever part contains $B$.~\qed\end{proof}
 \end{proof}
 \qedstep
 \begin{proof}\pf\ Step \stepref{s3connpartition} is the indecomposability definition in Section~\ref{def:indecomposable}.~\qed\end{proof}
\end{proof}

\begin{namedresult}{Computation}{C4}\label{comp:C4}
 \lean{Computations.C4.exceptional_star_data}
 \lean{Computations.C4.marked_star_data}
 \lean{LineSystem.E₆.latticeGraphIso}\lean{LineSystem.E₇.triangularGraphComplIso}
 \lean{LineSystem.Hyp24.nonempty_latticeGraph_iso_of_S_eq_E₆}
 \lean{LineSystem.Hyp24.nonempty_triangularGraph_compl_iso_of_S_eq_E₇}
 \lean{LineSystem.E₆.exists_isometry_star}\lean{LineSystem.E₇.exists_isometry_star}
 \lean{LineSystem.E₈.exists_isometry_star}
 \lean{LineSystem.Hyp24.orthGraph_isSRGWith_of_S_eq_E₆}\lean{LineSystem.Hyp24.orthGraph_isSRGWith_of_S_eq_E₇}
 \lean{LineSystem.Hyp24.orthGraph_isSRGWith_of_S_eq_E₈}
 \lean{LineSystem.Hyp24.card_Γa_of_S_eq_E₆}\lean{LineSystem.Hyp24.card_Γa_of_S_eq_E₇}
 \lean{LineSystem.Hyp24.card_Γa_of_S_eq_E₈}
 \lean{LineSystem.Hyp24.card_Δ_of_S_eq_E₆}\lean{LineSystem.Hyp24.card_Δ_of_S_eq_E₇}
 \lean{LineSystem.Hyp24.card_Δ_of_S_eq_E₈}
 \lean{LineSystem.isStarClosed_An}\lean{LineSystem.isStarClosed_Dn}
 \lean{LineSystem.isStarClosed_E₈}\lean{LineSystem.isStarClosed_E₇}\lean{LineSystem.isStarClosed_E₆}
 \lean{LineSystem.isStarClosed_E₈viaA₈}\lean{LineSystem.isStarClosed_E₈viaSteiner}
 \lean{LineSystem.isIndecomposable_An}\lean{LineSystem.isIndecomposable_Dn}
 \lean{LineSystem.isIndecomposable_E₈}\lean{LineSystem.isIndecomposable_E₇}
 \lean{LineSystem.isIndecomposable_E₆}\lean{LineSystem.isIndecomposable_E₈viaA₈}
 \lean{LineSystem.isIndecomposable_E₈viaSteiner}
 \lean{LineSystem.An.hyp24}\lean{LineSystem.Dn.hyp24}\lean{LineSystem.E₆.hyp24}
 \lean{LineSystem.E₇.hyp24}\lean{LineSystem.E₈.hyp24}
 \lean{LineSystem.An.orthGraphIso}\lean{LineSystem.Dn.orthGraphIso}
 \lean{LineSystem.E₆.orthGraphIso}\lean{LineSystem.E₇.orthGraphIso}\lean{LineSystem.E₈.orthGraphIso}
 \lean{LineSystem.An.card_Γa}\lean{LineSystem.Dn.card_Γa}\lean{LineSystem.E₆.card_Γa}
 \lean{LineSystem.E₇.card_Γa}\lean{LineSystem.E₈.card_Γa}
 \lean{LineSystem.E₆.card_Δ_div_two}\lean{LineSystem.E₇.card_Δ_div_two}
 \lean{LineSystem.E₈.card_Δ_div_two}\lean{LineSystem.E₆.orthGraph_isSRGWith}
 \lean{LineSystem.E₇.orthGraph_isSRGWith}\lean{LineSystem.E₈.orthGraph_isSRGWith}
 \uses{def:3.1,def:3.2,def:3.3,lem:sec3-integer,lem:sec3-spans,comp:C2,comp:C3,lem:sec3-transport,lem:sec3-connected-span,hyp:2.4,def:gamma-graph,def:2.1}
 \assume{We use the models in their specified ranges, and $X=\opk{Orth}(\Ga)$
 has the convention of Hypothesis~\ref{hyp:2.4}: distinct vertices are adjacent when
 their inner product is 0.}
 \prove{Every model is nonempty, star-closed, and indecomposable. The model
 $\typ A_1$ has one line and rank 1. For the marked star
 $\Za=\Ze_1-\Ze_2$, $\Zb=\Ze_2-\Ze_3$ the following data hold:
 \[
 \begin{array}{c|c|c|c|c}
 \text{model}&\text{rank}&\Ga\text{ or its graph}&|\Ga|&|\Delta|/2\\\hline
 \typ A_\ell\ (\ell\geq2)&\ell&\{\Ze_j-\Ze_3:4\leq j\leq\ell+1\}&\ell-2&\text{not needed}\\
 \typ D_\ell&\ell&\{\Ze_1+\Ze_2\}\cup\{-\Ze_3\pm\Ze_j:4\leq j\leq\ell\}&2\ell-5&\text{not needed}\\
 \typ E_6&6&\text{strongly regular }(9,4,1,2)&9&6\\
 \typ E_7&7&\text{strongly regular }(15,6,1,3)&15&15\\
 \typ E_8\text{ (ii)}&8&\text{strongly regular }(27,10,1,5)&27&36
 \end{array}
 \]
 The $\typ A$ graph is empty and the $\typ D$ graph is the friendship graph of $\ell-3$ triangles.
 The three exceptional rows hold for \emph{every} ordered pair $\Za,\Zb$ of roots
 with inner product $-1$ in the corresponding fixed model. The signed permutation
 $\sigma$ with $\sigma(\Ze_1)=\Ze_2$, $\sigma(\Ze_2)=-\Ze_1$, $\sigma(\Ze_3)=-\Ze_3$,
 fixing the other coordinate vectors, is an automorphism of $\typ D_\ell$, $\typ E_8$,
 $\typ E_7$ and $\typ E_6$ that carries $\Za,\Zb$ to the ordered star
 $\Za'=\Ze_1+\Ze_2$, $\Zb'=-\Ze_1+\Ze_3$. For this star, the graph $\opk{Orth}(\Ga)$ of
 $\typ E_6$ is the $3\times3$ lattice graph (adjacency in one common coordinate), and that
 of $\typ E_7$ is the disjointness graph on $\binom{[6]}2$, the complement of $T(6)$. Finally,
 the isometries of each exceptional model are transitive on its ordered stars: for
 $t\in\{6,7,8\}$ and $\Zv,\Zw\in\Phi(\typ E_t)$ with $(\Zv,\Zw)=-1$ there is a linear isometry
 $g$ of $\re^8$ with $g(\Phi(\typ E_t))=\Phi(\typ E_t)$, $g(\Za')=\Zv$ and $g(\Zb')=\Zw$.}
\end{namedresult}
\begin{proof}
 \pfsketch\ We will prove the classical assertions parametrically. Exceptional
 star data are finite tests on the complete integer root lists. Indecomposability
 has a short spanning-set proof, independent of any classification.
 \step{s3c4closed}{The five default models are star-closed.}
 \begin{proof}
  \step{s3c4Dclosed}{$\typ D_\ell$ is star-closed.}
  \begin{proof}
   \step{s3c4Dnorm}{\assume{$u,v$ are integral norm-2 roots with $u\cdot v=-1$.}
    \prove{$(u+v)\cdot(u+v)=2$.}}
   \begin{proof}\pf\ Bilinearity gives $2+2(-1)+2=2$.~\qed\end{proof}
   \qedstep
   \begin{proof}\pf\ An integral norm-2 vector has exactly two nonzero entries,
    each $\pm1$: this is \path{LineSystem.exists_dnRoot_of_inner_self}.
    Apply it to \stepref{s3c4Dnorm}.~\qed\end{proof}
  \end{proof}
  \step{s3c4Aclosed}{$\typ A_\ell$ is star-closed.}
  \begin{proof}\pf\ In \stepref{s3c4Dclosed} the sum also has coordinate sum zero,
   so its two signs are opposite. The Lean criterion is
   \path{LineSystem.exists_anRoot_of_inner_self_of_sum}.~\qed\end{proof}
  \step{s3c4Eclosed}{$\typ E_8$ is star-closed.}
  \begin{proof}\pf\ The existing certificate
   \leancite{Computations.C4.e₈_star_checked}{\texttt{Computations.\allowbreak C4.\allowbreak e\ensuremath{{}_8}\_\allowbreak star\_\allowbreak checked}} checks
   \path{Computations.checkStar} with dimension 8, target $-4$, the complete
   doubled root list, and its sound root decoder.
   Apply \path{LineSystem.IntegerModel.isStarClosed_of_checkStar}.~\qed\end{proof}
  \step{s3c4Sevenclosed}{$\typ E_7$ is star-closed.}
  \begin{proof}\pf\ Apply \path{LineSystem.IsStarClosed.orthogonalTo} to
   \stepref{s3c4Eclosed} and $\Zr_0$.~\qed\end{proof}
  \step{s3c4Sixclosed}{$\typ E_6$ is star-closed.}
  \begin{proof}\pf\ Apply the same named orthogonal-cut theorem to
   \stepref{s3c4Sevenclosed} and $\Zs_0$.~\qed\end{proof}
  \qedstep
  \begin{proof}\pf\ These five cases exhaust the default models.~\qed\end{proof}
 \end{proof}
 \step{s3c4indec}{The five default models are nonempty and indecomposable.}
 \begin{proof}
  \pflet{$B_m=\{\Ze_i-\Ze_m:1\leq i<m\}\cup\{\Ze_1+\Ze_m\}$ for $m\geq3$,
   and $\Zh=\frac12\mathbf1\in\re^8$.}
  \step{s3c4Aindec}{$\typ A_\ell$ has a nonempty connected spanning set.}
  \begin{proof}\pf\ The basis $\{\Ze_i-\Ze_{\ell+1}:1\leq i\leq\ell\}$ in
   Lemma~\ref{lem:sec3-spans} has product 1 between distinct members;
   for $\ell=1$ its graph has one vertex.~\qed\end{proof}
  \step{s3c4Bindec}{$B_m$ is a connected spanning set in $\typ D_m$.}
  \begin{proof}
   \step{s3c4Bconnect}{Its difference roots form a clique, and $\Ze_1+\Ze_m$
    has product $-1$ with $\Ze_2-\Ze_m$.}
   \begin{proof}\pf\ Compute the products from the coordinate basis, using $m\geq3$.~\qed\end{proof}
   \step{s3c4Bspan}{$B_m$ spans $\re^m$.}
   \begin{proof}\pf\ The difference of $\Ze_1+\Ze_m$ and $\Ze_1-\Ze_m$
    gives $2\Ze_m$; adding this recovered axis to each difference gives every other axis.~\qed\end{proof}
   \qedstep
   \begin{proof}\pf\ Use \stepref{s3c4Bconnect} and \stepref{s3c4Bspan}.~\qed\end{proof}
  \end{proof}
  \step{s3c4Eindec}{$B_8$ is a connected spanning subset of $\typ E_8$.}
  \begin{proof}\pf\ It lies in $\typ D_8\subseteq \typ E_8$ and has the properties in
   \stepref{s3c4Bindec}.~\qed\end{proof}
  \step{s3c4Sevenindec}{$B_6\cup\{\Ze_7+\Ze_8,\Zh\}$ is a connected spanning subset of $\typ E_7$.}
  \begin{proof}
   \step{s3c4Sevenjoin}{$\Zh$ has product 1 with both $\Ze_1+\Ze_6\in B_6$
    and $\Ze_7+\Ze_8$.}
   \begin{proof}\pf\ Each product is $1/2+1/2$.~\qed\end{proof}
   \qedstep
   \begin{proof}\pf\ Step \stepref{s3c4Bindec} spans the first six axes;
    adding $\Ze_7+\Ze_8$ spans the space in Lemma~\ref{lem:sec3-spans}.
    Step \stepref{s3c4Sevenjoin} connects the parts.~\qed\end{proof}
  \end{proof}
  \step{s3c4Sixindec}{$B_5\cup\{\Zh\}$ is a connected spanning subset of $\typ E_6$.}
  \begin{proof}\pf\ $B_5$ spans the first five axes by \stepref{s3c4Bindec};
   $\Zh$ has nonzero last three equal coordinates and has product 1 with
   $\Ze_1+\Ze_5$. Use the $\typ E_6$ span in Lemma~\ref{lem:sec3-spans}.~\qed\end{proof}
  \qedstep
  \begin{proof}\pf\ Apply Lemma~\ref{lem:sec3-connected-span} to each displayed nonempty set.~\qed\end{proof}
 \end{proof}
 \step{s3c4othermodels}{The other two $\typ E_8$ models are nonempty, star-closed, and indecomposable.}
 \begin{proof}\pf\ Transport \stepref{s3c4closed} and \stepref{s3c4indec}
  through C3, using Lemma~\ref{lem:sec3-transport}.~\qed\end{proof}
 \step{s3c4Agraph}{The marked $\typ A_\ell$ graph has the asserted vertices and no edges.}
 \begin{proof}
  \step{s3c4Avertices}{The equations $x_1=x_2$ and $x_2-x_3=1$ select precisely
   $\Ze_j-\Ze_3$ for $4\leq j\leq\ell+1$.}
  \begin{proof}\pf\ A difference root cannot have two equal nonzero entries at 1 and 2,
   so $x_1=x_2=0$ and $x_3=-1$.~\qed\end{proof}
  \qedstep
  \begin{proof}\pf\ Two distinct selected roots have product 1; their number is $\ell-2$.~\qed\end{proof}
 \end{proof}
 \step{s3c4Dgraph}{The marked $\typ D_\ell$ graph is the asserted friendship graph.}
 \begin{proof}
  \step{s3c4Dvertices}{The equations select $\Ze_1+\Ze_2$ and $-\Ze_3\pm\Ze_j$ for $j\geq4$.}
  \begin{proof}\pf\ Either $x_1=x_2\ne0$, forcing $x_1=x_2=1,x_3=0$;
   or $x_1=x_2=0$, forcing $x_3=-1$ and one further sign outside $[3]$.~\qed\end{proof}
  \step{s3c4Dcentre}{The vertex $\Ze_1+\Ze_2$ is adjacent to all the others.}
  \begin{proof}\pf\ Their supports are disjoint.~\qed\end{proof}
  \step{s3c4Dpairs}{Among the other vertices, adjacency pairs the two signs of the same $j$.}
  \begin{proof}\pf\ The product is 0 for those two signs, and 1 for different indices.~\qed\end{proof}
  \qedstep
  \begin{proof}\pf\ Steps \stepref{s3c4Dcentre} and \stepref{s3c4Dpairs} give
   $\ell-3$ triangles and $1+2(\ell-3)$ vertices.~\qed\end{proof}
 \end{proof}
 \step{s3c4exceptional}{\new\ Every ordered star in each default exceptional model
  has its stated row. We check all choices to avoid assuming transitivity.}
 \begin{proof}
  \pflet{$R$ be the sorted complete doubled integer root list of the model;
   for $a,b\in R$ with $a\cdot b=-4$, set
   $L(a,b)=\{v\in R:v\cdot a=0,\ v\cdot b=4\}$ and
   $D(a,b)=\{v\in R:v\cdot a=v\cdot b=0\}$.}
  A row test compares the two list lengths with $m,2\delta$; for each $u\in L$
  it counts $v\in L$ with $u\cdot v=0$ and compares with $d$; for each distinct
  $u,v\in L$ it counts $w\in L$ with $u\cdot w=v\cdot w=0$ and compares with
  $\lambda$ when $u\cdot v=0$, and with $\mu$ otherwise.
  \step{s3c4lists}{These lists represent exactly $\Ga$ and $\Delta$, respectively.}
  \begin{proof}\pf\ Divide raw dot products by $q=4$ using
   \path{LineSystem.inner_integerVector}; the third equation of $\Ga$ follows
   from $\Zc=-\Za-\Zb$. Completeness is the model's
   \texttt{mem\_rootList} theorem in \path{Computations/ModelTables.lean}.~\qed\end{proof}
  \step{s3c4census}{The complete finite tests give
  \[
  \begin{array}{c|r|r|r|r|r|r}
  R&\#(a,b)&|L|&|D|&d&\lambda&\mu\\\hline
  \typ E_6&1440&9&12&4&1&2\\
  \typ E_7&4032&15&30&6&1&3\\
  \typ E_8&13440&27&72&10&1&5.
  \end{array}
  \]}
  \begin{proof}\pf\ Exact integer evaluation gives \texttt{true} for
   \path{Computations.checkNoSolution} on $R\times R$, where a bad pair has dot
   product $-4$ and fails the specified row test. The lengths of the filtered
   ordered-pair lists are the three displayed counts. No real arithmetic is evaluated.~\qed\end{proof}
  \step{s3c4sound}{The row holds for every mathematical ordered star.}
  \begin{proof}\pf\ Apply \path{Computations.checkNoSolution_sound} to
   \stepref{s3c4census}, using the complete list identification in
   \stepref{s3c4lists}.~\qed\end{proof}
  \qedstep
  \begin{proof}\pf\ Adjacency is exactly a zero raw product; divide the length of each list $D(a,b)$
   by 2 to count the lines of $\Delta$. Thus \stepref{s3c4sound} gives the asserted parameters.~\qed\end{proof}
 \end{proof}
 \step{s3c4sigma}{\new\ The map $\sigma$ is an automorphism of $\typ D_\ell$, $\typ E_8$,
  $\typ E_7$ and $\typ E_6$, and $\sigma(\Za)=\Za'$, $\sigma(\Zb)=\Zb'$.}
 \begin{proof}
  \step{s3c4sigmaD}{$\sigma$ permutes $\Phi(\typ D_\ell)$ and the pair roots of model (ii).}
  \begin{proof}\pf\ A signed permutation of coordinates carries $s\Ze_i+t\Ze_j$ to a vector
   of the same form, in Definitions~\ref{def:3.1} and~\ref{def:3.2}.~\qed\end{proof}
  \step{s3c4sigmahalf}{$\sigma$ permutes the half roots of model (ii).}
  \begin{proof}\pf\ It permutes the entries of $\frac12\sum\epsilon(i)\Ze_i$ and changes
   two signs, so the product condition of Definition~\ref{def:3.2}(ii) is kept.~\qed\end{proof}
  \step{s3c4sigmacuts}{$\sigma$ permutes $\Phi(\typ E_7)$ and $\Phi(\typ E_6)$.}
  \begin{proof}\pf\ It fixes the coordinates 6, 7 and 8, so it keeps the equations of
   Definition~\ref{def:3.3}; use \stepref{s3c4sigmaD} and \stepref{s3c4sigmahalf}.~\qed\end{proof}
  \step{s3c4sigmastar}{$\sigma(\Za)=\Za'$ and $\sigma(\Zb)=\Zb'$.}
  \begin{proof}\pf\ $\sigma(\Ze_1-\Ze_2)=\Ze_2+\Ze_1$ and
   $\sigma(\Ze_2-\Ze_3)=-\Ze_1+\Ze_3$.~\qed\end{proof}
  \qedstep
  \begin{proof}\pf\ The map $\sigma$ is an orthogonal linear map; the four steps give the
   assertions.~\qed\end{proof}
 \end{proof}
 \step{s3c4labels}{For the ordered star $\Za'=\Ze_1+\Ze_2$,
  $\Zb'=-\Ze_1+\Ze_3$, the orthogonality graphs on $\Ga$ of $\typ E_6$ and $\typ E_7$
  are the lattice graph and the six-symbol disjointness graph.}
 \begin{proof}
  \pflet{$A_{\typ E_6}$ and $A_{\typ E_7}$ be the integer adjacency tables
   \texttt{e6Adjacency} and \texttt{e7Adjacency} of \path{Computations/C4E6Data.lean} and
   \path{Computations/C4E7Data.lean}, and let $A_{\rm lat}$ and $A_{\rm dis}$ be the
   tables of adjacency in one common coordinate on $\{0,1,2\}^2$ and of disjointness on
   $\binom{[6]}2$, in lexicographic order.}
  \step{s3c4Esixstar}{The roots $\Za',\Zb'$ are the marked star of the Lean term
   \leancite{LineSystem.E₆.hyp24}{\texttt{LineSystem.\allowbreak E\ensuremath{{}_6}.\allowbreak hyp24}}.}
  \begin{proof}\pf\ These are
   \leancite{LineSystem.E₆.hyp24_a}{\texttt{LineSystem.\allowbreak E\ensuremath{{}_6}.\allowbreak hyp24\_\allowbreak a}} and
   \leancite{LineSystem.E₆.hyp24_b}{\texttt{LineSystem.\allowbreak E\ensuremath{{}_6}.\allowbreak hyp24\_\allowbreak b}},
   in zero-based coordinates.~\qed\end{proof}
  \step{s3c4Esevenstar}{The roots $\Za',\Zb'$ are the marked star of the Lean term
   \leancite{LineSystem.E₇.hyp24}{\texttt{LineSystem.\allowbreak E\ensuremath{{}_7}.\allowbreak hyp24}}.}
  \begin{proof}\pf\ These are
   \leancite{LineSystem.E₇.hyp24_a}{\texttt{LineSystem.\allowbreak E\ensuremath{{}_7}.\allowbreak hyp24\_\allowbreak a}} and
   \leancite{LineSystem.E₇.hyp24_b}{\texttt{LineSystem.\allowbreak E\ensuremath{{}_7}.\allowbreak hyp24\_\allowbreak b}},
   in zero-based coordinates.~\qed\end{proof}
  \step{s3c4Esixtable}{The graph $\opk{Orth}(\Ga)$ of $\typ E_6$ for this star is
   isomorphic to the graph decoded from $A_{\typ E_6}$.}
  \begin{proof}\raggedright\pf\ This is
   \leancite{LineSystem.E₆.orthGraphIso}{\texttt{LineSystem.\allowbreak E\ensuremath{{}_6}.\allowbreak orthGraphIso}},
   by \stepref{s3c4Esixstar}. Its target
   \path{LineSystem.Hyp24.orthGraph} has adjacency at inner product 0, by
   \path{LineSystem.Hyp24.orthGraph_adj}.~\qed\end{proof}
  \step{s3c4Eseventable}{The graph $\opk{Orth}(\Ga)$ of $\typ E_7$ for this star is
   isomorphic to the graph decoded from $A_{\typ E_7}$.}
  \begin{proof}\raggedright\pf\ This is
   \leancite{LineSystem.E₇.orthGraphIso}{\texttt{LineSystem.\allowbreak E\ensuremath{{}_7}.\allowbreak orthGraphIso}},
   by \stepref{s3c4Esevenstar}, with the same adjacency
   \path{LineSystem.Hyp24.orthGraph_adj}.~\qed\end{proof}
  \step{s3c4grid}{The graph decoded from $A_{\typ E_6}$ is isomorphic to the lattice graph on $\{0,1,2\}^2$.}
  \begin{proof}\raggedright\pf\ Apply \path{Computations.nonempty_graphOfMatrix_iso} to
   \path{Computations.checkGraphIso} on $A_{\typ E_6}$, $A_{\rm lat}$ and a nine-entry permutation;
   in Lean it is the labelling \path{LineSystem.E₆.latticeLabel} of
   \path{Computations/C4Graphs.lean}.~\qed\end{proof}
  \step{s3c4duads}{The graph decoded from $A_{\typ E_7}$ is isomorphic to disjointness on $\binom{[6]}2$.}
  \begin{proof}\raggedright\pf\ Apply \path{Computations.nonempty_graphOfMatrix_iso} to
   \path{Computations.checkGraphIso} on $A_{\typ E_7}$, $A_{\rm dis}$ and a fifteen-entry permutation;
   in Lean it is the labelling \path{LineSystem.E₇.pairLabel} of the same file.~\qed\end{proof}
  \step{s3c4Esixlabel}{The $\typ E_6$ assertion holds.}
  \begin{proof}\pf\ Compose the isomorphisms of \stepref{s3c4Esixtable} and
   \stepref{s3c4grid}.~\qed\end{proof}
  \step{s3c4Esevenlabel}{The $\typ E_7$ assertion holds.}
  \begin{proof}\pf\ Compose the isomorphisms of \stepref{s3c4Eseventable} and
   \stepref{s3c4duads}.~\qed\end{proof}
  \qedstep
  \begin{proof}\pf\ Steps \stepref{s3c4Esixlabel} and \stepref{s3c4Esevenlabel} are the
   two assertions.~\qed\end{proof}
 \end{proof}
 \step{s3c4trans}{\new\ For $t\in\{6,7,8\}$ the isometries of $\re^8$ preserving
  $\Phi(\typ E_t)$ are transitive on the ordered stars of $\typ E_t$, as stated.}
 \begin{proof}\raggedright\pf\ This is
  \leancite{LineSystem.E₆.exists_isometry_star}{\texttt{LineSystem.\allowbreak E\ensuremath{{}_6}.\allowbreak exists\_\allowbreak isometry\_\allowbreak star}}
  and its analogues for $\typ E_7$ and $\typ E_8$, whose fixed star is $\Za',\Zb'$ by
  \leancite{LineSystem.E₈.hyp24_a}{\texttt{LineSystem.\allowbreak E\ensuremath{{}_8}.\allowbreak hyp24\_\allowbreak a}},
  \leancite{LineSystem.E₈.hyp24_b}{\texttt{LineSystem.\allowbreak E\ensuremath{{}_8}.\allowbreak hyp24\_\allowbreak b}}
  and their analogues; their trees of reflections are checked by
  \texttt{decide +kernel}.~\qed\end{proof}
 \step{s3c4ranks}{The ranks are $\ell$, $\ell$, 6, 7, 8, and $\typ A_1$ has one line and rank 1.}
 \begin{proof}\pf\ Read the ranks from Lemma~\ref{lem:sec3-spans}; the line of
  $\typ A_1$ is $\Ze_1-\Ze_2$ by Definition~\ref{def:3.1}.~\qed\end{proof}
 \qedstep
 \begin{proof}\pf\ Combine \stepref{s3c4closed}, \stepref{s3c4indec},
  \stepref{s3c4othermodels}, \stepref{s3c4Agraph}, \stepref{s3c4Dgraph},
  \stepref{s3c4exceptional}, \stepref{s3c4sigma}, \stepref{s3c4labels}, \stepref{s3c4trans},
  and \stepref{s3c4ranks}.
  None uses Theorem 2.13 or 3.5.~\qed\end{proof}
\end{proof}
\begingroup\raggedright
The two classical closure declarations are in \path{LineSystem/ClassicalModels.lean};
the three exceptional closure declarations are in \path{Computations/E8.lean}.
Computation C4 is proved in Lean, for the marked star and for every ordered star of the
exceptional models. The indecomposability theorems are in \path{Computations/C4.lean},
\path{Computations/C4Classical.lean} and \path{Computations/C4E6.lean} to
\path{Computations/C4E8.lean}; the marked stars \texttt{hyp24}, the isomorphisms
\texttt{orthGraphIso}, the counts \texttt{card\_\ensuremath{\Gamma}a} and \texttt{card\_\ensuremath{\Delta}\_div\_two}, and the
strongly regular parameters \texttt{orthGraph\_isSRGWith} are in
\path{Computations/C4An.lean}, \path{Computations/C4Dn.lean} and
\path{Computations/C4E6.lean} to \path{Computations/C4E8.lean}. The Lean star of $\typ A_\ell$
and of $\typ D_\ell$ is our $\Za,\Zb$ in zero-based coordinates, $\Ze_0-\Ze_1$, $\Ze_1-\Ze_2$.
The Lean star \texttt{hyp24} of $\typ E_6$, $\typ E_7$ and $\typ E_8$ is $\Za',\Zb'$, the image of
our star under $\sigma$; our star itself is \texttt{hyp24Simple}, and
\leancite{Computations.C4.marked_star_data}{\texttt{Computations.\allowbreak C4.\allowbreak marked\_\allowbreak star\_\allowbreak data}} records the table above for our star
in all five rows. For every ordered star, a reflection in a root of a star-closed system is an
automorphism, \leancite{LineSystem.rootReflection_mem_isometryGroup}{\texttt{LineSystem.\allowbreak rootReflection\_\allowbreak mem\_\allowbreak isometryGroup}} in
\path{LineSystem/Reflections.lean}, and two reflection trees for each exceptional model, checked
by \texttt{decide +kernel} in \path{Computations/C4E6Stars.lean} to
\path{Computations/C4E8Stars.lean}, give transitivity on ordered stars:
\leancite{LineSystem.E₆.exists_isometry_star}{\texttt{LineSystem.\allowbreak E\ensuremath{{}_6}.\allowbreak exists\_\allowbreak isometry\_\allowbreak star}}, \leancite{LineSystem.E₇.exists_isometry_star}{\texttt{LineSystem.\allowbreak E\ensuremath{{}_7}.\allowbreak exists\_\allowbreak isometry\_\allowbreak star}} and
\leancite{LineSystem.E₈.exists_isometry_star}{\texttt{LineSystem.\allowbreak E\ensuremath{{}_8}.\allowbreak exists\_\allowbreak isometry\_\allowbreak star}}. For every \texttt{H : Hyp24}
with \texttt{H.S} one of the three models, \leancite{LineSystem.Hyp24.orthGraph_isSRGWith_of_S_eq_E₆}{\texttt{LineSystem.\allowbreak Hyp24.\allowbreak orthGraph\_\allowbreak isSRGWith\_\allowbreak of\_\allowbreak S\_\allowbreak eq\_\allowbreak E\ensuremath{{}_6}}},
\leancite{LineSystem.Hyp24.card_Γa_of_S_eq_E₆}{\texttt{LineSystem.\allowbreak Hyp24.\allowbreak card\_\allowbreak \ensuremath{\Gamma}a\_\allowbreak of\_\allowbreak S\_\allowbreak eq\_\allowbreak E\ensuremath{{}_6}}} and
\leancite{LineSystem.Hyp24.card_Δ_of_S_eq_E₆}{\texttt{LineSystem.\allowbreak Hyp24.\allowbreak card\_\allowbreak \ensuremath{\Delta}\_\allowbreak of\_\allowbreak S\_\allowbreak eq\_\allowbreak E\ensuremath{{}_6}}}, with their $\typ E_7$ and $\typ E_8$
analogues, give the exceptional rows; Lean counts $\Delta$ as $12,30,72$ roots, twice the
entries of the last column. The bundle \leancite{Computations.C4.exceptional_star_data}{\texttt{Computations.\allowbreak C4.\allowbreak exceptional\_\allowbreak star\_\allowbreak data}}
collects them. The graph labels are the isomorphisms
\leancite{LineSystem.E₆.latticeGraphIso}{\texttt{LineSystem.\allowbreak E\ensuremath{{}_6}.\allowbreak latticeGraphIso}}, from the line graph of $K_{3,3}$, and
\leancite{LineSystem.E₇.triangularGraphComplIso}{\texttt{LineSystem.\allowbreak E\ensuremath{{}_7}.\allowbreak triangularGraphComplIso}}, from the complement of the line
graph of $K_6$, both for the star $\Za',\Zb'$, in \path{Computations/C4Graphs.lean}; their every-star
forms are \leancite{LineSystem.Hyp24.nonempty_latticeGraph_iso_of_S_eq_E₆}{\texttt{LineSystem.\allowbreak Hyp24.\allowbreak nonempty\_\allowbreak latticeGraph\_\allowbreak iso\_\allowbreak of\_\allowbreak S\_\allowbreak eq\_\allowbreak E\ensuremath{{}_6}}} and
\leancite{LineSystem.Hyp24.nonempty_triangularGraph_compl_iso_of_S_eq_E₇}{\texttt{LineSystem.\allowbreak Hyp24.\allowbreak nonempty\_\allowbreak triangularGraph\_\allowbreak compl\_\allowbreak iso\_\allowbreak of\_\allowbreak S\_\allowbreak eq\_\allowbreak E\ensuremath{{}_7}}}. In Lean the
transitivity replaces the transfer by $\sigma$. All these isomorphisms are for the orthogonality
graph; the bridge \path{LineSystem.IntegerModel.adjMatrix_eq_dotAdjacency} builds the graph at
inner product 1, and we do not use it. The script \path{blueprint/checks/sec3/check_sec3.py}
recomputes the finite data; no Lean proof evaluates a spectrum.\par
\endgroup
The source lists have the membership theorems
\leancite{LineSystem.E₆.mem_rootList}{\texttt{LineSystem.\allowbreak E\ensuremath{{}_6}.\allowbreak mem\_\allowbreak rootList}}, \leancite{LineSystem.E₇.mem_rootList}{\texttt{LineSystem.\allowbreak E\ensuremath{{}_7}.\allowbreak mem\_\allowbreak rootList}}, and
\leancite{LineSystem.E₈viaD₈.mem_rootList}{\texttt{LineSystem.\allowbreak E\ensuremath{{}_8}viaD\ensuremath{{}_8}.\allowbreak mem\_\allowbreak rootList}}. Sorting does not change membership;
the lists $R,L,D$ in C4 contain each raw vector once.

\begin{namedresult}{Lemma}{(independence of the perpendicular choices)}\label{lem:sec3-choices}
 \lean{LineSystem.E₈.nonempty_iso_E₇_orthogonalTo}
 \lean{LineSystem.E₈.nonempty_iso_E₆_orthogonalTo}
 \lean{LineSystem.E₈.exists_isometry_of_ordered_star}
 \lean{LineSystem.E₈.exists_isometry_of_root}
 \lean{LineSystem.E₈viaA₈.exists_isometry_of_ordered_stars}
 \lean{LineSystem.E₈viaA₈.exists_isometry_of_roots}
 \lean{LineSystem.E₈viaA₈.nonempty_iso_E₇_orthogonalTo}
 \lean{LineSystem.E₈viaA₈.nonempty_iso_E₆_orthogonalTo}
 \lean{LineSystem.E₈viaSteiner.exists_isometry_of_ordered_stars}
 \lean{LineSystem.E₈viaSteiner.exists_isometry_of_ordered_star}
 \lean{LineSystem.E₈viaSteiner.exists_isometry_of_roots}
 \lean{LineSystem.E₈viaSteiner.exists_isometry_of_root}
 \lean{LineSystem.E₈viaSteiner.nonempty_iso_E₇_orthogonalTo}
 \lean{LineSystem.E₈viaSteiner.nonempty_iso_E₆_orthogonalTo}
 \lean{LineSystem.T₇.nonempty_iso_E₇}\lean{LineSystem.T₆.nonempty_iso_E₆}
 \uses{def:3.2,def:3.3,comp:C2,comp:C3,comp:C4,lem:2.10,lem:2.12,def:indecomposable,facts:inner}
 \assume{$\typ E_8$ is the default model in $\re^8$; $\Zr\in\Phi(\typ E_8)$ and,
 for the star assertion, $\Zs\in\Phi(\typ E_8)$ with $(\Zr,\Zs)=-1$.}
 \prove{There is an automorphism of $\typ E_8$, induced by a linear isometry of
 $\re^8$, carrying $(\Zr,\Zs)$ to $(\Zr_0,\Zs_0)$.
 Every root $\Zr$ also admits such an automorphism carrying it to $\Zr_0$,
 without an initially chosen partner. Thus the cuts perpendicular to any line
 or any star are isometric to $\typ E_7$ or $\typ E_6$, respectively.
 The same assertions hold in the other two models, and $T_7\cong \typ E_7$, $T_6\cong \typ E_6$.}
\end{namedresult}
\begin{proof}
 \pfsketch\ \repair\ We fill Gap 5 before classification. C4 and Lemma 2.10
 identify the marked graphs; Lemma 2.12 extends their isomorphism to the ordered stars.
 \step{s3choicepair}{Any two ordered stars of $\typ E_8$ are carried to each other by an automorphism.}
 \begin{proof}
  \step{s3choicegraphs}{Their graphs on $\Ga$ are isomorphic.}
  \begin{proof}\pf\ C4 gives order 27, degree 10, and star-closure and
   indecomposability for each instance. Apply Lemma~\ref{lem:2.10}; the positive
   degree excludes an empty graph and degree $10<26$ excludes a universal vertex.~\qed\end{proof}
  \qedstep
  \begin{proof}\pf\ Lemma~\ref{lem:2.12} extends this graph isomorphism to a
   surjective span isometry carrying the ordered stars to each other.~\qed\end{proof}
 \end{proof}
 \step{s3choiceroot}{Every root has a partner of inner product $-1$.}
 \begin{proof}
  \step{s3choicenonorth}{Some root outside $\{\Zr,-\Zr\}$ has nonzero product with $\Zr$.}
  \begin{proof}\pf\ Otherwise this antipodal pair and its nonempty complement
   would be an orthogonal partition. This contradicts C4; the complement is
   nonempty since C2 gives 120 lines.~\qed\end{proof}
  \qedstep
  \begin{proof}\pf\ The product in \stepref{s3choicenonorth} is $1$ or $-1$
   by Section~\ref{facts:inner}; negate the partner when the product is 1.~\qed\end{proof}
 \end{proof}
 \step{s3choicecuts}{The default-model perpendicular cuts have the asserted isomorphisms.}
 \begin{proof}
  \step{s3choiceSeven}{The cut perpendicular to $\Zr$ maps onto $\typ E_7$.}
  \begin{proof}\pf\ Use \stepref{s3choiceroot} and \stepref{s3choicepair} to send
   $\Zr$ to $\Zr_0$; an isometry preserves the equation $(\Zx,\Zr)=0$.~\qed\end{proof}
  \step{s3choiceSix}{The cut perpendicular to the ordered star maps onto $\typ E_6$.}
  \begin{proof}\pf\ Step \stepref{s3choicepair} preserves the two equations of
   orthogonality to $\Zr,\Zs$; orthogonality to $-\Zr-\Zs$ is their sum.~\qed\end{proof}
  \qedstep
  \begin{proof}\pf\ These are the two cuts in Definition~\ref{def:3.3}.~\qed\end{proof}
 \end{proof}
 \step{s3choiceSteiner}{In model (iii), $T_7$ and $T_6$ are cuts of the required kinds.}
 \lean{LineSystem.T₇_eq_orthogonalTo}\lean{LineSystem.T₆_eq_orthogonalTo}
 \lean{LineSystem.inner_t₇Normal_t₆Normal}
 \begin{proof}\pf\ Take $\Za=\sqrt2\Ze_8$ and
  $\Zb=-(\Ze_5+\Ze_6+\Ze_7+\Ze_8)/\sqrt2$.
  Here $\Za$ is an axis root and $\Zb$ is a root by block 5678, and their product is
  $-1$. Their two perpendicular
  equations are exactly those defining $T_7,T_6$.~\qed\end{proof}
 \qedstep
 \begin{proof}\pf\ C3 carries any root or ordered star in either other model to
  the default model. Apply \stepref{s3choicecuts}, including the cuts identified
  by \stepref{s3choiceSteiner}.~\qed\end{proof}
\end{proof}

\begin{cgss}{Lemma}{3.4}\label{lem:3.4}
 \lean{LineSystem.numLines_An}\lean{LineSystem.numLines_Dn}
 \lean{LineSystem.numLines_E₈}\lean{LineSystem.numLines_E₈viaA₈}
 \lean{LineSystem.numLines_E₈viaSteiner}\lean{LineSystem.numLines_E₇}\lean{LineSystem.numLines_E₆}
 \uses{def:3.1,def:3.2,def:3.3,comp:C2,lem:sec3-choices}
 \assume{The model parameters are $\ell\geq1$ for $\typ A_\ell$ and $\ell\geq4$ for
 $\typ D_\ell$; the exceptional systems have the definitions above.}
 \prove{\repair\ With $|\typ A_\ell|$ in place of the printed $|A|$ (Error 13),
 \[
 |\typ A_\ell|=\tfrac12\ell(\ell+1),\quad |\typ D_\ell|=\ell(\ell-1),\quad
 |\typ E_8|=120,\quad |\typ E_7|=63,\quad |\typ E_6|=36.
 \]
 These are counts of lines, independent of the model and of the perpendicular choices.}
\end{cgss}
\begin{proof}
 \pfsketch\ We will use the signed counts of C2, and then justify the paper's
 use of model (iii) for the perpendicular systems.
 \step{s334classical}{The first two formulas hold.}
 \begin{proof}\pf\ Apply the first two counts of C2 in the stated parameter ranges.~\qed\end{proof}
 \step{s334eight}{All three definitions of $\typ E_8$ have 120 lines.}
 \begin{proof}\pf\ C2 gives $36+84=56+64=8+8\cdot14=120$.~\qed\end{proof}
 \step{s334seven}{Every perpendicular-line system has 63 lines.}
 \begin{proof}\pf\ Lemma~\ref{lem:sec3-choices} identifies it with $T_7$;
  C2 gives $7+8\cdot7=63$.~\qed\end{proof}
 \step{s334six}{Every perpendicular-star system has 36 lines.}
 \begin{proof}\pf\ Lemma~\ref{lem:sec3-choices} identifies it with $T_6$;
  C2 gives $4+8+4\cdot6=36$.~\qed\end{proof}
 \qedstep
 \begin{proof}\pf\ These four steps cover all five types and all allowed choices.
  The seven named count declarations already exist in \path{LineSystem/Models.lean}.~\qed\end{proof}
\end{proof}

\begin{cgss}{Theorem}{3.5}\label{thm:3.5}
 \lean{LineSystem.classification}\lean{LineSystem.RootType}
 \lean{LineSystem.RootType.model}\lean{LineSystem.RootType.rank}
 \uses{thm:2.13,def:line-system,def:3.1,def:3.2,def:3.3,comp:C2,comp:C3,comp:C4,lem:3.4,lem:sec3-spans,lem:sec3-transport,lem:sec3-restrict}
 \assume{\begin{pfenum}
  \item $V$ is a finite-dimensional real inner product space.
  \item $S$ is a nonempty line system in $V$, and $W=\operatorname{span}\Phi(S)$.
 \end{pfenum}}
 \prove{$S$ is indecomposable and star-closed if and only if it is isomorphic to
 exactly one of
 $\typ A_\ell$ ($\ell\geq1$), $\typ D_\ell$ ($\ell\geq4$), $\typ E_6$, $\typ E_7$, $\typ E_8$.
 The isomorphism is induced by a surjective linear isometry of spans.
 All three definitions of $\typ E_8$ give that same type.}
 \new\ In Lean the list is an index type \texttt{LineSystem.RootType}, with
 constructors for $\typ A_\ell$ with $1\leq\ell$ and $\typ D_\ell$ with $4\leq\ell$, and
 three constants for $\typ E_6$, $\typ E_7$, $\typ E_8$. Each type $t$ has an ambient
 dimension \texttt{t.dim}, a model \texttt{t.model}, a line system in the Euclidean space
 of that dimension, and a rank \texttt{t.rank}. The Lean statement then says that $S$ is
 indecomposable and star-closed if and only if there is a unique $t$ with an isomorphism
 from $S$ to \texttt{t.model}. Theorem 3.6 uses the same type, with $\typ A_7$, $\typ A_8$
 and $\typ D_8$ excluded.
\end{cgss}
\begin{proof}
 \pfsketch\ The models realize the five possibilities of Theorem 2.13. Its inputs
 from this section are C2, C4, and the span calculations, which have all been
 established without either classification theorem. We will include the one-line case.
 \new\ The uniqueness assertion of Theorem 2.13 is stated for systems spanning
 $\re^\ell$. We will restrict $S$ and the realizing model to their spans, map both spans
 to $\re^\ell$ through orthonormal bases, apply uniqueness there, and compose.
 \step{s335forward}{\assume{$S$ is indecomposable and star-closed.}
  \prove{There are a listed model $Y$ and a surjective linear isometry
  $F:W\to\operatorname{span}\Phi(Y)$ with $F(\Phi(S))=\Phi(Y)$.}}
 \begin{proof}
  \pflet{$\ell=\dim W$, and let $S|_W$ be the restriction of $S$ to $W$ in
   Lemma~\ref{lem:sec3-restrict}.}
  \step{s335restrict}{$S|_W$ satisfies the hypotheses of Theorem~\ref{thm:2.13},
   with $W$ as the ambient space of dimension $\ell$.}
  \begin{proof}\pf\ Lemma~\ref{lem:sec3-restrict} makes $S|_W$ a line system spanning
   $W$, nonempty by \asm[2]{0} and indecomposable and star-closed by \asm{1}; $W$ is
   finite-dimensional by \asm[1]{0}.~\qed\end{proof}
  \step{s335possibilities}{$S|_W$ lies in one of the five cases of Theorem 2.13.}
  \begin{proof}\pf\ Apply Theorem~\ref{thm:2.13} to \stepref{s335restrict}.~\qed\end{proof}
  \step{s335realized}{A listed model $Y$ of rank $\ell$ has the number of lines of that
   case and, when that case has more than one line, an ordered star with the data on
   $\Ga$ and $\Delta$ of that case.}
  \begin{proof}
   \step{s335A}{\case{Case 1 of Theorem 2.13.}}
   \begin{proof}\pf\ Take $Y=\typ A_\ell$: C4 gives rank $\ell$, the empty graph on
    $\ell-2$ vertices for $\ell\geq2$, and one line for $\ell=1$; Lemma~\ref{lem:3.4}
    gives $\ell(\ell+1)/2$ lines.~\qed\end{proof}
   \step{s335D}{\case{Case 2 of Theorem 2.13, so $\ell\geq4$.}}
   \begin{proof}\pf\ Take $Y=\typ D_\ell$: C4 gives rank $\ell$ and the friendship graph
    of $\ell-3$ triangles; Lemma~\ref{lem:3.4} gives $\ell(\ell-1)$ lines.~\qed\end{proof}
   \step{s335Eeight}{\case{Case 3 of Theorem 2.13, so $\ell=8$.}}
   \begin{proof}\pf\ Take $Y=\typ E_8$: C4 gives rank 8, $|\Ga|=27$ and
    $|\Delta|/2=36$; Lemma~\ref{lem:3.4} gives 120 lines.~\qed\end{proof}
   \step{s335Eseven}{\case{Case 4 of Theorem 2.13, so $\ell=7$.}}
   \begin{proof}\pf\ Take $Y=\typ E_7$: C4 gives rank 7, $|\Ga|=15$ and
    $|\Delta|/2=15$; Lemma~\ref{lem:3.4} gives 63 lines.~\qed\end{proof}
   \step{s335Esix}{\case{Case 5 of Theorem 2.13, so $\ell=6$.}}
   \begin{proof}\pf\ Take $Y=\typ E_6$: C4 gives rank 6, $|\Ga|=9$ and
    $|\Delta|/2=6$; Lemma~\ref{lem:3.4} gives 36 lines.~\qed\end{proof}
   \qedstep
   \begin{proof}\pf\ These are the five cases in \stepref{s335possibilities}.~\qed\end{proof}
  \end{proof}
  \pflet{$W_Y=\operatorname{span}\Phi(Y)$, and let $Y|_{W_Y}$ be the restriction of $Y$
   to $W_Y$ in Lemma~\ref{lem:sec3-restrict}.}
  \step{s335modelhyp}{$Y|_{W_Y}$ satisfies the hypotheses of
   Theorem~\ref{thm:2.13}, with $W_Y$ as the ambient space of dimension $\ell$.}
  \begin{proof}
   \step{s335modelprops}{$Y$ is nonempty, indecomposable and star-closed.}
   \begin{proof}\pf\ This is the first assertion of C4.~\qed\end{proof}
   \step{s335modeldim}{$\dim W_Y=\ell$.}
   \begin{proof}\pf\ The rank of $Y$ is $\ell$ by \stepref{s335realized}.~\qed\end{proof}
   \qedstep
   \begin{proof}\pf\ Lemma~\ref{lem:sec3-restrict} makes $Y|_{W_Y}$ a line system
    spanning $W_Y$ with the properties of \stepref{s335modelprops}; its dimension is
    \stepref{s335modeldim}.~\qed\end{proof}
  \end{proof}
  \step{s335samecase}{$Y|_{W_Y}$ lies in the same case of Theorem 2.13 as $S|_W$.}
  \begin{proof}\pf\ Lemma~\ref{lem:sec3-restrict} carries the number of lines and the
   marked-star data of \stepref{s335realized} from $Y$ to $Y|_{W_Y}$, whose dimension is
   $\ell$ by \stepref{s335modelhyp}; these are the data of the case of
   \stepref{s335possibilities}.~\qed\end{proof}
  \step{s335coords}{There are surjective linear isometries $\iota:W\to\re^\ell$
   and $\iota_Y:W_Y\to\re^\ell$; the sets $\iota(\Phi(S))$ and $\iota_Y(\Phi(Y))$ are the
   root sets of line systems $S'$ and $Y'$ in $\re^\ell$, which satisfy the hypotheses of
   Theorem~\ref{thm:2.13} and lie in the case of \stepref{s335possibilities}.}
  \begin{proof}
   \step{s335basis}{The isometries $\iota$ and $\iota_Y$ exist.}
   \begin{proof}\pf\ Take the representation \path{OrthonormalBasis.repr} of
    \path{stdOrthonormalBasis} of $W$ and of $W_Y$, reindexed to $\operatorname{Fin}\ell$
    by \path{OrthonormalBasis.reindex} along \path{finCongr}, using $\dim W=\ell$ and
    \stepref{s335modelhyp}.~\qed\end{proof}
   \step{s335imagelines}{$S'$ and $Y'$ are line systems in $\re^\ell$, and $\iota$ and
    $\iota_Y$ restrict to bijections of root sets that preserve inner products.}
   \begin{proof}\pf\ A linear isometry is injective, commutes with negation and
    preserves inner products, by \path{LinearIsometry.inner_map_map}; the conditions of
    the line-system definition in Section~\ref{def:line-system} are of these kinds.~\qed\end{proof}
   \step{s335imagespan}{$\Phi(S')$ and $\Phi(Y')$ span $\re^\ell$.}
   \begin{proof}\pf\ By \path{Submodule.span_image}, the span of an image is the image
    of the span, which is the whole of $\re^\ell$ for a surjective map; use
    \stepref{s335restrict} and \stepref{s335modelhyp}.~\qed\end{proof}
   \qedstep
   \begin{proof}\pf\ Lemma~\ref{lem:sec3-transport}, applied to the bijections of
    \stepref{s335imagelines}, carries nonemptiness, indecomposability, star-closure, the
    number of lines and the marked-star data from $S|_W$ and $Y|_{W_Y}$, which have them
    by \stepref{s335restrict}, \stepref{s335possibilities}, \stepref{s335modelhyp} and
    \stepref{s335samecase}; \stepref{s335imagespan} gives spanning.~\qed\end{proof}
  \end{proof}
  \step{s335orthogonal}{There is an orthogonal transformation $O$ of $\re^\ell$ with
   $O(\Phi(S'))=\Phi(Y')$.}
  \begin{proof}\pf\ Apply the uniqueness assertion of Theorem~\ref{thm:2.13}, for the
   ambient space $\re^\ell$, to \stepref{s335coords}.~\qed\end{proof}
  \qedstep
  \begin{proof}\pf\ Put $F=\iota_Y^{-1}\circ O\circ\iota$, a surjective linear isometry
   $W\to W_Y$ by \path{LinearIsometryEquiv.symm} and \path{LinearIsometryEquiv.trans};
   \stepref{s335coords} and \stepref{s335orthogonal} give
   $F(\Phi(S))=\Phi(Y)$.~\qed\end{proof}
 \end{proof}
 \step{s335reverse}{\assume{$S$ is isomorphic to a listed model.}
  \prove{$S$ is indecomposable and star-closed.}}
 \begin{proof}\pf\ C4 gives these properties for the model; apply
  Lemma~\ref{lem:sec3-transport}.~\qed\end{proof}
 \step{s335unique}{Different listed types are not isomorphic.}
 \begin{proof}
  \step{s335rank}{Isomorphic models have the same rank.}
  \begin{proof}\pf\ Apply Lemma~\ref{lem:sec3-transport}.~\qed\end{proof}
  \step{s335AD}{At a common rank $\ell\geq4$, the $\typ D$ count exceeds the $\typ A$ count by $\ell(\ell-3)/2>0$.}
  \begin{proof}\pf\ Subtract the formulas in Lemma~\ref{lem:3.4}.~\qed\end{proof}
  \step{s335exceptionalcounts}{At ranks $6,7,8$ the respective $(\typ A,\typ D,\typ E)$ counts are
   $(21,30,36)$, $(28,42,63)$, $(36,56,120)$, with strict inequalities in each row.}
  \begin{proof}\pf\ Substitute those ranks in Lemma~\ref{lem:3.4}.~\qed\end{proof}
  \qedstep
  \begin{proof}\pf\ Line count is an isomorphism invariant by
   Lemma~\ref{lem:sec3-transport}; \stepref{s335rank}--\stepref{s335exceptionalcounts}
   distinguish every pair of listed types.~\qed\end{proof}
 \end{proof}
 \step{s335atmost}{The system $S$ cannot be isomorphic to two different listed types.}
 \begin{proof}\pf\ Compose one span isometry with the inverse of the other from
  Lemma~\ref{lem:sec3-transport}. The resulting isometry between the two models
  contradicts \stepref{s335unique}.~\qed\end{proof}
 \qedstep
 \begin{proof}\pf\ Steps \stepref{s335forward}, \stepref{s335reverse}, and
  \stepref{s335atmost} prove the equivalence and uniqueness. C3 identifies the
  three $\typ E_8$ models independently.~\qed\end{proof}
\end{proof}

\begin{namedresult}{Lemma}{(a maximal extension in the span)}\label{lem:sec3-extension}
 \lean{LineSystem.exists_maximal_extension_span}
 \uses{def:line-system,def:indecomposable,def:2.2,lem:2.3,facts:inner}
 \assume{$S$ is a nonempty indecomposable line system in a finite-dimensional
 real inner product space, $W=\operatorname{span}\Phi(S)$, and $S$ is not maximal in $W$.}
 \prove{There is a line system $T$ in $W$ with
 $\Phi(S)\subsetneq\Phi(T)$, spanning $W$, that is maximal in $W$,
 indecomposable, and star-closed.}
\end{namedresult}
\begin{proof}
 \pfsketch\ \repair\ For the second omission in Gap 13 we will bound all candidate
 roots by their finite signatures against $\Phi(S)$, maximize the size of an extension,
 and prove indecomposability using the equality of spans.
 \pflet{$K=\{\Zx\in W:(\Zx,\Zx)=2,\ (\Zx,\Zr)\in\{-2,-1,0,1,2\}
 \text{ for every }\Zr\in\Phi(S)\}$.}
 \step{s3extfinite}{$K$ is finite.}
 \begin{proof}
  \step{s3extsignature}{The signature map $\Zx\mapsto((\Zx,\Zr))_{\Zr\in\Phi(S)}$ is injective on $K$.}
  \begin{proof}
   \step{s3extdifference}{Equal signatures make $\Zx-\Zy$ orthogonal to $W$.}
   \begin{proof}\pf\ Use bilinearity and the finite-span lemma in Section~\ref{facts:inner}.~\qed\end{proof}
   \qedstep
   \begin{proof}\pf\ Since $\Zx-\Zy\in W$, \stepref{s3extdifference} gives zero
    squared norm, so $\Zx=\Zy$ by Section~\ref{facts:inner}.~\qed\end{proof}
  \end{proof}
  \qedstep
  \begin{proof}\pf\ The signatures lie in a set of $5^{|\Phi(S)|}$ functions.
   Apply \path{Set.Finite.of_injOn} to \stepref{s3extsignature}.~\qed\end{proof}
 \end{proof}
 \step{s3extchoose}{Choose a largest-cardinality subset $\Psi\subseteq K$ containing
  $\Phi(S)$ that is the signed root set of a line system.}
 \begin{proof}\pf\ The eligible subsets form a nonempty finite family: it contains
  $\Phi(S)$ and is contained in the powerset of the finite $K$. Apply
  \path{Finset.exists_max_image} to their cardinalities.~\qed\end{proof}
 \pflet{$T$ be this line system with root set $\Psi$.}
 \step{s3extmaximal}{$T$ is maximal within $W$.}
 \begin{proof}
  \step{s3extcandidate}{\assume{$\Zx\in W$ has squared norm 2 and allowed products with $\Psi$.}
   \prove{$\Zx\in\Psi$.}}
  \begin{proof}
   \step{s3extinK}{$\Zx,-\Zx\in K$.}
   \begin{proof}\pf\ The candidate products include those with $\Phi(S)\subseteq\Psi$;
    the allowed product set is closed under negation.~\qed\end{proof}
   \step{s3extadjoin}{$\Psi\cup\{\Zx,-\Zx\}$ is eligible in \stepref{s3extchoose}.}
   \begin{proof}\pf\ Old-new products are allowed by \asm{2}, and the new-new
    products are $\pm2$; use \stepref{s3extinK} for containment in $K$.~\qed\end{proof}
   \qedstep
   \begin{proof}\pf\ Maximal cardinality in \stepref{s3extchoose} forbids a proper
    enlargement by \stepref{s3extadjoin}.~\qed\end{proof}
  \end{proof}
  \qedstep
  \begin{proof}\pf\ Step \stepref{s3extcandidate} is Definition~\ref{def:2.2}.~\qed\end{proof}
 \end{proof}
 \step{s3extspan}{$\operatorname{span}\Psi=W$.}
 \begin{proof}\pf\ $\Phi(S)\subseteq\Psi\subseteq W$ and $\Phi(S)$ spans $W$.~\qed\end{proof}
 \step{s3extindec}{$T$ is indecomposable.}
 \begin{proof}
  \step{s3extpartition}{\assume{$P\cup Q=\Psi$ and $P\perp Q$.}
   \prove{$P=\varnothing$ or $Q=\varnothing$.}}
  \begin{proof}
   \step{s3extoldpart}{All roots of $S$ lie in one of $P,Q$.}
   \begin{proof}\pf\ Intersect the partition with $\Phi(S)$ and use its indecomposability.~\qed\end{proof}
   \step{s3extempty}{The other part is empty.}
   \begin{proof}\pf\ Each root there is orthogonal to the spanning $\Phi(S)$,
    thus to itself by Section~\ref{facts:inner}; its norm would be zero.~\qed\end{proof}
   \qedstep
   \begin{proof}\pf\ Use \stepref{s3extempty}.~\qed\end{proof}
  \end{proof}
  \qedstep
  \begin{proof}\pf\ Apply the indecomposability definition in Section~\ref{def:indecomposable} to \stepref{s3extpartition}.~\qed\end{proof}
 \end{proof}
 \step{s3extproper}{$\Phi(S)\subsetneq\Psi$.}
 \begin{proof}\pf\ Equality would make $S$ maximal by \stepref{s3extmaximal},
  contrary to the hypothesis.~\qed\end{proof}
 \step{s3extclosed}{$T$ is star-closed.}
 \begin{proof}\pf\ Lemma~\ref{lem:2.3} applies to \stepref{s3extmaximal}.~\qed\end{proof}
 \qedstep
 \begin{proof}\pf\ Combine \stepref{s3extspan}, \stepref{s3extindec},
  \stepref{s3extproper}, and \stepref{s3extclosed}.~\qed\end{proof}
\end{proof}
The finiteness and maximum declarations used here are in
\path{Mathlib/Data/Set/Finite/Basic.lean} and \path{Mathlib/Data/Finset/Max.lean}.

\begin{namedresult}{Lemma}{(marked inclusions)}\label{lem:sec3-marked-inclusion}
 \lean{LineSystem.Hyp24.orthGraph_embedding_of_subset}
 \uses{hyp:2.4,lem:sec3-transport}
 \assume{$S,T$ are indecomposable star-closed line systems and a linear isometry
 of spans sends $\Phi(S)$ into $\Phi(T)$. Choose $\Za,\Zb\in\Phi(S)$ with
 $(\Za,\Zb)=-1$, and mark $T$ by their images.}
 \prove{The isometry induces an injection from the source $\Ga$ to the target
 $\Ga$ that preserves both adjacency and nonadjacency in their orthogonality graphs.}
\end{namedresult}
\begin{proof}
 \pfsketch\ \new\ We state the induced-embedding direction needed for C5 separately
 from the isomorphism assertion of Lemma 2.12. We will retain the same marked star
 and use the defining inner products.
 \step{s3markmembership}{A source member of $\Ga$ maps into the target $\Ga$.}
 \begin{proof}\pf\ The three defining products of Hypothesis~\ref{hyp:2.4}
  are preserved, and the image is a target root.~\qed\end{proof}
 \step{s3markinjection}{The induced map is injective.}
 \begin{proof}\pf\ A linear isometry is injective.~\qed\end{proof}
 \step{s3markedges}{For two distinct source vertices, the product is zero if and only if the image product is zero.}
 \begin{proof}\pf\ The isometry preserves that product.~\qed\end{proof}
 \qedstep
 \begin{proof}\pf\ The three steps give an induced graph embedding.~\qed\end{proof}
\end{proof}

\begin{namedresult}{Computation}{C5}\label{comp:C5}
 \lean{Computations.C5.inclusions_and_obstructions}
 \lean{LineSystem.not_embeds_An_Dn}
 \uses{def:3.1,def:3.2,def:3.3,comp:C3,comp:C4,lem:3.4,lem:sec3-spans,lem:sec3-choices,lem:sec3-marked-inclusion,lem:2.10}
 \assume{All inclusions are isometric inclusions of the signed models, in their spans.}
 \prove{There are proper inclusions $\typ A_8\subset \typ E_8$, $\typ D_8\subset \typ E_8$,
 and $\typ A_7\subset \typ E_7$, all of equal rank. There are no inclusions
 $\typ A_\ell\subset \typ D_\ell$ for $\ell\geq4$, $\typ A_6\subset \typ E_6$,
 $\typ D_6\subset \typ E_6$, or $\typ D_7\subset \typ E_7$.}
\end{namedresult}
\begin{proof}
 \pfsketch\ We will exhibit the inclusions, then obstruct induced marked graphs
 using coclique sizes and degrees. The $\typ A_\ell$ obstruction fills the first omission in Gap 13.
 \step{s3c5uniform}{The target marked graphs have the asserted types for every ordered star.}
 \begin{proof}
  \step{s3c5Dtransitive}{Signed coordinate permutations of $\typ D_\ell$ are transitive on ordered root pairs of product $-1$.}
  \begin{proof}
   \step{s3c5Dsupport}{Such roots share exactly one support position, with opposite signs there.}
   \begin{proof}\pf\ Equal supports give product $-2,0,2$; disjoint supports give
    0; a one-position intersection contributes its single sign product.~\qed\end{proof}
   \qedstep
   \begin{proof}\pf\ Send the unshared position of the first root to 1, the shared
    position to 2, and the unshared position of the second to 3. Choose signs
    to obtain $\Ze_1-\Ze_2,\Ze_2-\Ze_3$; extend the coordinate permutation and signs arbitrarily.~\qed\end{proof}
  \end{proof}
  \step{s3c5Euniform}{Every $\typ E_6$ marked graph is the lattice graph; every $\typ E_7$ marked graph is the disjointness graph on six symbols.}
  \begin{proof}\pf\ C4 gives constant order and parameters for every star.
   Lemma~\ref{lem:2.10} identifies its graph with the fixed-star graph, whose
   labelled type is certified in C4.~\qed\end{proof}
  \qedstep
  \begin{proof}\pf\ For $\typ D_\ell$ apply C4 after \stepref{s3c5Dtransitive}; for
   the two exceptional targets use \stepref{s3c5Euniform}.~\qed\end{proof}
 \end{proof}
 \step{s3c5Aeight}{$\typ A_8$ embeds properly in $\typ E_8$ with equal rank.}
 \begin{proof}\pf\ It is a subset of model (i); apply C3. Lemma~\ref{lem:sec3-spans}
  gives rank 8 on both sides, and Lemma~\ref{lem:3.4} gives $36<120$ lines.~\qed\end{proof}
 \step{s3c5Deight}{$\typ D_8$ embeds properly in $\typ E_8$ with equal rank.}
 \begin{proof}\pf\ Definition~\ref{def:3.2}(ii) includes these roots.
  Their ranks are 8 and their counts satisfy $56<120$ by the same two lemmas.~\qed\end{proof}
 \step{s3c5Aseven}{\new\ $\typ A_7$ embeds properly in $\typ E_7$ with equal rank. An explicit
  coordinate inclusion replaces the paper's embedding of marked graphs.}
 \begin{proof}
  \pflet{$\Zh=\frac12\mathbf1\in\re^8$.}
  \step{s3c5Ah}{$\typ A_7\subset \typ E_8\cap\Zh^\perp$.}
  \begin{proof}\pf\ Its roots are $\Ze_i-\Ze_j$ in $\re^8$, which are $\typ D_8$
   roots with product zero with $\Zh$. The all-positive half root $\Zh$ belongs to $\typ E_8$.~\qed\end{proof}
  \step{s3c5Amap}{An $\typ E_8$ automorphism carries this inclusion into $\typ E_7$.}
  \begin{proof}\pf\ Use the root assertion of Lemma~\ref{lem:sec3-choices}
   to send $\Zh$ to $\Zr_0$.~\qed\end{proof}
  \step{s3c5Aproper}{The inclusion has rank 7 on both sides and is proper.}
  \begin{proof}\pf\ Lemma~\ref{lem:sec3-spans} gives the ranks;
   Lemma~\ref{lem:3.4} gives $28<63$.~\qed\end{proof}
  \qedstep
  \begin{proof}\pf\ Combine \stepref{s3c5Amap} and \stepref{s3c5Aproper}.~\qed\end{proof}
 \end{proof}
 \step{s3c5AD}{\repair\ \assume{$\ell\geq4$ and an isometric inclusion $\typ A_\ell\subset \typ D_\ell$ exists.}
  \prove{FALSE.}}
 \begin{proof}
  \step{s3c5ADinduced}{The target friendship graph with $\ell-3$ triangles contains a coclique of size $\ell-2$.}
  \begin{proof}\pf\ Mark the source by its C4 star and apply
   Lemma~\ref{lem:sec3-marked-inclusion} and \stepref{s3c5uniform}.~\qed\end{proof}
  \step{s3c5friendbound}{A friendship graph with $t\geq1$ triangles has no coclique larger than $t$.}
  \begin{proof}
   \step{s3c5friendcentre}{A coclique containing the common vertex has size 1.}
   \begin{proof}\pf\ Every other vertex is adjacent to that vertex.~\qed\end{proof}
   \step{s3c5friendouter}{A coclique avoiding the common vertex has at most one vertex from each of the $t$ outer pairs.}
   \begin{proof}\pf\ Each outer pair is an edge.~\qed\end{proof}
   \qedstep
   \begin{proof}\pf\ These two cases give the bound $t$, since $t\geq1$.~\qed\end{proof}
  \end{proof}
  \qedstep
  \begin{proof}\pf\ Steps \stepref{s3c5ADinduced} and \stepref{s3c5friendbound}
   would give $\ell-2\leq\ell-3$.~\qed\end{proof}
 \end{proof}
 \step{s3c5Asix}{\assume{An isometric inclusion $\typ A_6\subset \typ E_6$ exists.}\prove{FALSE.}}
 \begin{proof}
  \step{s3c5gridfour}{The $3\times3$ lattice graph contains a 4-coclique.}
  \begin{proof}\pf\ C4 gives a 4-coclique for $\typ A_6$; use
   Lemma~\ref{lem:sec3-marked-inclusion} and \stepref{s3c5uniform}.~\qed\end{proof}
  \qedstep
  \begin{proof}\pf\ Among four grid positions two share one of the three rows,
   so those two are adjacent. This contradicts \stepref{s3c5gridfour}.~\qed\end{proof}
 \end{proof}
 \step{s3c5Dsix}{\assume{An isometric inclusion $\typ D_6\subset \typ E_6$ exists.}\prove{FALSE.}}
 \begin{proof}\pf\ The centre in the source marked graph has six neighbours by C4.
  Its image has at least six distinct neighbours by Lemma~\ref{lem:sec3-marked-inclusion};
  C4 gives target degree 4.~\qed\end{proof}
 \step{s3c5Dseven}{\assume{An isometric inclusion $\typ D_7\subset \typ E_7$ exists.}\prove{FALSE.}}
 \begin{proof}\pf\ The source centre has eight neighbours by C4; the induced
  injection of Lemma~\ref{lem:sec3-marked-inclusion} contradicts target degree 6.~\qed\end{proof}
 \qedstep
 \begin{proof}\pf\ The three positive inclusions and four obstructions give the assertion.~\qed\end{proof}
\end{proof}
\begingroup\raggedright
The coordinate inclusions already have Lean declarations
\leancite{LineSystem.An_Φ_subset_E₈viaA₈}{\texttt{LineSystem.\allowbreak An\_\allowbreak \ensuremath{\Phi}\_\allowbreak subset\_\allowbreak E\ensuremath{{}_8}viaA\ensuremath{{}_8}}} and \leancite{LineSystem.Dn_Φ_subset_E₈}{\texttt{LineSystem.\allowbreak Dn\_\allowbreak \ensuremath{\Phi}\_\allowbreak subset\_\allowbreak E\ensuremath{{}_8}}}.
C5 is proved in Lean, in \path{LineSystem/Inclusions.lean},
\path{Computations/C5Inclusions.lean}, \path{Computations/C5Dn.lean} and
\path{Computations/C5.lean}. An isometric inclusion is
\path{LineSystem.IsometricInclusion}; the lemma on an equal-rank proper inclusion is
\path{LineSystem.not_isMaximal_of_proper_isometric_inclusion}, and the marked-inclusion lemma
is \path{LineSystem.Hyp24.orthGraph_embedding_of_subset}. Lean proves
$\typ A_7\subset\typ E_7$ and $\typ A_8\subset\typ E_8$ by explicit integer coordinate rows,
\leancite{LineSystem.An.inclusionE₇}{\texttt{LineSystem.\allowbreak An.\allowbreak inclusionE\ensuremath{{}_7}}} and \leancite{LineSystem.An.inclusionE₈}{\texttt{LineSystem.\allowbreak An.\allowbreak inclusionE\ensuremath{{}_8}}}, not
through C3; $\typ D_8\subset\typ E_8$ is \leancite{LineSystem.Dn.inclusionE₈}{\texttt{LineSystem.\allowbreak Dn.\allowbreak inclusionE\ensuremath{{}_8}}}. Step
\texttt{s3c5Dtransitive} is \leancite{LineSystem.Dn.exists_isometry_star}{\texttt{LineSystem.\allowbreak Dn.\allowbreak exists\_\allowbreak isometry\_\allowbreak star}}, with
\leancite{LineSystem.Hyp24.nonempty_orthGraph_iso_of_S_eq_Dn}{\texttt{LineSystem.\allowbreak Hyp24.\allowbreak nonempty\_\allowbreak orthGraph\_\allowbreak iso\_\allowbreak of\_\allowbreak S\_\allowbreak eq\_\allowbreak Dn}}. The four obstructions are
\leancite{LineSystem.not_embeds_An_Dn}{\texttt{LineSystem.\allowbreak not\_\allowbreak embeds\_\allowbreak An\_\allowbreak Dn}} (for $\ell\geq4$), \leancite{LineSystem.not_embeds_An_E₆}{\texttt{LineSystem.\allowbreak not\_\allowbreak embeds\_\allowbreak An\_\allowbreak E\ensuremath{{}_6}}},
\leancite{LineSystem.not_embeds_Dn_E₆}{\texttt{LineSystem.\allowbreak not\_\allowbreak embeds\_\allowbreak Dn\_\allowbreak E\ensuremath{{}_6}}} and \leancite{LineSystem.not_embeds_Dn_E₇}{\texttt{LineSystem.\allowbreak not\_\allowbreak embeds\_\allowbreak Dn\_\allowbreak E\ensuremath{{}_7}}}, and
\leancite{Computations.C5.inclusions_and_obstructions}{\texttt{Computations.\allowbreak C5.\allowbreak inclusions\_\allowbreak and\_\allowbreak obstructions}} bundles the three proper
inclusions with them. There is no finite search: the grid bound of step \texttt{s3c5Asix} is the
pigeonhole argument above, and the friendship bound and the two degree contradictions are
proofs above as well.\par
\endgroup

\begin{remark}
 \repair\ The paper derives $\typ A_7\subset\typ E_7$ from the statement that a 5-coclique
 embeds in the triangular graph $T(6)$ (Error 5). With the orthogonality convention of C4,
 the graph of $\typ A_7$ is a 5-coclique and the graph of $\typ E_7$ is the disjointness
 graph $\comp{T(6)}$; the witness is a 5-coclique of $\comp{T(6)}$, equivalently a 5-clique
 of $T(6)$, and the five pairs $\{1,j\}$ with $2\leq j\leq6$ pairwise intersect. We record
 the correction only; C5 does not use this witness, since it proves the inclusion by an
 explicit coordinate inclusion.
\end{remark}

\begin{cgss}{Theorem}{3.6}\label{thm:3.6}
 \lean{LineSystem.classification_of_isMaximal}
 \uses{def:2.2,lem:2.3,thm:3.5,lem:3.4,comp:C4,comp:C5,lem:sec3-spans,lem:sec3-transport,lem:sec3-extension,lem:sec3-proper-inclusion}
 \assume{$S$ is a nonempty indecomposable line system in a finite-dimensional real
 inner product space, and $W=\operatorname{span}\Phi(S)$.}
 \prove{$S$ is maximal within $W$ if and only if it is isomorphic to one of
 \[
 \typ A_\ell\ (\ell\geq1,\ \ell\ne7,8),\qquad
 \typ D_\ell\ (\ell\geq4,\ \ell\ne8),\qquad \typ E_6,\ \typ E_7,\ \typ E_8.
 \]
 \repair\ Maximality is within the span (Implicit hypothesis 7), including for
 $\typ A_\ell\subset\re^{\ell+1}$.}
\end{cgss}
\begin{proof}
 \pfsketch\ A maximal system is star-closed. For the converse we will put any
 hypothetical proper extension inside an indecomposable star-closed system of the
 same rank, then exhaust its type using counts and C5.
 \step{s336forward}{\assume{$S$ is maximal within $W$.}
  \prove{$S$ has a type in the displayed list.}}
 \begin{proof}
  \step{s336closed}{$S$ is star-closed.}
  \begin{proof}\pf\ Apply Lemma~\ref{lem:2.3}.~\qed\end{proof}
  \step{s336classified}{$S$ has a type from Theorem 3.5.}
  \begin{proof}\pf\ Apply Theorem~\ref{thm:3.5} to \stepref{s336closed} and the standing hypotheses.~\qed\end{proof}
  \step{s336excluded}{Its type is not $\typ A_7$, $\typ A_8$, or $\typ D_8$.}
  \begin{proof}
   \step{s336exclAseven}{$\typ A_7$ is not maximal within its span.}
   \begin{proof}\pf\ Apply Lemma~\ref{lem:sec3-proper-inclusion} to the proper
    isometric inclusion $\typ A_7\subset\typ E_7$ of equal rank in C5.~\qed\end{proof}
   \step{s336exclAeight}{$\typ A_8$ is not maximal within its span.}
   \begin{proof}\pf\ Apply Lemma~\ref{lem:sec3-proper-inclusion} to the proper
    isometric inclusion $\typ A_8\subset\typ E_8$ of equal rank in C5.~\qed\end{proof}
   \step{s336exclDeight}{$\typ D_8$ is not maximal within its span.}
   \begin{proof}\pf\ Apply Lemma~\ref{lem:sec3-proper-inclusion} to the proper
    isometric inclusion $\typ D_8\subset\typ E_8$ of equal rank in C5.~\qed\end{proof}
   \qedstep
   \begin{proof}\pf\ A span isometry from $S$ to one of these models would carry
    maximality within $W$, from \asm{1}, to the model, by
    Lemma~\ref{lem:sec3-transport}; this contradicts \stepref{s336exclAseven},
    \stepref{s336exclAeight} or \stepref{s336exclDeight}.~\qed\end{proof}
  \end{proof}
  \qedstep
  \begin{proof}\pf\ Remove the three types of \stepref{s336excluded} from
   \stepref{s336classified}.~\qed\end{proof}
 \end{proof}
 \step{s336modelsmaximal}{Every retained model is maximal in its span.}
 \begin{proof}
  \step{s336contra}{\assume{$M$ is a retained model and is not maximal in its span.}
   \prove{FALSE.}}
  \begin{proof}
   \step{s336extension}{There is an indecomposable star-closed proper extension $T$ with the same rank as $M$.}
   \begin{proof}\pf\ C4 gives that $M$ is nonempty and indecomposable;
    apply Lemma~\ref{lem:sec3-extension} to $M$ in its span.~\qed\end{proof}
   \step{s336target}{There is a proper isometric inclusion of $M$ in a listed model
    of Theorem 3.5 of the same rank, with strictly more lines.}
   \begin{proof}
    \pflet{$W_M=\operatorname{span}\Phi(M)$.}
    \step{s336targethyp}{$T$ is a nonempty indecomposable star-closed line system in
     $W_M$, and $\operatorname{span}\Phi(T)=W_M$.}
    \begin{proof}\pf\ Lemma~\ref{lem:sec3-extension} gives these properties with the
     ambient space $W_M$, and $\Phi(M)\subseteq\Phi(T)$ makes $T$ nonempty.~\qed\end{proof}
    \step{s336targetiso}{There are a listed model $Y$ and a surjective linear isometry
     $G:W_M\to\operatorname{span}\Phi(Y)$ with $G(\Phi(T))=\Phi(Y)$.}
    \begin{proof}\pf\ Apply the forward direction of Theorem~\ref{thm:3.5} to $T$ in the
     finite-dimensional space $W_M$, using \stepref{s336targethyp}; its span is
     $W_M$.~\qed\end{proof}
    \step{s336targetmap}{$G$ is an isometric inclusion of $M$ in $Y$ with
     $G(\Phi(M))\subsetneq\Phi(Y)$.}
    \begin{proof}\pf\ $G$ is defined on $W_M=\operatorname{span}\Phi(M)$, it is injective,
     and $\Phi(M)\subsetneq\Phi(T)$ by \stepref{s336extension}; apply
     \stepref{s336targetiso}.~\qed\end{proof}
    \step{s336targetrank}{$Y$ has the rank of $M$.}
    \begin{proof}\pf\ $G$ is a linear bijection $W_M\to\operatorname{span}\Phi(Y)$ by
     \stepref{s336targetiso}.~\qed\end{proof}
    \qedstep
    \begin{proof}\pf\ A proper inclusion of finite antipodal sets strictly increases
     the number of pairs; use \stepref{s336targetmap} and \stepref{s336targetrank}.~\qed\end{proof}
   \end{proof}
   \step{s336E}{\case{$M=\typ E_r$ for $r\in\{6,7,8\}$.}}
   \begin{proof}\pf\ At those ranks the $(\typ A,\typ D,\typ E)$ line counts in
    Lemma~\ref{lem:3.4} are $(21,30,36)$, $(28,42,63)$, and $(36,56,120)$.
    The $\typ E$ entry is largest in each row, contrary to \stepref{s336target}.~\qed\end{proof}
   \step{s336A}{\case{$M=\typ A_\ell$, $\ell\ne7,8$.}}
   \begin{proof}
    \step{s336Asmall}{\case{$\ell\leq3$.}}
    \begin{proof}\pf\ Theorem~\ref{thm:3.5} has only $\typ A_\ell$ at this rank,
     so the strict count increase in \stepref{s336target} is impossible.~\qed\end{proof}
    \step{s336Alarge}{\case{$\ell\geq4$.}}
    \begin{proof}
     \step{s336AD}{The target cannot be $\typ D_\ell$.}
     \begin{proof}\pf\ C5 excludes $\typ A_\ell\subset \typ D_\ell$ for every $\ell\geq4$.~\qed\end{proof}
     \step{s336AE}{The target cannot be exceptional.}
     \begin{proof}\pf\ Rank equality and $\ell\ne7,8$ leave only $\typ E_6$;
      C5 excludes $\typ A_6\subset \typ E_6$.~\qed\end{proof}
     \qedstep
     \begin{proof}\pf\ The only remaining target is $\typ A_\ell$, with equal count,
      contrary to \stepref{s336target}.~\qed\end{proof}
    \end{proof}
    \qedstep
    \begin{proof}\pf\ Steps \stepref{s336Asmall} and \stepref{s336Alarge} exhaust positive ranks.~\qed\end{proof}
   \end{proof}
   \step{s336D}{\case{$M=\typ D_\ell$, $\ell\geq4$, $\ell\ne8$.}}
   \begin{proof}
    \step{s336DA}{The target cannot be $\typ A_\ell$ or $\typ D_\ell$.}
    \begin{proof}\pf\ Their counts are at most that of $\typ D_\ell$, since
     $\ell(\ell-1)-\ell(\ell+1)/2=\ell(\ell-3)/2>0$;
     this contradicts \stepref{s336target}.~\qed\end{proof}
    \step{s336DE}{The target cannot be exceptional.}
    \begin{proof}\pf\ With equal ranks and $\ell\ne8$, only $\typ D_6\subset \typ E_6$
     and $\typ D_7\subset \typ E_7$ remain; C5 excludes both.~\qed\end{proof}
    \qedstep
    \begin{proof}\pf\ Theorem~\ref{thm:3.5} has no further target types.~\qed\end{proof}
   \end{proof}
   \qedstep
   \begin{proof}\pf\ The retained model list is exhausted by
    \stepref{s336E}, \stepref{s336A}, and \stepref{s336D}.~\qed\end{proof}
  \end{proof}
  \qedstep
  \begin{proof}\pf\ Step \stepref{s336contra} rules out nonmaximality of any retained model.~\qed\end{proof}
 \end{proof}
 \step{s336reverse}{\assume{$S$ is isomorphic to a retained model.}
  \prove{$S$ is maximal within $W$.}}
 \begin{proof}\pf\ Transport \stepref{s336modelsmaximal} by
  Lemma~\ref{lem:sec3-transport}.~\qed\end{proof}
 \qedstep
 \begin{proof}\pf\ Steps \stepref{s336forward} and \stepref{s336reverse} give the equivalence.~\qed\end{proof}
\end{proof}

\begin{namedresult}{Definition}{(signings and the matrix convention)}\label{def:sec3-signing}
 \lean{LineSystem.Signing}\lean{LineSystem.signedAdj}\lean{LineSystem.nonempty_signing}
 \uses{def:line-system,facts:inner}
 \assume{$S$ is a line system, $W=\operatorname{span}\Phi(S)$, and $I$ is a finite index set.}
 \new\ We make ``one vector per line'' an explicit bijection.
 \pflet{A signing is a family $\sigma:I\to\Phi(S)$ for which
 $(i,\epsilon)\mapsto\epsilon\sigma(i)$ is a bijection
 $I\times\{1,-1\}\to\Phi(S)$. Put $N=|I|$,
 $M(i,j)=(\sigma(i),\sigma(j))$, and $C=M-2I_N$.}
 Thus $N$ is the number of lines. The diagonal of $C$ is zero, and its off-diagonal
 entries lie in $\{0,1,-1\}$ by Section~\ref{facts:inner}. A change of signs conjugates
 $C$ by a diagonal matrix of signs, and a reordering conjugates by a permutation matrix.
 We use $I_N$ for the identity matrix, to distinguish it from the index set $I$.
 Every line system has a signing indexed by $\{1,\dots,N\}$, where $N$ is its number of lines:
 choose one of the two vectors on each line (\path{LineSystem.nonempty_signing}). Thus the
 equations of Theorem 3.7, which hold for every signing, are not vacuous.
\end{namedresult}

\begin{namedresult}{Computation}{C6}\label{comp:C6}
 \lean{Computations.C6.coordinate_frame_matrices}
 \uses{def:3.1,def:3.2,def:3.3,lem:sec3-integer,comp:C2}
 \assume{$R$ is the complete, duplicate-free list of raw signed roots from
 Lemma~\ref{lem:sec3-integer}, with the coordinate cuts for $\typ E_7,\typ E_6$ from
 Definition~\ref{def:3.3}. Write
 $K=\sum_{v\in R}vv^T$ and let $J_m$ be the $m\times m$ all-one matrix.}
 \prove{The following are exact integer matrix identities:
 \[
 \begin{array}{c|c|c}
 \text{model}&q&K\\\hline
 \typ A_\ell&1&2((\ell+1)I_{\ell+1}-J_{\ell+1})\\
 \typ D_\ell&1&4(\ell-1)I_\ell\\
 \typ E_8\text{ (i)}&9&540I_9-60J_9\\
 \typ E_8\text{ (ii)}&4&240I_8\\
 \typ E_8\text{ (iii)}&2&120I_8\\
 \typ E_7&4&\operatorname{diag}(144I_6,72J_2)\\
 \typ E_6&4&\operatorname{diag}(96I_5,32J_3).
 \end{array}
 \]}
\end{namedresult}
\begin{proof}
 \pfsketch\ \new\ We replace spectral calculations by small coordinate frame
 matrices. The classical rows are parametric; the exceptional rows use the
 existing Gram checker on the columns of the raw root table.
 \step{s3c6A}{The $\typ A_\ell$ identity holds for every $\ell\geq1$.}
 \begin{proof}
  \step{s3c6Adiag}{$K(i,i)=2\ell$.}
  \begin{proof}\pf\ There are $\ell$ difference roots with entry 1 at $i$ and
   $\ell$ with entry $-1$ there.~\qed\end{proof}
  \step{s3c6Aoff}{For $i\ne j$, $K(i,j)=-2$.}
  \begin{proof}\pf\ Only $\Ze_i-\Ze_j$ and its negative contribute; both products are $-1$.~\qed\end{proof}
  \qedstep
  \begin{proof}\pf\ These entries are those of $2((\ell+1)I-J)$.~\qed\end{proof}
 \end{proof}
 \step{s3c6D}{The $\typ D_\ell$ identity holds for every $\ell\geq4$.}
 \begin{proof}
  \step{s3c6Ddiag}{$K(i,i)=4(\ell-1)$.}
  \begin{proof}\pf\ Each support $\{i,j\}$ with $j\ne i$ has four sign choices, each contributing 1.~\qed\end{proof}
  \step{s3c6Doff}{For $i\ne j$, $K(i,j)=0$.}
  \begin{proof}\pf\ The four signs on the only contributing support $\{i,j\}$
   give products $1,-1,-1,1$.~\qed\end{proof}
  \qedstep
  \begin{proof}\pf\ The entries in the two steps give $4(\ell-1)I$.~\qed\end{proof}
 \end{proof}
 \step{s3c6exceptional}{The five exceptional raw identities hold.}
 \begin{proof}
  \pflet{$P_i=(v_i)_{v\in R}$ be the $i$th column of the raw root table, in its
   list order. The list $P$ has length equal to the ambient dimension.}
  \step{s3c6one}{For model (i), \texttt{checkGram} accepts $P$ and $540I_9-60J_9$.}
  \begin{proof}\pf\ C6 checks \path{Computations.checkGram} with vector
   dimension 240 and nine columns, from the pair and signed-triple list of C2.~\qed\end{proof}
  \step{s3c6two}{For model (ii), \texttt{checkGram} accepts $P$ and $240I_8$.}
  \begin{proof}\pf\ C6 applies that checker with dimension 240 and eight columns
   of the doubled pair and even-sign list.~\qed\end{proof}
  \step{s3c6three}{For model (iii), \texttt{checkGram} accepts $P$ and $120I_8$.}
  \begin{proof}\pf\ C6 applies that checker with dimension 240 and eight columns
   of the axis and block-sign list.~\qed\end{proof}
  \step{s3c6seven}{For $\typ E_7$, \texttt{checkGram} accepts $P$ and $\operatorname{diag}(144I_6,72J_2)$.}
  \begin{proof}\pf\ C6 applies that checker with dimension 126 and eight columns
   of the doubled list filtered by $v_7=v_8$.~\qed\end{proof}
  \step{s3c6six}{For $\typ E_6$, \texttt{checkGram} accepts $P$ and $\operatorname{diag}(96I_5,32J_3)$.}
  \begin{proof}\pf\ C6 applies that checker with dimension 72 and eight columns
   of the doubled list filtered by $v_6=v_7=v_8$.~\qed\end{proof}
  \qedstep
  \begin{proof}\pf\ For each accepted table,
   \path{Computations.checkGram_sound} gives
   $K(i,j)=P_i\cdot P_j=\sum_{v\in R}v_iv_j$, including diagonal entries.~\qed\end{proof}
 \end{proof}
 \qedstep
 \begin{proof}\pf\ The two parametric cases and five column-Gram checks cover
  the seven rows.~\qed\end{proof}
\end{proof}

\begingroup\raggedright
C6 is proved in Lean, in \path{LineSystem/Frame.lean} and \path{Computations/C6Classical.lean},
\path{Computations/C6E8.lean}, \path{Computations/C6E7E6.lean} and \path{Computations/C6.lean}.
Lean states C6 as the frame property on the span, \path{LineSystem.IsFrame}, equivalent to the
column-Gram form by \path{LineSystem.isFrame_iff_gram} and transported by
\path{LineSystem.Iso.isFrame}. The two parametric rows are \path{Computations.C6.an_frame_matrix}
and \path{Computations.C6.dn_frame_matrix}, giving \path{LineSystem.isFrame_An} and
\path{LineSystem.isFrame_Dn}; the exceptional matrices are checked by \texttt{decide +kernel} over
structured indices, not by \path{Computations.checkGram}, and give
\leancite{LineSystem.isFrame_E₈}{\texttt{LineSystem.\allowbreak isFrame\_\allowbreak E\ensuremath{{}_8}}}, \leancite{LineSystem.isFrame_E₇}{\texttt{LineSystem.\allowbreak isFrame\_\allowbreak E\ensuremath{{}_7}}} and \leancite{LineSystem.isFrame_E₆}{\texttt{LineSystem.\allowbreak isFrame\_\allowbreak E\ensuremath{{}_6}}}.
The bundle is \path{Computations.C6.coordinate_frame_matrices}. The checker
\path{Computations.checkGram} is defined in \path{Computations/ListData.lean}, and its
soundness theorem \path{Computations.checkGram_sound} is in \path{Computations/Checkers.lean};
C6 does not use them. The frame lemma and Theorem~3.7 below are proved in Lean as well, the
frame lemma with the rank of the model (\path{LineSystem.frame_identity}) and with the rank of
the span (\path{LineSystem.frame_identity_span}). The three span lemmas are in Lean:
\texttt{LineSystem.\allowbreak Iso.\allowbreak spanIsometry} with the transport of
the predicates and of the marked sets, \texttt{LineSystem.\allowbreak restrictSpan} (with
\texttt{LineSystem.\allowbreak Iso.\allowbreak restrictSpan} and \texttt{LineSystem.\allowbreak Iso.\allowbreak ofSpanIsometry},
the interface of the final step of Theorem~3.5) in \path{LineSystem/SpanTransport.lean}, and
\texttt{LineSystem.\allowbreak exists\_\allowbreak maximal\_\allowbreak extension\_\allowbreak span} in
\path{LineSystem/SpanExtension.lean}; \path{LineSystem/SpanAudit.lean} prints their axioms. The Lean
proofs of C5 and C6 transport along \texttt{LineSystem.Iso} and use the model's rank.\par
Theorems~3.5 and~3.6 are in Lean. Theorem~3.5 is
\texttt{LineSystem.\allowbreak classification} in \path{LineSystem/Classification.lean}, with its
forward direction \texttt{LineSystem.\allowbreak exists\_\allowbreak iso\_\allowbreak rootType};
Theorem~3.6 is \texttt{LineSystem.\allowbreak classification\_\allowbreak of\_\allowbreak isMaximal}
in \path{LineSystem/ClassificationMaximal.lean}, with the form that Section~4 uses,
\texttt{LineSystem.\allowbreak IsMaximal.\allowbreak exists\_\allowbreak iso\_\allowbreak rootType}.
Both rest on the unconditional Theorem~2.13,
\texttt{LineSystem.\allowbreak five\_\allowbreak possibilities'}, whose inputs from Section~3 are
the packaged facts \texttt{LineSystem.\allowbreak modelInputs}, proved from C2 and C4
without either classification, in \path{LineSystem/ClassificationModels.lean}.
Lemma (a connected spanning set) is
\texttt{LineSystem.\allowbreak isIndecomposable\_\allowbreak of\_\allowbreak connected\_\allowbreak spanning}
in \path{LineSystem/ClassificationConnected.lean}, and the perpendicular choices of
Lemma (independence of the perpendicular choices) in the default model are
\leancite{LineSystem.E₈.nonempty_iso_E₇_orthogonalTo}{\texttt{LineSystem.\allowbreak E\ensuremath{{}_8}.\allowbreak nonempty\_\allowbreak iso\_\allowbreak E\ensuremath{{}_7}\_\allowbreak orthogonalTo}},
\leancite{LineSystem.E₈.nonempty_iso_E₆_orthogonalTo}{\texttt{LineSystem.\allowbreak E\ensuremath{{}_8}.\allowbreak nonempty\_\allowbreak iso\_\allowbreak E\ensuremath{{}_6}\_\allowbreak orthogonalTo}}
and
\leancite{LineSystem.E₈.exists_isometry_of_ordered_star}{\texttt{LineSystem.\allowbreak E\ensuremath{{}_8}.\allowbreak exists\_\allowbreak isometry\_\allowbreak of\_\allowbreak ordered\_\allowbreak star}}
in \path{LineSystem/ClassificationChoices.lean}. The Lean departs from the proofs above in three
places. The ambient space of Theorems~3.5 and~3.6 need not be finite-dimensional, since the
span of the finite root set is, and Theorem~2.13 is applied there. The forward direction of
Theorem~3.5 is shorter: the isomorphisms of Theorem~2.13 go directly to the named models and are
composed with \texttt{LineSystem.\allowbreak Iso.\allowbreak restrictSpan}, so steps
\texttt{s335coords} and \texttt{s335orthogonal} of the proof below are not needed. The retained list of
Theorem~3.6 is stated as an explicit disjunction of the five families.\par
The two remaining pieces are in Lean as well, and their modules are imported by
\path{LineGraphs.lean}. The bundles
\texttt{Computations.\allowbreak C2.\allowbreak model\_\allowbreak presentations} and
\texttt{Computations.\allowbreak C2.\allowbreak line\_\allowbreak counts} of
Lemma (integer realizations) and C2 are in \path{Computations/C2.lean}, with the cuts
\leancite{LineSystem.T₇}{\texttt{LineSystem.\allowbreak T\ensuremath{{}_7}}} and
\leancite{LineSystem.T₆}{\texttt{LineSystem.\allowbreak T\ensuremath{{}_6}}} of
Definition~3.3; Lean checks their 126 and 72 roots in the kernel, as for $\typ E_7$
and $\typ E_6$. The assertions of Lemma (independence of the perpendicular choices) for the
other two $\typ E_8$ models, with $T_7\cong\typ E_7$ and $T_6\cong\typ E_6$, are in
\path{LineSystem/ClassificationChoicesModels.lean}. In the final step the Lean departs from the
proof above: a C3 isomorphism is a bijection of roots, not a map of the ambient space, so Lean
carries the automorphisms of the default model across it through the span isometry of
Lemma (transport on spans), extended to the whole space (for model (i), from the zero-sum
hyperplane to $\re^9$). The file \path{Computations/C2ChoicesAudit.lean} prints the axioms of
all these declarations.\par
\endgroup

\begin{namedresult}{Lemma}{(the frame identity within the span)}\label{lem:sec3-frame}
 \lean{LineSystem.frame_identity}
 \uses{comp:C6,comp:C2,lem:sec3-spans,def:sec3-signing,lem:sec3-transport,facts:inner}
 \assume{$S$ is one of the listed models, or an isometric copy of it;
 $W=\operatorname{span}\Phi(S)$ has rank $\ell>0$ and $N$ lines.
 Let $\sigma:I\to\Phi(S)$ be any signing.}
 \prove{For every $\Zx\in W$,
 \[
 \sum_{\Zr\in\Phi(S)}(\Zr,\Zx)\Zr=\frac{4N}{\ell}\Zx,
 \qquad
 \sum_{i\in I}(\sigma(i),\Zx)\sigma(i)=\frac{2N}{\ell}\Zx.
 \]
 \repair\ These formulas are asserted on $W$, not on a larger coordinate space
 (Implicit hypothesis 7 and the model-span correction).}
\end{namedresult}
\begin{proof}
 \pfsketch\ We will divide C6 by the raw scale $q$, restrict to the known span,
 and halve the sum using antipodal pairs.
 \step{s3frameoperator}{For a coordinate model, the full signed-root operator is $K/q$.}
 \begin{proof}\pf\ Substitute $\Zr=c\Zv$ in the sum, using $c^2=1/q$.
  The exact no-multiplicity transfer is \leancite{LineSystem.IntegerModel.sum_Φ}{\texttt{LineSystem.\allowbreak IntegerModel.\allowbreak sum\_\allowbreak \ensuremath{\Phi}}}.~\qed\end{proof}
 \step{s3framerestrict}{The operator is scalar on $W$, with scalar $4N/\ell$.}
 \begin{proof}
  \step{s3frameA}{On the $\typ A_\ell$ span, $K/q=2(\ell+1)I$.}
  \begin{proof}\pf\ $J\Zx=0$ on the zero-sum span in Lemma~\ref{lem:sec3-spans};
   use the first row of C6.~\qed\end{proof}
  \step{s3frameD}{On the $\typ D_\ell$ span, $K/q=4(\ell-1)I$.}
  \begin{proof}\pf\ Its C6 matrix has that value on the whole coordinate space.~\qed\end{proof}
  \step{s3frameEi}{On model (i)'s span, $K/q=60I$.}
  \begin{proof}\pf\ Its span has sum zero, so $(540I-60J)\Zx/9=60\Zx$.~\qed\end{proof}
  \step{s3frameEii}{On model (ii)'s span, $K/q=60I$.}
  \begin{proof}\pf\ Divide $240I_8$ by $q=4$.~\qed\end{proof}
  \step{s3frameEiii}{On model (iii)'s span, $K/q=60I$.}
  \begin{proof}\pf\ Divide $120I_8$ by $q=2$.~\qed\end{proof}
  \step{s3frameSeven}{On the $\typ E_7$ span, $K/q=36I$.}
  \begin{proof}\pf\ On a last-coordinate pair $(t,t)$, $72J_2(t,t)=144(t,t)$;
   the first block is $144I_6$. Divide by 4.~\qed\end{proof}
  \step{s3frameSix}{On the $\typ E_6$ span, $K/q=24I$.}
  \begin{proof}\pf\ On $(t,t,t)$, $32J_3(t,t,t)=96(t,t,t)$;
   the first block is $96I_5$. Divide by 4.~\qed\end{proof}
  \qedstep
  \begin{proof}\pf\ Substitution of the C2 line counts gives $4N/\ell$ in each row.~\qed\end{proof}
 \end{proof}
 \step{s3framecopies}{The full-root formula holds for isometric copies in their spans.}
 \begin{proof}\pf\ Conjugate the full-root operator through the span isometry
  of Lemma~\ref{lem:sec3-transport}; reindex by its root bijection.~\qed\end{proof}
 \step{s3framehalf}{The full-root sum equals twice the signing sum.}
 \begin{proof}
  \step{s3framepair}{The roots $\Zr,-\Zr$ contribute the same summand.}
  \begin{proof}\pf\ $(-\Zr,\Zx)(-\Zr)=(\Zr,\Zx)\Zr$ by bilinearity.~\qed\end{proof}
  \qedstep
  \begin{proof}\pf\ Reindex the full-root sum by the signing bijection in
   Definition~\ref{def:sec3-signing} and use \stepref{s3framepair}.~\qed\end{proof}
 \end{proof}
 \qedstep
 \begin{proof}\pf\ The full formula follows from \stepref{s3frameoperator}--\stepref{s3framecopies};
  divide the equality in \stepref{s3framehalf} by 2 for the signing formula.~\qed\end{proof}
\end{proof}

\begin{namedresult}{Lemma}{(the Gram polynomial from a frame)}\label{lem:sec3-gram-polynomial}
 \lean{LineSystem.signedAdj_quadratic_of_frame}
 \uses{def:sec3-signing,facts:inner}
 \assume{\begin{pfenum}
  \item $S$ is a line system and $W=\operatorname{span}\Phi(S)$.
  \item $\sigma:I\to\Phi(S)$ is a signing.
  \item $t\in\re$ satisfies $\sum_i(\sigma(i),\Zx)\sigma(i)=t\Zx$ for every $\Zx\in W$.
 \end{pfenum}}
 \prove{With $M=2I_N+C$ as in Definition~\ref{def:sec3-signing},
 $M^2=tM$ and $(C+2I_N)(C-(t-2)I_N)=0$.}
\end{namedresult}
\begin{proof}
 \pfsketch\ \new\ We isolate the algebraic passage from a frame to the signed
 Gram polynomial. We will apply the frame equation to each root and take its product
 with each other root. This avoids choosing an ambient orthonormal basis.
 \step{s3gramentry}{For all $i,j\in I$, $(M^2)(i,j)=tM(i,j)$.}
 \begin{proof}
  \step{s3gramsum}{$(M^2)(i,j)=(\sigma(i),\sum_k(\sigma(k),\sigma(j))\sigma(k))$.}
  \begin{proof}\pf\ Expand matrix multiplication and then bilinearity from Section~\ref{facts:inner}.~\qed\end{proof}
  \step{s3gramapply}{The vector sum in \stepref{s3gramsum} is $t\sigma(j)$.}
  \begin{proof}\pf\ Apply the assumed frame identity to $\sigma(j)\in\Phi(S)\subseteq W$.~\qed\end{proof}
  \qedstep
  \begin{proof}\pf\ Substitute \stepref{s3gramapply} into \stepref{s3gramsum} and use scalar linearity.~\qed\end{proof}
 \end{proof}
 \step{s3grammatrix}{$M(M-tI_N)=0$.}
 \begin{proof}\pf\ Step \stepref{s3gramentry} gives $M^2=tM$ entry by entry;
  distribute multiplication by $M-tI_N$.~\qed\end{proof}
 \qedstep
 \begin{proof}\pf\ Substitute $M=C+2I_N$ in \stepref{s3grammatrix};
  then $M-tI_N=C-(t-2)I_N$.~\qed\end{proof}
\end{proof}

\begin{cgss}{Theorem}{3.7}\label{thm:3.7}
 \lean{LineSystem.signedAdj_quadratic_An}\lean{LineSystem.signedAdj_quadratic_Dn}
 \lean{LineSystem.signedAdj_quadratic_E₈}\lean{LineSystem.signedAdj_quadratic_E₇}
 \lean{LineSystem.signedAdj_quadratic_E₆}
 \uses{def:sec3-signing,lem:3.4,lem:sec3-spans,comp:C6,lem:sec3-frame,lem:sec3-gram-polynomial}
 \assume{$S$ is $\typ A_\ell$ with $\ell\geq1$, $\typ D_\ell$ with $\ell\geq4$, or one
 of $\typ E_8,\typ E_7,\typ E_6$, in any isometric realization. Choose one vector of squared
 norm 2 on each line, and let $2I_N+C$ be their Gram matrix.}
 \prove{For every choice and ordering of those vectors,
 \[
 \begin{array}{rcl}
 \typ A_\ell:&(C+2I_N)(C-(\ell-1)I_N)&=0,\\
 \typ D_\ell:&(C+2I_N)(C-2(\ell-2)I_N)&=0,\\
 \typ E_8:&(C+2I_N)(C-28I_N)&=0,\\
 \typ E_7:&(C+2I_N)(C-16I_N)&=0,\\
 \typ E_6:&(C+2I_N)(C-10I_N)&=0.
 \end{array}
 \]
 \repair\ The $\typ E_8$ constant is 28, correcting the printed 18 (Error 1).}
\end{cgss}
\begin{proof}
 \pfsketch\ We will use C6's frame identity within the span. The second factor
 has constant $\theta=2N/\ell-2$, with $\ell$ the rank, not the number of ambient coordinates.
 \step{s337general}{$(C+2I_N)(C-(2N/\ell-2)I_N)=0$.}
 \begin{proof}\pf\ Apply Lemma~\ref{lem:sec3-gram-polynomial} to the signing
  formula of Lemma~\ref{lem:sec3-frame}.~\qed\end{proof}
 \step{s337A}{For $\typ A_\ell$, $2N/\ell-2=\ell-1$.}
 \begin{proof}\pf\ Lemma~\ref{lem:3.4} gives $N=\ell(\ell+1)/2$;
  cancel the positive $\ell$.~\qed\end{proof}
 \step{s337D}{For $\typ D_\ell$, $2N/\ell-2=2(\ell-2)$.}
 \begin{proof}\pf\ Substitute $N=\ell(\ell-1)$ and cancel $\ell>0$.~\qed\end{proof}
 \step{s337Eeight}{For $\typ E_8$, $2N/\ell-2=28$.}
 \begin{proof}\pf\ The rank and count are 8 and 120 by
  Lemmas~\ref{lem:sec3-spans} and~\ref{lem:3.4}; $240/8-2=28$.~\qed\end{proof}
 \step{s337Eseven}{For $\typ E_7$, $2N/\ell-2=16$.}
 \begin{proof}\pf\ Its rank and count are 7 and 63, so $126/7-2=16$.~\qed\end{proof}
 \step{s337Esix}{For $\typ E_6$, $2N/\ell-2=10$.}
 \begin{proof}\pf\ Its rank and count are 6 and 36, so $72/6-2=10$.~\qed\end{proof}
 \qedstep
 \begin{proof}\pf\ Substitute the five constants into \stepref{s337general}.
  The frame lemma applies to every signing and every isometric copy, so no
  additional switching or eigenvalue calculation is needed.~\qed\end{proof}
\end{proof}
\endgroup

\subsection{The classification behind Remark 3.11}\label{sec:rem311}
\begingroup
\emergencystretch=3em
\newcommand{\rsq}[1]{(#1,#1)}
\newcommand{\Cyc}{\operatorname{Cyc}}

We will follow the ordinary proof in \path{cgss/proofs/remark311/remark311.tex}.
Its lemma labels are retained. Theorems A, B and C are auxiliary results;
Corollary D is the strengthened form of [CGSS, Theorem 3.8].
All spaces below are finite-dimensional real inner product spaces.
Coordinate vectors $\Zu_i$ are orthonormal and start at index 1;
simple roots $\alpha_i$ start at index 0. A graph called connected is nonempty.
All Lean declarations of ours, as opposed to library declarations, occur in
\verb|\lean| metadata.

\begingroup\raggedright
Every one of these declarations is in Lean. Theorem A and Corollary D
are in \path{LineGraphs/Cartan/}. Theorems B and C are in
\path{LineGraphs/RootSetClassification/}, and Remark (small ranks) and
Lemma (a Gram realization) in \path{LineGraphs/Cartan/SmallRanks.lean} and
\path{LineGraphs/Cartan/GramRealization.lean}. Their
modules are imported by \path{LineGraphs.lean}, and
\path{LineGraphs/RootSetClassification/Audit.lean} prints their axioms.\par
\endgroup

\begin{cgss}{Definition}{3.9}\label{def:roots}\label{def:3.9}
\lean{IsRootSet}\uses{facts:inner}
\assume{$V$ is a finite-dimensional real inner product space.}
A Root System is a finite set $\Phi\subseteq V\setminus\{0\}$ spanning $V$,
such that, for all $\alpha,\beta\in\Phi$,
\[
 \frac{2(\alpha,\beta)}{\rsq\alpha}\in\mathbb Z,
 \qquad s_\alpha\beta:=\beta-
       \frac{2(\alpha,\beta)}{\rsq\alpha}\alpha\in\Phi,
\]
and $t\alpha\in\Phi$ implies $t\in\{1,-1\}$ for every $t\in\mathbb R$.
It is indecomposable if it is nonempty and has no partition into two
nonempty mutually orthogonal parts. A similarity is a linear bijection
$f:V\to W$ carrying the root sets onto one another and satisfying
$(f\Zx,f\Zy)=a(\Zx,\Zy)$ for some $a>0$ and all $\Zx,\Zy$.
\repair\ Nonzero roots, nonemptiness in indecomposability, and the use of
$\Phi$ for the printed $S$ in condition (b) are explicit here.
\end{cgss}

\begin{namedresult}{Lemma}{(composition and inversion of similarities)}\label{lem:rem-similarity-laws}
\lean{IsRootSet.similarity_comp}\lean{IsRootSet.similarity_symm}
\uses{def:roots}
\assume{$\Phi,\Psi,\Omega$ are Root Systems in real inner product spaces.}
\prove{The inverse of any similarity $f:\Phi\to\Psi$ with factor $a>0$
is a similarity with factor $a^{-1}$. Given also a similarity
$g:\Psi\to\Omega$ with factor $b>0$, $g\circ f$ is a similarity with
factor $ab$.}
\end{namedresult}
\begin{proof}
\pfsketch\ Substitute the metric identities and compose or invert the
linear bijections and their equalities of root sets.
\step{rsimcompose}{\assume{$f:\Phi\to\Psi$, $g:\Psi\to\Omega$ are
similarities with factors $a,b>0$.}
\prove{$g\circ f$ is a linear bijection,
$(gf\Zx,gf\Zy)=ab(\Zx,\Zy)$, and $gf(\Phi)=\Omega$.}}
\begin{proof}\pf\ Substitute the two identities of
Definition~\ref{def:roots}; $ab>0$.~\qed\end{proof}
\step{rsiminvert}{\assume{$f:\Phi\to\Psi$ is a similarity with factor $a>0$.}
\prove{$f^{-1}$ is a linear bijection,
$(f^{-1}\Zx,f^{-1}\Zy)=a^{-1}(\Zx,\Zy)$, and $f^{-1}(\Psi)=\Phi$.}}
\begin{proof}\pf\ Evaluate the identity for $f$ on the two inverse
images and divide by $a>0$.~\qed\end{proof}
\qedstep
\begin{proof}\pf\ Composition and inversion of linear bijections,
with \stepref{rsimcompose}, \stepref{rsiminvert}, give the assertions.~\qed\end{proof}
\end{proof}

\begin{namedresult}{Definition}{(bases and Cartan conventions)}\label{def:rem-bases}
\lean{IsRootSet.IsBase}\lean{CartanMatrix.IsFiniteTypeModel}
\uses{def:roots}
\assume{$\Phi$ is a Root System in $V$.}
A root base is a linearly independent finite family in $\Phi$ in which every
root has integer coefficients all nonnegative or all nonpositive. Its matrix is
\begin{equation}\label{eq:pairing}
 M(i,j)=p(i,j):=2(\alpha_i,\alpha_j)/\rsq{\alpha_j}.
\end{equation}
The second root supplies the denominator. A finite integer matrix is finite
Cartan if its diagonal is 2, its other entries are nonpositive, its zeros
are symmetric, and some positive integer diagonal $D$ makes $DM$ symmetric
and positive on all nonzero integer vectors. The permitted models are
\begin{equation}\label{eq:list}
 \begin{gathered}
 \typ A_\ell\ (\ell\geq1),\quad \typ B_\ell\ (\ell\geq2),\quad
 \typ C_\ell\ (\ell\geq3),\quad \typ D_\ell\ (\ell\geq4),\\
 \typ E_6,\ \typ E_7,\ \typ E_8,\ \typ F_4,\ \typ G_2.
 \end{gathered}
\end{equation}
These mean the actual Mathlib matrices, with rank $\ell$ for
\path{CartanMatrix.A} despite its docstring. The model predicate includes
the displayed rank restrictions, and equality after transport to its index type.
Every root base is a real vector-space basis of $V$: its span contains
every root by the coefficient condition, and the roots span $V$ by
Definition~\ref{def:roots}; independence is part of the base definition.
\end{namedresult}

\paragraph{Library interface.}\label{facts:rem-library}
We use Mathlib at commit \path{082e2d37e8b0463410cdb532e111cd43d5a66174}
(4 September 2026), the version pinned in \path{lake-manifest.json}. The following
are the precise library packages used in the proofs; every vector space in
them is over $\mathbb R$ and finite dimensional, and all matrix index sets
are finite with classically chosen decidable equality.
\begingroup\raggedright
\begin{enumerate}
\item \path{Matrix.IsFiniteCartan} and \path{CartanMatrix.A},
\path{CartanMatrix.B}, \path{CartanMatrix.C}, \path{CartanMatrix.D},
\path{CartanMatrix.E}, and the matrices $\typ F_4,\typ G_2$ are in
\path{Mathlib/LinearAlgebra/Matrix/Cartan.lean}.
\path{Matrix.IsIndecomposable} in
\path{Mathlib/LinearAlgebra/Matrix/Block.lean} excludes a reindexing as
$\left(\begin{smallmatrix}A&B\\0&D\end{smallmatrix}\right)$ with both
diagonal block sizes positive. It imposes no nonempty-index hypothesis.
\item \path{Matrix.posDef_iff_dotProduct_mulVec} in
\path{Mathlib/LinearAlgebra/Matrix/PosDef.lean} converts positivity on
finitely supported vectors to the finite quadratic form, over both
$\mathbb Z$ and $\mathbb R$.
\path{Matrix.exists_mulVec_eq_zero_iff} in
\path{Mathlib/LinearAlgebra/Matrix/ToLinearEquiv.lean} identifies a nonzero
kernel vector with zero determinant over any commutative domain.
\path{RingHom.map_det} in
\path{Mathlib/LinearAlgebra/Matrix/Determinant/Basic.lean} commutes
determinants with casts. In the same file, \path{Matrix.det_apply'}
is the signed permutation expansion used for a leaf deletion.
\item \path{isCompact_sphere} in
\path{Mathlib/Topology/MetricSpace/ProperSpace.lean} and
\path{IsCompact.exists_isMaxOn} in
\path{Mathlib/Topology/Order/Compact.lean} give the maximum of a continuous
quadratic form on a nonempty unit sphere. Finite-dimensional Euclidean
spaces are proper. No Perron--Frobenius classification is a premise.
\path{Rat.denseRange_cast}, in
\path{Mathlib/Topology/Algebra/Order/Archimedean.lean}, gives density of
rationals in the reals; \path{DenseRange.piMap}, in
\path{Mathlib/Topology/NhdsWithin.lean}, gives the coordinatewise version.
\item \path{SimpleGraph.IsAcyclic.path_unique} in
\path{Mathlib/Combinatorics/SimpleGraph/Acyclic.lean} gives uniqueness of
paths in an acyclic graph. Connectedness supplies their existence.
The graph definitions are in \path{Connectivity/Connected.lean},
\path{Finite.lean}, and \path{AdjMatrix.lean} under
\path{Mathlib/Combinatorics/SimpleGraph/}.
\item \path{RootPairing} and \path{RootPairing.IsRootSystem} are in
\path{Mathlib/LinearAlgebra/RootSystem/Defs.lean}.
Crystallographic, reduced, and irreducible are defined respectively in
\path{IsValuedIn.lean}, \path{Reduced.lean}, and \path{Irreducible.lean}
in that directory. Irreducibility requires nontrivial root and coroot
spaces and no nonzero proper reflection-invariant submodule in either.
The real bilinear inner product $\texttt{inner}_{\mathtt l}$ is in
\path{Mathlib/Analysis/InnerProductSpace/Basic.lean}.
\path{LinearMap.IsPerfPair.of_injective} in
\path{Mathlib/LinearAlgebra/PerfectPairing/Basic.lean} makes it perfect
from injectivity in both arguments and finite dimensionality.
\item \path{RootPairing.nonempty_base} and \path{RootPairing.Base.mk'} in
\path{Mathlib/LinearAlgebra/RootSystem/BaseExists.lean} require a field
of characteristic zero, finitely many roots, spanning roots and coroots,
crystallographic pairings, and reducedness. The second constructor needs
root independence and membership of each root or its negative in the
additive monoid of the selected roots.
\path{Submodule.span_nat_eq_addSubmonoidClosure} in
\path{Mathlib/LinearAlgebra/Span/Defs.lean} identifies this monoid with
nonnegative integer combinations.
\item \path{RootPairing.Base.support_nonempty} in
\path{Mathlib/LinearAlgebra/RootSystem/Base.lean} needs nonempty roots
and $2\ne0$; \path{RootPairing.Base.toWeightBasis} and
\path{RootPairing.Base.toWeightBasis_apply} need spanning.
\path{Module.Basis.equiv} and \path{Module.Basis.equiv_apply} in
\path{Mathlib/LinearAlgebra/Basis/Defs.lean} extend a bijection of basis
indices to a linear equivalence with the stated values.
\item In \path{Mathlib/LinearAlgebra/RootSystem/CartanMatrix.lean},
\path{RootPairing.Base.cartanMatrix_isFiniteCartan} needs a characteristic
zero domain, finite roots, crystallographic pairings, decidable equality,
and spanning. \path{RootPairing.Base.cartanMatrix_isIndecomposable}
needs a characteristic zero domain, finite roots, crystallographic
pairings, reducedness, and irreducibility.
\path{RootPairing.Base.apply_mem_range_root_of_cartanMatrixEq} needs
finite, spanning, reduced, crystallographic source roots over a
characteristic zero domain; the target needs crystallographic pairings
and a base. Given equality of base matrices and a linear equivalence
taking one base to the other, it maps source roots into target roots.
\path{RootPairing.Base.equivOfCartanMatrixEq} packages both directions
with the basis equivalence above; it additionally requires finite,
spanning, reduced target roots.
\path{RootPairing.Base.algebraMap_cartanMatrixIn_apply} identifies the
integer Cartan entry cast to $\mathbb R$ with the real pairing.
\end{enumerate}
\par\endgroup
The elementary identities of \ref{facts:inner} remain available. Arithmetic
leaves below are substitutions in displayed finite sums or polynomial
identities (the Lean proofs may use \texttt{ring}, \texttt{norm\_num},
or \texttt{linarith}); they do not appeal to Python output.

\begin{namedresult}{Lemma}{(reflection identities)}\label{lem:rem-reflection}
\lean{IsRootSet.reflection_identities}\uses{def:roots,facts:inner}
\assume{$V$ is a real inner product space and $\alpha\ne0$.}
\prove{$s_\alpha$ is an orthogonal linear involution;
 $s_\alpha\alpha=-\alpha$; its $-1$ eigenspace is $\mathbb R\alpha$;
 $s_{-\alpha}=s_\alpha$; and for every linear isometry $f$ onto $W$,
 $s_{f\alpha}=f s_\alpha f^{-1}$. In a Root System these reflections
 permute the roots and $\Phi=-\Phi$.}
\end{namedresult}
\begin{proof}
\pfsketch\ We resolve a vector into its component along $\alpha$ and its
orthogonal component.
\step{rrsplit}{For $t=(\alpha,\Zx)/\rsq\alpha$ and $\Zy=\Zx-t\alpha$,
 $(\alpha,\Zy)=0$.}
\begin{proof}\pf\ Substitute $t$ using bilinearity in \ref{facts:inner}.~\qed\end{proof}
\step{rrformula}{$s_\alpha\Zx=-t\alpha+\Zy$.}
\begin{proof}\pf\ Substitute \stepref{rrsplit} in Definition~\ref{def:roots}.~\qed\end{proof}
\step{rrlinear}{$s_\alpha$ is linear.}
\begin{proof}\pf\ The formula in Definition~\ref{def:roots} is a difference
of linear maps, by \ref{facts:inner}.~\qed\end{proof}
\step{rrinvolution}{$s_\alpha(s_\alpha\Zx)=\Zx$.}
\begin{proof}\pf\ Apply \stepref{rrformula} to the components $-t,\Zy$.~\qed\end{proof}
\step{rrmetric}{$(s_\alpha\Zx,s_\alpha\Zw)=(\Zx,\Zw)$.}
\begin{proof}\pf\ Expand both sides by \stepref{rrsplit},
\stepref{rrformula} and \ref{facts:inner}; both are $tu\rsq\alpha+(\Zy,\Zv)$
for the decompositions $\Zx=t\alpha+\Zy$, $\Zw=u\alpha+\Zv$.~\qed\end{proof}
\step{rreigen}{$s_\alpha\alpha=-\alpha$ and
 $s_\alpha\Zx=-\Zx\iff\Zx\in\mathbb R\alpha$.}
\begin{proof}\pf\ In \stepref{rrformula}, equality to $-\Zx$ is exactly
$2\Zy=0$.~\qed\end{proof}
\step{rrsign}{$s_{-\alpha}=s_\alpha$.}
\begin{proof}\pf\ Substitute $-\alpha$ into Definition~\ref{def:roots}.~\qed\end{proof}
\step{rrconj}{$s_{f\alpha}(f\Zx)=f(s_\alpha\Zx)$.}
\begin{proof}\pf\ Substitute the isometry identities into
Definition~\ref{def:roots}.~\qed\end{proof}
\step{rrroots}{\assume{$\Phi$ is a Root System and $\alpha\in\Phi$.}
\prove{$s_\alpha$ permutes $\Phi$ and $-\alpha\in\Phi$.}}
\begin{proof}\pf\ Restrict \stepref{rrinvolution} and \stepref{rreigen}
using reflection closure in Definition~\ref{def:roots}.~\qed\end{proof}
\qedstep
\begin{proof}\pf\ Combine \stepref{rrlinear}--\stepref{rrroots}.~\qed\end{proof}
\end{proof}

\begin{namedresult}{Lemma}{(a null vector of a nonnegative form)}\label{lem:rem-null}
\lean{Matrix.mulVec_eq_zero_of_quadratic_eq_zero}
\uses{facts:inner}
\assume{\begin{pfenum}
\item $Q$ is real symmetric on a finite index set, and
 $q(\Zx):=\Zx^TQ\Zx\geq0$ for every real vector $\Zx$.
\item $q(\Zy)=0$.
\end{pfenum}}
\prove{$Q\Zy=0$.}
\end{namedresult}
\begin{proof}
\pfsketch\ A quadratic polynomial nonnegative on the line cannot have a
nonzero linear coefficient at a zero.
\step{rnpoly}{For every $\Zz,t$, $0\leq2t\Zz^TQ\Zy+t^2q(\Zz)$.}
\begin{proof}\pf\ Expand $q(\Zy+t\Zz)$ using \asm[1]{0}, \asm[2]{0}.~\qed\end{proof}
\step{rncoeff}{For every $\Zz$, $\Zz^TQ\Zy=0$.}
\begin{proof}
\pflet{$c=\Zz^TQ\Zy$, $d=q(\Zz)\geq0$.}
\step{rnchoice}{\assume{$c\ne0$.}
\prove{For $t=-c/(d+1)$,
 $2tc+t^2d=-c^2(d+2)/(d+1)^2<0$.}}
\begin{proof}\pf\ Substitute $t$, collect the common denominator,
and use $c^2>0$, $d+2>0$, $(d+1)^2>0$.~\qed\end{proof}
\qedstep
\begin{proof}\pf\ The inequality in \stepref{rnchoice} contradicts
\stepref{rnpoly}.~\qed\end{proof}
\end{proof}
\qedstep
\begin{proof}\pf\ Apply \stepref{rncoeff} to each coordinate unit vector. ~\qed\end{proof}
\end{proof}

\begin{namedresult}{Lemma}{(integer positivity)}\label{lem:integer-positive}
\lean{Matrix.posDef_real_of_posDef_int}
\uses{lem:rem-null,facts:rem-library}
\assume{\begin{pfenum}
\item $I$ is finite and $H:I\times I\to\mathbb Z$ is symmetric.
\item $\Zx^TH\Zx>0$ for every nonzero $\Zx\in\mathbb Z^I$.
\end{pfenum}}
\prove{The entrywise real cast of $H$ satisfies
$\Zy^TH\Zy>0$ for every nonzero $\Zy\in\mathbb R^I$.}
\end{namedresult}
\begin{proof}
\pfsketch\ A negative real value would yield a negative integer value
by density and scaling. A real null vector would force an integer null
vector through the determinant. Only the integer matrix is needed here.
\step{riposr}{The form is nonnegative on real vectors.}
\begin{proof}
\step{rinegative}{\assume{Some real vector has negative quadratic value.}
\prove{\textsc{False}.}}
\begin{proof}
\step{ricont}{The negative-value set of $\Zx\mapsto\Zx^TH\Zx$
is nonempty and open.}
\begin{proof}\pf\ A finite sum of products of coordinates is continuous;
use \asm{2} for nonemptiness.~\qed\end{proof}
\step{ridense}{\pick\ a rational vector $\Zv$ with $\Zv^TH\Zv<0$.}
\begin{proof}\pf\ Apply \path{DenseRange.piMap} to
\path{Rat.denseRange_cast} and \stepref{ricont}.~\qed\end{proof}
\step{riclear}{\pick\ an integer $k>0$ with $k\Zv\in\mathbb Z^I$.}
\begin{proof}\pf\ The product of the finitely many positive coordinate
denominators of $\Zv$ is such a $k$.~\qed\end{proof}
\step{riintnegative}{$(k\Zv)^TH(k\Zv)=k^2\Zv^TH\Zv<0$.}
\begin{proof}\pf\ Expand the quadratic form and use
\stepref{ridense}, \stepref{riclear}.~\qed\end{proof}
\qedstep
\begin{proof}\pf\ The vector in \stepref{riintnegative} is nonzero
and integral, contradicting \asm[2]{0} after the order-preserving
integer cast.~\qed\end{proof}
\end{proof}
\qedstep
\begin{proof}\pf\ The contradiction \stepref{rinegative} excludes
every negative value.~\qed\end{proof}
\end{proof}
\step{rinonzero}{\assume{$\Zy\ne0$ and $\Zy^TH\Zy=0$.}
\prove{\textsc{False}.}}
\begin{proof}
\step{rikernel}{$H\Zy=0$.}
\begin{proof}\pf\ Apply Lemma~\ref{lem:rem-null} to \stepref{riposr}
and \asm{1}.~\qed\end{proof}
\step{ridetr}{$\det H=0$ over $\mathbb R$.}
\begin{proof}\pf\ Apply \path{Matrix.exists_mulVec_eq_zero_iff}
to \stepref{rikernel} and $\Zy\ne0$ from \asm{1}.~\qed\end{proof}
\step{ridetz}{$\det H=0$ over $\mathbb Z$.}
\begin{proof}\pf\ Use \path{RingHom.map_det} and injectivity of the
integer cast in \stepref{ridetr}.~\qed\end{proof}
\step{riz}{\pick\ a nonzero $\Zv\in\mathbb Z^I$ with $H\Zv=0$.}
\begin{proof}\pf\ Apply \path{Matrix.exists_mulVec_eq_zero_iff}
over $\mathbb Z$ to \stepref{ridetz}.~\qed\end{proof}
\qedstep
\begin{proof}\pf\ The zero value $\Zv^TH\Zv=0$ from \stepref{riz}
contradicts \asm[2]{0}.~\qed\end{proof}
\end{proof}
\qedstep
\begin{proof}\pf\ Combine nonnegativity \stepref{riposr} and
the exclusion of nonzero null vectors \stepref{rinonzero}.~\qed\end{proof}
\end{proof}

\begin{namedresult}{Lemma}{(connected support)}\label{lem:connected}
\lean{Matrix.isIndecomposable_iff_connected_support}
\uses{facts:rem-library}
\assume{\begin{pfenum}
\item $I$ is finite and nonempty, $R$ is a ring, and $M:I\times I\to R$.
\item $M(i,j)=0\iff M(j,i)=0$ for every $i,j$.
\end{pfenum}}
\prove{The simple graph $i\sim j\iff i\ne j\text{ and }M(i,j)\ne0$
is connected if and only if $M$ is Mathlib-indecomposable.}
\end{namedresult}
\begin{proof}
\pfsketch\ A disconnected component gives two blocks; a block partition
cannot be crossed by a path.
\step{rcforward}{\assume{the support is disconnected.}
\prove{$M$ is decomposable.}}
\begin{proof}
\step{rccomponent}{There is a nonempty proper reachable component $P\subset I$.}
\begin{proof}\pf\ Unfold disconnectedness on the nonempty type $I$.~\qed\end{proof}
\step{rczero}{$M(i,j)=0$ whenever $i\in P$ and $j\notin P$.}
\begin{proof}\pf\ A nonzero entry would give an edge from $i$ to $j$,
contrary to the definition of $P$ in \stepref{rccomponent}.~\qed\end{proof}
\step{rcblocks}{Enumerating $P$ and $I\setminus P$ gives a block diagonal
matrix with both block sizes positive.}
\begin{proof}\pf\ Use \stepref{rczero} and \asm[2]{0} in the four
blocks of the two-set partition. ~\qed\end{proof}
\qedstep
\begin{proof}\pf\ The decomposition in \stepref{rcblocks} violates
\path{Matrix.IsIndecomposable}.~\qed\end{proof}
\end{proof}
\step{rcreverse}{\assume{$M$ has a nontrivial upper block triangular decomposition.}
\prove{Its support is disconnected.}}
\begin{proof}
\step{rcback}{Both cross blocks vanish.}
\begin{proof}\pf\ The lower block vanishes by hypothesis and the upper
one vanishes by \asm[2]{0}.~\qed\end{proof}
\step{rcpath}{Every path starting in the first block stays in that block.}
\begin{proof}\pf\ Induct on path length using \stepref{rcback}.~\qed\end{proof}
\qedstep
\begin{proof}\pf\ Choose one vertex in each positive-size block and apply
\stepref{rcpath}.~\qed\end{proof}
\end{proof}
\qedstep
\begin{proof}\pf\ Contraposition of \stepref{rcforward} and
\stepref{rcreverse} gives the equivalence. ~\qed\end{proof}
\end{proof}

\begin{namedresult}{Definition}{(weighted graphs)}\label{def:rem-weighted}
\lean{Matrix.maxRayleigh}\lean{WeightedGraph}\lean{WeightedGraph.Contains}
\uses{facts:rem-library}
\assume{$I$ is finite and nonempty.}
A weighted graph on $I$ has a real symmetric matrix $B$, zero diagonal,
and nonnegative entries. Its support joins $i,j$ when $B(i,j)>0$.
For any real symmetric matrix on $I$, write
$\lambda(B)=\max_{\rsq\Zx=1}\Zx^TB\Zx$.
For weighted graphs $Y,X$, write $Y\preccurlyeq X$ if an injection $e$
satisfies $B_Y(i,j)\leq B_X(e(i),e(j))$ for every $i,j$.
This specified containment is proper when $e$ is not onto or one inequality
is strict. Edges may be omitted and weights decreased.
\end{namedresult}

\begin{namedresult}{Lemma}{(edge-constant functions)}\label{lem:rem-edge-constant}
\lean{SimpleGraph.Connected.constant_of_edge_eq}
\uses{facts:rem-library}
\assume{\begin{pfenum}
\item $X$ is a connected simple graph on a vertex set $I$.
\item $S$ is any set and $f:I\to S$ satisfies $f(u)=f(v)$ whenever $u\sim v$.
\end{pfenum}}
\prove{For every $u,v\in I$, $f(u)=f(v)$. In particular there is
$c\in S$ with $f(v)=c$ for every $v$.}
\end{namedresult}
\begin{proof}
\pfsketch\ Induct on the walk supplied by connectedness.
\step{recwalk}{For every walk from $u$ to $v$, $f(u)=f(v)$.}
\begin{proof}
\step{recnil}{A walk of length zero has equal endpoint values.}
\begin{proof}\pf\ Its endpoints coincide by the definition of a walk.~\qed\end{proof}
\step{reccons}{An edge $u\sim w$ followed by a walk from $w$ to $v$
with $f(w)=f(v)$ gives $f(u)=f(v)$.}
\begin{proof}\pf\ Compose the edge equality in \asm[2]{0} with
the equality for the remaining walk.~\qed\end{proof}
\qedstep
\begin{proof}\pf\ Walk induction using \stepref{recnil},
\stepref{reccons}.~\qed\end{proof}
\end{proof}
\step{recpair}{For every $u,v$, $f(u)=f(v)$.}
\begin{proof}\pf\ Connectedness in \asm[1]{0} supplies a walk;
apply \stepref{recwalk}.~\qed\end{proof}
\qedstep
\begin{proof}\pf\ Choose $u$ by nonemptiness in connectedness and set
$c=f(u)$ in \stepref{recpair}.~\qed\end{proof}
\end{proof}

\begin{namedresult}{Lemma}{(propagation on a connected support)}\label{lem:rem-propagation}
\lean{WeightedGraph.eq_zero_of_nonneg_eigenvector_zero}
\uses{def:rem-weighted,lem:rem-edge-constant}
\assume{\begin{pfenum}
\item $B$ is a nonnegative real symmetric matrix on a finite nonempty set.
 The simple graph $i\sim j\iff i\ne j\text{ and }B(i,j)>0$ is connected.
\item $B\Zx=t\Zx$, all $\Zx(i)\geq0$, and $\Zx(i_0)=0$.
\end{pfenum}}
\prove{$\Zx=0$.}
\end{namedresult}
\begin{proof}
\pfsketch\ Show that having value zero is constant across every edge.
\step{rpneighbour}{\assume{$\Zx(i)=0$ and $B(i,j)>0$.}
\prove{$\Zx(j)=0$.}}
\begin{proof}
\step{rpsum}{$\sum_k B(i,k)\Zx(k)=0$ with every summand nonnegative.}
\begin{proof}\pf\ Evaluate \asm[2]{0} at $i$.~\qed\end{proof}
\step{rpterm}{$B(i,j)\Zx(j)=0$.}
\begin{proof}\pf\ A nonnegative summand of the zero sum
\stepref{rpsum} is zero. ~\qed\end{proof}
\qedstep
\begin{proof}\pf\ Cancel $B(i,j)>0$ in \stepref{rpterm}.~\qed\end{proof}
\end{proof}
\step{rppath}{Every vertex has value zero.}
\begin{proof}\pf\ By \stepref{rpneighbour} in both directions,
the indicator of $\{i:\Zx(i)=0\}$ agrees across every support edge.
Lemma~\ref{lem:rem-edge-constant} makes it constant; its value at $i_0$
is 1 by \asm[2]{0}.~\qed\end{proof}
\qedstep
\begin{proof}\pf\ Function extensionality applied to \stepref{rppath}.~\qed\end{proof}
\end{proof}

\begin{namedresult}{Lemma}{(Rayleigh quotients)}\label{lem:rayleigh}
\lean{Matrix.IsHermitian.maxRayleigh_eigenvalue}\lean{Matrix.IsHermitian.rayleigh_eq_iff}
\lean{Matrix.IsHermitian.quadratic_le_maxRayleigh}
\lean{Matrix.IsHermitian.exists_nonneg_max_eigenvector}\lean{Matrix.IsHermitian.max_eigenspace_simple}
\uses{def:rem-weighted,lem:rem-null,lem:rem-propagation,facts:rem-library}
\assume{$B$ is real symmetric on a finite nonempty set $I$.}
\prove{For every real vector $\Zy$,
$\Zy^TB\Zy\leq\lambda(B)\rsq\Zy$.
$\lambda(B)$ is an eigenvalue bounding every real eigenvalue.
A nonzero vector attains its Rayleigh quotient exactly when it is a
$\lambda(B)$-eigenvector. If $B\geq0$ there is a nonnegative unit maximizer.
If its support is also connected, this vector is strictly positive and
the $\lambda(B)$-eigenspace is one-dimensional.}
\end{namedresult}
\begin{proof}
\pfsketch\ Use a nonnegative form at the maximum, then absolute values
and the propagation lemma. The nonnegative conclusion does not require
connectedness and is stated for use with a disconnected smaller graph.
\step{rymax}{\pick\ a unit vector $\Zx$ attaining $\lambda=\lambda(B)$.}
\begin{proof}\pf\ Apply the compact maximum facts of
\ref{facts:rem-library} to the continuous quadratic form. ~\qed\end{proof}
\step{rybound}{For every $\Zy$, $\Zy^TB\Zy\leq\lambda\rsq\Zy$.}
\begin{proof}
\step{ryzero}{\case{$\Zy=0$.}}
\begin{proof}\pf\ Both sides vanish by substitution. ~\qed\end{proof}
\step{rynonzero}{\case{$\Zy\ne0$.}}
\begin{proof}\pf\ Apply \stepref{rymax} to $\Zy/\|\Zy\|$ and multiply
by $\|\Zy\|^2>0$. ~\qed\end{proof}
\qedstep
\begin{proof}\pf\ The cases \stepref{ryzero}, \stepref{rynonzero} exhaust
all vectors. ~\qed\end{proof}
\end{proof}
\step{ryattain}{For nonzero $\Zy$, equality in \stepref{rybound} holds
 if and only if $B\Zy=\lambda\Zy$.}
\begin{proof}
\step{rykernel}{Equality implies $(\lambda I-B)\Zy=0$.}
\begin{proof}\pf\ Apply Lemma~\ref{lem:rem-null} to the nonnegative form
\stepref{rybound}.~\qed\end{proof}
\step{rysub}{An eigenvector satisfies $\Zy^TB\Zy=\lambda\rsq\Zy$.}
\begin{proof}\pf\ Substitute the eigenvector equation. ~\qed\end{proof}
\qedstep
\begin{proof}\pf\ Combine \stepref{rykernel}, \stepref{rysub}.~\qed\end{proof}
\end{proof}
\step{ryeigen}{$\lambda$ is an eigenvalue and every real eigenvalue is
at most $\lambda$.}
\begin{proof}
\step{ryexists}{$B\Zx=\lambda\Zx$ with $\Zx\ne0$.}
\begin{proof}\pf\ Apply \stepref{ryattain} to \stepref{rymax}.~\qed\end{proof}
\step{ryother}{\assume{$B\Zv=t\Zv$ and $\Zv\ne0$.}
\prove{$t\leq\lambda$.}}
\begin{proof}\pf\ Substitute in \stepref{rybound} and cancel
$\rsq\Zv>0$.~\qed\end{proof}
\qedstep
\begin{proof}\pf\ Combine \stepref{ryexists}, \stepref{ryother}.~\qed\end{proof}
\end{proof}
\step{ryabs}{\assume{$B\geq0$.}
\prove{$|\Zx|$ is a nonnegative unit maximizer.}}
\begin{proof}
\step{ryabsnorm}{$\rsq{|\Zx|}=1$.}
\begin{proof}\pf\ Use $|x|^2=x^2$ in the coordinate sum of
\stepref{rymax}.~\qed\end{proof}
\step{ryabsform}{$|\Zx|^TB|\Zx|\geq\Zx^TB\Zx$.}
\begin{proof}\pf\ Sum $B(i,j)x(i)x(j)\leq B(i,j)|x(i)||x(j)|$
over ordered pairs, using $B(i,j)\geq0$.~\qed\end{proof}
\qedstep
\begin{proof}\pf\ Sandwich \stepref{ryabsform} between the maximum in
\stepref{rymax} and the bound \stepref{rybound}, using \stepref{ryabsnorm}.~\qed\end{proof}
\end{proof}
\step{rypositive}{\assume{$B\geq0$, its support is connected, and $\Zw\geq0$ is a nonzero maximizing eigenvector.}
\prove{$\Zw$ is positive everywhere.}}
\begin{proof}\pf\ A zero coordinate would force $\Zw=0$ by
Lemma~\ref{lem:rem-propagation}.~\qed\end{proof}
\step{rysimple}{\assume{$B\geq0$ has connected support, $\Zw>0$, and $B\Zw=\lambda\Zw$.}
\prove{Every $\lambda$-eigenvector is a multiple of $\Zw$.}}
\begin{proof}
\pflet{For a $\lambda$-eigenvector $\Zv$, let
$t=\max_i\Zv(i)/\Zw(i)$.}
\step{ryresidual}{$t\Zw-\Zv$ is nonnegative and has a zero coordinate.}
\begin{proof}\pf\ Use the attained finite maximum defining $t$ and
\stepref{rypositive}.~\qed\end{proof}
\step{ryresideigen}{$B(t\Zw-\Zv)=\lambda(t\Zw-\Zv)$.}
\begin{proof}\pf\ Linearity and the two eigenvector equations. ~\qed\end{proof}
\qedstep
\begin{proof}\pf\ Lemma~\ref{lem:rem-propagation} applied to
\stepref{ryresidual}, \stepref{ryresideigen} gives $\Zv=t\Zw$.~\qed\end{proof}
\end{proof}
\qedstep
\begin{proof}\pf\ Combine \stepref{ryattain}, \stepref{ryeigen},
\stepref{ryabs}, \stepref{rypositive}, \stepref{rysimple}.~\qed\end{proof}
\end{proof}

\begin{namedresult}{Lemma}{(Rayleigh quotients under reindexing)}\label{lem:rem-rayleigh-reindex}
\lean{Matrix.IsHermitian.maxRayleigh_reindex}
\uses{def:rem-weighted}
\assume{$B_I,B_J$ are real symmetric matrices on finite nonempty sets
$I,J$, and $e:I\to J$ is a bijection with
$B_J(e(i),e(j))=B_I(i,j)$ for all $i,j$.}
\prove{$\lambda(B_I)=\lambda(B_J)$. In particular isomorphic weighted
graphs have equal largest eigenvalues.}
\end{namedresult}
\begin{proof}
\pfsketch\ Reindex the unit sphere and the quadratic form.
\step{rrinorm}{The bijection $\Zx\mapsto\Zx\circ e$ preserves squared norms.}
\begin{proof}\pf\ Reindex the finite sum of squares along $e$.~\qed\end{proof}
\step{rriform}{$(\Zx\circ e)^TB_I(\Zx\circ e)=\Zx^TB_J\Zx$.}
\begin{proof}\pf\ Reindex the double sum along $e\times e$ using
the entry equality in \asm{0}.~\qed\end{proof}
\qedstep
\begin{proof}\pf\ The sets of values maximized in
Definition~\ref{def:rem-weighted} coincide by \stepref{rrinorm},
\stepref{rriform}.~\qed\end{proof}
\end{proof}

\begin{namedresult}{Lemma}{(strict monotonicity)}\label{lem:monotonicity}
\lean{WeightedGraph.maxRayleigh_mono}\lean{WeightedGraph.maxRayleigh_strictMono}
\lean{Matrix.IsHermitian.eigenvalue_eq_maxRayleigh_of_pos}
\uses{def:rem-weighted,lem:rayleigh,lem:rem-propagation,facts:inner}
\assume{\begin{pfenum}
\item $X$ is a finite connected weighted graph with nonnegative weights.
\item $Y$ is a finite nonempty weighted graph with nonnegative weights,
 and $e$ witnesses $Y\preccurlyeq X$.
\end{pfenum}}
\prove{$\lambda(Y)\leq\lambda(X)$, strictly for proper containment.
For any connected nonnegative real symmetric matrix, a strictly positive
eigenvector has eigenvalue equal to its largest eigenvalue.}
\end{namedresult}
\begin{proof}
\pfsketch\ Extend a nonnegative maximizer by zero; equality precludes
both missing vertices and increased entries.
\step{rmchoose}{\pick\ a nonnegative unit maximizer $\Zv$ of $Y$.}
\begin{proof}\pf\ Apply Lemma~\ref{lem:rayleigh} to $B_Y$.~\qed\end{proof}
\pflet{$\Zx(e(i))=\Zv(i)$ and $\Zx(j)=0$ off the image of $e$.}
\step{rmnorm}{$\rsq\Zx=1$.}
\begin{proof}\pf\ Reindex the sum of squares by the injection $e$ in
\asm[2]{0}, using \stepref{rmchoose}.~\qed\end{proof}
\step{rmcompare}{$\lambda(Y)\leq\Zx^TB_X\Zx\leq\lambda(X)$.}
\begin{proof}
\step{rmlower}{$\lambda(Y)\leq\Zx^TB_X\Zx$.}
\begin{proof}\pf\ Sum the entrywise inequalities in \asm[2]{0}
multiplied by the nonnegative $\Zv(i)\Zv(j)$ from \stepref{rmchoose}.~\qed\end{proof}
\step{rmupper}{$\Zx^TB_X\Zx\leq\lambda(X)$.}
\begin{proof}\pf\ Apply the definition of the maximum to
\stepref{rmnorm}.~\qed\end{proof}
\qedstep
\begin{proof}\pf\ Combine \stepref{rmlower}, \stepref{rmupper}.~\qed\end{proof}
\end{proof}
\step{rmequal}{\assume{$\lambda(Y)=\lambda(X)$.}
\prove{The specified containment is not proper.}}
\begin{proof}
\step{rmmax}{$\Zx$ is a maximizing eigenvector of $X$.}
\begin{proof}\pf\ Apply Lemma~\ref{lem:rayleigh} to equality in
\stepref{rmcompare}.~\qed\end{proof}
\step{rmpos}{$\Zx(j)>0$ for every vertex $j$ of $X$.}
\begin{proof}\pf\ Apply Lemma~\ref{lem:rem-propagation} to
\stepref{rmmax}, \stepref{rmnorm} and nonnegativity of $\Zx$.~\qed\end{proof}
\step{rmsurj}{$e$ is onto.}
\begin{proof}\pf\ A vertex outside the image has $\Zx(j)=0$ by definition,
contradicting \stepref{rmpos}.~\qed\end{proof}
\step{rmentries}{All entry inequalities are equalities.}
\begin{proof}\pf\ In the zero sum
$\sum_{i,j}(B_X(e(i),e(j))-B_Y(i,j))\Zv(i)\Zv(j)=0$
from \stepref{rmcompare}, every summand is nonnegative and both vector
factors are positive by \stepref{rmpos}.~\qed\end{proof}
\qedstep
\begin{proof}\pf\ Use the definition of properness with
\stepref{rmsurj}, \stepref{rmentries}.~\qed\end{proof}
\end{proof}
\step{rmposeigen}{\assume{$B$ is real symmetric, nonnegative, and has connected off-diagonal positive support; $B\Zv=t\Zv$ and $\Zv>0$.}
\prove{$t=\lambda(B)$.}}
\begin{proof}
\step{rmw}{\pick\ $\Zw>0$ with $B\Zw=\lambda(B)\Zw$.}
\begin{proof}\pf\ Apply Lemma~\ref{lem:rayleigh} to this matrix $B$.~\qed\end{proof}
\step{rmorth}{$(t-\lambda(B))(\Zv,\Zw)=0$.}
\begin{proof}\pf\ Use symmetry of $B$ in $(B\Zv,\Zw)=(\Zv,B\Zw)$
and \stepref{rmw}.~\qed\end{proof}
\qedstep
\begin{proof}\pf\ Cancel the positive sum $(\Zv,\Zw)$ in
\stepref{rmorth}.~\qed\end{proof}
\end{proof}
\qedstep
\begin{proof}\pf\ Combine \stepref{rmcompare}, the contrapositive of
\stepref{rmequal}, and \stepref{rmposeigen}.~\qed\end{proof}
\end{proof}

\begin{namedresult}{Lemma}{(symmetrization)}\label{lem:symmetrization}
\lean{Matrix.IsFiniteCartan.symmetrization}
\uses{def:rem-bases,def:rem-weighted,lem:integer-positive,lem:connected,lem:rayleigh,facts:rem-library}
\assume{\begin{pfenum}
\item $I$ is finite and nonempty, $M:I\times I\to\mathbb Z$ is finite
 Cartan and indecomposable.
\item $d:I\to\mathbb Z_{>0}$ symmetrizes $M$ as in its finite Cartan
 hypothesis, and $D=\operatorname{diag}(d)$.
\end{pfenum}}
\prove{$S=D^{1/2}MD^{-1/2}$ is symmetric positive definite and equals
$2I-B$, where $B(i,j)=\sqrt{M(i,j)M(j,i)}$ off the diagonal and $B(i,i)=0$.
The weighted graph of $B$ is connected, has edge weights
$1,\sqrt2,\sqrt3$, and $\lambda(B)<2$.}
\end{namedresult}
\begin{proof}
\pfsketch\ Congruence transfers positivity; a two-coordinate test bounds
the integer edge products.
\step{rsreal}{$DM$ is positive definite over $\mathbb R$.}
\begin{proof}\pf\ Apply Lemma~\ref{lem:integer-positive} and the
quadratic-form interface in \ref{facts:rem-library}.~\qed\end{proof}
\step{rscong}{$S=D^{-1/2}(DM)D^{-1/2}$ is symmetric positive definite.}
\begin{proof}\pf\ Substitute the invertible diagonal change of coordinates
in the form of \stepref{rsreal}.~\qed\end{proof}
\step{rsdiag}{$S(i,i)=2$ for every $i$.}
\begin{proof}\pf\ The diagonal factors cancel and $M(i,i)=2$ by
\asm[1]{0}.~\qed\end{proof}
\step{rsentry}{For $i\ne j$, $S(i,j)\leq0$ and
$S(i,j)^2=M(i,j)M(j,i)$.}
\begin{proof}\pf\ Substitute $S(i,j)=\sqrt{d(i)/d(j)}M(i,j)$ and
$d(i)M(i,j)=d(j)M(j,i)$ from \asm[2]{0}.~\qed\end{proof}
\step{rsbound}{For $i\ne j$, $S(i,j)^2<4$.}
\begin{proof}
\pflet{$b=S(i,j)$ and $\Zx=\Zu_i-(b/2)\Zu_j\ne0$.}
\step{rstest}{$\Zx^TS\Zx=2-b^2/2$.}
\begin{proof}\pf\ Expand the two coordinates using \stepref{rsdiag}
and symmetry in \stepref{rscong}.~\qed\end{proof}
\qedstep
\begin{proof}\pf\ Positivity in \stepref{rscong} and \stepref{rstest}
give $2-b^2/2>0$.~\qed\end{proof}
\end{proof}
\step{rsweights}{$M(i,j)M(j,i)\in\{0,1,2,3\}$.}
\begin{proof}\pf\ Apply integer interval arithmetic to
\stepref{rsentry}, \stepref{rsbound}.~\qed\end{proof}
\step{rsgraph}{$B=2I-S$ has precisely the required entries and weights.}
\begin{proof}\pf\ The nonpositive square root is unique in
\stepref{rsentry}; use \stepref{rsdiag}, \stepref{rsweights}.~\qed\end{proof}
\step{rsconn}{Its support is connected.}
\begin{proof}\pf\ Apply Lemma~\ref{lem:connected} to the identical
nonzero support in \stepref{rsgraph}.~\qed\end{proof}
\step{rslambda}{$\lambda(B)<2$.}
\begin{proof}
\step{rsmax}{\pick\ a unit maximizer $\Zv$ for $B$.}
\begin{proof}\pf\ Apply Lemma~\ref{lem:rayleigh}.~\qed\end{proof}
\qedstep
\begin{proof}\pf\ At \stepref{rsmax},
$\lambda(B)=2-\Zv^TS\Zv<2$ by \stepref{rscong}, \stepref{rsgraph}.~\qed\end{proof}
\end{proof}
\qedstep
\begin{proof}\pf\ Combine \stepref{rscong}, \stepref{rsgraph},
\stepref{rsconn}, \stepref{rslambda}.~\qed\end{proof}
\end{proof}

\subsubsection{Affine obstructions and labelled trees}

\begin{namedresult}{Definition}{(affine obstructions)}\label{def:affine}
\lean{WeightedGraph.IsAffineObstruction}
\lean{SimpleGraph.smithA}\lean{SimpleGraph.smithD}
\lean{SimpleGraph.smithE₆}\lean{SimpleGraph.smithE₇}\lean{SimpleGraph.smithE₈}
\uses{def:rem-weighted}
\assume{All lengths below are integers in the displayed ranges.}
The tree $T(a,b,c)$, $a,b,c\geq1$, has vertex set
\[
 V_{a,b,c}=\{*\}\sqcup(\{1\}\times[1,a])\sqcup
 (\{2\}\times[1,b])\sqcup(\{3\}\times[1,c]).
\]
Its only edges are $*\sim(r,1)$ and consecutive positions on an arm.
The affine obstructions are the following weighted graphs:
\begin{enumerate}
\item Unit cycles $\Cyc_m$, $m\geq3$.
\item $\widetilde{\typ D}_4=K_{1,4}$; and $\widetilde{\typ D}_\ell$,
$\ell\geq5$, a unit path of $\ell-4$ edges with two new leaves at each end.
\item $\widetilde{\typ E}_6=T(2,2,2)$,
$\widetilde{\typ E}_7=T(1,3,3)$, $\widetilde{\typ E}_8=T(1,2,5)$,
all with unit weights.
\item $X^{\rm end}_m$, $m\geq3$: a path on $m$ vertices with first and
last weights $\sqrt2$ and all intervening weights 1.
\item $X^{\rm fork}_k$, $k\geq1$: a path $v_0,\ldots,v_k$ with two
unit leaves at $v_0$, last path weight $\sqrt2$, and other path weights 1.
\item $X^{\rm mid}$: the path with weights $1,1,\sqrt2,1$.
\item $X^{\rm triple}$: the path with weights $\sqrt3,1$.
\end{enumerate}
Thus both edges of $X^{\rm end}_3$ have weight $\sqrt2$, and the weighted
edge of $X^{\rm fork}_1$ meets the branch vertex. The project model
\texttt{smithA} at parameter $\ell\geq2$ means $\Cyc_{\ell+1}$.
\end{namedresult}

\begin{namedresult}{Lemma}{(connected labelled graphs)}\label{lem:rem-labelled-connected}
\lean{WeightedGraph.connected_of_path_or_arm}
\lean{WeightedGraph.IsAffineObstruction.connected}
\uses{def:affine,def:rem-weighted}
\assume{$X$ is a finite weighted graph whose support admits a nonempty
path labelling, an arm labelling $T(a,b,c)$ with $a,b,c\geq1$, or is the
support of an obstruction in Definition~\ref{def:affine}.}
\prove{The support of $X$ is connected.}
\end{namedresult}
\begin{proof}
\pfsketch\ Give a walk from every vertex to one chosen vertex.
\step{rconnpath}{\case{The support has a path labelling.}}
\begin{proof}\pf\ Consecutive positions $i,i-1,\ldots,0$ give a walk
from every vertex to position 0, including the singleton path.~\qed\end{proof}
\step{rconnarm}{\case{The support has an arm labelling.}}
\begin{proof}\pf\ From $(r,t)$ follow $(r,t-1),\ldots,(r,1),*$;
the centre has the length-zero walk to itself.~\qed\end{proof}
\step{rconncycle}{\case{$X$ has the support of $\Cyc_m$.}}
\begin{proof}\pf\ Follow consecutive cycle positions down to 0;
each is an edge by Definition~\ref{def:affine}.~\qed\end{proof}
\step{rconnstar}{\case{$X$ has the support of $\widetilde{\typ D}_4$.}}
\begin{proof}\pf\ Every leaf has an edge to the centre of $K_{1,4}$.~\qed\end{proof}
\step{rconnD}{\case{$X$ has the support of $\widetilde{\typ D}_\ell$,
$\ell\geq5$.}}
\begin{proof}\pf\ A leaf first reaches its attached endpoint; from
any path position, consecutive path edges reach the first endpoint.~\qed\end{proof}
\step{rconnfork}{\case{$X$ has the support of $X^{\rm fork}_k$.}}
\begin{proof}\pf\ Its extra leaves reach $v_0$ in one edge; $v_i$
reaches it along $v_i,v_{i-1},\ldots,v_0$.~\qed\end{proof}
\qedstep
\begin{proof}\pf\ The $\widetilde{\typ E}$ supports are arm trees;
$X^{\rm end},X^{\rm mid},X^{\rm triple}$ have path supports.
Thus \stepref{rconnpath}--\stepref{rconnfork} exhaust
Definition~\ref{def:affine} and the two other hypotheses.~\qed\end{proof}
\end{proof}

\begin{namedresult}{Lemma}{(unit containment and support embeddings)}\label{lem:rem-unit-containment}
\lean{WeightedGraph.unit_contains_iff_support_embedding}
\uses{def:rem-weighted}
\assume{$X$ is a finite weighted graph with $B_X(u,v)\geq1$ on every
support edge, and $Y$ is a finite nonempty simple graph, given unit edge weights.}
\prove{$Y\preccurlyeq X$ if and only if there is an injection
$e:V(Y)\to V(X)$ sending every edge of $Y$ to a support edge of $X$.
The same injection witnesses either assertion; nonedges need not be preserved.}
\end{namedresult}
\begin{proof}
\pfsketch\ Check the containment inequality on edges and nonedges.
\step{rucforward}{\assume{$e$ witnesses $Y\preccurlyeq X$.}
\prove{$e$ preserves edges into the support.}}
\begin{proof}\pf\ On an edge the inequality is
$1\leq B_X(e(u),e(v))$, so that entry is positive.~\qed\end{proof}
\step{rucreverse}{\assume{$e$ is an injective edge-preserving map to the support.}
\prove{$e$ witnesses $Y\preccurlyeq X$.}}
\begin{proof}
\step{rucedge}{On an edge, $B_Y(u,v)=1\leq B_X(e(u),e(v))$.}
\begin{proof}\pf\ Apply the weight lower bound in \asm{0} to
the support edge supplied by \asm{1}.~\qed\end{proof}
\step{rucnonedge}{On a nonedge, $B_Y(u,v)=0\leq B_X(e(u),e(v))$.}
\begin{proof}\pf\ Apply nonnegativity from
Definition~\ref{def:rem-weighted}.~\qed\end{proof}
\qedstep
\begin{proof}\pf\ The cases \stepref{rucedge}, \stepref{rucnonedge}
give every inequality of containment.~\qed\end{proof}
\end{proof}
\qedstep
\begin{proof}\pf\ Combine \stepref{rucforward}, \stepref{rucreverse}.~\qed\end{proof}
\end{proof}

\begin{namedresult}{Lemma}{(a cycle gives a cycle subgraph)}\label{lem:rem-cycle-embedding}
\lean{SimpleGraph.Walk.IsCycle.exists_cycle_embedding}
\uses{def:affine,facts:rem-library}
\assume{$Z$ is a simple graph and $p$ is a closed walk of length $m$
which is a cycle.}
\prove{$m\geq3$, and the map sending $i\in[0,m-1]$ to the $i$th
vertex of $p$ is an injection preserving every edge of $\Cyc_m$.}
\end{namedresult}
\begin{proof}
\pfsketch\ Omit the repeated last vertex from the cycle walk.
\step{rcylength}{$m\geq3$ and the first $m$ vertices of $p$ are distinct.}
\begin{proof}
\step{rcythree}{$m\geq3$.}
\begin{proof}\raggedright\pf\ Apply
\path{SimpleGraph.Walk.isCycle_iff_isPath_tail_and_le_length}
in \path{Mathlib/Combinatorics/SimpleGraph/Paths.lean}.~\qed\end{proof}
\step{rcyinjective}{The first $m$ vertices are distinct.}
\begin{proof}\pf\ Apply \path{SimpleGraph.Walk.IsCycle.getVert_injOn'}
in that file to the positions $0,\ldots,m-1$.~\qed\end{proof}
\qedstep
\begin{proof}\pf\ Combine \stepref{rcythree}, \stepref{rcyinjective}.~\qed\end{proof}
\end{proof}
\step{rcyedges}{Successive positions and the pair $m-1,0$ map to edges.}
\begin{proof}\pf\ These are the walk edges of $p$, including its
last edge back to its start.~\qed\end{proof}
\qedstep
\begin{proof}\pf\ Apply \stepref{rcylength}, \stepref{rcyedges}
to the cycle definition in Definition~\ref{def:affine}.~\qed\end{proof}
\end{proof}

\begin{namedresult}{Lemma}{(affine eigenvectors)}\label{lem:affine}
\lean{WeightedGraph.IsAffineObstruction.maxRayleigh_eq_two}
\uses{def:affine,lem:monotonicity,lem:rem-labelled-connected}
\assume{$X$ is one of the graphs in Definition~\ref{def:affine}.}
\prove{$\lambda(X)=2$. The following positive vectors $\Zw$ satisfy
$B\Zw=2\Zw$: 1 on a cycle; 2 on each nonleaf and 1 on each leaf of
$\widetilde{\typ D}$; centre $h=3,4,6$ respectively on the three
$\widetilde{\typ E}$ trees and value $h(a+1-t)/(a+1)$ at position $t$
of an arm of length $a$; endpoints 1 and interiors $\sqrt2$ for
$X^{\rm end}$; values $2,\ldots,2,\sqrt2$ on the path and 1 on the
two extra leaves for $X^{\rm fork}$; $(1,2,3,2\sqrt2,\sqrt2)$ for
$X^{\rm mid}$; $(\sqrt3,2,1)$ for $X^{\rm triple}$.}
\end{namedresult}
\begin{proof}
\pfsketch\ Check each vertex type. This proves the unbounded families,
including their shortest instances.
\pflet{$\Zw$ is the vector specified in the statement for the graph at hand.}
\step{rafcycle}{On a cycle, each weighted neighbour sum is 2.}
\begin{proof}\pf\ Each vertex has two unit neighbours by
Definition~\ref{def:affine}.~\qed\end{proof}
\step{rafD}{The displayed vector solves the equations on
$\widetilde{\typ D}$.}
\begin{proof}
\step{rafDleaf}{At a leaf the sum is $2=2\cdot1$.}
\begin{proof}\pf\ Its unique neighbour has value 2 by the stated vector. ~\qed\end{proof}
\step{rafDcentre}{At the centre of $K_{1,4}$ the sum is $4=2\cdot2$.}
\begin{proof}\pf\ Add its four leaf values from the definition of $\Zw$.~\qed\end{proof}
\step{rafDfork}{At a fork for $\ell\geq5$ the sum is $2+1+1=2\cdot2$.}
\begin{proof}\pf\ Use its one path neighbour and two leaves in
Definition~\ref{def:affine}.~\qed\end{proof}
\step{rafDinternal}{At a path interior the sum is $2+2=2\cdot2$.}
\begin{proof}\pf\ Use its two path neighbours in Definition~\ref{def:affine}.~\qed\end{proof}
\qedstep
\begin{proof}\pf\ The leaf, centre, fork and interior cases
\stepref{rafDleaf}--\stepref{rafDinternal} exhaust the vertices. ~\qed\end{proof}
\end{proof}
\step{rafE}{The displayed vectors solve the equations on
$\widetilde{\typ E}_6,\widetilde{\typ E}_7,\widetilde{\typ E}_8$.}
\begin{proof}
\step{rafEarm}{At a noncentral vertex the sum is twice its value.}
\begin{proof}\pf\ The arm formula is an arithmetic progression with
value $h$ at position 0 and virtual value 0 at position $a+1$.~\qed\end{proof}
\step{rafEcentre}{For arms $a,b,c$ and centre value $h>0$, the central
eigenvector equation is equivalent to
$1/(a+1)+1/(b+1)+1/(c+1)=1$.}
\begin{proof}\pf\ Its neighbour sum is
$h(a/(a+1)+b/(b+1)+c/(c+1))$; cancel $h$ from equality to $2h$.~\qed\end{proof}
\step{rafE6}{The central equation holds for $(2,2,2)$.}
\begin{proof}\pf\ Substitute $1/3+1/3+1/3=1$ in
\stepref{rafEcentre}.~\qed\end{proof}
\step{rafE7}{The central equation holds for $(1,3,3)$.}
\begin{proof}\pf\ Substitute $1/2+1/4+1/4=1$ in \stepref{rafEcentre}.~\qed\end{proof}
\step{rafE8}{The central equation holds for $(1,2,5)$.}
\begin{proof}\pf\ Substitute $1/2+1/3+1/6=1$ in \stepref{rafEcentre}.~\qed\end{proof}
\qedstep
\begin{proof}\pf\ Combine \stepref{rafEarm} with the three central
cases \stepref{rafE6}--\stepref{rafE8}.~\qed\end{proof}
\end{proof}
\step{rafend}{The equations hold on $X^{\rm end}_m$.}
\begin{proof}
\step{rafendleaf}{At either endpoint the sum is $\sqrt2\sqrt2=2$.}
\begin{proof}\pf\ Substitute the endpoint edge and neighbouring value. ~\qed\end{proof}
\step{rafendinternal}{At an interior vertex each of its two sides
contributes $\sqrt2$.}
\begin{proof}
\step{rafendboundary}{A neighbouring endpoint contributes
$\sqrt2\cdot1=\sqrt2$.}
\begin{proof}\pf\ Substitute the endpoint value and edge weight. ~\qed\end{proof}
\step{rafendmiddle}{A neighbouring interior contributes
$1\cdot\sqrt2=\sqrt2$.}
\begin{proof}\pf\ Substitute the interior value and edge weight. ~\qed\end{proof}
\qedstep
\begin{proof}\pf\ Every neighbour is an endpoint or an interior vertex,
so \stepref{rafendboundary}, \stepref{rafendmiddle} apply to each side. ~\qed\end{proof}
\end{proof}
\qedstep
\begin{proof}\pf\ Add the two contributions in
\stepref{rafendinternal}, together with \stepref{rafendleaf}.~\qed\end{proof}
\end{proof}
\step{raffork}{The equations hold on $X^{\rm fork}_k$.}
\begin{proof}
\step{rafforkleaf}{An extra leaf has sum $2=2\cdot1$.}
\begin{proof}\pf\ Its neighbour is $v_0$ with value 2. ~\qed\end{proof}
\step{rafforklast}{At $v_k$ the sum is $\sqrt2\cdot2=2\sqrt2$.}
\begin{proof}\pf\ Substitute its unique edge. ~\qed\end{proof}
\step{rafforksmall}{For $k=1$, the sum at $v_0$ is
$1+1+\sqrt2\sqrt2=4$.}
\begin{proof}\pf\ Substitute all three incident edges from
Definition~\ref{def:affine}.~\qed\end{proof}
\step{rafforklarge}{For $k>1$, the sum at $v_0$ is $1+1+2=4$.}
\begin{proof}\pf\ The edge $v_0v_1$ has weight 1 and $v_1$ has value 2. ~\qed\end{proof}
\step{rafforkpenultimate}{For $k>1$, the sum at $v_{k-1}$ is
$2+\sqrt2\sqrt2=4$.}
\begin{proof}\pf\ Substitute the two incident path edges. ~\qed\end{proof}
\step{rafforkinterior}{At $v_i$ with $0<i<k-1$, the sum is $2+2=4$.}
\begin{proof}\pf\ Both neighbours have value 2 and both edges weight 1. ~\qed\end{proof}
\qedstep
\begin{proof}\pf\ The six cases \stepref{rafforkleaf}--\stepref{rafforkinterior}
exhaust the vertices for $k\geq1$.~\qed\end{proof}
\end{proof}
\step{raffinite}{The equations hold on $X^{\rm mid},X^{\rm triple}$.}
\begin{proof}
\step{rafmid0}{At position 0 of $X^{\rm mid}$ the sum is $2=2\cdot1$.}
\begin{proof}\pf\ Substitute its sole neighbour. ~\qed\end{proof}
\step{rafmid1}{At position 1 the sum is $1+3=2\cdot2$.}
\begin{proof}\pf\ Substitute its two unit edges. ~\qed\end{proof}
\step{rafmid2}{At position 2 the sum is $2+\sqrt2(2\sqrt2)=2\cdot3$.}
\begin{proof}\pf\ Substitute its unit and double-weight edges. ~\qed\end{proof}
\step{rafmid3}{At position 3 the sum is $3\sqrt2+\sqrt2=2(2\sqrt2)$.}
\begin{proof}\pf\ Substitute its double-weight and unit edges. ~\qed\end{proof}
\step{rafmid4}{At position 4 the sum is $2\sqrt2=2\cdot\sqrt2$.}
\begin{proof}\pf\ Substitute its sole neighbour. ~\qed\end{proof}
\step{raftriple0}{At position 0 of $X^{\rm triple}$ the sum is $2\sqrt3$.}
\begin{proof}\pf\ Its neighbour has value 2 and its edge weight is $\sqrt3$.~\qed\end{proof}
\step{raftriple1}{At position 1 the sum is $\sqrt3\sqrt3+1=2\cdot2$.}
\begin{proof}\pf\ Substitute its two neighbours. ~\qed\end{proof}
\step{raftriple2}{At position 2 the sum is $2=2\cdot1$.}
\begin{proof}\pf\ Substitute its sole neighbour. ~\qed\end{proof}
\qedstep
\begin{proof}\pf\ Positions \stepref{rafmid0}--\stepref{raftriple2}
exhaust both vertex sets. ~\qed\end{proof}
\end{proof}
\step{rafpositive}{Every displayed entry is positive.}
\begin{proof}\pf\ The constants are positive; on an arm $1\leq t\leq a$
gives $a+1-t\geq1$ in the displayed formula. ~\qed\end{proof}
\qedstep
\begin{proof}\pf\ Apply the positive-eigenvector assertion of
Lemma~\ref{lem:monotonicity} to \stepref{rafcycle}--\stepref{raffinite}
and \stepref{rafpositive}, with connectedness from
Lemma~\ref{lem:rem-labelled-connected}.~\qed\end{proof}
\end{proof}

\begin{namedresult}{Lemma}{(path labelling)}\label{lem:path-labelling}
\lean{SimpleGraph.nonempty_iso_pathGraph}\lean{SimpleGraph.exists_path_labelling_from}
\lean{SimpleGraph.path_labelling_reverse}
\uses{facts:rem-library}
\assume{$X$ is a finite nonempty connected acyclic simple graph on $m$
vertices and every degree is at most 2.}
\prove{There is a bijection $e:[0,m-1]\to V(X)$ with
$e(i)\sim e(j)\iff |i-j|=1$. If $w$ is any specified vertex of degree
at most 1, such a bijection can be chosen with $e(0)=w$.
For every path labelling $e$, its reversal $i\mapsto e(m-1-i)$
is also a path labelling.}
\end{namedresult}
\begin{proof}
\pfsketch\ Take a longest simple path; no edge can leave its support.
\step{rlpick}{\pick\ a simple path $v_0,\ldots,v_k$ of maximum length.}
\begin{proof}\pf\ Distinct-vertex lists in a finite set form a finite
set; singleton lists give a nonempty set of paths. ~\qed\end{proof}
\step{rlend0}{No neighbour of $v_0$ lies off this path.}
\begin{proof}\pf\ Prepending that neighbour would extend
\stepref{rlpick} as a simple path. ~\qed\end{proof}
\step{rlendk}{No neighbour of $v_k$ lies off this path.}
\begin{proof}\pf\ Appending that neighbour would extend
\stepref{rlpick} as a simple path. ~\qed\end{proof}
\step{rlinternal}{No neighbour of any $v_i$, $0<i<k$, lies off the path.}
\begin{proof}\pf\ Such a vertex already has the two distinct neighbours
$v_{i-1},v_{i+1}$, so a third violates the degree bound in \asm{0}.~\qed\end{proof}
\step{rlonto}{Every vertex belongs to the path.}
\begin{proof}
\step{rlcross}{A path from a vertex outside the support to $v_0$ would
have an edge with one endpoint inside and one outside.}
\begin{proof}\pf\ Take the first entry into the support in its finite
vertex sequence. ~\qed\end{proof}
\qedstep
\begin{proof}\pf\ Connectedness and \stepref{rlcross} contradict
\stepref{rlend0}, \stepref{rlendk}, \stepref{rlinternal}.~\qed\end{proof}
\end{proof}
\step{rlbijection}{$k+1=m$ and $e(i)=v_i$ is a bijection.}
\begin{proof}\pf\ The path is injective by \stepref{rlpick} and
surjective by \stepref{rlonto}.~\qed\end{proof}
\step{rledges}{The only edges are consecutive pairs.}
\begin{proof}
\step{rlconsecutive}{Every consecutive pair is adjacent.}
\begin{proof}\pf\ Unfold the path in \stepref{rlpick}.~\qed\end{proof}
\step{rlchord}{No two nonconsecutive positions are adjacent.}
\begin{proof}\pf\ Their path segment and their single edge would be
distinct simple paths, contradicting \path{SimpleGraph.IsAcyclic.path_unique}.~\qed\end{proof}
\qedstep
\begin{proof}\pf\ Combine \stepref{rlconsecutive}, \stepref{rlchord}
and the no-loop condition. ~\qed\end{proof}
\end{proof}
\step{rlreversal}{The reversal of every path labelling is a path labelling.}
\begin{proof}\pf\ The involution $i\mapsto m-1-i$ preserves
$|i-j|=1$, so composing with it preserves the defining edge equivalence.~\qed\end{proof}
\step{rlstart}{A prescribed vertex $w$ of degree at most 1 can be first.}
\begin{proof}
\step{rlone}{\case{$m=1$.}}
\begin{proof}\pf\ The bijection \stepref{rlbijection} has sole value $w$.~\qed\end{proof}
\step{rlmany}{\case{$m\geq2$.}}
\begin{proof}
\step{rlendpoint}{$w=v_0$ or $w=v_k$.}
\begin{proof}\pf\ All other positions have two neighbours by
\stepref{rledges}.~\qed\end{proof}
\step{rlforward}{\assume{$w=v_0$.}
\prove{$e$ is a path labelling with $e(0)=w$.}}
\begin{proof}\pf\ Its first value is $w$ by substitution. ~\qed\end{proof}
\step{rlreverse}{\assume{$w=v_k$.}
\prove{$i\mapsto e(m-1-i)$ is a path labelling whose first value is $w$.}}
\begin{proof}\pf\ Use \stepref{rlreversal}; the reversed first
value is $e(m-1)=v_k=w$.~\qed\end{proof}
\qedstep
\begin{proof}\pf\ Use \stepref{rlendpoint} to select
\stepref{rlforward} or \stepref{rlreverse}.~\qed\end{proof}
\end{proof}
\qedstep
\begin{proof}\pf\ Nonemptiness gives $m=1$ or $m\geq2$, the cases
\stepref{rlone}, \stepref{rlmany}.~\qed\end{proof}
\end{proof}
\qedstep
\begin{proof}\pf\ Combine \stepref{rlbijection}, \stepref{rledges},
\stepref{rlstart}, \stepref{rlreversal}.~\qed\end{proof}
\end{proof}

\begin{namedresult}{Lemma}{(three-arm labelling)}\label{lem:arm-labelling}
\lean{SimpleGraph.exists_three_arm_labelling}
\uses{def:affine,lem:path-labelling,facts:rem-library}
\assume{$X$ is a finite connected acyclic simple graph, $v$ has degree 3,
and all its other vertices have degree at most 2.}
\prove{There are positive integers $a\leq b\leq c$ and a bijection
$e:V_{a,b,c}\to V(X)$ with $e(*)=v$ whose edges are exactly the centre
edges and the consecutive arm edges of Definition~\ref{def:affine}.}
\end{namedresult}
\begin{proof}
\pfsketch\ Partition the other vertices by the first edge of the unique
path from $v$, then apply the endpoint version of path labelling to each part.
\step{raunique}{There is a unique simple path from $v$ to each vertex.}
\begin{proof}\pf\ Combine connectedness with
\path{SimpleGraph.IsAcyclic.path_unique}.~\qed\end{proof}
\pflet{Enumerate the neighbours of $v$ as $w_1,w_2,w_3$.
Let $P_r$ consist of vertices whose path from $v$ starts with $vw_r$.}
\step{rapartition}{$P_1,P_2,P_3$ are disjoint nonempty sets partitioning
$V(X)\setminus\{v\}$.}
\begin{proof}\pf\ Each nonconstant path in \stepref{raunique} has one
first edge, and $w_r\in P_r$.~\qed\end{proof}
\step{raprefix}{For every $r$ and $z\in P_r$, the segment from $w_r$ to
$z$ stays inside $P_r$.}
\begin{proof}\pf\ Each prefix is the unique path to its endpoint by
\stepref{raunique}.~\qed\end{proof}
\step{ranocross}{There is no edge between distinct $P_r,P_s$.}
\begin{proof}\pf\ Such an edge would give two paths from $v$ to an
endpoint, contradicting \stepref{raunique}; their prefixes cannot meet
after $v$ by that same uniqueness. ~\qed\end{proof}
\step{raattachment}{The only neighbour of $v$ in $P_r$ is $w_r$.}
\begin{proof}\pf\ The unique path to any neighbour of $v$ is its
single edge, by \stepref{raunique}.~\qed\end{proof}
\step{raarm}{For each $r\in\{1,2,3\}$ the induced graph on $P_r$ admits
a path labelling starting at $w_r$.}
\begin{proof}
\step{raarmconn}{The induced graph on $P_r$ is connected and acyclic.}
\begin{proof}\pf\ Use \stepref{raprefix} for reachability; an induced
cycle would be a cycle in $X$.~\qed\end{proof}
\step{raarmdegree}{All its degrees are at most 2, and $w_r$ has degree
at most 1 within it.}
\begin{proof}\pf\ Deleting $v$ removes the edge $vw_r$ from a vertex
with degree at most 2; every remaining vertex also differs from $v$.~\qed\end{proof}
\qedstep
\begin{proof}\pf\ Apply Lemma~\ref{lem:path-labelling} to
\stepref{raarmconn}, \stepref{raarmdegree}, and $w_r\in P_r$.~\qed\end{proof}
\end{proof}
\step{raassemble}{These three bijections and $*\mapsto v$ give an arm
bijection with all the asserted edges.}
\begin{proof}\pf\ Take their disjoint union using \stepref{rapartition},
\stepref{ranocross}, \stepref{raattachment}, \stepref{raarm}.~\qed\end{proof}
\step{raorder}{The three positive lengths may be ordered increasingly.}
\begin{proof}\pf\ Permute the three labelled parts of
\stepref{raassemble} by a sorting permutation. ~\qed\end{proof}
\qedstep
\begin{proof}\pf\ Use \stepref{raassemble}, \stepref{raorder}.~\qed\end{proof}
\end{proof}

\begin{namedresult}{Lemma}{(unweighted trees without obstructions)}\label{lem:unweighted-trees}
\lean{SimpleGraph.dynkin_labelling_of_no_affine}
\uses{def:affine,lem:path-labelling,lem:arm-labelling,lem:rem-cycle-embedding,facts:rem-library}
\assume{$X$ is a finite connected simple graph containing no cycle,
no $\widetilde{\typ D}_\ell$ ($\ell\geq4$), and no
$\widetilde{\typ E}_6,\widetilde{\typ E}_7,\widetilde{\typ E}_8$ as a
subgraph (the subgraph need not be induced).}
\prove{$X$ has a path labelling, or an arm labelling of $T(1,1,c)$,
$c\geq1$, or of $T(1,2,2)$, $T(1,2,3)$, $T(1,2,4)$.}
\end{namedresult}
\begin{proof}
\pfsketch\ Exclude high degrees and two branch vertices, then truncate
the ordered arms.
\step{rutree}{$X$ is acyclic.}
\begin{proof}\pf\ A closed walk satisfying the cycle predicate would
embed $\Cyc_m$ for $m\geq3$ by Lemma~\ref{lem:rem-cycle-embedding},
contrary to \asm{0}; this is the definition of acyclicity.~\qed\end{proof}
\step{rudegree}{Every degree is at most 3.}
\begin{proof}\pf\ Four distinct neighbours of one vertex would embed
$\widetilde{\typ D}_4$ by Definition~\ref{def:affine}.~\qed\end{proof}
\step{ruone}{\assume{There are distinct degree-3 vertices $u,v$.}
\prove{\textsc{False}.}}
\begin{proof}
\pfsketch\ \new\ We may use any two branch vertices: additional edges
at internal path vertices are omitted from the obstruction subgraph.
\pick\ a simple path between $u,v$, by connectedness in \asm{0}.
\step{ruextra}{Each endpoint has two distinct neighbours off the path.}
\begin{proof}\pf\ An additional neighbour on the path would contradict
path uniqueness from \stepref{rutree}; the degree is 3. ~\qed\end{proof}
\step{rudistinct}{The four selected neighbours are mutually distinct.}
\begin{proof}\pf\ A common neighbour of the endpoints would give a
second simple path between them, contradicting \stepref{rutree}.~\qed\end{proof}
\step{ruembed}{The selected path and four incident edges embed
$\widetilde{\typ D}_{k+4}$, where $k\geq1$ is the path length.}
\begin{proof}\pf\ Use the injective vertex list from
\stepref{ruextra}, \stepref{rudistinct} in Definition~\ref{def:affine}.~\qed\end{proof}
\qedstep
\begin{proof}\pf\ Step \stepref{ruembed} contradicts \asm{0}.~\qed\end{proof}
\end{proof}
\step{rupath}{\case{There is no degree-3 vertex.}}
\begin{proof}\pf\ Apply Lemma~\ref{lem:path-labelling} using
\stepref{rudegree}, \stepref{rutree}.~\qed\end{proof}
\step{rubranch}{\case{There is a degree-3 vertex.}}
\begin{proof}
\step{ruarms}{\pick\ an arm bijection with $1\leq a\leq b\leq c$.}
\begin{proof}\pf\ Apply Lemma~\ref{lem:arm-labelling} to
\stepref{rutree}, \stepref{rudegree}, \stepref{ruone}.~\qed\end{proof}
\step{rua}{$a=1$.}
\begin{proof}\pf\ Otherwise the restriction to the first two positions
of each arm in \stepref{ruarms} embeds $T(2,2,2)$, forbidden by \asm{0}.~\qed\end{proof}
\step{rub}{$b\leq2$.}
\begin{proof}\pf\ If $b\geq3$, restricting the arms to lengths $1,3,3$
in \stepref{ruarms} embeds a forbidden $\widetilde{\typ E}_7$.~\qed\end{proof}
\step{rubone}{\assume{$b=1$.}
\prove{The labelling is $T(1,1,c)$.}}
\begin{proof}\pf\ Substitute in \stepref{ruarms}, using \stepref{rua}.~\qed\end{proof}
\step{rubtwo}{\assume{$b=2$.}
\prove{$c\in\{2,3,4\}$.}}
\begin{proof}\pf\ A value $c\geq5$ embeds $T(1,2,5)$ by restriction;
$2\leq c\leq4$ gives the asserted integer alternatives. ~\qed\end{proof}
\qedstep
\begin{proof}\pf\ The positive integer $b\leq2$ in \stepref{rub} is
1 or 2, so \stepref{rubone}, \stepref{rubtwo} exhaust this case. ~\qed\end{proof}
\end{proof}
\qedstep
\begin{proof}\pf\ The cases \stepref{rupath}, \stepref{rubranch} are
exhaustive. ~\qed\end{proof}
\end{proof}

\begin{namedresult}{Theorem}{(weighted dichotomy)}\label{thm:dichotomy}
\lean{WeightedGraph.affine_or_dynkin}
\uses{def:affine,def:rem-weighted,lem:unweighted-trees,lem:path-labelling,lem:affine,lem:monotonicity,lem:rem-unit-containment,lem:rem-rayleigh-reindex,lem:rem-labelled-connected}
\assume{$X$ is a finite connected weighted graph whose edge weights belong
to $\{1,\sqrt2,\sqrt3\}$.}
\prove{Exactly one alternative holds: $X$ contains an affine obstruction;
or it admits a weight-preserving path or arm bijection to one of
\[
 \begin{gathered}
 \typ A_\ell\ (\ell\geq1),\quad
 \typ{BC}_\ell\ (\ell\geq2),\quad
 T(1,1,\ell-3)=\typ D_\ell\ (\ell\geq4),\\
 T(1,2,\ell-4)=\typ E_\ell\ (\ell=6,7,8),\quad
 \typ F_4,\quad\typ G_2.
 \end{gathered}
\]
Here $\typ A$ is a unit path, $\typ{BC}$ has one terminal $\sqrt2$ edge,
$\typ D,\typ E$ have unit edges, $\typ F_4$ has path weights
$1,\sqrt2,1$, and $\typ G_2$ is one $\sqrt3$ edge.
Every graph in the second alternative has largest eigenvalue less than 2.}
\end{namedresult}
\begin{proof}
\pfsketch\ Retain the labelling of the unit support and specify the
injection for every excluded weighted pattern.
\step{rdrecognize}{\assume{$X$ contains no affine obstruction.}
\prove{$X$ admits one of the listed weight-preserving Dynkin labellings.}}
\begin{proof}
\step{rdsupport}{The support of $X$ has one of the path or arm labellings in
Lemma~\ref{lem:unweighted-trees}.}
\begin{proof}\pf\ By Lemma~\ref{lem:rem-unit-containment}, an unweighted
obstruction embedded in the support would be contained in $X$, since
all its edge weights are at least 1. Thus \asm{1} supplies all hypotheses
of Lemma~\ref{lem:unweighted-trees}.~\qed\end{proof}
\step{rdtriple}{\assume{There is a $\sqrt3$ edge with endpoints $u,v$.}
\prove{Those endpoints give a weight-preserving $\typ G_2$ labelling of $X$.}}
\begin{proof}
\step{rdthird}{\assume{There is a vertex outside $\{u,v\}$.}
\prove{\textsc{False}.}}
\begin{proof}
\step{rdexit}{There is a support edge from $\{u,v\}$ to its complement.}
\begin{proof}\pf\ Take the first exit from $\{u,v\}$ on a connecting
walk to the vertex in \asm{3}, using connectedness in \asm{0}.~\qed\end{proof}
\step{rdtripleembed}{$X^{\rm triple}\preccurlyeq X$.}
\begin{proof}\pf\ Keep the edge in \asm{2} and lower the edge in
\stepref{rdexit} to 1; their three distinct vertices give the injection
in Definition~\ref{def:affine}.~\qed\end{proof}
\qedstep
\begin{proof}\pf\ Step \stepref{rdtripleembed} contradicts \asm{1}.~\qed\end{proof}
\end{proof}
\qedstep
\begin{proof}\pf\ Step \stepref{rdthird} leaves exactly the two distinct
endpoints; their edge has weight $\sqrt3$ by \asm{2}.~\qed\end{proof}
\end{proof}
\step{rdbranch}{\assume{There is no $\sqrt3$ edge, and the support has
an arm labelling from \stepref{rdsupport}.}
\prove{This labelling identifies $X$ with a listed $\typ D$ or $\typ E$ graph.}}
\begin{proof}
\step{rdfork}{\assume{There is a $\sqrt2$ edge.}
\prove{\textsc{False}.}}
\begin{proof}
\pick\ the arm containing that edge; permute the arm labels to call it
arm 1, and let $k\geq1$ be the outer endpoint position of that edge.
\step{rdforkmap}{The map from $v_0,\ldots,v_k$ to
$e(*),e(1,1),\ldots,e(1,k)$, and from the two extra leaves to
$e(2,1),e(3,1)$, is injective.}
\begin{proof}\pf\ These positions are distinct in the arm bijection
of \asm{2}, including when $k=1$.~\qed\end{proof}
\step{rdforkweights}{This map witnesses $X^{\rm fork}_k\preccurlyeq X$.}
\begin{proof}\pf\ Retain its last $\sqrt2$ edge and lower every other
selected edge to 1 in \stepref{rdforkmap}; use the weight bounds in \asm{0}.~\qed\end{proof}
\qedstep
\begin{proof}\pf\ Step \stepref{rdforkweights} contradicts \asm{1}.~\qed\end{proof}
\end{proof}
\step{rdbranchunit}{All edges have weight 1.}
\begin{proof}\pf\ The weight alternatives in \asm{0}, absence of
$\sqrt3$ in \asm{2}, and \stepref{rdfork} exclude both other weights.~\qed\end{proof}
\step{rdbranchD}{\case{The labelling is $T(1,1,c)$, $c\geq1$.}
\prove{It is a weight-preserving $\typ D_{c+3}$ labelling.}}
\begin{proof}\pf\ Set $\ell=c+3\geq4$ in the theorem's definition of
$\typ D_\ell$, using \stepref{rdbranchunit}.~\qed\end{proof}
\step{rdbranchE}{\case{The labelling is $T(1,2,c)$, $c\in\{2,3,4\}$.}
\prove{It is a weight-preserving $\typ E_{c+4}$ labelling.}}
\begin{proof}\pf\ Set $\ell=c+4\in\{6,7,8\}$ in the theorem's definition
of $\typ E_\ell$, using \stepref{rdbranchunit}.~\qed\end{proof}
\qedstep
\begin{proof}\pf\ The two arm alternatives in \stepref{rdsupport}
are exactly \stepref{rdbranchD}, \stepref{rdbranchE}.~\qed\end{proof}
\end{proof}
\step{rdpath}{\assume{There is no $\sqrt3$ edge, and the support has
a path labelling $e:[0,\ell-1]\to V(X)$.}
\prove{$X$ has a weight-preserving $\typ A$, $\typ{BC}$ or $\typ F_4$ labelling.}}
\begin{proof}
\step{rdtwo}{There is at most one $\sqrt2$ edge.}
\begin{proof}
\step{rdtwocontra}{\assume{Two such edges have positions $r<s$,
where position $i$ denotes the edge $e(i)e(i+1)$.}
\prove{\textsc{False}.}}
\begin{proof}
\step{rdendmap}{The map $i\mapsto e(r+i)$,
$0\leq i\leq s+1-r$, is injective.}
\begin{proof}\pf\ Restrict the path bijection of \asm{2} to this interval.~\qed\end{proof}
\step{rdendembed}{This map witnesses $X^{\rm end}_{s-r+2}\preccurlyeq X$.}
\begin{proof}\pf\ Retain the first and last $\sqrt2$ edges of \asm{4}
and lower intervening weights to 1, using \stepref{rdendmap};
$s-r+2\geq3$ includes consecutive multiple edges.~\qed\end{proof}
\qedstep
\begin{proof}\pf\ Step \stepref{rdendembed} contradicts \asm{1}.~\qed\end{proof}
\end{proof}
\qedstep
\begin{proof}\pf\ Two distinct multiple edges would have ordered
positions forbidden by \stepref{rdtwocontra}.~\qed\end{proof}
\end{proof}
\step{rdunit}{\case{There is no $\sqrt2$ edge.}
\prove{$e$ is a weight-preserving $\typ A_\ell$ labelling.}}
\begin{proof}\pf\ Every edge has weight 1 by \asm{0}, \asm{2},
and this case.~\qed\end{proof}
\step{rddouble}{\case{There is a $\sqrt2$ edge.}
\prove{$X$ has a weight-preserving $\typ{BC}$ or $\typ F_4$ labelling.}}
\begin{proof}
\step{rdorient}{\pick\ a path direction in which the two sides of the
multiple edge have $1\leq p\leq q$ vertices; its position is $p-1,p$.}
\begin{proof}\pf\ The reversal assertion of Lemma~\ref{lem:path-labelling}
interchanges the two sides, so one of the two directions orders their sizes.~\qed\end{proof}
\step{rdterminal}{\case{$p=1$.}
\prove{This is a $\typ{BC}_{p+q}$ labelling.}}
\begin{proof}\pf\ The unique multiple edge is terminal by
\stepref{rdorient}; all others are unit by \stepref{rdtwo}, \asm{2}.~\qed\end{proof}
\step{rdF}{\case{$p=q=2$.}
\prove{This is a $\typ F_4$ labelling.}}
\begin{proof}\pf\ Its three weights are $1,\sqrt2,1$ by
\stepref{rdtwo}, \stepref{rdorient}, \asm{2}.~\qed\end{proof}
\step{rdmiddle}{\case{$p\geq2,q\geq3$.}
\prove{\textsc{False}.}}
\begin{proof}
\step{rdmidvertices}{The positions $p+2,p+1,p,p-1,p-2$ are distinct
and lie in $[0,p+q-1]$.}
\begin{proof}\pf\ The least position satisfies $p-2\geq0$, the greatest
satisfies $p+2\leq p+q-1$, and consecutive listed positions differ by 1.~\qed\end{proof}
\step{rdmidweights}{Their path weights in this order are $1,1,\sqrt2,1$.}
\begin{proof}\pf\ Use the unique multiple edge at $p-1,p$ from
\stepref{rdtwo}, \stepref{rdorient}, with no $\sqrt3$ edge by \asm{2}.~\qed\end{proof}
\qedstep
\begin{proof}\pf\ Steps \stepref{rdmidvertices}, \stepref{rdmidweights}
embed $X^{\rm mid}$, contrary to \asm{1}.~\qed\end{proof}
\end{proof}
\qedstep
\begin{proof}\pf\ The integer alternatives \stepref{rdterminal},
\stepref{rdF}, \stepref{rdmiddle} exhaust $1\leq p\leq q$.~\qed\end{proof}
\end{proof}
\qedstep
\begin{proof}\pf\ A $\sqrt2$ edge either exists or does not;
apply \stepref{rddouble} or \stepref{rdunit}.~\qed\end{proof}
\end{proof}
\qedstep
\begin{proof}\pf\ If a $\sqrt3$ edge exists, use \stepref{rdtriple}.
Otherwise \stepref{rdsupport} supplies an arm or a path labelling,
the hypotheses of \stepref{rdbranch} or \stepref{rdpath}.~\qed\end{proof}
\end{proof}
\step{rdbelow}{Every listed Dynkin graph has largest eigenvalue less than 2.}
\begin{proof}
\step{rdAone}{$\typ A_1$ is a proper subgraph of $\Cyc_3$.}
\begin{proof}\pf\ Map its sole vertex to any cycle vertex. ~\qed\end{proof}
\step{rdAmany}{For $\ell\geq2$, $\typ A_\ell$ is a proper subgraph of
$\Cyc_{\ell+1}$.}
\begin{proof}\pf\ Delete one vertex of that cycle. ~\qed\end{proof}
\step{rdDfour}{$\typ D_4$ is a proper subgraph of $\widetilde{\typ D}_4$.}
\begin{proof}\pf\ Delete one leaf of $K_{1,4}$.~\qed\end{proof}
\step{rdDmany}{For $\ell\geq5$, $\typ D_\ell$ is a proper subgraph of
$\widetilde{\typ D}_\ell$.}
\begin{proof}\pf\ Delete one leaf at one fork; the other fork has arms
$1,1,(\ell-4)+1$ by Definition~\ref{def:affine}.~\qed\end{proof}
\step{rdE6}{$\typ E_6$ embeds properly in $T(2,2,2)$.}
\begin{proof}\pf\ Restrict one arm to length 1. ~\qed\end{proof}
\step{rdE7}{$\typ E_7$ embeds properly in $T(1,3,3)$.}
\begin{proof}\pf\ Restrict one length-3 arm to length 2. ~\qed\end{proof}
\step{rdE8}{$\typ E_8$ embeds properly in $T(1,2,5)$.}
\begin{proof}\pf\ Restrict the length-5 arm to length 4. ~\qed\end{proof}
\step{rdBC}{$\typ{BC}_\ell$ embeds properly in $X^{\rm end}_{\ell+1}$.}
\begin{proof}\pf\ Delete its first endpoint, retaining the last
$\sqrt2$ edge; this includes $\ell=2$.~\qed\end{proof}
\step{rdFsub}{$\typ F_4$ embeds properly in $X^{\rm mid}$.}
\begin{proof}\pf\ Delete the first endpoint of the path with weights
$1,1,\sqrt2,1$.~\qed\end{proof}
\step{rdGsub}{$\typ G_2$ embeds properly in $X^{\rm triple}$.}
\begin{proof}\pf\ Delete the endpoint incident only with the unit edge. ~\qed\end{proof}
\qedstep
\begin{proof}\pf\ Apply Lemmas~\ref{lem:affine}, \ref{lem:monotonicity}
to each containment \stepref{rdAone}--\stepref{rdGsub};
the larger obstruction is connected by Lemma~\ref{lem:rem-labelled-connected}.
Lemma~\ref{lem:rem-rayleigh-reindex} transfers the bound along every
weight-preserving labelling.~\qed\end{proof}
\end{proof}
\step{rdexclusive}{No graph admitting a listed Dynkin labelling contains an obstruction.}
\begin{proof}\pf\ Such a containment would give $2\leq\lambda$ by
Lemmas~\ref{lem:affine}, \ref{lem:monotonicity}, contradicting
\stepref{rdbelow} after transfer by Lemma~\ref{lem:rem-rayleigh-reindex};
its labelled support is connected by Lemma~\ref{lem:rem-labelled-connected}.~\qed\end{proof}
\qedstep
\begin{proof}\pf\ Combine \stepref{rdrecognize}, \stepref{rdbelow},
\stepref{rdexclusive} with excluded middle for obstruction containment. ~\qed\end{proof}
\end{proof}

\begin{namedresult}{Definition}{(the indexed matrices)}\label{def:rem-matrices}
\lean{CartanMatrix.modelEntries}\uses{facts:rem-library,def:rem-bases}
\assume{The rank is in the corresponding range of \eqref{eq:list};
indices are $0,\ldots,\ell-1$.}
All diagonals are 2 and unspecified entries are zero. Unit edges have
both entries $-1$. For $\typ A$ the edges are consecutive pairs.
For $\typ B$ they are consecutive pairs except that the last forward
entry is $-2$; $\typ C$ is its transpose.
For $\typ D$ the edges are $\{i,i+1\}$ for $0\leq i\leq\ell-3$,
and $\{\ell-3,\ell-1\}$. For $\typ E$ they are
$\{0,2\},\{1,3\}$ and $\{i,i+1\}$ for $2\leq i\leq\ell-2$.
The two remaining matrices are
\[
 \typ F_4=\begin{pmatrix}2&-1&0&0\\-1&2&-2&0\\0&-1&2&-1\\0&0&-1&2\end{pmatrix},
 \qquad \typ G_2=\begin{pmatrix}2&-3\\-1&2\end{pmatrix}.
\]
The following positive integer diagonal vectors will be used as symmetrizers:
\begin{equation}\label{eq:symmetrizers}
\begin{array}{c|c}
\typ A,\typ D,\typ E&(1,\ldots,1)\\
\typ B&(1,\ldots,1,2)\\
\typ C&(2,\ldots,2,1)\\
\typ F_4&(1,1,2,2)\\
\typ G_2&(1,3).
\end{array}
\end{equation}
These are the defining expressions in the Cartan file in
\ref{facts:rem-library}, not new matrix conventions.
\end{namedresult}

\begin{namedresult}{Lemma}{(supports of the named matrices)}\label{lem:rem-model-supports}
\lean{CartanMatrix.model_support_labelling}
\lean{CartanMatrix.model_support_connected}
\uses{def:rem-matrices,def:affine,def:rem-weighted,lem:rem-labelled-connected,facts:rem-library}
\assume{$N$ is $\typ A_\ell$ with $\ell\geq1$, $\typ B_\ell$ or
$\typ C_\ell$ with $\ell\geq2$, $\typ D_\ell$ with $\ell\geq4$,
$\typ E_\ell$ with $\ell\geq5$, $\typ F_4$, or $\typ G_2$, using the
actual Mathlib definitions on $\mathrm{Fin}\,\ell$.}
\pflet{$W_N(i,i)=0$ and
$W_N(i,j)=\sqrt{N(i,j)N(j,i)}$ for $i\ne j$.}
\prove{$W_N$ is a weighted graph with the same support as the off-diagonal
nonzero entries of $N$. The following maps are weight-preserving
bijections onto it, with exactly the displayed edges, and its support is connected:
\[
\begin{array}{c|c|c}
N&\text{domain and map}&\text{successive weights}\\ \hline
\typ A_\ell&[0,\ell-1],\ i\mapsto i&1,\ldots,1\\
\typ B_\ell,\typ C_\ell&[0,\ell-1],\ i\mapsto i&1,\ldots,1,\sqrt2\\
\typ F_4&[0,3],\ i\mapsto i&1,\sqrt2,1\\
\typ G_2&[0,1],\ i\mapsto i&\sqrt3\\
\typ D_\ell&T(1,1,\ell-3),\ h_D&\text{all }1\\
\typ E_\ell&T(1,2,\ell-4),\ h_E&\text{all }1
\end{array}
\]
Here the arm bijections are
\begin{equation}\label{eq:rem-support-maps}
\begin{aligned}
h_D(*)&=\ell-3,&h_D(1,1)&=\ell-2,&h_D(2,1)&=\ell-1,\\
h_D(3,t)&=\ell-3-t&&&(1\leq t\leq\ell-3),\\
h_E(*)&=3,&h_E(1,1)&=1,&h_E(2,1)&=2,\\
h_E(2,2)&=0,&h_E(3,t)&=3+t&&(1\leq t\leq\ell-4).
\end{aligned}
\end{equation}
All diagonals of $N$ are 2, all unit-edge entries are $(-1,-1)$, and
all nonedges have zero entries in both directions. In the displayed
path orders the exceptional forward/reverse pairs are $(-2,-1)$ on
the last $\typ B$ edge, $(-1,-2)$ on the last $\typ C$ edge,
$(-2,-1)$ on the middle $\typ F$ edge, and $(-3,-1)$ on the $\typ G$ edge.}
\end{namedresult}
\begin{proof}
\pfsketch\ Check the matrix definitions on their own index types.
The arm maps are the inverses of the index orders used to recognize
arbitrary matrices; this lemma establishes their properties independently.
\step{rmsA}{The identity is the asserted unit-path labelling for $\typ A_\ell$.}
\begin{proof}\pf\ The off-diagonal nonzero clause of
\path{CartanMatrix.A} is $i+1=j$ or $j+1=i$, with entry $-1$.~\qed\end{proof}
\step{rmsB}{The identity is the asserted path labelling for $\typ B_\ell$.}
\begin{proof}\pf\ In \path{CartanMatrix.B}, the forward consecutive
entry is $-2$ exactly at $j=\ell-1$; all reverse consecutive entries
are $-1$, and every other off-diagonal entry vanishes.~\qed\end{proof}
\step{rmsC}{The identity is the asserted path labelling for $\typ C_\ell$.}
\begin{proof}\pf\ In \path{CartanMatrix.C}, the reverse consecutive
entry is $-2$ exactly at $i=\ell-1$; the other consecutive entries
are $-1$, and the other off-diagonal entries vanish.~\qed\end{proof}
\step{rmsF}{The identity is the asserted path labelling for $\typ F_4$.}
\begin{proof}\pf\ Read the consecutive pairs $(-1,-1),(-2,-1),(-1,-1)$
and all zero nonedges in \texttt{CartanMatrix.F}\textsubscript{4}.~\qed\end{proof}
\step{rmsG}{The identity is the asserted path labelling for $\typ G_2$.}
\begin{proof}\pf\ Its only off-diagonal pair is $(-3,-1)$ in
\texttt{CartanMatrix.G}\textsubscript{2}.~\qed\end{proof}
\step{rmsD}{$h_D$ is the asserted unit-arm bijection.}
\begin{proof}
\step{rmsDbij}{$h_D$ is bijective.}
\begin{proof}\pf\ Its centre and short arms cover $\ell-3,\ell-2,\ell-1$;
$t\mapsto\ell-3-t$ bijects arm 3 onto $[0,\ell-4]$.
These ranges are disjoint and exhaust $\mathrm{Fin}\,\ell$.~\qed\end{proof}
\step{rmsDentries}{The nonzero off-diagonal unordered pairs of
$\typ D_\ell$ are
$\{i,i+1\}$ for $0\leq i\leq\ell-4$ and
$\{\ell-3,\ell-2\},\{\ell-3,\ell-1\}$, all unit pairs.}
\begin{proof}\pf\ For $i<j$, the clauses of \path{CartanMatrix.D}
reduce to $j=i+1$, $j+2<\ell$, or
$i=\ell-3$, $j\in\{\ell-2,\ell-1\}$.
The reversed clauses give the reversed entries.~\qed\end{proof}
\step{rmsDcentre}{The three centre edges map to
$\{\ell-3,\ell-2\},\{\ell-3,\ell-1\},\{\ell-3,\ell-4\}$.}
\begin{proof}\pf\ Substitute the centre and first arm positions
in \eqref{eq:rem-support-maps}.~\qed\end{proof}
\step{rmsDarm}{The remaining arm edges map exactly to
$\{i,i+1\}$ for $0\leq i$ and $i+5\leq\ell$.}
\begin{proof}\pf\ On arm 3, consecutive $t,t+1$ give consecutive
$\ell-3-t,\ell-4-t$; the interval is empty at $\ell=4$.
The two length-1 arms have no internal edges.~\qed\end{proof}
\qedstep
\begin{proof}\pf\ The edge images in \stepref{rmsDcentre},
\stepref{rmsDarm} exhaust \stepref{rmsDentries}; bijectivity
\stepref{rmsDbij} therefore also excludes every other edge.~\qed\end{proof}
\end{proof}
\step{rmsE}{$h_E$ is the asserted unit-arm bijection for every $\ell\geq5$.}
\begin{proof}
\step{rmsEbij}{$h_E$ is bijective.}
\begin{proof}\pf\ The centre and first two arms cover $3,1,2,0$;
$t\mapsto3+t$ bijects arm 3 onto $[4,\ell-1]$,
disjoint from those four values.~\qed\end{proof}
\step{rmsEentries}{The nonzero off-diagonal unordered pairs are
$\{0,2\},\{1,3\}$ and $\{i,i+1\}$ for $2\leq i\leq\ell-2$,
all unit pairs.}
\begin{proof}\pf\ These are exactly the nonzero clauses in
\path{CartanMatrix.E}, valid in the whole assumed range $\ell\geq5$.~\qed\end{proof}
\step{rmsEcentre}{The centre edges map to $\{3,1\},\{3,2\},\{3,4\}$.}
\begin{proof}\pf\ Substitute the first position on each arm
in \eqref{eq:rem-support-maps}.~\qed\end{proof}
\step{rmsEarm}{The remaining edges map to $\{2,0\}$ and
$\{i,i+1\}$ for $4\leq i\leq\ell-2$.}
\begin{proof}\pf\ The two positions of arm 2 map to $2,0$;
consecutive positions on arm 3 map to consecutive integers starting at 4.
Arm 1 has no internal edge.~\qed\end{proof}
\qedstep
\begin{proof}\pf\ The edge images in \stepref{rmsEcentre},
\stepref{rmsEarm} exhaust \stepref{rmsEentries}; use
\stepref{rmsEbij} for the reverse edge implication.~\qed\end{proof}
\end{proof}
\step{rmsweights}{The weight and support assertions follow from these entries.}
\begin{proof}\pf\ The products of the displayed nonzero pairs are
$1,2,3$; every other off-diagonal product is zero.
Their nonnegative square roots define $W_N$.~\qed\end{proof}
\step{rmsconnected}{Every displayed support is connected.}
\begin{proof}\pf\ Apply Lemma~\ref{lem:rem-labelled-connected}
to the path or arm bijections in \stepref{rmsA}--\stepref{rmsE}.~\qed\end{proof}
\qedstep
\begin{proof}\pf\ Combine \stepref{rmsA}--\stepref{rmsconnected};
the diagonal clause of each named definition is 2.~\qed\end{proof}
\end{proof}

\begin{namedresult}{Definition}{(integer matrix invariants)}\label{def:rem-invariants}
\lean{CartanMatrix.supportInvariants}
\uses{}
\assume{$N:I\times I\to\mathbb Z$ is a matrix on a finite set with
$N(i,j)=0\iff N(j,i)=0$.}
Let $d_N(i)=|\{j:j\ne i,\ N(i,j)\ne0\}|$ be its support degree.
Define the following nonnegative integers, with cardinalities regarded in $\mathbb Z$:
\begin{equation}\label{eq:rem-invariants}
\begin{aligned}
n(N)&=|I|,\\
t(N)&=|\{(i,j):N(i,j)N(j,i)=3\}|,\\
u(N)&=|\{(i,j):N(i,j)N(j,i)=2\}|,\\
b(N)&=|\{i:d_N(i)=3\}|,\\
a_{\rm in}(N)&=|\{(i,j):N(i,j)=-2,\ d_N(j)=1\}|,\\
a_{\rm out}(N)&=|\{(i,j):N(i,j)=-2,\ d_N(i)=1\}|,\\
s(N)&=|\{i:d_N(i)=1,\ \exists j,\ N(i,j)\ne0,\ d_N(j)=3\}|.
\end{aligned}
\end{equation}
Write $\mathcal I(N)=(n,t,u,b,a_{\rm in},a_{\rm out},s)(N)$.
Pairs in $t,u,a_{\rm in},a_{\rm out}$ are ordered; $s$ counts leaves,
not pairs.
\end{namedresult}

\begin{namedresult}{Lemma}{(invariance and values)}\label{lem:rem-invariants}
\lean{CartanMatrix.supportInvariants_reindex}
\lean{CartanMatrix.model_supportInvariants}
\uses{def:rem-invariants,lem:rem-model-supports,def:rem-bases}
\assume{$N$ is an integer matrix on a finite set $I$ with symmetric
zero relation.}
\prove{For every bijection $e:J\to I$, the matrix
$N_e(i,j)=N(e(i),e(j))$ satisfies $\mathcal I(N_e)=\mathcal I(N)$.
For every permitted named matrix, the values are:
\begin{equation}\label{eq:rem-invariant-table}
\begin{array}{c|rrrrrrr}
N&n&t&u&b&a_{\rm in}&a_{\rm out}&s\\ \hline
\typ A_\ell\ (\ell\geq1)&\ell&0&0&0&0&0&0\\
\typ B_2&2&0&2&0&1&1&0\\
\typ B_\ell\ (\ell\geq3)&\ell&0&2&0&1&0&0\\
\typ C_\ell\ (\ell\geq3)&\ell&0&2&0&0&1&0\\
\typ D_4&4&0&0&1&0&0&3\\
\typ D_\ell\ (\ell\geq5)&\ell&0&0&1&0&0&2\\
\typ E_6&6&0&0&1&0&0&1\\
\typ E_7&7&0&0&1&0&0&1\\
\typ E_8&8&0&0&1&0&0&1\\
\typ F_4&4&0&2&0&0&0&0\\
\typ G_2&2&2&0&0&0&0&0
\end{array}
\end{equation}}
\end{namedresult}
\begin{proof}
\pfsketch\ Transport finite counting sets along the bijection, then
count on the independently labelled supports of the named matrices.
\step{rivreindex}{All seven integers are preserved by reindexing.}
\begin{proof}
\step{rivdegree}{$d_{N_e}(i)=d_N(e(i))$ for every $i$.}
\begin{proof}\pf\ The map $j\mapsto e(j)$ bijects the neighbour sets
in Definition~\ref{def:rem-invariants}.~\qed\end{proof}
\step{rivvertices}{The integers $n,b,s$ are preserved.}
\begin{proof}\pf\ Transport their defining vertex sets by $e$,
using \stepref{rivdegree} and the entry equality defining $N_e$.~\qed\end{proof}
\step{rivpairs}{The integers $t,u,a_{\rm in},a_{\rm out}$ are preserved.}
\begin{proof}\pf\ The map $(i,j)\mapsto(e(i),e(j))$ bijects each
defining pair set, by \stepref{rivdegree} and the entry equality.~\qed\end{proof}
\qedstep
\begin{proof}\pf\ Combine \stepref{rivvertices}, \stepref{rivpairs}.~\qed\end{proof}
\end{proof}
\step{rivA}{The $\typ A$ row has only its size nonzero.}
\begin{proof}\pf\ Lemma~\ref{lem:rem-model-supports} gives a unit
path, with no degree-3 vertex or $-2$ entry; use
Definition~\ref{def:rem-invariants}.~\qed\end{proof}
\step{rivBtwo}{The $\typ B_2$ row has $u=2$ and
$(a_{\rm in},a_{\rm out})=(1,1)$.}
\begin{proof}\pf\ The sole double edge has two ordered products 2;
both endpoints of its sole $-2$ entry are leaves, by
Lemma~\ref{lem:rem-model-supports}.~\qed\end{proof}
\step{rivB}{For $\ell\geq3$, the $\typ B_\ell$ row has $u=2$ and
$(a_{\rm in},a_{\rm out})=(1,0)$.}
\begin{proof}\pf\ The sole $-2$ entry points from position $\ell-2$
of degree 2 to the leaf $\ell-1$ in Lemma~\ref{lem:rem-model-supports};
the path has no branch vertex.~\qed\end{proof}
\step{rivC}{For $\ell\geq3$, the $\typ C_\ell$ row has $u=2$ and
$(a_{\rm in},a_{\rm out})=(0,1)$.}
\begin{proof}\pf\ The sole $-2$ entry points from the leaf $\ell-1$
to position $\ell-2$ of degree 2 in Lemma~\ref{lem:rem-model-supports};
the path has no branch vertex.~\qed\end{proof}
\step{rivDfour}{For $\typ D_4$, $b=1$ and $s=3$.}
\begin{proof}\pf\ The unit tree $T(1,1,1)$ in
Lemma~\ref{lem:rem-model-supports} has one degree-3 centre adjacent
to all three leaves.~\qed\end{proof}
\step{rivDmany}{For $\typ D_\ell$, $\ell\geq5$, $b=1$ and $s=2$.}
\begin{proof}\pf\ In $T(1,1,\ell-3)$ the third arm has length at least 2,
so only the first two leaves meet the centre, by
Lemma~\ref{lem:rem-model-supports}.~\qed\end{proof}
\step{rivE}{For each of $\typ E_6,\typ E_7,\typ E_8$, $b=1$ and $s=1$.}
\begin{proof}\pf\ The arm lengths $1,2,\ell-4$ in
Lemma~\ref{lem:rem-model-supports} have exactly one length equal to 1;
all edges are unit.~\qed\end{proof}
\step{rivF}{For $\typ F_4$, $u=2$ and
$(a_{\rm in},a_{\rm out})=(0,0)$.}
\begin{proof}\pf\ The sole $-2$ entry joins the two degree-2 vertices
of the four-vertex path in Lemma~\ref{lem:rem-model-supports}.~\qed\end{proof}
\step{rivG}{For $\typ G_2$, $n=t=2$ and
$u=b=a_{\rm in}=a_{\rm out}=s=0$.}
\begin{proof}\pf\ Lemma~\ref{lem:rem-model-supports} gives two vertices,
one product-3 edge, no $-2$ entry and no branch vertex.~\qed\end{proof}
\qedstep
\begin{proof}\pf\ In the counts \stepref{rivA}--\stepref{rivG}, each
matrix has its stated size, and the listed supports and entries
exclude all uncounted patterns; together with \stepref{rivreindex}
these give \eqref{eq:rem-invariant-table}.~\qed\end{proof}
\end{proof}

\begin{namedresult}{Theorem}{A}\label{thm:A}
\lean{Matrix.IsFiniteCartan.classification}
\lean{CartanMatrix.IsFiniteTypeModel.type_eq_of_reindex_eq}
\lean{CartanMatrix.IsFiniteTypeModel.isFiniteCartan}
\lean{CartanMatrix.IsFiniteTypeModel.isIndecomposable}
\uses{def:rem-matrices,lem:rem-model-supports,lem:rem-invariants,def:rem-invariants,lem:symmetrization,thm:dichotomy,lem:affine,lem:monotonicity,lem:connected,lem:rayleigh,lem:rem-rayleigh-reindex,lem:path-labelling,facts:rem-library,def:rem-bases}
\assume{$I$ is a finite nonempty set and $M:I\times I\to\mathbb Z$
satisfies \path{Matrix.IsFiniteCartan} and \path{Matrix.IsIndecomposable}.}
\prove{There are $\ell>0$ and a bijection $e:[0,\ell-1]\to I$ such that
$M(e(i),e(j))$ is a matrix in \eqref{eq:list}. Its type and rank are
unique; the bijection need not be unique. Conversely every listed matrix
is finite Cartan and indecomposable.}
\end{namedresult}
\noindent\repair\ The empty matrix satisfies both library predicates.
It is excluded here, or may be added as a separate rank-zero alternative.
\begin{proof}
\pfsketch\ Classify the weighted support, orient its one possible multiple
edge, and distinguish the resulting matrices by intrinsic invariants.
\step{rAsupport}{The symmetrized support of $M$ has a Dynkin labelling.}
\begin{proof}
\step{rAbound}{The symmetrization has $\lambda<2$.}
\begin{proof}\pf\ Apply Lemma~\ref{lem:symmetrization}.~\qed\end{proof}
\step{rAnoobstruction}{It contains no affine obstruction.}
\begin{proof}\pf\ Lemmas~\ref{lem:affine}, \ref{lem:monotonicity}
would give $2\leq\lambda$, contradicting \stepref{rAbound}.~\qed\end{proof}
\qedstep
\begin{proof}\pf\ Use the no-obstruction alternative of
Theorem~\ref{thm:dichotomy} with \stepref{rAnoobstruction}.~\qed\end{proof}
\end{proof}
\step{rAentries}{The possible ordered entry pairs are
\[
\begin{array}{c|l}
\text{unit edge}&(-1,-1)\\
\text{double edge}&(-1,-2)\text{ or }(-2,-1)\\
\text{triple edge}&(-1,-3)\text{ or }(-3,-1).
\end{array}
\]}
\begin{proof}
\step{rAunit}{Product 1 forces $(-1,-1)$.}
\begin{proof}\pf\ The two negative integers are negatives of positive
divisors of 1. ~\qed\end{proof}
\step{rAdouble}{Product 2 forces $(-1,-2)$ or $(-2,-1)$.}
\begin{proof}\pf\ The positive divisors of 2 are 1 and 2. ~\qed\end{proof}
\step{rAtriple}{Product 3 forces $(-1,-3)$ or $(-3,-1)$.}
\begin{proof}\pf\ The positive divisors of 3 are 1 and 3. ~\qed\end{proof}
\qedstep
\begin{proof}\pf\ Apply these three arithmetic cases to
Lemma~\ref{lem:symmetrization}.~\qed\end{proof}
\end{proof}
\step{rAlabelling}{There are $\ell>0$, a listed type of rank $\ell$,
and a bijection $e:\mathrm{Fin}\,\ell\to I$ such that
$M(e(i),e(j))$ equals that named matrix at every $i,j$.}
\begin{proof}
\step{rApath}{For a unit path, its path bijection gives $\typ A_\ell$.}
\begin{proof}\pf\ Substitute its edge equivalence and
\stepref{rAentries} in Definition~\ref{def:rem-matrices}.~\qed\end{proof}
\step{rAD}{For $T(1,1,\ell-3)$, use its arm bijection $f$ and set
\[
 \begin{split}
 e(\ell-3)&=f(*),\quad e(\ell-2)=f(1,1),\quad e(\ell-1)=f(2,1),\\
 e(i)&=f(3,\ell-3-i)\quad(0\leq i\leq\ell-4).
 \end{split}
\]
Then $M(e(i),e(j))=\typ D_\ell(i,j)$.}
\begin{proof}
\step{rADbij}{The formula for $e$ is a bijection.}
\begin{proof}\pf\ It is $f\circ h_D^{-1}$, a composition of bijections by
Lemma~\ref{lem:rem-model-supports}. ~\qed\end{proof}
\step{rADedges}{The edges become $\{i,i+1\}$ for $0\leq i\leq\ell-3$
and $\{\ell-3,\ell-1\}$.}
\begin{proof}\pf\ Transport the edges of the arm labelling $f$ along
$h_D^{-1}$ from Lemma~\ref{lem:rem-model-supports}.~\qed\end{proof}
\qedstep
\begin{proof}\pf\ Use \stepref{rADbij}, \stepref{rADedges},
\stepref{rAentries}, and Definition~\ref{def:rem-matrices}.~\qed\end{proof}
\end{proof}
\step{rAE}{For $T(1,2,\ell-4)$, use its arm bijection $f$ and set
\[
 \begin{split}
 e(3)&=f(*),\quad e(1)=f(1,1),\quad e(2)=f(2,1),\quad e(0)=f(2,2),\\
 e(i)&=f(3,i-3)\quad(4\leq i\leq\ell-1).
 \end{split}
\]
Then $M(e(i),e(j))=\typ E_\ell(i,j)$.}
\begin{proof}
\step{rAEbij}{The formula for $e$ is a bijection.}
\begin{proof}\pf\ It is $f\circ h_E^{-1}$, a composition of bijections by
Lemma~\ref{lem:rem-model-supports}. ~\qed\end{proof}
\step{rAEedges}{Its edges become $\{0,2\},\{1,3\}$ and
$\{i,i+1\}$ for $2\leq i\leq\ell-2$.}
\begin{proof}\pf\ Transport the edges of $f$ along $h_E^{-1}$
from Lemma~\ref{lem:rem-model-supports}. ~\qed\end{proof}
\qedstep
\begin{proof}\pf\ Use \stepref{rAEbij}, \stepref{rAEedges},
\stepref{rAentries}, and Definition~\ref{def:rem-matrices}.~\qed\end{proof}
\end{proof}
\step{rABC}{For a $\typ{BC}_\ell$ support, there is a path bijection $e$
for which $M(e(i),e(j))$ is $\typ B_\ell(i,j)$ or, when $\ell\geq3$,
$\typ C_\ell(i,j)$.}
\begin{proof}
\step{rABCorient}{\pick\ a path labelling putting the terminal multiple
edge between positions $\ell-2,\ell-1$.}
\begin{proof}\pf\ Reverse the given path when necessary, by the
reversal assertion of Lemma~\ref{lem:path-labelling}.~\qed\end{proof}
\step{rAB}{\assume{its last forward entry is $-2$.}
\prove{The matrix is $\typ B_\ell$.}}
\begin{proof}\pf\ Use \stepref{rAentries} and Definition~\ref{def:rem-matrices}.~\qed\end{proof}
\step{rAC}{\assume{that entry is $-1$ and $\ell\geq3$.}
\prove{It is $\typ C_\ell$.}}
\begin{proof}\pf\ Use \stepref{rAentries} and Definition~\ref{def:rem-matrices}.~\qed\end{proof}
\step{rAtwo}{\assume{that entry is $-1$ and $\ell=2$.}
\prove{Swapping the indices
gives $\typ B_2$.}}
\begin{proof}\pf\ Transpose the two-entry pair in \stepref{rAentries}.~\qed\end{proof}
\qedstep
\begin{proof}\pf\ The entry orders and rank alternatives in
\stepref{rAB}, \stepref{rAC}, \stepref{rAtwo} are exhaustive. ~\qed\end{proof}
\end{proof}
\step{rAF}{For a $\typ F_4$ support, there is a path bijection $e$
with $M(e(i),e(j))=\typ F_4(i,j)$ for every $i,j$.}
\begin{proof}
\step{rAForient}{\pick\ the direction in which the central forward entry is $-2$.}
\begin{proof}\pf\ Reversal in Lemma~\ref{lem:path-labelling}
interchanges the two entries of the central edge in \stepref{rAentries}.~\qed\end{proof}
\qedstep
\begin{proof}\pf\ The two unit edges and the central pair in
\stepref{rAForient} are exactly the entries in
Lemma~\ref{lem:rem-model-supports}; diagonals and nonedges agree.~\qed\end{proof}
\end{proof}
\step{rAG}{For a $\typ G_2$ support, there is a bijection $e$ with
$M(e(i),e(j))=\typ G_2(i,j)$ for every $i,j$.}
\begin{proof}
\step{rAGorient}{\pick\ the endpoint order with forward entry $-3$.}
\begin{proof}\pf\ One of the two orders has that entry by \stepref{rAentries}.~\qed\end{proof}
\qedstep
\begin{proof}\pf\ Step \stepref{rAGorient} gives the off-diagonal pair
$(-3,-1)$, with diagonal 2, exactly as in
Lemma~\ref{lem:rem-model-supports}.~\qed\end{proof}
\end{proof}
\qedstep
\begin{proof}\pf\ The alternatives in \stepref{rAsupport} are precisely
\stepref{rApath}--\stepref{rAG}.~\qed\end{proof}
\end{proof}
\step{rAconverse}{Every listed matrix is finite Cartan and indecomposable.}
\begin{proof}
\pflet{For a listed matrix $M$ take its positive integer symmetrizer from \eqref{eq:symmetrizers}.}

\step{rAunitbalance}{On each unit edge $d(i)=d(j)$ and both entries are $-1$.}
\begin{proof}\pf\ Read Definition~\ref{def:rem-matrices} and
\eqref{eq:symmetrizers} on the unit edges. ~\qed\end{proof}
\step{rABbalance}{On the $\typ B$ multiple edge the two products are
$1(-2)=2(-1)$.}
\begin{proof}\pf\ Substitute its two entries and symmetrizer values. ~\qed\end{proof}
\step{rACbalance}{On the $\typ C$ multiple edge they are $2(-1)=1(-2)$.}
\begin{proof}\pf\ Substitute its two entries and symmetrizer values. ~\qed\end{proof}
\step{rAFbalance}{On the $\typ F$ multiple edge they are $1(-2)=2(-1)$.}
\begin{proof}\pf\ Substitute its two entries and symmetrizer values. ~\qed\end{proof}
\step{rAGbalance}{On the $\typ G$ edge they are $1(-3)=3(-1)$.}
\begin{proof}\pf\ Substitute its two entries and symmetrizer values. ~\qed\end{proof}
\step{rAnonedge}{On every nonedge both products vanish.}
\begin{proof}\pf\ Use the unspecified-entry rule in
Definition~\ref{def:rem-matrices}.~\qed\end{proof}
\step{rAsymmetric}{$DM$ is symmetric and $D^{1/2}MD^{-1/2}=2I-B$
for the corresponding Dynkin graph.}
\begin{proof}\pf\ Combine the entry equalities
\stepref{rAunitbalance}--\stepref{rAnonedge}, including the diagonal 2,
with the labelled weighted supports of Lemma~\ref{lem:rem-model-supports}. ~\qed\end{proof}
\step{rApositive}{$2I-B$ is positive definite over $\mathbb R$.}
\begin{proof}
\step{rAlambda}{$\lambda(B)<2$.}
\begin{proof}\pf\ Lemma~\ref{lem:rem-model-supports} gives a
weight-preserving Dynkin labelling; transfer the bound in
Theorem~\ref{thm:dichotomy} by Lemma~\ref{lem:rem-rayleigh-reindex}.~\qed\end{proof}
\step{rAquadratic}{For $\Zx\ne0$,
$\Zx^T(2I-B)\Zx\geq(2-\lambda(B))\rsq\Zx>0$.}
\begin{proof}\pf\ Use the quadratic-form bound in Lemma~\ref{lem:rayleigh}
and \stepref{rAlambda}.~\qed\end{proof}
\qedstep
\begin{proof}\pf\ Symmetry from \stepref{rAsymmetric} and
\stepref{rAquadratic} are positive definiteness.~\qed\end{proof}
\end{proof}
\step{rAinteger}{$DM$ is positive on nonzero integer vectors.}
\begin{proof}\pf\ Substitute $D^{1/2}\Zx\ne0$ in
\stepref{rApositive}, using \stepref{rAsymmetric}.~\qed\end{proof}
\step{rAfinite}{The matrix is finite Cartan.}
\begin{proof}\pf\ Assemble its four defining fields from
Definition~\ref{def:rem-matrices}, \stepref{rAsymmetric}, \stepref{rAinteger}
and \path{Matrix.posDef_iff_dotProduct_mulVec}.~\qed\end{proof}
\step{rAindecomp}{The matrix is indecomposable.}
\begin{proof}\pf\ Apply Lemma~\ref{lem:connected} to the symmetric zero relation
and connected support in Lemma~\ref{lem:rem-model-supports}.~\qed\end{proof}
\qedstep
\begin{proof}\pf\ Combine \stepref{rAfinite}, \stepref{rAindecomp}.~\qed\end{proof}
\end{proof}
\step{rAunique}{\assume{$P,Q$ are two permitted named matrices equal
after reindexing by a bijection.}
\prove{Their types and ranks agree.}}
\begin{proof}
\step{rAcounts}{$\mathcal I(P)=\mathcal I(Q)$; in particular their ranks agree.}
\begin{proof}\pf\ Apply Lemma~\ref{lem:rem-invariants} to the
reindexing equality in \asm{1}.~\qed\end{proof}
\step{rAuniqueG}{\case{$t(P)=2$.}
\prove{$P,Q$ both have type $\typ G_2$.}}
\begin{proof}\pf\ Only the $\typ G_2$ row in
\eqref{eq:rem-invariant-table} has $t=2$; use \stepref{rAcounts}.~\qed\end{proof}
\step{rAuniqueF}{\case{$u(P)=2$ and
$(a_{\rm in}(P),a_{\rm out}(P))=(0,0)$.}
\prove{$P,Q$ both have type $\typ F_4$.}}
\begin{proof}\pf\ Among the four double-edge rows of
\eqref{eq:rem-invariant-table}, only $\typ F_4$ has this pair;
use \stepref{rAcounts}.~\qed\end{proof}
\step{rAuniqueBC}{\case{$u(P)=2$ and
$(a_{\rm in}(P),a_{\rm out}(P))\ne(0,0)$.}
\prove{$P,Q$ have the same $\typ B$ or $\typ C$ type and rank.}}
\begin{proof}
\step{rAuniqueBtwo}{\case{$n(P)=2$.}}
\begin{proof}\pf\ Only $\typ B_2$ among the permitted double-edge
rows has this size in \eqref{eq:rem-invariant-table};
apply \stepref{rAcounts}.~\qed\end{proof}
\step{rAuniqueBmany}{\case{$n(P)\geq3$ and $a_{\rm in}(P)=1$.}}
\begin{proof}\pf\ The pair is $(1,0)$ only for the $\typ B_\ell$
row at these ranks in \eqref{eq:rem-invariant-table};
apply \stepref{rAcounts}.~\qed\end{proof}
\step{rAuniqueCmany}{\case{$n(P)\geq3$ and $a_{\rm in}(P)=0$.}}
\begin{proof}\pf\ Under \asm{2}, the pair is $(0,1)$,
which singles out $\typ C_\ell$ in \eqref{eq:rem-invariant-table};
apply \stepref{rAcounts}.~\qed\end{proof}
\qedstep
\begin{proof}\pf\ The three eligible rows give exactly
\stepref{rAuniqueBtwo}--\stepref{rAuniqueCmany}.~\qed\end{proof}
\end{proof}
\step{rAuniqueA}{\case{$t(P)=u(P)=b(P)=0$.}
\prove{$P,Q$ both have type $\typ A$ at the same rank.}}
\begin{proof}\pf\ Only the $\typ A_\ell$ row has these three zero
counts in \eqref{eq:rem-invariant-table}; use \stepref{rAcounts}.~\qed\end{proof}
\step{rAuniqueDE}{\case{$t(P)=u(P)=0$ and $b(P)=1$.}
\prove{$P,Q$ have the same $\typ D$ or $\typ E$ type and rank.}}
\begin{proof}
\step{rAuniqueDfour}{\case{$s(P)=3$.}}
\begin{proof}\pf\ Only the $\typ D_4$ row has $s=3$ in
\eqref{eq:rem-invariant-table}; use \stepref{rAcounts}.~\qed\end{proof}
\step{rAuniqueDmany}{\case{$s(P)=2$.}}
\begin{proof}\pf\ Only the $\typ D_\ell$, $\ell\geq5$, row has $s=2$
in \eqref{eq:rem-invariant-table}; size in \stepref{rAcounts} fixes $\ell$.~\qed\end{proof}
\step{rAuniqueE}{\case{$s(P)=1$.}}
\begin{proof}\pf\ The remaining rows in \eqref{eq:rem-invariant-table}
are $\typ E_6,\typ E_7,\typ E_8$; their sizes $6,7,8$ distinguish
them by \stepref{rAcounts}.~\qed\end{proof}
\qedstep
\begin{proof}\pf\ The branched rows give precisely the three values
in \stepref{rAuniqueDfour}--\stepref{rAuniqueE}.~\qed\end{proof}
\end{proof}
\qedstep
\begin{proof}\pf\ The table exhausts the cases
\stepref{rAuniqueG}--\stepref{rAuniqueDE}; all parameter equalities
follow from the size equality in \stepref{rAcounts}.~\qed\end{proof}
\end{proof}
\qedstep
\begin{proof}\pf\ Combine \stepref{rAlabelling}, \stepref{rAconverse},
\stepref{rAunique}, composing the two labellings when comparing
two models of the original matrix.~\qed\end{proof}
\end{proof}

\begin{namedresult}{Remark}{(small ranks)}\label{rem:small}
\lean{CartanMatrix.small_rank_identifications}
\uses{def:rem-matrices,lem:rem-model-supports,facts:rem-library,lem:affine,lem:symmetrization,def:rem-bases,lem:connected,lem:monotonicity}
\assume{$\ell$ is a natural number.}
Here the symbols $\typ A_\ell,\typ B_\ell,\typ C_\ell,\typ D_\ell,\typ E_\ell$
denote Mathlib's generalized definitions at all natural ranks. An old-index
order $(i_0,\dots,i_{\ell-1})$ lists, for each new index $k$, its old index
$i_k$: the identification with $N$ means $N(k,l)=M(i_k,i_l)$.
\prove{$\typ B_1=\typ C_1=\typ D_1=\typ A_1=[2]$,
$\typ D_2=2I_2$, and the following old-index orders give the identifications:
$\typ C_2\cong\typ B_2$ by $(1,0)$;
$\typ D_3\cong\typ A_3$ by $(1,0,2)$;
$\typ E_1=\typ A_1$, $\typ E_2=2I_2$;
$\typ E_3\cong\typ A_2\oplus\typ A_1$ by $(0,2,1)$;
$\typ E_4\cong\typ A_4$ by $(0,2,3,1)$;
$\typ E_5\cong\typ D_5$ by $(0,2,3,1,4)$.
The matrices $\typ E_\ell$, $\ell\geq9$, are not finite Cartan.
All five parameterized definitions at rank 0 give the empty matrix.}
\end{namedresult}
\begin{proof}
\pfsketch\ The small identities are entrywise substitutions; the long
$\typ E$ arm contains an affine obstruction.
\step{rsmallone}{All stated rank-1 matrices equal $[2]$.}
\begin{proof}\pf\ The only index pair is diagonal in the definitions in
\ref{facts:rem-library}.~\qed\end{proof}
\step{rsmallDtwo}{$\typ D_2=2I_2$.}
\begin{proof}\pf\ Apply \path{CartanMatrix.D_two}.~\qed\end{proof}
\step{rsmallBC}{The order $(1,0)$ sends $\typ C_2$ to $\typ B_2$.}
\begin{proof}\pf\ Substitute in \path{CartanMatrix.B_two},
\path{CartanMatrix.C_two}.~\qed\end{proof}
\step{rsmallDthree}{The order $(1,0,2)$ sends $\typ D_3$ to $\typ A_3$.}
\begin{proof}\pf\ Apply \path{CartanMatrix.D_three'} in
\path{Mathlib/LinearAlgebra/Matrix/Cartan.lean}; its transposition
of indices 0 and 1 is the order $(1,0,2)$.~\qed\end{proof}
\step{rsmallEtwo}{$\typ E_2=2I_2$.}
\begin{proof}\pf\ None of the off-diagonal edge clauses of
\path{CartanMatrix.E} has indices less than 2. ~\qed\end{proof}
\step{rsmallEthree}{At rank 3 its sole edge is $\{0,2\}$.}
\begin{proof}\pf\ Evaluate the edge clauses of \path{CartanMatrix.E}.~\qed\end{proof}
\step{rsmallEfour}{At rank 4 its path order is $(0,2,3,1)$.}
\begin{proof}\pf\ Evaluate those edge clauses at indices less than 4. ~\qed\end{proof}
\step{rsmallEfive}{At rank 5 its edges in order $(0,2,3,1,4)$ are
exactly the $\typ D_5$ edges.}
\begin{proof}\pf\ Substitute that order into Definition~\ref{def:rem-matrices}.~\qed\end{proof}
\step{rsmallElarge}{For $\ell\geq9$ the leading nine vertices contain
$T(1,2,5)$.}
\begin{proof}\pf\ Restrict $h_E$ in Lemma~\ref{lem:rem-model-supports} to
the first five positions of its long arm and retain both short arms.~\qed\end{proof}
\step{rsmallnotfinite}{These $\typ E_\ell$ are not finite Cartan.}
\begin{proof}
\step{rsmallconnected}{Their support is connected.}
\begin{proof}\pf\ Apply the connectedness assertion of
Lemma~\ref{lem:rem-model-supports}, which includes every $\ell\geq9$.~\qed\end{proof}
\step{rsmallindecomposable}{The matrices are indecomposable.}
\begin{proof}\pf\ Apply Lemma~\ref{lem:connected} to
\stepref{rsmallconnected} and their symmetric entries. ~\qed\end{proof}
\step{rsmalllower}{$\lambda\geq2$ for their unit support.}
\begin{proof}\pf\ Apply Lemmas~\ref{lem:affine}, \ref{lem:monotonicity}
to \stepref{rsmallElarge}.~\qed\end{proof}
\qedstep
\begin{proof}\pf\ A finite Cartan hypothesis with
\stepref{rsmallindecomposable} would give $\lambda<2$ by
Lemma~\ref{lem:symmetrization}, contradicting \stepref{rsmalllower}.~\qed\end{proof}
\end{proof}
\step{rsmallempty}{At rank 0 there is only one matrix.}
\begin{proof}\pf\ Its index type has no elements. ~\qed\end{proof}
\qedstep
\begin{proof}\pf\ The entry checks \stepref{rsmallone}--\stepref{rsmallEfive},
\stepref{rsmallnotfinite}, \stepref{rsmallempty} give the claims. ~\qed\end{proof}
\end{proof}

\paragraph{Further linear algebra interfaces.}\label{facts:rem-dimension}
\begingroup\raggedright
\path{Matrix.IsHermitian.eigenvectorBasis} and
\path{Matrix.IsHermitian.mulVec_eigenvectorBasis}, in
\path{Mathlib/Analysis/Matrix/Spectrum.lean}, give an orthonormal
eigenbasis for a real symmetric matrix on a finite index type.
\path{LinearMap.finrank_range_add_finrank_ker}, in
\path{Mathlib/LinearAlgebra/FiniteDimensional/Lemmas.lean}, is rank-nullity;
\path{finrank_span_singleton}, in
\path{Mathlib/LinearAlgebra/FiniteDimensional/Basic.lean}, gives dimension
1 for the span of a nonzero vector.
\path{LinearIndependent.span_eq_top_of_card_eq_finrank'}, in the former
file, makes an independent family of the right cardinality spanning.
\path{linearIndependent_of_top_le_span_of_card_eq_finrank}, in
\path{Mathlib/LinearAlgebra/Dimension/OrzechProperty.lean}, proves the
converse implication for a spanning family of the right cardinality.
All applications below are over a field, with finite-dimensional spaces.
\path{Matrix.inv_def}, in
\path{Mathlib/LinearAlgebra/Matrix/NonsingularInverse.lean}, expresses
the inverse as the inverse determinant times the adjugate.
\par\endgroup

\begin{namedresult}{Lemma}{(a Gram realization)}\label{lem:rem-gram}
\lean{Matrix.exists_gram_realization}
\uses{facts:rem-dimension,facts:inner,lem:rem-null}
\assume{$Q$ is real symmetric on a finite set $I$ and
$\Zx^TQ\Zx\geq0$ for every real vector $\Zx$.}
\prove{There is a finite-dimensional Euclidean space $V$ and spanning
vectors $(\Zv_i)_{i\in I}$ with $(\Zv_i,\Zv_j)=Q(i,j)$.
For every family $(\Zw_i)_{i\in I}$ with Gram matrix $Q$ in any
real inner product space, its relations are exactly $\ker Q$ and
$\dim\operatorname{span}\{\Zw_i:i\in I\}=|I|-\dim\ker Q$.
In particular $\dim V=|I|-\dim\ker Q$ for the constructed realization.}
\end{namedresult}
\begin{proof}
\pfsketch\ \new\ Use the orthonormal eigenbasis to construct the
square root map explicitly. This replaces the ordinary proof's quotient
inner-product construction.
\step{rgdiag}{\pick\ an orthonormal eigenbasis $(\Zb_k)$ with eigenvalues $t_k$.}
\begin{proof}\pf\ Apply the eigenbasis declarations in
\ref{facts:rem-dimension}.~\qed\end{proof}
\step{rgnonneg}{$t_k\geq0$ for every $k$.}
\begin{proof}\pf\ Evaluate the nonnegative form at the unit eigenvector
$\Zb_k$ of \stepref{rgdiag}.~\qed\end{proof}
\pflet{$T\Zx=\sum_k\sqrt{t_k}(\Zb_k,\Zx)\Zb_k$,
$V=\operatorname{range}T$, and $\Zv_i=T\Zu_i$ regarded in $V$.}
\step{rgform}{$(T\Zx,T\Zy)=\Zx^TQ\Zy$.}
\begin{proof}\pf\ Expand in the orthonormal basis \stepref{rgdiag} and
use $(\sqrt{t_k})^2=t_k$ from \stepref{rgnonneg}.~\qed\end{proof}
\step{rgentries}{$(\Zv_i,\Zv_j)=Q(i,j)$.}
\begin{proof}\pf\ Specialize \stepref{rgform} to $\Zu_i,\Zu_j$.~\qed\end{proof}
\step{rgspan}{The vectors $\Zv_i$ span $V$.}
\begin{proof}\pf\ Apply the linear map $T$ to the coordinate expansion
$\Zx=\sum_i\Zx(i)\Zu_i$.~\qed\end{proof}
\step{rgrelations}{For any family $(\Zw_i)$ with Gram matrix $Q$,
$\sum_i x(i)\Zw_i=0$ if and only if $Qx=0$.}
\begin{proof}
\step{rgnorm}{$\|\sum_i x(i)\Zw_i\|^2=x^TQx$.}
\begin{proof}\pf\ Bilinearity in \ref{facts:inner} and the Gram entries. ~\qed\end{proof}
\step{rgzero}{The sum is zero if and only if $x^TQx=0$.}
\begin{proof}\pf\ Apply definiteness in \ref{facts:inner} to
\stepref{rgnorm}.~\qed\end{proof}
\qedstep
\begin{proof}\pf\ Lemma~\ref{lem:rem-null} and substitution of $Qx=0$
give the equivalence in \stepref{rgzero}.~\qed\end{proof}
\end{proof}
\step{rgdimension}{For every realizing family $(\Zw_i)$,
$\dim\operatorname{span}\{\Zw_i:i\in I\}=|I|-\dim\ker Q$.}
\begin{proof}\pf\ Apply rank-nullity from \ref{facts:rem-dimension}
to $x\mapsto\sum_i x(i)\Zw_i$: its range is the span by
\ref{facts:inner}, and its kernel is \stepref{rgrelations}.~\qed\end{proof}
\qedstep
\begin{proof}\pf\ Combine \stepref{rgentries}, \stepref{rgspan},
\stepref{rgrelations}, \stepref{rgdimension}.~\qed\end{proof}
\end{proof}

\begin{cgss}{Theorem}{3.8}\label{thm:3.8}\label{cor:D}
\lean{SimpleGraph.smith}\lean{SimpleGraph.smith_gram_delete_basis}
\lean{SimpleGraph.smith_exists_gram}
\lean{SimpleGraph.smithA_maxRayleigh}\lean{SimpleGraph.smithD_maxRayleigh}
\lean{SimpleGraph.smithE₆_maxRayleigh}\lean{SimpleGraph.smithE₇_maxRayleigh}
\lean{SimpleGraph.smithE₈_maxRayleigh}
\uses{def:affine,thm:dichotomy,lem:affine,lem:monotonicity,lem:rem-gram,lem:rem-edge-constant,lem:rem-rayleigh-reindex,lem:rayleigh,facts:rem-dimension,def:rem-weighted}
\assume{$X$ is a finite connected simple graph on $m$ vertices and $B$
is its real adjacency matrix.}
\prove{\textup{(Corollary D of the ordinary proof.)}
$\lambda(B)=2$ if and only if $X$ is isomorphic to a cycle
$\Cyc_s$ ($s\geq3$), $\widetilde{\typ D}_\ell$ ($\ell\geq4$), or one of
$\widetilde{\typ E}_6,\widetilde{\typ E}_7,\widetilde{\typ E}_8$.
For any such graph, $2I-B$ is the Gram matrix of spanning vectors of norm
$\sqrt2$ in dimension $m-1$. In every realization, deleting any one
vector leaves a basis of their span.}
\end{cgss}
\noindent\repair\ The ambient space in the deletion assertion is the span.
The assertion includes every deletion, thus includes the marked vertices
of [CGSS, Figure 1]. Its $A_n$ is our cycle $\Cyc_{n+1}$.
\begin{proof}
\pfsketch\ Strict monotonicity gives equality only at the unweighted
obstructions; a positive null vector controls every deletion.
\step{rDclassification}{$\lambda(B)=2$ exactly for graphs isomorphic
to the listed unweighted obstructions.}
\begin{proof}
\step{rDunit}{Only unweighted obstructions can be contained in $X$.}
\begin{proof}\pf\ A weight $\sqrt2$ or $\sqrt3$ cannot be at most an
adjacency entry 0 or 1, by Definition~\ref{def:rem-weighted}.~\qed\end{proof}
\step{rDforward}{\assume{$\lambda(B)=2$.}
\prove{$X$ is isomorphic to a listed unweighted obstruction.}}
\begin{proof}
\step{rDdynkin}{$X$ has no listed Dynkin labelling.}
\begin{proof}\pf\ Such a labelling would give $\lambda(B)<2$ by
Theorem~\ref{thm:dichotomy} and Lemma~\ref{lem:rem-rayleigh-reindex},
contradicting \asm{2}.~\qed\end{proof}
\step{rDobstruction}{\pick\ an obstruction $Y$ and a containment
$e:Y\preccurlyeq X$.}
\begin{proof}\pf\ Theorem~\ref{thm:dichotomy} and
\stepref{rDdynkin} give the obstruction alternative.~\qed\end{proof}
\step{rDobstructionUnit}{$Y$ is one of the listed unweighted graphs.}
\begin{proof}\pf\ Apply \stepref{rDunit} to \stepref{rDobstruction}.~\qed\end{proof}
\step{rDproper}{The containment $e$ is not proper.}
\begin{proof}\pf\ Otherwise Lemmas~\ref{lem:affine}, \ref{lem:monotonicity}
give $2=\lambda(Y)<\lambda(B)=2$, contradicting \asm{2}.~\qed\end{proof}
\step{rDequal}{$e$ is a graph isomorphism from $Y$ onto $X$.}
\begin{proof}\pf\ By Definition~\ref{def:rem-weighted},
\stepref{rDproper} says $e$ is onto and every entry agrees.
Both graphs have unit edge weights by \stepref{rDobstructionUnit}.~\qed\end{proof}
\qedstep
\begin{proof}\pf\ Use \stepref{rDobstructionUnit}, \stepref{rDequal}.~\qed\end{proof}
\end{proof}
\step{rDbackward}{\assume{$X$ is isomorphic to a listed unweighted graph $Y$.}
\prove{$\lambda(B)=2$.}}
\begin{proof}
\step{rDlisted}{$\lambda(Y)=2$.}
\begin{proof}\pf\ Apply Lemma~\ref{lem:affine}.~\qed\end{proof}
\step{rDtransport}{$\lambda(B)=\lambda(Y)$.}
\begin{proof}\raggedright\pf\ The graph isomorphism in \asm{2}
preserves unit adjacency entries; apply Lemma~\ref{lem:rem-rayleigh-reindex}.~\qed\end{proof}
\qedstep
\begin{proof}\pf\ Combine \stepref{rDlisted}, \stepref{rDtransport}.~\qed\end{proof}
\end{proof}
\qedstep
\begin{proof}\pf\ Steps \stepref{rDforward}, \stepref{rDbackward}
are the two directions of the equivalence.~\qed\end{proof}
\end{proof}
\step{rDgeometry}{\assume{$X$ is isomorphic to a listed graph.}
\prove{The stated realization and every-deletion basis assertion hold.}}
\begin{proof}
\step{rDpositivevector}{\pick\ $\Zw>0$ with $B\Zw=2\Zw$.}
\begin{proof}\pf\ Apply Lemma~\ref{lem:rayleigh} to the connected
adjacency matrix, with $\lambda(B)=2$ from \stepref{rDclassification}
and \asm{1}.~\qed\end{proof}
\pflet{$Q=2I-B$.}
\step{rDform}{For every real vector $\Zx$,
\begin{equation}\label{eq:affine-form}
 \Zx^TQ\Zx=\sum_{\{i,j\}\in E(X)}w(i)w(j)
 \left(\frac{x(i)}{w(i)}-\frac{x(j)}{w(j)}\right)^2.
\end{equation}}
\begin{proof}
\step{rDcross}{The mixed term for each edge is $-2x(i)x(j)$.}
\begin{proof}\pf\ Expand that edge's square in \eqref{eq:affine-form}.~\qed\end{proof}
\step{rDdiagonal}{The coefficient of $x(i)^2$ on the right is 2.}
\begin{proof}\pf\ It is $\sum_{j\sim i}w(j)/w(i)=2$ by the chosen
eigenvector equation. ~\qed\end{proof}
\qedstep
\begin{proof}\pf\ Compare the coefficients in \stepref{rDcross},
\stepref{rDdiagonal} with the quadratic form of $2I-B$.~\qed\end{proof}
\end{proof}
\step{rDpsd}{The form of $Q$ is nonnegative.}
\begin{proof}\pf\ Every summand in \stepref{rDform} is a nonnegative
square times a positive number. ~\qed\end{proof}
\step{rDkernel}{$\ker Q=\mathbb R\Zw$.}
\begin{proof}
\step{rDratios}{\assume{$Q\Zx=0$.}
\prove{$x(i)/w(i)=x(j)/w(j)$ on every edge.}}
\begin{proof}\pf\ Each nonnegative summand in the zero sum
\stepref{rDform} must vanish. ~\qed\end{proof}
\step{rDconstant}{\assume{$Q\Zx=0$.}
\prove{$\Zx$ is a scalar multiple of $\Zw$.}}
\begin{proof}\pf\ Apply
Lemma~\ref{lem:rem-edge-constant} to $i\mapsto x(i)/w(i)$ on the
connected adjacency graph $X$, using \stepref{rDratios}.~\qed\end{proof}
\step{rDreverse}{Every multiple of $\Zw$ is in $\ker Q$.}
\begin{proof}\pf\ Use $Q\Zw=2\Zw-B\Zw=0$.~\qed\end{proof}
\qedstep
\begin{proof}\pf\ Combine \stepref{rDconstant}, \stepref{rDreverse}.~\qed\end{proof}
\end{proof}
\step{rDrealize}{There is a spanning realization of $Q$ in dimension
$m-1$, and each realizing vector has squared norm 2.}
\begin{proof}
\step{rDgram}{A realization exists with dimension $m-\dim\ker Q$.}
\begin{proof}\pf\ Apply Lemma~\ref{lem:rem-gram} to \stepref{rDpsd}.~\qed\end{proof}
\step{rDrank}{$\dim\ker Q=1$.}
\begin{proof}\pf\ Use \stepref{rDkernel}, $\Zw\ne0$, and
\path{finrank_span_singleton}.~\qed\end{proof}
\step{rDnorm}{Every realizing vector has squared norm 2.}
\begin{proof}\pf\ Its squared norm is the diagonal entry $Q(i,i)=2$.~\qed\end{proof}
\qedstep
\begin{proof}\pf\ Combine \stepref{rDgram}, \stepref{rDrank},
\stepref{rDnorm}.~\qed\end{proof}
\end{proof}
\step{rDdelete}{For any realization and any index $i_0$, the remaining
vectors form a basis of its span.}
\begin{proof}
\step{rDrelation}{A relation omitting $i_0$ has coefficient vector $t\Zw$
and zero $i_0$-coordinate.}
\begin{proof}\pf\ Use Lemma~\ref{lem:rem-gram} and \stepref{rDkernel}.~\qed\end{proof}
\step{rDindependent}{All coefficients of this relation vanish.}
\begin{proof}\pf\ In \stepref{rDrelation}, $tw(i_0)=0$ and $w(i_0)>0$
force $t=0$.~\qed\end{proof}
\step{rDomittedspan}{The omitted vector $\Zv_{i_0}$ is in the span
of the remaining vectors.}
\begin{proof}\pf\ By Lemma~\ref{lem:rem-gram} and $Q\Zw=0$,
$\sum_i w(i)\Zv_i=0$; solving with $w(i_0)>0$ gives
$\Zv_{i_0}=-\sum_{i\ne i_0}(w(i)/w(i_0))\Zv_i$.~\qed\end{proof}
\qedstep
\begin{proof}\pf\ The independent remaining family of
\stepref{rDindependent} spans the full realizing span by
\stepref{rDomittedspan}, so it is a basis.~\qed\end{proof}
\end{proof}
\qedstep
\begin{proof}\pf\ Combine \stepref{rDrealize}, \stepref{rDdelete}.~\qed\end{proof}
\end{proof}
\qedstep
\begin{proof}\pf\ Combine \stepref{rDclassification}, \stepref{rDgeometry}.~\qed\end{proof}
\end{proof}

\subsubsection{Coordinate systems and their bases}\label{sec:rem-models}

\begin{namedresult}{Definition}{(classical coordinate models)}\label{def:rem-classical}
\lean{RootSetModel.An}\lean{RootSetModel.Bn}\lean{RootSetModel.Cn}\lean{RootSetModel.Dn}
\uses{def:roots,def:rem-bases}
\assume{$\ell\geq1$ for $\typ A$, $\ell\geq2$ for $\typ B,\typ C$,
and $\ell\geq4$ for $\typ D$.}
In standard coordinates put
\[
\begin{aligned}
\Phi(\typ A_\ell)&=\{\Zu_i-\Zu_j:1\leq i\ne j\leq\ell+1\},&
 V(\typ A_\ell)&=\{\Zx:\sum_i x(i)=0\},\\
\Phi(\typ B_\ell)&=\{\pm\Zu_i\}\cup\{\pm\Zu_i\pm\Zu_j:i<j\},&
 V(\typ B_\ell)&=\mathbb R^\ell,\\
\Phi(\typ C_\ell)&=\{\pm2\Zu_i\}\cup\{\pm\Zu_i\pm\Zu_j:i<j\},&
 V(\typ C_\ell)&=\mathbb R^\ell,\\
\Phi(\typ D_\ell)&=\{\pm\Zu_i\pm\Zu_j:i<j\},&
 V(\typ D_\ell)&=\mathbb R^\ell.
\end{aligned}
\]
The indices in the last three rows range from 1 to $\ell$.
The proposed bases are
\begin{equation}\label{eq:classical-bases}
\begin{aligned}
\typ A:\quad&\alpha_i=\Zu_{i+1}-\Zu_{i+2}&&(0\leq i<\ell);\\
\typ B,\typ C,\typ D:\quad&\alpha_i=\Zu_{i+1}-\Zu_{i+2}&&(0\leq i<\ell-1),
\end{aligned}
\end{equation}
with last root $\Zu_\ell$, $2\Zu_\ell$, or $\Zu_{\ell-1}+\Zu_\ell$,
respectively. We retain the $\typ C_2$ model for comparison with $\typ B_2$.
\end{namedresult}

\begin{namedresult}{Lemma}{(signed coordinate reflections)}\label{lem:rem-coordinate-reflections}
\lean{RootSetModel.coordinate_reflection}
\uses{lem:rem-reflection,facts:inner,def:roots}
\assume{$\Zu_1,\ldots,\Zu_m$ is the standard orthonormal basis of
$\mathbb R^m$, $i\ne j$, and $c\ne0$.}
\prove{Reflection in $\Zu_i-\Zu_j$ interchanges coordinates $i,j$;
reflection in $\Zu_i+\Zu_j$ interchanges and negates both;
reflection in $c\Zu_i$ negates coordinate $i$ and fixes the others.
Signed coordinate permutations preserve each set consisting of all axis
roots of a fixed length, all two-coordinate roots with coefficients
$\pm1$, or all half-coordinate sign vectors. Coordinate permutations
preserve all differences and the zero-sum subspace.}
\end{namedresult}
\begin{proof}
\pfsketch\ Evaluate the reflection formula in each coordinate.
\step{rcoorddiff}{The difference reflection has entries $x(j),x(i)$
at $i,j$ and $x(k)$ elsewhere.}
\begin{proof}\pf\ Substitute $(\Zx,\Zu_i-\Zu_j)=x(i)-x(j)$ and
$\|\Zu_i-\Zu_j\|^2=2$ in Definition~\ref{def:roots}.~\qed\end{proof}
\step{rcoordsum}{The sum reflection has entries $-x(j),-x(i)$
at $i,j$ and $x(k)$ elsewhere.}
\begin{proof}\pf\ Substitute $(\Zx,\Zu_i+\Zu_j)=x(i)+x(j)$ and
squared norm 2 in the reflection formula. ~\qed\end{proof}
\step{rcoordaxis}{The axis reflection has entry $-x(i)$ at $i$
and $x(k)$ elsewhere.}
\begin{proof}\pf\ Substitute $(\Zx,c\Zu_i)=cx(i)$ and squared norm
$c^2$ in the reflection formula. ~\qed\end{proof}
\step{rcoordpatterns}{Signed permutations preserve the three asserted
coordinate patterns.}
\begin{proof}\pf\ They preserve support size and the absolute value
of every nonzero coordinate. ~\qed\end{proof}
\step{rcoordpermutation}{Coordinate permutations preserve differences
and the zero-sum subspace.}
\begin{proof}\pf\ Reindex their two nonzero coordinates and the coordinate sum. ~\qed\end{proof}
\qedstep
\begin{proof}\pf\ Combine \stepref{rcoorddiff}--\stepref{rcoordpermutation}.~\qed\end{proof}
\end{proof}

\begin{namedresult}{Lemma}{(classical Root Systems)}\label{lem:classical}
\lean{RootSetModel.classical_isRootSet}\lean{RootSetModel.classical_base}
\lean{RootSetModel.classical_cartanMatrix}
\uses{def:rem-classical,def:rem-matrices,lem:rem-coordinate-reflections,lem:rem-reflection,facts:inner,facts:rem-dimension,def:rem-bases,def:roots}
\assume{The parameters, spaces, root sets and candidate bases are those
of Definition~\ref{def:rem-classical}.}
\prove{Each displayed set is a Root System in its specified space.
The displayed family is a root base and has exactly Mathlib's Cartan
matrix of that type. The dimension of each space is $\ell$.}
\end{namedresult}
\begin{proof}
\pfsketch\ Check reflection closure by coordinate operations and prove
the base property by nonnegative telescoping expansions.
\step{rclfinite}{Each root set is finite, contains no zero, and is reduced.}
\begin{proof}
\step{rclfiniteonly}{Each set is finite and excludes zero.}
\begin{proof}\pf\ Its finite index and sign choices in
Definition~\ref{def:rem-classical} have a specified nonzero coordinate. ~\qed\end{proof}
\step{rclaxisreduce}{On an axis the only listed vectors are one sign pair.}
\begin{proof}\pf\ Read the axis clause in Definition~\ref{def:rem-classical}.~\qed\end{proof}
\step{rclpairreduce}{On a two-coordinate line the absolute values of
both nonzero coordinates are 1.}
\begin{proof}\pf\ Read the two-coordinate clause in
Definition~\ref{def:rem-classical}.~\qed\end{proof}
\qedstep
\begin{proof}\pf\ Proportional vectors have identical supports, and
\stepref{rclaxisreduce}, \stepref{rclpairreduce} force the scalar to be
$1$ or $-1$; combine with \stepref{rclfiniteonly}.~\qed\end{proof}
\end{proof}
\step{rclreflection}{Each set is closed under all its root reflections.}
\begin{proof}
\step{rclrefA}{In $\typ A$ each reflection permutes coordinates.}
\begin{proof}\pf\ Apply the difference case of
Lemma~\ref{lem:rem-coordinate-reflections}.~\qed\end{proof}
\step{rclrefB}{In $\typ B$ every reflection is a signed permutation
preserving both its patterns.}
\begin{proof}\pf\ Apply the axis, sum and difference formulas of
Lemma~\ref{lem:rem-coordinate-reflections} to its listed roots. ~\qed\end{proof}
\step{rclrefC}{In $\typ C$ every reflection is a signed permutation
preserving both its patterns.}
\begin{proof}\pf\ Apply that lemma with axis scalar $c=2$, and its sum
and difference formulas, to the $\typ C$ list. ~\qed\end{proof}
\step{rclrefD}{In $\typ D$ every reflection is a signed permutation
preserving its two-coordinate pattern.}
\begin{proof}\pf\ Apply the sum and difference formulas of that lemma
to the $\typ D$ list. ~\qed\end{proof}
\qedstep
\begin{proof}\pf\ Use \stepref{rclrefA}--\stepref{rclrefD} and
$s_{-\alpha}=s_\alpha$ from Lemma~\ref{lem:rem-reflection}.~\qed\end{proof}
\end{proof}
\step{rclinteger}{Every root pairing is integral.}
\begin{proof}
\step{rclpairlengthtwo}{Pairing against a two-coordinate root is an
integer inner product.}
\begin{proof}\pf\ All coordinates are integers and its squared norm is
2 in \eqref{eq:pairing}.~\qed\end{proof}
\step{rclpairB}{Pairing against $\pm\Zu_i$ is $\pm2x(i)\in\mathbb Z$.}
\begin{proof}\pf\ Substitute squared norm 1 in \eqref{eq:pairing}.~\qed\end{proof}
\step{rclpairC}{Pairing against $\pm2\Zu_i$ is $\pm x(i)\in\mathbb Z$.}
\begin{proof}\pf\ Substitute squared norm 4 in \eqref{eq:pairing}.~\qed\end{proof}
\qedstep
\begin{proof}\pf\ The denominator-root cases
\stepref{rclpairlengthtwo}, \stepref{rclpairB}, \stepref{rclpairC}
exhaust all four root lists. ~\qed\end{proof}
\end{proof}
\step{rclexpansions}{Every listed root has coefficients of one sign in
the proposed family.}
\begin{proof}
\step{rclexpdiff}{For $i<j$ in all four types,
$\Zu_i-\Zu_j=\sum_{k=i-1}^{j-2}\alpha_k$.}
\begin{proof}\pf\ Telescope the consecutive differences in
\eqref{eq:classical-bases}.~\qed\end{proof}
\step{rclexpBaxis}{In $\typ B$, $\Zu_i=\sum_{k=i-1}^{\ell-1}\alpha_k$.}
\begin{proof}\pf\ Telescope to the endpoint $\alpha_{\ell-1}=\Zu_\ell$.~\qed\end{proof}
\step{rclexpBsum}{In $\typ B$, for $i<j$,
$\Zu_i+\Zu_j=\sum_{k=i-1}^{j-2}\alpha_k+
2\sum_{k=j-1}^{\ell-1}\alpha_k$.}
\begin{proof}\pf\ Add \stepref{rclexpdiff} to twice \stepref{rclexpBaxis}
at $j$.~\qed\end{proof}
\step{rclexpCaxis}{In $\typ C$,
$2\Zu_i=2\sum_{k=i-1}^{\ell-2}\alpha_k+\alpha_{\ell-1}$.}
\begin{proof}\pf\ Telescope to $\alpha_{\ell-1}=2\Zu_\ell$.~\qed\end{proof}
\step{rclexpCsum}{In $\typ C$, for $i<j$,
$\Zu_i+\Zu_j=\sum_{k=i-1}^{j-2}\alpha_k+
2\sum_{k=j-1}^{\ell-2}\alpha_k+\alpha_{\ell-1}$.}
\begin{proof}\pf\ Add \stepref{rclexpdiff} to \stepref{rclexpCaxis} at $j$.~\qed\end{proof}
\step{rclexpDsum}{In $\typ D$, for $i<j\leq\ell-1$,
\[
\Zu_i+\Zu_j=\sum_{k=i-1}^{j-2}\alpha_k+
2\sum_{k=j-1}^{\ell-3}\alpha_k+\alpha_{\ell-2}+\alpha_{\ell-1}.
\]}
\begin{proof}\pf\ Telescope using
$\alpha_{\ell-2}+\alpha_{\ell-1}=2\Zu_{\ell-1}$
from \eqref{eq:classical-bases}.~\qed\end{proof}
\step{rclexpDlast}{In $\typ D$, for $i<\ell$,
$\Zu_i+\Zu_\ell=\sum_{k=i-1}^{\ell-3}\alpha_k+\alpha_{\ell-1}$.}
\begin{proof}\pf\ Telescope using only
$\alpha_{\ell-1}=\Zu_{\ell-1}+\Zu_\ell$.~\qed\end{proof}
\step{rclexpnegative}{The negative of each displayed root has the
negative of its displayed coefficient vector.}
\begin{proof}\pf\ Negate each linear combination. ~\qed\end{proof}
\qedstep
\begin{proof}\pf\ The root patterns in Definition~\ref{def:rem-classical}
are exactly the vectors in \stepref{rclexpdiff}--\stepref{rclexpDlast}
and their negatives; all displayed coefficients are nonnegative integers. ~\qed\end{proof}
\end{proof}
\step{rclspan}{The roots span the specified $\ell$-dimensional spaces.}
\begin{proof}
\step{rclspanA}{The differences $\Zu_i-\Zu_{\ell+1}$,
$1\leq i\leq\ell$, form a real basis of the zero-sum space.}
\begin{proof}\pf\ The unique coefficients of a zero-sum vector are its
first $\ell$ coordinates, by coordinate substitution. ~\qed\end{proof}
\step{rclspanB}{The $\typ B$ axis roots span $\mathbb R^\ell$.}
\begin{proof}\pf\ They include its standard basis. ~\qed\end{proof}
\step{rclspanC}{The $\typ C$ axis roots span $\mathbb R^\ell$.}
\begin{proof}\pf\ Each standard vector is half a listed axis root. ~\qed\end{proof}
\step{rclspanD}{The $\typ D$ roots span $\mathbb R^\ell$.}
\begin{proof}\pf\ For each $i$, choose $j\ne i$ and use
$2\Zu_i=(\Zu_i+\Zu_j)+(\Zu_i-\Zu_j)$.~\qed\end{proof}
\qedstep
\begin{proof}\pf\ Combine \stepref{rclspanA}--\stepref{rclspanD}.~\qed\end{proof}
\end{proof}
\step{rclbase}{The proposed family is a root base.}
\begin{proof}
\step{rclsimplein}{Every proposed simple root belongs to its root list.}
\begin{proof}\pf\ Substitute the simple-root formulas in
Definition~\ref{def:rem-classical}.~\qed\end{proof}
\step{rclsimplesspan}{The $\ell$ simple roots span their ambient space.}
\begin{proof}\pf\ Combine \stepref{rclexpansions} with \stepref{rclspan}.~\qed\end{proof}
\step{rclindependent}{They are linearly independent.}
\begin{proof}\pf\ Apply the spanning-cardinality lemma of
\ref{facts:rem-dimension} to \stepref{rclsimplesspan} and dimension
$\ell$ in \stepref{rclspan}.~\qed\end{proof}
\qedstep
\begin{proof}\pf\ Use the base definition with \stepref{rclsimplein},
\stepref{rclindependent}, \stepref{rclexpansions}.~\qed\end{proof}
\end{proof}
\step{rclcartan}{The base matrices have the asserted entries.}
\begin{proof}
\step{rcldiffgram}{Consecutive differences have squared norm 2,
inner product $-1$ at consecutive indices, and 0 at nonconsecutive indices.}
\begin{proof}\pf\ Expand
$(\Zu_{i+1}-\Zu_{i+2},\Zu_{j+1}-\Zu_{j+2})$
using orthonormality. ~\qed\end{proof}
\step{rclBgram}{The $\typ B$ last root has squared norm 1, inner product
$-1$ with its predecessor, and zero with the other differences.}
\begin{proof}\pf\ Substitute $\alpha_{\ell-1}=\Zu_\ell$.~\qed\end{proof}
\step{rclCgram}{The $\typ C$ last root has squared norm 4, inner product
$-2$ with its predecessor, and zero with the other differences.}
\begin{proof}\pf\ Substitute $\alpha_{\ell-1}=2\Zu_\ell$.~\qed\end{proof}
\step{rclDgram}{The $\typ D$ last root has squared norm 2, product $-1$
with $\alpha_{\ell-3}$, and zero with all other simple roots.}
\begin{proof}\pf\ Expand the product with $\Zu_{\ell-1}+\Zu_\ell$
using orthonormality, including $1-1=0$ at $\alpha_{\ell-2}$.~\qed\end{proof}
\qedstep
\begin{proof}\pf\ Divide the columns in \stepref{rcldiffgram}--\stepref{rclDgram}
by half their squared lengths, as in \eqref{eq:pairing}, and compare
Definition~\ref{def:rem-matrices}.~\qed\end{proof}
\end{proof}
\qedstep
\begin{proof}\pf\ Assemble Definition~\ref{def:roots} from
\stepref{rclfinite}, \stepref{rclreflection}, \stepref{rclinteger},
\stepref{rclspan}, and add \stepref{rclbase}, \stepref{rclcartan}.~\qed\end{proof}
\end{proof}

\begin{namedresult}{Definition}{(exceptional coordinate data)}\label{def:rem-exceptional}
\lean{RootSetModel.E₆}\lean{RootSetModel.E₇}\lean{RootSetModel.E₈}
\lean{RootSetModel.F₄}\lean{RootSetModel.G₂}
\uses{def:roots,def:rem-bases}
\assume{Coordinates have the standard real inner product.}
In $\mathbb R^8$ define
\begin{equation}\label{eq:E-lattice}
 L=\{\Zx\in\mathbb Z^8:\sum_i x(i)\in2\mathbb Z\}
 \ \cup\
 \{\Zx\in(\mathbb Z+\tfrac12)^8:\sum_i x(i)\in2\mathbb Z\}.
\end{equation}
Put
\[
\begin{aligned}
\Phi(\typ E_8)&=\{\pm\Zu_i\pm\Zu_j:i<j\}\\
&\quad\cup\{\tfrac12\sum_{i=1}^8\epsilon(i)\Zu_i:
                \epsilon(i)\in\{1,-1\},\ \prod_i\epsilon(i)=1\},\\
V(\typ E_7)&=\{\Zx:x(7)+x(8)=0\},\\
V(\typ E_6)&=\{\Zx:x(6)=x(7)=-x(8)\},\\
\Phi(\typ E_j)&=\Phi(\typ E_8)\cap V(\typ E_j)\quad(j=6,7),
\qquad V(\typ E_8)=\mathbb R^8.
\end{aligned}
\]
The ordered $\typ E_8$ family is
\begin{equation}\label{eq:E-base}
\begin{split}
\alpha_0&=\tfrac12(\Zu_1+\Zu_8-\Zu_2-\Zu_3-\Zu_4-\Zu_5-\Zu_6-\Zu_7),\\
\alpha_1&=\Zu_1+\Zu_2,\qquad
\alpha_i=\Zu_i-\Zu_{i-1}\quad(2\leq i\leq7).
\end{split}
\end{equation}
Use its first $j$ vectors for $\typ E_j$.
In $\mathbb R^4$ put
\begin{equation}\label{eq:F-roots}
\Phi(\typ F_4)=\{\pm\Zu_i\pm\Zu_j:i<j\}\cup\{\pm\Zu_i\}
\cup\{\tfrac12\sum_{i=1}^4\epsilon(i)\Zu_i:\epsilon(i)\in\{1,-1\}\}
\end{equation}
with family
\begin{equation}\label{eq:F-base}
(\alpha_0,\alpha_1,\alpha_2,\alpha_3)=
(\Zu_2-\Zu_3,\ \Zu_3-\Zu_4,\ \Zu_4,\
\tfrac12(\Zu_1-\Zu_2-\Zu_3-\Zu_4)).
\end{equation}
In the zero-sum plane $V\subset\mathbb R^3$ put
\begin{equation}\label{eq:G-roots}
\Phi(\typ G_2)=\{\Zu_i-\Zu_j:i\ne j\}\cup
\{\pm(2\Zu_i-\Zu_j-\Zu_k):\{i,j,k\}=\{1,2,3\}\},
\end{equation}
with family
\begin{equation}\label{eq:G-base}
\alpha_0=\Zu_1-2\Zu_2+\Zu_3,\qquad \alpha_1=\Zu_2-\Zu_3.
\end{equation}
All sign choices are independent except for the displayed $\typ E_8$
product condition.
\end{namedresult}

\begin{namedresult}{Lemma}{(the even lattice)}\label{lem:even-lattice}
\lean{RootSetModel.E₈_lattice_even}\lean{RootSetModel.E₈_eq_norm_two}
\uses{def:rem-exceptional,facts:inner}
\assume{$L\subset\mathbb R^8$ is the set in \eqref{eq:E-lattice}.}
\prove{$L$ is an additive subgroup, all its inner products are integers,
all its squared norms are even integers, and its squared-norm-2 vectors
are exactly $\Phi(\typ E_8)$. Every nonzero element has squared norm at least 2.}
\end{namedresult}
\begin{proof}
\pfsketch\ Describe two cosets of the integer even-sum subgroup and
minimize the sum of the eight coordinate squares.
\pflet{$L_0=\{\Zx\in\mathbb Z^8:\sum_i x(i)\in2\mathbb Z\}$ and
$\Zv=\tfrac12\sum_i\Zu_i$.}
\step{relcosets}{$L=L_0\cup(L_0+\Zv)$ and $2\Zv\in L_0$.}
\begin{proof}\pf\ Subtracting $\Zv$ changes the coordinate sum by 4;
$2\Zv$ has integer coordinates and sum 8. ~\qed\end{proof}
\step{relgroup}{$L$ is an additive subgroup.}
\begin{proof}
\step{rel00}{$L_0+L_0\subseteq L_0$.}
\begin{proof}\pf\ Addition preserves integer coordinates and even sums. ~\qed\end{proof}
\step{rel01}{$L_0+(L_0+\Zv)\subseteq L_0+\Zv$.}
\begin{proof}\pf\ Translate the inclusion \stepref{rel00} by $\Zv$.~\qed\end{proof}
\step{rel11}{$(L_0+\Zv)+(L_0+\Zv)\subseteq L_0$.}
\begin{proof}\pf\ Use \stepref{rel00} and $2\Zv\in L_0$ from
\stepref{relcosets}.~\qed\end{proof}
\step{relneg0}{$-L_0=L_0$.}
\begin{proof}\pf\ Negation preserves integer coordinates and even sums. ~\qed\end{proof}
\step{relneg1}{$-(L_0+\Zv)=L_0+\Zv$.}
\begin{proof}\pf\ Write $-\Zx-\Zv=(-\Zx-2\Zv)+\Zv$ using
\stepref{relcosets}, \stepref{relneg0}.~\qed\end{proof}
\qedstep
\begin{proof}\pf\ The two-coset cases \stepref{rel00}--\stepref{relneg1},
commutativity of addition, and $0\in L_0$ give the subgroup laws. ~\qed\end{proof}
\end{proof}
\step{relintegral}{All inner products in $L$ are integers.}
\begin{proof}
\step{relint00}{For $\Zx,\Zy\in L_0$, $(\Zx,\Zy)\in\mathbb Z$.}
\begin{proof}\pf\ Sum products of their integer coordinates. ~\qed\end{proof}
\step{relintv}{For $\Zx\in L_0$, $(\Zx,\Zv)=\sum_i x(i)/2\in\mathbb Z$.}
\begin{proof}\pf\ Use the even-sum definition of $L_0$.~\qed\end{proof}
\step{relintvv}{$\rsq\Zv=2\in\mathbb Z$.}
\begin{proof}\pf\ Sum eight copies of $1/4$.~\qed\end{proof}
\qedstep
\begin{proof}\pf\ Expand
$(\Zx+a\Zv,\Zy+b\Zv)$ for $a,b\in\{0,1\}$ using
\stepref{relint00}, \stepref{relintv}, \stepref{relintvv}.~\qed\end{proof}
\end{proof}
\step{releven}{Every squared norm in $L$ is an even integer.}
\begin{proof}
\step{releven0}{For $\Zx\in L_0$, $\sum_i x(i)^2$ is even.}
\begin{proof}\pf\ Each integer satisfies $x(i)^2\equiv x(i)\pmod2$,
and $\sum_i x(i)$ is even. ~\qed\end{proof}
\step{releven1}{For $\Zx\in L_0$,
$\|\Zx+\Zv\|^2=\sum_i x(i)(x(i)+1)+2$ is even.}
\begin{proof}\pf\ Expand the eight coordinate squares; a product of
consecutive integers is even. ~\qed\end{proof}
\qedstep
\begin{proof}\pf\ Use the coset partition \stepref{relcosets} with
\stepref{releven0}, \stepref{releven1}.~\qed\end{proof}
\end{proof}
\step{relminimum}{A nonzero element of $L$ has squared norm at least 2.}
\begin{proof}\pf\ A positive even integer from \stepref{releven}
is at least 2, using definiteness in \ref{facts:inner}.~\qed\end{proof}
\step{relroots0}{The squared-norm-2 elements of $L_0$ are precisely
the two-coordinate roots $\pm\Zu_i\pm\Zu_j$, $i<j$.}
\begin{proof}
\step{relintbound}{Every nonzero integer coordinate contributes at least
1, and any coordinate of absolute value at least 2 contributes at least 4.}
\begin{proof}\pf\ Square the corresponding integer inequalities. ~\qed\end{proof}
\step{relintsupport}{A sum of integer squares equal to 2 has exactly
two nonzero coordinates, both of absolute value 1.}
\begin{proof}\pf\ Apply \stepref{relintbound} to the finite sum of squares. ~\qed\end{proof}
\qedstep
\begin{proof}\pf\ The patterns in \stepref{relintsupport} have even
coordinate sum and squared norm 2 by substitution. ~\qed\end{proof}
\end{proof}
\step{relroots1}{The squared-norm-2 elements in $L_0+\Zv$ are precisely
the half-coordinate roots with an even number of minus signs.}
\begin{proof}
\step{relhalfbound}{Each of their eight absolute coordinates is at least $1/2$.}
\begin{proof}\pf\ Each coordinate is an integer plus $1/2$.~\qed\end{proof}
\step{relhalfequal}{Squared norm 2 forces all absolute coordinates to be $1/2$.}
\begin{proof}\pf\ The eight lower bounds $1/4$ from
\stepref{relhalfbound} already sum to 2. ~\qed\end{proof}
\step{relhalfparity}{For $t$ minus signs, the coordinate sum is $4-t$.}
\begin{proof}\pf\ Sum $8-t$ copies of $1/2$ and $t$ copies of $-1/2$.~\qed\end{proof}
\qedstep
\begin{proof}\pf\ The sum in \stepref{relhalfparity} is even exactly for
even $t$; all these patterns have squared norm 2 by \stepref{relhalfequal}.~\qed\end{proof}
\end{proof}
\qedstep
\begin{proof}\pf\ Combine \stepref{relgroup}, \stepref{relintegral},
\stepref{releven}, \stepref{relminimum}, \stepref{relroots0},
\stepref{relroots1} with the root list in Definition~\ref{def:rem-exceptional}.~\qed\end{proof}
\end{proof}

\begin{namedresult}{Lemma}{(minimal vectors and an obtuse basis)}\label{lem:minimal-base}
\lean{RootSetModel.coefficients_one_sign_of_minimal}
\uses{facts:inner,def:rem-bases}
\assume{\begin{pfenum}
\item $L$ is an additive subgroup of a finite-dimensional real inner
 product space, $b>0$, and every nonzero $\Zx\in L$ has $\rsq\Zx\geq b$.
\item $(\alpha_i)_{i<m}$ is a real basis generating $L$ over $\mathbb Z$,
 and $(\alpha_i,\alpha_j)\leq0$ for $i\ne j$.
\end{pfenum}}
\prove{Every element of $L$ of squared norm $b$ has integer basis
coefficients all nonnegative or all nonpositive.}
\end{namedresult}
\begin{proof}
\pfsketch\ Split a coefficient vector into disjoint positive and negative parts.
\pflet{For $\Zx=\sum_i c(i)\alpha_i\in L$ with $\rsq\Zx=b$, put
$\Zv=\sum_i\max(c(i),0)\alpha_i$ and
$\Zw=\sum_i\max(-c(i),0)\alpha_i$.}
\step{rmbsplit}{$\Zx=\Zv-\Zw$, with $\Zv,\Zw\in L$.}
\begin{proof}\pf\ Use the integer generation in \asm[2]{0} and
$c=\max(c,0)-\max(-c,0)$.~\qed\end{proof}
\step{rmbcross}{$(\Zv,\Zw)\leq0$.}
\begin{proof}\pf\ Expand the inner product; same-index coefficients
have product zero, and all other summands are nonpositive by \asm[2]{0}.~\qed\end{proof}
\step{rmbone}{At least one of $\Zv,\Zw$ is zero.}
\begin{proof}\pf\ If both were nonzero, \asm[1]{0}, \stepref{rmbsplit},
\stepref{rmbcross} would give
$b=\rsq\Zv+\rsq\Zw-2(\Zv,\Zw)\geq2b>b$.~\qed\end{proof}
\qedstep
\begin{proof}\pf\ Independence in \asm[2]{0} makes every coefficient
of the zero part in \stepref{rmbone} vanish. ~\qed\end{proof}
\end{proof}

\begin{namedresult}{Lemma}{(the three E systems)}\label{lem:E}
\lean{RootSetModel.E_isRootSet}\lean{RootSetModel.E_base}
\lean{RootSetModel.E_cartanMatrix}
\uses{def:rem-exceptional,def:rem-matrices,lem:even-lattice,lem:minimal-base,lem:rem-reflection,facts:inner,def:rem-bases,def:roots}
\assume{$j\in\{6,7,8\}$, with $L$, $V(\typ E_j)$, $\Phi(\typ E_j)$,
and the first $j$ vectors of \eqref{eq:E-base} as in
Definition~\ref{def:rem-exceptional}.}
\prove{$\Phi(\typ E_j)$ is a Root System in $V(\typ E_j)$; those $j$
vectors are a root base and their Cartan matrix is
\path{CartanMatrix.E} at rank $j$. They are also an integer basis for
$L\cap V(\typ E_j)$, and $\dim V(\typ E_j)=j$.}
\end{namedresult}
\begin{proof}
\pfsketch\ Invert the coordinate matrix, then apply the minimum-norm
lemma in each subspace.
\step{rEcoefficient}{The coefficients of every $\Zx\in\mathbb R^8$ are
\begin{equation}\label{eq:E-coeff}
\begin{split}
c(0)&=2x(8),\qquad c(1)=\tfrac12\sum_{i=1}^8x(i)+2x(8),\\
c(2)&=c(1)+x(8)-x(1),\\
c(k)&=\sum_{i=k}^7x(i)+(8-k)x(8)\quad(3\leq k\leq7).
\end{split}
\end{equation}
In particular these formulas are the inverse of the coordinate map.}
\begin{proof}
\step{rEcoord8}{The eighth coordinate of $\sum_i c(i)\alpha_i$ is $x(8)$.}
\begin{proof}\pf\ It is $c(0)/2$ by \eqref{eq:E-base}; use \eqref{eq:E-coeff}.~\qed\end{proof}
\step{rEcoord1}{Its first coordinate is $x(1)$.}
\begin{proof}\pf\ Substitute in $c(0)/2+c(1)-c(2)$.~\qed\end{proof}
\step{rEcoord2}{Its second coordinate is $x(2)$.}
\begin{proof}\pf\ Substitute in $-c(0)/2+c(1)+c(2)-c(3)$.~\qed\end{proof}
\step{rEcoordmiddle}{For $3\leq i\leq6$, its $i$th coordinate is $x(i)$.}
\begin{proof}\pf\ Telescope $-c(0)/2+c(i)-c(i+1)$ using \eqref{eq:E-coeff}.~\qed\end{proof}
\step{rEcoord7}{Its seventh coordinate is $x(7)$.}
\begin{proof}\pf\ Substitute in $-c(0)/2+c(7)$.~\qed\end{proof}
\step{rEunique}{These coefficients are unique.}
\begin{proof}\pf\ The same coordinate equations solve successively for
$c(0),c(7),c(6),c(5),c(4),c(3)$ and then $c(1),c(2)$.~\qed\end{proof}
\qedstep
\begin{proof}\pf\ Combine the exhaustive coordinate checks
\stepref{rEcoord8}--\stepref{rEcoord7} and \stepref{rEunique}.~\qed\end{proof}
\end{proof}
\step{rElattice}{For $\Zx\in L$, all coefficients in \eqref{eq:E-coeff}
are integers.}
\begin{proof}
\step{rEint0}{$c(0)$ is an integer.}
\begin{proof}\pf\ Twice an integer or half-integer is integral. ~\qed\end{proof}
\step{rEint1}{$c(1)$ is an integer.}
\begin{proof}\pf\ Half the even coordinate sum is integral, and so is
$2x(8)$ by \stepref{rEint0}.~\qed\end{proof}
\step{rEint2}{$c(2)$ is an integer.}
\begin{proof}\pf\ Use \stepref{rEint1} and the integer difference
$x(8)-x(1)$ of two coordinates in the same coset. ~\qed\end{proof}
\step{rEintk}{For $3\leq k\leq7$, $c(k)$ is an integer.}
\begin{proof}\pf\ Write it as $\sum_{i=k}^7(x(i)+x(8))$;
each summand is an integer in either coset of \eqref{eq:E-lattice}.~\qed\end{proof}
\qedstep
\begin{proof}\pf\ The index cases \stepref{rEint0}--\stepref{rEintk}
are exhaustive. ~\qed\end{proof}
\end{proof}
\step{rEsubspaces}{All displayed simple roots lie in $L$ and have squared
norm 2. The first $j$ form a real basis of $V(\typ E_j)$ and generate
$L\cap V(\typ E_j)$ over $\mathbb Z$.}
\begin{proof}
\step{rEseven}{$x(7)+x(8)=0$ if and only if $c(7)=0$.}
\begin{proof}\pf\ Substitute $k=7$ in \eqref{eq:E-coeff}.~\qed\end{proof}
\step{rEsix}{The two $\typ E_6$ subspace equations hold if and only if
$c(6)=c(7)=0$.}
\begin{proof}\pf\ Use $c(6)=x(6)+x(7)+2x(8)$ and
\stepref{rEseven}.~\qed\end{proof}
\step{rEsimple}{Every $\alpha_i$ lies in $L$ and has squared norm 2.}
\begin{proof}\pf\ For $\alpha_0$, its six minus signs and eight
half-coordinates satisfy \eqref{eq:E-lattice}; each other vector has
two unit coordinates and even sum. ~\qed\end{proof}
\qedstep
\begin{proof}\pf\ Restrict the inverse bijection \stepref{rEcoefficient}
by \stepref{rEseven}, \stepref{rEsix}; use \stepref{rElattice},
\stepref{rEsimple}, and subgroup closure for integer generation. ~\qed\end{proof}
\end{proof}
\step{rEgram}{The simple Gram matrix is $\typ E_j$.}
\begin{proof}
\step{rEgram01}{$(\alpha_0,\alpha_1)=0$.}
\begin{proof}\pf\ The contributions in coordinates 1 and 2 are
$1/2-1/2=0$.~\qed\end{proof}
\step{rEgram02}{$(\alpha_0,\alpha_2)=-1$.}
\begin{proof}\pf\ The contributions are $-1/2-1/2=-1$.~\qed\end{proof}
\step{rEgram0k}{For $3\leq k\leq7$, $(\alpha_0,\alpha_k)=0$.}
\begin{proof}\pf\ The two coefficients of $\alpha_0$ are both $-1/2$
and cancel in the difference $\alpha_k$.~\qed\end{proof}
\step{rEgram12}{$(\alpha_1,\alpha_2)=0$.}
\begin{proof}\pf\ Expand $(\Zu_1+\Zu_2,\Zu_2-\Zu_1)=1-1$.~\qed\end{proof}
\step{rEgram13}{$(\alpha_1,\alpha_3)=-1$.}
\begin{proof}\pf\ Expand $(\Zu_1+\Zu_2,\Zu_3-\Zu_2)=-1$.~\qed\end{proof}
\step{rEgram1k}{For $k\geq4$, $(\alpha_1,\alpha_k)=0$.}
\begin{proof}\pf\ The supports are disjoint. ~\qed\end{proof}
\step{rEgramjk}{For distinct $i,k\geq2$, the product is $-1$ when
$|i-k|=1$ and zero otherwise.}
\begin{proof}\pf\ Expand $(\Zu_i-\Zu_{i-1},\Zu_k-\Zu_{k-1})$.~\qed\end{proof}
\qedstep
\begin{proof}\pf\ Combine \stepref{rEgram01}--\stepref{rEgramjk},
symmetry, and squared norms 2 from \stepref{rEsubspaces} with
Definition~\ref{def:rem-matrices}.~\qed\end{proof}
\end{proof}
\step{rEbase}{Every norm-2 vector in $L\cap V(\typ E_j)$ has integer
coefficients of one sign.}
\begin{proof}\pf\ Apply Lemma~\ref{lem:minimal-base} with $b=2$, using
Lemma~\ref{lem:even-lattice}, \stepref{rEsubspaces}, \stepref{rEgram}.~\qed\end{proof}
\step{rEreflection}{For $\alpha,\beta\in\Phi(\typ E_j)$,
$s_\alpha\beta\in\Phi(\typ E_j)$.}
\begin{proof}
\step{rEreflectL}{$s_\alpha\beta=\beta-(\alpha,\beta)\alpha\in L$.}
\begin{proof}\pf\ Use integer inner products and subgroup closure in
Lemma~\ref{lem:even-lattice}, since $\rsq\alpha=2$.~\qed\end{proof}
\step{rEreflectV}{$s_\alpha\beta\in V(\typ E_j)$.}
\begin{proof}\pf\ The reflection formula is a linear combination of
$\alpha,\beta$ in that subspace. ~\qed\end{proof}
\step{rEreflectnorm}{$\|s_\alpha\beta\|^2=2$.}
\begin{proof}\pf\ Apply Lemma~\ref{lem:rem-reflection}.~\qed\end{proof}
\qedstep
\begin{proof}\pf\ Lemma~\ref{lem:even-lattice} identifies the vectors in
\stepref{rEreflectL}--\stepref{rEreflectnorm} with the stated root list. ~\qed\end{proof}
\end{proof}
\step{rErootset}{The displayed set is a Root System.}
\begin{proof}
\step{rEfinite}{It is finite and contains no zero.}
\begin{proof}\pf\ Its finite coordinate patterns all have squared norm 2
by Lemma~\ref{lem:even-lattice}.~\qed\end{proof}
\step{rEreduced}{It is reduced.}
\begin{proof}\pf\ Equal positive squared norms force a proportionality
scalar to satisfy $t^2=1$.~\qed\end{proof}
\step{rEintegral}{Its pairings are integers.}
\begin{proof}\pf\ Equation~\eqref{eq:pairing} equals the integer inner
product of Lemma~\ref{lem:even-lattice}.~\qed\end{proof}
\qedstep
\begin{proof}\pf\ Assemble Definition~\ref{def:roots} using
\stepref{rEfinite}, \stepref{rEreduced}, \stepref{rEintegral},
\stepref{rEreflection}, and the spanning roots of \stepref{rEsubspaces}.~\qed\end{proof}
\end{proof}
\qedstep
\begin{proof}\pf\ Combine \stepref{rErootset}, \stepref{rEsubspaces},
\stepref{rEbase}, \stepref{rEgram}, and denominator 2 in \eqref{eq:pairing}.~\qed\end{proof}
\end{proof}

\begin{namedresult}{Lemma}{(the F system)}\label{lem:F}
\lean{RootSetModel.F₄_isRootSet}\lean{RootSetModel.F₄_base}
\lean{RootSetModel.F₄_cartanMatrix}
\uses{def:rem-exceptional,def:rem-matrices,lem:rem-coordinate-reflections,lem:rem-reflection,facts:inner,def:rem-bases}
\assume{$\Phi(\typ F_4)\subset\mathbb R^4$ and $(\alpha_i)_{i<4}$ are
the set and family in \eqref{eq:F-roots}, \eqref{eq:F-base}.}
\prove{This is a Root System with the displayed root base and Cartan
matrix $\typ F_4$. Its first two simple roots have squared length 2
and its last two squared length 1.}
\end{namedresult}
\begin{proof}
\pfsketch\ Only half-coordinate reflections need a new closure check;
an inverse coordinate formula proves the base property.
\step{rFbasic}{The set is finite, nonzero, reduced, spans $\mathbb R^4$,
and its two-coordinate roots have squared length 2, all others length squared 1.}
\begin{proof}
\step{rFfinite}{The set is finite and has no zero.}
\begin{proof}\pf\ Its three finite coordinate patterns have a nonzero
coordinate in \eqref{eq:F-roots}.~\qed\end{proof}
\step{rFlength}{The squared lengths are $1+1=2$, $1$, and $4/4=1$.}
\begin{proof}\pf\ Sum coordinate squares in the three patterns. ~\qed\end{proof}
\step{rFreduced}{Proportional roots differ only by sign.}
\begin{proof}\pf\ Support sizes 2, 1, 4 separate the patterns; within
each pattern the fixed absolute nonzero coordinates force $|t|=1$.~\qed\end{proof}
\step{rFspan}{The roots span $\mathbb R^4$.}
\begin{proof}\pf\ The axis pattern contains the standard basis. ~\qed\end{proof}
\qedstep
\begin{proof}\pf\ Combine \stepref{rFfinite}--\stepref{rFspan}.~\qed\end{proof}
\end{proof}
\step{rFreflection}{All root reflections preserve the set.}
\begin{proof}
\step{rFintegerref}{Reflections in an axis or two-coordinate root preserve it.}
\begin{proof}\pf\ Apply Lemma~\ref{lem:rem-coordinate-reflections} to
the three complete signed patterns in \eqref{eq:F-roots}.~\qed\end{proof}
\pflet{$\Zv=(\Zu_1+\Zu_2+\Zu_3+\Zu_4)/2$.}
\step{rFvformula}{$s_\Zv\Zx=\Zx-(\sum_i x(i))\Zv$.}
\begin{proof}\pf\ Substitute $\rsq\Zv=1$ in the reflection formula. ~\qed\end{proof}
\step{rFvaxis}{For an axis root $\epsilon\Zu_i$, its image is
$\epsilon(\Zu_i-\Zv)$, a half-coordinate root.}
\begin{proof}\pf\ Substitute its coordinate sum $\epsilon$ in
\stepref{rFvformula}.~\qed\end{proof}
\step{rFvopposite}{For $\epsilon(\Zu_i-\Zu_j)$ its image is itself.}
\begin{proof}\pf\ Its coordinate sum is zero in \stepref{rFvformula}.~\qed\end{proof}
\step{rFvequal}{For $\epsilon(\Zu_i+\Zu_j)$ its image is
$-\epsilon(\Zu_k+\Zu_l)$ on the complementary pair.}
\begin{proof}\pf\ Substitute its coordinate sum $2\epsilon$ in
\stepref{rFvformula}.~\qed\end{proof}
\step{rFvhalf0}{A half-coordinate root with zero plus signs maps to its negative.}
\begin{proof}\pf\ It is $-\Zv$, with sum $-2$, in \stepref{rFvformula}.~\qed\end{proof}
\step{rFvhalf1}{A half-coordinate root with one plus sign maps to the
positive axis at that sign.}
\begin{proof}\pf\ Its sum is $-1$, so \stepref{rFvformula} adds $\Zv$.~\qed\end{proof}
\step{rFvhalf2}{A half-coordinate root with two plus signs is fixed.}
\begin{proof}\pf\ Its sum is zero in \stepref{rFvformula}.~\qed\end{proof}
\step{rFvhalf3}{A half-coordinate root with three plus signs maps to
the negative axis at its minus sign.}
\begin{proof}\pf\ Its sum is 1, so \stepref{rFvformula} subtracts $\Zv$.~\qed\end{proof}
\step{rFvhalf4}{A half-coordinate root with four plus signs maps to its negative.}
\begin{proof}\pf\ It is $\Zv$, with sum 2, in \stepref{rFvformula}.~\qed\end{proof}
\step{rFvclosed}{$s_\Zv$ preserves the whole set.}
\begin{proof}\pf\ The cases \stepref{rFvaxis}--\stepref{rFvhalf4}
exhaust \eqref{eq:F-roots}.~\qed\end{proof}
\step{rFhalfref}{Reflection in any half-coordinate root preserves the set.}
\begin{proof}
\step{rFsignmap}{A diagonal sign isometry taking $\Zv$ to that root
permutes $\Phi(\typ F_4)$.}
\begin{proof}\pf\ Apply the signed-pattern assertion of
Lemma~\ref{lem:rem-coordinate-reflections}.~\qed\end{proof}
\qedstep
\begin{proof}\pf\ Conjugate \stepref{rFvclosed} by \stepref{rFsignmap}
using Lemma~\ref{lem:rem-reflection}.~\qed\end{proof}
\end{proof}
\qedstep
\begin{proof}\pf\ Combine \stepref{rFintegerref}, \stepref{rFhalfref}
over the three root patterns. ~\qed\end{proof}
\end{proof}
\step{rFintegrality}{Every root pairing is an integer.}
\begin{proof}
\step{rFpairlongint}{Pairing an integer root against a long root is integral.}
\begin{proof}\pf\ Its value is an inner product of integer coordinates
by \eqref{eq:pairing} and \stepref{rFbasic}.~\qed\end{proof}
\step{rFpairlonghalf}{Pairing a half-coordinate root against a long
root is integral.}
\begin{proof}\pf\ The value is a sum of two signed halves by \eqref{eq:pairing}.~\qed\end{proof}
\step{rFpairaxis}{Pairing any root against an axis root is integral.}
\begin{proof}\pf\ Its value is twice a signed coordinate by
\eqref{eq:pairing}; all coordinates are half-integers. ~\qed\end{proof}
\step{rFpairhalfint}{Pairing an integer root against a half-coordinate
root is integral.}
\begin{proof}\pf\ Twice the inner product is a signed sum of the
integer root's coordinates. ~\qed\end{proof}
\step{rFpairhalfhalf}{Pairing two half-coordinate roots is integral.}
\begin{proof}\pf\ If their signs differ in $t$ places, the value is
$2(4-2t)/4=2-t$.~\qed\end{proof}
\qedstep
\begin{proof}\pf\ The denominator and numerator patterns in
\stepref{rFpairlongint}--\stepref{rFpairhalfhalf} are exhaustive. ~\qed\end{proof}
\end{proof}
\step{rFcoeff}{The inverse coefficients in the displayed family are
\begin{equation}\label{eq:F-coeff}
\begin{aligned}
c(0)&=x(1)+x(2),&c(1)&=2x(1)+x(2)+x(3),\\
c(2)&=3x(1)+x(2)+x(3)+x(4),&c(3)&=2x(1).
\end{aligned}
\end{equation}}
\begin{proof}
\step{rFcoord1}{The first coordinate of the combination is $c(3)/2=x(1)$.}
\begin{proof}\pf\ Substitute \eqref{eq:F-base}, \eqref{eq:F-coeff}.~\qed\end{proof}
\step{rFcoord2}{The second is $c(0)-c(3)/2=x(2)$.}
\begin{proof}\pf\ Substitute \eqref{eq:F-coeff}.~\qed\end{proof}
\step{rFcoord3}{The third is $-c(0)+c(1)-c(3)/2=x(3)$.}
\begin{proof}\pf\ Substitute \eqref{eq:F-coeff}.~\qed\end{proof}
\step{rFcoord4}{The fourth is $-c(1)+c(2)-c(3)/2=x(4)$.}
\begin{proof}\pf\ Substitute \eqref{eq:F-coeff}.~\qed\end{proof}
\qedstep
\begin{proof}\pf\ These four triangular equations uniquely solve for
$c(3),c(0),c(1),c(2)$, giving the asserted inverse. ~\qed\end{proof}
\end{proof}
\step{rFbase}{The displayed family is a root base.}
\begin{proof}
\step{rFbasehalf}{Orient a half-coordinate root to have $x(1)=1/2$.
Its coefficients are integers in $[0,1],[0,2],[0,3],\{1\}$.}
\begin{proof}\pf\ Substitute the three independent signed halves in
\eqref{eq:F-coeff}.~\qed\end{proof}
\step{rFbasefirst}{For an integer root with $x(1)\ne0$, orient it with
$x(1)=1$; all four coefficients are nonnegative integers.}
\begin{proof}\pf\ In \eqref{eq:F-coeff} at most one of $x(2),x(3),x(4)$
is nonzero, and it is $1$ or $-1$.~\qed\end{proof}
\step{rFbasezero}{For $x(1)=0$, orient the first nonzero coordinate
among 2,3,4 positively; the four coefficients are nonnegative integers.}
\begin{proof}
\step{rFprefix}{Every prefix sum of coordinates 2,3,4 is nonnegative.}
\begin{proof}\pf\ It begins with zeros, then 1, and has at most one
further nonzero term, bounded below by $-1$.~\qed\end{proof}
\qedstep
\begin{proof}\pf\ For $x(1)=0$, \eqref{eq:F-coeff} consists of those
three prefix sums and zero; apply \stepref{rFprefix}.~\qed\end{proof}
\end{proof}
\step{rFbaseindependent}{The four simple roots are linearly independent.}
\begin{proof}\pf\ Their coordinate map has the inverse \stepref{rFcoeff}.~\qed\end{proof}
\qedstep
\begin{proof}\pf\ The root patterns in \stepref{rFbasehalf},
\stepref{rFbasefirst}, \stepref{rFbasezero} exhaust the roots up to sign;
use \stepref{rFbaseindependent} and their membership in \eqref{eq:F-roots}.~\qed\end{proof}
\end{proof}
\step{rFmatrix}{The simple squared lengths are $(2,2,1,1)$ and their
Cartan matrix is $\typ F_4$.}
\begin{proof}
\step{rFgram01}{$(\alpha_0,\alpha_1)=-1$.}
\begin{proof}\pf\ Expand the product of $\Zu_2-\Zu_3$ and $\Zu_3-\Zu_4$.~\qed\end{proof}
\step{rFgram12}{$(\alpha_1,\alpha_2)=-1$.}
\begin{proof}\pf\ Expand the product of $\Zu_3-\Zu_4$ and $\Zu_4$.~\qed\end{proof}
\step{rFgram23}{$(\alpha_2,\alpha_3)=-1/2$.}
\begin{proof}\pf\ Read the fourth coordinate of $\alpha_3$.~\qed\end{proof}
\step{rFgram02}{$(\alpha_0,\alpha_2)=0$.}
\begin{proof}\pf\ Their supports are disjoint. ~\qed\end{proof}
\step{rFgram03}{$(\alpha_0,\alpha_3)=0$.}
\begin{proof}\pf\ The equal second and third coordinates of $\alpha_3$
cancel. ~\qed\end{proof}
\step{rFgram13}{$(\alpha_1,\alpha_3)=0$.}
\begin{proof}\pf\ The equal third and fourth coordinates of $\alpha_3$
cancel. ~\qed\end{proof}
\qedstep
\begin{proof}\pf\ Use \stepref{rFbasic} for diagonal entries and
\stepref{rFgram01}--\stepref{rFgram13} in \eqref{eq:pairing};
compare Definition~\ref{def:rem-matrices}.~\qed\end{proof}
\end{proof}
\qedstep
\begin{proof}\pf\ Assemble the Root System conditions from
\stepref{rFbasic}, \stepref{rFreflection}, \stepref{rFintegrality};
add \stepref{rFbase}, \stepref{rFmatrix}.~\qed\end{proof}
\end{proof}

\begin{namedresult}{Lemma}{(the G system)}\label{lem:G}
\lean{RootSetModel.G₂_isRootSet}\lean{RootSetModel.G₂_base}
\lean{RootSetModel.G₂_cartanMatrix}
\uses{def:rem-exceptional,def:rem-matrices,lem:rem-coordinate-reflections,lem:rem-reflection,facts:inner,def:rem-bases,def:roots}
\assume{$V$ is the zero-sum plane in $\mathbb R^3$, with root set
\eqref{eq:G-roots} and family \eqref{eq:G-base}.}
\prove{This is a Root System with the displayed root base and Cartan
matrix $\typ G_2$. Its simple roots 0 and 1 have squared lengths 6 and 2,
respectively.}
\end{namedresult}
\begin{proof}
\pfsketch\ A long-root reflection is a permutation followed by global
negation. Six coefficient pairs give one root from every sign pair.
\step{rGbasic}{There are six short roots of squared length 2 and six long
roots of squared length 6; they are finite, nonzero and reduced and span $V$.}
\begin{proof}
\step{rGshort}{The six distinct-coordinate differences have squared norm 2.}
\begin{proof}\pf\ There are $3\cdot2$ ordered pairs and each squared
norm is $1+1$.~\qed\end{proof}
\step{rGlong}{The six signed roots $(2,-1,-1)$ under coordinate
permutations have squared norm 6.}
\begin{proof}\pf\ Choose the doubled coordinate and a global sign;
their squared norm is $4+1+1$.~\qed\end{proof}
\step{rGreduced}{Proportional listed vectors differ by sign.}
\begin{proof}\pf\ Short and long supports have different sizes; in
either pattern the ratios of nonzero coordinate magnitudes are fixed. ~\qed\end{proof}
\step{rGspan}{The differences span the zero-sum plane.}
\begin{proof}\pf\ A vector in that plane is
$x(1)(\Zu_1-\Zu_3)+x(2)(\Zu_2-\Zu_3)$.~\qed\end{proof}
\qedstep
\begin{proof}\pf\ Combine \stepref{rGshort}--\stepref{rGspan}.~\qed\end{proof}
\end{proof}
\pflet{$\beta=2\Zu_1-\Zu_2-\Zu_3$.}
\step{rGlongformula}{For $\Zx\in V$,
$(\Zx,\beta)=3x(1)$ and
$s_\beta\Zx=(-x(1),-x(3),-x(2))$.}
\begin{proof}
\step{rGinner}{$2x(1)-x(2)-x(3)=3x(1)$.}
\begin{proof}\pf\ Substitute the zero-sum equation of $V$.~\qed\end{proof}
\qedstep
\begin{proof}\pf\ Substitute \stepref{rGinner}, $\rsq\beta=6$ in the
reflection formula, and use $x(1)+x(2)+x(3)=0$.~\qed\end{proof}
\end{proof}
\step{rGreflection}{The set is closed under every root reflection.}
\begin{proof}
\step{rGshortref}{Each short reflection preserves it.}
\begin{proof}\pf\ It is a coordinate interchange by
Lemma~\ref{lem:rem-coordinate-reflections}; both lists are permutation-invariant. ~\qed\end{proof}
\step{rGbetaref}{$s_\beta$ preserves it.}
\begin{proof}\pf\ Its formula in \stepref{rGlongformula} is a
coordinate interchange followed by global negation, preserving both lists. ~\qed\end{proof}
\step{rGlongref}{Each long reflection preserves it.}
\begin{proof}\pf\ Every long root is a signed coordinate permutation
of $\beta$ using only a global sign; conjugate \stepref{rGbetaref}
by Lemma~\ref{lem:rem-reflection}.~\qed\end{proof}
\qedstep
\begin{proof}\pf\ The short and long cases \stepref{rGshortref},
\stepref{rGlongref} exhaust \eqref{eq:G-roots}.~\qed\end{proof}
\end{proof}
\step{rGinteger}{All pairings are integral.}
\begin{proof}
\step{rGshortpair}{Against a short root the pairing is an integer inner product.}
\begin{proof}\pf\ All root coordinates are integers and its norm squared
is 2 by \stepref{rGbasic}.~\qed\end{proof}
\step{rGbetapair}{Against $\beta$ the pairing is $x(1)\in\mathbb Z$.}
\begin{proof}\pf\ Divide twice the inner product in
\stepref{rGlongformula} by 6. ~\qed\end{proof}
\step{rGlongpair}{Against every long root it is integral.}
\begin{proof}\pf\ Transport \stepref{rGbetapair} by the global sign
and coordinate permutation carrying $\beta$ to that root. ~\qed\end{proof}
\qedstep
\begin{proof}\pf\ Combine \stepref{rGshortpair}, \stepref{rGlongpair}.~\qed\end{proof}
\end{proof}
\step{rGcoeff}{The six positive coefficient pairs are
\begin{equation}\label{eq:G-coeff}
(0,1),\ (1,1),\ (1,2),\ (1,0),\ (1,3),\ (2,3).
\end{equation}
They give all roots up to sign.}
\begin{proof}
\step{rGc01}{$\alpha_1=\Zu_2-\Zu_3$.}
\begin{proof}\pf\ Use \eqref{eq:G-base}.~\qed\end{proof}
\step{rGc11}{$\alpha_0+\alpha_1=\Zu_1-\Zu_2$.}
\begin{proof}\pf\ Add the coordinates in \eqref{eq:G-base}.~\qed\end{proof}
\step{rGc12}{$\alpha_0+2\alpha_1=\Zu_1-\Zu_3$.}
\begin{proof}\pf\ Add one more $\alpha_1$ to \stepref{rGc11}.~\qed\end{proof}
\step{rGc10}{$\alpha_0=\Zu_1-2\Zu_2+\Zu_3$.}
\begin{proof}\pf\ Use \eqref{eq:G-base}.~\qed\end{proof}
\step{rGc13}{$\alpha_0+3\alpha_1=\Zu_1+\Zu_2-2\Zu_3$.}
\begin{proof}\pf\ Substitute \eqref{eq:G-base}.~\qed\end{proof}
\step{rGc23}{$2\alpha_0+3\alpha_1=2\Zu_1-\Zu_2-\Zu_3$.}
\begin{proof}\pf\ Substitute \eqref{eq:G-base}.~\qed\end{proof}
\qedstep
\begin{proof}\pf\ The first three vectors and their negatives enumerate
the short pattern; the last three and their negatives enumerate the
long pattern of \eqref{eq:G-roots}.~\qed\end{proof}
\end{proof}
\step{rGmatrix}{The simple Gram matrix is
$\left(\begin{smallmatrix}6&-3\\-3&2\end{smallmatrix}\right)$ and
the Cartan matrix is $\typ G_2$.}
\begin{proof}\pf\ Sum coordinate products in \eqref{eq:G-base} and
divide columns by 3 and 1 according to \eqref{eq:pairing}.~\qed\end{proof}
\step{rGindependent}{The two simple roots are independent.}
\begin{proof}\pf\ Their Gram determinant from \stepref{rGmatrix} is
$12-9=3\ne0$.~\qed\end{proof}
\qedstep
\begin{proof}\pf\ Assemble Definition~\ref{def:roots} from
\stepref{rGbasic}, \stepref{rGreflection}, \stepref{rGinteger};
\stepref{rGcoeff}, \stepref{rGindependent} give the base, and
\stepref{rGmatrix} gives its matrix. ~\qed\end{proof}
\end{proof}

\begin{namedresult}{Lemma}{(indecomposability from a base)}\label{lem:rem-base-connected}
\lean{IsRootSet.isIndecomposable_of_connected_base}
\uses{def:roots,def:rem-bases,lem:rem-edge-constant,facts:inner}
\assume{\begin{pfenum}
\item $\Phi$ is a Root System in $V$ with a nonempty root base $\Delta$.
\item The graph on $\Delta$ joining distinct nonorthogonal roots is connected.
\end{pfenum}}
\prove{$\Phi$ is indecomposable.}
\end{namedresult}
\begin{proof}
\pfsketch\ An orthogonal partition puts all simple roots on the same
side, leaving a vector orthogonal to a spanning set on the other side.
\step{rbcpartition}{\assume{$\Phi=P\sqcup Q$ is a partition into
mutually orthogonal parts.}
\prove{$P=\varnothing$ or $Q=\varnothing$.}}
\begin{proof}
\step{rbcside}{Adjacent base roots belong to the same part.}
\begin{proof}\pf\ Roots in different parts are orthogonal by the partition
hypothesis. ~\qed\end{proof}
\step{rbcwhole}{All base roots belong to one part.}
\begin{proof}\pf\ Apply
Lemma~\ref{lem:rem-edge-constant} to the function recording the part
containing a simple root, on the connected nonorthogonality graph
of $\Delta$, using \stepref{rbcside}.~\qed\end{proof}
\step{rbcother}{The other part is empty.}
\begin{proof}
\step{rbcperp}{A root in the other part is orthogonal to every vector in $V$.}
\begin{proof}\pf\ Expand any vector in the spanning base, using
\stepref{rbcwhole} and bilinearity. ~\qed\end{proof}
\qedstep
\begin{proof}\pf\ Pair such a root with itself in \stepref{rbcperp};
definiteness contradicts its nonzero status. ~\qed\end{proof}
\end{proof}
\qedstep
\begin{proof}\pf\ Step \stepref{rbcother} makes one part empty.~\qed\end{proof}
\end{proof}
\qedstep
\begin{proof}\pf\ The nonempty base and \stepref{rbcpartition}
give indecomposability in Definition~\ref{def:roots}.~\qed\end{proof}
\end{proof}

\begin{namedresult}{Lemma}{(the models are indecomposable)}\label{lem:models-irreducible}
\lean{RootSetModel.isIndecomposable}\lean{RootSetModel.Cn_isIndecomposable}
\uses{lem:classical,lem:E,lem:F,lem:G,lem:rem-model-supports,lem:rem-base-connected}
\assume{The model is $\typ A_\ell$ ($\ell\geq1$), $\typ B_\ell$
or $\typ C_\ell$ ($\ell\geq2$), $\typ D_\ell$ ($\ell\geq4$),
$\typ E_j$ ($j=6,7,8$), $\typ F_4$, or $\typ G_2$, in its specified space.}
\prove{It is indecomposable.}
\end{namedresult}
\begin{proof}
\pfsketch\ Apply a single general base-connectivity lemma to each model,
so this proof never reuses local steps for a second Root System.
\step{rmiA}{The $\typ A$ base graph is a connected nonempty path.}
\begin{proof}\pf\ Its Cartan matrix is $\typ A_\ell$ by
Lemma~\ref{lem:classical}; Lemma~\ref{lem:rem-model-supports}
identifies its connected support with the base nonorthogonality graph.~\qed\end{proof}
\step{rmiBC}{Each $\typ B,\typ C$ base graph is a connected nonempty path.}
\begin{proof}\pf\ Use the two matrix identities in
Lemma~\ref{lem:classical} and the path supports of
Lemma~\ref{lem:rem-model-supports}, including $\typ C_2$.~\qed\end{proof}
\step{rmiD}{The $\typ D$ base graph is the connected tree $T(1,1,\ell-3)$.}
\begin{proof}\pf\ Combine Lemma~\ref{lem:classical} with $h_D$ and
connectedness in Lemma~\ref{lem:rem-model-supports}.~\qed\end{proof}
\step{rmiE}{Each $\typ E$ base graph is the connected tree $T(1,2,j-4)$.}
\begin{proof}\pf\ Combine Lemma~\ref{lem:E} with $h_E$ and
connectedness in Lemma~\ref{lem:rem-model-supports}.~\qed\end{proof}
\step{rmiF}{The $\typ F_4$ base graph is a connected four-vertex path.}
\begin{proof}\pf\ Use the Cartan matrix in Lemma~\ref{lem:F}
and its connected support in Lemma~\ref{lem:rem-model-supports}.~\qed\end{proof}
\step{rmiG}{The $\typ G_2$ base graph is a connected single edge.}
\begin{proof}\pf\ Use the Cartan matrix in Lemma~\ref{lem:G}
and its connected support in Lemma~\ref{lem:rem-model-supports}.~\qed\end{proof}
\qedstep
\begin{proof}\pf\ Apply Lemma~\ref{lem:rem-base-connected}
to the nonempty bases and connected graphs in
\stepref{rmiA}--\stepref{rmiG}.~\qed\end{proof}
\end{proof}

\begin{namedresult}{Lemma}{(root, length and line counts)}\label{lem:counts}
\lean{RootSetModel.root_counts}\lean{RootSetModel.length_counts}
\lean{RootSetModel.numLines}\lean{RootSetModel.one_length_iff}
\lean{RootSetModel.lengthRatio}\lean{RootSetModel.maxLengthSq_eq_lengthRatio_mul}
\uses{def:rem-classical,def:rem-exceptional,lem:classical,lem:E,lem:F,lem:G,lem:rem-reflection}
\assume{The coordinate models are in their construction ranges:
$\ell\geq1$ for $\typ A$, $\ell\geq2$ for $\typ B,\typ C$,
$\ell\geq4$ for $\typ D$, and $j\in\{6,7,8\}$ for $\typ E_j$.}
\prove{The counts are those in the following table. The two-length
models have squared-length ratio 2 in types $\typ B,\typ C,\typ F$ and
3 in type $\typ G$; all other models have one length.}
\begin{center}
\begin{tabular}{c|r|r|r|r}
\label{tab:roots}
Type&one length&short&long&lines\\ \hline
$\typ A_\ell$&$\ell(\ell+1)$&&&$\ell(\ell+1)/2$\\
$\typ B_\ell$&&$2\ell$&$2\ell(\ell-1)$&$\ell^2$\\
$\typ C_\ell$&&$2\ell(\ell-1)$&$2\ell$&$\ell^2$\\
$\typ D_\ell$&$2\ell(\ell-1)$&&&$\ell(\ell-1)$\\
$\typ E_6$&72&&&36\\
$\typ E_7$&126&&&63\\
$\typ E_8$&240&&&120\\
$\typ F_4$&&24&24&24\\
$\typ G_2$&&6&6&6
\end{tabular}
\end{center}
\end{namedresult}
\begin{proof}
\pfsketch\ Count disjoint coordinate patterns, then divide by the
two roots on each line.
\step{rctA}{The $\typ A_\ell$ set has $\ell(\ell+1)$ roots.}
\begin{proof}\pf\ The ordered distinct-coordinate pair is recovered
from its coefficients $1,-1$; count $(\ell+1)\ell$ pairs. ~\qed\end{proof}
\step{rctpairs}{There are $4\binom\ell2=2\ell(\ell-1)$ signed
two-coordinate roots.}
\begin{proof}\pf\ Their support and the two signs give a bijection with
two-element coordinate subsets times $\{1,-1\}^2$.~\qed\end{proof}
\step{rctaxes}{There are $2\ell$ axis roots in $\typ B$ and in $\typ C$.}
\begin{proof}\pf\ The coordinate and sign are uniquely recovered from
each such root in Definition~\ref{def:rem-classical}.~\qed\end{proof}
\step{rctE8int}{The $\typ E_8$ integer pattern has 112 roots.}
\begin{proof}\pf\ Evaluate \stepref{rctpairs} at $\ell=8$.~\qed\end{proof}
\step{rctE8half}{Its half-coordinate pattern has 128 roots.}
\begin{proof}\pf\ Choose seven signs freely; their product determines
the eighth in Definition~\ref{def:rem-exceptional}.~\qed\end{proof}
\step{rctE7int}{The $\typ E_7$ integer pattern has $60+2=62$ roots.}
\begin{proof}
\step{rctE7first}{There are 60 with both coordinates among the first six.}
\begin{proof}\pf\ Evaluate \stepref{rctpairs} at $\ell=6$.~\qed\end{proof}
\step{rctE7last}{There are two using coordinates 7 or 8.}
\begin{proof}\pf\ The equation $x(7)+x(8)=0$ forces the pair
$\pm(\Zu_7-\Zu_8)$ in a two-coordinate root. ~\qed\end{proof}
\qedstep
\begin{proof}\pf\ The disjoint cases \stepref{rctE7first},
\stepref{rctE7last} exhaust the integer pattern. ~\qed\end{proof}
\end{proof}
\step{rctE7half}{The $\typ E_7$ half-coordinate pattern has 64 roots.}
\begin{proof}\pf\ Its last two signs are opposite, the first six have
product $-1$, and five of the first six signs and sign 8 are free. ~\qed\end{proof}
\step{rctE6int}{The $\typ E_6$ integer pattern has 40 roots.}
\begin{proof}
\step{rctE6support}{Its nonzero coordinates are among the first five.}
\begin{proof}\pf\ A nonzero coordinate among 6,7,8 would force all
three to be nonzero, contrary to support size 2. ~\qed\end{proof}
\qedstep
\begin{proof}\pf\ Apply \stepref{rctpairs} at $\ell=5$ to
\stepref{rctE6support}.~\qed\end{proof}
\end{proof}
\step{rctE6half}{The $\typ E_6$ half-coordinate pattern has 32 roots.}
\begin{proof}\pf\ The first five signs are free; their product is
sign 8, and signs 6,7 are its negative. ~\qed\end{proof}
\step{rctFlong}{There are 24 long $\typ F_4$ roots.}
\begin{proof}\pf\ Apply \stepref{rctpairs} at $\ell=4$.~\qed\end{proof}
\step{rctFshort}{There are $8+16=24$ short $\typ F_4$ roots.}
\begin{proof}\pf\ The disjoint axis and half-coordinate patterns have
$2\cdot4$ and $2^4$ independent index/sign choices. ~\qed\end{proof}
\step{rctG}{There are six roots of each length in $\typ G_2$.}
\begin{proof}\pf\ Count $3\cdot2$ ordered distinct differences and
$3\cdot2$ choices of doubled coordinate and global sign in \eqref{eq:G-roots}.~\qed\end{proof}
\step{rctlengths}{The squared lengths are 2 for $\typ A,\typ D,\typ E$;
1 and 2 for $\typ B,\typ F$; 2 and 4 for $\typ C$; 2 and 6 for $\typ G$.}
\begin{proof}\pf\ Sum the coordinate squares in
Definitions~\ref{def:rem-classical}, \ref{def:rem-exceptional}.~\qed\end{proof}
\step{rctlines}{Each line contains exactly two roots.}
\begin{proof}\pf\ Reducedness and closure under negation, from the
model lemmas and Lemma~\ref{lem:rem-reflection}, give precisely a sign pair. ~\qed\end{proof}
\qedstep
\begin{proof}\pf\ Add disjoint pattern counts
\stepref{rctA}--\stepref{rctG}, separate by \stepref{rctlengths},
and divide by 2 using \stepref{rctlines}.~\qed\end{proof}
\end{proof}

\subsubsection{The Mathlib bridge and classification}

\begin{namedresult}{Lemma}{(the finite-root-set pairing)}\label{lem:pairing-construction}
\lean{IsRootSet.toRootPairing}\lean{IsRootSet.isIrreducible_iff}
\uses{def:roots,lem:rem-reflection,facts:inner,facts:rem-library,def:rem-bases}
\assume{$\Phi$ is a Root System in a finite-dimensional real inner
product space $V$.}
\prove{It defines a finite reduced crystallographic root pairing $P$
over $\mathbb R$ with root and coroot spaces both $V$, and $P$ is a
root system in Mathlib's sense. Its pairing is \eqref{eq:pairing}.
If $\Phi\ne\varnothing$, its indecomposability is equivalent to
$P$ being Mathlib-irreducible.}
\end{namedresult}
\begin{proof}
\pfsketch\ Construct the structure on the prescribed index set $\Phi$.
For irreducibility a reflection-invariant subspace assigns each root
to itself or to its orthogonal complement.
\pflet{Use root indices $\Phi$, root embedding $\alpha\mapsto\alpha$,
coroot map $\alpha^\vee=2\alpha/\rsq\alpha$, and bilinear form the real inner product.}
\step{rpcperfect}{This bilinear pairing is perfect.}
\begin{proof}
\step{rpcinjective}{The map from $V$ to its dual is injective.}
\begin{proof}\pf\ If $(\Zx-\Zy,\Zz)=0$ for every $\Zz$, choose
$\Zz=\Zx-\Zy$ and apply definiteness in \ref{facts:inner}.~\qed\end{proof}
\step{rpcflip}{The flipped map is injective.}
\begin{proof}\pf\ Rewrite it as the same map by real inner-product
symmetry and use \stepref{rpcinjective}.~\qed\end{proof}
\qedstep
\begin{proof}\pf\ Apply \path{LinearMap.IsPerfPair.of_injective} to
\stepref{rpcinjective}, \stepref{rpcflip}, using finite dimensionality. ~\qed\end{proof}
\end{proof}
\step{rpccoroot}{The coroot map is injective.}
\begin{proof}
\step{rpctwice}{$2\alpha^\vee/\rsq{\alpha^\vee}=\alpha$ for $\alpha\ne0$.}
\begin{proof}\pf\ Substitute $\rsq{\alpha^\vee}=4/\rsq\alpha$
using bilinearity and $\rsq\alpha>0$.~\qed\end{proof}
\qedstep
\begin{proof}\pf\ A map with the left inverse \stepref{rpctwice} is injective. ~\qed\end{proof}
\end{proof}
\step{rpcself}{$(\alpha,\alpha^\vee)=2$.}
\begin{proof}\pf\ Cancel $\rsq\alpha>0$ in the coroot formula. ~\qed\end{proof}
\step{rpcreflection}{The restricted map $s_\alpha$ is a permutation of
the root indices.}
\begin{proof}\pf\ Apply Lemma~\ref{lem:rem-reflection} and reflection
closure in Definition~\ref{def:roots}.~\qed\end{proof}
\step{rpccorootidentity}{$(s_\alpha\beta)^\vee
=\beta^\vee-(\alpha,\beta^\vee)\alpha^\vee$.}
\begin{proof}
\step{rpcnorm}{$(s_\alpha\beta)^\vee=s_\alpha(\beta^\vee)$.}
\begin{proof}\pf\ The orthogonal linear map of
Lemma~\ref{lem:rem-reflection} commutes with normalization by squared norm. ~\qed\end{proof}
\qedstep
\begin{proof}\pf\ Substitute the reflection formula and
$\alpha^\vee=2\alpha/\rsq\alpha$ in \stepref{rpcnorm}.~\qed\end{proof}
\end{proof}
\step{rpcstructure}{These data define a root pairing $P$ with
$P.p(i,j)=2(\alpha_i,\alpha_j)/\rsq{\alpha_j}$.}
\begin{proof}\pf\ Fill the structure fields of \path{RootPairing}
using \stepref{rpcperfect}, root inclusion, \stepref{rpccoroot},
\stepref{rpcself}, \stepref{rpcreflection}, \stepref{rpccorootidentity}
and the defining root reflection identity. ~\qed\end{proof}
\step{rpcreflectionmaps}{The root and coroot reflection at index
$\alpha$ are both $s_\alpha$ on every vector of $V$.}
\begin{proof}\pf\ The root formula subtracts
$(\Zx,\alpha^\vee)\alpha$; the coroot formula subtracts
$(\alpha,\Zx)\alpha^\vee$. Both equal
$2(\alpha,\Zx)\alpha/\rsq\alpha$ by the coroot definition. ~\qed\end{proof}
\step{rpcproperties}{$P$ is finite, crystallographic, reduced and spanning.}
\begin{proof}
\step{rpcfinite}{The index type is finite.}
\begin{proof}\pf\ It is the finite set $\Phi$ of Definition~\ref{def:roots}.~\qed\end{proof}
\step{rpccrystal}{All its pairings are integral.}
\begin{proof}\pf\ Apply crystallographic integrality in
Definition~\ref{def:roots} with the two roots interchanged to
\stepref{rpcstructure}.~\qed\end{proof}
\step{rpcreduced}{A dependent pair of roots is an equal or opposite pair.}
\begin{proof}
\step{rpcdependence}{For a dependent pair $\alpha,\beta$, choose
$(a,b)\ne(0,0)$ with $a\alpha+b\beta=0$.}
\begin{proof}\pf\ Unfold linear dependence for the two-vector family. ~\qed\end{proof}
\step{rpcsecond}{$b\ne0$.}
\begin{proof}\pf\ Otherwise \stepref{rpcdependence} and $\alpha\ne0$
would force $a=0$.~\qed\end{proof}
\step{rpcproportional}{$\beta=(-a/b)\alpha$.}
\begin{proof}\pf\ Solve the relation \stepref{rpcdependence} using
\stepref{rpcsecond}.~\qed\end{proof}
\qedstep
\begin{proof}\pf\ Apply reducedness in Definition~\ref{def:roots}
to the proportionality \stepref{rpcproportional}.~\qed\end{proof}
\end{proof}
\step{rpcrootspan}{The roots span $V$.}
\begin{proof}\pf\ This is the spanning condition in Definition~\ref{def:roots}.~\qed\end{proof}
\step{rpccorootspan}{The coroots span $V$.}
\begin{proof}\pf\ Each root is a nonzero scalar multiple of its coroot,
so \stepref{rpcrootspan} lies in the coroot span. ~\qed\end{proof}
\qedstep
\begin{proof}\pf\ Assemble the Mathlib predicates from
\stepref{rpcfinite}--\stepref{rpccorootspan} using \ref{facts:rem-library}.~\qed\end{proof}
\end{proof}
\step{rpcirreducible}{\assume{$\Phi$ is nonempty and indecomposable.}
\prove{$P$ is irreducible.}}
\begin{proof}
\step{rpcnontrivial}{$V$ is nontrivial.}
\begin{proof}\pf\ A chosen root is nonzero by Definition~\ref{def:roots}.~\qed\end{proof}
\step{rpcinvariant}{Every reflection-invariant subspace $W$ is $0$ or $V$.}
\begin{proof}
\step{rpcalternative}{For every root $\alpha$, either $\alpha\in W$
or $\alpha\in W^\perp$.}
\begin{proof}
\step{rpcdetect}{\assume{some $\Zw\in W$ has $(\Zw,\alpha)\ne0$.}
\prove{$\alpha\in W$.}}
\begin{proof}\pf\ The vector $\Zw-s_\alpha\Zw$ is a nonzero scalar
multiple of $\alpha$ in $W$.~\qed\end{proof}
\step{rpcundetect}{\assume{no such $\Zw$ exists.}
\prove{$\alpha\in W^\perp$.}}
\begin{proof}\pf\ Unfold the definition of the orthogonal complement. ~\qed\end{proof}
\qedstep
\begin{proof}\pf\ Excluded middle for the detecting vector gives
\stepref{rpcdetect} or \stepref{rpcundetect}.~\qed\end{proof}
\end{proof}
\step{rpcpartition}{The roots in $W$ and the roots in $W^\perp$
partition $\Phi$ into orthogonal parts.}
\begin{proof}\pf\ Use \stepref{rpcalternative} and
$W\cap W^\perp=\{0\}$, excluding zero roots. ~\qed\end{proof}
\step{rpconeempty}{One of those parts is empty.}
\begin{proof}\pf\ Apply indecomposability to \stepref{rpcpartition}.~\qed\end{proof}
\step{rpcallW}{\assume{all roots lie in $W$.}
\prove{$W=V$.}}
\begin{proof}\pf\ Use root spanning in Definition~\ref{def:roots}.~\qed\end{proof}
\step{rpcallperp}{\assume{all roots lie in $W^\perp$.}
\prove{$W=0$.}}
\begin{proof}\pf\ Each $\Zw\in W$ is orthogonal to the spanning root
set, thus to itself; use definiteness. ~\qed\end{proof}
\qedstep
\begin{proof}\pf\ Use \stepref{rpconeempty} to select
\stepref{rpcallW} or \stepref{rpcallperp}.~\qed\end{proof}
\end{proof}
\qedstep
\begin{proof}\pf\ Both copies of $V$ have the same orthogonal reflection
maps by \stepref{rpcreflectionmaps}; \stepref{rpcnontrivial},
\stepref{rpcinvariant} supply both irreducibility fields. ~\qed\end{proof}
\end{proof}
\step{rpcconverse}{\assume{$P$ is irreducible and
$\Phi=A\sqcup B$ is a partition into nonempty mutually orthogonal parts.}
\prove{\textsc{False}.}}
\begin{proof}
\pflet{$W=\operatorname{span}A$.}
\step{rpcWnonzero}{$W\ne0$.}
\begin{proof}\pf\ A member of nonempty $A$ is a nonzero vector in $W$.~\qed\end{proof}
\step{rpcWproper}{$W\ne V$.}
\begin{proof}\pf\ A nonzero root in $B$ is orthogonal to $W$ and
cannot be orthogonal to itself. ~\qed\end{proof}
\step{rpcWA}{A reflection at a root of $A$ preserves $W$.}
\begin{proof}\pf\ Its formula subtracts a scalar multiple of an element
of $A$ from a vector of $W$.~\qed\end{proof}
\step{rpcWB}{A reflection at a root of $B$ fixes $W$ pointwise.}
\begin{proof}\pf\ Orthogonality to $A$ extends linearly to $W$ in the
reflection formula. ~\qed\end{proof}
\qedstep
\begin{proof}\pf\ The invariant nonzero proper subspace in
\stepref{rpcWnonzero}--\stepref{rpcWB} contradicts irreducibility. ~\qed\end{proof}
\end{proof}
\qedstep
\begin{proof}\pf\ Combine \stepref{rpcstructure}, \stepref{rpcproperties},
\stepref{rpcirreducible}, \stepref{rpcconverse}.~\qed\end{proof}
\end{proof}
\noindent\new\ We use the direct structure constructor.
\path{RootPairing.ofBilinear} takes all reflective vectors of a form,
and would not select this prescribed finite root set over $\mathbb R$.

\begin{namedresult}{Lemma}{(one common similarity factor)}\label{lem:similarity}
\lean{IsRootSet.similarity_of_cartanMatrix_eq}
\uses{def:rem-bases,lem:rem-edge-constant,facts:inner,facts:rem-library,def:roots}
\assume{\begin{pfenum}
\item $V,W$ are finite-dimensional real inner product spaces and
 $(\alpha_i)_{i\in I}$, $(\beta_i)_{i\in I}$ are real vector-space bases
 indexed by the same finite nonempty set $I$.
\item The matrices $p_\alpha(i,j)=2(\alpha_i,\alpha_j)/\rsq{\alpha_j}$
 and $p_\beta(i,j)=2(\beta_i,\beta_j)/\rsq{\beta_j}$ agree, and their
 common nonzero off-diagonal support is connected.
\end{pfenum}}
\prove{The linear equivalence $f:V\to W$ taking $\alpha_i$ to $\beta_i$
satisfies $(f\Zx,f\Zy)=a(\Zx,\Zy)$ for a single $a>0$.
For any Root Systems $\Phi\subset V$, $\Psi\subset W$ with
$f(\Phi)=\Psi$, this is a similarity of those Root Systems.}
\end{namedresult}
\begin{proof}
\pfsketch\ Ratios of squared simple-root lengths propagate along the
connected support; bilinearity extends the resulting Gram identity.
\pflet{$q(i)=\rsq{\alpha_i}$, $r(i)=\rsq{\beta_i}$, and $M=p_\alpha=p_\beta$.}
\step{rsimpositive}{$q(i)>0$ and $r(i)>0$ for every $i$.}
\begin{proof}\pf\ A vector in a linearly independent family is nonzero;
use the bases in \asm[1]{0} and definiteness. ~\qed\end{proof}
\step{rsimbalance}{$M(i,j)q(j)=M(j,i)q(i)$ and
$M(i,j)r(j)=M(j,i)r(i)$.}
\begin{proof}\pf\ Use symmetry of each inner product in \eqref{eq:pairing}.~\qed\end{proof}
\step{rsimedge}{On a support edge, $r(i)/q(i)=r(j)/q(j)$.}
\begin{proof}\pf\ Divide the two equalities in \stepref{rsimbalance}
by their nonzero entries and the positive lengths of \stepref{rsimpositive}.~\qed\end{proof}
\step{rsimconstant}{There is $a>0$ with $r(i)=a q(i)$ for all $i$.}
\begin{proof}\pf\ Apply
Lemma~\ref{lem:rem-edge-constant} to $i\mapsto r(i)/q(i)$ on the common
Cartan support, using \stepref{rsimedge} and connectedness in \asm[2]{0};
\stepref{rsimpositive} gives a positive value at any vertex. ~\qed\end{proof}
\step{rsimgram}{$(\beta_i,\beta_j)=a(\alpha_i,\alpha_j)$ for every $i,j$.}
\begin{proof}\pf\ Substitute \stepref{rsimconstant} in
$(\beta_i,\beta_j)=M(i,j)r(j)/2$ and \eqref{eq:pairing}.~\qed\end{proof}
\step{rsimmap}{The basis correspondence extends to a linear equivalence $f$.}
\begin{proof}\pf\ Apply \path{Module.Basis.equiv} to the two real
bases and the identity index equivalence. ~\qed\end{proof}
\step{rsimall}{$(f\Zx,f\Zy)=a(\Zx,\Zy)$ for all vectors.}
\begin{proof}\pf\ Expand both vectors in their finite bases and apply
\stepref{rsimgram} term by term. ~\qed\end{proof}
\qedstep
\begin{proof}\pf\ Combine \stepref{rsimconstant}, \stepref{rsimmap},
\stepref{rsimall}; root-set equality supplies exactly the last clause
of the similarity definition. ~\qed\end{proof}
\end{proof}

\begin{namedresult}{Lemma}{(orthogonal components)}\label{lem:rem-components}
\lean{IsRootSet.orthogonal_components}\lean{IsRootSet.component_isRootSet}
\lean{IsRootSet.isIndecomposable_iff_connected}\lean{IsRootSet.subsingleton_of_eq_empty}
\lean{IsRootSet.eq_empty_iff_subsingleton}
\lean{IsRootSet.isEmpty_connectedComponent_of_eq_empty}
\uses{def:roots,lem:rem-reflection,facts:inner}
\assume{$\Phi$ is a Root System in $V$. Join distinct roots when their
inner product is nonzero.}
\prove{Each connected component is an indecomposable Root System in its
own span. Different component spans are orthogonal, and $V$ is their
orthogonal direct sum. In particular, a nonempty $\Phi$ is indecomposable
if and only if this graph is connected. The empty Root System has $V=0$
and no components.}
\end{namedresult}
\begin{proof}
\pfsketch\ Reflections preserve the component of their reflecting root;
orthogonality separates different components and their spans.
\step{roorth}{Roots in distinct components are orthogonal.}
\begin{proof}\pf\ A nonzero product would give an edge between the
components by the graph definition. ~\qed\end{proof}
\step{roreflect}{A component $A$ is closed under reflections in its roots.}
\begin{proof}
\pflet{Fix $\alpha,\beta\in A$.}
\step{rozero}{\assume{$(\alpha,\beta)=0$.}
\prove{$s_\alpha\beta=\beta\in A$.}}
\begin{proof}\pf\ Substitute zero in the reflection formula. ~\qed\end{proof}
\step{rononzero}{\assume{$(\alpha,\beta)\ne0$.}
\prove{$s_\alpha\beta\in A$.}}
\begin{proof}
\step{roproduct}{$(\alpha,s_\alpha\beta)=-(\alpha,\beta)\ne0$.}
\begin{proof}\pf\ Expand the reflection formula by bilinearity. ~\qed\end{proof}
\qedstep
\begin{proof}\pf\ The reflected root equals $\alpha$ or is joined to it
by \stepref{roproduct}, so it lies in the same component. ~\qed\end{proof}
\end{proof}
\qedstep
\begin{proof}\pf\ The cases \stepref{rozero}, \stepref{rononzero} exhaust
the inner-product values. ~\qed\end{proof}
\end{proof}
\step{rorootset}{Each component is a Root System in its span.}
\begin{proof}\pf\ Finiteness, nonzero roots, integrality and reducedness
are inherited from $\Phi$; \stepref{roreflect} gives closure and its
own span gives the spanning field of Definition~\ref{def:roots}.~\qed\end{proof}
\step{roindecomp}{Each component is indecomposable.}
\begin{proof}\pf\ An orthogonal partition could have no graph edge
crossing its parts by definition, contrary to connectedness of the component. ~\qed\end{proof}
\step{rospans}{Distinct component spans are orthogonal.}
\begin{proof}\pf\ Expand finite linear combinations using
\stepref{roorth} and bilinearity. ~\qed\end{proof}
\step{rodirect}{The sum of those spans is direct and equals $V$.}
\begin{proof}
\step{roallspan}{The sum equals $V$.}
\begin{proof}\pf\ Every root lies in a component, and $\Phi$ spans $V$.~\qed\end{proof}
\step{roindependent}{\assume{$\sum_A \Zx_A=0$ with each $\Zx_A$ in its
component span.}
\prove{Every $\Zx_A=0$.}}
\begin{proof}\pf\ Pair the relation with $\Zx_A$;
\stepref{rospans} leaves $\rsq{\Zx_A}=0$, so definiteness applies. ~\qed\end{proof}
\qedstep
\begin{proof}\pf\ Combine \stepref{roallspan}, \stepref{roindependent}.~\qed\end{proof}
\end{proof}
\step{roconnected}{For nonempty $\Phi$, indecomposability is equivalent
to connectedness of its graph.}
\begin{proof}\pf\ Multiple components give an orthogonal partition by
\stepref{roorth}; one component is indecomposable by \stepref{roindecomp}.~\qed\end{proof}
\step{roempty}{\assume{$\Phi$ is empty.}
\prove{$V=0$.}}
\begin{proof}\pf\ Use its spanning condition and the span of the empty set. ~\qed\end{proof}
\qedstep
\begin{proof}\pf\ Combine \stepref{rorootset}, \stepref{roindecomp},
\stepref{rospans}, \stepref{rodirect}, \stepref{roconnected}, \stepref{roempty}.~\qed\end{proof}
\end{proof}

\begin{namedresult}{Lemma}{(indecomposable classification)}\label{lem:rem-indecomp-classification}
\lean{IsRootSet.classification_of_isIndecomposable}\lean{RootSetModel.two_lengths_iff}
\lean{RootSetModel.squaredLengths_eq}\lean{RootSetModel.B₂_C₂_similarity}
\uses{def:roots,def:rem-bases,lem:rem-similarity-laws,lem:pairing-construction,thm:A,lem:classical,lem:E,lem:F,lem:G,lem:connected,lem:similarity,lem:counts,facts:rem-library}
\assume{$\Phi$ is a nonempty indecomposable Root System in a
finite-dimensional real inner product space $V$.}
\prove{After one common positive rescaling, $\Phi$ is isometric to exactly
one coordinate Root System in \eqref{eq:list}. Exactly $\typ B,\typ C,
\typ F,\typ G$ have two lengths. The additional $\typ C_2$ coordinate
model is similar to $\typ B_2$.}
\end{namedresult}
\begin{proof}
\pfsketch\ Apply Theorem A to a Mathlib base. Construct the basis map
and use the library equivalence of root pairings. Counts distinguish
the models without a conjugacy theorem for bases.
\step{rBpairing}{The pairing $P$ constructed from $\Phi$ is finite,
spanning, reduced, crystallographic and irreducible over $\mathbb R$.}
\begin{proof}\pf\ Apply Lemma~\ref{lem:pairing-construction} to \asm{0}.~\qed\end{proof}
\step{rBbase}{\pick\ a Mathlib base $b$ for $P$.}
\begin{proof}\pf\ Apply \path{RootPairing.nonempty_base} to
\stepref{rBpairing}; $\mathbb R$ is a characteristic zero field. ~\qed\end{proof}
\step{rBsupport}{The finite base support is nonempty.}
\begin{proof}\pf\ Apply \path{RootPairing.Base.support_nonempty} to
\stepref{rBbase}, the nonempty root set and $2\ne0$ in $\mathbb R$.~\qed\end{proof}
\step{rBfinite}{Its integer Cartan matrix is finite Cartan.}
\begin{proof}\pf\ Apply \path{RootPairing.Base.cartanMatrix_isFiniteCartan}
using \stepref{rBpairing} and classically chosen decidable equality. ~\qed\end{proof}
\step{rBindecomp}{That matrix is indecomposable.}
\begin{proof}\pf\ Apply \path{RootPairing.Base.cartanMatrix_isIndecomposable}
using \stepref{rBpairing}.~\qed\end{proof}
\step{rBtype}{\pick\ a type, rank $\ell$, and reindexing from its
model indices to the support of $b$.}
\begin{proof}\pf\ Apply Theorem~\ref{thm:A} to
\stepref{rBsupport}, \stepref{rBfinite}, \stepref{rBindecomp}.~\qed\end{proof}
\step{rBmodel}{There is a coordinate Root System $\Psi$ with a root base
having precisely this reindexed matrix.}
\begin{proof}\pf\ Select the matching case of Lemmas~\ref{lem:classical},
\ref{lem:E}, \ref{lem:F}, \ref{lem:G}, using the rank restrictions
in \stepref{rBtype}.~\qed\end{proof}
\step{rBmodelpairing}{Its pairing $P_2$ is finite, spanning, reduced and
crystallographic over $\mathbb R$.}
\begin{proof}\pf\ Apply Lemma~\ref{lem:pairing-construction} to
the Root System in \stepref{rBmodel}.~\qed\end{proof}
\step{rBmodelbase}{The displayed root base gives a Mathlib base $b_2$.}
\begin{proof}
\step{rBmonoid}{Every model root or its negative lies in the additive
monoid of the selected simple roots.}
\begin{proof}\pf\ Use the nonnegative integer expansions in
\stepref{rBmodel} and \path{Submodule.span_nat_eq_addSubmonoidClosure}.~\qed\end{proof}
\qedstep
\begin{proof}\pf\ Apply \path{RootPairing.Base.mk'} using
\stepref{rBmodelpairing}, root independence in \stepref{rBmodel},
and \stepref{rBmonoid}.~\qed\end{proof}
\end{proof}
\pflet{Let $e:b.\mathrm{support}\to b_2.\mathrm{support}$ be the index
bijection induced by \stepref{rBtype} and the model base parametrization.}
\step{rBcartan}{For all base indices $i,j$,
$b_2.\mathrm{cartanMatrix}(e(i),e(j))=b.\mathrm{cartanMatrix}(i,j)$.}
\begin{proof}
\step{rBrealcartan}{The real pairings of the two bases agree through $e$.}
\begin{proof}\pf\ The entry equality in \stepref{rBtype} and the
model pairing matrix in \stepref{rBmodel} identify both real
pairings with the same cast named matrix, using
\path{RootPairing.Base.algebraMap_cartanMatrixIn_apply}.~\qed\end{proof}
\qedstep
\begin{proof}\pf\ Rewrite both real pairings in
\stepref{rBrealcartan} as casts of integer Cartan entries by
\path{RootPairing.Base.algebraMap_cartanMatrixIn_apply};
injectivity of $\mathbb Z\to\mathbb R$ gives the equality.~\qed\end{proof}
\end{proof}
\step{rBmap}{\pick\ the linear equivalence
$f=b.\mathrm{toWeightBasis.equiv}\ b_2.\mathrm{toWeightBasis}\ e$.}
\begin{proof}\pf\ Use the spanning pairings
\stepref{rBpairing}, \stepref{rBmodelpairing} in the weight-basis
and basis-equivalence declarations of \ref{facts:rem-library}.~\qed\end{proof}
\step{rBsimple}{$f(\alpha_i)=\beta_{e(i)}$ for every simple root.}
\begin{proof}\pf\ Apply \path{Module.Basis.equiv_apply} and
\path{RootPairing.Base.toWeightBasis_apply} to \stepref{rBmap}.~\qed\end{proof}
\step{rBequivalence}{There is a root-pairing equivalence
$P\simeq P_2$ whose map on root spaces is $f$.}
\begin{proof}\pf\ Apply \path{RootPairing.Base.equivOfCartanMatrixEq}
to $b,b_2,e$ and \stepref{rBcartan}, using the source hypotheses in
\stepref{rBpairing} and finite, spanning, reduced, crystallographic
target hypotheses in \stepref{rBmodelpairing}.
Its defining root-space map is exactly \stepref{rBmap}.~\qed\end{proof}
\step{rBonto}{$f(\Phi)=\Psi$.}
\begin{proof}\pf\ The index equivalence and root compatibility of
\stepref{rBequivalence} give a bijection of the root sets.~\qed\end{proof}
\step{rBconnected}{The common base matrix has connected nonzero support.}
\begin{proof}\pf\ Apply Lemma~\ref{lem:connected} to
\stepref{rBsupport}, \stepref{rBfinite}, \stepref{rBindecomp}.~\qed\end{proof}
\step{rBsimilarity}{$f$ is a similarity using one positive factor.}
\begin{proof}
\step{rBreindexedbases}{The families
$\alpha_i=b.\mathrm{toWeightBasis}(i)$ and
$\widehat\beta_i=b_2.\mathrm{toWeightBasis}(e(i))$, indexed by
$b.\mathrm{support}$, are real vector-space bases.}
\begin{proof}\pf\ Reindex the second weight basis along $e$;
the first is the weight basis from \stepref{rBmap}.~\qed\end{proof}
\step{rBmetric}{For some $a>0$, $(f\Zx,f\Zy)=a(\Zx,\Zy)$ for all vectors.}
\begin{proof}\pf\ Apply Lemma~\ref{lem:similarity} to
\stepref{rBreindexedbases}, \stepref{rBcartan} cast to real pairings,
and \stepref{rBconnected}; \stepref{rBsimple} identifies its basis
map with $f$.~\qed\end{proof}
\qedstep
\begin{proof}\pf\ Steps \stepref{rBmetric}, \stepref{rBonto}
are the similarity conditions of Definition~\ref{def:roots}.~\qed\end{proof}
\end{proof}
\step{rBunique}{No two distinct models in \eqref{eq:list} are similar.}
\begin{proof}
\step{rBinvariants}{Similarity preserves rank, root count, number of
lengths, their ratio, and the numbers of short and long roots.}
\begin{proof}\pf\ A bijection scaling every squared length by the same
$a>0$ preserves these definitions. ~\qed\end{proof}
\step{rBonevsTwo}{One-length and two-length models are distinguished.}
\begin{proof}\pf\ Use the lengths in Lemma~\ref{lem:counts} with
\stepref{rBinvariants}.~\qed\end{proof}
\step{rBAD}{At equal rank, the $\typ A$ and $\typ D$ root counts differ
throughout the permitted $\typ D$ range.}
\begin{proof}\pf\ The equality $\ell(\ell+1)=2\ell(\ell-1)$ with
$\ell>0$ implies $\ell=3$, excluded by $\ell\geq4$.~\qed\end{proof}
\step{rBE6}{At rank 6, the counts $42,60,72$ distinguish $\typ A,\typ D,\typ E$.}
\begin{proof}\pf\ Evaluate Lemma~\ref{lem:counts} at rank 6. ~\qed\end{proof}
\step{rBE7}{At rank 7, the counts $56,84,126$ distinguish these types.}
\begin{proof}\pf\ Evaluate Lemma~\ref{lem:counts} at rank 7. ~\qed\end{proof}
\step{rBE8}{At rank 8, the counts $72,112,240$ distinguish these types.}
\begin{proof}\pf\ Evaluate Lemma~\ref{lem:counts} at rank 8. ~\qed\end{proof}
\step{rBGunique}{Squared-length ratio 3 distinguishes $\typ G_2$.}
\begin{proof}\pf\ All other two-length models have ratio 2 by
Lemma~\ref{lem:counts}.~\qed\end{proof}
\step{rBFunique}{At rank 4, $\typ F_4$ has 48 roots whereas
$\typ B_4,\typ C_4$ have 32.}
\begin{proof}\pf\ Evaluate Lemma~\ref{lem:counts}.~\qed\end{proof}
\step{rBBCunique}{For $\ell\geq3$, $\typ B_\ell,\typ C_\ell$ have
different numbers of short roots.}
\begin{proof}\pf\ Their counts $2\ell$ and $2\ell(\ell-1)$ from
Lemma~\ref{lem:counts} differ when $\ell\geq3$.~\qed\end{proof}
\qedstep
\begin{proof}\pf\ The cases \stepref{rBonevsTwo}--\stepref{rBBCunique}
separate every pair at equal rank, and \stepref{rBinvariants} separates
unequal ranks. ~\qed\end{proof}
\end{proof}
\step{rBCtwo}{The map with matrix
$\left(\begin{smallmatrix}1&1\\-1&1\end{smallmatrix}\right)$
is a similarity from $\typ B_2$ onto $\typ C_2$.}
\begin{proof}
\step{rBCaxis1}{$\Zu_1$ maps to $\Zu_1-\Zu_2$.}
\begin{proof}\pf\ Read column 1 of the displayed matrix. ~\qed\end{proof}
\step{rBCaxis2}{$\Zu_2$ maps to $\Zu_1+\Zu_2$.}
\begin{proof}\pf\ Read column 2 of the displayed matrix. ~\qed\end{proof}
\step{rBCplus}{$\Zu_1+\Zu_2$ maps to $2\Zu_1$.}
\begin{proof}\pf\ Add \stepref{rBCaxis1}, \stepref{rBCaxis2}.~\qed\end{proof}
\step{rBCminus}{$\Zu_1-\Zu_2$ maps to $-2\Zu_2$.}
\begin{proof}\pf\ Subtract \stepref{rBCaxis2} from \stepref{rBCaxis1}.~\qed\end{proof}
\step{rBCmetric}{The transpose times the matrix is $2I$.}
\begin{proof}\pf\ Compute each entry of $T^TT$ as the dot product of two columns of the displayed matrix $T$. ~\qed\end{proof}
\qedstep
\begin{proof}\pf\ The four images and their negatives exhaust
$\Phi(\typ C_2)$, while \stepref{rBCmetric} gives scale 2 and invertibility. ~\qed\end{proof}
\end{proof}
\qedstep
\begin{proof}\pf\ Use \stepref{rBsimilarity}, \stepref{rBunique},
\stepref{rBCtwo} and the length assertion of
Lemma~\ref{lem:counts}. Two similarities from $\Phi$ to models
compose, after inversion of the first, to a similarity between the
models by Lemma~\ref{lem:rem-similarity-laws}; \stepref{rBunique}
then forces the same model.~\qed\end{proof}
\end{proof}

\begin{namedresult}{Theorem}{B}\label{thm:B}
\lean{IsRootSet.classification}
\uses{def:roots,lem:rem-indecomp-classification,lem:rem-components,lem:models-irreducible,def:rem-bases}
\assume{$\Phi$ is a Root System in a finite-dimensional real inner
product space $V$.}
\prove{If $\Phi$ is nonempty and indecomposable, it is similar, by one
common positive scale, to exactly one model in \eqref{eq:list}.
Exactly the $\typ B,\typ C,\typ F,\typ G$ models have two lengths;
$\typ C_2$ is similar to $\typ B_2$ and is omitted from that list.
In general $V$ is the orthogonal direct sum of the component spans, and
each nonzero component is similar to a listed indecomposable model,
with an independent scale on each component. $\Phi$ is empty exactly
when $V$ is the zero space, and then there are no components.}
\end{namedresult}
\begin{proof}
\pfsketch\ \new\ The indecomposable assertion is a separate lemma so
that applying it to each component does not reuse local proof steps.
\step{rBmain}{The nonempty indecomposable assertion and the model length
and rank-2 assertions hold.}
\begin{proof}\pf\ Apply Lemma~\ref{lem:rem-indecomp-classification}.~\qed\end{proof}
\step{rBdecompose}{Every nonempty component is an indecomposable Root
System in its span, and these spans form an orthogonal direct sum.}
\begin{proof}\pf\ Apply Lemma~\ref{lem:rem-components} to $\Phi$.~\qed\end{proof}
\step{rBclassifycomponents}{Each component is similar to a unique listed
model using a single positive scale for that component.}
\begin{proof}\pf\ Apply Lemma~\ref{lem:rem-indecomp-classification}
separately to each Root System supplied by \stepref{rBdecompose}.~\qed\end{proof}
\step{rBmodelindecomp}{All listed models are indecomposable.}
\begin{proof}\pf\ Apply Lemma~\ref{lem:models-irreducible}.~\qed\end{proof}
\qedstep
\begin{proof}\pf\ Combine \stepref{rBmain}--\stepref{rBmodelindecomp};
the empty case is the final clause of Lemma~\ref{lem:rem-components}, and
conversely the zero space has no nonzero vector, so no root.~\qed\end{proof}
\end{proof}

\subsubsection{Angles, lines and the normalized matrix}

\paragraph{Angle convention.}\label{facts:rem-angles}
\begingroup\raggedright
For two distinct lines the angle lies in $(0,90^\circ]$ and its cosine
is $|(\alpha,\beta)|/(\|\alpha\|\|\beta\|)$ for nonzero representatives.
We use \path{Real.cos_pi_div_three}, \path{Real.cos_pi_div_four},
\path{Real.cos_pi_div_six}, and \path{Real.injOn_cos}, all in
\path{Mathlib/Analysis/SpecialFunctions/Trigonometric/Basic.lean}.
Their values and injectivity on $[0,\pi]$ identify squared cosines
$1/4,1/2,3/4$ with $60^\circ,45^\circ,30^\circ$, respectively.
\par\endgroup

\begin{namedresult}{Lemma}{(angles and fixed-length subsystems)}\label{lem:angles}
\lean{IsRootSet.pairing_product_mem}\lean{IsRootSet.angle_length_cases}
\lean{IsRootSet.fixedLength_isRootSet}
\lean{IsRootSet.fixedLength_components_simplyLaced}
\uses{def:roots,lem:rem-reflection,thm:B,lem:counts,facts:inner,facts:rem-angles}
\assume{$\Phi$ is a Root System in a finite-dimensional real inner product space.}
\prove{For every $\alpha,\beta\in\Phi$ which are neither perpendicular
nor proportional,
\[
\frac{4(\alpha,\beta)^2}{\rsq\alpha\,\rsq\beta}\in\{1,2,3\}.
\]
The longer-to-shorter length ratio and acute line angle are exactly one of
$(1,60^\circ)$, $(\sqrt2,45^\circ)$, $(\sqrt3,30^\circ)$.
For every occurring length $a>0$, the roots of length $a$ form a Root
System in their own span, and it is an orthogonal sum of systems similar
to $\typ A_\ell$ ($\ell\geq1$), $\typ D_\ell$ ($\ell\geq4$),
$\typ E_6,\typ E_7,\typ E_8$.}
\end{namedresult}
\noindent\repair\ The hypotheses exclude proportional pairs; fixed-length
subsystems need not be indecomposable or span the original space.
\begin{proof}
\pfsketch\ Bound the product of two nonzero integer pairings by 4.
Reflections preserve length, then Theorem B classifies each component.
\step{rangpair}{\assume{$\alpha,\beta\in\Phi$ are neither perpendicular
nor proportional.}
\prove{Their pairing product, length ratio and angle have the stated alternatives.}}
\begin{proof}
\pflet{$q=\rsq\alpha>0$, $r=\rsq\beta>0$, $h=(\alpha,\beta)\ne0$,
$u=2h/q$, and $v=2h/r$.}
\step{ranginteger}{$u,v$ are nonzero integers of the same sign.}
\begin{proof}\pf\ Use the two integrality conditions in
Definition~\ref{def:roots} and the positive denominators $q,r$.~\qed\end{proof}
\step{rangstrict}{$h^2<qr$.}
\begin{proof}
\step{rangresidual}{$\beta-(h/q)\alpha\ne0$.}
\begin{proof}\pf\ Otherwise $\beta$ would be proportional to $\alpha$,
contrary to \asm{1}.~\qed\end{proof}
\step{rangnorm}{$\|\beta-(h/q)\alpha\|^2=r-h^2/q>0$.}
\begin{proof}
\step{rangnormidentity}{$\|\beta-(h/q)\alpha\|^2=r-h^2/q$.}
\begin{proof}\pf\ Expand by bilinearity in \ref{facts:inner} and
substitute the definitions of $q,r,h$.~\qed\end{proof}
\step{rangnormpositive}{$\|\beta-(h/q)\alpha\|^2>0$.}
\begin{proof}\pf\ Apply definiteness in \ref{facts:inner} to
\stepref{rangresidual}.~\qed\end{proof}
\qedstep
\begin{proof}\pf\ Combine \stepref{rangnormidentity},
\stepref{rangnormpositive}.~\qed\end{proof}
\end{proof}
\qedstep
\begin{proof}\pf\ Multiply \stepref{rangnorm} by $q>0$.~\qed\end{proof}
\end{proof}
\step{rangproducts}{$uv=4h^2/(qr)\in\{1,2,3\}$.}
\begin{proof}\pf\ It is a positive integer less than 4 by
\stepref{ranginteger}, \stepref{rangstrict}.~\qed\end{proof}
\step{rangcases}{The length and angle alternatives follow.}
\begin{proof}
\step{rangone}{\assume{$uv=1$.}
\prove{$|u|=|v|=1$, $q=r$, and the angle is $60^\circ$.}}
\begin{proof}
\step{rangoneinteger}{$|u|=|v|=1$.}
\begin{proof}\pf\ A product of two positive integers is 1 only when
both are 1, applied to \stepref{ranginteger} and $uv=1$.~\qed\end{proof}
\step{rangonelength}{$q=r$.}
\begin{proof}\pf\ Substitute \stepref{rangoneinteger} in
$u/v=r/q$, using the equal signs in \stepref{ranginteger}.~\qed\end{proof}
\step{rangoneangle}{The angle is $60^\circ$.}
\begin{proof}\pf\ The squared cosine is $h^2/(qr)=uv/4=1/4$;
apply \ref{facts:rem-angles}.~\qed\end{proof}
\qedstep
\begin{proof}\pf\ Combine \stepref{rangoneinteger},
\stepref{rangonelength}, \stepref{rangoneangle}.~\qed\end{proof}
\end{proof}
\step{rangtwo}{\assume{$uv=2$.}
\prove{$\{|u|,|v|\}=\{1,2\}$, the squared-length
ratio is 2, and the angle is $45^\circ$.}}
\begin{proof}
\step{rangtwointeger}{$(|u|,|v|)$ is $(1,2)$ or $(2,1)$.}
\begin{proof}\pf\ The only positive divisors of 2 are 1 and 2,
applied to \stepref{ranginteger} and $uv=2$.~\qed\end{proof}
\step{rangtwolength}{$\max(q,r)/\min(q,r)=2$.}
\begin{proof}\pf\ Substitute the two alternatives
\stepref{rangtwointeger} in $u/v=r/q$, with $q,r>0$.~\qed\end{proof}
\step{rangtwoangle}{The angle is $45^\circ$.}
\begin{proof}\pf\ The squared cosine is $h^2/(qr)=uv/4=1/2$;
apply \ref{facts:rem-angles}.~\qed\end{proof}
\qedstep
\begin{proof}\pf\ Combine \stepref{rangtwointeger},
\stepref{rangtwolength}, \stepref{rangtwoangle}.~\qed\end{proof}
\end{proof}
\step{rangthree}{\assume{$uv=3$.}
\prove{$\{|u|,|v|\}=\{1,3\}$, the squared-length
ratio is 3, and the angle is $30^\circ$.}}
\begin{proof}
\step{rangthreeinteger}{$(|u|,|v|)$ is $(1,3)$ or $(3,1)$.}
\begin{proof}\pf\ The only positive divisors of 3 are 1 and 3,
applied to \stepref{ranginteger} and $uv=3$.~\qed\end{proof}
\step{rangthreelength}{$\max(q,r)/\min(q,r)=3$.}
\begin{proof}\pf\ Substitute the two alternatives
\stepref{rangthreeinteger} in $u/v=r/q$, with $q,r>0$.~\qed\end{proof}
\step{rangthreeangle}{The angle is $30^\circ$.}
\begin{proof}\pf\ The squared cosine is $h^2/(qr)=uv/4=3/4$;
apply \ref{facts:rem-angles}.~\qed\end{proof}
\qedstep
\begin{proof}\pf\ Combine \stepref{rangthreeinteger},
\stepref{rangthreelength}, \stepref{rangthreeangle}.~\qed\end{proof}
\end{proof}
\qedstep
\begin{proof}\pf\ Take positive square roots of the squared-length
ratios in the exhaustive cases \stepref{rangone}--\stepref{rangthree}.~\qed\end{proof}
\end{proof}
\qedstep
\begin{proof}\pf\ Combine \stepref{rangproducts}, \stepref{rangcases}.~\qed\end{proof}
\end{proof}
\step{rangfixed}{For each occurring $a>0$, the set
$\Psi=\{\gamma\in\Phi:\|\gamma\|=a\}$ is a Root System in its span.}
\begin{proof}
\step{rangfixedbasic}{$\Psi$ is finite, nonzero, reduced, integral,
and spans its own span.}
\begin{proof}\pf\ Restrict the corresponding conditions of
Definition~\ref{def:roots} to $\Psi$.~\qed\end{proof}
\step{rangfixedroot}{For $\gamma,\delta\in\Psi$, $s_\gamma\delta\in\Phi$.}
\begin{proof}\pf\ Apply reflection closure of $\Phi$.~\qed\end{proof}
\step{rangfixednorm}{$\|s_\gamma\delta\|=a$.}
\begin{proof}\pf\ Apply Lemma~\ref{lem:rem-reflection} and $\|\delta\|=a$.~\qed\end{proof}
\step{rangfixedspan}{$s_\gamma\delta$ lies in $\operatorname{span}\Psi$.}
\begin{proof}\pf\ Its defining formula is a linear combination of
$\gamma,\delta\in\Psi$.~\qed\end{proof}
\qedstep
\begin{proof}\pf\ Assemble Definition~\ref{def:roots} from
\stepref{rangfixedbasic}--\stepref{rangfixedspan}.~\qed\end{proof}
\end{proof}
\step{rangADE}{Every nonempty component of this subsystem has type
$\typ A,\typ D$ or $\typ E$.}
\begin{proof}\pf\ Apply Theorem~\ref{thm:B} to \stepref{rangfixed};
its other models have two lengths by Lemma~\ref{lem:counts}, impossible
in the fixed-length subset. ~\qed\end{proof}
\qedstep
\begin{proof}\pf\ Combine \stepref{rangpair},
\stepref{rangfixed}, \stepref{rangADE}.~\qed\end{proof}
\end{proof}

\begin{namedresult}{Lemma}{(unit-line orthogonal projections)}\label{lem:rem-projectors}
\lean{IsRootSet.lineProjector_identities}
\uses{facts:inner}
\assume{$V$ is a finite-dimensional real inner product space,
$\Zv$ is a unit vector, and $P_\Zv\Zx=(\Zv,\Zx)\Zv$ is the rank-one orthogonal projection onto $\mathbb R\Zv$.}
\prove{$P_{-\Zv}=P_\Zv$; $P_\Zv$ is symmetric; its trace is 1;
and for every orthogonal linear bijection $f$,
$fP_\Zv f^{-1}=P_{f\Zv}$.}
\end{namedresult}
\begin{proof}
\pfsketch\ Expand the rank-one operator and its diagonal.
\step{rprsign}{$P_{-\Zv}\Zx=P_\Zv\Zx$.}
\begin{proof}\pf\ The two minus signs cancel by bilinearity. ~\qed\end{proof}
\step{rprsymmetry}{$(P_\Zv\Zx,\Zy)=(\Zx,P_\Zv\Zy)$.}
\begin{proof}\pf\ Both sides are $(\Zv,\Zx)(\Zv,\Zy)$ by
\ref{facts:inner}.~\qed\end{proof}
\step{rprtrace}{In an orthonormal basis the trace is $\sum_i v(i)^2=1$.}
\begin{proof}\pf\ The $i$th diagonal entry is $(\Zv,\Zu_i)^2$;
sum and use $\|\Zv\|=1$.~\qed\end{proof}
\step{rprconjugate}{$fP_\Zv f^{-1}\Zx=(f\Zv,\Zx)f\Zv$.}
\begin{proof}\pf\ Substitute the orthogonal projection formula and the isometry
identity $(\Zv,f^{-1}\Zx)=(f\Zv,\Zx)$.~\qed\end{proof}
\qedstep
\begin{proof}\pf\ Operator extensionality applied to
\stepref{rprsign}, \stepref{rprsymmetry}, \stepref{rprtrace},
\stepref{rprconjugate} gives all four identities. ~\qed\end{proof}
\end{proof}

\begin{namedresult}{Lemma}{(the line frame)}\label{lem:tight-frame}
\lean{IsRootSet.frame_identity}
\uses{def:roots,lem:rem-reflection,lem:rem-components,lem:rem-edge-constant,lem:rem-projectors,facts:inner}
\assume{\begin{pfenum}
\item $\Phi$ is a nonempty indecomposable Root System in an
 $n$-dimensional real inner product space $V$.
\item Its line set is $\mathcal L$, $|\mathcal L|=N$, and $\Zv_L$ is
 any unit representative of each line.
\end{pfenum}}
\prove{$n>0$ and
\begin{equation}\label{eq:frame}
\mathcal F:=\sum_{L\in\mathcal L}P_{\Zv_L}=\frac Nn I_V.
\end{equation}
The identity is independent of all representative signs.}
\end{namedresult}
\begin{proof}
\pfsketch\ The symmetric frame operator commutes with every root
reflection. It acts by one scalar on all roots, and its trace determines it.
\step{rfrdimension}{$n>0$.}
\begin{proof}\pf\ The space contains a nonzero root from \asm[1]{0}.~\qed\end{proof}
\step{rfrsign}{The operator $\mathcal F$ is independent of signs and symmetric.}
\begin{proof}\pf\ Sum the sign and symmetry identities of
Lemma~\ref{lem:rem-projectors}.~\qed\end{proof}
\step{rfrcommute}{$\mathcal F$ commutes with each root reflection $s_\alpha$.}
\begin{proof}
\step{rfrpermute}{$s_\alpha$ permutes $\mathcal L$.}
\begin{proof}\pf\ It is an orthogonal root permutation by
Lemma~\ref{lem:rem-reflection}.~\qed\end{proof}
\step{rfrconjugation}{$s_\alpha\mathcal F s_\alpha^{-1}=\mathcal F$.}
\begin{proof}\pf\ Reindex the sum by \stepref{rfrpermute}, using
projector conjugation and sign independence in Lemma~\ref{lem:rem-projectors}.~\qed\end{proof}
\qedstep
\begin{proof}\pf\ Multiply \stepref{rfrconjugation} on the right by
$s_\alpha$.~\qed\end{proof}
\end{proof}
\step{rfreigen}{For every root $\alpha$ there is a unique scalar
$t_\alpha$ with $\mathcal F\alpha=t_\alpha\alpha$.}
\begin{proof}
\step{rfrminus}{$s_\alpha(\mathcal F\alpha)=-\mathcal F\alpha$.}
\begin{proof}\pf\ Use \stepref{rfrcommute} and
$s_\alpha\alpha=-\alpha$ from Lemma~\ref{lem:rem-reflection}.~\qed\end{proof}
\qedstep
\begin{proof}\pf\ The $-1$ eigenspace in
Lemma~\ref{lem:rem-reflection} is $\mathbb R\alpha$; uniqueness uses
$\alpha\ne0$.~\qed\end{proof}
\end{proof}
\step{rfredge}{Nonorthogonal roots have the same scalar.}
\begin{proof}\pf\ Symmetry in \stepref{rfrsign} gives
$(t_\alpha-t_\beta)(\alpha,\beta)=0$ by \stepref{rfreigen};
cancel the nonzero inner product. ~\qed\end{proof}
\step{rfrconstant}{There is a single real $t$ with $t_\alpha=t$
for all roots.}
\begin{proof}\pf\ Apply Lemma~\ref{lem:rem-edge-constant} to
$\alpha\mapsto t_\alpha$ on the nonorthogonality graph of all roots,
connected by Lemma~\ref{lem:rem-components}, using \stepref{rfredge}.~\qed\end{proof}
\step{rfrscalar}{$\mathcal F=tI_V$.}
\begin{proof}\pf\ Expand a vector in the spanning roots and use
\stepref{rfreigen}, \stepref{rfrconstant}.~\qed\end{proof}
\step{rfrtrace}{$\operatorname{tr}\mathcal F=N$.}
\begin{proof}\pf\ Sum the diagonal entries and the unit traces of
Lemma~\ref{lem:rem-projectors}.~\qed\end{proof}
\step{rfrvalue}{$t=N/n$.}
\begin{proof}\pf\ The trace of \stepref{rfrscalar} is $nt$;
use \stepref{rfrtrace} and divide by $n>0$ from \stepref{rfrdimension}.~\qed\end{proof}
\qedstep
\begin{proof}\pf\ Substitute \stepref{rfrvalue} in \stepref{rfrscalar}
and use \stepref{rfrdimension}, \stepref{rfrsign}.~\qed\end{proof}
\end{proof}

\begin{namedresult}{Lemma}{(the normalized matrix convention)}\label{lem:C-convention}
\lean{IsRootSet.signedAdj_quadratic_iff_frame}\lean{IsRootSet.lineFrame_signs}
\lean{IsRootSet.signedAdj_signs}
\uses{lem:rem-projectors,facts:inner}
\assume{\begin{pfenum}
\item $N$ lines span an $n$-dimensional real inner product space, $n>0$.
\item In an orthonormal coordinate basis, $U$ has a chosen unit vector on
 each line as its columns, $G=U^TU$, and $\theta\in\mathbb R$.
\end{pfenum}}
\pflet{\begin{equation}\label{eq:C-convention}C=2G-2I_N.\end{equation}}
\prove{$(C+2I_N)(C-\theta I_N)=0$ if and only if
$UU^T=(N/n)I_n$ and $\theta=2N/n-2$.
Both conditions are independent of the choices of signs.}
\end{namedresult}
\noindent\repair\ This extends the squared-norm-2 representative
convention preceding Theorem 3.7. Unequal-length roots are normalized
line by line for this matrix; that operation is not asserted to preserve
the Root System. Entries may be $\pm\sqrt2,\pm\sqrt3$.
\begin{proof}
\pfsketch\ Reduce the polynomial to $G^2=tG$, then use surjectivity of
$U$ to cancel on both sides and use the trace.
\pflet{$t=(\theta+2)/2$.}
\step{rCexpand}{$(C+2I)(C-\theta I)=4(G^2-tG)$.}
\begin{proof}\pf\ Expand using \eqref{eq:C-convention} and the definition
of $t$.~\qed\end{proof}
\step{rCfactor}{$G^2-tG=U^T(UU^T-tI)U$.}
\begin{proof}\pf\ Substitute $G=U^TU$ and use matrix associativity. ~\qed\end{proof}
\step{rCsurjective}{$U$ is surjective onto the ambient coordinate space.}
\begin{proof}\pf\ Its columns span by \asm[1]{0}.~\qed\end{proof}
\step{rCforward}{\assume{$G^2-tG=0$.}
\prove{$UU^T=tI$.}}
\begin{proof}
\step{rCtest}{For any $\Zx,\Zy$ there exist $\Zv,\Zw$ with
$U\Zv=\Zx$, $U\Zw=\Zy$.}
\begin{proof}\pf\ Apply \stepref{rCsurjective} to each of $\Zx,\Zy$.~\qed\end{proof}
\step{rCbilinear}{$\Zx^T(UU^T-tI)\Zy=0$ for every $\Zx,\Zy$.}
\begin{proof}\pf\ Substitute \stepref{rCtest} into \stepref{rCfactor}
under $G^2-tG=0$.~\qed\end{proof}
\qedstep
\begin{proof}\pf\ Evaluate \stepref{rCbilinear} on every pair of
coordinate unit vectors to obtain each zero entry. ~\qed\end{proof}
\end{proof}
\step{rCreverse}{\assume{$UU^T=tI$.}
\prove{$G^2-tG=0$.}}
\begin{proof}\pf\ Substitute in \stepref{rCfactor}.~\qed\end{proof}
\step{rCtrace}{\assume{$UU^T=tI$.}
\prove{$t=N/n$.}}
\begin{proof}
\step{rCunittrace}{$\operatorname{tr}(UU^T)=N$.}
\begin{proof}\pf\ Sum the squared coordinates of each unit column in
\asm[2]{0}, using Lemma~\ref{lem:rem-projectors}.~\qed\end{proof}
\qedstep
\begin{proof}\pf\ The assumed equality gives $nt=N$ by
\stepref{rCunittrace}; cancel $n>0$.~\qed\end{proof}
\end{proof}
\step{rCsigns}{Changing column signs leaves $UU^T$ unchanged and
conjugates both $G$ and $C$ by a diagonal sign matrix.}
\begin{proof}
\pflet{$D_\epsilon=\operatorname{diag}(\epsilon)$ with $\epsilon(i)^2=1$.}
\step{rCsignsquare}{$D_\epsilon^T=D_\epsilon$ and $D_\epsilon^2=I$.}
\begin{proof}\pf\ Multiply diagonal entries. ~\qed\end{proof}
\step{rCsignframe}{$(UD_\epsilon)(UD_\epsilon)^T=UU^T$.}
\begin{proof}\pf\ Apply associativity and \stepref{rCsignsquare}.~\qed\end{proof}
\step{rCsignGram}{$(UD_\epsilon)^T(UD_\epsilon)=D_\epsilon G D_\epsilon$.}
\begin{proof}\pf\ Use transpose of a product and $G=U^TU$.~\qed\end{proof}
\qedstep
\begin{proof}\pf\ Substitute \stepref{rCsignGram} in
\eqref{eq:C-convention}, and use \stepref{rCsignsquare},
\stepref{rCsignframe}.~\qed\end{proof}
\end{proof}
\qedstep
\begin{proof}\pf\ Combine \stepref{rCexpand}, \stepref{rCforward},
\stepref{rCreverse}, \stepref{rCtrace}, $t=(\theta+2)/2$, and
the conjugation invariance in \stepref{rCsigns}.~\qed\end{proof}
\end{proof}

\begin{namedresult}{Lemma}{(three lines at forty-five degrees)}\label{lem:rem-three-lines}
\lean{IsRootSet.not_contains_three_equiangular_lines}\lean{IsRootSet.fortyFiveTriple_properties}
\uses{def:roots,lem:angles,facts:inner,facts:rem-angles}
\assume{$\Zu_1,\Zu_2,\Zu_3$ are orthonormal in $\mathbb R^3$ and $c=1/\sqrt2$.}
\pflet{$\Zv_1=\Zu_1$, $\Zv_2=c\Zu_1+c\Zu_2$, and
$\Zv_3=c\Zu_1+(1-c)\Zu_2+\sqrt{2c-1}\Zu_3$.}
\prove{These are unit vectors on three distinct lines pairwise at
$45^\circ$. There are no Root System $\Phi$ in any Euclidean space
and $\alpha,\beta,\gamma\in\Phi$ satisfying all three equations
\[
\begin{aligned}
4(\alpha,\beta)^2&=2\rsq\alpha\,\rsq\beta,\\
4(\alpha,\gamma)^2&=2\rsq\alpha\,\rsq\gamma,\\
4(\beta,\gamma)^2&=2\rsq\beta\,\rsq\gamma.
\end{aligned}
\]}
\end{namedresult}
\begin{proof}
\pfsketch\ Check their Gram matrix, then choose a shortest root on the
three lines in a putative containing Root System.
\step{rtlradicand}{$c^2=1/2$, $1/2<c<1$, and $2c-1>0$.}
\begin{proof}\pf\ Use $1<\sqrt2<2$ and $(\sqrt2)^2=2$ in the definition of $c$.~\qed\end{proof}
\step{rtlnorm1}{$\rsq{\Zv_1}=1$.}
\begin{proof}\pf\ Use orthonormality. ~\qed\end{proof}
\step{rtlnorm2}{$\rsq{\Zv_2}=2c^2=1$.}
\begin{proof}\pf\ Sum its coordinate squares using \stepref{rtlradicand}.~\qed\end{proof}
\step{rtlnorm3}{$\rsq{\Zv_3}=c^2+(1-c)^2+(2c-1)=1$.}
\begin{proof}\pf\ Sum its coordinate squares and use \stepref{rtlradicand}.~\qed\end{proof}
\step{rtlip12}{$(\Zv_1,\Zv_2)=c$.}
\begin{proof}\pf\ Read the first coordinate of $\Zv_2$.~\qed\end{proof}
\step{rtlip13}{$(\Zv_1,\Zv_3)=c$.}
\begin{proof}\pf\ Read the first coordinate of $\Zv_3$.~\qed\end{proof}
\step{rtlip23}{$(\Zv_2,\Zv_3)=c^2+c(1-c)=c$.}
\begin{proof}\pf\ Sum their coordinate products. ~\qed\end{proof}
\step{rtlangle}{The three lines are distinct and pairwise at $45^\circ$.}
\begin{proof}\pf\ Use the unit norms \stepref{rtlnorm1}--\stepref{rtlnorm3},
the inner products \stepref{rtlip12}--\stepref{rtlip23}, and
$0<c<1$ in \ref{facts:rem-angles}.~\qed\end{proof}
\step{rtlimpossible}{\assume{A Root System $\Phi$ and roots
$\alpha,\beta,\gamma\in\Phi$ satisfy the three equations in the statement.}
\prove{\textsc{False}.}}
\begin{proof}
\step{rtlpairs}{The three root pairs are nonperpendicular, nonproportional,
and have line angle $45^\circ$.}
\begin{proof}\pf\ Their squared cosines are $1/2$ by \asm{1}
and the positive squared norms of nonzero roots in
Definition~\ref{def:roots}; use \ref{facts:rem-angles}.~\qed\end{proof}
\step{rtlminimum}{\pick\ a permutation of the three roots with $\alpha$
of minimum squared length $b=\rsq\alpha>0$.}
\begin{proof}\pf\ A finite nonempty set of three squared lengths
has a minimum; the three equations are unchanged by the permutation.~\qed\end{proof}
\step{rtlbeta}{$\rsq\beta=2b$.}
\begin{proof}\pf\ Apply Lemma~\ref{lem:angles} to the $45^\circ$ pair
$\alpha,\beta$ from \stepref{rtlpairs} and minimality in
\stepref{rtlminimum}.~\qed\end{proof}
\step{rtlgamma}{$\rsq\gamma=2b$.}
\begin{proof}\pf\ Apply Lemma~\ref{lem:angles} to the $45^\circ$ pair
$\alpha,\gamma$ from \stepref{rtlpairs} and minimality in
\stepref{rtlminimum}.~\qed\end{proof}
\qedstep
\begin{proof}\pf\ The equal lengths in \stepref{rtlbeta},
\stepref{rtlgamma} contradict Lemma~\ref{lem:angles} for the
$45^\circ$ pair $\beta,\gamma$.~\qed\end{proof}
\end{proof}
\qedstep
\begin{proof}\pf\ Combine \stepref{rtlnorm1}--\stepref{rtlnorm3},
\stepref{rtlangle}, \stepref{rtlimpossible}.~\qed\end{proof}
\end{proof}

\begin{namedresult}{Theorem}{C}\label{thm:C}
\lean{IsRootSet.remark_lines}\lean{RootSetModel.Bn_Cn_same_lines}
\lean{RootSetModel.signedAdj_quadratic_Bn}\lean{RootSetModel.signedAdj_quadratic_Cn}
\lean{RootSetModel.signedAdj_quadratic_F₄}\lean{RootSetModel.signedAdj_quadratic_G₂}
\lean{IsRootSet.exists_rootRepresentatives}
\uses{lem:angles,lem:counts,lem:models-irreducible,lem:tight-frame,lem:C-convention,lem:rem-three-lines,def:rem-classical,def:rem-exceptional,facts:rem-angles,def:roots}
\assume{\begin{pfenum}
\item Root Systems are as in Definition~\ref{def:roots}; every line
 matrix is $C=2U^TU-2I$ with unit columns.
\item For the classical line assertions $n\geq2$, and the models are
 those of Definitions~\ref{def:rem-classical}, \ref{def:rem-exceptional}.
\end{pfenum}}
\prove{Nonperpendicular, nonproportional root pairs have exactly the
length-ratio/angle alternatives in Lemma~\ref{lem:angles}. Each occurring
fixed-length subset is a Root System in its own span and an orthogonal
sum of systems similar to $\typ A_\ell$ ($\ell\geq1$), $\typ D_\ell$
($\ell\geq4$), $\typ E_6,\typ E_7,\typ E_8$.
The $\typ B_n,\typ C_n$ models have the same $n^2$ lines in dimension $n$,
with allowed angles $90^\circ,60^\circ,45^\circ$, and
$(C+2I)(C-2(n-1)I)=0$.
Type $\typ F_4$ has 24 lines in dimension 4, with those three angles,
and $(C+2I)(C-10I)=0$.
Type $\typ G_2$ has 6 lines in dimension 2, with angles
$90^\circ,60^\circ,30^\circ$, and $(C+2I)(C-4I)=0$.
For $n=2$ the $60^\circ$ classical angle does not occur; for $n\geq3$
all three classical angles occur. All listed angles occur in $\typ F_4$
and $\typ G_2$. Three lines pairwise at $45^\circ$ exist in $\mathbb R^3$
and are contained in no Root System in any ambient dimension.}
In Lean the matrix equations hold for every choice of one root on each line, with a unit vector
on the same line. Such a choice of roots exists, indexed by $\{1,\dots,N\}$ for the number $N$ of
lines (\texttt{IsRootSet.exists\_rootRepresentatives}), so the equations are not vacuous.
\end{namedresult}
\begin{proof}
\pfsketch\ Use the common line frame for the matrix equations, and give
witness pairs for each claimed occurring angle.
\step{rCgeneral}{The root-pair and fixed-length assertions hold.}
\begin{proof}\pf\ Apply Lemma~\ref{lem:angles}.~\qed\end{proof}
\step{rCclassicalsame}{The $\typ B_n$ and $\typ C_n$ line sets coincide.}
\begin{proof}
\step{rCsameaxes}{Their axis roots lie on exactly the same axes.}
\begin{proof}\pf\ The vectors $\Zu_i$ and $2\Zu_i$ are nonzero
multiples in Definition~\ref{def:rem-classical}.~\qed\end{proof}
\step{rCsamepairs}{Their two-coordinate roots are identical.}
\begin{proof}\pf\ Compare their second clauses in
Definition~\ref{def:rem-classical}.~\qed\end{proof}
\qedstep
\begin{proof}\pf\ The two patterns in \stepref{rCsameaxes},
\stepref{rCsamepairs} exhaust both root sets. ~\qed\end{proof}
\end{proof}
\step{rCcounts}{The line counts and dimensions are
$(n^2,n)$, $(24,4)$, $(6,2)$ in the classical, $\typ F$, and $\typ G$ cases.}
\begin{proof}\pf\ Apply Lemma~\ref{lem:counts} and the coordinate-space
definitions. ~\qed\end{proof}
\step{rCallowed}{The allowed angles are the asserted ones.}
\begin{proof}
\step{rCallowedBF}{In $\typ B,\typ C,\typ F$ a nonorthogonal distinct-line
pair has angle $60^\circ$ or $45^\circ$.}
\begin{proof}\pf\ Lemma~\ref{lem:counts} gives equal lengths or ratio
$\sqrt2$; apply Lemma~\ref{lem:angles}.~\qed\end{proof}
\step{rCallowedG}{In $\typ G$ such a pair has angle $60^\circ$ or $30^\circ$.}
\begin{proof}\pf\ Lemma~\ref{lem:counts} gives equal lengths or ratio
$\sqrt3$; apply Lemma~\ref{lem:angles}.~\qed\end{proof}
\qedstep
\begin{proof}\pf\ Adjoin the perpendicular case $90^\circ$ to
\stepref{rCallowedBF}, \stepref{rCallowedG}.~\qed\end{proof}
\end{proof}
\step{rCoccurrence}{All claimed angles occur, with the stated rank-2 exception.}
\begin{proof}
\step{rCright}{The axes $\mathbb R\Zu_1,\mathbb R\Zu_2$ give $90^\circ$
in $\typ B_n,\typ C_n,\typ F_4$.}
\begin{proof}\pf\ Their product is zero, and both lines occur in the
coordinate lists. ~\qed\end{proof}
\step{rCfortyfive}{The lines of $\Zu_1,\Zu_1+\Zu_2$ give $45^\circ$
in those models.}
\begin{proof}\pf\ Their squared cosine is $1/(1\cdot2)=1/2$;
use \ref{facts:rem-angles}.~\qed\end{proof}
\step{rCsixty}{For $n\geq3$, the lines of $\Zu_1+\Zu_2,\Zu_1+\Zu_3$
give $60^\circ$ in $\typ B_n,\typ C_n$ and also in $\typ F_4$.}
\begin{proof}\pf\ Their squared cosine is $1/(2\cdot2)=1/4$.~\qed\end{proof}
\step{rCtwoangles}{For $n=2$ the classical models have no $60^\circ$ pair.}
\begin{proof}
\step{rCtwoaxes}{The two axes are perpendicular.}
\begin{proof}\pf\ Their inner product is zero. ~\qed\end{proof}
\step{rCtwodiagonals}{The two diagonal lines are perpendicular.}
\begin{proof}\pf\ $(\Zu_1+\Zu_2,\Zu_1-\Zu_2)=1-1=0$.~\qed\end{proof}
\step{rCtwomixed}{Every axis-diagonal pair has squared cosine $1/2$.}
\begin{proof}\pf\ Its absolute inner product is 1 and its squared norms
are 1 and 2. ~\qed\end{proof}
\qedstep
\begin{proof}\pf\ These three pair types exhaust the four lines
in Definition~\ref{def:rem-classical}.~\qed\end{proof}
\end{proof}
\step{rCGsixty}{The short roots $\Zu_1-\Zu_2,\Zu_1-\Zu_3$ give
$60^\circ$ in $\typ G_2$.}
\begin{proof}\pf\ Their squared cosine is $1/(2\cdot2)=1/4$.~\qed\end{proof}
\step{rCGright}{The roots $2\Zu_1-\Zu_2-\Zu_3,\Zu_2-\Zu_3$
give $90^\circ$.}
\begin{proof}\pf\ Their inner product is $-1+1=0$.~\qed\end{proof}
\step{rCGthirty}{The roots $2\Zu_1-\Zu_2-\Zu_3,\Zu_1-\Zu_2$
give $30^\circ$.}
\begin{proof}\pf\ Their squared cosine is $3^2/(6\cdot2)=3/4$.~\qed\end{proof}
\qedstep
\begin{proof}\pf\ Combine \stepref{rCright}--\stepref{rCGthirty}
using \ref{facts:rem-angles}.~\qed\end{proof}
\end{proof}
\step{rCequations}{All three asserted matrix equations hold.}
\begin{proof}
\step{rCframes}{The frame identity holds for each model.}
\begin{proof}\pf\ Apply Lemma~\ref{lem:tight-frame} to each
indecomposable Root System of Lemma~\ref{lem:models-irreducible}.~\qed\end{proof}
\step{rCeqBC}{For $\typ B_n$ and $\typ C_n$, $\theta=2(n-1)$.}
\begin{proof}\pf\ Substitute $N=n^2$ from \stepref{rCcounts}
and the frame \stepref{rCframes} in Lemma~\ref{lem:C-convention}.~\qed\end{proof}
\step{rCeqF}{For $\typ F_4$, $\theta=10$.}
\begin{proof}\pf\ Substitute $N=24$, dimension 4, and
\stepref{rCframes} in Lemma~\ref{lem:C-convention}.~\qed\end{proof}
\step{rCeqG}{For $\typ G_2$, $\theta=4$.}
\begin{proof}\pf\ Substitute $N=6$, dimension 2, and
\stepref{rCframes} in Lemma~\ref{lem:C-convention}.~\qed\end{proof}
\qedstep
\begin{proof}\pf\ Combine \stepref{rCeqBC}, \stepref{rCeqF}, \stepref{rCeqG}.~\qed\end{proof}
\end{proof}
\step{rCexample}{The asserted three-line example exists and cannot lie
in any Root System.}
\begin{proof}\pf\ Apply Lemma~\ref{lem:rem-three-lines}.~\qed\end{proof}
\qedstep
\begin{proof}\pf\ Combine \stepref{rCgeneral}, \stepref{rCclassicalsame},
\stepref{rCcounts}, \stepref{rCallowed}, \stepref{rCoccurrence},
\stepref{rCequations}, \stepref{rCexample}.~\qed\end{proof}
\end{proof}

\begin{namedresult}{Corollary}{(the equal-length matrix constants)}\label{cor:rem-constants}
\lean{RootSetModel.signedAdj_quadratic_simplyLaced}
\uses{lem:counts,lem:models-irreducible,lem:tight-frame,lem:C-convention}
\assume{The model is $\typ A_n$ ($n\geq1$), $\typ D_n$ ($n\geq4$),
or $\typ E_6,\typ E_7,\typ E_8$, in its span; $C=2U^TU-2I$
uses unit representatives of its lines.}
\prove{$(C+2I)(C-\theta I)=0$, with respective constants
$n-1,2(n-2),10,16,28$.}
\end{namedresult}
\begin{proof}
\pfsketch\ \repair\ The last constant is 28, correcting the printed
18 in Theorem 3.7.
\step{reqframe}{Each model satisfies the line frame identity.}
\begin{proof}\pf\ Apply Lemma~\ref{lem:tight-frame} using
Lemma~\ref{lem:models-irreducible}.~\qed\end{proof}
\step{reqA}{For $\typ A_n$, the second constant is $n-1$.}
\begin{proof}\pf\ Use $N=n(n+1)/2$ from Lemma~\ref{lem:counts}
in Lemma~\ref{lem:C-convention} with \stepref{reqframe}.~\qed\end{proof}
\step{reqD}{For $\typ D_n$, it is $2(n-2)$.}
\begin{proof}\pf\ Substitute $N=n(n-1)$ from Lemma~\ref{lem:counts}
in Lemma~\ref{lem:C-convention} with \stepref{reqframe}.~\qed\end{proof}
\step{reqE6}{For $\typ E_6$, it is $2(36/6)-2=10$.}
\begin{proof}\pf\ Substitute the $\typ E_6$ count in
Lemma~\ref{lem:C-convention} with \stepref{reqframe}.~\qed\end{proof}
\step{reqE7}{For $\typ E_7$, it is $2(63/7)-2=16$.}
\begin{proof}\pf\ Substitute the $\typ E_7$ count in
Lemma~\ref{lem:C-convention} with \stepref{reqframe}.~\qed\end{proof}
\step{reqE8}{For $\typ E_8$, it is $2(120/8)-2=28$.}
\begin{proof}\pf\ Substitute the $\typ E_8$ count in
Lemma~\ref{lem:C-convention} with \stepref{reqframe}.~\qed\end{proof}
\qedstep
\begin{proof}\pf\ The cases \stepref{reqA}--\stepref{reqE8}
exhaust the listed models. ~\qed\end{proof}
\end{proof}

\subsubsection{Finite certificates from the ordinary proof}

\begin{namedresult}{Lemma}{(leaf determinant expansion)}\label{lem:rem-leaf-det}
\lean{Matrix.det_eq_of_leaf}
\uses{facts:rem-library}
\assume{$M$ is a square matrix over a commutative ring, on a finite
index set; $v\ne w$, $M(v,v)=2$, and every off-diagonal entry in
row or column $v$ is zero except possibly the entry at $w$.}
\prove{Writing $p=M(v,w)M(w,v)$ and subscripts for principal deletions,
\begin{equation}\label{eq:leaf-det}
\det M=2\det M_{-v}-p\det M_{-v,-w}.
\end{equation}}
\end{namedresult}
\begin{proof}
\pfsketch\ Separate permutations fixing $v$ from those interchanging
$v,w$. This applies to a leaf of any support, not only a path.
\step{rldfixed}{Permutation terms fixing $v$ sum to $2\det M_{-v}$.}
\begin{proof}\pf\ Restrict those permutations to the complement of $v$
in \path{Matrix.det_apply'}; its diagonal factor is 2. ~\qed\end{proof}
\step{rldmoving}{A nonzero permutation term not fixing $v$ interchanges
$v$ and $w$.}
\begin{proof}
\step{rldcolumn}{Column $v$ must choose row $w$.}
\begin{proof}\pf\ All other available off-diagonal entries of column
$v$ vanish by \asm{0}.~\qed\end{proof}
\step{rldrow}{Row $v$ must choose column $w$.}
\begin{proof}\pf\ Column $v$ is already used by row $w$ in
\stepref{rldcolumn}, and all remaining entries of row $v$ except at $w$ vanish. ~\qed\end{proof}
\qedstep
\begin{proof}\pf\ Combine \stepref{rldcolumn}, \stepref{rldrow}.~\qed\end{proof}
\end{proof}
\step{rldsum}{The moving terms sum to $-p\det M_{-v,-w}$.}
\begin{proof}\pf\ Delete the transposition in \stepref{rldmoving}
from each permutation; it contributes sign $-1$ and entry product $p$
in \path{Matrix.det_apply'}.~\qed\end{proof}
\qedstep
\begin{proof}\pf\ Add the exhaustive permutation classes
\stepref{rldfixed}, \stepref{rldsum}.~\qed\end{proof}
\end{proof}

\begin{namedresult}{Lemma}{(leading determinants)}\label{comp:rem-minors}
\lean{CartanMatrix.model_leading_determinants}
\uses{def:rem-matrices,rem:small,lem:rem-leaf-det,def:rem-bases,facts:rem-library}
\assume{$M$ is a named matrix in \eqref{eq:list};
$\Delta_k$ is its leading $k$-row principal determinant and $\Delta_0=1$.}
\prove{The complete lists are
\begin{equation}\label{eq:minors}
\begin{array}{c|l}
\typ A_\ell&\Delta_k=k+1\\
\typ B_\ell,\typ C_\ell&\Delta_k=k+1\ (k<\ell),\quad\Delta_\ell=2\\
\typ D_\ell&\Delta_k=k+1\ (k<\ell),\quad\Delta_\ell=4\\
\typ E_6&(2,4,6,5,4,3)\\
\typ E_7&(2,4,6,5,4,3,2)\\
\typ E_8&(2,4,6,5,4,3,2,1)\\
\typ F_4&(2,3,2,1)\\
\typ G_2&(2,1).
\end{array}
\end{equation}}
\end{namedresult}
\begin{proof}
\pfsketch\ Apply the leaf formula with the actual deleted neighbour,
then solve the path recurrence.
\step{rmdpath}{For a path extension with last edge product $p$,
$\Delta_k=2\Delta_{k-1}-p\Delta_{k-2}$.}
\begin{proof}\pf\ Apply Lemma~\ref{lem:rem-leaf-det} to its last vertex. ~\qed\end{proof}
\step{rmdA}{The unit-path determinants are $k+1$.}
\begin{proof}
\step{rmdAbase}{$\Delta_0=1$ and $\Delta_1=2$.}
\begin{proof}\pf\ Evaluate the empty and one-entry determinants. ~\qed\end{proof}
\step{rmdAstep}{\assume{$\Delta_{k-1}=k$ and $\Delta_{k-2}=k-1$.}
\prove{$\Delta_k=k+1$.}}
\begin{proof}\pf\ Substitute in \stepref{rmdpath} with $p=1$.~\qed\end{proof}
\qedstep
\begin{proof}\pf\ Induction using \stepref{rmdAbase}, \stepref{rmdAstep}.~\qed\end{proof}
\end{proof}
\step{rmdBC}{The $\typ B,\typ C$ final determinants are 2.}
\begin{proof}\pf\ Their last edge product is 2, so
\stepref{rmdpath}, \stepref{rmdA} give $2\ell-2(\ell-1)=2$.~\qed\end{proof}
\step{rmdD}{The $\typ D_\ell$ final determinant is 4.}
\begin{proof}
\step{rmdDleaf}{Deleting $v=\ell-1$ leaves a unit path of determinant $\ell$.}
\begin{proof}\pf\ Use Definition~\ref{def:rem-matrices} and \stepref{rmdA}.~\qed\end{proof}
\step{rmdDneighbour}{Deleting also $w=\ell-3$ leaves a path of
$\ell-3$ vertices and an isolated vertex, of determinant $2(\ell-2)$.}
\begin{proof}\pf\ Factor the isolated diagonal entry 2 in the
determinant expansion and apply \stepref{rmdA}.~\qed\end{proof}
\qedstep
\begin{proof}\pf\ Lemma~\ref{lem:rem-leaf-det} with $p=1$,
\stepref{rmdDleaf}, \stepref{rmdDneighbour} gives $2\ell-2(\ell-2)=4$.~\qed\end{proof}
\end{proof}
\step{rmdE}{The $\typ E$ prefix determinants are $2,4,6,5,4,3,2,1$.}
\begin{proof}
\step{rmdE1}{The first determinant is 2.}
\begin{proof}\pf\ Evaluate the single diagonal entry. ~\qed\end{proof}
\step{rmdE2}{The second is 4.}
\begin{proof}\pf\ The leading two vertices are isolated, giving $2\cdot2$.~\qed\end{proof}
\step{rmdE3}{The third is 6.}
\begin{proof}\pf\ Its single unit edge contributes determinant 3 and
its isolated vertex contributes 2. ~\qed\end{proof}
\step{rmdE4}{The fourth is 5.}
\begin{proof}\pf\ The order in Remark~\ref{rem:small} makes it a
four-vertex unit path; apply \stepref{rmdA}.~\qed\end{proof}
\step{rmdErest}{The next four determinants are $4,3,2,1$.}
\begin{proof}\pf\ Its leading $k$-row matrix is exactly
$\typ E_k$ by \path{CartanMatrix.E}; for $k=5,6,7,8$, apply
\path{CartanMatrix.E_det}, which gives $9-k$ for every $k\geq3$.~\qed\end{proof}
\qedstep
\begin{proof}\pf\ Concatenate \stepref{rmdE1}--\stepref{rmdErest}.~\qed\end{proof}
\end{proof}
\step{rmdF}{The $\typ F$ determinants are $2,3,2,1$.}
\begin{proof}
\step{rmdFprefix}{The first three determinants are $2,3,2$.}
\begin{proof}\pf\ Substitute the edge products $1,2$ in
\stepref{rmdpath}, starting from $\Delta_0=1,\Delta_1=2$.~\qed\end{proof}
\step{rmdFfinal}{The fourth determinant is 1.}
\begin{proof}\pf\ Apply \texttt{CartanMatrix.F}\textsubscript{4}\texttt{\_det} in
\path{Mathlib/LinearAlgebra/Matrix/Cartan.lean}.~\qed\end{proof}
\qedstep
\begin{proof}\pf\ Combine \stepref{rmdFprefix}, \stepref{rmdFfinal}.~\qed\end{proof}
\end{proof}
\step{rmdG}{The $\typ G$ determinants are $2,1$.}
\begin{proof}\pf\ The first diagonal entry is 2; its final
determinant is 1 by \texttt{CartanMatrix.G}\textsubscript{2}\texttt{\_det} in
\path{Mathlib/LinearAlgebra/Matrix/Cartan.lean}.~\qed\end{proof}
\qedstep
\begin{proof}\pf\ Proper classical prefixes are unit paths; combine
\stepref{rmdA}--\stepref{rmdG}.~\qed\end{proof}
\end{proof}

\begin{namedresult}{Lemma}{(rational positivity certificates)}\label{comp:rem-ldl}
\lean{RootSetModel.exceptional_ldl}\lean{RootSetModel.symmetrizer_eq_modelSymmetrizer}
\uses{def:rem-matrices}
\assume{$M$ is $\typ E_6,\typ E_7,\typ E_8,\typ F_4$ or $\typ G_2$,
and $H=\operatorname{diag}(d)M$ with $d$ from \eqref{eq:symmetrizers}.}
\prove{There is a lower triangular $L$ with diagonal 1 such that
\begin{equation}\label{eq:ldl}
H=L\operatorname{diag}(q)L^T,\qquad
\Zx^TH\Zx=\sum_i q(i)((L^T\Zx)(i))^2,
\end{equation}
with every $q(i)>0$. The certificate data are:
\[
\begin{array}{c|l}
\typ E_8&q=(2,2,3/2,5/6,4/5,3/4,2/3,1/2)\\
\typ F_4&q=(2,3/2,4/3,1)\\
\typ G_2&q=(2,3/2).
\end{array}
\]
For $\typ E_8$ the nonzero strict lower entries are
\[
\begin{split}
L(2,0)&=-1/2,\quad L(3,1)=-1/2,\quad L(3,2)=-2/3,\\
L(4,3)&=-6/5,\quad L(5,4)=-5/4,\quad
L(6,5)=-4/3,\quad L(7,6)=-3/2.
\end{split}
\]
Use the leading six or seven entries and rows for $\typ E_6,\typ E_7$.
For $\typ F_4$ use $L(1,0)=-1/2$, $L(2,1)=-4/3$, $L(3,2)=-3/2$.
For $\typ G_2$ use $L(1,0)=-3/2$; all unspecified strict lower entries
are zero. These data prove $H$ positive definite.}
\end{namedresult}
\begin{proof}
\pfsketch\ These are bounded rational matrix identities, followed by
back substitution in a triangular matrix.
\step{rldlE}{The eight-row data satisfy
$H(i,j)=\sum_k L(i,k)q(k)L(j,k)$ for all $i,j<8$.}
\begin{proof}\pf\ Substitute the displayed rational entries and the
$\typ E_8$ entries of Definition~\ref{def:rem-matrices} in that finite sum. ~\qed\end{proof}
\step{rldlprefix}{The same equality holds for its six- and seven-row prefixes.}
\begin{proof}\pf\ Lower triangularity makes every omitted summand
zero in the prefix entries of \stepref{rldlE}.~\qed\end{proof}
\step{rldlF}{The equality holds for the four-row $\typ F$ data.}
\begin{proof}\pf\ Substitute those rational entries and $d=(1,1,2,2)$
in the finite entry sum. ~\qed\end{proof}
\step{rldlG}{The equality holds for the two-row $\typ G$ data.}
\begin{proof}\pf\ Substitute those rational entries and $d=(1,3)$
in the finite entry sum. ~\qed\end{proof}
\step{rldlform}{The quadratic identity in \eqref{eq:ldl} holds.}
\begin{proof}\pf\ Multiply the identities
\stepref{rldlE}--\stepref{rldlG} on the left and right by $\Zx^T,\Zx$.~\qed\end{proof}
\step{rldlinjective}{\assume{$L^T\Zx=0$.}
\prove{$\Zx=0$.}}
\begin{proof}\pf\ Backward induction on the coordinates of the upper
triangular matrix $L^T$ with diagonal 1. ~\qed\end{proof}
\qedstep
\begin{proof}\pf\ For $\Zx\ne0$, \stepref{rldlinjective} makes at least
one square in \stepref{rldlform} positive; every displayed $q(i)$ is
positive, so the sum is positive. ~\qed\end{proof}
\end{proof}

\begin{namedresult}{Lemma}{(rational coordinate frame)}\label{comp:rem-rational-frame}
\lean{RootSetModel.rational_frame}\lean{RootSetModel.exists_rat_simpleRootMatrix}
\lean{RootSetModel.exists_rat_inv_gram_simpleRootMatrix}
\uses{lem:classical,lem:E,lem:F,lem:G,lem:models-irreducible,lem:counts,lem:tight-frame,lem:rem-projectors,facts:inner,facts:rem-dimension}
\assume{A coordinate model from the construction ranges has $N$ root
lines and rank $n$; the matrix $A$ has its simple roots as columns in the
given coordinate space, which may be larger than its span.}
\prove{\begin{equation}\label{eq:rational-frame}
\frac12\sum_{\alpha\in\Phi}\frac{\alpha\alpha^T}{\rsq\alpha}
=\frac Nn\,A(A^TA)^{-1}A^T.
\end{equation}
All entries in this equation are rational.}
\end{namedresult}
\begin{proof}
\pfsketch\ Both sides act by $N/n$ on the root span and vanish on its
orthogonal complement.
\step{rrfprojection}{$P=A(A^TA)^{-1}A^T$ is the orthogonal projection
onto the column span.}
\begin{proof}
\step{rrfinvertible}{$A^TA$ is invertible.}
\begin{proof}\pf\ A null coefficient vector has combination of squared
norm zero, thus zero; the model bases are independent by
Lemmas~\ref{lem:classical}, \ref{lem:E}, \ref{lem:F}, \ref{lem:G}.~\qed\end{proof}
\step{rrffix}{$P(Ax)=Ax$.}
\begin{proof}\pf\ Cancel $(A^TA)^{-1}(A^TA)$ using \stepref{rrfinvertible}.~\qed\end{proof}
\step{rrfkill}{\assume{$A^T\Zy=0$.}
\prove{$P\Zy=0$.}}
\begin{proof}\pf\ Substitute in the definition of $P$.~\qed\end{proof}
\qedstep
\begin{proof}\pf\ The column span and its orthogonal complement are
the two summands tested in \stepref{rrffix}, \stepref{rrfkill}.~\qed\end{proof}
\end{proof}
\step{rrfleft}{The left side is the sum of the unit-line orthogonal projections.}
\begin{proof}\pf\ Each line has exactly its two opposite roots by
Lemma~\ref{lem:counts}; both summands are its orthogonal projection by
Lemma~\ref{lem:rem-projectors}.~\qed\end{proof}
\step{rrfequality}{The two operators agree on the ambient coordinate space.}
\begin{proof}\pf\ Lemma~\ref{lem:tight-frame} gives $(N/n)I$ on the
root span for \stepref{rrfleft}, and every root projection annihilates
the orthogonal complement; use \stepref{rrfprojection}.~\qed\end{proof}
\step{rrfrational}{All entries are rational.}
\begin{proof}\pf\ The model roots have rational coordinates and
positive rational squared norms; an invertible rational matrix has
rational inverse by \path{Matrix.inv_def}.~\qed\end{proof}
\qedstep
\begin{proof}\pf\ Combine \stepref{rrfequality}, \stepref{rrfrational}.~\qed\end{proof}
\end{proof}

\begin{cgss}{Remark}{3.11}\label{rem:3.11}
\uses{thm:A,thm:B,thm:C,cor:rem-constants}
Theorems A, B and C provide the finite-type classification, its passage
to Euclidean Root Systems, and all the unequal-length line assertions
of this remark. Corollary~\ref{cor:rem-constants} also recovers the
equal-length equations with the corrected $\typ E_8$ constant.
The classification does not use [CGSS, Theorem 3.10] as a premise.
Theorem 3.12 is proved separately (Section~\ref{sec:thm312}).
\end{cgss}
\endgroup

\subsection{Lines contained in a Root System}\label{sec:thm312}
\begin{mergedproofsection}
We will follow the ordinary proof in
\path{cgss/proofs/thm312/thm312.tex}. The construction is the explicit set
$\Psi(T)$ of Definition~\hyperref[def:312-integral-roots]{(the integer-span construction)}.
All spaces are real inner product spaces; finite dimension is stated
where needed. Entries of a vector or a matrix are function values.
The reflection $\tau_{\Zr}$ is $s_{\Zr}$ of
Definition~\hyperref[def:3.9]{3.9}. We will write $q(\Zx)=(\Zx,\Zx)$ and
$\re\Zx=\opk{span}_{\re}\{\Zx\}$. For a real scalar $a$,
$a\in\ints$ and
$a\in2\ints$ mean $a=z$ and $a=2z$, respectively, for some $z\in\ints$,
with the integer cast into $\re$ understood.

The proof format is that of \cite{L95}. Numbered assumptions of a result
are cited as \asm[1]{0}, and assumptions of an enclosing step as
\asm{1}, \asm{2}, and so on. Each proof ends with a \textsc{Qed} step;
each leaf gives its facts. The tags \new\ and \repair\ identify new
arguments and repairs to the printed paper. The metadata
\verb|\lean| and \verb|\uses| record the Lean declarations and the dependencies.

\subsubsection{Definitions and library facts}

\begin{namedresult}{Definition}{(angles and compatible colourings)}\label{def:312-angles}
\lean{LineSystem.HasSquaredCosine}
\lean{LineSystem.HasAllowedAngles}
\lean{LineSystem.HasBipartition}\uses{}
\assume{$V$ is a real inner product space, $S$ is a set of real
subspaces of $V$, and $\ell,m$ are real subspaces.}
For $a\in\re$, let $\opk{HasSquaredCosine}(\ell,m,a)$ mean
\[
 \forall\Zx\in\ell\ \forall\Zy\in m,\qquad
 (\Zx,\Zy)^2=a\,q(\Zx)q(\Zy).
\]
Let $\opk{HasAllowedAngles}(S)$ mean that for all distinct
$\ell,m\in S$ there is $a\in\{0,1/4,1/2\}$ with this property.
Let $\opk{HasBipartition}(S)$ mean that there is a function
$b:S\to\mathrm{Bool}$ such that, for all distinct $\ell,m\in S$,
\[
\begin{array}{ll}
b(\ell)=b(m)&\Longrightarrow
 \opk{HasSquaredCosine}(\ell,m,0)\ \lor\
 \opk{HasSquaredCosine}(\ell,m,1/4),\\
b(\ell)\ne b(m)&\Longrightarrow
 \opk{HasSquaredCosine}(\ell,m,0)\ \lor\
 \opk{HasSquaredCosine}(\ell,m,1/2).
\end{array}
\]
Neither colour fibre is required to be nonempty.
\end{namedresult}

\begin{namedresult}{Definition}{(lines of a vector set)}\label{def:312-root-lines}
\lean{LineSystem.rootLines}\uses{facts:312-linear}
\assume{$V$ is a real inner product space and $T\sbs V$.}
Define
\[
 \cL(T):=\{\ell:\exists\Zr\in T,\ \ell=\re\Zr\}.
\]
This is \path{LineSystem.rootLines}. When $0\notin T$, each member
has dimension one.
\end{namedresult}

\begin{namedresult}{Definition}{(the integer-span construction)}\label{def:312-integral-roots}
\lean{LineSystem.integralSpanRoots}\uses{}
\assume{$V$ is a real inner product space and $T\sbs V$.}
Put $L(T)=\opk{span}_{\ints}(T)$ and
\[
\begin{split}
 \Psi(T)={}&\{\Zx\in L(T):q(\Zx)=2\}\\
 &{}\cup\{\Zx\in L(T):q(\Zx)=4\text{ and }
                  (\Zx,\Zy)\in2\ints\text{ for every }\Zy\in L(T)\}.
\end{split}
\]
This is \path{LineSystem.integralSpanRoots}. The evenness condition
quantifies over the entire fixed integer span.
\end{namedresult}

\paragraph{Inner products and square roots.}\label{facts:312-inner}\uses{facts:inner}
\begingroup\raggedright
All library citations were checked in Mathlib at commit
\path{082e2d37e8b0463410cdb532e111cd43d5a66174} (4 September 2026),
the version pinned in \path{lake-manifest.json}.
We will use \hyperref[facts:inner]{Elementary identities}, together with
\path{real_inner_self_eq_norm_sq} and \path{abs_real_inner_le_norm} in
\path{Mathlib/Analysis/InnerProductSpace/Basic.lean}.
These give $(\Zx,\Zy)^2\leq q(\Zx)q(\Zy)$.
For positive real $a$, $\sqrt a>0$ and $(\sqrt a)^2=a$ are
\path{Real.sqrt_pos} and \path{Real.sq_sqrt} in
\path{Mathlib/Analysis/Real/Sqrt.lean}.
\par\endgroup

\paragraph{Spans and finite-dimensional linear algebra.}\label{facts:312-linear}\uses{}
\begingroup\raggedright
\path{finrank_eq_one_iff'} in
\path{Mathlib/LinearAlgebra/Dimension/FreeAndStrongRankCondition.lean}
gives a nonzero generator of a one-dimensional subspace.
\path{eq_span_singleton_of_mem_of_finrank_eq_one} and
\path{finrank_span_singleton} in
\path{Mathlib/LinearAlgebra/FiniteDimensional/Basic.lean}
identify its span and give the converse dimension assertion.
\path{Submodule.mem_span_singleton} in
\path{Mathlib/LinearAlgebra/Span/Defs.lean} expresses membership as a scalar multiple.
\path{Submodule.subset_span}, \path{Submodule.span_mono}, and
\path{Submodule.span_le} in the same file give membership of generators,
monotonicity, and the criterion for a span to lie in a subspace.
\path{Submodule.span_span_coe_preimage} there says that the preimage of
a generating set under the inclusion of its span spans the whole subtype.
\path{exists_linearIndependent} in
\path{Mathlib/LinearAlgebra/LinearIndependent/Lemmas.lean}
chooses an independent subset with the same real span.
\path{LinearIndependent.finite} in
\path{Mathlib/LinearAlgebra/Dimension/Finite.lean}
makes that subset finite in finite dimension.
\path{Submodule.span_induction} and \path{Submodule.span_induction}\textsubscript{2} in
\path{Mathlib/LinearAlgebra/Span/Defs.lean}
give induction on a span, also over $\ints$.
\path{Submodule.mem_span_range_iff_exists_fun} in
\path{Mathlib/LinearAlgebra/Finsupp/LinearCombination.lean}
expresses membership in a finite span by a finite sum.
\path{Int.cast_smul_eq_zsmul} in
\path{Mathlib/Algebra/Module/NatInt.lean} identifies multiplication by
a real-cast integer with integer scalar multiplication.
\par\endgroup

\paragraph{Finite sets and cardinalities.}\label{facts:312-finite}\uses{}
\begingroup\raggedright
\path{Set.Finite.of_injOn}, \path{Set.Finite.image},
\path{Set.Finite.subset}, and \path{Set.Finite.preimage} in
\path{Mathlib/Data/Set/Finite/Basic.lean} give finiteness from an
injection into a finite set, finiteness of images and subsets, and
finiteness of a preimage when the map is injective on that preimage.
A function set $A^I$ is finite when $A$ and $I$ are finite.
\par\endgroup

\paragraph{Integer arithmetic and logarithms.}\label{facts:312-arithmetic}\uses{}
Finite sums of integers are integers; finite sums of even integers are even.
We will use field arithmetic, factorization of differences of squares,
and integer inequalities in the arithmetic leaves.
\begingroup\raggedright
\path{zpow_add_one}\textsubscript{0} in
\path{Mathlib/Algebra/GroupWithZero/Basic.lean} gives the successor
power identity for a nonzero base, including base 2.
\path{Int.log} is the integer logarithm, and
\path{Int.zpow_log_le_self}, \path{Int.lt_zpow_succ_log_self},
\path{Int.zpow_le_iff_le_log}, and \path{Int.lt_zpow_iff_log_lt}
in \path{Mathlib/Data/Int/Log.lean} apply to base $2>1$ and $x>0$.
\path{Int.odd_iff} and \path{Int.not_even_iff_odd} in
\path{Mathlib/Algebra/Ring/Int/Parity.lean} identify the two parities,
also for negative integers. In particular $k+1$ and $k-1$ have
parity opposite to $k$, by reduction modulo 2.
\par\endgroup

\paragraph{Orthogonal complements.}\label{facts:312-orthogonal}\uses{}
\begingroup\raggedright
\path{stdOrthonormalBasis} in
\path{Mathlib/Analysis/InnerProductSpace/PiL2.lean}
gives an orthonormal basis indexed by the finite dimension.
\path{Submodule.inner_right_of_mem_orthogonal} in
\path{Mathlib/Analysis/InnerProductSpace/Orthogonal.lean} gives
$(\Zx,\Zy)=0$ for $\Zx\in W$ and $\Zy\in W^\perp$.
\path{Orthonormal.inner_eq_zero} in
\path{Mathlib/Analysis/InnerProductSpace/Orthonormal.lean} gives zero
inner products at distinct indices of an orthonormal family.
\path{Submodule.isCompl_orthogonal} and
\path{Submodule.HasOrthogonalProjection.ofCompleteSpace} in
\path{Mathlib/Analysis/InnerProductSpace/Projection/Basic.lean}
give $W+W^\perp=V$. Here $W$ is complete by the theorem
\path{FiniteDimensional.complete} in
\path{Mathlib/Topology/Algebra/Module/FiniteDimension.lean}.
\par\endgroup
These facts apply to the induced inner product on a subspace, including
the zero subspace.

\subsubsection{Representatives and finiteness}

\begin{namedresult}{Lemma}{(representatives of a line)}\label{lem:312-representatives}
\lean{LineSystem.exists_mem_line_sqNorm_eq}
\uses{facts:312-inner}
\assume{$V$ is a real inner product space, $\ell$ is a real subspace,
$\Zr\in\ell\diff\{0\}$, and $h>0$.}
\prove{There is $\Zt\in\ell$ with $q(\Zt)=h$.}
\end{namedresult}
\begin{proof}
\pfsketch\ We will rescale the given nonzero vector.
\step{t312repPositive}{$q(\Zr)>0$.}
\begin{proof}\pf\ Positive definiteness in the inner-product facts,
\asm{0}.~\qed\end{proof}
\pflet{$\Zt=(\sqrt h/\sqrt{q(\Zr)})\Zr$.}
\step{t312repPick}{$\Zt\in\ell$.}
\begin{proof}\pf\ The subspace $\ell$ contains $\Zr$ by \asm{0}
and is closed under real scalar multiplication.~\qed\end{proof}
\step{t312repNorm}{$q(\Zt)=h$.}
\begin{proof}\pf\ Bilinearity and the square-root facts give
$q(\Zt)=h\,q(\Zr)/q(\Zr)=h$, using $h>0$ in \asm{0} and
\stepref{t312repPositive}.~\qed\end{proof}
\qedstep
\begin{proof}\pf\ The witness $\Zt$ satisfies \stepref{t312repPick}
and \stepref{t312repNorm}.~\qed\end{proof}
\end{proof}

\begin{namedresult}{Lemma}{(squared cosine from representatives)}\label{lem:312-squared-cosine}
\lean{LineSystem.hasSquaredCosine_iff}
\lean{LineSystem.HasSquaredCosine.unique}
\uses{def:312-angles,facts:312-inner,facts:312-linear}
\assume{\begin{pfenum}
 \item $\ell$ and $m$ are subspaces of dimension one in the
 real inner product space $V$.
 \item $\Zr\in\ell\diff\{0\}$ and $\Zs\in m\diff\{0\}$.
\end{pfenum}}
\prove{For every $a\in\re$,
\[
 \opk{HasSquaredCosine}(\ell,m,a)
 \ \Longleftrightarrow\ (\Zr,\Zs)^2=a\,q(\Zr)q(\Zs).
\]
For all $a,a'\in\re$, if $\opk{HasSquaredCosine}(\ell,m,a)$
and $\opk{HasSquaredCosine}(\ell,m,a')$, then $a=a'$.}
\end{namedresult}
\begin{proof}
\pfsketch\ We will write vectors on the two lines as scalar multiples
of their representatives, then cancel their positive squared norms.
\step{t312repProductPositive}{$q(\Zr)q(\Zs)>0$.}
\begin{proof}\pf\ Positive definiteness and \asm[2]{0}.~\qed\end{proof}
\step{t312repForward}{\assume{$a\in\re$ and $\opk{HasSquaredCosine}(\ell,m,a)$.}
 \prove{$(\Zr,\Zs)^2=a\,q(\Zr)q(\Zs)$.}}
\begin{proof}\pf\ Substitute \asm[2]{0} in Definition~\hyperref[def:312-angles]{(angles and compatible colourings)}.~\qed\end{proof}
\step{t312repReverse}{\assume{$a\in\re$ and $(\Zr,\Zs)^2=a\,q(\Zr)q(\Zs)$.}
 \prove{$\opk{HasSquaredCosine}(\ell,m,a)$.}}
\begin{proof}
 \step{t312repMultiples}{For every $\Zx\in\ell$ and $\Zy\in m$ there
 are $u,v\in\re$ with $\Zx=u\Zr$ and $\Zy=v\Zs$.}
 \begin{proof}\pf\ Apply the singleton-span facts to each nonzero
 representative in \asm[2]{0}.~\qed\end{proof}
 \qedstep
 \begin{proof}\pf\ Multiply \asm{1} by $u^2v^2$ and use bilinearity
 with \stepref{t312repMultiples}; this includes $u=0$ or $v=0$.~\qed\end{proof}
\end{proof}
\step{t312repUnique}{\assume{$a,a'\in\re$,
$\opk{HasSquaredCosine}(\ell,m,a)$ and
$\opk{HasSquaredCosine}(\ell,m,a')$.}\prove{$a=a'$.}}
\begin{proof}\pf\ Apply \stepref{t312repForward} at $a$ and $a'$.
The resulting equalities give
$(a-a')q(\Zr)q(\Zs)=0$; cancel the positive product in
\stepref{t312repProductPositive}.~\qed\end{proof}
\qedstep
\begin{proof}\pf\ Steps \stepref{t312repForward} and \stepref{t312repReverse}
give the equivalence; \stepref{t312repUnique} gives uniqueness.~\qed\end{proof}
\end{proof}

For lines, $\opk{HasSquaredCosine}(\ell,m,a)$ with
$a=0,1/4,1/2$ means the unoriented angle is
$90^\circ,60^\circ,45^\circ$, respectively. We will use the
division-free predicate and Lemma~\hyperref[lem:312-squared-cosine]{(squared cosine from representatives)};
no trigonometric fact is a premise below.

\begin{namedresult}{Lemma}{(inner-product coordinates)}\label{lem:312-coordinates}
\lean{LineSystem.inner_coordinates_injective}
\uses{facts:312-inner,facts:312-linear}
\assume{\begin{pfenum}
 \item $V$ is a real inner product space, $W$ a real subspace,
 and $(\Zb_i)_{i\in I}$ a family spanning $W$.
 \item $\Zx,\Zy\in W$ and $(\Zx,\Zb_i)=(\Zy,\Zb_i)$ for every $i\in I$.
\end{pfenum}}
\prove{$\Zx=\Zy$. Equivalently, $\Zx\mapsto((\Zx,\Zb_i))_{i\in I}$ is
injective on $W$.}
\end{namedresult}
\begin{proof}
\pfsketch\ We will show that the difference is orthogonal to its own span.
\step{t312coordBasis}{$(\Zx-\Zy,\Zb_i)=0$ for every $i\in I$.}
\begin{proof}\pf\ Subtract the equalities of \asm[2]{0} by bilinearity.~\qed\end{proof}
\step{t312coordSpan}{$(\Zx-\Zy,\Zw)=0$ for every $\Zw\in W$.}
\begin{proof}\pf\ Real span induction from \stepref{t312coordBasis}:
the assertion holds at zero and is preserved by addition and real
scalar multiplication, by bilinearity.~\qed\end{proof}
\step{t312coordSelf}{$q(\Zx-\Zy)=0$.}
\begin{proof}\pf\ Substitute $\Zw=\Zx-\Zy\in W$ in \stepref{t312coordSpan}.~\qed\end{proof}
\qedstep
\begin{proof}\pf\ Positive definiteness applied to \stepref{t312coordSelf}.~\qed\end{proof}
\end{proof}

\begin{namedresult}{Lemma}{(finitely many lines)}\label{lem:312-finite}
\lean{LineSystem.finite_of_hasAllowedAngles}
\uses{def:312-angles,lem:312-representatives,lem:312-coordinates,facts:312-inner,facts:312-linear,facts:312-finite,facts:312-arithmetic}
\assume{\begin{pfenum}
 \item $V$ is a finite-dimensional real inner product space and
 $S$ is a set of subspaces of dimension one.
 \item $\opk{HasAllowedAngles}(S)$.
\end{pfenum}}
\prove{$S$ is finite.}
\end{namedresult}
\begin{proof}
\pfsketch\ We will encode all unit representatives by their inner
products with a finite spanning family chosen from those representatives.
\pflet{\[
\begin{gathered}
U=\{\Zx:\exists\ell\in S,\ \Zx\in\ell\text{ and }q(\Zx)=1\},
\qquad W=\opk{span}_{\re}U,\\
\Theta=\{-1,-1/\sqrt2,-1/2,0,1/2,1/\sqrt2,1\}.
\end{gathered}
\]}
\step{t312finiteBasis}{\pick\ a finite family $(\Zb_i)_{i\in I}$ spanning $W$
with every $\Zb_i\in U$.}
\begin{proof}\pf\ The independent-subset and finite-dimensional facts
in \hyperref[facts:312-linear]{Spans and finite-dimensional linear algebra}, applied to $U$, give this family. The index set
is empty if $U$ is empty.~\qed\end{proof}
\step{t312finiteValues}{For every $\Zx\in U$ and $i\in I$,
$(\Zx,\Zb_i)\in\Theta$.}
\begin{proof}
 \step{t312finiteLines}{\pick\ $\ell,m\in S$ with $\Zx\in\ell$,
 $\Zb_i\in m$, and $q(\Zx)=q(\Zb_i)=1$.}
 \begin{proof}\pf\ The definition of $U$ and \stepref{t312finiteBasis}.~\qed\end{proof}
 \pflet{$p=(\Zx,\Zb_i)$.}
 \step{t312finiteSame}{\case{$\ell=m$.}\prove{$p\in\{-1,1\}$.}}
 \begin{proof}
  \step{t312finiteScalar}{There is $t\in\re$ with $\Zx=t\Zb_i$ and $t^2=1$.}
  \begin{proof}\pf\ The singleton-span facts give the scalar; the two
  unit norms give its square by bilinearity.~\qed\end{proof}
  \qedstep
  \begin{proof}\pf\ Factor $t^2-1$ in \stepref{t312finiteScalar};
  then $p=t\in\{-1,1\}$.~\qed\end{proof}
 \end{proof}
 \step{t312finiteDistinct}{\case{$\ell\ne m$.}\prove{$p\in\Theta$.}}
 \begin{proof}
  \step{t312finiteSquare}{$p^2\in\{0,1/4,1/2\}$.}
  \begin{proof}\pf\ Substitute the unit representatives in
  \asm[2]{0}, using Definition~\hyperref[def:312-angles]{(angles and compatible colourings)}.~\qed\end{proof}
  \step{t312finiteZero}{\case{$p^2=0$.}\prove{$p=0$.}}
  \begin{proof}\pf\ A real square is zero only at zero.~\qed\end{proof}
  \step{t312finiteQuarter}{\case{$p^2=1/4$.}\prove{$p\in\{-1/2,1/2\}$.}}
  \begin{proof}\pf\ Factor $(p-1/2)(p+1/2)=0$.~\qed\end{proof}
  \step{t312finiteHalf}{\case{$p^2=1/2$.}\prove{$p\in\{-1/\sqrt2,1/\sqrt2\}$.}}
  \begin{proof}\pf\ The square-root facts give $(1/\sqrt2)^2=1/2$;
  factor $(p-1/\sqrt2)(p+1/\sqrt2)=0$.~\qed\end{proof}
  \qedstep
  \begin{proof}\pf\ Step \stepref{t312finiteSquare} makes
  \stepref{t312finiteZero}, \stepref{t312finiteQuarter}, and
  \stepref{t312finiteHalf} exhaustive.~\qed\end{proof}
 \end{proof}
 \qedstep
 \begin{proof}\pf\ The lines are equal or distinct; apply
 \stepref{t312finiteSame} and \stepref{t312finiteDistinct}.~\qed\end{proof}
\end{proof}
\step{t312finiteU}{$U$ is finite.}
\begin{proof}\pf\ By Lemma~\hyperref[lem:312-coordinates]{(inner-product coordinates)} and
\stepref{t312finiteBasis}, the coordinate map is injective on $U\sbs W$.
Step \stepref{t312finiteValues} puts its image in the finite set
$\Theta^I$; use the finite-set facts in \hyperref[facts:312-finite]{Finite sets and cardinalities}.~\qed\end{proof}
\step{t312finiteOnto}{$\{\re\Zx:\Zx\in U\}=S$.}
\begin{proof}
 \step{t312finiteImage}{For $\Zx\in U$, $\re\Zx\in S$.}
 \begin{proof}\pf\ The witness line in the definition of $U$ equals
 $\re\Zx$ by the singleton-span facts, since $q(\Zx)=1$.~\qed\end{proof}
 \step{t312finiteSurj}{For $\ell\in S$, some $\Zx\in U$ satisfies $\ell=\re\Zx$.}
 \begin{proof}\pf\ A nonzero generator exists by the dimension-one
 facts. Lemma~\hyperref[lem:312-representatives]{(representatives of a line)}, with $h=1$, normalizes it;
 the singleton-span facts identify its line.~\qed\end{proof}
 \qedstep
 \begin{proof}\pf\ The two inclusions are \stepref{t312finiteImage}
 and \stepref{t312finiteSurj}.~\qed\end{proof}
\end{proof}
\qedstep
\begin{proof}\pf\ The image of the finite set in \stepref{t312finiteU}
is finite; use \stepref{t312finiteOnto}.~\qed\end{proof}
\end{proof}

\begin{namedresult}{Lemma}{(a bounded part of an integer span)}\label{lem:312-integer-span}
\lean{LineSystem.inner_int_of_mem_intSpan}
\lean{LineSystem.inner_even_of_mem_intSpan}
\lean{LineSystem.finite_bounded_intSpan}
\uses{lem:312-coordinates,facts:312-inner,facts:312-linear,facts:312-finite,facts:312-arithmetic}
\assume{\begin{pfenum}
 \item $V$ is a finite-dimensional real inner product space and $T\sbs V$ is finite.
 \item $q(\Zt)\in\{2,4\}$ for $\Zt\in T$, and
 $(\Zt,\Zu)\in\ints$ for $\Zt,\Zu\in T$.
\end{pfenum}}
\pflet{$L=\opk{span}_{\ints}T$, $W=\opk{span}_{\re}T$,
$B=\{\Zx\in L:q(\Zx)\leq4\}$.}
\prove{\begin{pfenum}
 \item $L\sbs W$, and $(\Zx,\Zy)\in\ints$ for all $\Zx,\Zy\in L$.
 \item For every $\Zt\in T$ such that $(\Zt,\Zu)\in2\ints$ for all
 $\Zu\in T$, we have $(\Zt,\Zx)\in2\ints$ for all $\Zx\in L$.
 \item $B$ is finite.
\end{pfenum}}
\end{namedresult}
\begin{proof}
\pfsketch\ \new\ We will expand integer combinations and use bounded
inner-product coordinates against $T$ itself. This simplifies the
ordinary draft's basis argument; no independence of $T$ is assumed.
\step{t312intSums}{Every $\Zx\in L$ has an expression
$\Zx=\sum_{\Zt\in T}a(\Zt)\Zt$ with $a:T\to\ints$.}
\begin{proof}\pf\ Apply the finite-span formula in
\hyperref[facts:312-linear]{Spans and finite-dimensional linear algebra} over $\ints$ to the inclusion $T\to V$.~\qed\end{proof}
\step{t312intReal}{$L\sbs W$.}
\begin{proof}\pf\ Cast the integer coefficients in
\stepref{t312intSums} to real coefficients.~\qed\end{proof}
\step{t312intPair}{For all $\Zx,\Zy\in L$, $(\Zx,\Zy)\in\ints$.}
\begin{proof}
 \step{t312intExpand}{For coefficients $a,b:T\to\ints$,
 \[
 \left(\sum_{\Zt\in T}a(\Zt)\Zt,\sum_{\Zu\in T}b(\Zu)\Zu\right)
 =\sum_{\Zt,\Zu\in T}a(\Zt)b(\Zu)(\Zt,\Zu).
 \]}
 \begin{proof}\pf\ Bilinearity, iterated over the two finite sums.~\qed\end{proof}
 \qedstep
 \begin{proof}\pf\ Each summand in \stepref{t312intExpand} is an integer
 by \asm[2]{0}; use \stepref{t312intSums} and closure under finite sums.~\qed\end{proof}
\end{proof}
\step{t312intEven}{\assume{$\Zt\in T$ and $(\Zt,\Zu)\in2\ints$ for every $\Zu\in T$.}
 \prove{$(\Zt,\Zx)\in2\ints$ for every $\Zx\in L$.}}
\begin{proof}\pf\ Bilinearity and \stepref{t312intSums} give
$(\Zt,\Zx)=\sum_{\Zu\in T}a(\Zu)(\Zt,\Zu)$; every summand is even by
\asm{1}, so the sum is even.~\qed\end{proof}
\step{t312intBound}{For $\Zx\in B$ and $\Zt\in T$,
$(\Zx,\Zt)\in\{-4,-3,\ldots,4\}$.}
\begin{proof}
 \step{t312intCS}{$(\Zx,\Zt)^2\leq16$.}
 \begin{proof}\pf\ Cauchy--Schwarz in \hyperref[facts:312-inner]{Inner products and square roots} gives
 $(\Zx,\Zt)^2\leq q(\Zx)q(\Zt)\leq4\cdot4$,
 by the definition of $B$, $\Zt\in T$, and \asm[2]{0}.~\qed\end{proof}
 \step{t312intDiscrete}{$(\Zx,\Zt)\in\ints$.}
 \begin{proof}\pf\ Step \stepref{t312intPair}, since $\Zx\in L$ and
 $\Zt\in T\sbs L$ by \path{Submodule.subset_span}.~\qed\end{proof}
 \qedstep
 \begin{proof}\pf\ An integer whose square is at most 16 lies between
 $-4$ and 4; apply \stepref{t312intCS} and \stepref{t312intDiscrete}.~\qed\end{proof}
\end{proof}
\step{t312intFinite}{$B$ is finite.}
\begin{proof}\pf\ The family indexed by $T$ spans $W$ by definition.
Lemma~\hyperref[lem:312-coordinates]{(inner-product coordinates)} makes
$\Zx\mapsto((\Zx,\Zt))_{\Zt\in T}$ injective on $W$, which contains
$B$ by \stepref{t312intReal}. Step \stepref{t312intBound} puts its
image in the finite function set $\{-4,-3,\ldots,4\}^{T}$, since
$T$ is finite by \asm[1]{0}. Apply \path{Set.Finite.of_injOn}.~\qed\end{proof}
\qedstep
\begin{proof}\pf\ Steps \stepref{t312intReal}, \stepref{t312intPair},
\stepref{t312intEven}, and \stepref{t312intFinite}. For $T=\varnothing$
the same proof uses the empty coordinate family and gives $B=\{0\}$.~\qed\end{proof}
\end{proof}

\subsubsection{The explicit construction}

\begin{namedresult}{Lemma}{(reflection adjointness)}\label{lem:312-reflection}
\lean{IsRootSet.reflection_self_adjoint}
\uses{lem:rem-reflection}
\assume{$V$ is a real inner product space and $\Zr\in V\diff\{0\}$.}
\prove{For all $\Zx,\Zy\in V$,
$(\tau_{\Zr}\Zx,\Zy)=(\Zx,\tau_{\Zr}\Zy)$.}
\end{namedresult}
\begin{proof}
\pfsketch\ We will use the metric identity and involution in
Lemma~\hyperref[lem:rem-reflection]{(reflection identities)}.
\step{t312reflAdjoint}{$(\tau_{\Zr}\Zx,\Zy)=(\Zx,\tau_{\Zr}\Zy)$ for all $\Zx,\Zy\in V$.}
\begin{proof}\pf\ Apply the metric identity in Lemma~\hyperref[lem:rem-reflection]{(reflection identities)}
to $\Zx,\tau_{\Zr}\Zy$; its involution identity changes the left side to
$(\tau_{\Zr}\Zx,\Zy)$. Its nonzero hypothesis is \asm{0}.~\qed\end{proof}
\qedstep
\begin{proof}\pf\ Step \stepref{t312reflAdjoint}.~\qed\end{proof}
\end{proof}

\begin{namedresult}{Lemma}{(integral proportional vectors)}\label{lem:312-proportional}
\lean{LineSystem.eq_one_or_neg_one_of_integral_proportional}
\uses{facts:312-inner,facts:312-arithmetic}
\assume{\begin{pfenum}
 \item $V$ is a real inner product space, $\Zr,\Zs\in V$,
 and $q(\Zr),q(\Zs)\in\{2,4\}$.
 \item $(\Zr,\Zs)\in\ints$ and $\Zs=t\Zr$ for $t\in\re$.
\end{pfenum}}
\prove{$t=1$ or $t=-1$.}
\end{namedresult}
\begin{proof}
\pfsketch\ \new\ We will use $t^2=1$ for equal lengths. For unequal
lengths the integer inner product would have square 8, simplifying
the ordinary draft's irrational-square-root argument.
\step{t312propNorm}{$q(\Zs)=t^2q(\Zr)$.}
\begin{proof}\pf\ Bilinearity and $\Zs=t\Zr$ in \asm[2]{0}.~\qed\end{proof}
\step{t312propEqual}{\case{$q(\Zr)=q(\Zs)$.}\prove{$t\in\{1,-1\}$.}}
\begin{proof}\pf\ Divide \stepref{t312propNorm} by $q(\Zr)>0$ and
factor $t^2-1$.~\qed\end{proof}
\step{t312propSquare}{$(\Zr,\Zs)^2=q(\Zr)q(\Zs)$.}
\begin{proof}\pf\ Bilinearity gives $(\Zr,\Zs)=tq(\Zr)$.
Square this equality and substitute \stepref{t312propNorm}.~\qed\end{proof}
\step{t312propUnequal}{\case{$q(\Zr)\ne q(\Zs)$.}\prove{False.}}
\begin{proof}
 \step{t312propSquareEight}{$(\Zr,\Zs)^2=8$.}
 \begin{proof}\pf\ The two unequal members of $\{2,4\}$ in \asm[1]{0}
 have product 8; substitute in \stepref{t312propSquare}.~\qed\end{proof}
 \step{t312propInteger}{\pick\ $z\in\ints$ with $(\Zr,\Zs)=z$.}
 \begin{proof}\pf\ Integrality in \asm[2]{0}.~\qed\end{proof}
 \qedstep
 \begin{proof}\pf\ Steps \stepref{t312propSquareEight} and
 \stepref{t312propInteger} give $z^2=8$. Thus $4<|z|^2<9$, so
 $2<|z|<3$, impossible for an integer.~\qed\end{proof}
\end{proof}
\qedstep
\begin{proof}\pf\ Equality or inequality of the squared lengths gives
\stepref{t312propEqual} or \stepref{t312propUnequal}.
The unequal-length case is impossible.~\qed\end{proof}
\end{proof}

\begin{namedresult}{Proposition}{(the integer-span Root System)}\label{prop:312-rootset}
\lean{LineSystem.isRootSet_integralSpanRoots}
\uses{def:3.9,def:312-integral-roots,lem:312-integer-span,lem:rem-reflection,lem:312-reflection,lem:312-proportional,facts:312-inner,facts:312-linear,facts:312-finite,facts:312-arithmetic}
\assume{\begin{pfenum}
 \item $V$ is a finite-dimensional real inner product space,
 $T\sbs V$ is finite, and $\opk{span}_{\re}T=V$.
 \item $q(\Zt)\in\{2,4\}$ for every $\Zt\in T$.
 \item $(\Zt,\Zu)\in\ints$ for every $\Zt,\Zu\in T$.
 \item If $\Zt\in T$ and $q(\Zt)=4$, then $(\Zt,\Zu)\in2\ints$
 for every $\Zu\in T$.
\end{pfenum}}
\prove{$\opk{IsRootSet}(\Psi(T))$, $T\sbs\Psi(T)$, and
$q(\Zr)\in\{2,4\}$ for every $\Zr\in\Psi(T)$.}
\end{namedresult}
\begin{proof}
\pfsketch\ \repair\ This fills part of Gap 11. The explicit set in
Definition~\hyperref[def:312-integral-roots]{(the integer-span construction)} replaces a reflection closure,
which would need stages, a union and a leastness argument.
We will check each Root System
requirement, keeping the evenness condition invariant under reflection.
\pflet{$L=\opk{span}_{\ints}T$ and $\Psi=\Psi(T)$.}
\step{t312psiLengths}{$\Psi\sbs L$, every $\Zr\in\Psi$ has
$q(\Zr)\in\{2,4\}$, and $0\notin\Psi$.}
\begin{proof}\pf\ Definition~\hyperref[def:312-integral-roots]{(the integer-span construction)} and $q(0)=0$.~\qed\end{proof}
\step{t312psiSeeds}{$T\sbs\Psi$.}
\begin{proof}
 \step{t312psiSeedShort}{\assume{$\Zt\in T$ and $q(\Zt)=2$.}\prove{$\Zt\in\Psi$.}}
 \begin{proof}\pf\ By \path{Submodule.subset_span}, $T\sbs L$.
 The membership and length in \asm{2} give the first clause of
 Definition~\hyperref[def:312-integral-roots]{(the integer-span construction)}.~\qed\end{proof}
 \step{t312psiSeedLong}{\assume{$\Zt\in T$ and $q(\Zt)=4$.}\prove{$\Zt\in\Psi$.}}
 \begin{proof}
  \step{t312psiSeedEven}{$(\Zt,\Zx)\in2\ints$ for every $\Zx\in L$.}
  \begin{proof}\pf\ Lemma~\hyperref[lem:312-integer-span]{(a bounded part of an integer span)}, part 2,
  applied with \asm[1]{0}, \asm[2]{0}, and \asm[3]{0};
  the vector belongs to $T$ and has squared norm 4 by \asm{2},
  so \asm[4]{0} supplies its even pairings with $T$.~\qed\end{proof}
  \qedstep
  \begin{proof}\pf\ The second clause of Definition~\hyperref[def:312-integral-roots]{(the integer-span construction)}
  follows from $\Zt\in L$, the case length, and \stepref{t312psiSeedEven}.~\qed\end{proof}
 \end{proof}
 \qedstep
 \begin{proof}\pf\ Hypothesis \asm[2]{0} makes
 \stepref{t312psiSeedShort} and \stepref{t312psiSeedLong} exhaustive.~\qed\end{proof}
\end{proof}
\step{t312psiFinite}{$\Psi$ is finite.}
\begin{proof}\pf\ By \stepref{t312psiLengths}, it is contained in
$\{\Zx\in L:q(\Zx)\leq4\}$, finite by part 3 of
Lemma~\hyperref[lem:312-integer-span]{(a bounded part of an integer span)} with \asm[1]{0}, \asm[2]{0}, and
\asm[3]{0}. Apply \path{Set.Finite.subset}.~\qed\end{proof}
\step{t312psiSpan}{$\opk{span}_{\re}\Psi=V$.}
\begin{proof}\pf\ \path{Submodule.span_mono}, \stepref{t312psiSeeds}, and \asm[1]{0}.~\qed\end{proof}
\step{t312psiCoefficient}{For $\Zr\in\Psi$ and $\Zx\in L$,
$2(\Zr,\Zx)/q(\Zr)\in\ints$.}
\begin{proof}
 \step{t312psiCoeffShort}{\case{$q(\Zr)=2$.}}
 \begin{proof}\pf\ The coefficient is $(\Zr,\Zx)$, an integer by
 Lemma~\hyperref[lem:312-integer-span]{(a bounded part of an integer span)} and \stepref{t312psiLengths}.~\qed\end{proof}
 \step{t312psiCoeffLong}{\case{$q(\Zr)=4$.}}
 \begin{proof}\pf\ The coefficient is $(\Zr,\Zx)/2$, an integer by the
 evenness clause of Definition~\hyperref[def:312-integral-roots]{(the integer-span construction)}.~\qed\end{proof}
 \qedstep
 \begin{proof}\pf\ Step \stepref{t312psiLengths} makes
 \stepref{t312psiCoeffShort} and \stepref{t312psiCoeffLong} exhaustive.~\qed\end{proof}
\end{proof}
\step{t312psiLattice}{For $\Zr\in\Psi$ and $\Zx\in L$, $\tau_{\Zr}\Zx\in L$.}
\begin{proof}\pf\ Step \stepref{t312psiCoefficient} gives an integer $z$
equal to the reflection coefficient. By \path{Int.cast_smul_eq_zsmul},
the term $(z:\re)\Zr$ is the integer multiple $z\Zr$.
Step \stepref{t312psiLengths} puts $\Zr$ in $L$; the integer submodule
$L$ is closed under integer multiples and subtraction.~\qed\end{proof}
\step{t312psiClosed}{For $\Zr,\Zs\in\Psi$, $\tau_{\Zr}\Zs\in\Psi$.}
\begin{proof}
 \pflet{$\Zu=\tau_{\Zr}\Zs$.}
 \step{t312psiImageNorm}{$\Zu\in L$ and $q(\Zu)=q(\Zs)$.}
 \begin{proof}\pf\ Step \stepref{t312psiLattice} and
 the metric identity of Lemma~\hyperref[lem:rem-reflection]{(reflection identities)} at
 $\Zs,\Zs$, using the nonzero assertion in \stepref{t312psiLengths}.~\qed\end{proof}
 \step{t312psiImageShort}{\case{$q(\Zs)=2$.}\prove{$\Zu\in\Psi$.}}
 \begin{proof}\pf\ The first clause of Definition~\hyperref[def:312-integral-roots]{(the integer-span construction)},
 with \stepref{t312psiImageNorm}.~\qed\end{proof}
 \step{t312psiImageLong}{\case{$q(\Zs)=4$.}\prove{$\Zu\in\Psi$.}}
 \begin{proof}
  \step{t312psiImageEven}{For every $\Zx\in L$, $(\Zu,\Zx)\in2\ints$.}
  \begin{proof}\pf\ Lemma~\hyperref[lem:312-reflection]{(reflection adjointness)} gives
  $(\Zu,\Zx)=(\Zs,\tau_{\Zr}\Zx)$. This is even by the long-root
  clause for $\Zs$ and \stepref{t312psiLattice}.~\qed\end{proof}
  \qedstep
  \begin{proof}\pf\ Apply the second clause of
  Definition~\hyperref[def:312-integral-roots]{(the integer-span construction)} to \stepref{t312psiImageNorm}
  and \stepref{t312psiImageEven}.~\qed\end{proof}
 \end{proof}
 \qedstep
 \begin{proof}\pf\ Step \stepref{t312psiLengths} makes
 \stepref{t312psiImageShort} and \stepref{t312psiImageLong} exhaustive.~\qed\end{proof}
\end{proof}
\step{t312psiReduced}{For $\Zr,t\Zr\in\Psi$, $t\in\{1,-1\}$.}
\begin{proof}\pf\ Apply Lemma~\hyperref[lem:312-proportional]{(integral proportional vectors)}; its length
hypotheses are \stepref{t312psiLengths}, and its inner-product hypothesis
is part 1 of Lemma~\hyperref[lem:312-integer-span]{(a bounded part of an integer span)}, applied with
\asm[1]{0}, \asm[2]{0}, and \asm[3]{0} and the inclusion
$\Psi\sbs L$ from \stepref{t312psiLengths}.~\qed\end{proof}
\qedstep
\begin{proof}\pf\ Definition~\hyperref[def:3.9]{3.9} has six requirements:
\stepref{t312psiFinite}, \stepref{t312psiLengths}, \stepref{t312psiSpan},
\stepref{t312psiCoefficient} restricted to $\Psi\sbs L$,
\stepref{t312psiReduced}, and \stepref{t312psiClosed} supply them.
The additional conclusions are \stepref{t312psiSeeds} and \stepref{t312psiLengths}.~\qed\end{proof}
\end{proof}

\begin{namedresult}{Lemma}{(completion of the seeds)}\label{lem:312-completion}
\lean{LineSystem.exists_spanning_integral_seed_extension}
\uses{facts:312-inner,facts:312-linear,facts:312-finite,facts:312-arithmetic,facts:312-orthogonal}
\assume{\begin{pfenum}
 \item $V$ is a finite-dimensional real inner product space and $T\sbs V$ is finite.
 \item $q(\Zt)\in\{2,4\}$ for every $\Zt\in T$.
 \item $(\Zt,\Zu)\in\ints$ for every $\Zt,\Zu\in T$.
 \item If $\Zt\in T$ and $q(\Zt)=4$, then $(\Zt,\Zu)\in2\ints$
 for every $\Zu\in T$.
\end{pfenum}}
\prove{There is a finite set $T'\sps T$ spanning $V$ with
$q(\Zt)\in\{2,4\}$ for $\Zt\in T'$, all inner products in $T'$
integral, and $(\Zt,\Zu)\in2\ints$ whenever
$\Zt,\Zu\in T'$ and $q(\Zt)=4$.}
\end{namedresult}
\begin{proof}
\pfsketch\ We will add scaled orthonormal vectors of the orthogonal
complement. The diagonal inner product of a new vector is 2.
\pflet{$W=\opk{span}_{\re}T$.}
\step{t312completeBasis}{\pick\ a finite orthonormal basis
$(\Ze_i)_{i\in I}$ of $W^\perp$, viewed in $V$.}
\begin{proof}\pf\ The orthogonal-complement facts in
\hyperref[facts:312-orthogonal]{Orthogonal complements}.~\qed\end{proof}
\pflet{$\Zu_i=\sqrt2\Ze_i$, $E=\{\Zu_i:i\in I\}$, and $T'=T\cup E$.}
\step{t312completeNorm}{$q(\Zu_i)=2$ for every $i\in I$.}
\begin{proof}\pf\ Orthonormality in \stepref{t312completeBasis} and
bilinearity with $(\sqrt2)^2=2$.~\qed\end{proof}
\step{t312completePairs}{All inner products in $T'$ are integers.}
\begin{proof}
 \step{t312completeOldOld}{\assume{$\Zt,\Zu\in T$.}\prove{$(\Zt,\Zu)\in\ints$.}}
 \begin{proof}\pf\ Hypothesis \asm[3]{0}.~\qed\end{proof}
 \step{t312completeOldNew}{\assume{$\Zt\in T$ and $i\in I$.}\prove{$(\Zt,\Zu_i)=0$.}}
 \begin{proof}\pf\ Here $\Zt\in W$ and $\Ze_i\in W^\perp$ by
 \stepref{t312completeBasis}; \path{Submodule.subset_span} gives the
 first membership. Apply \path{Submodule.inner_right_of_mem_orthogonal}
 and multiply the zero inner product by $\sqrt2$.~\qed\end{proof}
 \step{t312completeNewOld}{\assume{$\Zt\in T$ and $i\in I$.}\prove{$(\Zu_i,\Zt)=0$.}}
 \begin{proof}\pf\ Symmetry and \stepref{t312completeOldNew}.~\qed\end{proof}
 \step{t312completeNewDiagonal}{\assume{$i\in I$.}\prove{$(\Zu_i,\Zu_i)=2$.}}
 \begin{proof}\pf\ Step \stepref{t312completeNorm}.~\qed\end{proof}
 \step{t312completeNewDistinct}{\assume{$i,j\in I$ and $i\ne j$.}\prove{$(\Zu_i,\Zu_j)=0$.}}
 \begin{proof}\pf\ Apply \path{Orthonormal.inner_eq_zero} to
 \stepref{t312completeBasis} and $i\ne j$ in \asm{2};
 bilinearity preserves zero after scaling.~\qed\end{proof}
 \qedstep
 \begin{proof}\pf\ Membership in $T\cup E$ gives the old/old,
 old/new, new/old, and new/new possibilities. The last splits into
 equal and distinct indices. The preceding five steps give integers
 in every case.~\qed\end{proof}
\end{proof}
\step{t312completeEven}{For $\Zt,\Zu\in T'$ with $q(\Zt)=4$,
$(\Zt,\Zu)\in2\ints$.}
\begin{proof}
 \step{t312completeLongOld}{$\Zt\in T$.}
 \begin{proof}\pf\ Step \stepref{t312completeNorm} excludes $\Zt\in E$;
 use $\Zt\in T\cup E$.~\qed\end{proof}
 \step{t312completeEvenOld}{\case{$\Zu\in T$.}}
 \begin{proof}\pf\ Hypothesis \asm[4]{0} and \stepref{t312completeLongOld}.~\qed\end{proof}
 \step{t312completeEvenNew}{\case{$\Zu\in E$.}}
 \begin{proof}\pf\ Step \stepref{t312completeLongOld} and
 \path{Submodule.subset_span} put $\Zt$ in $W$.
 The definition of $E$ puts $\Zu$ in $W^\perp$.
 The inner product is zero by
 \path{Submodule.inner_right_of_mem_orthogonal}.~\qed\end{proof}
 \qedstep
 \begin{proof}\pf\ The two cases cover $\Zu\in T\cup E$.~\qed\end{proof}
\end{proof}
\step{t312completeSpan}{$\opk{span}_{\re}T'=V$.}
\begin{proof}
 \step{t312completeW}{$W\sbs\opk{span}_{\re}T'$.}
 \begin{proof}\pf\ \path{Submodule.span_mono} and $T\sbs T'$.~\qed\end{proof}
 \step{t312completePerp}{$W^\perp\sbs\opk{span}_{\re}T'$.}
 \begin{proof}\pf\ Since $\sqrt2\ne0$, each basis vector is
 $\Ze_i=(1/\sqrt2)\Zu_i\in\opk{span}_{\re}T'$.
 The basis in \stepref{t312completeBasis} spans $W^\perp$.~\qed\end{proof}
 \qedstep
 \begin{proof}\pf\ The orthogonal decomposition $W+W^\perp=V$,
 with \stepref{t312completeW} and \stepref{t312completePerp}.~\qed\end{proof}
\end{proof}
\qedstep
\begin{proof}\pf\ The set $E$ is a finite image of $I$, so $T'$ is
finite and contains $T$. The lengths follow from \asm[2]{0} and
\stepref{t312completeNorm}; the other properties are
\stepref{t312completePairs}, \stepref{t312completeEven}, and
\stepref{t312completeSpan}. The construction allows an empty basis.~\qed\end{proof}
\end{proof}

\begin{namedresult}{Lemma}{(seeds from a colouring)}\label{lem:312-seeds}
\lean{LineSystem.exists_integral_seeds_of_hasBipartition}
\uses{def:312-angles,def:312-root-lines,lem:312-representatives,lem:312-finite,facts:312-inner,facts:312-linear,facts:312-finite,facts:312-arithmetic}
\assume{\begin{pfenum}
 \item $V$ is a finite-dimensional real inner product space,
 and $S$ is a set of subspaces of dimension one.
 \item $\opk{HasBipartition}(S)$.
\end{pfenum}}
\pflet{$W=\opk{span}_{\re}(\bigcup_{\ell\in S}\ell)$.}
\prove{There is a finite set $T\sbs V$ such that
\begin{pfenum}
 \item $\cL(T)=S$ and $\opk{span}_{\re}T=W$;
 \item $q(\Zt)\in\{2,4\}$ for every $\Zt\in T$;
 \item $(\Zt,\Zu)\in\ints$ for every $\Zt,\Zu\in T$; and
 \item $(\Zt,\Zu)\in2\ints$ whenever $\Zt,\Zu\in T$ and $q(\Zt)=4$.
\end{pfenum}}
\end{namedresult}
\begin{proof}
\pfsketch\ We will give the two colour classes squared lengths 2 and 4
and check each ordered pair of colours.
\step{t312seedsColour}{\pick\ a compatible colouring $b:S\to\mathrm{Bool}$.}
\begin{proof}\pf\ Hypothesis \asm[2]{0} and Definition~\hyperref[def:312-angles]{(angles and compatible colourings)}.~\qed\end{proof}
\step{t312seedsAllowed}{$\opk{HasAllowedAngles}(S)$ and $S$ is finite.}
\begin{proof}
 \step{t312seedsEqualColour}{\assume{$\ell,m\in S$, $\ell\ne m$ and $b(\ell)=b(m)$.}
 \prove{There is $a\in\{0,1/4,1/2\}$ with
 $\opk{HasSquaredCosine}(\ell,m,a)$.}}
 \begin{proof}\pf\ The equal-colour clause of \stepref{t312seedsColour}
 gives a parameter in $\{0,1/4\}$ with the required predicate.~\qed\end{proof}
 \step{t312seedsDifferentColour}{\assume{$\ell,m\in S$, $\ell\ne m$ and $b(\ell)\ne b(m)$.}
 \prove{There is $a\in\{0,1/4,1/2\}$ with
 $\opk{HasSquaredCosine}(\ell,m,a)$.}}
 \begin{proof}\pf\ The different-colour clause of \stepref{t312seedsColour}
 gives a parameter in $\{0,1/2\}$ with the required predicate.~\qed\end{proof}
 \qedstep
 \begin{proof}\pf\ Equality or inequality of the colours exhausts all
 pairs. Definition~\hyperref[def:312-angles]{(angles and compatible colourings)} gives allowed angles, and
 Lemma~\hyperref[lem:312-finite]{(finitely many lines)} gives finiteness.~\qed\end{proof}
\end{proof}
\step{t312seedsChoose}{\pick\ $\Zt_\ell\in\ell$ for every $\ell\in S$, with
$q(\Zt_\ell)=2$ when $b(\ell)=\mathrm{false}$ and
$q(\Zt_\ell)=4$ when $b(\ell)=\mathrm{true}$.}
\begin{proof}\pf\ Choose a nonzero generator by the dimension-one facts,
then apply Lemma~\hyperref[lem:312-representatives]{(representatives of a line)} with $h=2$ or $h=4$.
Choice on the empty set gives the empty family.~\qed\end{proof}
\pflet{$T=\{\Zt_\ell:\ell\in S\}$.}
\step{t312seedsLines}{$T$ is finite, $\cL(T)=S$, and
$\opk{span}_{\re}T=W$.}
\begin{proof}
 \step{t312seedsFinite}{$T$ is finite.}
 \begin{proof}\pf\ It is an image of the finite set in \stepref{t312seedsAllowed}.~\qed\end{proof}
 \step{t312seedsLineEq}{$\re\Zt_\ell=\ell$ for every $\ell\in S$.}
 \begin{proof}\pf\ Step \stepref{t312seedsChoose} gives a nonzero member;
 use the singleton-span facts in \hyperref[facts:312-linear]{Spans and finite-dimensional linear algebra}.~\qed\end{proof}
 \step{t312seedsSpanForward}{$\opk{span}_{\re}T\sbs W$.}
 \begin{proof}\pf\ Every seed belongs to its line in $S$ by
 \stepref{t312seedsChoose}; use the definition of $W$ and \path{Submodule.span_mono}.~\qed\end{proof}
 \step{t312seedsSpanReverse}{$W\sbs\opk{span}_{\re}T$.}
 \begin{proof}\pf\ Step \stepref{t312seedsLineEq} puts every line in
 $S$ in the span of $T$, by \path{Submodule.subset_span}
 and \path{Submodule.span_le}. Apply \path{Submodule.span_le}
 again to their union.~\qed\end{proof}
 \qedstep
 \begin{proof}\pf\ Definition~\hyperref[def:312-root-lines]{(lines of a vector set)} and
 \stepref{t312seedsLineEq} give $\cL(T)=S$. The other assertions are
 \stepref{t312seedsFinite}, \stepref{t312seedsSpanForward}, and
 \stepref{t312seedsSpanReverse}.~\qed\end{proof}
\end{proof}
\step{t312seedsPairs}{For $\ell,m\in S$, $(\Zt_\ell,\Zt_m)$ is an integer,
and is even if $q(\Zt_\ell)=4$.}
\begin{proof}
 \pflet{$p=(\Zt_\ell,\Zt_m)$.}
 \step{t312seedsDiagonal}{\case{$\ell=m$.}}
 \begin{proof}\pf\ Step \stepref{t312seedsChoose} gives $p=2$ or $p=4$;
 both are even integers.~\qed\end{proof}
 \step{t312seedsFF}{\case{$\ell\ne m$ and $b(\ell)=b(m)=\mathrm{false}$.}}
 \begin{proof}
  \step{t312seedsFFSquare}{$p^2\in\{0,1\}$.}
  \begin{proof}\pf\ The norms in \stepref{t312seedsChoose} have product
  4; substitute them in the equal-colour clause of \stepref{t312seedsColour}.~\qed\end{proof}
  \qedstep
  \begin{proof}\pf\ Factoring the two possible square equations in
  \stepref{t312seedsFFSquare} gives $p\in\{0,1,-1\}$.
  The evenness implication has false hypothesis, since $q(\Zt_\ell)=2$.~\qed\end{proof}
 \end{proof}
 \step{t312seedsTT}{\case{$\ell\ne m$ and $b(\ell)=b(m)=\mathrm{true}$.}}
 \begin{proof}
  \step{t312seedsTTSquare}{$p^2\in\{0,4\}$.}
  \begin{proof}\pf\ The norms in \stepref{t312seedsChoose} have product
  16; use the equal-colour clause of \stepref{t312seedsColour}.~\qed\end{proof}
  \qedstep
  \begin{proof}\pf\ Factoring the square equations in
  \stepref{t312seedsTTSquare} gives $p\in\{0,2,-2\}$, all even integers.~\qed\end{proof}
 \end{proof}
 \step{t312seedsFT}{\case{$\ell\ne m$, $b(\ell)=\mathrm{false}$ and $b(m)=\mathrm{true}$.}}
 \begin{proof}
  \step{t312seedsFTSquare}{$p^2\in\{0,4\}$.}
  \begin{proof}\pf\ The norms have product 8 by \stepref{t312seedsChoose};
  use the different-colour clause of \stepref{t312seedsColour}.~\qed\end{proof}
  \qedstep
  \begin{proof}\pf\ Factoring \stepref{t312seedsFTSquare} gives
  $p\in\{0,2,-2\}$, all even integers.~\qed\end{proof}
 \end{proof}
 \step{t312seedsTF}{\case{$\ell\ne m$, $b(\ell)=\mathrm{true}$ and $b(m)=\mathrm{false}$.}}
 \begin{proof}
  \step{t312seedsTFSquare}{$p^2\in\{0,4\}$.}
  \begin{proof}\pf\ The norms have product 8 by \stepref{t312seedsChoose};
  use the different-colour clause of \stepref{t312seedsColour}.~\qed\end{proof}
  \qedstep
  \begin{proof}\pf\ Factoring \stepref{t312seedsTFSquare} gives
  $p\in\{0,2,-2\}$, all even integers.~\qed\end{proof}
 \end{proof}
 \qedstep
 \begin{proof}\pf\ The diagonal and the four ordered Boolean pairs
 exhaust the possibilities, as in \stepref{t312seedsDiagonal},
 \stepref{t312seedsFF}, \stepref{t312seedsTT}, \stepref{t312seedsFT},
 and \stepref{t312seedsTF}.~\qed\end{proof}
\end{proof}
\qedstep
\begin{proof}\pf\ Step \stepref{t312seedsLines} supplies part 1 and
finiteness; \stepref{t312seedsChoose} supplies part 2; and
\stepref{t312seedsPairs} supplies parts 3 and 4 on the image $T$.~\qed\end{proof}
\end{proof}

\begin{namedresult}{Corollary}{(normalized containment)}\label{cor:312-containment}
\lean{LineSystem.exists_rootSet_of_hasBipartition}
\lean{LineSystem.exists_rootSet_in_span_of_hasBipartition}
\uses{def:3.9,def:312-angles,def:312-root-lines,def:312-integral-roots,lem:312-seeds,lem:312-completion,prop:312-rootset,facts:312-linear,facts:312-inner,facts:312-finite}
\assume{\begin{pfenum}
 \item $V$ is a finite-dimensional real inner product space,
 and $S$ is a set of subspaces of dimension one.
 \item $\opk{HasBipartition}(S)$.
\end{pfenum}}
\pflet{$W=\opk{span}_{\re}(\bigcup_{\ell\in S}\ell)$,
and let $\iota:W\to V$ be the inclusion.}
\prove{\begin{pfenum}
 \item There is $\Phi\sbs V$ with $\opk{IsRootSet}(\Phi)$,
 $S\sbs\cL(\Phi)$, and $q(\Zr)\in\{2,4\}$ for $\Zr\in\Phi$.
 \item There is $\Phi_W\sbs W$ with $\opk{IsRootSet}(\Phi_W)$
 in the induced inner product on $W$, $S\sbs\cL(\iota(\Phi_W))$,
 and $q(\Zr)\in\{2,4\}$ for $\Zr\in\Phi_W$.
\end{pfenum}}
\end{namedresult}
\begin{proof}
\pfsketch\ We will apply the construction in $V$ after completing
seeds. For the span assertion we will apply the same proposition to the
seeds regarded as vectors in the separate ambient type $W$.
\step{t312containSeeds}{\pick\ a finite set $T$ such that
$\cL(T)=S$, $\opk{span}_{\re}T=W$, each squared length in
$T$ is 2 or 4, all inner products in $T$ are integers, and each
length-4 vector in $T$ pairs evenly with every vector in $T$.}
\begin{proof}\pf\ Lemma~\hyperref[lem:312-seeds]{(seeds from a colouring)} applies to
\asm[1]{0} and \asm[2]{0}.~\qed\end{proof}
\step{t312containAmbient}{There is a set $\Phi$ satisfying part 1.}
\begin{proof}
 \step{t312containExtend}{\pick\ a finite $T'\sps T$ spanning $V$,
 with every squared length 2 or 4, all inner products integral,
 and every length-4 vector pairing evenly with all of $T'$.}
 \begin{proof}\pf\ Apply Lemma~\hyperref[lem:312-completion]{(completion of the seeds)} to the
 conditions in \stepref{t312containSeeds} and the finite dimension
 in \asm[1]{0}.~\qed\end{proof}
 \step{t312containRoot}{$\Phi=\Psi(T')$ is a Root System in $V$,
 contains $T'$, and has squared lengths 2 or 4.}
 \begin{proof}\pf\ Proposition~\hyperref[prop:312-rootset]{(the integer-span Root System)} with
 \stepref{t312containExtend}.~\qed\end{proof}
 \qedstep
 \begin{proof}\pf\ Definition~\hyperref[def:312-root-lines]{(lines of a vector set)} sends
 $T\sbs T'\sbs\Phi$ to $S=\cL(T)\sbs\cL(\Phi)$, by
 \stepref{t312containSeeds} and \stepref{t312containRoot}.~\qed\end{proof}
\end{proof}
\step{t312containSpan}{There is a set $\Phi_W$ satisfying part 2.}
\begin{proof}
 \pflet{$T_W=\{\Zw\in W:\iota\Zw\in T\}$.}
 \step{t312containLiftFinite}{$T_W$ is finite and $\iota(T_W)=T$.}
 \begin{proof}\pf\ The inclusion is injective, and every member of
 $T$ lies in $W$ by \stepref{t312containSeeds}; it restricts to a
 bijection $T_W\to T$. Finiteness follows from
 \path{Set.Finite.preimage}, applied to the finite $T$ and the
 injectivity of $\iota$ on $T_W$.~\qed\end{proof}
 \step{t312containLiftSpan}{$\opk{span}_{\re}T_W=W$ as a subspace of $W$.}
 \begin{proof}\pf\ Rewrite $W=\opk{span}_{\re}T$ using
 \stepref{t312containSeeds}. The assertion is
 \path{Submodule.span_span_coe_preimage} for the inclusion $\iota$.~\qed\end{proof}
 \step{t312containLiftData}{Every squared length in $T_W$ is 2 or 4,
 all its inner products are integers, and each length-4 vector pairs
 evenly with all of $T_W$.}
 \begin{proof}\pf\ The induced inner product on $W$ is the ambient
 inner product; substitute $\iota$ in the three conditions for $T$
 in \stepref{t312containSeeds}.~\qed\end{proof}
 \step{t312containLiftRoot}{$\Phi_W=\Psi(T_W)$ is a Root System in $W$,
 contains $T_W$, and has squared lengths 2 or 4.}
 \begin{proof}\pf\ Apply Proposition~\hyperref[prop:312-rootset]{(the integer-span Root System)} in the
 finite-dimensional inner product space $W$, with
 \stepref{t312containLiftFinite}, \stepref{t312containLiftSpan}, and
 \stepref{t312containLiftData}.~\qed\end{proof}
 \qedstep
 \begin{proof}\pf\ By \stepref{t312containLiftFinite} and
 \stepref{t312containLiftRoot}, $T\sbs\iota(\Phi_W)$.
 Definition~\hyperref[def:312-root-lines]{(lines of a vector set)} and \stepref{t312containSeeds}
 give $S\sbs\cL(\iota(\Phi_W))$.~\qed\end{proof}
\end{proof}
\qedstep
\begin{proof}\pf\ Steps \stepref{t312containAmbient} and \stepref{t312containSpan}.~\qed\end{proof}
\end{proof}

\subsubsection{Lengths and the necessary colouring}

\begin{namedresult}{Lemma}{(lengths on root lines)}\label{lem:312-lengths}
\lean{IsRootSet.sqNorm_eq_of_four_mul_inner_sq_eq}
\lean{IsRootSet.sqNorm_eq_two_mul_or_of_two_mul_inner_sq_eq}
\uses{def:3.9,facts:312-inner,lem:angles,facts:rem-angles}
\assume{\begin{pfenum}
 \item $V$ is a finite-dimensional real inner product space and
 $\opk{IsRootSet}(\Phi)$.
 \item $\Zr,\Zs\in\Phi$.
\end{pfenum}}
\prove{\begin{pfenum}
 \item If $4(\Zr,\Zs)^2=q(\Zr)q(\Zs)$, then $q(\Zr)=q(\Zs)$.
 \item If $2(\Zr,\Zs)^2=q(\Zr)q(\Zs)$, then
 $q(\Zs)=2q(\Zr)$ or $q(\Zr)=2q(\Zs)$.
\end{pfenum}}
\end{namedresult}
\begin{proof}
\pfsketch\ We will apply the pairwise conclusion of
Lemma~\hyperref[lem:angles]{(angles and fixed-length subsystems)}.
Only that conclusion is used; the fixed-length classification in the same
lemma is not needed. We retain the division-free interfaces above.
\step{t312lengthPositive}{$q(\Zr),q(\Zs)>0$.}
\begin{proof}\pf\ Nonzero roots in Definition~\hyperref[def:3.9]{3.9}, \asm[2]{0},
and positive definiteness in
\hyperref[facts:312-inner]{Inner products and square roots}.~\qed\end{proof}
\step{t312lengthPair}{\assume{$k\in\{2,4\}$ and
$k(\Zr,\Zs)^2=q(\Zr)q(\Zs)$.}
\prove{$\Zr,\Zs$ are neither perpendicular nor proportional.}}
\begin{proof}
 \step{t312lengthNonzero}{$(\Zr,\Zs)\ne0$.}
 \begin{proof}\pf\ The right side of \asm{1} is positive by
 \stepref{t312lengthPositive}; its left side cannot vanish.~\qed\end{proof}
 \step{t312lengthNotProportional}{\assume{$\Zs=t\Zr$ for a real $t$.}
 \prove{False.}}
 \begin{proof}
  \step{t312lengthProportionalIdentity}{$(\Zr,\Zs)^2=q(\Zr)q(\Zs)$.}
  \begin{proof}\pf\ Bilinearity in
  \hyperref[facts:312-inner]{Inner products and square roots} gives
  $(\Zr,\Zs)=tq(\Zr)$ and $q(\Zs)=t^2q(\Zr)$ from \asm{2}.~\qed\end{proof}
  \qedstep
  \begin{proof}\pf\ Substitute \stepref{t312lengthProportionalIdentity}
  into \asm{1} and cancel $q(\Zr)q(\Zs)>0$ from
  \stepref{t312lengthPositive}. This gives $k=1$, contrary to
  $k\in\{2,4\}$ in \asm{1}.~\qed\end{proof}
 \end{proof}
 \qedstep
 \begin{proof}\pf\ Steps \stepref{t312lengthNonzero} and
 \stepref{t312lengthNotProportional} give the two exclusions.~\qed\end{proof}
\end{proof}
\step{t312lengthQuarter}{\assume{$4(\Zr,\Zs)^2=q(\Zr)q(\Zs)$.}
 \prove{$q(\Zr)=q(\Zs)$.}}
\begin{proof}
 \step{t312lengthQuarterAngle}{The acute angle between their lines is $60^\circ$.}
 \begin{proof}\pf\ Divide \asm{1} by the positive product in
 \stepref{t312lengthPositive}. The squared cosine is $1/4$;
 use the \hyperref[facts:rem-angles]{Angle convention}.~\qed\end{proof}
 \step{t312lengthQuarterRatio}{The two roots have equal length.}
 \begin{proof}\pf\ Apply Lemma~\hyperref[lem:angles]{(angles and fixed-length subsystems)}
 to \asm[1]{0}, \asm[2]{0}, and \stepref{t312lengthPair} with $k=4$
 and \asm{1}. Step \stepref{t312lengthQuarterAngle} selects the
 length ratio $1$.~\qed\end{proof}
 \qedstep
 \begin{proof}\pf\ Square the equality of lengths in
 \stepref{t312lengthQuarterRatio}, using
 \hyperref[facts:312-inner]{Inner products and square roots}.~\qed\end{proof}
\end{proof}
\step{t312lengthHalf}{\assume{$2(\Zr,\Zs)^2=q(\Zr)q(\Zs)$.}
 \prove{$q(\Zs)=2q(\Zr)$ or $q(\Zr)=2q(\Zs)$.}}
\begin{proof}
 \step{t312lengthHalfAngle}{The acute angle between their lines is $45^\circ$.}
 \begin{proof}\pf\ Divide \asm{1} by the positive product in
 \stepref{t312lengthPositive}. The squared cosine is $1/2$;
 use the \hyperref[facts:rem-angles]{Angle convention}.~\qed\end{proof}
 \step{t312lengthHalfRatio}{The longer-to-shorter length ratio is $\sqrt2$.}
 \begin{proof}\pf\ Apply Lemma~\hyperref[lem:angles]{(angles and fixed-length subsystems)}
 to \asm[1]{0}, \asm[2]{0}, and \stepref{t312lengthPair} with $k=2$
 and \asm{1}. Step \stepref{t312lengthHalfAngle} selects the
 length ratio $\sqrt2$.~\qed\end{proof}
 \qedstep
 \begin{proof}\pf\ Square the ratio in \stepref{t312lengthHalfRatio}.
 The identity $(\sqrt2)^2=2$ in
 \hyperref[facts:312-inner]{Inner products and square roots} gives
 the two stated alternatives according to which root is longer.~\qed\end{proof}
\end{proof}
\qedstep
\begin{proof}\pf\ Steps \stepref{t312lengthQuarter} and
\stepref{t312lengthHalf} give the two conclusions.~\qed\end{proof}
\end{proof}

\begin{namedresult}{Lemma}{(integer logarithm shifts)}\label{lem:312-log}
\lean{Int.log_two_mul}
\lean{Int.log_two_div}
\uses{facts:312-arithmetic}
\assume{$x\in\re$ and $x>0$.}
\pflet{$\eta(y)=\opk{Int.log}(2,y)$ for real $y>0$.}
\prove{$\eta(2x)=\eta(x)+1$ and $\eta(x/2)=\eta(x)-1$.}
\end{namedresult}
\begin{proof}
\pfsketch\ We will multiply the defining interval by 2, then apply
the doubling identity at $x/2$.
\step{t312logDouble}{For every real $y>0$, $\eta(2y)=\eta(y)+1$.}
\begin{proof}
 \step{t312logBounds}{$2^{\eta(y)}\leq y<2^{\eta(y)+1}$.}
 \begin{proof}\pf\ \path{Int.zpow_log_le_self} and
 \path{Int.lt_zpow_succ_log_self} with base $2>1$ and $y>0$
 in \stepref{t312logDouble}.~\qed\end{proof}
 \step{t312logDoubleBounds}{$2^{\eta(y)+1}\leq2y<2^{(\eta(y)+1)+1}$.}
 \begin{proof}\pf\ Multiply \stepref{t312logBounds} by $2>0$ and
 apply \path{zpow_add_one}\textsubscript{0} at base $2\ne0$.~\qed\end{proof}
 \step{t312logDoubleLower}{$\eta(y)+1\leq\eta(2y)$.}
 \begin{proof}\pf\ Apply \path{Int.zpow_le_iff_le_log} at $2y>0$
 to the lower bound of \stepref{t312logDoubleBounds}.~\qed\end{proof}
 \step{t312logDoubleUpper}{$\eta(2y)<\eta(y)+2$.}
 \begin{proof}\pf\ Apply \path{Int.lt_zpow_iff_log_lt} at $2y>0$
 to the upper bound of \stepref{t312logDoubleBounds}.~\qed\end{proof}
 \qedstep
 \begin{proof}\pf\ The adjacent integer bounds in
 \stepref{t312logDoubleLower} and \stepref{t312logDoubleUpper}.~\qed\end{proof}
\end{proof}
\step{t312logHalf}{$\eta(x/2)=\eta(x)-1$.}
\begin{proof}\pf\ Apply \stepref{t312logDouble} to $x/2>0$, using
\asm{0}: $\eta(x)=\eta(2(x/2))=\eta(x/2)+1$.
Rearrange this integer equality.~\qed\end{proof}
\qedstep
\begin{proof}\pf\ Step \stepref{t312logDouble} at $x>0$ and
\stepref{t312logHalf}.~\qed\end{proof}
\end{proof}

The two logarithm lemmas are declarations in
namespace \path{Int}; they concern positive real numbers, independently
of line systems.

\begin{namedresult}{Lemma}{(colouring selected root lines)}\label{lem:312-necessary}
\lean{LineSystem.hasBipartition_of_subset_rootLines}
\uses{def:3.9,def:312-angles,def:312-root-lines,lem:312-squared-cosine,lem:312-lengths,lem:312-log,facts:312-inner,facts:312-linear,facts:312-arithmetic}
\assume{\begin{pfenum}
 \item $V$ is a finite-dimensional real inner product space,
 $\opk{IsRootSet}(\Phi)$, and $S\sbs\cL(\Phi)$.
 \item $\opk{HasAllowedAngles}(S)$.
\end{pfenum}}
\prove{$\opk{HasBipartition}(S)$.}
\end{namedresult}
\begin{proof}
\pfsketch\ We will colour a line by the parity of the integer logarithm
of its squared root length. Orthogonal pairs impose no constraint.
\step{t312necessaryChoose}{\pick\ $\Zr_\ell\in\Phi$ with $\ell=\re\Zr_\ell$
for every $\ell\in S$.}
\begin{proof}\pf\ Containment in \asm[1]{0} and
Definition~\hyperref[def:312-root-lines]{(lines of a vector set)} give the witnesses.~\qed\end{proof}
\pflet{$\delta(\ell)=q(\Zr_\ell)$, $\eta(x)=\opk{Int.log}(2,x)$,
and $b(\ell)=\mathrm{true}$ exactly when $\eta(\delta(\ell))$ is odd.}
\step{t312necessaryLines}{Every $\ell\in S$ has dimension one,
and $\delta(\ell)>0$.}
\begin{proof}
 \step{t312necessaryPositive}{$\Zr_\ell\ne0$ and $\delta(\ell)>0$.}
 \begin{proof}\pf\ Definition~\hyperref[def:3.9]{3.9}, \stepref{t312necessaryChoose},
 and positive definiteness.~\qed\end{proof}
 \step{t312necessaryDimension}{$\dim\ell=1$.}
 \begin{proof}\pf\ Apply \path{finrank_span_singleton} to
 \stepref{t312necessaryPositive} and \stepref{t312necessaryChoose}.~\qed\end{proof}
 \qedstep
 \begin{proof}\pf\ Steps \stepref{t312necessaryPositive}
 and \stepref{t312necessaryDimension}.~\qed\end{proof}
\end{proof}
\step{t312necessaryQuarter}{\assume{$\ell,m\in S$, $\ell\ne m$, and $\opk{HasSquaredCosine}(\ell,m,1/4)$.}
 \prove{$b(\ell)=b(m)$.}}
\begin{proof}
 \step{t312necessaryEqualLength}{$\delta(\ell)=\delta(m)$.}
 \begin{proof}\pf\ Lemma~\hyperref[lem:312-squared-cosine]{(squared cosine from representatives)}, applied to
 \asm{1} and the nonzero representatives from \stepref{t312necessaryChoose}
 and \stepref{t312necessaryLines}, gives
 $4(\Zr_\ell,\Zr_m)^2=\delta(\ell)\delta(m)$.
 Apply part 1 of Lemma~\hyperref[lem:312-lengths]{(lengths on root lines)} to these roots and
 the Root System in \asm[1]{0}.~\qed\end{proof}
 \qedstep
 \begin{proof}\pf\ Equal inputs in \stepref{t312necessaryEqualLength}
 give equal integer logarithms and equal Boolean colours.~\qed\end{proof}
\end{proof}
\step{t312necessaryHalf}{\assume{$\ell,m\in S$, $\ell\ne m$, and $\opk{HasSquaredCosine}(\ell,m,1/2)$.}
 \prove{$b(\ell)\ne b(m)$.}}
\begin{proof}
 \step{t312necessaryRatio}{$\delta(m)=2\delta(\ell)$ or $\delta(m)=\delta(\ell)/2$.}
 \begin{proof}\pf\ Lemma~\hyperref[lem:312-squared-cosine]{(squared cosine from representatives)}, applied to
 \asm{1} and the representatives in \stepref{t312necessaryChoose}
 and \stepref{t312necessaryLines}, gives
 $2(\Zr_\ell,\Zr_m)^2=\delta(\ell)\delta(m)$.
 Part 2 of Lemma~\hyperref[lem:312-lengths]{(lengths on root lines)}, using \asm[1]{0},
 gives $\delta(m)=2\delta(\ell)$ or $\delta(\ell)=2\delta(m)$.
 Divide the second equality by 2.~\qed\end{proof}
 \step{t312necessaryDouble}{\case{$\delta(m)=2\delta(\ell)$.}
 \prove{$b(\ell)\ne b(m)$.}}
 \begin{proof}
  \step{t312necessaryPlus}{$\eta(\delta(m))=\eta(\delta(\ell))+1$.}
  \begin{proof}\pf\ The doubling identity in Lemma~\hyperref[lem:312-log]{(integer logarithm shifts)}
  and positivity in \stepref{t312necessaryLines}.~\qed\end{proof}
  \qedstep
  \begin{proof}\pf\ Consecutive integers have opposite parities by
  \hyperref[facts:312-arithmetic]{Integer arithmetic and logarithms}; apply \stepref{t312necessaryPlus} and the definition of $b$.~\qed\end{proof}
 \end{proof}
 \step{t312necessaryHalved}{\case{$\delta(m)=\delta(\ell)/2$.}
 \prove{$b(\ell)\ne b(m)$.}}
 \begin{proof}
  \step{t312necessaryMinus}{$\eta(\delta(m))=\eta(\delta(\ell))-1$.}
  \begin{proof}\pf\ The halving identity in Lemma~\hyperref[lem:312-log]{(integer logarithm shifts)}
  and positivity in \stepref{t312necessaryLines}.~\qed\end{proof}
  \qedstep
  \begin{proof}\pf\ Integers differing by 1 have opposite parities by
  \hyperref[facts:312-arithmetic]{Integer arithmetic and logarithms}; apply \stepref{t312necessaryMinus} and the definition of $b$.~\qed\end{proof}
 \end{proof}
 \qedstep
 \begin{proof}\pf\ Step \stepref{t312necessaryRatio} makes
 \stepref{t312necessaryDouble} and \stepref{t312necessaryHalved} exhaustive.~\qed\end{proof}
\end{proof}
\step{t312necessaryEqual}{\assume{$\ell,m\in S$, $\ell\ne m$, and $b(\ell)=b(m)$.}
 \prove{$\opk{HasSquaredCosine}(\ell,m,0)$ or
 $\opk{HasSquaredCosine}(\ell,m,1/4)$.}}
\begin{proof}\pf\ Hypothesis \asm[2]{0} and Definition~\hyperref[def:312-angles]{(angles and compatible colourings)}
give a parameter in $\{0,1/4,1/2\}$ satisfying the predicate.
Step \stepref{t312necessaryHalf} and the colour equality in \asm{1}
exclude the predicate at $1/2$. The other two parameters give the
stated disjunction.~\qed\end{proof}
\step{t312necessaryDifferent}{\assume{$\ell,m\in S$, $\ell\ne m$, and $b(\ell)\ne b(m)$.}
 \prove{$\opk{HasSquaredCosine}(\ell,m,0)$ or
 $\opk{HasSquaredCosine}(\ell,m,1/2)$.}}
\begin{proof}\pf\ Hypothesis \asm[2]{0} and Definition~\hyperref[def:312-angles]{(angles and compatible colourings)}
give a parameter in $\{0,1/4,1/2\}$ satisfying the predicate.
Step \stepref{t312necessaryQuarter} and the colour inequality in \asm{1}
exclude the predicate at $1/4$. The other two parameters give the
stated disjunction.~\qed\end{proof}
\qedstep
\begin{proof}\pf\ Steps \stepref{t312necessaryEqual} and
\stepref{t312necessaryDifferent} give the two clauses of
Definition~\hyperref[def:312-angles]{(angles and compatible colourings)}; $b$ is the required witness. For $S=\varnothing$ it is the empty function.~\qed\end{proof}
\end{proof}

\begin{cgss}{Theorem}{3.12}\label{thm:3.12}
\lean{LineSystem.subset_rootSet_iff_bipartition}
\uses{def:3.9,def:312-angles,def:312-root-lines,cor:312-containment,lem:312-necessary}
\assume{\begin{pfenum}
 \item $V$ is a finite-dimensional real inner product space.
 The set $S$ consists of real subspaces of $V$, with
 $\dim\ell=1$ for every $\ell\in S$.
 \item $\opk{HasAllowedAngles}(S)$.
\end{pfenum}}
\prove{\[
 \bigl(\exists\Phi\sbs V,\quad
 \opk{IsRootSet}(\Phi)\ \text{and}\ S\sbs\cL(\Phi)\bigr)
 \quad\Longleftrightarrow\quad \opk{HasBipartition}(S).
\]}
\end{cgss}
\begin{proof}
\pfsketch\ \repair\ The paper gives no proof (Gap 11).
We will use the normalized construction for sufficiency and the
integer-logarithm colouring for necessity.
\step{t312mainNecessary}{\assume{There is a Root System $\Phi$ in $V$ with $S\sbs\cL(\Phi)$.}
 \prove{$\opk{HasBipartition}(S)$.}}
\begin{proof}\pf\ Lemma~\hyperref[lem:312-necessary]{(colouring selected root lines)} with the finite
dimension in \asm[1]{0}, the containment and Root System in \asm{1},
and the angle hypothesis \asm[2]{0}.~\qed\end{proof}
\step{t312mainSufficient}{\assume{$\opk{HasBipartition}(S)$.}
 \prove{There is a Root System $\Phi$ in $V$ with $S\sbs\cL(\Phi)$.}}
\begin{proof}\pf\ Part 1 of Corollary~\hyperref[cor:312-containment]{(normalized containment)}, with
\asm[1]{0} and \asm{1}.~\qed\end{proof}
\qedstep
\begin{proof}\pf\ Steps \stepref{t312mainNecessary} and
\stepref{t312mainSufficient} are the two implications.~\qed\end{proof}
\end{proof}

\paragraph{Conventions in the statement.}
\repair\ A bipartition means the two fibres of a Boolean colouring;
either fibre may be empty. This is required even for a single line,
or for two lines at $60^\circ$, which must have the same colour.
The ordinary proof records this as Implicit hypothesis 8 of
\path{cgss/errata.md}. The word \emph{and} between the two cross-part
angles in the printed statement is read as \emph{or} (Error 13, item 3).
Finiteness is a conclusion of Lemma~\hyperref[lem:312-finite]{(finitely many lines)};
spanning and nonemptiness of $S$ are not hypotheses.
Corollary~\hyperref[cor:312-containment]{(normalized containment)} separately states the stronger
normalized construction, in both the span and the given ambient space.
Allowing a containing Root System in a larger finite-dimensional inner
product space gives the same condition: apply
Lemma~\hyperref[lem:312-necessary]{(colouring selected root lines)} in that space to the isometric images
of the selected lines and pull back its colouring.

The Lean declaration has argument \path{S : Set (Submodule Real V)},
not an argument of the bundled type \path{LineSystem V}, which permits
only the $60^\circ$ and $90^\circ$ cases. Its ambient assumptions are
\path{NormedAddCommGroup V}, \path{InnerProductSpace Real V},
and \path{FiniteDimensional Real V}; the line hypothesis is
\path{Module.finrank Real} equal to 1 on each member of $S$.
The docstring records the empty-fibre convention and the
errata references above. The names in the metadata are the settled
ones from \path{cgss/notation.md}, Section 10, supplemented by the
separately stated helper lemmas.

\subsubsection{The bridge from Section 2}

\begin{namedresult}{Lemma}{(the lines of a line system)}\label{lem:312-bridge}
\lean{LineSystem.hasAllowedAngles_rootLines}
\lean{LineSystem.hasBipartition_rootLines}
\uses{def:line-system,def:312-root-lines,def:312-angles,lem:312-squared-cosine,facts:312-inner,facts:312-linear}
\assume{$V$ is a finite-dimensional real inner product space and
$R:\opk{LineSystem}(V)$ is as in Definition~\hyperref[def:line-system]{(line systems of Section 2)}.}
\prove{$\opk{HasAllowedAngles}(\cL(R.\Phi))$ and
$\opk{HasBipartition}(\cL(R.\Phi))$, witnessed by the constant
colouring with value $\mathrm{false}$.}
\end{namedresult}
\begin{proof}
\pfsketch\ We will exclude the two endpoint inner products on distinct
lines and then use the constant colouring.
\step{t312bridgeDimension}{Every $\ell\in\cL(R.\Phi)$ has dimension one.}
\begin{proof}\pf\ A representative has squared norm 2 by
Definition~\hyperref[def:line-system]{(line systems of Section 2)}, so is nonzero; apply
\path{finrank_span_singleton} and Definition~\hyperref[def:312-root-lines]{(lines of a vector set)}.~\qed\end{proof}
\step{t312bridgeAngles}{\assume{$\ell,m\in\cL(R.\Phi)$ and $\ell\ne m$.}
 \prove{$\opk{HasSquaredCosine}(\ell,m,0)$ or
 $\opk{HasSquaredCosine}(\ell,m,1/4)$.}}
\begin{proof}
 \step{t312bridgePick}{\pick\ $\Zr,\Zs\in R.\Phi$ with $\ell=\re\Zr$, $m=\re\Zs$.}
 \begin{proof}\pf\ Definition~\hyperref[def:312-root-lines]{(lines of a vector set)}.~\qed\end{proof}
 \step{t312bridgePositive}{\assume{$(\Zr,\Zs)=2$.}\prove{False.}}
 \begin{proof}
  \step{t312bridgeDifference}{$q(\Zr-\Zs)=0$ and $\Zr=\Zs$.}
  \begin{proof}\pf\ Bilinearity and the squared norms in
  Definition~\hyperref[def:line-system]{(line systems of Section 2)} give $q(\Zr-\Zs)=2+2-2\cdot2=0$;
  positive definiteness gives the equality.~\qed\end{proof}
  \qedstep
  \begin{proof}\pf\ Step \stepref{t312bridgeDifference} makes the spans
  in \stepref{t312bridgePick} equal, contrary to \asm{1}.~\qed\end{proof}
 \end{proof}
 \step{t312bridgeNegative}{\assume{$(\Zr,\Zs)=-2$.}\prove{False.}}
 \begin{proof}
  \step{t312bridgeSum}{$q(\Zr+\Zs)=0$ and $\Zr=-\Zs$.}
  \begin{proof}\pf\ Bilinearity gives $q(\Zr+\Zs)=2+2+2(-2)=0$;
  positive definiteness gives the equality.~\qed\end{proof}
  \qedstep
  \begin{proof}\pf\ Multiplication by $-1$ does not change a real span;
  \stepref{t312bridgeSum} contradicts the distinct spans in \asm{1}.~\qed\end{proof}
 \end{proof}
 \step{t312bridgeSmall}{$(\Zr,\Zs)\in\{-1,0,1\}$.}
 \begin{proof}\pf\ Exclude the two endpoint values in
 Definition~\hyperref[def:line-system]{(line systems of Section 2)} by \stepref{t312bridgePositive} and
 \stepref{t312bridgeNegative}.~\qed\end{proof}
 \qedstep
 \begin{proof}\pf\ Step \stepref{t312bridgeSmall} gives
 $(\Zr,\Zs)^2\in\{0,1\}$. The squared norms are 2 by
 Definition~\hyperref[def:line-system]{(line systems of Section 2)}. Lemma~\hyperref[lem:312-squared-cosine]{(squared cosine from representatives)},
 using \stepref{t312bridgeDimension} and \stepref{t312bridgePick},
 gives the predicate at $0$ or $1/4$.~\qed\end{proof}
\end{proof}
\step{t312bridgeAllowed}{$\opk{HasAllowedAngles}(\cL(R.\Phi))$.}
\begin{proof}\pf\ Step \stepref{t312bridgeAngles} supplies an allowed
parameter in Definition~\hyperref[def:312-angles]{(angles and compatible colourings)}.~\qed\end{proof}
\step{t312bridgeColour}{The constant function $b(\ell)=\mathrm{false}$ is compatible.}
\begin{proof}
 \step{t312bridgeSameColour}{For distinct $\ell,m$, the equal-colour angle clause holds.}
 \begin{proof}\pf\ Step \stepref{t312bridgeAngles}.~\qed\end{proof}
 \step{t312bridgeCrossColour}{There are no $\ell,m$ with $b(\ell)\ne b(m)$.}
 \begin{proof}\pf\ The definition of the constant function.~\qed\end{proof}
 \qedstep
 \begin{proof}\pf\ Steps \stepref{t312bridgeSameColour} and
 \stepref{t312bridgeCrossColour} give the two clauses of
 Definition~\hyperref[def:312-angles]{(angles and compatible colourings)}.~\qed\end{proof}
\end{proof}
\qedstep
\begin{proof}\pf\ Steps \stepref{t312bridgeAllowed} and \stepref{t312bridgeColour}.~\qed\end{proof}
\end{proof}

The exact computational examples remain in the ordinary proof and its
\path{blueprint/checks/thm312/check_thm312.py}. They are not premises of any result here.
The construction and colouring need no two-length assertion for a reducible
Root System. The colouring uses only the pairwise angle and length conclusion
of Lemma~\hyperref[lem:angles]{(angles and fixed-length subsystems)};
that conclusion is proved there from integrality, before the application
of classification to fixed-length subsystems.
The proof permits unrelated scales on orthogonal components, and
restricts angles only on the selected set of lines.
\end{mergedproofsection}

\section{Graphs with least eigenvalue \texorpdfstring{$-2$}{-2}}\label{sec:4}
\begingroup
\emergencystretch=3em

We will follow Section 4 of [CGSS]. The paper's graphs $G$ and $H$ are our graphs $X$ and $Y$.
All graphs are finite and simple. For a graph $X$ we write $A=A(X)$ for its adjacency matrix,
$J$ for the all-ones matrix and $\lambda_{\min}(X)$ for its least eigenvalue; the least
eigenvalue of $X$ is at least $-2$ exactly when $2I+A$ is positive semidefinite, and we will
state the results of this section in that form, as the Lean statements do. The coordinates
$\Ze_1,\dots,\Ze_\ell$ are those of Definition~\ref{def:3.1}. For a vector
$\Zx=\sum_i\Zx(i)\Ze_i$ the \emph{support} of $\Zx$ is $\supp(\Zx)=\{i:\Zx(i)\ne0\}$. For a graph $Y$,
the line graph $L(Y)$ has the edges of $Y$ as its vertices, two of them adjacent when they are
distinct and share an end.

\begingroup\raggedright
\medskip\noindent\textbf{Facts from Mathlib and the project.}\label{facts:sec4}
We will use the following declarations of Mathlib (commit \path{082e2d37}):
\path{Matrix.posSemidef_iff_eq_sum_vecMulVec} (a positive semidefinite real matrix is a sum of
matrices $\Zv\Zv^T$), \path{Matrix.posSemidef_gram} (a Gram matrix is positive semidefinite),
\path{Matrix.rank_eq_finrank_span_cols}, \path{Submodule.map_span},
\path{LinearMap.finrank_range_of_inj}, \path{Submodule.finrank_le}, \path{stdOrthonormalBasis}
with its field \path{OrthonormalBasis.repr},
\path{Matrix.det_gram_ne_zero_iff_linearIndependent},
\path{Matrix.posDef_gram_iff_linearIndependent},
\path{LinearIndependent.fintype_card_le_finrank}, \path{finrank_euclideanSpace_fin}, the
definition \path{Matrix.rank} (the dimension of the range), \path{Submodule.span_mono},
\path{Submodule.finrank_mono}, \path{Submodule.finrank_sup_add_finrank_inf_eq},
\path{finrank_span_singleton}, \path{SimpleGraph.adjMatrix_mulVec_const_apply_of_regular}, for a
real symmetric matrix its eigenvalues \path{Matrix.IsHermitian.eigenvalues} with
\path{Matrix.IsHermitian.spectral_theorem}, \path{Matrix.IsHermitian.trace_eq_sum_eigenvalues}
and \path{Matrix.IsHermitian.rank_eq_card_non_zero_eigs}, \path{Matrix.trace_mul_comm}, the
Cauchy--Schwarz inequality \path{Finset.sum_mul_sq_le_sq_mul_sq}, the predicate
\path{SimpleGraph.IsBipartite} (a 2-colouring) with
\path{SimpleGraph.two_colorable_iff_forall_loop_even} and
\path{SimpleGraph.Coloring.even_length_iff_congr}, and the structure
\path{SimpleGraph.IsSRGWith}. From the
project we use the bridge \path{SimpleGraph.posSemidef_two_add_adjMatrix_iff} ($2I+A$ is
positive semidefinite exactly when every eigenvalue of $X$ is at least $-2$) and
\path{SimpleGraph.leastEigenvalue_eq_neg_two_iff} (for a graph with a vertex,
$\lambda_{\min}(X)=-2$ exactly when $2I+A$ is positive semidefinite and singular), both in
\path{LineGraphs/Spectral/Positivity.lean}; and the sets
$N_\Zr=\{\Zx\in\Phi:(\Zx,\Zr)=1\}$ of \path{LineGraphs/Computations/NeighborCertificates.lean}.
\par
\endgroup

\subsection{Theorem 4.1}

\begin{namedresult}{Definition}{(representations)}\label{def:represented}
 \lean{SimpleGraph.IsRepresentation}\lean{SimpleGraph.IsRepresentedIn}
 \uses{def:line-system,def:3.2}
 \assume{$X$ is a graph, $V$ is a real inner product space and $S$ is a line system in $V$.}
 \pflet{A \emph{representation} of $X$ in $V$ is a family $(\Zx_u)_{u\in V(X)}$ of vectors of
 $V$ whose Gram matrix is $2I+A(X)$; that is, $(\Zx_u,\Zx_u)=2$, $(\Zx_u,\Zx_w)=1$ when $u$ and
 $w$ are adjacent, and $(\Zx_u,\Zx_w)=0$ when $u\ne w$ are not adjacent. The graph $X$ is
 \emph{represented in $S$} if it has a representation with $\Zx_u\in\Phi(S)$ for every $u$.}
 \new\ The paper speaks of a set of vectors and of the graph it represents; we will index the
 vectors by the vertices, so that the graph is given first, as in Lean. For $S=\typ E_8$ we take
 model (ii) of Definition~\ref{def:3.2} in $\re^8$.
\end{namedresult}

\begin{namedresult}{Lemma}{(Gram factorization)}\label{lem:4-gramfactor}
 \lean{Matrix.rank_gram_eq_finrank_span}\lean{Matrix.PosSemidef.exists_gram_euclideanSpace}
 \uses{facts:sec4}
 \assume{$M$ is a real symmetric matrix indexed by a finite set $V$, and $r\ge0$.}
 \prove{\begin{pfenum}
  \item If $M$ is the Gram matrix of a family $(\Zx_u)_{u\in V}$ in a real inner product space,
   then $\rk M$ is the dimension of the span of the $\Zx_u$.
  \item If $M$ is positive semidefinite and $\rk M\le r$, then $M$ is the Gram matrix of a
   family in $\re^r$, that is, in \texttt{EuclideanSpace}~$\re$~\texttt{(Fin r)}.
 \end{pfenum}}
 \new\ Mathlib factors a positive semidefinite matrix through a sum of an unspecified number of
 terms; the sections below need the dimension.
\end{namedresult}
\begin{proof}
 \pfsketch\ We will compare the columns of $M$ with the span, and then move the span into
 $\re^r$ by an isometry.
 \step{s4gfspan}{Assertion 1 holds.}
 \begin{proof}\pf\ Let $W$ be the span of the $\Zx_u$ and $T\Zy=((\Zx_w,\Zy))_{w\in V}$ for
  $\Zy\in W$. The column $u$ of $M$ is $T\Zx_u$, so the span of the columns is $T(W)$, and $T$
  is injective on $W$, since $T\Zy=0$ gives $(\Zy,\Zy)=0$. So $\dim T(W)=\dim W$ by
  \path{Submodule.map_span} and \path{LinearMap.finrank_range_of_inj}; apply
  \path{Matrix.rank_eq_finrank_span_cols}.~\qed\end{proof}
 \step{s4gfsum}{\assume{$M$ is positive semidefinite.}\prove{$M$ is the Gram matrix of a family
  $(\Zy_u)$ in $\re^m$ for some $m$.}}
 \begin{proof}\pf\ By \path{Matrix.posSemidef_iff_eq_sum_vecMulVec},
  $M=\sum_{k=1}^m\Zv_k\Zv_k^T$ for some $m$ and vectors $\Zv_k\in\re^V$. Put
  $\Zy_u=\sum_k\Zv_k(u)\Ze_k$; then $(\Zy_u,\Zy_w)=\sum_k\Zv_k(u)\Zv_k(w)$ is the $uw$-entry of
  $M$.~\qed\end{proof}
 \step{s4gfcoord}{\assume{$M$ is positive semidefinite.}\prove{$M$ is the Gram matrix of a
  family in $\re^{\dim W}$, where $W$ is the span of the $\Zy_u$ of \stepref{s4gfsum}.}}
 \begin{proof}\pf\ The coordinates in the orthonormal basis \path{stdOrthonormalBasis} of $W$,
  that is, its \path{OrthonormalBasis.repr}, form a linear isometry of $W$ onto
  $\re^{\dim W}$; they preserve inner products.~\qed\end{proof}
 \step{s4gfpad}{\assume{$k\le r$.}\prove{Padding with $r-k$ zero coordinates maps $\re^k$
  into $\re^r$ and preserves inner products.}}
 \begin{proof}\pf\ The padded coordinates contribute $0$ to each inner
  product.~\qed\end{proof}
 \qedstep
 \begin{proof}\pf\ Assertion 1 is \stepref{s4gfspan}. For assertion 2, $\dim W=\rk M\le r$ by
  \stepref{s4gfspan}; apply \stepref{s4gfcoord} and then \stepref{s4gfpad} with
  $k=\dim W$.~\qed\end{proof}
\end{proof}

\begin{namedresult}{Lemma}{(the Gram matrix)}\label{lem:4-gram}
 \lean{SimpleGraph.exists_representation_iff_posSemidef}
 \uses{def:represented,lem:4-gramfactor,facts:sec4}
 \assume{$X$ is a graph on $n$ vertices.}
 \prove{\begin{pfenum}
  \item $X$ has a representation in $\re^n$ if and only if $2I+A(X)$ is positive semidefinite,
   that is, if and only if every eigenvalue of $X$ is at least $-2$.
  \item If $n\ge1$ and $X$ has a representation in a space of dimension less than $n$, then
   $\lambda_{\min}(X)=-2$.
 \end{pfenum}}
\end{namedresult}
\begin{proof}
 \pfsketch\ We will factor the matrix into a Gram matrix, and then count independent vectors.
 \step{s4gramfwd}{\assume{$(\Zx_u)$ is a representation of $X$ in a real inner product
  space $V$.}
  \prove{$2I+A(X)$ is positive semidefinite.}}
 \begin{proof}\pf\ $2I+A(X)$ is the Gram matrix of $(\Zx_u)$ by \asm{1} and
  Definition~\ref{def:represented}; apply \path{Matrix.posSemidef_gram}.~\qed\end{proof}
 \step{s4gramback}{\assume{$2I+A(X)$ is positive semidefinite.}
  \prove{$X$ has a representation in $\re^n$.}}
 \begin{proof}\pf\ $2I+A(X)$ has rank at most $n$; apply Lemma~\ref{lem:4-gramfactor},
  assertion 2, with $r=n$, and Definition~\ref{def:represented}.~\qed\end{proof}
 \step{s4gramone}{Assertion 1 holds.}
 \begin{proof}\pf\ \stepref{s4gramfwd} and \stepref{s4gramback} give the first equivalence, and
  \path{SimpleGraph.posSemidef_two_add_adjMatrix_iff} the second.~\qed\end{proof}
 \step{s4gramsing}{\assume{$n\ge1$ and $(\Zx_u)$ is a representation of $X$ in a space $V$ of
  dimension less than $n$.}\prove{$\det(2I+A(X))=0$.}}
 \begin{proof}\pf\ The family $(\Zx_u)$ has $n$ members in a space of dimension less than $n$,
  so it is not linearly independent by \path{LinearIndependent.fintype_card_le_finrank}. Its
  Gram matrix $2I+A(X)$ is singular by
  \path{Matrix.det_gram_ne_zero_iff_linearIndependent}.~\qed\end{proof}
 \qedstep
 \begin{proof}\pf\ Assertion 1 is \stepref{s4gramone}. For assertion 2, $2I+A(X)$ is positive
  semidefinite by \stepref{s4gramfwd} and
  singular by \stepref{s4gramsing}; apply
  \path{SimpleGraph.leastEigenvalue_eq_neg_two_iff}.~\qed\end{proof}
\end{proof}

\begin{namedresult}{Lemma}{(the lines of a representation)}\label{lem:4-lines}
 \lean{SimpleGraph.IsRepresentation.isIndecomposable_of_connected}
 \uses{def:represented,def:line-system,def:indecomposable,lem:sec3-connected-span}
 \assume{$(\Zx_u)_{u\in V(X)}$ is a representation of a graph $X$ in a finite-dimensional real
 inner product space $V$, and $\Phi=\{\pm\Zx_u:u\in V(X)\}$.}
 \prove{\begin{pfenum}
  \item $u\mapsto\Zx_u$ is injective, and $\Zx_u\ne-\Zx_w$ for all $u,w$.
  \item $\Phi$ is a line system $S_\Zx$ in $V$, with $|V(X)|$ lines.
  \item If $X$ is connected, then $S_\Zx$ is indecomposable.
 \end{pfenum}}
\end{namedresult}
\begin{proof}
 \pfsketch\ We will compare inner products with the squared norm 2.
 \step{s4linesinj}{Assertion 1 holds.}
 \begin{proof}\pf\ For $u\ne w$, $(\Zx_u,\Zx_w)\in\{0,1\}$ by Definition~\ref{def:represented},
  while $(\Zx_u,\Zx_u)=2$ and $(\Zx_u,-\Zx_u)=-2$; so $\Zx_w\ne\Zx_u$ and $\Zx_w\ne-\Zx_u$. For
  $u=w$, $\Zx_u=-\Zx_u$ would give $\Zx_u=0$, against $(\Zx_u,\Zx_u)=2$.~\qed\end{proof}
 \step{s4linessys}{Assertion 2 holds.}
 \begin{proof}\pf\ $\Phi=-\Phi$, every member has squared norm 2, and the inner products of
  members are $\pm$ those of the $\Zx_u$, which lie in $\{0,1,2\}$; this is
  Definition~\ref{def:line-system}. The lines are the pairs $\{\Zx_u,-\Zx_u\}$, which are
  distinct by \stepref{s4linesinj}.~\qed\end{proof}
 \qedstep
 \begin{proof}\pf\ Take $B=\{\Zx_u\}$ in Lemma~\ref{lem:sec3-connected-span}: $B$ spans the span
  of $\Phi$, and distinct members of $B$ have nonzero inner product exactly when their vertices
  are adjacent, so the graph of that lemma is a copy of $X$, which is connected.~\qed\end{proof}
\end{proof}

\begin{namedresult}{Lemma}{(transport of representations)}\label{lem:4-transport}
 \lean{SimpleGraph.IsRepresentedIn.of_isometricInclusion}
 \uses{def:represented,def:line-iso}
 \assume{$S$ and $T$ are line systems, $f:\Phi(S)\to\Phi(T)$ preserves inner products (for
 instance an isomorphism, or the inclusion of $\Phi(S)$ in $\Phi(T)$ when $\Phi(S)\sbs\Phi(T)$),
 and $X$ is represented in $S$ by $(\Zx_u)$.}
 \prove{$X$ is represented in $T$ by $(f(\Zx_u))$.}
\end{namedresult}
\begin{proof}
 \pfsketch\ The Gram matrix does not change.
 \qedstep
 \begin{proof}\pf\ $(f(\Zx_u),f(\Zx_w))=(\Zx_u,\Zx_w)$ for all $u,w$ by \asm{0}, and
  $f(\Zx_u)\in\Phi(T)$; apply Definition~\ref{def:represented}.~\qed\end{proof}
\end{proof}

\begin{namedresult}{Lemma}{(subfamilies of a representation)}\label{lem:4-sub}
 \lean{SimpleGraph.IsRepresentation.induce}
 \uses{def:represented}
 \assume{$(\Zx_u)_{u\in V(X)}$ is a representation of a graph $X$ in $V$, and $U\sbs V(X)$.}
 \prove{$(\Zx_u)_{u\in U}$ is a representation in $V$ of the subgraph $X[U]$ induced on $U$.}
\end{namedresult}
\begin{proof}
 \pfsketch\ The Gram matrix of the subfamily is a principal submatrix.
 \qedstep
 \begin{proof}\pf\ Two vertices of $U$ are adjacent in $X[U]$ exactly when they are adjacent in
  $X$, so $(\Zx_u,\Zx_w)$ for $u,w\in U$ is the $uw$-entry of $2I+A(X[U])$; apply
  Definition~\ref{def:represented}.~\qed\end{proof}
\end{proof}

\begin{namedresult}{Lemma}{(a root system containing the lines)}\label{lem:4-rootsystem}
 \lean{SimpleGraph.IsRepresentation.exists_isRepresentedIn_rootType}
 \uses{def:represented,lem:4-lines,lem:4-transport,lem:sec3-extension,lem:sec3-transport,lem:sec3-spans,def:2.2,lem:2.3,thm:3.5,thm:3.6,def:3.1,def:3.3}
 \assume{$X$ is a connected graph with a representation $(\Zx_u)$ in a finite-dimensional real
 inner product space $V$, and the span of the $\Zx_u$ has dimension $k$.}
 \prove{\begin{pfenum}
  \item $X$ is represented in the model of one of $\typ A_k$, $\typ D_k$ ($k\ge4$), $\typ E_6$,
   $\typ E_7$ and $\typ E_8$, and the rank of that model is $k$.
  \item $X$ is represented in $\typ A_\ell$ for some $\ell\ge1$, in $\typ D_\ell$ for some
   $\ell\ge4$, or in $\typ E_8$.
 \end{pfenum}}
\end{namedresult}
\begin{proof}
 \pfsketch\ This is the paper's ``Such a set of lines is contained in a root system. If the
 graph is connected, then the root system is indecomposable.'' \repair\ The paper does not say
 why a set that is not maximal lies in an indecomposable star-closed set of the same dimension;
 we will use the fill of Gap 13 (errata), which is Lemma~\ref{lem:sec3-extension}. We will then
 place $\typ E_6$ and $\typ E_7$ inside $\typ E_8$. The first of the two sentences holds without
 connectedness: the vectors of squared norm 2 in the integral span of the $\Zx_u$ form a
 star-closed line system in which $X$ is represented. We use this system for Remark 4.14; in
 Lean it is \path{SimpleGraph.IsRepresentation.latticeRoots}, with the lemmas
 \path{isStarClosed_latticeRoots} and \path{isRepresentedIn_latticeRoots} in the same
 namespace. The claim of the introduction of [CGSS], that the lines lie in one of the root
 systems, needs connectedness (errata, Implicit hypothesis 1).
 \pflet{$S=S_\Zx$ be the line system of Lemma~\ref{lem:4-lines} and $W$ the span of $\Phi(S)$,
 which has dimension $k$.}
 \step{s4rsnonempty}{$S$ is nonempty and indecomposable.}
 \begin{proof}\pf\ A connected graph has a vertex, and Lemma~\ref{lem:4-lines}, assertion 3,
  applies by \asm{0}.~\qed\end{proof}
 \step{s4rsmax}{There is an indecomposable line system $T$ in $W$, maximal in $W$ and spanning
  $W$, with $\Phi(S)\sbs\Phi(T)$.}
 \begin{proof}
  \step{s4rsmaxyes}{\case{$S$ is maximal in $W$.}\prove{$T=S$ will do.}}
  \begin{proof}\pf\ $S$ is indecomposable by \stepref{s4rsnonempty}, maximal by \asm{2}, and
   spans $W$ by the definition of $W$.~\qed\end{proof}
  \step{s4rsmaxno}{\case{$S$ is not maximal in $W$.}\prove{Such a $T$ exists.}}
  \begin{proof}\pf\ Apply Lemma~\ref{lem:sec3-extension} to $S$, with
   \stepref{s4rsnonempty} and \asm{2}.~\qed\end{proof}
  \qedstep
  \begin{proof}\pf\ The two cases are exhaustive.~\qed\end{proof}
 \end{proof}
 \step{s4rstype}{$T$ is isomorphic to the model of $\typ A_\ell$ ($\ell\ge1$), $\typ D_\ell$
  ($\ell\ge4$), $\typ E_6$, $\typ E_7$ or $\typ E_8$, of rank $k$.}
 \begin{proof}\pf\ Apply Theorem~\ref{thm:3.6} to $T$ of \stepref{s4rsmax}; its list is part of
  this one. The isomorphism preserves rank by Lemma~\ref{lem:sec3-transport}, and $T$ has rank
  $k$ since it spans $W$.~\qed\end{proof}
 \step{s4rsinT}{$X$ is represented in $T$ by $(\Zx_u)$.}
 \begin{proof}\pf\ Lemma~\ref{lem:4-transport}, with the inclusion $\Phi(S)\sbs\Phi(T)$ of
  \stepref{s4rsmax}.~\qed\end{proof}
 \step{s4rsone}{Assertion 1 holds.}
 \begin{proof}\pf\ Apply Lemma~\ref{lem:4-transport} to \stepref{s4rsinT} and the isomorphism
  of \stepref{s4rstype}. The model of $\typ A_\ell$ or $\typ D_\ell$ has rank $\ell$ by
  Lemma~\ref{lem:sec3-spans}, so $\ell=k$.~\qed\end{proof}
 \step{s4rsE}{$\Phi(\typ E_6)\sbs\Phi(\typ E_8)$ and $\Phi(\typ E_7)\sbs\Phi(\typ E_8)$.}
 \begin{proof}\pf\ Definition~\ref{def:3.3}.~\qed\end{proof}
 \qedstep
 \begin{proof}\pf\ Assertion 1 is \stepref{s4rsone}. For assertion 2, in the cases $\typ E_6$
  and $\typ E_7$ apply Lemma~\ref{lem:4-transport} once more to the inclusions of
  \stepref{s4rsE}.~\qed\end{proof}
\end{proof}

\begin{namedresult}{Lemma}{(the line graph)}\label{lem:4-linegraph}
 \lean{SimpleGraph.lineGraph_isRepresentation}
 \uses{def:represented}
 \assume{$Y$ is a graph with vertex set $\{1,\dots,n\}$.}
 \prove{The vectors $\Ze_i+\Ze_j$ of $\re^n$, for the edges $\{i,j\}$ of $Y$, form a
 representation of $L(Y)$ in $\re^n$.}
\end{namedresult}
\begin{proof}
 \pfsketch\ The inner product counts common ends.
 \step{s4lgip}{For edges $e,f$ of $Y$, $(\sum_{i\in e}\Ze_i,\sum_{j\in f}\Ze_j)=|e\cap f|$.}
 \begin{proof}\pf\ Expand; $(\Ze_i,\Ze_j)$ is 1 if $i=j$ and 0
  otherwise.~\qed\end{proof}
 \qedstep
 \begin{proof}\pf\ $|e\cap e|=2$; for $e\ne f$, $|e\cap f|$ is 1 when $e$ and $f$ share an end,
  that is, when they are adjacent in $L(Y)$, and 0 otherwise. Apply
  \stepref{s4lgip}.~\qed\end{proof}
\end{proof}


\begin{cgss}{Theorem}{4.1}\label{thm:4.1}
 \lean{SimpleGraph.isRepresentedIn_An_iff}
 \uses{def:represented,def:3.1,lem:4-linegraph}
 \assume{$X$ is a graph and $n\ge1$.}
 \prove{$X$ is represented in $\typ A_n$ if and only if $X\cong L(Y)$ for a bipartite graph $Y$
  on the vertex set $\{1,\dots,n+1\}$.}
 \new\ The paper's ``any such graph can be represented by a subset of $\typ A_n$'' needs $Y$ to
 have at most $n+1$ vertices; we will state it with that bound, as the Lean statement does.
\end{cgss}
\begin{proof}
 \pfsketch\ We will read off $Y$ from the supports of the vectors, and show that each coordinate
 has one sign.
 \step{s41fwd}{\assume{$X$ is represented in $\typ A_n$ by $(\Zx_u)$.}
  \prove{$X\cong L(Y)$ for a bipartite graph $Y$ on $\{1,\dots,n+1\}$.}}
 \begin{proof}
  \step{s41pq}{\pick\ $p(u)\ne q(u)$ in $\{1,\dots,n+1\}$ with
   $\Zx_u=\Ze_{p(u)}-\Ze_{q(u)}$ for each vertex $u$.}
  \begin{proof}\pf\ $\Zx_u\in\Phi(\typ A_n)$ by \asm{1}; apply
   Definition~\ref{def:3.1}.~\qed\end{proof}
  \step{s41sign}{For all vertices $u,w$, $p(u)\ne q(w)$.}
  \begin{proof}\pf\ If $p(u)=q(w)=i$, then $u\ne w$, and
   $(\Zx_u,\Zx_w)=(\Ze_i-\Ze_{q(u)},\Ze_{p(w)}-\Ze_i)=-1-[q(u)=p(w)]<0$, since $p(w)\ne i$ and
   $q(u)\ne i$; this contradicts Definition~\ref{def:represented}.~\qed\end{proof}
  \step{s41ip}{For $u\ne w$, $(\Zx_u,\Zx_w)=[p(u)=p(w)]+[q(u)=q(w)]$.}
  \begin{proof}\pf\ Expand $(\Ze_{p(u)}-\Ze_{q(u)},\Ze_{p(w)}-\Ze_{q(w)})$; the two cross terms
   vanish by \stepref{s41sign}.~\qed\end{proof}
  \step{s41inj}{The supports $\{p(u),q(u)\}$ of distinct vertices are distinct.}
  \begin{proof}\pf\ If $u\ne w$ have the same support, then $p(u)=p(w)$ and $q(u)=q(w)$ by
   \stepref{s41sign}, and $(\Zx_u,\Zx_w)=2$ by \stepref{s41ip}, against
   Definition~\ref{def:represented}.~\qed\end{proof}
  \pflet{$Y$ be the graph on $\{1,\dots,n+1\}$ with edges $\{p(u),q(u)\}$, and
   $P=\{p(u):u\in V(X)\}$.}
  \step{s41bip}{$Y$ is bipartite, with parts $P$ and its complement.}
  \begin{proof}\pf\ Each edge $\{p(u),q(u)\}$ has one end $p(u)$ in $P$ and the other end
   $q(u)$ outside $P$, by \stepref{s41sign}.~\qed\end{proof}
  \step{s41iso}{$u\mapsto\{p(u),q(u)\}$ is an isomorphism from $X$ to $L(Y)$.}
  \begin{proof}\pf\ It is a bijection onto the edges of $Y$ by \stepref{s41inj}. For $u\ne w$,
   the edges share an end exactly when $p(u)=p(w)$ or $q(u)=q(w)$, by \stepref{s41sign}, that is,
   when $(\Zx_u,\Zx_w)=1$ by \stepref{s41ip}, which is adjacency of $u$ and
   $w$.~\qed\end{proof}
  \qedstep
  \begin{proof}\pf\ \stepref{s41bip} and \stepref{s41iso}.~\qed\end{proof}
 \end{proof}
 \step{s41back}{\assume{$Y$ is bipartite on $\{1,\dots,n+1\}$, with parts $P$ and $Q$, and
  $X\cong L(Y)$.}\prove{$X$ is represented in $\typ A_n$.}}
 \begin{proof}
  \pflet{$\Zx_e=\Ze_p-\Ze_q$ for the edge $e=\{p,q\}$ of $Y$ with $p\in P$ and $q\in Q$.}
  \step{s41backip}{For edges $e\ne f$, $(\Zx_e,\Zx_f)$ is 1 when $e$ and $f$ share an end, and
   0 otherwise; and $(\Zx_e,\Zx_e)=2$.}
  \begin{proof}\pf\ Write $e=\{p,q\}$ and $f=\{p',q'\}$ with $p,p'\in P$. Then $p\ne q'$ and
   $q\ne p'$, since $P\cap Q=\varnothing$, and $(\Zx_e,\Zx_f)=[p=p']+[q=q']$, which is 2 for
   $e=f$ and at most 1 for $e\ne f$.~\qed\end{proof}
  \qedstep
  \begin{proof}\pf\ $\Zx_e\in\Phi(\typ A_n)$ by Definition~\ref{def:3.1}, and $(\Zx_e)$
   represents $L(Y)$ by \stepref{s41backip}; compose with the isomorphism of
   \asm{1}.~\qed\end{proof}
 \end{proof}
 \qedstep
 \begin{proof}\pf\ \stepref{s41fwd} and \stepref{s41back}.~\qed\end{proof}
\end{proof}

\begin{namedresult}{Remark}{(after Theorem 4.1)}\label{rem:4-after41}
 \lean{LineSystem.An.card_nonneg_subset_le}\lean{LineSystem.An.nonneg_subset_maximal_iff}
 \lean{LineSystem.An.exists_nonneg_subset_card_eq}
 \lean{LineSystem.An.isRepresentation_lineGraph_completeBipartite_sigmaP}
 \uses{def:represented,def:3.1,lem:4-sub}
 \assume{$n\ge1$ and $N=n+1$. A set $\Sigma\sbs\Phi(\typ A_n)$ has \emph{nonnegative inner
 products} if $(\Zx,\Zy)\ge0$ for all $\Zx,\Zy\in\Sigma$, and is \emph{maximal} if no subset of
 $\Phi(\typ A_n)$ that properly contains it has nonnegative inner products. For
 $P\sbs\{1,\dots,N\}$ let $\Sigma_P=\{\Ze_p-\Ze_q:p\in P,\ q\notin P\}$.}
 \prove{\begin{pfenum}
  \item A set $\Sigma$ with nonnegative inner products is maximal if and only if
   $\Sigma=\Sigma_P$ for some $P$ with $1\le|P|\le n$.
  \item For such $P$, with $m=|P|$, the vectors $\Ze_p-\Ze_q$, indexed by the edges $\{p,q\}$ of
   the complete bipartite graph $Y_P$ with parts $P$ and $\{1,\dots,N\}\setminus P$, represent
   $L(Y_P)\cong L(K_{m,N-m})$.
  \item Every $\Sigma$ with nonnegative inner products has $|\Sigma|\le\lfloor\frac14N^2\rfloor$,
   and $|\Sigma_P|=\lfloor\frac14N^2\rfloor$ for $|P|=\lfloor\frac12N\rfloor$.
 \end{pfenum}}
\end{namedresult}
\begin{proof}
 \pfsketch\ \new\ The paper states the remark without proof. We will show, as in the proof of
 Theorem 4.1, that each coordinate occurs with one sign in $\Sigma$, and then count.
 \pflet{$\Zx=\Ze_{p(\Zx)}-\Ze_{q(\Zx)}$ for $\Zx\in\Phi(\typ A_n)$, by
  Definition~\ref{def:3.1}.}
 \step{s4raip}{\assume{$P\sbs\{1,\dots,N\}$, $p,p'\in P$ and $q,q'\notin P$.}
  \prove{$(\Ze_p-\Ze_q,\Ze_{p'}-\Ze_{q'})=[p=p']+[q=q']$; in particular $\Sigma_P$ has
  nonnegative inner products.}}
 \begin{proof}\pf\ Expand; the cross terms vanish since $p\ne q'$ and $q\ne
  p'$.~\qed\end{proof}
 \step{s4rasign}{\assume{$\Sigma$ has nonnegative inner products.}\prove{$\Sigma\sbs\Sigma_P$
  for $P=\{p(\Zx):\Zx\in\Sigma\}$.}}
 \begin{proof}\pf\ We need $q(\Zx)\notin P$ for $\Zx\in\Sigma$. If $q(\Zx)=p(\Zy)=i$ with
  $\Zy\in\Sigma$, then $p(\Zx)\ne i$ and $q(\Zy)\ne i$, so
  $(\Zx,\Zy)=(\Ze_{p(\Zx)}-\Ze_i,\Ze_i-\Ze_{q(\Zy)})=-1-[p(\Zx)=q(\Zy)]<0$, against
  \asm{1}.~\qed\end{proof}
 \step{s4ramax}{\assume{$\Sigma$ is maximal with nonnegative inner products.}
  \prove{$\Sigma=\Sigma_P$ for some $P$ with $1\le|P|\le n$.}}
 \begin{proof}
  \step{s4ramaxne}{$\Sigma$ is nonempty.}
  \begin{proof}\pf\ $\{\Ze_1-\Ze_2\}$ has nonnegative inner products, as $N\ge2$, and
   properly contains $\varnothing$; apply \asm{1}.~\qed\end{proof}
  \step{s4ramaxeq}{$\Sigma=\Sigma_P$ for $P=\{p(\Zx):\Zx\in\Sigma\}$.}
  \begin{proof}\pf\ $\Sigma\sbs\Sigma_P$ by \stepref{s4rasign}, and $\Sigma_P$ has nonnegative
   inner products by \stepref{s4raip}; apply \asm{1}.~\qed\end{proof}
  \qedstep
  \begin{proof}\pf\ By \stepref{s4ramaxne}, $P$ of \stepref{s4ramaxeq} is nonempty, and its
   complement contains $q(\Zx)$ for $\Zx\in\Sigma=\Sigma_P$; so $1\le|P|\le
   N-1=n$.~\qed\end{proof}
 \end{proof}
 \step{s4raback}{\assume{$1\le|P|\le n$.}\prove{$\Sigma_P$ is maximal with nonnegative inner
  products.}}
 \begin{proof}\pf\ $\Sigma_P$ has nonnegative inner products by \stepref{s4raip}. A root
  $\Zz=\Ze_a-\Ze_b\notin\Sigma_P$ has $b\in P$ or $a\notin P$; pick $p\in P$ and $q\notin P$, by
  \asm{1}. If $b\in P$, then $(\Zz,\Ze_b-\Ze_q)=-1-[a=q]<0$; if $a\notin P$, then
  $(\Zz,\Ze_p-\Ze_a)=-1-[b=p]<0$. Thus no larger set has nonnegative inner
  products.~\qed\end{proof}
 \step{s4rarep}{Assertion 2 holds.}
 \begin{proof}\pf\ For edges $e=\{p,q\}$ and $f=\{p',q'\}$ of $Y_P$ with $p,p'\in P$,
  \stepref{s4raip} gives the inner product $[p=p']+[q=q']=|e\cap f|$, which is 2 for $e=f$, and
  for $e\ne f$ is 1 when $e$ and $f$ share an end and 0 otherwise; apply
  Definition~\ref{def:represented}. $Y_P\cong K_{m,N-m}$.~\qed\end{proof}
 \step{s4racard}{$|\Sigma_P|=|P|(N-|P|)$.}
 \begin{proof}\pf\ $(p,q)\mapsto\Ze_p-\Ze_q$ is injective on $P\times(\{1,\dots,N\}\setminus
  P)$.~\qed\end{proof}
 \step{s4rabound}{\assume{$\Sigma$ has nonnegative inner products.}
  \prove{$|\Sigma|\le\lfloor\frac14N^2\rfloor$.}}
 \begin{proof}
  \pflet{$P$ be as in \stepref{s4rasign}, and $m=|P|$.}
  \step{s4rabsub}{$|\Sigma|\le m(N-m)$.}
  \begin{proof}\pf\ $\Sigma\sbs\Sigma_P$ by \asm{1} and \stepref{s4rasign}, and
   $|\Sigma_P|=m(N-m)$ by \stepref{s4racard}.~\qed\end{proof}
  \step{s4rabsq}{$m(N-m)\le\lfloor\frac14N^2\rfloor$.}
  \begin{proof}\pf\ $4m(N-m)=N^2-(N-2m)^2\le N^2$, and $m(N-m)$ is an
   integer.~\qed\end{proof}
  \qedstep
  \begin{proof}\pf\ \stepref{s4rabsub} and \stepref{s4rabsq}.~\qed\end{proof}
 \end{proof}
 \step{s4raeq}{\assume{$m=\lfloor\frac12N\rfloor$.}\prove{$m(N-m)=\lfloor\frac14N^2\rfloor$.}}
 \begin{proof}\pf\ For $N=2k$ both sides are $k^2$; for $N=2k+1$ both are
  $k^2+k$.~\qed\end{proof}
 \qedstep
 \begin{proof}\pf\ Assertion 1 is \stepref{s4ramax} and \stepref{s4raback}, and assertion 2 is
  \stepref{s4rarep}. Assertion 3 is \stepref{s4rabound}, with \stepref{s4racard} and
  \stepref{s4raeq} for $|P|=\lfloor\frac12N\rfloor$, which lies between 1 and $n$ as
  $N\ge2$.~\qed\end{proof}
\end{proof}
The paper states assertions 1 to 3 in the form ``The maximal subsets of $\typ A_n$ with
nonnegative inner products represent the graphs $L(K_{m,n+1-m})$, and have maximum cardinality
$\lfloor\frac14(n+1)^2\rfloor$, attained for $m=\lfloor\frac12(n+1)\rfloor$.'' A set
$\Sigma$ with nonnegative inner products lies in $\Sigma_P$ for $P=\{p(\Zx):\Zx\in\Sigma\}$, by
the second step of the proof. Thus, by assertion 2 and Lemma~\ref{lem:4-sub}, the graph that
$\Sigma$ represents is $L(Y)$, where $Y$ is the subgraph of $Y_P$ whose edges are the
$\{p,q\}$ with $\Ze_p-\Ze_q\in\Sigma$, and assertion 3 bounds its number of vertices.

\subsection{Theorems 4.2 to 4.9}

\begin{namedresult}{Definition}{(cocktail party graphs)}\label{def:CP}
 \lean{SimpleGraph.cocktailParty}
 \uses{}
 \assume{$m\ge0$ is an integer.}
 \pflet{The \emph{cocktail party graph} $CP(m)$ has vertex set $\{1,\dots,m\}\times\{0,1\}$, and
 $(i,s)$ is adjacent to $(j,t)$ exactly when $i\ne j$.}
 Thus $CP(m)$ is $K_{2m}$ with a perfect matching deleted, and it is $(2m-2)$-regular on $2m$
 vertices; $CP(0)$ has no vertices.
\end{namedresult}

\begin{namedresult}{Definition}{(generalized line graphs)}\label{def:glg}
 \lean{SimpleGraph.generalizedLineGraph}\lean{SimpleGraph.IsGeneralizedLineGraph}
 \uses{def:CP}
 \assume{$Y$ is a graph with vertex set $\{1,\dots,m\}$ and $a_1,\dots,a_m\ge0$ are integers.}
 \pflet{The \emph{generalized line graph} $L(Y;a_1,\dots,a_m)$ has as vertices the edges of $Y$
 and the triples $(i,\alpha,s)$ with $1\le i\le m$, $1\le\alpha\le a_i$ and $s\in\{0,1\}$. Two
 edges are adjacent when they are adjacent in $L(Y)$; $(i,\alpha,s)$ and $(j,\beta,t)$ are
 adjacent when $i=j$ and $\alpha\ne\beta$; and $(i,\alpha,s)$ is adjacent to the edge $e$ when
 $i\in e$. A graph is a \emph{generalized line graph} if it is isomorphic to some
 $L(Y;a_1,\dots,a_m)$.}
 For each $i$, the triples $(i,\alpha,s)$ span a copy of $CP(a_i)$, and $L(Y;0,\dots,0)\cong
 L(Y)$.
\end{namedresult}

\begin{namedresult}{Lemma}{(the representation of a generalized line graph)}\label{lem:4-glgrep}
 \lean{SimpleGraph.generalizedLineGraph_isRepresentedIn_Dn}
 \uses{def:glg,def:represented,def:3.1,lem:4-linegraph}
 \assume{$Y$ and $a_1,\dots,a_m$ are as in Definition~\ref{def:glg}, $n=\sum_{i=1}^m(1+a_i)$,
 and the orthonormal basis of $\re^n$ is written $\Ze_{i,\alpha}$, for $1\le i\le m$ and
 $0\le\alpha\le a_i$.}
 \prove{The vectors $\Zx_e=\Ze_{i,0}+\Ze_{j,0}$ for the edges $e=\{i,j\}$ of $Y$, and
 $\Zx_{(i,\alpha,s)}=\Ze_{i,0}+(-1)^s\Ze_{i,\alpha}$, lie in $\Phi(\typ D_n)$ and represent
 $L(Y;a_1,\dots,a_m)$.}
\end{namedresult}
\begin{proof}
 \pfsketch\ \repair\ We read the paper's ``$1\le\alpha\le\alpha_i$'' as $1\le\alpha\le a_i$
 (errata, Error 13, item 5). We will compute the three kinds of inner products.
 \step{s4glgee}{For edges $e,f$, $(\Zx_e,\Zx_f)=|e\cap f|$.}
 \begin{proof}\pf\ Lemma~\ref{lem:4-linegraph}, in the coordinates $\Ze_{i,0}$, gives 2 for
  $e=f$, 1 when $e\ne f$ share an end, and 0 otherwise; this is $|e\cap f|$.~\qed\end{proof}
 \step{s4glgce}{$(\Zx_{(i,\alpha,s)},\Zx_e)$ is 1 if $i\in e$ and 0 otherwise.}
 \begin{proof}\pf\ $\Ze_{i,\alpha}$ with $\alpha\ge1$ does not occur in $\Zx_e$, and
  $(\Ze_{i,0},\Zx_e)=[i\in e]$.~\qed\end{proof}
 \step{s4glgcc}{$(\Zx_{(i,\alpha,s)},\Zx_{(j,\beta,t)})$ is 2 if the triples are equal, 0 if
  $i=j$ and $\alpha=\beta$ and $s\ne t$, 1 if $i=j$ and $\alpha\ne\beta$, and 0 if $i\ne j$.}
 \begin{proof}\pf\ The inner product is
  $[i=j]\bigl(1+(-1)^{s+t}[\alpha=\beta]\bigr)$.~\qed\end{proof}
 \qedstep
 \begin{proof}\pf\ Each vector is a sum or difference of two distinct basis vectors, so lies in
  $\Phi(\typ D_n)$ by Definition~\ref{def:3.1}. By \stepref{s4glgee}, \stepref{s4glgce} and
  \stepref{s4glgcc}, the Gram matrix is $2I$ plus the adjacency matrix of
  Definition~\ref{def:glg}.~\qed\end{proof}
\end{proof}


\begin{namedresult}{Lemma}{(the normal form in $\typ D_n$)}\label{lem:4-Dnform}
 \lean{SimpleGraph.IsRepresentation.exists_normal_form_Dn}
 \uses{def:represented,def:3.1,def:glg,lem:4-glgrep,lem:4-transport}
 \assume{$n\ge1$, and $(\Zx_u)_{u\in V(X)}$ represents a graph $X$ in $\typ D_n$.}
 \prove{There are a set $C\sbs\{1,\dots,n\}$, integers $a_c\ge0$ for $c\in C$, distinct
 coordinates $p_{c,\alpha}$ for $c\in C$ and $1\le\alpha\le a_c$, which are the members of
 $\{1,\dots,n\}\setminus C$, and signs $\epsilon_1,\dots,\epsilon_n\in\{1,-1\}$ such that, with
 $D$ the diagonal isometry $\Ze_i\mapsto\epsilon_i\Ze_i$ and $\Zy_u=D\Zx_u$:
 \begin{pfenum}
  \item $(\Zy_u)$ represents $X$ in $\typ D_n$;
  \item for each $c\in C$, $1\le\alpha\le a_c$ and $s\in\{0,1\}$ exactly one vertex $u$ has
   $\Zy_u=\Ze_c+(-1)^s\Ze_{p_{c,\alpha}}$, and every vertex $u$ not of this form has
   $\Zy_u=\Ze_i+\Ze_j$ with $i\ne j$ in $C$, distinct such vertices having distinct supports;
  \item $X\cong L(Y;(a_c)_{c\in C})$, where $Y$ is the graph on $C$ whose edges are the supports
   of the vertices of assertion 2 not of this form.
 \end{pfenum}
 In particular $|C|+\sum_{c\in C}a_c=n$.}
 \new\ This is the proof of assertion 1 of Theorem 4.2 below, with the coordinates kept; the
 remark after that theorem counts them.
\end{namedresult}
\begin{proof}
 \pfsketch\ We will follow the paper: we will find the pairs $\Ze_c\pm\Ze_p$, then choose the
 signs of the other coordinates. For $n<4$, $\typ D_n$ is the coordinate model
 $\{\pm\Ze_i\pm\Ze_j:i\ne j\}$, as in the remark of Definition~\ref{def:3.1}.
 \repair\ The paper leaves both steps as assertions, among them that the set of pairs is
 unique and maximal (errata, Gap 19); we will prove what we need.
 \step{s42form}{Each $\Zx_u$ is $\sigma\Ze_i+\tau\Ze_j$ with $i\ne j$ and
  $\sigma,\tau\in\{1,-1\}$; its support has two members.}
 \begin{proof}\pf\ \asm{0} and Definition~\ref{def:3.1}.~\qed\end{proof}
 \step{s42same}{\assume{$u\ne w$ and $\supp(\Zx_u)=\supp(\Zx_w)=\{i,j\}$.}
  \prove{$(\Zx_u,\Zx_w)=0$, and $\Zx_u,\Zx_w$ agree in exactly one coordinate.}}
 \begin{proof}\pf\ The inner product is $\sigma\sigma'+\tau\tau'\in\{-2,0,2\}$ by
  \stepref{s42form}, and lies in $\{0,1\}$ by \asm{0} and Definition~\ref{def:represented}; so it is 0, and
  exactly one of $\sigma\sigma'$, $\tau\tau'$ is 1.~\qed\end{proof}
 \pflet{A \emph{pair} is a set $\{u,w\}$ as in \stepref{s42same}. Its \emph{centre} is the
  coordinate in which $\Zx_u$ and $\Zx_w$ agree and its \emph{petal} the other one; let $\Pi$
  be the set of petals of all pairs.}
 \step{s42petal}{\assume{$p\in\Pi$ is the petal of the pair $\{u,w\}$, with centre $c$.}
  \prove{$u$ and $w$ are the only vertices $z$ with $p\in\supp(\Zx_z)$.}}
 \begin{proof}
  \step{s42petalsupp}{\assume{$p\in\supp(\Zx_z)$.}\prove{$\supp(\Zx_z)=\{c,p\}$.}}
  \begin{proof}\pf\ Otherwise $\Zx_z=\rho\Ze_p+\ldots$ has no $\Ze_c$ term, so
   $(\Zx_z,\Zx_u)=\rho\tau$ and $(\Zx_z,\Zx_w)=-\rho\tau$, where $\pm\tau$ are the coefficients
   of $\Ze_p$ in $\Zx_u,\Zx_w$; one of them is $-1$, against \asm{0} and
   Definition~\ref{def:represented}.~\qed\end{proof}
  \qedstep
  \begin{proof}\pf\ If $z\ne u,w$, then $z$, $u$ and $w$ have pairwise the support $\{c,p\}$ by
   \stepref{s42petalsupp}, so pairwise agree in exactly one of two coordinates by
   \stepref{s42same}. Their three sign vectors in $\{1,-1\}^2$ would then be pairwise adjacent in
   the graph on $\{1,-1\}^2$ joining vectors that differ in exactly one coordinate; that graph
   is $C_4$, which has no triangle.~\qed\end{proof}
 \end{proof}
 \step{s42centre}{No centre lies in $\Pi$, and every petal is the petal of exactly one
  pair.}
 \begin{proof}\pf\ The pair of a petal is the set of vertices whose support contains it, by
  \stepref{s42petal}. If the centre $c$ of the pair $\{u,w\}$ were the petal of a pair, that
  pair would be $\{u,w\}$ by the same step, whose petal is not $c$.~\qed\end{proof}
 \pflet{$C=\{1,\dots,n\}\setminus\Pi$.}
 \step{s42onesign}{\assume{$i\in C$.}\prove{All $\Zx_u$ with $i\in\supp(\Zx_u)$ have the same
  coefficient $\epsilon_i$ at $\Ze_i$.}}
 \begin{proof}\pf\ Suppose $\Zx_u$ and $\Zx_w$ have coefficients $1$ and $-1$ at $\Ze_i$. If
  their supports are equal, $\{u,w\}$ is a pair that disagrees at $i$, so $i\in\Pi$ by the
  definition of petals; otherwise $(\Zx_u,\Zx_w)=-1$. The first contradicts \asm{1}, and the second \asm{0}
  and Definition~\ref{def:represented}.~\qed\end{proof}
 \pflet{$\epsilon_i=1$ for the $i\in C$ in no support, and $\epsilon_p=1$ for $p\in\Pi$; let
  $D$ be the diagonal isometry $\Ze_i\mapsto\epsilon_i\Ze_i$, and $\Zy_u=D\Zx_u$.}
 \step{s42flip}{$(\Zy_u)$ represents $X$ in $\typ D_n$; every $\Zy_u$ is $\Ze_c\pm\Ze_p$ for a
  pair with centre $c$ and petal $p$, or $\Ze_i+\Ze_j$ with $i,j\in C$.}
 \begin{proof}\pf\ $D$ is an isometry that preserves $\Phi(\typ D_n)$, so
  Lemma~\ref{lem:4-transport} applies. A vertex in a pair has support $\{c,p\}$, and
  $c\in C$ by \stepref{s42centre}, so its $\Ze_c$ coefficient becomes 1 by
  \stepref{s42onesign}. Every other support lies in $C$ by \stepref{s42petal}, and both
  coefficients become 1 by \stepref{s42onesign}.~\qed\end{proof}
 \pflet{$a_c$ be the number of pairs with centre $c$, for $c\in C$; number the petals of centre
  $c$ as $p_{c,1},\dots,p_{c,a_c}$; and let $Y$ be the graph on $C$ whose edges are the supports
  of the vertices in no pair.}
 \step{s42petals}{The $p_{c,\alpha}$ are distinct, and they are the members of
  $\{1,\dots,n\}\setminus C$.}
 \begin{proof}\pf\ Each petal has one pair by \stepref{s42centre}, and so one centre, which lies
  in $C$ by the same step; and $\{1,\dots,n\}\setminus C=\Pi$.~\qed\end{proof}
 \step{s42two}{Assertion 2 holds.}
 \begin{proof}\pf\ The vertices of the pair with petal $p_{c,\alpha}$ have
  $\Zy_u=\Ze_c\pm\Ze_{p_{c,\alpha}}$ by \stepref{s42flip}, one with each sign by
  \stepref{s42same}, and no other vertex has $p_{c,\alpha}$ in its support by
  \stepref{s42petal}. The vertices in no pair have $\Zy_u=\Ze_i+\Ze_j$ with $i,j\in C$ by
  \stepref{s42flip}, and distinct ones have distinct supports, since equal supports make a
  pair.~\qed\end{proof}
 \step{s42bij}{The map $\psi$ sending a vertex in no pair to its support, and the vertex $u$
  with $\Zy_u=\Ze_c+(-1)^s\Ze_{p_{c,\alpha}}$ to $(c,\alpha,s)$, is a bijection from $V(X)$
  to the vertex set of $L(Y;(a_c)_{c\in C})$.}
 \begin{proof}\pf\ By \stepref{s42two}, every triple $(c,\alpha,s)$ is the image of exactly one
  vertex, and the other vertices map injectively onto the edges of $Y$, by the definition of
  $Y$.~\qed\end{proof}
 \step{s42three}{Assertion 3 holds.}
 \begin{proof}\pf\ Renumber $C$ as $\{1,\dots,m\}$. The vectors $\Zy_u$ are those of
  Lemma~\ref{lem:4-glgrep} for $\psi(u)$, with $\Ze_{c,0}=\Ze_c$ and
  $\Ze_{c,\alpha}=\Ze_{p_{c,\alpha}}$, by \stepref{s42two} and \stepref{s42petals}. Both
  families have Gram matrix $2I$ plus the adjacency matrix, of $X$ by \stepref{s42flip} and of
  $L(Y;a)$ by that lemma; so $\psi$ of \stepref{s42bij} preserves adjacency.~\qed\end{proof}
 \qedstep
 \begin{proof}\pf\ Assertion 1 is part of \stepref{s42flip}, and assertions 2 and 3 are
  \stepref{s42two} and \stepref{s42three}. The last sentence counts $\{1,\dots,n\}$ by
  \stepref{s42petals}.~\qed\end{proof}
\end{proof}


\begin{cgss}{Theorem}{4.2}\label{thm:4.2}
 \lean{SimpleGraph.isRepresentedIn_Dn_iff}
 \uses{def:represented,def:3.1,def:glg,lem:4-glgrep,lem:4-Dnform}
 \assume{$X$ is a graph and $n\ge1$.}
 \prove{\begin{pfenum}
  \item If $X$ is represented in $\typ D_n$, then $X$ is a generalized line graph.
  \item $L(Y;a_1,\dots,a_m)$ is represented in $\typ D_n$ for $n=\sum_i(1+a_i)$.
 \end{pfenum}}
\end{cgss}
\begin{proof}
 \pfsketch\ Assertion 2 is Lemma~\ref{lem:4-glgrep}. For assertion 1 we will follow the paper:
 we will find the pairs $\Ze_c\pm\Ze_p$, then choose the signs of the other coordinates; this is
 Lemma~\ref{lem:4-Dnform}, which also fills Gap 19 (errata).
 \repair\ The paper's ``whose lines from a subset'' is ``form'' (errata, Error 13, item 4).
 \step{s42glg}{Assertion 2 holds.}
 \begin{proof}\pf\ Lemma~\ref{lem:4-glgrep}.~\qed\end{proof}
 \step{s42one}{\assume{$(\Zx_u)$ represents $X$ in $\typ D_n$.}
  \prove{$X$ is a generalized line graph.}}
 \begin{proof}\pf\ Lemma~\ref{lem:4-Dnform}, assertion 3, with \asm{1}, and $n\ge1$ by
  \asm{0}; then Definition~\ref{def:glg}.~\qed\end{proof}
 \qedstep
 \begin{proof}\pf\ \stepref{s42one} and \stepref{s42glg}.~\qed\end{proof}
\end{proof}

\begin{namedresult}{Remark}{(after Theorem 4.2)}\label{rem:4-after42}
 \lean{LineSystem.Dn.card_nonneg_subset_le}
 \lean{LineSystem.Dn.exists_isRepresentation_generalizedLineGraph_top_of_maximal}
 \lean{LineSystem.Dn.exists_nonneg_subset_card_eq}
 \uses{def:represented,def:3.1,def:glg,lem:4-Dnform,lem:4-linegraph,rem:4-after41}
 \assume{$n\ge1$, and $\Sigma\sbs\Phi(\typ D_n)$ has nonnegative inner products, in the sense of
 \hyperref[rem:4-after41]{the remark after Theorem 4.1}, with $\typ D_n$ for $\typ A_n$.}
 \prove{\begin{pfenum}
  \item If $\Sigma$ is maximal, then there are $m\ge1$ and $a_1,\dots,a_m\ge0$ with
   $\sum_{i=1}^m(1+a_i)=n$ such that the vectors of $\Sigma$, suitably indexed, represent
   $L(K_m;a_1,\dots,a_m)$.
  \item If $n\ge4$, then $|\Sigma|\le\frac12n(n-1)$; the set $\{\Ze_i+\Ze_j:i<j\}$ has
   nonnegative inner products and $\frac12n(n-1)$ members, and represents $L(K_n)$.
 \end{pfenum}}
\end{namedresult}
\begin{proof}
 \pfsketch\ \new\ The paper states the remark without proof. We will apply
 Lemma~\ref{lem:4-Dnform}, which is the proof of Theorem 4.2 with the coordinates
 kept, and count. \new\ The bound fails for $n=2$ and $n=3$, where the paper does not define
 $\typ D_n$: the sets $\{\Ze_1\pm\Ze_2\}$ and $\{\Ze_1\pm\Ze_2,\Ze_1\pm\Ze_3\}$ are larger than
 $\frac12n(n-1)$.
 \step{s4rdrep}{The family $(\Zx)_{\Zx\in\Sigma}$ represents in $\typ D_n$ the graph $X$ on
  $\Sigma$ in which $\Zx\sim\Zy$ when $(\Zx,\Zy)=1$.}
 \begin{proof}\pf\ For $\Zx\ne\Zy$ in $\Sigma$, $(\Zx,\Zy)$ is an integer at least 0 by
  \asm{0} and Definition~\ref{def:3.1}, and $(\Zx,\Zy)<2$ since
  $(\Zx-\Zy,\Zx-\Zy)=4-2(\Zx,\Zy)>0$; so it is 0 or 1. Apply
  Definition~\ref{def:represented}.~\qed\end{proof}
 \step{s4rdnf}{\pick\ $C$, $(a_c)$, $(p_{c,\alpha})$, $D$, $(\Zy_\Zx)$ and $Y$ as in
  Lemma~\ref{lem:4-Dnform} for the representation of \stepref{s4rdrep}, and let $m=|C|$. Then
  $m+\sum_ca_c=n$ and $m\ge1$.}
 \begin{proof}\pf\ The equation is the last sentence of the lemma. If $m=0$, then
  $C=\varnothing$, so the equation gives $n=0$, against \asm{0}.~\qed\end{proof}
 \step{s4rdcount}{$|\Sigma|=|E(Y)|+2\sum_ca_c\le\frac12m(m-1)+2(n-m)$.}
 \begin{proof}\pf\ Lemma~\ref{lem:4-Dnform}, assertion 2, gives two members of $\Sigma$ for each
  $p_{c,\alpha}$ and one for each edge of $Y$, a graph on the $m$ vertices of $C$; and $\sum_c
  a_c=n-m$ by \stepref{s4rdnf}.~\qed\end{proof}
 \step{s4rdmax}{\assume{$\Sigma$ is maximal.}\prove{$Y$ is complete.}}
 \begin{proof}
  \step{s4rdadd}{\assume{$i\ne j$ lie in $C$.}\prove{$\Zz=D(\Ze_i+\Ze_j)$ lies in
   $\Phi(\typ D_n)$, and $(\Zz,\Zx)\ge0$ for every $\Zx\in\Sigma$.}}
  \begin{proof}
   \step{s4rdaddroot}{$\Zz=\epsilon_i\Ze_i+\epsilon_j\Ze_j$, and $\Zz\in\Phi(\typ D_n)$.}
   \begin{proof}\pf\ $D$ is the diagonal isometry $\Ze_k\mapsto\epsilon_k\Ze_k$, and $i\ne j$ by
    \asm{2}; apply Definition~\ref{def:3.1}.~\qed\end{proof}
   \step{s4rdaddip}{\assume{$\Zx\in\Sigma$.}\prove{$(\Zz,\Zx)=(\Ze_i+\Ze_j,\Zy_\Zx)$.}}
   \begin{proof}\pf\ $D$ is an isometry with $D^2=I$, so
    $(D(\Ze_i+\Ze_j),\Zx)=(\Ze_i+\Ze_j,D\Zx)$, and $D\Zx=\Zy_\Zx$.~\qed\end{proof}
   \step{s4rdaddnn}{\assume{$\Zx\in\Sigma$.}\prove{$(\Ze_i+\Ze_j,\Zy_\Zx)\ge0$.}}
   \begin{proof}\pf\ By Lemma~\ref{lem:4-Dnform}, assertion 2, $\Zy_\Zx$ is
    $\Ze_c\pm\Ze_{p_{c,\alpha}}$ with $p_{c,\alpha}\notin C$, which gives $[c\in\{i,j\}]$ by
    \asm{2}, or $\Ze_k+\Ze_l$, which gives $|\{i,j\}\cap\{k,l\}|$.~\qed\end{proof}
   \qedstep
   \begin{proof}\pf\ \stepref{s4rdaddroot}, \stepref{s4rdaddip} and
    \stepref{s4rdaddnn}.~\qed\end{proof}
  \end{proof}
  \qedstep
  \begin{proof}\pf\ For $i\ne j$ in $C$, $\Zz\in\Sigma$ by \stepref{s4rdadd} and \asm{1}. Its
   image $\Zy_\Zz=\Ze_i+\Ze_j$ has no petal $p_{c,\alpha}$ in its support, so it is of the second
   kind of Lemma~\ref{lem:4-Dnform}, assertion 2, and $\{i,j\}$ is an edge of
   $Y$.~\qed\end{proof}
 \end{proof}
 \step{s4rdone}{Assertion 1 holds.}
 \begin{proof}\pf\ Renumber $C$ as $\{1,\dots,m\}$. By Lemma~\ref{lem:4-Dnform}, assertion 3,
  and \stepref{s4rdmax}, $X\cong L(K_m;a_1,\dots,a_m)$, and $\sum_i(1+a_i)=n$ by
  \stepref{s4rdnf}; index $\Sigma$ by this isomorphism.~\qed\end{proof}
 \step{s4rdineq}{\assume{$n\ge4$ and $1\le m\le n$.}
  \prove{$\frac12m(m-1)+2(n-m)\le\frac12n(n-1)$.}}
 \begin{proof}\pf\ The difference of the two sides is $\frac12(n-m)(n+m-5)$, and $n+m\ge5$ by
  \asm{1}.~\qed\end{proof}
 \step{s4rdK}{$\{\Ze_i+\Ze_j:i<j\}$ has $\frac12n(n-1)$ members and nonnegative inner products,
  and represents $L(K_n)$.}
 \begin{proof}\pf\ Lemma~\ref{lem:4-linegraph} for $Y=K_n$, whose Gram matrix $2I+A$ has no
  negative entry; the vectors lie in $\Phi(\typ D_n)$ by
  Definition~\ref{def:3.1}.~\qed\end{proof}
 \qedstep
 \begin{proof}\pf\ Assertion 1 is \stepref{s4rdone}. For assertion 2, \stepref{s4rdcount} and
  \stepref{s4rdineq}, with $1\le m\le n$ by \stepref{s4rdnf}, give the bound, and
  \stepref{s4rdK} the example.~\qed\end{proof}
\end{proof}
The paper states the remark as ``A maximal subset of $\typ D_n$ with nonnegative inner products
represents $L(K_m;\mathbf a)$ for some $m$ and $\mathbf a$; the maximum cardinality is
$\frac12n(n-1)$, attained by $L(K_n)$.'' Assertion 1 holds for every $n\ge1$. For $n=4$ the
maximum 6 is attained also by sets with $(m;a_1)=(1;3)$; these represent $L(K_1;3)=CP(3)$,
which is isomorphic to $L(K_4)$, so what differs is the set and its parameters, not the graph.
For $n\ge5$ the maximum is attained only by sets with $m=n$, which represent $L(K_n)$, since
$\frac12(n-m)(n+m-5)>0$ for $1\le m<n$.


\begin{namedresult}{Lemma}{(the symmetric square)}\label{lem:4-symsq}
 \lean{EuclideanSpace.inner_symSq}
 \uses{}
 \assume{$\ell\ge1$, and for $\Zv\in\re^\ell$ the vector $\sigma(\Zv)\in\re^{\ell(\ell+1)/2}$
 has the coordinates $\Zv(i)^2$ for $1\le i\le\ell$ and $\sqrt2\,\Zv(i)\Zv(j)$ for\\ $1\le i<j\le\ell$.}
 \prove{$(\sigma(\Zv),\sigma(\Zw))=(\Zv,\Zw)^2$ for all $\Zv,\Zw\in\re^\ell$.}
 \new\ The paper uses the projections onto the lines, in the space of symmetric maps with the
 trace inner product; the coordinates of that space give $\sigma$, and $\re^{36}$ is easier in
 Lean.
\end{namedresult}
\begin{proof}
 \pfsketch\ We will expand the square.
 \qedstep
 \begin{proof}\pf\ $(\Zv,\Zw)^2=\bigl(\sum_i\Zv(i)\Zw(i)\bigr)^2=\sum_i\Zv(i)^2\Zw(i)^2+2\sum_{i<j}\Zv(i)\Zv(j)\Zw(i)\Zw(j)$, which
  is $(\sigma(\Zv),\sigma(\Zw))$.~\qed\end{proof}
\end{proof}

\begin{namedresult}{Lemma}{(valencies from the sets $N_\Zr$)}\label{lem:4-valency}
 \lean{SimpleGraph.IsRepresentedIn.degree_le_of_card_N_le}
 \uses{def:2.1,def:represented,lem:4-lines,def:line-system,facts:sec4}
 \assume{\begin{pfenum}
  \item $S$ is a star-closed line system, $b\ge0$, and $|N_\Zr|\le2b$ for every
   $\Zr\in\Phi(S)$, where $N_\Zr=\{\Zx\in\Phi(S):(\Zx,\Zr)=1\}$.
  \item $X$ is represented in $S$ by $(\Zx_u)$, and $u$ is a vertex of $X$.
 \end{pfenum}}
 \prove{$u$ has at most $b$ neighbours.}
 \repair\ The paper cites ``Lemma 3.7'' for such bounds (errata, Error 6), and Theorem 3.7 gives
 only the number $|N_\Zr|$; the pairing below is the fill of Gap 14 (errata), with LN6 of
 \path{cgss/prerequisites.md}.
\end{namedresult}
\begin{proof}
 \pfsketch\ We will pair the vectors with inner product 1 with $\Zr=\Zx_u$.
 \pflet{$\Zr=\Zx_u$.}
 \step{s4valpair}{\assume{$\Zx\in N_\Zr$.}\prove{$\Zr-\Zx\in N_\Zr$, $\Zr-\Zx\ne\Zx$ and
  $(\Zx,\Zr-\Zx)=-1$.}}
 \begin{proof}\pf\ Since $(\Zx,-\Zr)=-1$ and $S$ is star-closed by \asm[1]{0}, the vector
  $\Zx-\Zr$ and its negative lie in $\Phi(S)$. Further $(\Zr-\Zx,\Zr)=2-1=1$ and
  $(\Zx,\Zr-\Zx)=1-2=-1$, which also gives $\Zr-\Zx\ne\Zx$.~\qed\end{proof}
 \step{s4valnbr}{$w\mapsto\Zx_w$ maps the neighbours of $u$ injectively into $N_\Zr$, and its
  image meets each pair $\{\Zx,\Zr-\Zx\}$ at most once.}
 \begin{proof}\pf\ $(\Zx_w,\Zr)=1$ for a neighbour $w$ by Definition~\ref{def:represented}, and
  the map is injective by Lemma~\ref{lem:4-lines}, assertion 1. Two neighbours have inner product
  0 or 1, not the $-1$ of \stepref{s4valpair}.~\qed\end{proof}
 \qedstep
 \begin{proof}\pf\ By \stepref{s4valpair}, $\Zx\mapsto\Zr-\Zx$ is an involution of $N_\Zr$
  without fixed points, so $N_\Zr$ is the union of at most $b$ disjoint pairs by
  \asm[1]{0}; apply \stepref{s4valnbr}.~\qed\end{proof}
\end{proof}

\begin{namedresult}{Computation}{C20}\label{comp:C20}
 \lean{LineSystem.An.card_N}\lean{LineSystem.Dn.card_N}\lean{LineSystem.E₆.card_N}
 \lean{LineSystem.E₇.card_N}\lean{LineSystem.E₈.card_N}
 \uses{def:3.1,def:3.2,def:3.3,thm:3.5,lem:4-valency}
 \assume{$\Zr$ is a vector of the model of $\typ A_\ell$, $\typ D_\ell$, $\typ E_6$,
 $\typ E_7$ or $\typ E_8$, and $X$ is represented in that model.}
 \prove{$|N_\Zr|$ is $2(\ell-1)$, $4(\ell-2)$, 20, 32 or 56 respectively, and every vertex of
 $X$ has at most $\ell-1$, $2(\ell-2)$, 10, 16 or 28 neighbours respectively.}
\end{namedresult}
\begin{proof}
 \pfsketch\ The counts are LN6 of \path{cgss/prerequisites.md}. We will cite their Lean
 proofs, which are parametric in $\ell$ for $\typ A_\ell$ and $\typ D_\ell$ and use
 \texttt{decide} for the three exceptional models.
 \step{s4c20count}{The counts of $|N_\Zr|$ hold.}
 \begin{proof}\pf\ \path{LineSystem.An.card_N}, \path{LineSystem.Dn.card_N}, and
  \leancite{LineSystem.E₆.card_N}{\texttt{LineSystem.E\ensuremath{{}_6}.card\_N}} and its
  analogues for $\typ E_7$ and $\typ E_8$.~\qed\end{proof}
 \qedstep
 \begin{proof}\pf\ The models are star-closed by Theorem~\ref{thm:3.5}; apply
  Lemma~\ref{lem:4-valency} with half of each count.~\qed\end{proof}
\end{proof}

\begin{namedresult}{Computation}{C7}\label{comp:C7}
 \lean{Computations.C7.isRepresentedIn_E₈}
 \uses{def:3.1,def:3.2,def:represented}
 \assume{In $\re^9$, let $\Zc=\frac13\sum_{l=1}^9\Ze_l$, $\Zy_h=\Ze_h-\Ze_9$ for $1\le h\le8$,
 and $\Zz_{ij}=\Zc-\Ze_i-\Ze_j-\Ze_9$ for $1\le i<j\le8$.}
 \prove{These 36 vectors lie in model (i) of $\typ E_8$ and represent a graph $X_{36}$ on 36
 vertices in which $\Zy_1$ has 28 neighbours.}
\end{namedresult}
\begin{proof}
 \pfsketch\ We will compute the inner products; the Lean proof is by \texttt{decide} on the
 integer vectors $3\Zy_h$ and $3\Zz_{ij}$.
 \step{s4c7mem}{The vectors lie in model (i).}
 \begin{proof}\pf\ $\Zy_h\in\Phi(\typ A_8)$, and $\Zz_{ij}$ is
  $\frac13\sum_l\Ze_l-\sum_{l\in T}\Ze_l$ with $T=\{i,j,9\}$; apply
  Definition~\ref{def:3.2}, (i).~\qed\end{proof}
 \step{s4c7ip}{$(\Zy_h,\Zy_k)=1$ and $(\Zz_{ij},\Zz_{kl})=|\{i,j\}\cap\{k,l\}|$ for distinct
  vectors, and $(\Zy_h,\Zz_{ij})$ is 0 if $h\in\{i,j\}$ and 1 otherwise.}
 \begin{proof}\pf\ $(\Zy_h,\Zc)=0$ and $(\Zc,\Zc)=1$, and $(\Zc,\Ze_l)=\frac13$. Then
  $(\Zy_h,\Zy_k)=1+[h=k]$, $(\Zy_h,\Zz_{ij})=1-[h\in\{i,j\}]$ and
  $(\Zz_{ij},\Zz_{kl})=1-2+|\{i,j,9\}\cap\{k,l,9\}|$.~\qed\end{proof}
 \qedstep
 \begin{proof}\pf\ Each vector has squared norm 2 by \stepref{s4c7ip}, and the other inner
  products lie in $\{0,1\}$; so they represent a graph by Definition~\ref{def:represented}. The
  neighbours of $\Zy_1$ are the 7 vectors $\Zy_k$ and the $\binom72=21$ vectors $\Zz_{ij}$ with
  $1\notin\{i,j\}$.~\qed\end{proof}
\end{proof}

\begin{cgss}{Theorem}{4.3}\label{thm:4.3}
 \lean{SimpleGraph.IsRepresentedIn.card_le_36}\lean{SimpleGraph.IsRepresentedIn.degree_le_28}
 \lean{SimpleGraph.exists_isRepresentedIn_E₈_card_36}
 \uses{def:represented,lem:4-symsq,comp:C20,lem:4-transport,comp:C7,thm:3.5,facts:sec4}
 \assume{$X$ is represented in $\typ E_8$ by $(\Zx_u)$.}
 \prove{\begin{pfenum}
  \item $X$ has at most 36 vertices.
  \item Every vertex of $X$ has at most 28 neighbours.
  \item There is a graph represented in $\typ E_8$ with 36 vertices and a vertex with 28
   neighbours.
 \end{pfenum}}
\end{cgss}
\begin{proof}
 \pfsketch\ We will follow the paper: the Gram matrix of the symmetric squares is
 nonsingular. \repair\ The bound on the valency needs the pairing of Gap 14 (errata).
 \step{s43pd}{$4I+A(X)$ is positive definite.}
 \begin{proof}\pf\ $2I+A(X)$ is the Gram matrix of $(\Zx_u)$, so positive semidefinite by
  \path{Matrix.posSemidef_gram}; for $\Zz\ne0$,
  $\Zz^T(4I+A(X))\Zz=2|\Zz|^2+\Zz^T(2I+A(X))\Zz>0$.~\qed\end{proof}
 \step{s43gram}{The Gram matrix of $(\sigma(\Zx_u))$ in $\re^{36}$ is $4I+A(X)$.}
 \begin{proof}\pf\ Lemma~\ref{lem:4-symsq} with $\ell=8$: the entries are $(\Zx_u,\Zx_w)^2$,
  which are 4, 1 and 0 by Definition~\ref{def:represented}.~\qed\end{proof}
 \step{s43one}{Assertion 1 holds.}
 \begin{proof}\pf\ The family $(\sigma(\Zx_u))$ is linearly independent by \stepref{s43pd},
  \stepref{s43gram} and \path{Matrix.posDef_gram_iff_linearIndependent}; apply
  \path{LinearIndependent.fintype_card_le_finrank} and
  \path{finrank_euclideanSpace_fin}.~\qed\end{proof}
 \step{s43two}{Assertion 2 holds.}
 \begin{proof}\pf\ Computation~\ref{comp:C20}, for $\typ E_8$.~\qed\end{proof}
 \step{s43three}{Assertion 3 holds.}
 \begin{proof}\pf\ The graph $X_{36}$ of Computation~\ref{comp:C7} is represented in model (i);
  models (i) and (ii) are isomorphic by Theorem~\ref{thm:3.5}, so it is represented in model
  (ii) by Lemma~\ref{lem:4-transport}.~\qed\end{proof}
 \qedstep
 \begin{proof}\pf\ \stepref{s43one}, \stepref{s43two} and \stepref{s43three}.~\qed\end{proof}
\end{proof}


\begin{namedresult}{Lemma}{(trace and rank)}\label{lem:4-tracerank}
 \lean{Matrix.IsHermitian.trace_sq_le_rank_mul_sum_entry_sq}
 \uses{facts:sec4}
 \assume{$N$ is a real symmetric matrix indexed by a finite set $V$.}
 \prove{$(\tr N)^2\le\rk N\cdot\sum_{u,w\in V}N(u,w)^2$.}
 \new\ This is the inequality behind [LS73] Theorem 3.6; we will name it, since Lemma EQ1 and the
 remark after Theorem 4.4 both use it.
 \begingroup\raggedright\new\ In Lean the entrywise form stated here is
 \path{Matrix.IsHermitian.trace_sq_le_rank_mul_sum_entry_sq}, in
 \path{LineGraphs/Representation/Theorem44.lean}, which Theorem 4.4 and the remark after it use.
 The eigenvalue form $(\tr N)^2\le\rk N\cdot\sum_{u\in V}\lambda_u^2$ of the steps below is
 \path{Matrix.IsHermitian.trace_sq_le_rank_mul_sum_sq}, in
 \path{LineGraphs/Spectral/Applications.lean}; the name of the entrywise form differs, since
 that one was taken first.\par\endgroup
\end{namedresult}
\begin{proof}
 \pfsketch\ We will apply the Cauchy--Schwarz inequality to the nonzero eigenvalues.
 \pflet{$\lambda_u$, for $u\in V$, be the eigenvalues \path{Matrix.IsHermitian.eigenvalues} of
 $N$, and $Z=\{u\in V:\lambda_u\ne0\}$.}
 \step{s4trtr}{$\tr N=\sum_{u\in Z}\lambda_u$ and $\rk N=|Z|$.}
 \begin{proof}\pf\ \path{Matrix.IsHermitian.trace_eq_sum_eigenvalues}, with the zero terms
  dropped, and \path{Matrix.IsHermitian.rank_eq_card_non_zero_eigs}.~\qed\end{proof}
 \step{s4trsq}{$\sum_{u,w}N(u,w)^2=\sum_{u\in Z}\lambda_u^2$.}
 \begin{proof}\pf\ By \path{Matrix.IsHermitian.spectral_theorem}, $N=UDU^T$ with $U$ orthogonal
  and $D$ diagonal with the entries $\lambda_u$. As $N$ is symmetric,
  $\sum_{u,w}N(u,w)^2=\tr(NN^T)=\tr(UD^2U^T)=\tr(D^2U^TU)=\sum_u\lambda_u^2$, where the third
  equality is \path{Matrix.trace_mul_comm}; drop the zero terms.~\qed\end{proof}
 \qedstep
 \begin{proof}\pf\ By \path{Finset.sum_mul_sq_le_sq_mul_sq} for the functions $\lambda$ and 1 on
  $Z$, $\bigl(\sum_{u\in Z}\lambda_u\bigr)^2\le|Z|\sum_{u\in Z}\lambda_u^2$; apply
  \stepref{s4trtr} and \stepref{s4trsq}.~\qed\end{proof}
\end{proof}

\begin{namedresult}{Lemma}{(the relative bound, EQ1)}\label{lem:4-EQ1}
 \lean{InnerProductSpace.card_le_relativeBound}
 \uses{lem:4-gramfactor,lem:4-tracerank}
 \assume{The unit vectors $\Zv_1,\dots,\Zv_n$ of a real inner product space span a subspace of
 dimension $k$, where $1\le k\le7$, and $(\Zv_i,\Zv_j)=\pm\frac13$ for $i\ne j$.}
 \prove{$n\le k\cdot\frac89/(1-\frac k9)$, and in particular $n\le28$.}
\end{namedresult}
\begin{proof}
 \pfsketch\ This is [LS73] Theorem 3.6, which [CGSS] cites as [15]; we will prove it from the
 trace of the Gram matrix.
 \pflet{$M$ be the Gram matrix of the $\Zv_i$.}
 \step{s4eq1tr}{$\tr M=n$, $\tr M^2=n+\frac19n(n-1)$ and $\rk M=k$.}
 \begin{proof}\pf\ The diagonal entries of $M$ are 1 and the others $\pm\frac13$, and
  $\tr M^2=\sum_{i,j}(\Zv_i,\Zv_j)^2$; the rank is Lemma~\ref{lem:4-gramfactor},
  assertion 1.~\qed\end{proof}
 \step{s4eq1cs}{$n^2\le k\bigl(n+\frac19n(n-1)\bigr)$.}
 \begin{proof}\pf\ Lemma~\ref{lem:4-tracerank} for $M$, where
  $\sum_{i,j}M(i,j)^2=\tr M^2$ as $M$ is symmetric; apply \stepref{s4eq1tr}.~\qed\end{proof}
 \qedstep
 \begin{proof}\pf\ Divide \stepref{s4eq1cs} by $n$: $n(1-\frac k9)\le\frac89k$, with
  $1-\frac k9>0$. The bound $\frac89k/(1-\frac k9)=8k/(9-k)$ increases with $k$, and is 28 for
  $k=7$.~\qed\end{proof}
\end{proof}

\begin{namedresult}{Lemma}{(lines spanning eight dimensions, EQ2)}\label{lem:4-EQ2}
 \lean{InnerProductSpace.card_le_fourteen_of_finrank_eq_eight}
 \uses{def:represented,lem:4-rootsystem,lem:sec3-spans,comp:C20}
 \assume{$U$ is a real inner product space of dimension 8, spanned by unit vectors
 $\Zv_1,\dots,\Zv_n$ with $(\Zv_i,\Zv_j)=\pm\frac13$ for $i\ne j$.}
 \prove{$n\le14$.}
 \repair\ This is the cone construction of \path{cgss/prerequisites.md} (EQ2), which fills
 the unstated part of [LS73] (errata, Gap 7).
\end{namedresult}
\begin{proof}
 \pfsketch\ We will put the lines on a cone around a root $\Zr$ in $U\oplus\re\Ze$ and count
 the pairs of $N_\Zr$.
 \pflet{$\Zr=\sqrt2\,\Ze$ and $\Zu_i=\sqrt{3/2}\,\Zv_i+\Ze/\sqrt2$ in $U\oplus\re\Ze$, where $\Ze$
 is a unit vector orthogonal to $U$.}
 \step{s4eq2ip}{$(\Zr,\Zr)=(\Zu_i,\Zu_i)=2$, $(\Zu_i,\Zr)=1$, and
  $(\Zu_i,\Zu_j)=\frac32(\Zv_i,\Zv_j)+\frac12\in\{0,1\}$ for $i\ne j$.}
 \begin{proof}\pf\ Expand, with $(\Zv_i,\Ze)=0$.~\qed\end{proof}
 \pflet{$X'$ be the graph on $\{0,1,\dots,n\}$ in which 0 is adjacent to every $i$, and $i\ne j$
 are adjacent when $(\Zu_i,\Zu_j)=1$.}
 \step{s4eq2rep}{$\Zr,\Zu_1,\dots,\Zu_n$ represent the connected graph $X'$, and span a space of
  dimension 9.}
 \begin{proof}\pf\ Definition~\ref{def:represented} and \stepref{s4eq2ip}; 0 is adjacent to
  every other vertex. The span contains $\Ze=\Zr/\sqrt2$ and each
  $\Zv_i=(\Zu_i-\Ze/\sqrt2)/\sqrt{3/2}$, so it is $U\oplus\re\Ze$.~\qed\end{proof}
 \step{s4eq2type}{$X'$ is represented in $\typ A_9$ or in $\typ D_9$.}
 \begin{proof}\pf\ Lemma~\ref{lem:4-rootsystem}, assertion 1, with $k=9$; the models of
  $\typ E_6$, $\typ E_7$, $\typ E_8$ have ranks 6, 7, 8 by
  Lemma~\ref{lem:sec3-spans}.~\qed\end{proof}
 \qedstep
 \begin{proof}\pf\ The vertex 0 has $n$ neighbours, at most $9-1=8$ in $\typ A_9$ and at most
  $2(9-2)=14$ in $\typ D_9$, by Computation~\ref{comp:C20} and \stepref{s4eq2type}.~\qed\end{proof}
\end{proof}

\begin{namedresult}{Computation}{C8}\label{comp:C8}
 \lean{Computations.C8.quadratic_bound}
 \uses{}
 \assume{$n\le28$ and $d\le n-1$ are nonnegative integers with $n<2d+4$ and
 $0\le(d-2n+16)+14(2n-2d-4)/(2d+4-n)$.}
 \prove{$d\le16$, or $n=28$ and $d=17$.}
\end{namedresult}
\begin{proof}
 \pfsketch\ We will clear the positive denominator; the Lean proof is \texttt{decide} over the
 finitely many pairs.
 \qedstep
 \begin{proof}\pf\ Multiply by $2d+4-n>0$:
  $(d-2n+16)(2d+4-n)+14(2n-2d-4)\ge0$. Among the pairs $(n,d)$ with $n\le28$, $d\le n-1$ and
  $n<2d+4$, this holds only for $d\le16$ and for $(28,17)$; this is a finite check of the 406
  pairs with $d<n\le28$.~\qed\end{proof}
\end{proof}

\begin{namedresult}{Lemma}{(the matrices of a regular graph)}\label{lem:4-regmat}
 \lean{SimpleGraph.IsRegularOfDegree.posSemidef_centered}
 \lean{SimpleGraph.IsRegularOfDegree.rank_centered}
 \lean{SimpleGraph.IsRegularOfDegree.posSemidef_equiangular}
 \lean{SimpleGraph.IsRegularOfDegree.rank_equiangular}
 \uses{facts:sec4}
 \assume{\begin{pfenum}
  \item $X$ is $d$-regular on $n\ge1$ vertices, $Q=2I+A(X)$ is positive semidefinite, and $\Zj$
   is the all-ones vector.
  \item $t=(d+2)/(2n)$, $M=\frac12Q-tJ$ and $E=\frac23Q-\frac13J$.
 \end{pfenum}}
 \prove{\begin{pfenum}
  \item $Q\Zj=(d+2)\Zj$, and $J\Zz=0$ and $(Q\Zz,\Zj)=0$ for every $\Zz\perp\Zj$.
  \item $M\Zj=0$ and $M\Zz=\frac12Q\Zz$ for $\Zz\perp\Zj$; $M$ is positive semidefinite, and
   $\rk M=\rk Q-1$.
  \item If $n<2d+4$, then $E\Zj=\frac13(2d+4-n)\Zj$ and $E\Zz=\frac23Q\Zz$ for $\Zz\perp\Zj$; $E$
   is positive semidefinite, and $\rk E=\rk Q$.
 \end{pfenum}}
 \new\ The paper reads these facts off the eigenvalues of $A$; we will compare ranges instead, and
 state the facts once, for Theorem 4.4 and for the remark after it.
\end{namedresult}
\begin{proof}
 \pfsketch\ We will split $\re^n$ as $\re\Zj\oplus\Zj^\perp$.
 \pflet{$W=Q(\Zj^\perp)$ be the image of $\Zj^\perp$ under $Q$.}
 \step{s4rmj}{$A\Zj=d\Zj$, $Q\Zj=(d+2)\Zj$ and $J\Zj=n\Zj$.}
 \begin{proof}\pf\ \path{SimpleGraph.adjMatrix_mulVec_const_apply_of_regular} with \asm[1]{0};
  then $Q\Zj=2\Zj+A\Zj$, and the rows of $J$ sum to $n$.~\qed\end{proof}
 \step{s4rmperp}{\assume{$\Zz\perp\Zj$.}\prove{$J\Zz=0$ and $(Q\Zz,\Zj)=0$.}}
 \begin{proof}\pf\ Every entry of $J\Zz$ is $(\Zz,\Zj)=0$. Further $Q$ is symmetric, so
  $(Q\Zz,\Zj)=(\Zz,Q\Zj)=(d+2)(\Zz,\Zj)=0$ by \stepref{s4rmj}.~\qed\end{proof}
 \step{s4rmsplit}{Every $\Zy\in\re^n$ is $c\Zj+\Zz$ with $c=(\Zy,\Zj)/n$ and $\Zz\perp\Zj$.}
 \begin{proof}\pf\ $(\Zy-c\Zj,\Zj)=(\Zy,\Zj)-cn=0$, as $(\Zj,\Zj)=n\ge1$ by
  \asm[1]{0}.~\qed\end{proof}
 \step{s4rmrange}{The range of $Q$ is $\re\Zj+W$.}
 \begin{proof}\pf\ By \stepref{s4rmsplit} and \stepref{s4rmj}, $Q(c\Zj+\Zz)=c(d+2)\Zj+Q\Zz$ for
  $\Zz\perp\Zj$, and $d+2\ne0$.~\qed\end{proof}
 \step{s4rmcap}{$\re\Zj\cap W=0$.}
 \begin{proof}\pf\ If $c\Zj\in W$, then $cn=(c\Zj,\Zj)=0$ by \stepref{s4rmperp}, and $n\ge1$ by
  \asm[1]{0}; so $c=0$.~\qed\end{proof}
 \step{s4rmrk}{$\rk Q=1+\dim W$.}
 \begin{proof}\pf\ By the definition \path{Matrix.rank}, $\rk Q$ is the dimension of the range
  of $Q$, which is $\re\Zj+W$ by \stepref{s4rmrange}. Apply
  \path{Submodule.finrank_sup_add_finrank_inf_eq} with \stepref{s4rmcap}, and
  \path{finrank_span_singleton}, as $\Zj\ne0$.~\qed\end{proof}
 \step{s4rmtwo}{Assertion 2 holds.}
 \begin{proof}
  \step{s4rmMact}{$M\Zj=0$, and $M\Zz=\frac12Q\Zz$ for $\Zz\perp\Zj$.}
  \begin{proof}\pf\ $M\Zj=\frac12(d+2)\Zj-tn\Zj=0$ by \stepref{s4rmj}, since
   $tn=\frac12(d+2)$ by \asm[2]{0}; and $J\Zz=0$ by \stepref{s4rmperp}.~\qed\end{proof}
  \step{s4rmMpsd}{$M$ is positive semidefinite.}
  \begin{proof}\pf\ $M$ is symmetric, as $Q$ and $J$ are. For $\Zy=c\Zj+\Zz$ as in
   \stepref{s4rmsplit}, $(\Zy,M\Zy)=(c\Zj+\Zz,\frac12Q\Zz)=\frac12(\Zz,Q\Zz)\ge0$, by
   \stepref{s4rmMact}, \stepref{s4rmperp} and \asm[1]{0}.~\qed\end{proof}
  \step{s4rmMrk}{$\rk M=\rk Q-1$.}
  \begin{proof}\pf\ By \stepref{s4rmsplit} and \stepref{s4rmMact}, $M(c\Zj+\Zz)=\frac12Q\Zz$, so
   the range of $M$ is $W$; apply \stepref{s4rmrk} and the definition
   \path{Matrix.rank}.~\qed\end{proof}
  \qedstep
  \begin{proof}\pf\ \stepref{s4rmMact}, \stepref{s4rmMpsd} and \stepref{s4rmMrk}.~\qed\end{proof}
 \end{proof}
 \step{s4rmthree}{\assume{$n<2d+4$.}\prove{The conclusion of assertion 3 holds.}}
 \begin{proof}
  \pflet{$b=\frac13(2d+4-n)$, which is positive by \asm{1}.}
  \step{s4rmEact}{$E\Zj=b\Zj$, and $E\Zz=\frac23Q\Zz$ for $\Zz\perp\Zj$.}
  \begin{proof}\pf\ $E\Zj=\frac23(d+2)\Zj-\frac13n\Zj$ by \stepref{s4rmj}, and $J\Zz=0$ by
   \stepref{s4rmperp}.~\qed\end{proof}
  \step{s4rmEpsd}{$E$ is positive semidefinite.}
  \begin{proof}\pf\ $E$ is symmetric. For $\Zy=c\Zj+\Zz$ as in \stepref{s4rmsplit},
   $(\Zy,E\Zy)=(c\Zj+\Zz,cb\Zj+\frac23Q\Zz)=c^2bn+\frac23(\Zz,Q\Zz)\ge0$, by \stepref{s4rmEact}
   and \stepref{s4rmperp}, with $(\Zz,\Zj)=0$, $b>0$ and \asm[1]{0}.~\qed\end{proof}
  \step{s4rmErk}{$\rk E=\rk Q$.}
  \begin{proof}\pf\ By \stepref{s4rmsplit} and \stepref{s4rmEact}, $E(c\Zj+\Zz)=cb\Zj+\frac23Q\Zz$
   with $b\ne0$, so the range of $E$ is $\re\Zj+W$, which is the range of $Q$ by
   \stepref{s4rmrange}.~\qed\end{proof}
  \qedstep
  \begin{proof}\pf\ \stepref{s4rmEact}, \stepref{s4rmEpsd} and \stepref{s4rmErk}.~\qed\end{proof}
 \end{proof}
 \qedstep
 \begin{proof}\pf\ Assertion 1 is \stepref{s4rmj} and \stepref{s4rmperp}; assertions 2 and 3 are
  \stepref{s4rmtwo} and \stepref{s4rmthree}.~\qed\end{proof}
\end{proof}

\begin{cgss}{Theorem}{4.4}\label{thm:4.4}
 \lean{SimpleGraph.IsRepresentedIn.card_le_28_of_regular}
 \lean{SimpleGraph.IsRepresentedIn.degree_le_16_of_regular}
 \uses{def:represented,lem:4-gramfactor,lem:4-regmat,lem:4-sub,lem:4-symsq,lem:4-EQ1,lem:4-EQ2,lem:4-tracerank,lem:4-rootsystem,lem:sec3-spans,comp:C8,comp:C20,facts:sec4}
 \assume{\begin{pfenum}
  \item $X$ is $d$-regular on $n\ge1$ vertices.
  \item $X$ is represented in $\typ E_8$ by $(\Zx_u)$.
 \end{pfenum}}
 \prove{$n\le28$ and $d\le16$.}
 \begingroup\raggedright\new\ In Lean the theorem is proved as
 \path{SimpleGraph.IsRepresentedIn.card_le_28_of_regular_of_classification} and
 \path{SimpleGraph.IsRepresentedIn.degree_le_16_of_regular_of_classification}, in
 \path{LineGraphs/Representation/Theorem44.lean}, with the classification of the maximal
 indecomposable line systems of the nine-dimensional space of Lemma~\ref{lem:4-EQ2} as an
 explicit hypothesis. The settled names above, without the hypothesis, are proved in
 \path{LineGraphs/Representation/Theorem44Unconditional.lean}, by
 discharging it with Theorem~3.6 (\path{LineSystem.IsMaximal.exists_iso_rootType}).\par\endgroup
\end{cgss}
\begin{proof}
 \pfsketch\ We will follow the paper, with these repairs. The equiangular facts come from
 Lemmas EQ1 and EQ2 in place of [15] (errata, Error 8 and Gap 7, decision 1 of
 \path{cgss/prerequisites.md}). \repair\ The case $n=28$, $d=12$ needs its own treatment (errata,
 Error 15), and the final valency bound needs the pairing of Gap 14 (errata). The printed Gram
 matrix is that of the vectors $\sigma(\Zq_u)$, not of projections (errata, Error 13, item 6), and
 ``each eigenvalue $\alpha\ne d$'' concerns the action on $\Zj^\perp$ (item 7); we will use that
 action through quadratic forms, without counting eigenvalues.
 \pflet{$Q=2I+A(X)$, $\Zj$ the all-ones vector, $t=(d+2)/(2n)$, $M=\frac12Q-tJ$ and
 $E=\frac23Q-\frac13J$, as in Lemma~\ref{lem:4-regmat}.}
 \step{s44spec}{$Q$ is positive semidefinite of rank at most 8.}
 \begin{proof}\pf\ $Q$ is the Gram matrix of $n$ vectors of
  $\re^8$ by \asm[2]{0}, so it is positive semidefinite by \path{Matrix.posSemidef_gram}, and of rank
  at most 8 by Lemma~\ref{lem:4-gramfactor}, assertion 1, since a subspace of $\re^8$ has
  dimension at most 8 (\path{Submodule.finrank_le}).~\qed\end{proof}
 \step{s44small}{\case{$n\le17$.}\prove{$n\le28$ and $d\le16$.}}
 \begin{proof}\pf\ $d\le n-1\le16$ by \asm[1]{0} and \asm{1}.~\qed\end{proof}
 \step{s44q}{$M$ is the Gram matrix of vectors $\Zq_u\in\re^7$, with
  $(\Zq_u,\Zq_u)=1-t$, $(\Zq_u,\Zq_w)=\frac12-t$ for adjacent $u,w$, and $-t$ otherwise.}
 \begin{proof}\pf\ By Lemma~\ref{lem:4-regmat}, assertion 2, with \asm[1]{0} and
  \stepref{s44spec}, $M$ is positive semidefinite of rank $\rk Q-1\le7$. Apply
  Lemma~\ref{lem:4-gramfactor}, assertion 2, with $r=7$; the entries are those of
  $\frac12Q-tJ$.~\qed\end{proof}
 \step{s44P}{$P=(1-t)^2I+(\frac12-t)^2A+t^2(J-I-A)$ is positive semidefinite, and
  $P\Zz=(1-2t)\Zz+(\frac14-t)A\Zz$ for $\Zz\perp\Zj$.}
 \begin{proof}\pf\ By Lemma~\ref{lem:4-symsq} with $\ell=7$ and \stepref{s44q}, $P$ is the Gram
  matrix of the $\sigma(\Zq_u)\in\re^{28}$; apply \path{Matrix.posSemidef_gram}. For
  $\Zz\perp\Zj$, $J\Zz=0$ by Lemma~\ref{lem:4-regmat}, assertion 1, and $(1-t)^2-t^2=1-2t$,
  $(\frac12-t)^2-t^2=\frac14-t$.~\qed\end{proof}
 \step{s44big}{\case{$n\ge18$ and $n\ge2d+4$.}\prove{$n\le28$ and $d\le12$.}}
 \begin{proof}
  \step{s44bigpd}{$P$ is positive definite.}
  \begin{proof}\pf\ For $\Zz\perp\Zj$, $(\Zz,A\Zz)=(\Zz,Q\Zz)-2|\Zz|^2\ge-2|\Zz|^2$ by
   \stepref{s44spec}, and $\frac14-t\ge0$ by \asm{1}; so by \stepref{s44P},
   $(\Zz,P\Zz)\ge(1-2t-2(\frac14-t))|\Zz|^2=\frac12|\Zz|^2$. Further $P\Zj=p\Zj$ with
   $p=(1-t)^2+(\frac12-t)^2d+t^2(n-1-d)>0$, as $1-t\ge\frac34$ since $t\le\frac14$. $P$ is
   symmetric, being a combination of $I$, $A$ and $J$; so $(P\Zz,\Zj)=(\Zz,P\Zj)=0$ for
   $\Zz\perp\Zj$, and $P$ maps $\Zj^\perp$ into itself. So
   $(\Zy,P\Zy)=pc^2n+(\Zz,P\Zz)>0$ for $\Zy=c\Zj+\Zz\ne0$.~\qed\end{proof}
  \qedstep
  \begin{proof}\pf\ The $n$ vectors $\sigma(\Zq_u)$ are linearly independent in $\re^{28}$ by
   \stepref{s44bigpd} and \path{Matrix.posDef_gram_iff_linearIndependent}, so $n\le28$ by
   \path{LinearIndependent.fintype_card_le_finrank}; and $d\le(n-4)/2\le12$ by
   \asm{1}.~\qed\end{proof}
 \end{proof}
 \step{s44rest}{\case{$n\ge18$ and $n<2d+4$.}\prove{$n\le28$ and $d\le16$.}}
 \begin{proof}
  \pflet{$B=2(2n-2d-4)/(2d+4-n)$, and $\Zb_k=(\sqrt2\,\Zq_u(k))_{u\in V(X)}\in\re^{V(X)}$ for
   $1\le k\le7$, with the $\Zq_u$ of \stepref{s44q}.}
  \step{s44bound}{$(\Zz,Q\Zz)\le(B+2)|\Zz|^2$ for every $\Zz\perp\Zj$.}
  \begin{proof}\pf\ By \stepref{s44P}, $0\le(\Zz,P\Zz)=(1-2t)|\Zz|^2+(\frac14-t)(\Zz,A\Zz)$,
   where $\frac14-t=(n-2d-4)/(4n)<0$ by \asm{1}. Divide by it: $(\Zz,A\Zz)\le
   (1-2t)/(t-\frac14)\cdot|\Zz|^2=B|\Zz|^2$, since $1-2t=(n-d-2)/n$; add
   $2|\Zz|^2$.~\qed\end{proof}
  \step{s44cols}{$2M=\sum_{k=1}^7\Zb_k\Zb_k^T$, and $\Zb_k\perp\Zj$ for every $k$.}
  \begin{proof}\pf\ The $uw$-entry of $2M$ is $2(\Zq_u,\Zq_w)=\sum_k\Zb_k(u)\Zb_k(w)$ by
   \stepref{s44q}; and $\sum_k(\Zb_k,\Zj)^2=(\Zj,2M\Zj)=0$ by Lemma~\ref{lem:4-regmat},
   assertion 2.~\qed\end{proof}
  \step{s44ck}{$|\Zb_k|^2\le B+2$ for every $k$.}
  \begin{proof}
   \step{s44ckM}{$(\Zb_k,2M\Zb_k)=\sum_l(\Zb_l,\Zb_k)^2\ge|\Zb_k|^4$.}
   \begin{proof}\pf\ \stepref{s44cols}; drop the terms $l\ne k$.~\qed\end{proof}
   \step{s44ckQ}{$(\Zb_k,2M\Zb_k)=(\Zb_k,Q\Zb_k)\le(B+2)|\Zb_k|^2$.}
   \begin{proof}\pf\ $2M\Zb_k=Q\Zb_k$ by Lemma~\ref{lem:4-regmat}, assertion 2, as $\Zb_k\perp\Zj$ by
    \stepref{s44cols}; then \stepref{s44bound}.~\qed\end{proof}
   \qedstep
   \begin{proof}\pf\ If $\Zb_k=0$ there is nothing to prove. Otherwise divide
    $|\Zb_k|^4\le(B+2)|\Zb_k|^2$, from \stepref{s44ckM} and \stepref{s44ckQ}, by
    $|\Zb_k|^2>0$.~\qed\end{proof}
  \end{proof}
  \step{s44trace}{$0\le(d-2n+16)+7B$.}
  \begin{proof}\pf\ The diagonal entries of $2M$ are $2(1-t)$ by \stepref{s44q}, so
   $\sum_k|\Zb_k|^2=\tr(2M)=2n-d-2$ by \stepref{s44cols}; this is at most $7(B+2)$ by
   \stepref{s44ck}. \new\ The Lean proof of this step and of
   \stepref{s44bound}, in \path{LineGraphs/Representation/Theorem44.lean}, is shorter: it
   bounds the nonzero eigenvalues of $M$, at most 7 of them, through the positive semidefinite
   $P$ of \stepref{s44P}, in place of the columns $\Zb_k$.~\qed\end{proof}
  \step{s44equi}{$E$ is the Gram matrix of $n$ unit vectors $\Zw_u$ in
   $\re^8$ with $(\Zw_u,\Zw_w)=\pm\frac13$ for $u\ne w$.}
  \begin{proof}\pf\ By Lemma~\ref{lem:4-regmat}, assertion 3, with \asm[1]{0},
   \stepref{s44spec} and \asm{1}, $E$ is positive semidefinite of rank $\rk Q\le8$. Apply
   Lemma~\ref{lem:4-gramfactor}, assertion 2, with $r=8$; the entries of $E$ are 1 on the
   diagonal and $\frac23-\frac13$ or $-\frac13$ elsewhere.~\qed\end{proof}
  \step{s44n}{$n\le28$.}
  \begin{proof}\pf\ The span of the $\Zw_u$ of \stepref{s44equi} has dimension $k=\rk E\le8$
   by Lemma~\ref{lem:4-gramfactor}, assertion 1. If $k=8$, then $n\le14$ by Lemma~\ref{lem:4-EQ2}; otherwise $k\le7$ and $n\le28$ by
   Lemma~\ref{lem:4-EQ1}.~\qed\end{proof}
  \step{s44C8}{$d\le16$, or $n=28$ and $d=17$.}
  \begin{proof}\pf\ Computation~\ref{comp:C8}, with \stepref{s44trace}, \stepref{s44n} and
   \asm{1}.~\qed\end{proof}
  \step{s44rankE}{\assume{$n=28$.}\prove{The $\Zw_u$ of \stepref{s44equi} span a space of
   dimension at most 7.}}
  \begin{proof}\pf\ \repair\ This is the split at $d=12$ of Error 15 (errata): here $28<2d+4$
   by \asm{1} and \asm{2}, so $d>12$, and the case $d=12$ is in \stepref{s44big}. The span lies
   in $\re^8$, so its dimension is at most 8 by \path{Submodule.finrank_le}; it is not 8, since
   then Lemma~\ref{lem:4-EQ2} would give $n\le14$, against \asm{2}.~\qed\end{proof}
  \step{s44rankQ}{\assume{$n=28$.}\prove{$\rk Q\le7$.}}
  \begin{proof}\pf\ $\rk Q=\rk E$ by Lemma~\ref{lem:4-regmat}, assertion 3, with \asm[1]{0},
   \stepref{s44spec} and \asm{1}; and $\rk E$ is the dimension of the span of the $\Zw_u$, by
   Lemma~\ref{lem:4-gramfactor}, assertion 1, and \stepref{s44equi}. Apply \stepref{s44rankE}
   with \asm{2}.~\qed\end{proof}
  \step{s44d}{\assume{$n=28$.}\prove{$d\le16$.}}
  \begin{proof}
   \new\ The Lean proof, \path{SimpleGraph.not_regular_28_17_of_rank_le_seven}, closes the case
   $(n,d)=(28,17)$ left by \stepref{s44C8} without the components: there $t=\frac{19}{56}$, and
   $M$ has the entries $\frac{37}{56}$ on the diagonal, $\frac9{56}$ at the 17 neighbours and
   $-\frac{19}{56}$ at the other 10 vertices of each row, so $\tr M=\frac{37}2$ and
   $\sum_{u,w}M(u,w)^2=\frac{227}4$. As $\rk M=\rk Q-1\le6$ by Lemma~\ref{lem:4-regmat},
   assertion 2, and \stepref{s44rankQ}, Lemma~\ref{lem:4-tracerank} would give
   $(\frac{37}2)^2=342.25\le6\cdot\frac{227}4=340.5$, which is false. This is the argument of the
   remark after the theorem, and it needs neither Lemma~\ref{lem:4-rootsystem} nor
   Computation~\ref{comp:C20} here; we keep the route of the paper below.
   \pflet{$u$ be a vertex of $X$, $X_u$ the connected component of $X$ that contains $u$, and
    $k$ the dimension of the span of the $\Zx_w$ with $w\in V(X_u)$.}
   \step{s44dsub}{$X_u$ is connected, the $\Zx_w$ with $w\in V(X_u)$ represent it in $\re^8$, and
    $u$ has $d$ neighbours in $X_u$.}
   \begin{proof}\pf\ A component is connected and is the subgraph induced on its vertices, so
    Lemma~\ref{lem:4-sub} applies to \asm[2]{0}. The component contains every neighbour of $u$,
    and $u$ has $d$ neighbours by \asm[1]{0}.~\qed\end{proof}
   \step{s44ddim}{$k\le7$.}
   \begin{proof}\pf\ The span of the $\Zx_w$ with $w\in V(X_u)$ lies in the span of all the
    $\Zx_w$ by \path{Submodule.span_mono}, and the latter has dimension $\rk Q$ by
    Lemma~\ref{lem:4-gramfactor}, assertion 1; so $k\le\rk Q$ by \path{Submodule.finrank_mono},
    and $\rk Q\le7$ by \stepref{s44rankQ} with \asm{2}.~\qed\end{proof}
   \step{s44dtype}{$X_u$ is represented in the model of $\typ A_k$, of $\typ D_k$ with $k\ge4$,
    of $\typ E_6$ or of $\typ E_7$.}
   \begin{proof}\pf\ Lemma~\ref{lem:4-rootsystem}, assertion 1, applies by \stepref{s44dsub};
    the model of $\typ E_8$ is excluded, since it has rank $8>k$ by Lemma~\ref{lem:sec3-spans}
    and \stepref{s44ddim}.~\qed\end{proof}
   \qedstep
   \begin{proof}\pf\ By Computation~\ref{comp:C20} and \stepref{s44dtype}, $u$ has at most
    $k-1$, $2(k-2)$, 10 or 16 neighbours in $X_u$; these are at most 16 by \stepref{s44ddim}.
    Apply \stepref{s44dsub}.~\qed\end{proof}
  \end{proof}
  \qedstep
  \begin{proof}\pf\ $n\le28$ by \stepref{s44n}. If $d>16$, then $n=28$ by \stepref{s44C8}, and
   \stepref{s44d} gives $d\le16$; so $d\le16$.~\qed\end{proof}
 \end{proof}
 \qedstep
 \begin{proof}\pf\ The cases \stepref{s44small}, \stepref{s44big} and \stepref{s44rest} are
  exhaustive.~\qed\end{proof}
\end{proof}

\begin{namedresult}{Computation}{C24}\label{comp:C24}
 \lean{Computations.C24.trace_sq_gt}
 \uses{}
 \assume{$X$ is a 16-regular graph on 28 vertices, and $N=2I+A(X)-\frac9{14}J$.}
 \prove{$\tr N=38$, $\sum_{u,w}N(u,w)^2=236$, and $38^2>6\cdot236$.}
 \new\ Section 12 of \path{cgss/prerequisites.md} states C24 as $(2m-16)^2>(27-m)(192-4m)$
 for $21\le m\le27$, where $m$ is the
 multiplicity of the eigenvalue $-2$. We will count entries in its place, which needs no
 multiplicities. With Lemma (trace and rank), C24 gives only
 $\rk N\ge7$. The remark after Theorem 4.4 gets $\rk(2I+A)=8$ from this, from assertion 2 of
 Lemma (the matrices of a regular graph), which gives $\rk(2I+A)=\rk N+1$, and from the bound
 $\rk(2I+A)\le8$ of a representation in $\typ E_8$.
\end{namedresult}
\begin{proof}
 \pfsketch\ We will count the entries of each kind; the Lean proof counts them with the degrees
 and closes with \texttt{norm\_num}.
 \step{s4c24row}{Each row of $N$ has the entry $\frac{19}{14}$ on the diagonal, 16 entries
  $\frac5{14}$ and 11 entries $-\frac9{14}$.}
 \begin{proof}\pf\ The entries of $2I+A(X)$ in a row are 2 on the diagonal, 1 at the 16
  neighbours by \asm{0}, and 0 at the other 11 vertices; subtract $\frac9{14}$.~\qed\end{proof}
 \qedstep
 \begin{proof}\pf\ By \stepref{s4c24row}, $\tr N=28\cdot\frac{19}{14}=38$ and
  $\sum_{u,w}N(u,w)^2=\frac{28}{196}(361+16\cdot25+11\cdot81)=\frac{1652}7=236$; and
  $38^2=1444>1416=6\cdot236$.~\qed\end{proof}
\end{proof}

\begin{namedresult}{Remark}{(after Theorem 4.4)}\label{rem:4-after44}
 \lean{SimpleGraph.not_regular_28_16_isRepresentedIn_E₈}
 \lean{SimpleGraph.regular_28_12_and_27_16_isRepresentedIn_E₈}
 \uses{def:represented,lem:4-gramfactor,lem:4-regmat,lem:4-tracerank,lem:4-EQ2,comp:C24,lem:4-linegraph,def:glg,comp:C10,def:3.2,thm:4.4,ex:4.5,ex:4.6,facts:sec4}
 \assume{$X$ is a graph.}
 \prove{\begin{pfenum}
  \item If $X$ is 16-regular on 28 vertices, then $X$ is not represented in $\typ E_8$.
  \item The bound $n\le28$ of Theorem~\ref{thm:4.4} is attained by $L(K_8)$ and by the Chang
   graphs I, II and III, which are 12-regular on 28 vertices, and the bound $d\le16$ by the
   Schl\"afli graph, which is 16-regular on 27 vertices; all five are represented in
   $\typ E_8$.
 \end{pfenum}}
\end{namedresult}
\begin{proof}
 \pfsketch\ \repair\ The paper gives no proof of assertion 1; we will follow EQ3 of
 \path{cgss/prerequisites.md}, the repair of Gap 7 (errata). The rank of $2I+A$ is 8, so the 28
 equiangular lines of the proof of Theorem 4.4 span 8 dimensions, which Lemma EQ2 forbids.
 \new\ We will bound the rank by Lemma~\ref{lem:4-tracerank} and Computation~\ref{comp:C24}, in
 place of the count of eigenvalues in EQ3. Assertion 2 cites the examples below.
 \new\ The Lean proof, \leancite{SimpleGraph.not_regular_28_16_isRepresentedIn_E₈_of_classification}{\texttt{SimpleGraph.\allowbreak not\_\allowbreak regular\_\allowbreak 28\_\allowbreak 16\_\allowbreak isRepresentedIn\_\allowbreak E\ensuremath{{}_8}\_\allowbreak of\_\allowbreak classification}},
 runs in the other order: EQ2 first, which gives $\rk Q\le7$
 (\path{SimpleGraph.IsRepresentedIn.rank_le_seven_of_regular_of_classification}, the argument of
 the proof of Theorem~\ref{thm:4.4} for $n=28$), so that $\rk N\le6$; then C24 and
 Lemma~\ref{lem:4-tracerank} give $38^2\le6\cdot236$, which is false. It carries the same
 classification hypothesis as Theorem~\ref{thm:4.4}; the settled name, without it, is
 proved in \path{LineGraphs/Representation/Theorem44Unconditional.lean} in the same way.
 \step{s4r44one}{\assume{\begin{pfenum}
   \item $X$ is 16-regular on 28 vertices.
   \item $X$ is represented in $\typ E_8$ by $(\Zx_u)$.
  \end{pfenum}}\prove{FALSE.}}
 \begin{proof}
  \pflet{$Q=2I+A(X)$, $M=\frac12Q-\frac9{28}J$ and $E=\frac23Q-\frac13J$, the matrices of
   Lemma~\ref{lem:4-regmat} for $d=16$ and $n=28$, where $t=\frac{18}{56}=\frac9{28}$.}
  \step{s4r44psd}{$Q$ is positive semidefinite of rank at most 8.}
  \begin{proof}\pf\ $Q$ is the Gram matrix of the $\Zx_u\in\re^8$ by \asm[2]{1}, so it is
   positive semidefinite by \path{Matrix.posSemidef_gram}, and of rank at most 8 by
   Lemma~\ref{lem:4-gramfactor}, assertion 1, and \path{Submodule.finrank_le}.~\qed\end{proof}
  \step{s4r44M}{$\rk M\ge7$.}
  \begin{proof}\pf\ $2M$ is the matrix $N$ of Computation~\ref{comp:C24}, by \asm[1]{1}. By
   Lemma~\ref{lem:4-tracerank} and that computation, $38^2\le\rk(2M)\cdot236$ and
   $38^2>6\cdot236$; so $\rk M=\rk(2M)\ge7$.~\qed\end{proof}
  \step{s4r44rk}{$\rk Q=8$.}
  \begin{proof}\pf\ $\rk Q=\rk M+1\ge8$ by Lemma~\ref{lem:4-regmat}, assertion 2, with
   \asm[1]{1}, \stepref{s4r44psd} and \stepref{s4r44M}; and $\rk Q\le8$ by
   \stepref{s4r44psd}.~\qed\end{proof}
  \step{s4r44E}{$E$ is positive semidefinite of rank 8.}
  \begin{proof}\pf\ As $28<2\cdot16+4$, Lemma~\ref{lem:4-regmat}, assertion 3, applies with
   \asm[1]{1} and \stepref{s4r44psd}; it makes $E$ positive semidefinite with $\rk E=\rk Q$, and
   $\rk Q=8$ by \stepref{s4r44rk}.~\qed\end{proof}
  \step{s4r44w}{$E$ is the Gram matrix of 28 unit vectors $\Zw_u$ of $\re^8$ with
   $(\Zw_u,\Zw_w)=\pm\frac13$ for $u\ne w$.}
  \begin{proof}\pf\ Lemma~\ref{lem:4-gramfactor}, assertion 2 with $r=8$, applies by
   \stepref{s4r44E}; the entries of $E$ are 1 on the diagonal and $\frac23-\frac13$ or
   $-\frac13$ elsewhere.~\qed\end{proof}
  \step{s4r44span}{The $\Zw_u$ of \stepref{s4r44w} span a space of dimension 8.}
  \begin{proof}\pf\ Its dimension is $\rk E$ by Lemma~\ref{lem:4-gramfactor}, assertion 1, and
   \stepref{s4r44w}; and $\rk E=8$ by \stepref{s4r44E}.~\qed\end{proof}
  \qedstep
  \begin{proof}\pf\ Lemma~\ref{lem:4-EQ2}, for the span of \stepref{s4r44span} and the $\Zw_u$
   of \stepref{s4r44w}, gives $28\le14$.~\qed\end{proof}
 \end{proof}
 \step{s4r44two}{Assertion 2 holds.}
 \begin{proof}
  \step{s4r44K8}{$L(K_8)$ is 12-regular on 28 vertices and is represented in $\typ E_8$.}
  \begin{proof}\pf\ $K_8$ has $\binom82=28$ edges, and each has degree $7+7-2=12$ in
   $L(K_8)\cong L(K_8;0,\dots,0)$, by Computation~\ref{comp:C10} and Definition~\ref{def:glg}.
   Lemma~\ref{lem:4-linegraph} gives the representation by the $\Ze_i+\Ze_j$, which lie in
   $\Phi(\typ D_8)\sbs\Phi(\typ E_8)$ by Definition~\ref{def:3.2}, (ii).~\qed\end{proof}
  \qedstep
  \begin{proof}\pf\ Example~\ref{ex:4.5} gives the Chang graphs, with parameters
   $(28,12,6,4)$, and Example~\ref{ex:4.6} the Schl\"afli graph, with parameters $(27,16,10,8)$,
   each represented in $\typ E_8$; add \stepref{s4r44K8}.~\qed\end{proof}
 \end{proof}
 \qedstep
 \begin{proof}\pf\ \stepref{s4r44one} and \stepref{s4r44two}.~\qed\end{proof}
\end{proof}

The paper introduces the graphs of Examples 4.5 to 4.9 as strongly regular graphs ``which are
not line graphs''. We will prove this in the proof of Remark 4.14, as Corollary (the seven
graphs are not line graphs), since the argument uses the properties of the seven graphs that we
establish there.
We will use model (ii) of Definition 3.2 and the vector
$\Zc=\frac12\sum_{k=1}^8\Ze_k$ of $\Phi(\typ E_8)$. For a graph $X$ and $Y\sbs V(X)$, the
\emph{switched graph} $X^Y$ keeps the pairs inside $Y$ and inside $V(X)\setminus Y$, and
complements the pairs between them (SW1 of \path{cgss/prerequisites.md}).

\begin{namedresult}{Lemma}{(switching by vectors, SW4)}\label{lem:4-switch}
 \lean{SimpleGraph.IsRepresentation.switch}
 \uses{def:2.1,def:represented,def:line-system}
 \assume{\begin{pfenum}
  \item $X$ is represented in a star-closed line system $S$ by $(\Zx_u)$.
  \item $\Zh\in\Phi(S)$, $(\Zh,\Zx_u)=1$ for every $u$, and $Y\sbs V(X)$.
 \end{pfenum}}
 \prove{The family with $\Zh-\Zx_u$ for $u\in Y$ and $\Zx_u$ for $u\notin Y$ represents $X^Y$
 in $S$.}
\end{namedresult}
\begin{proof}
 \pfsketch\ We will compute the new inner products.
 \step{s4swmem}{$\Zh-\Zx_u\in\Phi(S)$.}
 \begin{proof}\pf\ $(\Zx_u,-\Zh)=-1$ by \asm[2]{0}, so $\Zx_u-\Zh$ and its negative lie in
  $\Phi(S)$, which is star-closed by \asm[1]{0}.~\qed\end{proof}
 \step{s4swip}{For $u,w\in Y$, $(\Zh-\Zx_u,\Zh-\Zx_w)=(\Zx_u,\Zx_w)$; and for $u\in Y$, $w\notin
  Y$, $(\Zh-\Zx_u,\Zx_w)=1-(\Zx_u,\Zx_w)$.}
 \begin{proof}\pf\ Expand, with $(\Zh,\Zh)=2$ and $(\Zh,\Zx_u)=(\Zh,\Zx_w)=1$ by
  \asm[2]{0}.~\qed\end{proof}
 \qedstep
 \begin{proof}\pf\ By \stepref{s4swip}, the Gram matrix keeps the entries inside $Y$ and outside
  $Y$, and replaces $0,1$ by $1,0$ between them; that is $2I+A(X^Y)$. The vectors lie in $\Phi(S)$
  by \stepref{s4swmem}.~\qed\end{proof}
\end{proof}

\begin{namedresult}{Computation}{C9}\label{comp:C9}
 \lean{SimpleGraph.changI_isSRGWith}\lean{SimpleGraph.changII_isSRGWith}
 \lean{SimpleGraph.changIII_isSRGWith}\lean{SimpleGraph.schlafli_isSRGWith}
 \lean{SimpleGraph.shrikhande_isSRGWith}\lean{SimpleGraph.clebsch_isSRGWith}
 \lean{SimpleGraph.petersen_isSRGWith}
 \lean{SimpleGraph.clebsch.foldedFiveCubeComplIso}\lean{SimpleGraph.schlafli.gq24ComplIso}
 \lean{SimpleGraph.petersen.kneserIso}
 \lean{SimpleGraph.clebsch_isSwitchingEquivalent_shrikhande}
 \lean{SimpleGraph.clebsch_isSwitchingEquivalent_lineGraph_completeBipartite}
 \uses{def:represented,def:414-folded,def:414-gq24}
 \assume{The graphs of Examples 4.5 to 4.9 are given by the integer vectors $2\Zx_u$ of their
 representations. The folded 5-cube $\square_5$ is the graph on the subsets of
 $\{1,\dots,5\}$ of even size, adjacent when their symmetric difference has 4 elements, and the
 double-six graph $Q_{27}$ is the graph on the symbols $a_i$, $b_i$ and $c_{ij}=c_{ji}$
 ($i\ne j$ in $\{1,\dots,6\}$) with $a_i\sim b_j$ for $i\ne j$, $a_i\sim c_{jl}$ and
 $b_i\sim c_{jl}$ for $i\in\{j,l\}$, and $c_{ij}\sim c_{lm}$ for disjoint $\{i,j\}$, $\{l,m\}$;
 these are \hyperref[def:414-folded]{the folded 5-cube} and \hyperref[def:414-gq24]{the
 double-six graph} of the proof of Remark 4.14.}
 \prove{The Chang graphs are strongly regular with parameters $(28,12,6,4)$, the Schl\"afli
 graph $(27,16,10,8)$, the Shrikhande graph $(16,6,2,2)$, the Clebsch graph $(16,10,6,6)$ and
 the Petersen graph of Example 4.9 $(10,3,0,1)$; the last is isomorphic to the Kneser graph of
 the 2-subsets of a 5-set; the Clebsch graph is switching equivalent to $L(K_{4,4})$ and to the
 Shrikhande graph. The Clebsch graph is isomorphic to $\comp{\square_5}$, and the Schl\"afli
 graph to $\comp{Q_{27}}$.}
\end{namedresult}
\begin{proof}
 \pfsketch\ These are finite checks on at most 28 vertices; the Lean proofs are by
 \texttt{decide} on the adjacency matrices $A(u,w)=\frac14(2\Zx_u,2\Zx_w)$, with explicit
 isomorphisms and switching sets as certificates. The two isomorphisms of the last sentence
 are the certificates of \path{cgss/proofs/clebsch/clebsch.tex}, Section 6, found by SageMath; in Lean
 they are definitions, checked by \texttt{decide}.
 \qedstep
 \begin{proof}\pf\ For each graph, \texttt{decide} checks the degree of every vertex and the
  number of common neighbours of every pair; for the Petersen graph it checks a given
  bijection, for the Clebsch graph given switching sets and bijections, and for the Clebsch and
  Schl\"afli graphs given bijections onto $\comp{\square_5}$ and $\comp{Q_{27}}$.~\qed\end{proof}
\end{proof}

\begin{cgss}{Example}{4.5}\label{ex:4.5}
 \lean{SimpleGraph.changI_isRepresentedIn_E₈}\lean{SimpleGraph.changII_isRepresentedIn_E₈}
 \lean{SimpleGraph.changIII_isRepresentedIn_E₈}
 \uses{lem:4-linegraph,lem:4-switch,comp:C9,def:3.2,thm:3.5}
 \assume{$T(8)=L(K_8)$ is represented by the vectors $\Ze_i+\Ze_j$, and the switching sets are
 $Y_{\rm I}=\{\Ze_i+\Ze_{4+i}:1\le i\le4\}$, $Y_{\rm II}=\{\Ze_i+\Ze_{i+1}:1\le i\le8\}$ (indices
 modulo 8) and $Y_{\rm III}=\{\Ze_1+\Ze_2,\Ze_2+\Ze_3,\Ze_3+\Ze_1\}\cup\{\Ze_4+\Ze_5,\Ze_5+\Ze_6,
 \Ze_6+\Ze_7,\Ze_7+\Ze_8,\Ze_8+\Ze_4\}$.}
 \prove{Replacing $\Ze_i+\Ze_j$ by $\Zc-\Ze_i-\Ze_j$ on $Y_{\rm I}$, $Y_{\rm II}$ or $Y_{\rm III}$
 gives vectors of $\Phi(\typ E_8)$ representing the Chang graphs I, II and III, which are
 strongly regular with parameters $(28,12,6,4)$.}
\end{cgss}
\begin{proof}
 \pfsketch\ We will switch with $\Zh=\Zc$.
 \step{s45base}{$T(8)$ is represented in $\typ E_8$ by the $\Ze_i+\Ze_j$, and
  $(\Zc,\Ze_i+\Ze_j)=1$.}
 \begin{proof}\pf\ Lemma~\ref{lem:4-linegraph}; the vectors lie in $\Phi(\typ D_8)\sbs
  \Phi(\typ E_8)$ by Definition~\ref{def:3.2}, (ii).~\qed\end{proof}
 \qedstep
 \begin{proof}\pf\ $\typ E_8$ is star-closed by Theorem~\ref{thm:3.5}; apply
  Lemma~\ref{lem:4-switch} to \stepref{s45base}, and Computation~\ref{comp:C9} for the
  parameters.~\qed\end{proof}
\end{proof}

\begin{cgss}{Example}{4.6}\label{ex:4.6}
 \lean{SimpleGraph.schlafli_isRepresentedIn_E₈}
 \uses{def:represented,def:3.2,comp:C9}
 \assume{$\Zx_{ij}=\Ze_i+\Ze_j$ for $1\le i<j\le6$, and $\Zy_i=\Zc-\Ze_i-\Ze_7$ and
 $\Zz_i=\Zc-\Ze_i-\Ze_8$ for $1\le i\le6$.}
 \prove{These 27 vectors lie in $\Phi(\typ E_8)$ and represent the Schl\"afli graph, which is
 strongly regular with parameters $(27,16,10,8)$.}
 \repair\ The paper prints $\frac13\sum\Ze_k$ in the second family, which gives squared norm
 $14/9$; we use $\Zc$ (errata, Error 2).
\end{cgss}
\begin{proof}
 \pfsketch\ We will check membership and the inner products.
 \step{s46mem}{The vectors lie in $\Phi(\typ E_8)$.}
 \begin{proof}\pf\ $\Zx_{ij}\in\Phi(\typ D_8)$; $\Zy_i$ and $\Zz_i$ have entries $\pm\frac12$ with
  exactly two negative ones, so they lie in model (ii) of Definition~\ref{def:3.2}.~\qed\end{proof}
 \step{s46cc}{For 2-subsets $A,B$ of $\{1,\dots,8\}$,
  $(\Zc-\sum_{a\in A}\Ze_a,\Zc-\sum_{b\in B}\Ze_b)=|A\cap B|$ and
  $(\Ze_i+\Ze_j,\Zc-\sum_{b\in B}\Ze_b)=1-|\{i,j\}\cap B|$.}
 \begin{proof}\pf\ Expand, with $(\Zc,\Zc)=2$ and $(\Zc,\Ze_k)=\frac12$.~\qed\end{proof}
 \step{s46ip}{Distinct vectors have inner products in $\{0,1\}$.}
 \begin{proof}\pf\ $(\Zx_{ij},\Zx_{kl})=|\{i,j\}\cap\{k,l\}|$; by \stepref{s46cc},
  $(\Zy_i,\Zy_k)=(\Zz_i,\Zz_k)=1$ for $i\ne k$, $(\Zy_i,\Zz_k)=[i=k]$, and
  $(\Zx_{ij},\Zy_k)=(\Zx_{ij},\Zz_k)=1-[k\in\{i,j\}]$.~\qed\end{proof}
 \qedstep
 \begin{proof}\pf\ By \stepref{s46mem} and \stepref{s46ip} the vectors represent a graph
  (Definition~\ref{def:represented}); Computation~\ref{comp:C9} gives its
  parameters.~\qed\end{proof}
\end{proof}

\begin{cgss}{Example}{4.7}\label{ex:4.7}
 \lean{SimpleGraph.shrikhande_isRepresentedIn_E₈}
 \lean{SimpleGraph.shrikhande.changIIndex}\lean{SimpleGraph.shrikhande.changIIndex_checked}
 \lean{SimpleGraph.shrikhande.changIEmbedding}
 \lean{SimpleGraph.shrikhande.nonempty_iso_induce_changI}
 \uses{lem:4-linegraph,lem:4-switch,lem:4-sub,comp:C9,def:3.2,thm:3.5,ex:4.5}
 \assume{$L(K_{4,4})$ is represented by the vectors $\Ze_i+\Ze_j$ with $1\le i\le4<j\le8$, and
 $Y=\{\Ze_i+\Ze_{4+i}:1\le i\le4\}$.}
 \prove{Replacing $\Ze_i+\Ze_j$ by $\Zc-\Ze_i-\Ze_j$ on $Y$ gives vectors of $\Phi(\typ E_8)$
 representing the Shrikhande graph, strongly regular with parameters $(16,6,2,2)$; it is an
 induced subgraph of the Chang graph I.}
 In Lean the last assertion is \path{SimpleGraph.shrikhande.nonempty_iso_induce_changI}: the
 Shrikhande graph is isomorphic to the subgraph of the Chang graph I induced on the sixteen
 vertices given by \path{SimpleGraph.shrikhande.changIIndex}, an induced embedding that the
 kernel checks (\texttt{decide +kernel}).
\end{cgss}
\begin{proof}
 \pfsketch\ We will switch with $\Zh=\Zc$, as in Example 4.5.
 \step{s47rep}{The switched vectors lie in $\Phi(\typ E_8)$ and represent the Shrikhande graph.}
 \begin{proof}\pf\ Lemma~\ref{lem:4-linegraph} for $K_{4,4}$ gives a representation in
  $\Phi(\typ D_8)\sbs\Phi(\typ E_8)$ (Definition~\ref{def:3.2}) with $(\Zc,\Ze_i+\Ze_j)=1$; $\typ
  E_8$ is star-closed by Theorem~\ref{thm:3.5}; apply Lemma~\ref{lem:4-switch}.~\qed\end{proof}
 \qedstep
 \begin{proof}\pf\ Computation~\ref{comp:C9} gives the parameters. The switched vectors are among
  those of the Chang graph I in Example~\ref{ex:4.5}, since $Y=Y_{\rm I}$, and a subfamily of a
  representation represents the induced subgraph (Lemma~\ref{lem:4-sub}).~\qed\end{proof}
\end{proof}

\begin{cgss}{Example}{4.8}\label{ex:4.8}
 \lean{SimpleGraph.clebsch_isRepresentedIn_E₈}
 \lean{SimpleGraph.clebsch.schlafliIndex}\lean{SimpleGraph.clebsch.schlafliIndex_checked}
 \lean{SimpleGraph.clebsch.schlafliEmbedding}
 \lean{SimpleGraph.clebsch.nonempty_iso_induce_schlafli}
 \uses{def:represented,lem:4-sub,def:3.2,comp:C9,ex:4.6}
 \assume{The vectors are $\Ze_i+\Ze_5$, $\Ze_i+\Ze_6$, $\Zc-\Ze_i-\Ze_7$ and $\Zc-\Ze_i-\Ze_8$,
 for $1\le i\le4$.}
 \prove{These 16 vectors are the vectors of Example 4.6 with inner product 1 with $\Ze_5+\Ze_6$;
 they represent the Clebsch graph, strongly regular with parameters $(16,10,6,6)$ and switching
 equivalent to $L(K_{4,4})$ and to the Shrikhande graph.}
 In Lean the first two assertions together are
 \path{SimpleGraph.clebsch.nonempty_iso_induce_schlafli}: the Clebsch graph is isomorphic to the
 subgraph of the Schl\"afli graph induced on the neighbours of the vertex whose vector is
 $\Ze_5+\Ze_6$, which are the vertices whose vectors have inner product 1 with it. The kernel
 checks that \path{SimpleGraph.clebsch.schlafliIndex} is an induced embedding onto these
 neighbours.
\end{cgss}
\begin{proof}
 \pfsketch\ We will read off the neighbours of $\Ze_5+\Ze_6$ in Example 4.6.
 \step{s48nbr}{The vectors of Example~\ref{ex:4.6} with inner product 1 with $\Ze_5+\Ze_6$ are
  the 16 listed.}
 \begin{proof}\pf\ Expand: $(\Ze_5+\Ze_6,\Ze_k+\Ze_l)=|\{5,6\}\cap\{k,l\}|$, which is 1 exactly for $\{k,l\}=\{i,5\}$ or
  $\{i,6\}$ with $i\le4$; and $(\Ze_5+\Ze_6,\Zc-\Ze_k-\Ze_7)=1-|\{5,6\}\cap\{k\}|$, which is 1
  exactly for $k\le4$, and the same for $\Ze_8$.~\qed\end{proof}
 \qedstep
 \begin{proof}\pf\ A subfamily of the representation of Example~\ref{ex:4.6} represents the
  induced subgraph (Lemma~\ref{lem:4-sub}); Computation~\ref{comp:C9} gives the
  parameters and the switching equivalences.~\qed\end{proof}
\end{proof}

\begin{cgss}{Example}{4.9}\label{ex:4.9}
 \lean{SimpleGraph.petersen_isRepresentedIn_E₈}
 \uses{lem:4-linegraph,lem:4-switch,comp:C9,def:3.2,thm:3.5}
 \assume{$L(K_{3,3})$ together with an isolated vertex is represented by $\Ze_4+\Ze_8$ and the
 $\Ze_i+\Ze_j$ with $1\le i\le3$ and $5\le j\le7$, and $Y=\{\Ze_i+\Ze_{4+i}:1\le i\le3\}$.}
 \prove{Replacing $\Ze_i+\Ze_j$ by $\Zc-\Ze_i-\Ze_j$ on $Y$ gives vectors of $\Phi(\typ E_8)$
 representing the Petersen graph, strongly regular with parameters $(10,3,0,1)$.}
\end{cgss}
\begin{proof}
 \pfsketch\ We will switch with $\Zh=\Zc$, as in Example 4.5.
 \step{s49base}{The ten vectors represent $L(K_{3,3})$ with an isolated vertex, in
  $\Phi(\typ E_8)$, and each has inner product 1 with $\Zc$.}
 \begin{proof}\pf\ Lemma~\ref{lem:4-linegraph} for the graph with the edges $\{i,j\}$ and
  $\{4,8\}$, whose line graph is $L(K_{3,3})$ with an isolated vertex; the vectors lie in
  $\Phi(\typ D_8)\sbs\Phi(\typ E_8)$ by Definition~\ref{def:3.2}.~\qed\end{proof}
 \qedstep
 \begin{proof}\pf\ $\typ E_8$ is star-closed by Theorem~\ref{thm:3.5}; apply
  Lemma~\ref{lem:4-switch} to \stepref{s49base}, and Computation~\ref{comp:C9} for the
  parameters and the isomorphism with the Petersen graph.~\qed\end{proof}
\end{proof}

\subsection{Theorems 4.10 and 4.11, and Remark 4.12}

\begin{cgss}{Theorem}{4.10}\label{thm:4.10}
 \lean{SimpleGraph.hoffman}
 \uses{lem:4-gram,lem:4-rootsystem,lem:4-transport,thm:4.2,def:3.1,def:glg,def:represented}
 \assume{$X$ is a connected graph, and $2I+A(X)$ is positive semidefinite.}
 \prove{$X$ is a generalized line graph, or $X$ is represented in $\typ E_8$.}
 \repair\ The paper does not assume that $X$ is connected, and the theorem is false without it:
 two disjoint Petersen graphs are neither (errata, Implicit hypothesis 1). We assume only that
 the least eigenvalue is at least $-2$; the paper assumes that it is $-2$.
\end{cgss}
\begin{proof}
 \pfsketch\ ``This follows from the Theorems 3.6 and 4.2, and the inclusion $\typ
 A_n\subseteq\typ D_{n+1}$.''
 \step{s410rep}{$X$ is represented in $\typ A_\ell$ ($\ell\ge1$), in $\typ D_\ell$
  ($\ell\ge4$), or in $\typ E_8$.}
 \begin{proof}\pf\ $X$ has a representation in $\re^{|V(X)|}$ by Lemma~\ref{lem:4-gram} and
  \asm{0}; apply Lemma~\ref{lem:4-rootsystem}, assertion 2.~\qed\end{proof}
 \step{s410A}{\case{$X$ is represented in $\typ A_\ell$.}\prove{$X$ is a generalized line
  graph.}}
 \begin{proof}\pf\ $\Phi(\typ A_\ell)\sbs\Phi(\typ D_{\ell+1})$ by Definition~\ref{def:3.1} and
  its remark on the coordinate model for small $\ell$, so
  $X$ is represented in $\typ D_{\ell+1}$ by Lemma~\ref{lem:4-transport}; apply
  Theorem~\ref{thm:4.2}, assertion 1.~\qed\end{proof}
 \step{s410D}{\case{$X$ is represented in $\typ D_\ell$.}\prove{$X$ is a generalized line
  graph.}}
 \begin{proof}\pf\ Theorem~\ref{thm:4.2}, assertion 1.~\qed\end{proof}
 \qedstep
 \begin{proof}\pf\ The cases of \stepref{s410rep} are exhaustive, and in the third $X$ is
  represented in $\typ E_8$.~\qed\end{proof}
\end{proof}

\begin{namedresult}{Computation}{C10}\label{comp:C10}
 \lean{SimpleGraph.generalizedLineGraph_degree_edge}
 \lean{SimpleGraph.generalizedLineGraph_degree_petal}
 \uses{def:glg}
 \assume{$Y$ is a graph on $\{1,\dots,m\}$, $\deg(i)$ is the degree of $i$ in $Y$, and
 $a_1,\dots,a_m\ge0$ are integers.}
 \prove{In $L(Y;a_1,\dots,a_m)$ the edge $\{i,j\}$ has degree $2(a_i+a_j)+\deg(i)+\deg(j)-2$, and the
  vertex $(i,\alpha,s)$ has degree $2(a_i-1)+\deg(i)$.}
\end{namedresult}
\begin{proof}
 \pfsketch\ We will count the neighbours of each kind; the Lean proof is parametric.
 \step{s4c10edge}{The edge $\{i,j\}$ has $\deg(i)+\deg(j)-2$ neighbours among the edges and
  $2a_i+2a_j$ among the triples.}
 \begin{proof}\pf\ The edges other than $\{i,j\}$ on $i$ or on $j$ number $(\deg(i)-1)+(\deg(j)-1)$,
  since no other edge is on both. The triples adjacent to it are the $(i,\alpha,s)$ and the
  $(j,\beta,t)$, by Definition~\ref{def:glg}.~\qed\end{proof}
 \step{s4c10petal}{$(i,\alpha,s)$ has $\deg(i)$ neighbours among the edges and $2(a_i-1)$ among
  the triples.}
 \begin{proof}\pf\ Its neighbours are the edges on $i$ and the $(i,\beta,t)$ with
  $\beta\ne\alpha$, by Definition~\ref{def:glg}.~\qed\end{proof}
 \qedstep
 \begin{proof}\pf\ Add the two counts in \stepref{s4c10edge} and in
  \stepref{s4c10petal}.~\qed\end{proof}
\end{proof}

\begin{namedresult}{Lemma}{(regular connected generalized line graphs)}\label{lem:4-regglg}
 \lean{SimpleGraph.IsGeneralizedLineGraph.lineGraph_or_cocktailParty}
 \uses{def:glg,def:CP,comp:C10}
 \assume{$X\cong L(Y;a_1,\dots,a_m)$ is connected and regular, with $Y$ on $\{1,\dots,m\}$
 and $\deg(i)$ the degree of $i$ in $Y$.}
 \prove{$X\cong L(Y')$ for some graph $Y'$, or $X\cong CP(a)$ for some $a$.}
\end{namedresult}
\begin{proof}
 \pfsketch\ This is the paper's first paragraph of the proof of Theorem 4.11. \repair\ The paper
 says that $Y$ is connected, which need not hold (errata, Gap 18), since vertices $i$ of $Y$
 with $\deg(i)=0$ and $a_i=0$ do not appear in $X$; we will split on the vertices with $a_i>0$
 instead.
 \step{s4rgzero}{\case{$a_i=0$ for every $i$.}\prove{$X\cong L(Y)$.}}
 \begin{proof}\pf\ $L(Y;0,\dots,0)\cong L(Y)$ by Definition~\ref{def:glg}.~\qed\end{proof}
 \step{s4rgiso}{\case{$a_i>0$ and $\deg(i)=0$ for some $i$.}\prove{$X\cong CP(a_i)$.}}
 \begin{proof}\pf\ No edge of $Y$ contains $i$, so by Definition~\ref{def:glg} the $2a_i\ge2$
  triples $(i,\alpha,s)$ have all their neighbours among themselves. Since $X$ is connected,
  they are all its vertices, and they span a copy of $CP(a_i)$.~\qed\end{proof}
 \step{s4rgedge}{\case{$a_i>0$ and $\deg(i)>0$ for some $i$.}\prove{This case does not occur.}}
 \begin{proof}\pf\ Pick an edge $\{i,j\}$. By regularity and C10 (Computation~\ref{comp:C10}),
  $2(a_i+a_j)+\deg(i)+\deg(j)-2=2(a_i-1)+\deg(i)$, so $2a_j+\deg(j)=0$; but $\deg(j)\ge1$.~\qed\end{proof}
 \qedstep
 \begin{proof}\pf\ Either every $a_i$ is 0, or some $a_i>0$, and then $\deg(i)=0$ or $\deg(i)>0$; so the
  three cases are exhaustive.~\qed\end{proof}
\end{proof}

\begin{cgss}{Theorem}{4.11}\label{thm:4.11}
 \lean{SimpleGraph.hoffman_rayChaudhuri}
 \uses{thm:4.10,lem:4-regglg,def:glg,def:CP,def:represented}
 \assume{$X$ is a connected regular graph, and $2I+A(X)$ is positive semidefinite.}
 \prove{$X\cong L(Y)$ for some graph $Y$, or $X\cong CP(m)$ for some $m$, or $X$ is represented
 in $\typ E_8$.}
 \new\ As in Theorem 4.10, we assume only that the least eigenvalue is at least $-2$.
\end{cgss}
\begin{proof}
 \pfsketch\ We will apply Theorem 4.10, and then the lemma on regular generalized line graphs.
 \step{s411hoff}{$X$ is a generalized line graph, or $X$ is represented in $\typ E_8$.}
 \begin{proof}\pf\ Theorem~\ref{thm:4.10}, with \asm{0}.~\qed\end{proof}
 \step{s411glg}{\case{$X\cong L(Y;a_1,\dots,a_m)$.}\prove{$X\cong L(Y')$ for some $Y'$ or
  $X\cong CP(a)$ for some $a$.}}
 \begin{proof}\pf\ Lemma~\ref{lem:4-regglg}, with \asm{0} and \asm{1}.~\qed\end{proof}
 \qedstep
 \begin{proof}\pf\ The cases of \stepref{s411hoff} are exhaustive, by
  Definition~\ref{def:glg}.~\qed\end{proof}
\end{proof}

\begin{namedresult}{Definition}{(claws)}\label{def:claw}
 \lean{SimpleGraph.IsClaw}\lean{SimpleGraph.IsFourClaw}
 \lean{SimpleGraph.ClawFree}\lean{SimpleGraph.FourClawFree}
 \uses{}
 \assume{$X$ is a graph and $k\ge1$.}
 \pflet{A \emph{$k$-claw} of $X$ with \emph{centre} $u$ is a set of $k$ pairwise non-adjacent
 neighbours of $u$; with $u$, they induce a copy of $K_{1,k}$. The graph $X$ is \emph{claw-free}
 if it has no 3-claw, and \emph{four-claw-free} if it has no 4-claw.}
 In Lean, \texttt{X.IsClaw u T} and \texttt{X.IsFourClaw u T} will say that the finite set $T$ is
 a 3-claw, or a 4-claw, of $X$ with centre $u$. Following \path{cgss/notation.md}, the
 predicates \texttt{ClawFree} and \texttt{FourClawFree} will say instead that there is no induced
 copy of the star graph \path{SimpleGraph.starGraph} on 4 or 5 vertices; the lemma below proves
 that the two readings agree. A 4-claw contains a 3-claw with the same centre.
\end{namedresult}

\begin{namedresult}{Lemma}{(claws and stars)}\label{lem:4-clawstar}
 \lean{SimpleGraph.clawFree_iff_forall_not_isClaw}
 \lean{SimpleGraph.fourClawFree_iff_forall_not_isFourClaw}
 \lean{SimpleGraph.Iso.clawFree_iff}\lean{SimpleGraph.Iso.fourClawFree_iff}
 \uses{def:claw}
 \assume{$X$ is a graph, $k\ge1$, and $K_{1,k}$ is the star graph \path{SimpleGraph.starGraph}
 on $\{0,\dots,k\}$ with centre 0.}
 \prove{\begin{pfenum}
  \item $X$ has an induced copy of $K_{1,k}$ (\path{SimpleGraph.IsIndContained}) if and only if
   $X$ has a $k$-claw.
  \item If $X\cong X'$, then $X$ has an induced copy of $K_{1,k}$ if and only if $X'$ has one.
 \end{pfenum}}
 \new\ The Lean predicates of Definition~\ref{def:claw} forbid an induced star, while our
 arguments count claws; the Lean names are the cases $k=3$ and $k=4$.
\end{namedresult}
\begin{proof}
 \pfsketch\ We will send the centre of the star to the centre of the claw.
 \step{s4csfwd}{\assume{$f$ embeds $K_{1,k}$ into $X$ as an induced subgraph.}
  \prove{$\{f(1),\dots,f(k)\}$ is a $k$-claw of $X$ with centre $f(0)$.}}
 \begin{proof}\pf\ By \path{SimpleGraph.starGraph_adj}, each $l$ with $1\le l\le k$ is adjacent
  to 0 in $K_{1,k}$, and no two of them are adjacent. As $f$ is injective and preserves
  adjacency and non-adjacency by \asm{1}, the $f(l)$ are $k$ distinct neighbours of $f(0)$, and
  no two of them are adjacent.~\qed\end{proof}
 \step{s4csbwd}{\assume{$T=\{t_1,\dots,t_k\}$ is a $k$-claw of $X$ with centre $u$.}
  \prove{The map $f$ with $f(0)=u$ and $f(l)=t_l$ for $1\le l\le k$ embeds $K_{1,k}$ into $X$ as
  an induced subgraph.}}
 \begin{proof}\pf\ $u\notin T$, since $u$ is adjacent to every member of $T$ by \asm{1} and not to
  itself; so $f$ is injective. By \asm{1} and Definition~\ref{def:claw}, $f(0)$ is adjacent to
  every $f(l)$ with $l\ge1$, and no two $f(l)$ with $l\ge1$ are adjacent; these are the
  adjacencies of $K_{1,k}$, by \path{SimpleGraph.starGraph_adj}.~\qed\end{proof}
 \step{s4csiso}{\assume{$e$ is an isomorphism from $X$ to $X'$.}\prove{Assertion 2 holds.}}
 \begin{proof}\pf\ Compose an embedding of $K_{1,k}$ with $e$ or with its inverse, by
  \path{SimpleGraph.Iso.isIndContained}, \path{SimpleGraph.Iso.isIndContained'} and
  \path{SimpleGraph.IsIndContained.trans}.~\qed\end{proof}
 \qedstep
 \begin{proof}\pf\ Assertion 1 is \stepref{s4csfwd} and \stepref{s4csbwd}; assertion 2 is
  \stepref{s4csiso}.~\qed\end{proof}
\end{proof}

\begin{namedresult}{Lemma}{(claws in a generalized line graph)}\label{lem:4-claws}
 \lean{SimpleGraph.generalizedLineGraph_not_isClaw_petal}
 \lean{SimpleGraph.generalizedLineGraph_isClaw_edge_iff}
 \lean{SimpleGraph.generalizedLineGraph_isFourClaw_edge_iff}
 \uses{def:glg,def:claw}
 \assume{$Y$ is a graph on $\{1,\dots,m\}$, $\deg(i)$ is the degree of $i$ in $Y$,
 $a_1,\dots,a_m\ge0$ are integers, and $X=L(Y;a_1,\dots,a_m)$.}
 In the Lean names, as in those of Computation C10, a \emph{petal} of $X$ is a vertex
 $(i,\alpha,s)$ of the cocktail party graph attached at $i$; it is not the petal of the proof of
 Theorem 4.2, which is a coordinate.
 \prove{\begin{pfenum}
  \item No vertex $(i,\alpha,s)$ of $X$ is the centre of a 3-claw.
  \item The edge $\{i,j\}$ of $Y$ is the centre of a 3-claw of $X$ if and only if $a_i\ge1$ and
   $a_j+\deg(j)\ge2$, or $a_j\ge1$ and $a_i+\deg(i)\ge2$.
  \item The edge $\{i,j\}$ of $Y$ is the centre of a 4-claw of $X$ if and only if $a_i\ge1$ and
   $a_j\ge1$.
 \end{pfenum}}
\end{namedresult}
\begin{proof}
 \pfsketch\ We will split the neighbours of an edge into those at its two ends.
 \pflet{For an edge $e=\{i,j\}$ of $Y$, $N_i(e)$ is the set of the edges $f\ne e$ of $Y$ with
 $i\in f$ together with the triples $(i,\alpha,s)$, and $N_j(e)$ is defined in the same way.}
 \step{s4clnbr}{The neighbours of $e=\{i,j\}$ in $X$ are the members of $N_i(e)\cup N_j(e)$,
  and $|N_j(e)|=\deg(j)-1+2a_j$.}
 \begin{proof}\pf\ By Definition~\ref{def:glg}, $e$ is adjacent to the edges that share an end
  with it and to the triples $(l,\alpha,s)$ with $l\in e$. The edges $f\ne e$ on $j$ number
  $\deg(j)-1$, and the triples $(j,\alpha,s)$ number $2a_j$.~\qed\end{proof}
 \step{s4clside}{\assume{$T\sbs N_i(e)$ is a set of pairwise non-adjacent vertices with
  $|T|\ge2$.}\prove{$T=\{(i,\alpha,0),(i,\alpha,1)\}$ for some $\alpha$; in particular
  $a_i\ge1$.}}
 \begin{proof}\pf\ By Definition~\ref{def:glg}, two edges on $i$ are adjacent, a triple
  $(i,\alpha,s)$ is adjacent to every edge on $i$, and $(i,\alpha,s)$ and $(i,\beta,r)$ are
  adjacent when $\alpha\ne\beta$. So $T$ consists of triples with one value of $\alpha$, and
  there are two of these.~\qed\end{proof}
 \step{s4clcross}{A triple $(i,\alpha,s)$ is adjacent to no member of $N_j(e)$.}
 \begin{proof}\pf\ An edge $f\ne e$ on $j$ does not contain $i$, since $e$ is the only edge on
  both $i$ and $j$; and a triple $(j,\beta,r)$ has $j\ne i$. Apply
  Definition~\ref{def:glg}.~\qed\end{proof}
 \step{s4clone}{Assertion 1 holds.}
 \begin{proof}\pf\ By Definition~\ref{def:glg}, the neighbours of $(i,\alpha,s)$ are the edges on
  $i$, each adjacent to all the others, and the triples $(i,\beta,r)$ with $\beta\ne\alpha$. So a
  3-claw would consist of three triples $(i,\beta,r)$ with one value of $\beta$, and there are
  only two of these.~\qed\end{proof}
 \step{s4cltwo}{Assertion 2 holds.}
 \begin{proof}
  \step{s4cltwoif}{\assume{$a_i\ge1$ and $a_j+\deg(j)\ge2$.}\prove{$e=\{i,j\}$ is the centre of a
   3-claw.}}
  \begin{proof}\pf\ Since $\deg(j)\ge1$, \asm{2} gives $\deg(j)-1+2a_j\ge1$, so $N_j(e)$ has a
   member $z$ by \stepref{s4clnbr}. The vertices $(i,1,0)$, $(i,1,1)$ and $z$ are neighbours of
   $e$ by \stepref{s4clnbr}; the first two are not adjacent by Definition~\ref{def:glg}, and
   neither is adjacent to $z$ by \stepref{s4clcross}.~\qed\end{proof}
  \step{s4cltwoonly}{\assume{$T$ is a 3-claw with centre $e=\{i,j\}$.}\prove{$a_i\ge1$ and
   $a_j+\deg(j)\ge2$, or $a_j\ge1$ and $a_i+\deg(i)\ge2$.}}
  \begin{proof}\pf\ $T\sbs N_i(e)\cup N_j(e)$ by \stepref{s4clnbr}, so one of the two meets $T$
   in at least two members; by the symmetry of $i$ and $j$, say $N_i(e)$. By
   \stepref{s4clside}, $a_i\ge1$ and $|T\cap N_i(e)|=2$, so $N_j(e)$ meets $T$. Thus
   $\deg(j)-1+2a_j\ge1$ by \stepref{s4clnbr}, which gives $a_j+\deg(j)\ge2$, as
   $\deg(j)\ge1$.~\qed\end{proof}
  \qedstep
  \begin{proof}\pf\ \stepref{s4cltwoif}, with $i$ and $j$ exchanged for the second alternative,
   and \stepref{s4cltwoonly}.~\qed\end{proof}
 \end{proof}
 \step{s4clthree}{Assertion 3 holds.}
 \begin{proof}
  \step{s4clthreeif}{\assume{$a_i\ge1$ and $a_j\ge1$.}\prove{$e=\{i,j\}$ is the centre of a
   4-claw.}}
  \begin{proof}\pf\ The triples $(i,1,0)$, $(i,1,1)$, $(j,1,0)$ and $(j,1,1)$ are neighbours of $e$
   by \stepref{s4clnbr}. The two at $i$ are not adjacent, nor are the two at $j$, by
   Definition~\ref{def:glg}; and those at $i$ are not adjacent to those at $j$ by
   \stepref{s4clcross}.~\qed\end{proof}
  \step{s4clthreeonly}{\assume{$T$ is a 4-claw with centre $e=\{i,j\}$.}\prove{$a_i\ge1$ and
   $a_j\ge1$.}}
  \begin{proof}\pf\ $T\sbs N_i(e)\cup N_j(e)$ by \stepref{s4clnbr}, and each of $T\cap N_i(e)$ and
   $T\cap N_j(e)$ has at most two members by \stepref{s4clside}. So each has exactly two, and
   \stepref{s4clside} gives $a_i\ge1$ and $a_j\ge1$.~\qed\end{proof}
  \qedstep
  \begin{proof}\pf\ \stepref{s4clthreeif} and \stepref{s4clthreeonly}.~\qed\end{proof}
 \end{proof}
 \qedstep
 \begin{proof}\pf\ \stepref{s4clone}, \stepref{s4cltwo} and \stepref{s4clthree}.~\qed\end{proof}
\end{proof}

\begin{cgss}{Remark}{4.12}\label{rem:4.12}
 \lean{SimpleGraph.IsGeneralizedLineGraph.clawFree_iff}
 \lean{SimpleGraph.generalizedLineGraph_fourClawFree_iff}
 \lean{SimpleGraph.lineGraph_clawFree}\lean{SimpleGraph.cocktailParty_clawFree}
 \lean{SimpleGraph.generalizedLineGraph_clawFree_iff}
 \lean{SimpleGraph.hoffman_of_clawFree}
 \uses{def:claw,lem:4-clawstar,lem:4-claws,def:glg,def:CP,thm:4.3,thm:4.10}
 \assume{$X$ is a graph.}
 \prove{\begin{pfenum}
  \item Line graphs and cocktail party graphs are claw-free.
  \item Let $Y$ be a graph on $\{1,\dots,m\}$, with $\deg(i)$ the degree of $i$ in $Y$, and let
   $a_1,\dots,a_m\ge0$ be integers. Then $L(Y;a_1,\dots,a_m)$ is claw-free if and only if $a_j=0$
   and $\deg(j)=1$ for every vertex $i$ of $Y$ with $a_i\ge1$ and every neighbour $j$ of $i$.
  \item If $X$ is a connected generalized line graph, then $X$ is claw-free if and only if
   $X\cong L(Y')$ for some graph $Y'$ or $X\cong L(K_{1,s};t,0,\dots,0)$ for some $s,t\ge0$,
   where the vertex 1 of $K_{1,s}$ is its centre. The graph $L(K_{1,s};t,0,\dots,0)$ is
   $K_{s+2t}$ with $t$ pairwise disjoint edges deleted.
  \item If $X$ is connected and claw-free, $2I+A(X)$ is positive semidefinite, and $X$ has more
   than 36 vertices or a vertex with more than 28 neighbours, then $X\cong L(Y')$ for some graph
   $Y'$ or $X\cong L(K_{1,s};t,0,\dots,0)$ for some $s,t\ge0$.
  \item For $Y$ and $a_1,\dots,a_m$ as in assertion 2, $L(Y;a_1,\dots,a_m)$ is four-claw-free if
   and only if $a_i=0$ or $a_j=0$ for every edge $\{i,j\}$ of $Y$.
 \end{pfenum}}
 \new\ The paper's $L(K_{1,n};a,0,\dots,0)$ is our $L(K_{1,s};t,0,\dots,0)$, since $n$ and $a$
 are spoken for (\path{cgss/notation.md}), and its ``sufficiently large'' is made precise by
 Theorem 4.3. The remark on Ray-Chaudhuri's theorem is history, and needs no proof.
\end{cgss}
\begin{proof}
 \pfsketch\ \repair\ The paper says ``it is easy to see''; these are the full arguments that
 Gap 9 (errata) asks for. We will read everything off the lemma on claws.
 \step{s412two}{Assertion 2 holds.}
 \begin{proof}\pf\ Claw-freeness is the absence of 3-claws, by Definition~\ref{def:claw} and, for
  the Lean predicate, Lemma~\ref{lem:4-clawstar}, assertion 1. By Lemma~\ref{lem:4-claws},
  assertion 1, only edges of $Y$ can be centres of 3-claws, and by its assertion 2 the edge $\{i,j\}$ is not one exactly when $a_i\ge1$ implies
  $a_j+\deg(j)\le1$, and the same with $i$ and $j$ exchanged. As $\deg(j)\ge1$,
  $a_j+\deg(j)\le1$ means that $a_j=0$ and $\deg(j)=1$.~\qed\end{proof}
 \step{s412one}{Assertion 1 holds.}
 \begin{proof}\pf\ Claw-freeness is preserved by isomorphisms, by
  Lemma~\ref{lem:4-clawstar}, assertion 2. A line graph $L(Y')$, with the
  vertices of $Y'$ numbered, is isomorphic to $L(Y';0,\dots,0)$ by Definition~\ref{def:glg}, and
  the condition of \stepref{s412two} holds for it trivially. By Definitions~\ref{def:CP}
  and~\ref{def:glg}, $CP(t)$ is isomorphic to $L(Y_1;t)$ for the graph $Y_1$ with one vertex and
  no edges, and the condition holds trivially again.~\qed\end{proof}
 \step{s412three}{\assume{$X$ is connected, and $X\cong L(Y;a_1,\dots,a_m)$ for $Y$ and
  $a_1,\dots,a_m$ as in assertion 2.}\prove{$X$ is claw-free if and only if $X\cong L(Y')$ for some graph $Y'$
  or $X\cong L(K_{1,s};t,0,\dots,0)$ for some $s,t\ge0$.}}
 \begin{proof}
  \step{s412ifK}{$L(K_{1,s};t,0,\dots,0)$ is claw-free.}
  \begin{proof}\pf\ The only vertex of $K_{1,s}$ with a nonzero number attached is its centre 1,
   and every neighbour of 1 has number 0 and degree 1; apply \stepref{s412two}.~\qed\end{proof}
  \step{s412only}{\assume{$X$ is claw-free, and $a_i\ge1$ for some $i$.}\prove{$X\cong
   L(K_{1,s};t,0,\dots,0)$ with $s=\deg(i)$ and $t=a_i$.}}
  \begin{proof}
   \pflet{$U$ be the set of the triples $(i,\alpha,r)$ and of the edges of $Y$ on $i$.}
   \step{s412leaf}{Every neighbour $j$ of $i$ in $Y$ has $a_j=0$ and $\deg(j)=1$.}
   \begin{proof}\pf\ $X$ is claw-free by \asm{2} and $X\cong L(Y;a_1,\dots,a_m)$ by \asm{1}, so
    $L(Y;a_1,\dots,a_m)$ is claw-free by Lemma~\ref{lem:4-clawstar}, assertion 2. Apply
    \stepref{s412two}, with \asm{2}.~\qed\end{proof}
   \step{s412closed}{Every neighbour in $L(Y;a_1,\dots,a_m)$ of a member of $U$ lies in $U$.}
   \begin{proof}\pf\ By Definition~\ref{def:glg}, the neighbours of a triple $(i,\alpha,r)$ are
    edges on $i$ and triples at $i$. The neighbours of an edge $\{i,j\}$ are the edges on $i$ or
    on $j$ and the triples at $i$ or at $j$; by \stepref{s412leaf}, the only edge on $j$ is
    $\{i,j\}$, and there are no triples at $j$.~\qed\end{proof}
   \step{s412all}{$U$ is the whole vertex set of $L(Y;a_1,\dots,a_m)$.}
   \begin{proof}\pf\ $U$ is nonempty, as $a_i\ge1$, and closed under adjacency by
    \stepref{s412closed}; and $X$ is connected by \asm{1}.~\qed\end{proof}
   \qedstep
   \begin{proof}\pf\ Number the neighbours of $i$ as $j_1,\dots,j_s$. By \stepref{s412all}, the
    map sending $\{i,j_k\}$ to the edge $\{1,k+1\}$ of $K_{1,s}$, and $(i,\alpha,r)$ to
    $(1,\alpha,r)$, is a bijection onto the vertices of $L(K_{1,s};t,0,\dots,0)$. It preserves
    adjacency by Definition~\ref{def:glg}: in both graphs the edges are pairwise adjacent and
    adjacent to every triple, and two triples are adjacent exactly when their second coordinates
    differ.~\qed\end{proof}
  \end{proof}
  \qedstep
  \begin{proof}\pf\ If $X\cong L(Y')$, then $X$ is claw-free by \stepref{s412one}; if $X\cong
   L(K_{1,s};t,0,\dots,0)$, by \stepref{s412ifK}. Conversely, let $X$ be claw-free. If every
   $a_i$ is 0, then $X\cong L(Y)$ by Definition~\ref{def:glg}; otherwise apply
   \stepref{s412only}.~\qed\end{proof}
 \end{proof}
 \step{s412clique}{$L(K_{1,s};t,0,\dots,0)$ is $K_{s+2t}$ with $t$ pairwise disjoint edges
  deleted.}
 \begin{proof}\pf\ It has $s$ edges and $2t$ triples, and by Definition~\ref{def:glg} two of its
  vertices are non-adjacent exactly when they are $(1,\alpha,0)$ and $(1,\alpha,1)$ for some
  $\alpha$; these $t$ pairs are disjoint.~\qed\end{proof}
 \step{s412four}{\assume{$X$ is connected and claw-free, $2I+A(X)$ is positive semidefinite, and
  $X$ has more than 36 vertices or a vertex with more than 28 neighbours.}\prove{$X\cong L(Y')$
  for some graph $Y'$, or $X\cong L(K_{1,s};t,0,\dots,0)$ for some $s,t\ge0$.}}
 \begin{proof}\pf\ By Theorem~\ref{thm:4.10}, whose hypotheses hold by \asm{1}, $X$ is a
  generalized line graph or is represented in $\typ E_8$; \asm{1} with Theorem~\ref{thm:4.3},
  assertions 1 and 2, excludes the second. Thus $X\cong L(Y;a_1,\dots,a_m)$ for some $Y$ and
  $a_1,\dots,a_m$ (Definition~\ref{def:glg}), and \stepref{s412three} applies, as $X$ is
  connected and claw-free by \asm{1}.~\qed\end{proof}
 \step{s412five}{Assertion 5 holds.}
 \begin{proof}\pf\ Four-claw-freeness is the absence of 4-claws, by
  Lemma~\ref{lem:4-clawstar}, assertion 1. A 4-claw contains a 3-claw with the same centre, so
  no triple is the centre of a 4-claw, by Lemma~\ref{lem:4-claws}, assertion 1; apply its
  assertion 3 to the edges.~\qed\end{proof}
 \qedstep
 \begin{proof}\pf\ Assertions 1, 2 and 5 are \stepref{s412one}, \stepref{s412two} and
  \stepref{s412five}. Assertion 3 is \stepref{s412three} and \stepref{s412clique}, since a
  generalized line graph is isomorphic to some $L(Y;a_1,\dots,a_m)$ (Definition~\ref{def:glg}).
  Assertion 4 is \stepref{s412four}.~\qed\end{proof}
\end{proof}

\subsection{Theorem 4.13 and Remark 4.14}

A graph $X$ is \emph{strongly regular} if it is regular, any two adjacent vertices have
$a$ common neighbours and any two distinct non-adjacent vertices have $c$ common
neighbours, for some $a$ and $c$; this is Mathlib's \path{SimpleGraph.IsSRGWith}, which
allows complete graphs.

\begin{namedresult}{Lemma}{(the line graph case, LG1)}\label{lem:4-LG1}
 \lean{SimpleGraph.lineGraph_isSRGWith_cases}
 \uses{def:glg,comp:C10,facts:sec4}
 \assume{$Y$ is a connected graph with at least one edge, and $L(Y)$ is strongly regular.}
 \prove{$Y$ is a cycle, or $Y\cong K_n$, $Y\cong K_{h,h}$ or $Y\cong K_{1,h}$ for some $n$ or $h$.}
\end{namedresult}
\begin{proof}
 \pfsketch\ \repair\ We will fill the two steps of the paper named in Gap 9 (errata): the values
 0 and 1 of the constant, and the bipartite case, after the proof of [BCN] Theorem 3.12.4. We
 use the corrected count of Error 7 (errata).
 \step{s4lgdeg}{$Y$ is regular, or $Y$ is bipartite and its degree is constant on each part.}
 \begin{proof}
  \step{s4lgdeg2}{\assume{$j$ and $k$ are neighbours of a vertex $i$ of $Y$.}\prove{$\deg
   j=\deg k$.}}
  \begin{proof}\pf\ The edges $\{i,j\}$ and $\{i,k\}$ have degrees $\deg i+\deg j-2$ and $\deg
   i+\deg k-2$ in $L(Y)\cong L(Y;0,\dots,0)$, by Computation~\ref{comp:C10} and
   Definition~\ref{def:glg}; these are equal, as $L(Y)$ is regular by \asm{0}.~\qed\end{proof}
  \step{s4lgwalk}{\assume{$P$ is a walk of even length in $Y$ from $u$ to $w$.}\prove{$\deg
   u=\deg w$.}}
  \begin{proof}\pf\ Induction on the length of $P$: two vertices two apart on $P$ are equal, or
   are neighbours of the vertex between them and have equal degrees by
   \stepref{s4lgdeg2}.~\qed\end{proof}
  \step{s4lgnb}{\case{$Y$ is not bipartite.}\prove{$Y$ is regular.}}
  \begin{proof}\pf\ By \path{SimpleGraph.two_colorable_iff_forall_loop_even}, with
   \path{SimpleGraph.IsBipartite}, $Y$ has a closed walk $C$ of odd length at some vertex $x$. Let
   $u,w$ be vertices; as $Y$ is connected by \asm{0}, there are walks $P$ from $u$ to $w$ and
   $R$ from $u$ to $x$. The walks $P$ and $R,C,R^{-1},P$ have lengths of opposite parity, so one
   is even, and \stepref{s4lgwalk} gives $\deg u=\deg w$.~\qed\end{proof}
  \step{s4lgbc}{\case{$Y$ is bipartite.}\prove{The degree of $Y$ is constant on each part of a
   2-colouring of $Y$.}}
  \begin{proof}\pf\ \path{SimpleGraph.IsBipartite} gives a colouring of $Y$ by
   $\mathrm{Fin}\,2$; recolour it by $\mathrm{Fin}\,2\simeq\mathrm{Bool}$ with
   \path{SimpleGraph.recolorOfEquiv} and \path{finTwoEquiv}, which keeps the colour classes, as
   \path{SimpleGraph.Coloring.even_length_iff_congr} takes a colouring by $\mathrm{Bool}$. Two
   vertices of the same colour are joined by a walk, as $Y$ is connected by \asm{0}, and it has
   even length by that lemma; apply \stepref{s4lgwalk}.~\qed\end{proof}
  \qedstep
  \begin{proof}\pf\ The two cases are exhaustive.~\qed\end{proof}
 \end{proof}
 \step{s4lgbip}{\case{$Y$ is bipartite with degrees $p\ne q$ on its parts.}\prove{$Y\cong
  K_{1,h}$.}}
 \begin{proof}\pf\ Two edges on a common vertex of degree $p$ have $p-2$ common neighbours in
  $L(Y)$, since their other ends lie in one part and are not adjacent; likewise $q-2$. Both
  kinds of pairs cannot occur, as $p\ne q$; so one degree, say $q$, is 1, and a connected graph
  with such a part is a star.~\qed\end{proof}
 \step{s4lgsmall}{\case{$Y$ is regular of valency $h\le2$.}\prove{$Y\cong K_{1,1}$ or $Y$ is a
  cycle.}}
 \begin{proof}\pf\ A connected 1-regular graph is $K_2$, and a connected 2-regular graph is a
  cycle.~\qed\end{proof}
 \step{s4lgbig}{\case{$Y$ is regular of valency $h\ge3$.}\prove{$Y\cong K_n$ or
  $Y\cong K_{h,h}$.}}
 \begin{proof}
  \step{s4lgcommon}{Two edges $\{i,j\}$, $\{i,k\}$ have $h-1$ common neighbours in $L(Y)$ if
   $j\sim k$ and $h-2$ otherwise.}
  \begin{proof}\pf\ The common neighbours are the $h-2$ other edges on $i$, and $\{j,k\}$ when
   it is an edge.~\qed\end{proof}
  \step{s4lgKn}{\case{Some $j\sim k$ with a common neighbour $i$.}\prove{$Y\cong K_n$.}}
  \begin{proof}\pf\ By \stepref{s4lgcommon} and the constant $a$, every path $j,i,k$
   closes to a triangle; so neighbourhoods are cliques, and a connected such graph is
   complete.~\qed\end{proof}
  \step{s4lgtf}{\case{$Y$ is triangle-free.}\prove{$Y\cong K_{h,h}$.}}
  \begin{proof}
   \pflet{$c$ be the constant of $L(Y)$ for non-adjacent vertices: two disjoint edges are joined by exactly $c$
    edges of $Y$.}
   \step{s4lgdisj}{Every edge $\{i,j\}$ is disjoint from some edge, and $c\le2$.}
   \begin{proof}\pf\ Pick $k\sim i$ with $k\ne j$, and $l\sim k$ with $l\ne i$, as $h\ge3$. Then
    $l\ne j$, else $i,j,k$ is a triangle; so $\{k,l\}$ is disjoint from $\{i,j\}$. Of the edges joining
    $\{i,j\}$ to a disjoint edge, at most one leaves $i$ and at most one leaves $j$, by
    triangle-freeness.~\qed\end{proof}
   \step{s4lgc0}{$c\ne0$.}
   \begin{proof}\pf\ If $c=0$, two vertices of $L(Y)$ at distance 2 would have no common
    neighbour; so the connected graph $L(Y)$ is complete, and $Y$ has no two disjoint edges,
    against \stepref{s4lgdisj}.~\qed\end{proof}
   \step{s4lgc1}{$c\ne1$.}
   \begin{proof}\pf\ Fix an edge $\{i,j\}$. An edge $\{i,k\}$, $k\ne j$, and an edge
    $\{k,l\}$, $l\ne i$, give a disjoint edge $\{k,l\}$ ($l\ne j$ by triangle-freeness); with the
    same count at $j$, there are $2(h-1)^2$ such pairs. If $c=1$, each disjoint edge arises
    once, so $|E(Y)|=1+2(h-1)+2(h-1)^2$, and $|V(Y)|=2|E(Y)|/h=4h-4+2/h$ is not an integer for
    $h\ge3$.~\qed\end{proof}
   \step{s4lgc2}{\assume{$c=2$ and $\{i,j\}$ is an edge.}\prove{$Y\cong K_{h,h}$, with parts
    $N(j)$ and $N(i)$.}}
   \begin{proof}\pf\ $N(i)$, $N(j)$ are disjoint cocliques, by triangle-freeness. For
    $x\in N(j)\setminus\{i\}$ and $y\in N(i)\setminus\{j\}$, the disjoint edges $\{x,j\}$ and
    $\{i,y\}$ are joined by $\{i,j\}$ and one more edge, which can only be $\{x,y\}$. A vertex
    $k\notin N(i)\cup N(j)$ has a neighbour $l\notin\{i,j\}$, and the two edges joining
    $\{k,l\}$ to $\{i,j\}$ put $k$ in $N(i)\cup N(j)$.~\qed\end{proof}
   \qedstep
   \begin{proof}\pf\ $c\in\{0,1,2\}$ by \stepref{s4lgdisj}; apply \stepref{s4lgc0},
    \stepref{s4lgc1} and \stepref{s4lgc2}.~\qed\end{proof}
  \end{proof}
  \qedstep
  \begin{proof}\pf\ Either some path of length 2 closes to a triangle, or $Y$ is
   triangle-free.~\qed\end{proof}
 \end{proof}
 \qedstep
 \begin{proof}\pf\ By \stepref{s4lgdeg}, $Y$ is regular or bipartite with two distinct degrees;
  the cases \stepref{s4lgbip}, \stepref{s4lgsmall} and \stepref{s4lgbig} are
  exhaustive.~\qed\end{proof}
\end{proof}

\begin{cgss}{Theorem}{4.13}\label{thm:4.13}
 \lean{SimpleGraph.seidel}
 \uses{thm:4.11,lem:4-LG1,def:CP,def:represented,facts:sec4}
 \assume{\begin{pfenum}
  \item $X$ is a connected strongly regular graph, and $2I+A(X)$ is positive semidefinite.
  \item $-2$ is an eigenvalue of $X$.
 \end{pfenum}}
 \prove{$X\cong L(K_n)$, $X\cong L(K_{n,n})$ or $X\cong CP(n)$ for some $n$, or $X$ is
 represented in $\typ E_8$.}
 \new\ The two hypotheses together are the paper's ``least eigenvalue $-2$'', read literally,
 so this is the printed theorem. We state them apart because Theorems 4.10 and 4.11 read the
 phrase as the lower bound of the first hypothesis alone, and Theorem 4.13 fails under that
 reading: for $n>36$, $K_n$ is a counterexample (errata, Implicit hypothesis 2).
\end{cgss}
\begin{proof}
 \pfsketch\ We will apply Theorem 4.11, and then the line graph case LG1.
 \step{s413hrc}{$X\cong L(Y)$, or $X\cong CP(m)$, or $X$ is represented in $\typ E_8$.}
 \begin{proof}\pf\ Theorem~\ref{thm:4.11}, with \asm[1]{0}.~\qed\end{proof}
 \step{s413L}{\case{$X\cong L(Y)$.}\prove{$X\cong L(K_n)$ or $X\cong L(K_{n,n})$.}}
 \begin{proof}
  \step{s413Y}{We may take $Y$ connected, with an edge, and $Y$ is a cycle, $K_n$, $K_{h,h}$ or
   $K_{1,h}$.}
  \begin{proof}\pf\ Remove the isolated vertices of $Y$; since $X$ is connected with a vertex,
   the rest is connected with an edge. Apply Lemma~\ref{lem:4-LG1}.~\qed\end{proof}
  \step{s413Kn}{$X$ is neither a complete graph nor $C_5$, and $X$ is not a cycle $C_p$ with
   $p\ge6$.}
  \begin{proof}\pf\ $\det(2I+A)\ne0$ for $K_n$, since $2I+A=I+J$ is positive definite, and
   $\det(2I+A(C_5))=4$; so neither has the eigenvalue $-2$ of \asm[2]{0}, which is a root of
   $\det(\mu I-A)$. In $C_p$ with $p\ge6$, vertices at distance
   2 and 3 have 1 and 0 common neighbours.~\qed\end{proof}
  \qedstep
  \begin{proof}\pf\ $L(K_{1,h})\cong K_h$ is excluded by \stepref{s413Kn}. If $Y=C_p$, then
   $X\cong C_p$, and $p=4$ by \stepref{s413Kn}, since $C_3$ is complete; and
   $C_4=L(K_{2,2})$. The rest is $L(K_n)$ or $L(K_{h,h})$, by \stepref{s413Y}.~\qed\end{proof}
 \end{proof}
 \qedstep
 \begin{proof}\pf\ The cases of \stepref{s413hrc} are exhaustive, and the other two are in the
  conclusion.~\qed\end{proof}
\end{proof}

\begin{namedresult}{Remark}{(after the examples)}\label{rem:4-large}
 \lean{SimpleGraph.hoffman_of_large}\lean{SimpleGraph.hoffman_rayChaudhuri_of_large}
 \lean{SimpleGraph.seidel_of_large}
 \uses{thm:4.3,thm:4.4,thm:4.10,thm:4.11,thm:4.13,def:glg,def:CP}
 \assume{$X$ is a connected graph, and $2I+A(X)$ is positive semidefinite.}
 \prove{\begin{pfenum}
  \item If $X$ has more than 36 vertices, or a vertex with more than 28 neighbours, then $X$ is a
   generalized line graph.
  \item If $X$ is $d$-regular on $n$ vertices, with $n>28$ or $d>16$, then $X\cong L(Y)$ for some
   graph $Y$ or $X\cong CP(m)$ for some $m$.
  \item If $X$ is strongly regular of valency $k$ on $n$ vertices, with $n>28$ or $k>16$, and $-2$
   is an eigenvalue of $X$, then $X\cong L(K_m)$, $X\cong L(K_{m,m})$ or $X\cong CP(m)$ for
   some $m$.
 \end{pfenum}}
 \new\ The paper places this remark before Theorem 4.10 and speaks of ``any subsequent
 theorem''; we will state it after Theorems 4.10, 4.11 and 4.13, with their hypotheses (errata,
 Implicit hypotheses 1 and 2).
\end{namedresult}
\begin{proof}
 \pfsketch\ We will apply the three theorems, and exclude $\typ E_8$ by Theorems 4.3 and 4.4.
 \step{s4lrone}{\assume{$X$ has more than 36 vertices, or a vertex with more than 28
  neighbours.}\prove{$X$ is a generalized line graph.}}
 \begin{proof}\pf\ By Theorem~\ref{thm:4.10} and \asm{0}, $X$ is a generalized line graph or is
  represented in $\typ E_8$; Theorem~\ref{thm:4.3}, assertions 1 and 2, with \asm{1}, excludes
  the second.~\qed\end{proof}
 \step{s4lrtwo}{\assume{$X$ is $d$-regular on $n$ vertices, and $n>28$ or $d>16$.}\prove{$X\cong
  L(Y)$ for some graph $Y$ or $X\cong CP(m)$ for some $m$.}}
 \begin{proof}\pf\ By Theorem~\ref{thm:4.11}, with \asm{0} and \asm{1}, $X\cong L(Y)$, $X\cong
  CP(m)$ or $X$ is represented in $\typ E_8$; Theorem~\ref{thm:4.4} excludes the third, by
  \asm{1}, as $n\ge1$ for a connected graph, by \asm{0}.~\qed\end{proof}
 \step{s4lrthree}{\assume{$X$ is strongly regular of valency $k$ on $n$ vertices, $n>28$ or
  $k>16$, and $-2$ is an eigenvalue of $X$.}\prove{$X\cong L(K_m)$, $X\cong L(K_{m,m})$ or
  $X\cong CP(m)$ for some $m$.}}
 \begin{proof}\pf\ By Theorem~\ref{thm:4.13}, with \asm{0} and \asm{1}, $X$ is one of these or is
  represented in $\typ E_8$; Theorem~\ref{thm:4.4} excludes the last, as $X$ is $k$-regular on
  $n$ vertices by \asm{1}, and $n\ge1$ for a connected graph, by \asm{0}.~\qed\end{proof}
 \qedstep
 \begin{proof}\pf\ \stepref{s4lrone}, \stepref{s4lrtwo} and \stepref{s4lrthree}.~\qed\end{proof}
\end{proof}

\medskip\noindent\textbf{What remains of Section 4.}

Remark 4.14 is stated below, and its proof is in Section~\ref{sec:rem414}.
The remarks after Theorems 4.1 and 4.2, the last pieces of the section, are stated and proved
after those theorems.

\begin{cgss}{Remark}{4.14}\label{rem:4.14}
 \lean{SimpleGraph.seidel_seven}
 \uses{thm:4.11,def:CP,def:represented}
 \assume{$X$ is a connected strongly regular graph with least eigenvalue $-2$, represented in
 $\typ E_8$, and isomorphic to none of $L(K_m)$, $L(K_{m,m})$ and $CP(m)$.}
 \prove{$X$ is isomorphic to one of the Chang graphs I, II and III, the Schl\"afli graph, the
 Shrikhande graph, the Clebsch graph and the Petersen graph; and these seven graphs are
 pairwise non-isomorphic and have these properties.}
 \repair\ The paper gives no proof, and ``the seven graphs'' needs the Chang graphs to be
 pairwise non-isomorphic (errata, Gap 12). We will give its proof, by the
 route of \path{cgss/prerequisites.md} (Section 7.3), in Section~\ref{sec:rem414}.
\end{cgss}

\endgroup

% Section 4, continued: the proof of Remark 4.14 (Seidel's seven graphs).
% Route: cgss/prerequisites.md, Section 7.3 (decision 3 of Section 13).
% Macros of this file, in case macros.tex does not provide them.
% \Ecal_t: the edge set E_t of the graph \Delta_t on {1,...,8} (t = 6, 7, 8).
\providecommand{\Ecal}{\mathcal E}
% The 4-clique profile of a graph, as a multiset.
\providecommand{\cprof}{\opk{prof}_4}
% The symmetric difference of two sets.
\providecommand{\sd}{\mathbin{\triangle}}

\subsection{The proof of Remark 4.14}\label{sec:rem414}
\begingroup
\emergencystretch=3em

We will prove Remark 4.14 by the route of \path{cgss/prerequisites.md}, Section 7.3, which
follows the computer-free part of the proof of [BCN] Theorem 3.12.2 in the language of line
systems, and replaces the enumeration of [BCN] Proposition 3.11.2 by three switching arguments.
A graph $X$ as in the remark is represented by a family $(\Zx_u)$ in $\typ E_8$. We will first
pin down its parameters (Computation C18), and then study the line system $R$ of the vectors of
squared norm 2 in the $\ints$-span of the $\Zx_u$. It is of type $\typ E_6$, $\typ E_7$ or
$\typ E_8$, and the vector $\Zw=\frac1{k+2}\sum_u\Zx_u$ has integral inner products with it. The
simple roots of Computation C21 then exclude two parameter sets. In the three cases where
$(\Zw,\Zw)=2$, $\Zw$ is a root, and $X$ is a switching of $T(5)$, of $T(6)$ with an isolated
vertex, or of $T(8)$, by the edge set of a regular graph; Computation C22 names the result. The
two remaining parameter sets are those of the Clebsch and the Schl\"afli graphs, which are
unique. We will write $n$, $k$, $a$, $c$ for the parameters of a strongly regular graph, as in
the rest of the blueprint. A result tagged \new\ departs from Section 7.3 of
\path{cgss/prerequisites.md}, or from the prerequisite that it cites, and the departure is
stated after the tag; the untagged results follow it.

\begingroup\raggedright
\medskip\noindent\textbf{Facts from Mathlib and the project.}\label{facts:414}
We will use the following declarations of Mathlib (commit \path{082e2d37}):
\path{SimpleGraph.IsSRGWith} with its fields,
\path{SimpleGraph.IsSRGWith.matrix_eq} ($A^2=kI+aA+c\,A(\comp X)$),
\path{SimpleGraph.one_add_adjMatrix_add_compl_adjMatrix_eq_of_one} with
\path{SimpleGraph.compl_adjMatrix_eq_adjMatrix_compl} ($I+A+A(\comp X)=J$),
\path{SimpleGraph.IsSRGWith.param_eq} ($k(k-a-1)=c(n-k-1)$ for $0<n$, in $\mathbb N$ with truncated
subtraction),
\path{SimpleGraph.IsSRGWith.compl} (the complement has parameters
$(n,n-k-1,n-2k+c-2,n-2k+a)$),
\path{SimpleGraph.IsSRGWith.card_neighborFinset_union_of_not_adj} ($|N(u)\cup N(w)|=2k-c$ for
distinct non-adjacent $u,w$),
\path{SimpleGraph.IsRegularOfDegree.compl},
\path{SimpleGraph.sum_degrees_eq_twice_card_edges}, \path{SimpleGraph.IsNClique.map} (an
embedding maps cliques to cliques), \path{compl_compl} ($\comp{\comp X}=X$),
\path{Matrix.exists_mulVec_eq_zero_iff} (a square matrix is singular exactly when it has a
nonzero null vector), \path{Matrix.inv_def} and \path{Matrix.mul_adjugate},
\path{Matrix.det_gram_ne_zero_iff_linearIndependent},
\path{LinearMap.IsProj.trace} with \path{LinearMap.isProj_range_iff_isIdempotentElem} (an
idempotent has rank equal to its trace), \path{abs_real_inner_le_norm} (the Cauchy--Schwarz
inequality), \path{exists_linearIndependent} (a spanning set contains a basis) and
\path{Submodule.finrank_le}. From the project we use
\path{SimpleGraph.leastEigenvalue_eq_neg_two_iff} of
\path{LineGraphs/Spectral/Positivity.lean} ($\lambda_{\min}(X)=-2$ exactly when $2I+A$ is
positive semidefinite and singular), the checkers \path{Computations.checkGraphIso}, \path{Computations.cliqueProfile} and
\path{Computations.cliqueCount} with their soundness theorems
\path{Computations.nonempty_graphOfMatrix_iso}, \path{Computations.checkNonisomorphic_sound} and
\path{Computations.cliqueCount_sound},
and the invariance \path{SimpleGraph.Iso.cliqueProfile_eq}, in
\path{LineGraphs/Computations/CliqueCertificates.lean}. Finally, from the Lean of Computation
C4 we use the transitivity on ordered stars: for $t\in\{6,7,8\}$ and roots $\Za,\Zb$ of
$\typ E_t$ with $(\Za,\Zb)=-1$, there is a linear isometry $g$ of $\re^8$ with
$g(\Phi(\typ E_t))=\Phi(\typ E_t)$ that carries the fixed star $(\Za_0,\Zb_0)$ of Lean to
$(\Za,\Zb)$; this is \leancite{LineSystem.E₈.exists_isometry_star}{\texttt{LineSystem.E\ensuremath{{}_8}.exists\_isometry\_star}}
and its analogues for $\typ E_7$ and $\typ E_6$, in
\path{LineGraphs/Computations/C4E8Stars.lean} and its analogues, whose trees of reflections are
checked by \texttt{decide +kernel}. It replaces LN5 of \path{cgss/prerequisites.md}.
\par
\endgroup

\medskip
Throughout, $\Zc=\frac12\sum_{i=1}^8\Ze_i$ is the root of model (ii) of $\typ E_8$ fixed before
Lemma~\ref{lem:4-switch}, and for a 2-subset $P=\{i,j\}$ of $\{1,\dots,8\}$ we write
$\Ze_P=\Ze_i+\Ze_j$. We will write $\Ecal_6$ for the set of 2-subsets of $\{1,\dots,5\}$,
$\Ecal_7$ for the set of 2-subsets of $\{1,\dots,6\}$ together with $\varepsilon=\{7,8\}$, and
$\Ecal_8$ for the set of all 2-subsets of $\{1,\dots,8\}$; $\Delta_t$ is the graph on
$\{1,\dots,8\}$ with edge set $\Ecal_t$. Thus $L(\Delta_6)\cong T(5)$, $L(\Delta_7)$ is $T(6)$ together
with the isolated vertex $\varepsilon$, and $L(\Delta_8)=T(8)$. For a graph $Z$ and a vertex $i$ we
write $\deg_Z(i)$ for its degree, and for a set $H$ of 2-subsets we write $\deg_H(i)$ for the
number of members of $H$ that contain $i$.

\begin{namedresult}{Definition}{(switching)}\label{def:414-switch}
 \lean{SimpleGraph.switch}
 \uses{}
 \assume{$X$ is a graph and $U\sbs V(X)$.}
 \pflet{The \emph{switched graph} $X^U$ has vertex set $V(X)$, and distinct $u,w$ are adjacent
 in $X^U$ when either they are adjacent in $X$ and both or neither lie in $U$, or they are not
 adjacent in $X$ and exactly one of them lies in $U$. For a graph $Y$ and $H\sbs E(Y)$ we
 write $L(Y)^H$ for the switching of the line graph $L(Y)$ with respect to $H$.}
 This is SW1 of \path{cgss/prerequisites.md}, with $U$ in place of its $Y$, since $Y$ is a
 graph here.
\end{namedresult}

\begin{namedresult}{Lemma}{(switching)}\label{lem:414-switch}
 \lean{SimpleGraph.switch_compl}\lean{SimpleGraph.Iso.switch}\lean{SimpleGraph.Iso.lineGraphSwitch}
 \uses{def:414-switch}
 \assume{$X$ and $X'$ are graphs, $U\sbs V(X)$, $Y$ and $Y'$ are graphs and $H\sbs E(Y)$.}
 \prove{\begin{pfenum}
  \item $X^U=X^{V(X)\setminus U}$.
  \item If $\sigma$ is an isomorphism from $X$ to $X'$, then $\sigma$ is an isomorphism from
   $X^U$ to $X'^{\sigma(U)}$.
  \item If $\sigma$ is an isomorphism from $Y$ to $Y'$, then $e\mapsto\sigma(e)$ is an
   isomorphism from $L(Y)^H$ to $L(Y')^{\sigma(H)}$.
 \end{pfenum}}
\end{namedresult}
\begin{proof}
 \pfsketch\ We will read each assertion off Definition~\ref{def:414-switch}.
 \step{s414swcompl}{Assertion 1 holds.}
 \begin{proof}\pf\ A pair has exactly one member in $U$ if and only if it has exactly one member
  in $V(X)\setminus U$; apply Definition~\ref{def:414-switch}.~\qed\end{proof}
 \step{s414swiso}{Assertion 2 holds.}
 \begin{proof}\pf\ For distinct $u,w$, $\sigma$ preserves adjacency, and $u\in U$ exactly when
  $\sigma(u)\in\sigma(U)$; so the condition of Definition~\ref{def:414-switch} for $u,w$ in $X$
  and $U$ is that for $\sigma(u),\sigma(w)$ in $X'$ and $\sigma(U)$.~\qed\end{proof}
 \step{s414swline}{$e\mapsto\sigma(e)$ is an isomorphism from $L(Y)$ to $L(Y')$.}
 \begin{proof}\pf\ It is a bijection from $E(Y)$ onto $E(Y')$, since $\sigma$ is an isomorphism,
  and $|\sigma(e)\cap\sigma(f)|=|e\cap f|$ since $\sigma$ is injective; two edges are adjacent in
  the line graph when they are distinct and meet.~\qed\end{proof}
 \qedstep
 \begin{proof}\pf\ Assertions 1 and 2 are \stepref{s414swcompl} and \stepref{s414swiso};
  assertion 3 is assertion 2 applied to the isomorphism of \stepref{s414swline} and the set
  $H$.~\qed\end{proof}
\end{proof}

\begin{namedresult}{Lemma}{(isomorphisms and complements)}\label{lem:414-iso}
 \lean{SimpleGraph.Iso.compl}\lean{SimpleGraph.IsSRGWith.of_iso}
 \uses{}
 \assume{$X$ and $X'$ are graphs, and $\sigma$ is an isomorphism from $X$ to $X'$.}
 \prove{\begin{pfenum}
  \item $\sigma$ is an isomorphism from $\comp X$ to $\comp{X'}$.
  \item If $X$ is strongly regular with parameters $(n,k,a,c)$, then so is $X'$.
 \end{pfenum}}
 Mathlib has the complement of an embedding, \path{SimpleGraph.Embedding.complEquiv}, but neither
 assertion.
\end{namedresult}
\begin{proof}
 \pfsketch\ An isomorphism is a bijection that preserves adjacency and non-adjacency.
 \step{s414isocompl}{Assertion 1 holds.}
 \begin{proof}\pf\ $\sigma$ is a bijection, and for distinct $u,w$ the vertices $\sigma(u)$,
  $\sigma(w)$ are distinct and adjacent in $X'$ exactly when $u,w$ are adjacent in $X$; distinct
  vertices are adjacent in the complement exactly when they are not adjacent in the
  graph.~\qed\end{proof}
 \step{s414isosrg}{Assertion 2 holds.}
 \begin{proof}\pf\ $|V(X')|=|V(X)|$, $\sigma$ maps the neighbours of $u$ onto those of $\sigma(u)$,
  and the common neighbours of $u,w$ onto those of $\sigma(u),\sigma(w)$.~\qed\end{proof}
 \qedstep
 \begin{proof}\pf\ \stepref{s414isocompl} and \stepref{s414isosrg}.~\qed\end{proof}
\end{proof}

\subsubsection*{Parameters}

\begin{namedresult}{Lemma}{(cocktail party parameters)}\label{lem:414-CP}
 \lean{SimpleGraph.IsRegularOfDegree.nonempty_iso_cocktailParty}
 \uses{def:CP,facts:414}
 \assume{$m\ge1$, and $X$ is a $(2m-2)$-regular graph on $2m$ vertices.}
 \prove{$X\cong CP(m)$.}
\end{namedresult}
\begin{proof}
 \pfsketch\ The complement of $X$ is a perfect matching.
 \step{s414cpcompl}{The complement $\comp X$ is 1-regular.}
 \begin{proof}\pf\ \path{SimpleGraph.IsRegularOfDegree.compl}: its degree is
  $2m-1-(2m-2)=1$.~\qed\end{proof}
 \step{s414cpmatch}{\pick\ vertices $p_1,q_1,\dots,p_m,q_m$, which are all the vertices of $X$,
  such that the edges of $\comp X$ are the pairs $\{p_i,q_i\}$.}
 \begin{proof}\pf\ By \stepref{s414cpcompl} each vertex lies in exactly one edge of $\comp X$, so
  the edges of $\comp X$ partition $V(X)$ into pairs; there are $2m/2=m$ of
  them.~\qed\end{proof}
 \qedstep
 \begin{proof}\pf\ The map $(i,0)\mapsto p_i$, $(i,1)\mapsto q_i$ is a bijection by
  \stepref{s414cpmatch}. Distinct vertices are adjacent in $X$ exactly when they are not an edge
  of $\comp X$, that is, when their indices differ; this is adjacency in $CP(m)$
  (Definition~\ref{def:CP}).~\qed\end{proof}
\end{proof}

\begin{namedresult}{Lemma}{(strongly regular graphs with eigenvalue $-2$)}\label{lem:414-srg}
 \lean{SimpleGraph.IsSRGWith.params_of_det_two_add_adjMatrix_eq_zero}
 \lean{SimpleGraph.IsSRGWith.rank_two_add_adjMatrix_of_det_eq_zero}
 \uses{facts:414}
 \assume{\begin{pfenum}
  \item $X$ is connected and strongly regular with parameters $(n,k,a,c)$.
  \item $\det(2I+A)=0$, where $A=A(X)$.
 \end{pfenum}}
 \prove{\begin{pfenum}
  \item $1\le k\le n-2$, $a\le k-1$ and $1\le c\le k$.
  \item $2k-c\le n-2$.
  \item $k-2a+c=4$.
  \item $(a-c+4)(\rk(2I+A)-1)=2n-k-2$.
 \end{pfenum}}
 \new\ This is SR4 of \path{cgss/prerequisites.md}, in a form without the spectral theorem:
 the multiplicity $m_1$ of the second eigenvalue is replaced by $\rk(2I+A)-1$, through a
 spectral idempotent. Assertion 2 is the counting condition that the search of C18 needs
 (\path{cgss/certificates/c18_parameter_search.sage}).
\end{namedresult}
\begin{proof}
 \pfsketch\ We will derive the linear relation from a null vector of $2I+A$, and then find an
 idempotent multiple of $2I+A-\frac{k+2}nJ$.
 \pflet{$Q=2I+A$, and let $J$ and $\Zj$ be the all-ones matrix and vector.}
 \step{s414srgnc}{$k\ge1$, and $X$ is not complete.}
 \begin{proof}\pf\ If $k=0$ then $Q=2I$, and if $X$ is complete then $Q=I+J$, which is positive
  definite since $(\Zz,(I+J)\Zz)=|\Zz|^2+(\Zz,\Zj)^2$; in both cases $\det Q\ne0$, against
  \asm[2]{0}.~\qed\end{proof}
 \step{s414srgpairs}{There are an edge $uw$ and distinct non-adjacent vertices $u',w'$.}
 \begin{proof}\pf\ $k\ge1$ and $X$ is not complete, by \stepref{s414srgnc}.~\qed\end{proof}
 \step{s414srgak}{$a\le k-1$.}
 \begin{proof}\pf\ The $a$ common neighbours of the edge $uw$ of \stepref{s414srgpairs} lie in
  $N(u)\setminus\{w\}$, which has $k-1$ elements.~\qed\end{proof}
 \step{s414srgck}{$c\le k$.}
 \begin{proof}\pf\ The $c$ common neighbours of $u',w'$ of \stepref{s414srgpairs} lie in $N(u')$,
  which has $k$ elements.~\qed\end{proof}
 \step{s414srgunion}{$2k-c\le n-2$.}
 \begin{proof}\pf\ $N(u')\cup N(w')$, for $u',w'$ of \stepref{s414srgpairs}, has $2k-c$ elements by
  \path{SimpleGraph.IsSRGWith.card_neighborFinset_union_of_not_adj}, and lies in
  $V(X)\setminus\{u',w'\}$, since $u'\not\sim w'$.~\qed\end{proof}
 \step{s414srgkn}{$k\le n-2$.}
 \begin{proof}\pf\ $k\le2k-c$ by \stepref{s414srgck}; apply \stepref{s414srgunion}.~\qed\end{proof}
 \step{s414srgc}{$c\ge1$.}
 \begin{proof}\pf\ Take non-adjacent distinct $u',w'$ (\stepref{s414srgnc}) and a shortest path
  $u'=u_0,u_1,\dots,u_d=w'$ in the connected graph $X$, so $d\ge2$. Then $u_0$ and $u_2$ are
  distinct and non-adjacent, by minimality, and $u_1$ is a common
  neighbour.~\qed\end{proof}
 \step{s414srgcompl}{$A(\comp X)=J-I-A$.}
 \begin{proof}\raggedright\pf\ \path{SimpleGraph.one_add_adjMatrix_add_compl_adjMatrix_eq_of_one} and
  \path{SimpleGraph.compl_adjMatrix_eq_adjMatrix_compl}.~\qed\end{proof}
 \step{s414srgmat}{$A^2=kI+aA+c(J-I-A)$, $AJ=JA=kJ$ and $J^2=nJ$.}
 \begin{proof}\pf\ \path{SimpleGraph.IsSRGWith.matrix_eq} gives $A^2=kI+aA+c\,A(\comp X)$; apply
  \stepref{s414srgcompl}. $AJ=JA=kJ$ by regularity, and $J^2=nJ$ since each entry of $J^2$ is a
  sum of $n$ ones.~\qed\end{proof}
 \step{s414srgnull}{\pick\ $\Zz\ne0$ with $Q\Zz=0$; then $A\Zz=-2\Zz$ and $J\Zz=0$.}
 \begin{proof}\pf\ \path{Matrix.exists_mulVec_eq_zero_iff} and \asm[2]{0} give $\Zz$. Then
  $k(\Zj,\Zz)=(A\Zj,\Zz)=(\Zj,A\Zz)=-2(\Zj,\Zz)$ by \stepref{s414srgmat} and the symmetry of $A$,
  so $(\Zj,\Zz)=0$, since $k+2>0$; and $J\Zz=(\Zj,\Zz)\Zj$.~\qed\end{proof}
 \step{s414srglin}{$k-2a+c=4$.}
 \begin{proof}\pf\ Apply both sides of \stepref{s414srgmat} to $\Zz$ of \stepref{s414srgnull}:
  $4\Zz=k\Zz-2a\Zz+c(0-\Zz+2\Zz)=(k-2a+c)\Zz$, and $\Zz\ne0$.~\qed\end{proof}
 \pflet{$\eta=a-c+2$.}
 \step{s414srgtheta}{$k-c=2\eta$, $\eta\ge0$, and $(A-\eta I)Q=cJ$.}
 \begin{proof}\pf\ The first is \stepref{s414srglin}, and $\eta\ge0$ follows with $c\le k$
  (\stepref{s414srgck}). By \stepref{s414srgmat},
  $(A-\eta I)(A+2I)=A^2+(2-\eta)A-2\eta I=(k-c-2\eta)I+(a-c+2-\eta)A+cJ$, and both
  coefficients vanish.~\qed\end{proof}
 \step{s414srgcn}{$cn=(k+2)(k-\eta)$.}
 \begin{proof}\pf\ Apply \stepref{s414srgtheta} to $\Zj$: $Q\Zj=(k+2)\Zj$ and
  $(A-\eta I)\Zj=(k-\eta)\Zj$ by \stepref{s414srgmat}, so
  $(k-\eta)(k+2)\Zj=cJ\Zj=cn\Zj$.~\qed\end{proof}
 \pflet{$E=Q-\frac{k+2}nJ$.}
 \step{s414srgPJ}{$PJ=JP=0$, and $E=Q(I-\frac1nJ)$.}
 \begin{proof}\pf\ $QJ=JQ=(k+2)J$ and $J^2=nJ$ by \stepref{s414srgmat}.~\qed\end{proof}
 \step{s414srgidem}{$E^2=(\eta+2)E$.}
 \begin{proof}\pf\
  \[
  \begin{split}
   E^2&=PQ-\tfrac{k+2}nPJ=PQ=Q^2-\tfrac{(k+2)^2}nJ\\
      &=(A-\eta I)Q+(\eta+2)Q-\tfrac{(k+2)^2}nJ
       =(\eta+2)Q+\Bigl(c-\tfrac{(k+2)^2}n\Bigr)J\\
      &=(\eta+2)Q-(\eta+2)\tfrac{k+2}nJ=(\eta+2)E,
  \end{split}
  \]
  where the first line uses \stepref{s414srgPJ} and $JQ=(k+2)J$, the second
  \stepref{s414srgtheta}, and the third \stepref{s414srgcn}, which gives
  $cn-(k+2)^2=-(k+2)(\eta+2)$.~\qed\end{proof}
 \step{s414srgrkP}{$(\eta+2)\rk E=2n-k-2$.}
 \begin{proof}\pf\ $\eta+2\ge2$ by \stepref{s414srgtheta}, so $E/(\eta+2)$ is idempotent by
  \stepref{s414srgidem}, and its rank is its trace (\path{LinearMap.IsProj.trace} and
  \path{LinearMap.isProj_range_iff_isIdempotentElem}). Further $\tr E=\tr Q-\frac{k+2}n\tr J=2n-(k+2)$,
  since $\tr A=0$.~\qed\end{proof}
 \step{s414srgrkQ}{$\rk Q=\rk E+1$.}
 \begin{proof}\pf\ The image of $E=Q(I-\frac1nJ)$ lies in that of $Q$, and so does
  $\Zj=Q\Zj/(k+2)$; conversely $Q=E+\frac{k+2}nJ$ has image in $\opk{im}E+\re\Zj$. Finally
  $\Zj\notin\opk{im}E$: $JP=0$ (\stepref{s414srgPJ}) gives $(\Zj,E\Zy)=0$ for every $\Zy$, while
  $(\Zj,\Zj)=n>0$.~\qed\end{proof}
 \qedstep
 \begin{proof}\pf\ Assertion 1 is \stepref{s414srgnc}, \stepref{s414srgkn}, \stepref{s414srgak},
  \stepref{s414srgck} and \stepref{s414srgc}, assertion 2 is \stepref{s414srgunion} and assertion 3 is
  \stepref{s414srglin}. For assertion 4, $\eta+2=a-c+4$, and \stepref{s414srgrkP} and
  \stepref{s414srgrkQ} give $(\eta+2)(\rk Q-1)=2n-k-2$.~\qed\end{proof}
\end{proof}

\begin{namedresult}{Computation}{C18}\label{comp:C18}
 \lean{Computations.C18.srg_parameters}
 \uses{}
 \assume{$n,k,a,c,\rho$ are integers with $1\le k\le n-2$, $n\le36$, $k\le16$, $0\le a\le k-1$,
 $1\le c\le k$ and $0\le\rho\le8$, such that
 \begin{pfenum}
  \item $k(k-a-1)=c(n-k-1)$,
  \item $k-2a+c=4$,
  \item $2k-c\le n-2$, and
  \item $(a-c+4)(\rho-1)=2n-k-2$.
 \end{pfenum}}
 \prove{$(n,k,a,c,\rho)$ is one of the following 15 tuples:
 $(2m,2m-2,2m-4,2m-2,m+1)$ for $2\le m\le7$; $(9,4,1,2,5)$; $(10,6,3,4,5)$;
 $(10,3,0,1,6)$; $(15,8,4,4,6)$; $(16,10,6,6,6)$; $(16,6,2,2,7)$; $(21,10,5,4,7)$;
 $(27,16,10,8,7)$; and $(28,12,6,4,8)$.}
 This is Lemma P of \path{cgss/prerequisites.md} (SR5), with $\rho=1+m_1$. The first six are
 the parameters of $CP(m)$, the next two those of $L_2(3)$ and $T(5)$.
\end{namedresult}
\begin{proof}
 \pfsketch\ This is a finite search. The Lean proof is one \texttt{Bool} search over the
 valencies with a soundness lemma, as for the checker \path{Computations.allSRGParameters} of
 \path{LineGraphs/Computations/BoundedSearch.lean}, whose pruning ($a<k$, $0<c\le k$, $k+1\le n$)
 matches our bounds; \texttt{interval\_cases} would make about $1.5\cdot10^5$ goals.
 \qedstep
 \begin{proof}\pf\ For each $k\le16$, $a<k$ and $1\le c\le k$, condition 1 determines $n$, and
  condition 4 determines $\rho$, since $a-c+4=0$ would give $2n=k+2$, against $k+2\le n$ and
  $n\ge3$; the search keeps the tuples satisfying all four
  conditions and the bounds. It finds the 15 tuples, and no other (SageMath,
  \path{cgss/certificates/c18_parameter_search.sage}).
  Without condition 3 the tuple $(21,16,12,12,7)$ would also survive.~\qed\end{proof}
\end{proof}

\begin{namedresult}{Lemma}{(strongly regular line graphs)}\label{lem:414-linegraph}
 \lean{SimpleGraph.IsSRGWith.nonempty_iso_lineGraph_complete_or_completeBipartite}
 \uses{lem:4-LG1,lem:414-iso,facts:414}
 \assume{\begin{pfenum}
  \item $X$ is a connected graph, strongly regular with parameters $(n,k,a,c)$.
  \item $\det(2I+A(X))=0$.
  \item $Y$ is a graph and $X\cong L(Y)$.
 \end{pfenum}}
 \prove{$X\cong L(K_m)$ or $X\cong L(K_{m,m})$ for some $m\ge2$.}
 \new\ This is the line graph case of Theorem 4.13, stated as a result so that we can cite it;
 its argument is that of the proof of Theorem 4.13, through Lemma LG1.
\end{namedresult}
\begin{proof}
 \pfsketch\ We will remove the isolated vertices of $Y$ and apply Lemma LG1; the eigenvalue $-2$
 excludes complete graphs and $C_5$.
 \pflet{$Y'$ be $Y$ with its isolated vertices removed.}
 \step{s414lgsame}{$L(Y')=L(Y)$.}
 \begin{proof}\pf\ Removing isolated vertices keeps the edges and their
  intersections.~\qed\end{proof}
 \step{s414lgedge}{$Y'$ has an edge.}
 \begin{proof}\pf\ $X$ has a vertex, being connected, so $L(Y)\cong X$ (\asm[3]{0}) has one; it is
  an edge of $Y$, and of $Y'$ by \stepref{s414lgsame}.~\qed\end{proof}
 \step{s414lgconn}{$Y'$ is connected.}
 \begin{proof}\pf\ Every vertex of $Y'$ lies on an edge, and a path in the connected graph
  $L(Y')\cong X$ (\stepref{s414lgsame}, \asm[3]{0}) between two edges gives a walk in $Y'$ between
  their ends.~\qed\end{proof}
 \step{s414lgcases}{$Y'$ is a cycle, or $Y'\cong K_m$, $Y'\cong K_{q,q}$ or $Y'\cong K_{1,q}$
  for some $m$ or $q$.}
 \begin{proof}\pf\ Lemma~\ref{lem:4-LG1}, with \stepref{s414lgedge} and \stepref{s414lgconn},
  since $L(Y')$ is strongly regular by \asm[1]{0}, \asm[3]{0}, \stepref{s414lgsame} and
  Lemma~\ref{lem:414-iso}, assertion 2.~\qed\end{proof}
 \step{s414lgKn}{$X$ is not complete.}
 \begin{proof}\pf\ If $X$ is complete on $n\ge1$ vertices, then $2I+A=I+J$, which is positive
  definite since $(\Zz,(I+J)\Zz)=|\Zz|^2+(\Zz,\Zj)^2$; so its determinant is not 0, against
  \asm[2]{0}.~\qed\end{proof}
 \step{s414lgCfivedet}{$\det(2I+A(C_5))=4$.}
 \begin{proof}\pf\ \texttt{decide} on the $5\times5$ integer matrix.~\qed\end{proof}
 \step{s414lgC5}{$X\not\cong C_5$.}
 \begin{proof}\pf\ An isomorphism $X\to C_5$ conjugates $2I+A(X)$ into $2I+A(C_5)$ by a
  permutation matrix, which keeps the determinant; apply \stepref{s414lgCfivedet} and
  \asm[2]{0}.~\qed\end{proof}
 \step{s414lgCp}{$X\not\cong C_p$ for $p\ge6$.}
 \begin{proof}\pf\ In $C_p$ with $p\ge6$ there are non-adjacent vertices at distance 2 and at
  distance 3, with 1 and 0 common neighbours; so $C_p$ is not strongly regular, against
  \asm[1]{0}.~\qed\end{proof}
 \step{s414lgcycle}{\case{$Y'$ is a cycle $C_p$.}\prove{$X\cong L(K_{2,2})$.}}
 \begin{proof}\pf\ $L(C_p)\cong C_p$, and $p\ge3$. The cases $p=3$ (complete), $p=5$ and $p\ge6$
  are excluded by \stepref{s414lgKn}, \stepref{s414lgC5} and \stepref{s414lgCp}. So $p=4$, and
  $C_4\cong L(C_4)=L(K_{2,2})$.~\qed\end{proof}
 \step{s414lgstar}{\case{$Y'\cong K_{1,q}$.}\prove{This case does not occur.}}
 \begin{proof}\pf\ Any two edges of $K_{1,q}$ meet, so $L(K_{1,q})\cong K_q$ is complete, against
  \stepref{s414lgKn}.~\qed\end{proof}
 \qedstep
 \begin{proof}\pf\ By \stepref{s414lgcases} the cases \stepref{s414lgcycle},
  \stepref{s414lgstar}, $Y'\cong K_m$ and $Y'\cong K_{q,q}$ are exhaustive. In the last two,
  $X\cong L(K_m)$ or $L(K_{q,q})$ by \asm[3]{0} and \stepref{s414lgsame}. Further $m\ge2$, since
  $L(K_m)$ has no vertex for $m\le1$ while the connected graph $X$ has one; and $q\ge2$, since
  $L(K_{0,0})$ has no vertex and $L(K_{1,1})$ is complete (\stepref{s414lgKn}).~\qed\end{proof}
\end{proof}

\subsubsection*{The root system of a representation}

\begin{namedresult}{Lemma}{(the lattice roots)}\label{lem:414-lattice}
 \lean{SimpleGraph.IsRepresentation.latticeRoots}
 \lean{SimpleGraph.IsRepresentation.isStarClosed_latticeRoots}
 \lean{SimpleGraph.IsRepresentation.isIndecomposable_latticeRoots}
 \lean{SimpleGraph.IsRepresentation.span_latticeRoots}
 \lean{SimpleGraph.IsRepresentation.isRepresentedIn_latticeRoots}
 \uses{def:represented,def:line-system,def:indecomposable,def:2.1,lem:4-lines,lem:sec3-connected-span,facts:414}
 \assume{$X$ is a connected graph with a representation $(\Zx_u)$ in a finite-dimensional real
 inner product space $V$; $\Lambda$ is the $\ints$-span and $W$ the span of the $\Zx_u$; and
 $\Phi_R=\{\Zy\in \Lambda:(\Zy,\Zy)=2\}$.}
 \prove{\begin{pfenum}
  \item $(\Zy,\Zz)\in\ints$ for all $\Zy,\Zz\in \Lambda$.
  \item $\Phi_R$ is the set of vectors of a line system $R$ in $V$, which is nonempty,
   star-closed and indecomposable and spans $W$.
  \item $X$ is represented in $R$ by $(\Zx_u)$.
 \end{pfenum}}
 \new\ \path{cgss/prerequisites.md} (LN3) uses the star closure of the $\Zx_u$. The larger set
 $\Phi_R$ is star-closed as well, and it lies in $\Lambda$ by definition, which is what the apex
 argument needs; we do not need $R\sbs\typ E_8$.
\end{namedresult}
\begin{proof}
 \pfsketch\ Inner products in $\Lambda$ are integers, so norms and Cauchy--Schwarz do the rest.
 \step{s414latint}{Assertion 1 holds.}
 \begin{proof}\pf\ $\Zy=\sum_u y_u\Zx_u$ and $\Zz=\sum_w z_w\Zx_w$ with integers $y_u,z_w$, and
  $(\Zy,\Zz)=\sum_{u,w}y_uz_w(\Zx_u,\Zx_w)$, where each $(\Zx_u,\Zx_w)\in\{0,1,2\}$ by
  Definition~\ref{def:represented}.~\qed\end{proof}
 \step{s414latvals}{$\Phi_R=-\Phi_R$, and $(\Zy,\Zz)\in\{-2,-1,0,1,2\}$ for
  $\Zy,\Zz\in\Phi_R$.}
 \begin{proof}\pf\ $\Lambda=-\Lambda$. By \path{abs_real_inner_le_norm},
  $|(\Zy,\Zz)|\le|\Zy||\Zz|=2$, and $(\Zy,\Zz)$ is an integer by
  \stepref{s414latint}.~\qed\end{proof}
 \step{s414latmem}{$\Zx_u\in\Phi_R$ for every $u$, and $\Phi_R\sbs \Lambda\sbs W$.}
 \begin{proof}\pf\ $(\Zx_u,\Zx_u)=2$ by Definition~\ref{def:represented}; $\Lambda\sbs W$ since integer
  combinations are combinations.~\qed\end{proof}
 \step{s414latbasis}{\pick\ vertices $u_1,\dots,u_\rho$ such that $\Zx_{u_1},\dots,\Zx_{u_\rho}$
  form a basis of $W$.}
 \begin{proof}\pf\ \path{exists_linearIndependent}, since the $\Zx_u$ span $W$ and $W$ is
  finite-dimensional.~\qed\end{proof}
 \step{s414latinj}{The map $\iota:\Zy\mapsto((\Zy,\Zx_{u_i}))_{i\le\rho}$ is injective on $W$.}
 \begin{proof}\pf\ It is linear, and a vector of $W$ orthogonal to the basis of
  \stepref{s414latbasis} is orthogonal to $W$, so to itself, and is 0.~\qed\end{proof}
 \step{s414latfin}{$\Phi_R$ is finite.}
 \begin{proof}\pf\ By \stepref{s414latvals} and \stepref{s414latmem}, $\iota$ maps $\Phi_R\sbs W$
  into the finite set $\{-2,\dots,2\}^\rho$; apply \stepref{s414latinj}.~\qed\end{proof}
 \step{s414latstar}{$\Phi_R$ is star-closed.}
 \begin{proof}\pf\ For $\Zy,\Zz\in\Phi_R$ with $(\Zy,\Zz)=-1$, $\Zy+\Zz\in \Lambda$ and
  $(\Zy+\Zz,\Zy+\Zz)=2-2+2=2$; this is Definition~\ref{def:2.1}.~\qed\end{proof}
 \step{s414latspan}{$\Phi_R$ spans $W$.}
 \begin{proof}\pf\ By \stepref{s414latmem}, $\{\Zx_u\}\sbs\Phi_R\sbs W$, and the $\Zx_u$ span
  $W$.~\qed\end{proof}
 \step{s414latind}{$R$ is nonempty and indecomposable.}
 \begin{proof}\pf\ Take $B=\{\Zx_u\}$ in Lemma~\ref{lem:sec3-connected-span}: $B$ is nonempty
  since $X$ has a vertex, it spans the span of $\Phi_R$ by \stepref{s414latspan}, and distinct
  members of $B$ have nonzero inner product exactly when their vertices are adjacent
  (Lemma~\ref{lem:4-lines}, assertion 1, and Definition~\ref{def:represented}), so its graph is a
  copy of the connected graph $X$.~\qed\end{proof}
 \qedstep
 \begin{proof}\pf\ Assertion 1 is \stepref{s414latint}. By \stepref{s414latvals},
  \stepref{s414latfin} and Definition~\ref{def:line-system}, $\Phi_R$ is the set of vectors of a
  line system $R$; \stepref{s414latstar}, \stepref{s414latspan} and \stepref{s414latind} give the
  rest of assertion 2. Assertion 3 is \stepref{s414latmem} with
  Definition~\ref{def:represented}.~\qed\end{proof}
\end{proof}

\begin{namedresult}{Lemma}{(the apex)}\label{lem:414-apex}
 \lean{SimpleGraph.IsRepresentation.apex}\lean{SimpleGraph.IsRepresentation.inner_apex}
 \lean{SimpleGraph.IsRepresentation.inner_apex_self}
 \lean{SimpleGraph.IsRepresentation.inner_apex_mem_int}
 \lean{SimpleGraph.IsRepresentation.apex_mem_span}
 \uses{def:represented}
 \assume{$X$ is a $k$-regular graph on $n\ge1$ vertices with a representation $(\Zx_u)$ in a
 real inner product space $V$; $\Lambda$ is the $\ints$-span and $W$ the span of the $\Zx_u$; and
 $\Zw=\frac1{k+2}\sum_u\Zx_u$.}
 \prove{\begin{pfenum}
  \item $(\Zw,\Zx_v)=1$ for every vertex $v$.
  \item $(\Zw,\Zw)=n/(k+2)$.
  \item $(\Zw,\Zy)\in\ints$ for every $\Zy\in \Lambda$.
  \item $\Zw\in W$.
 \end{pfenum}}
 \new\ Step $\langle1\rangle5$ of \path{cgss/prerequisites.md}, Section 7.3, asks for integral
 inner products with the roots of $R$ only; we prove them on the whole $\ints$-span, and add
 assertion 4, which the switching from the apex needs.
\end{namedresult}
\begin{proof}
 \pfsketch\ This is step $\langle1\rangle5$ of \path{cgss/prerequisites.md}, Section 7.3; the
 row sums of $2I+A$ are $k+2$.
 \step{s414apexone}{Assertion 1 holds.}
 \begin{proof}\pf\ $(\Zw,\Zx_v)=\frac1{k+2}\sum_u(\Zx_u,\Zx_v)=\frac1{k+2}(2+k)$, since by
  Definition~\ref{def:represented} the terms are 2 for $u=v$, 1 for the $k$ neighbours of $v$ and 0
  otherwise.~\qed\end{proof}
 \step{s414apextwo}{Assertion 2 holds.}
 \begin{proof}\pf\ $(\Zw,\Zw)=\frac1{k+2}\sum_v(\Zw,\Zx_v)=\frac n{k+2}$ by
  \stepref{s414apexone}.~\qed\end{proof}
 \step{s414apexthree}{Assertion 3 holds.}
 \begin{proof}\pf\ For $\Zy=\sum_vz_v\Zx_v$ with integers $z_v$, $(\Zw,\Zy)=\sum_vz_v$ by
  \stepref{s414apexone}.~\qed\end{proof}
 \qedstep
 \begin{proof}\pf\ \stepref{s414apexone}, \stepref{s414apextwo}, \stepref{s414apexthree}, and
  $\Zw$ is a combination of the $\Zx_u$.~\qed\end{proof}
\end{proof}

\begin{namedresult}{Lemma}{(integrality)}\label{lem:414-dual}
 \lean{InnerProductSpace.det_gram_mul_inner_self_mem_int}
 \uses{facts:414}
 \assume{$W$ is a real inner product space with a basis $\Zb_1,\dots,\Zb_t$ whose Gram matrix
 $B$ has integer entries, and $\Zw\in W$ has $(\Zw,\Zb_i)\in\ints$ for every $i$.}
 \prove{$\det(B)\,(\Zw,\Zw)$ is an integer.}
 \new\ Step $\langle1\rangle6$ of \path{cgss/prerequisites.md}, Section 7.3, writes
 $(\Zw,\Zw)=\Za^TB^{-1}\Za$ and cites LA6, which lists tools only; we state the integrality as a
 result and prove it with the adjugate.
\end{namedresult}
\begin{proof}
 \pfsketch\ We will write $(\Zw,\Zw)$ with the inverse of $B$, and clear the denominator with
 the adjugate.
 \pflet{$\Za\in\ints^t$ with $\Za(i)=(\Zw,\Zb_i)$, and $\beta\in\re^t$ with
 $\Zw=\sum_i\beta(i)\Zb_i$.}
 \step{s414dualeq}{$\Za=B\beta$ and $(\Zw,\Zw)=\beta^T\Za$.}
 \begin{proof}\pf\ $\Za(i)=\sum_j\beta(j)(\Zb_j,\Zb_i)$ and
  $(\Zw,\Zw)=\sum_i\beta(i)(\Zw,\Zb_i)$, by bilinearity.~\qed\end{proof}
 \step{s414dualinv}{$\det B\ne0$.}
 \begin{proof}\raggedright\pf\ The $\Zb_i$ are linearly independent; apply
  \path{Matrix.det_gram_ne_zero_iff_linearIndependent}.~\qed\end{proof}
 \step{s414dualadj}{$\det(B)(\Zw,\Zw)=\Za^T\opk{adj}(B)\,\Za$.}
 \begin{proof}\pf\
  $\det(B)(\Zw,\Zw)=\det(B)\,\Za^T\beta=\det(B)\,\Za^TB^{-1}\Za=\Za^T\opk{adj}(B)\,\Za$, by
  \stepref{s414dualeq}, then $\beta=B^{-1}\Za$ with \stepref{s414dualinv}, then
  $B^{-1}=\det(B)^{-1}\opk{adj}(B)$ (\path{Matrix.inv_def}).~\qed\end{proof}
 \qedstep
 \begin{proof}\pf\ The entries of $\opk{adj}(B)$ are signed minors of the integer matrix $B$, and
  $\Za$ is integral; apply \stepref{s414dualadj}.~\qed\end{proof}
\end{proof}

\begin{namedresult}{Computation}{C21}\label{comp:C21}
 \lean{LineSystem.E₆.simpleRoot}\lean{LineSystem.E₆.simpleRoot_mem}
 \lean{LineSystem.E₆.gram_simpleRoot}\lean{LineSystem.E₆.det_gram_simpleRoot}
 \lean{LineSystem.E₆.simpleBasis}\lean{LineSystem.E₆.gram_simpleRoot_eq_cartan}
 \lean{LineSystem.E₇.simpleRoot}\lean{LineSystem.E₇.simpleRoot_mem}
 \lean{LineSystem.E₇.gram_simpleRoot}\lean{LineSystem.E₇.det_gram_simpleRoot}
 \lean{LineSystem.E₇.simpleBasis}\lean{LineSystem.E₇.gram_simpleRoot_eq_cartan}
 \lean{LineSystem.E₈.simpleRoot}\lean{LineSystem.E₈.simpleRoot_mem}
 \lean{LineSystem.E₈.gram_simpleRoot}\lean{LineSystem.E₈.det_gram_simpleRoot}
 \lean{LineSystem.E₈.simpleBasis}\lean{LineSystem.E₈.gram_simpleRoot_eq_cartan}
 \lean{Computations.C21.e6_gram_det}\lean{Computations.C21.e7_gram_det}
 \lean{Computations.C21.e8_gram_det}
 \uses{def:3.2,def:3.3,lem:sec3-spans}
 \assume{$t\in\{6,7,8\}$.}
 \prove{There are roots $\Zs_1,\dots,\Zs_t\in\Phi(\typ E_t)$ which form a basis of the span of
 $\Phi(\typ E_t)$, and whose Gram matrix $B_t$ has integer entries and determinant 3, 2 and 1
 for $t=6,7,8$. Up to a relabelling, $B_t$ is Mathlib's \texttt{CartanMatrix.E}~$t$.}
 This is LN7 of \path{cgss/prerequisites.md}.
\end{namedresult}
\begin{proof}
 \pfsketch\ The roots are those of the SageMath certificate. Lean checks their membership, their
 Gram matrix and a scaled $LDL^T$ factorization of it by \texttt{decide +kernel}; see
 \path{LineGraphs/Computations/C21E6.lean} and its analogues.
 \qedstep
 \begin{proof}\pf\ \leancite{LineSystem.E₆.simpleRoot_mem}{\texttt{LineSystem.E\ensuremath{{}_6}.simpleRoot\_mem}},
  \leancite{LineSystem.E₆.gram_simpleRoot}{\texttt{E\ensuremath{{}_6}.gram\_simpleRoot}} and
  \leancite{LineSystem.E₆.det_gram_simpleRoot}{\texttt{E\ensuremath{{}_6}.det\_gram\_simpleRoot}}
  give the roots, their integer Gram matrix and its determinant; the Gram determinant is not 0,
  so the $t$ roots are independent, and they span since the span has dimension $t$
  (Lemma~\ref{lem:sec3-spans}), which is
  \leancite{LineSystem.E₆.simpleBasis}{\texttt{E\ensuremath{{}_6}.simpleBasis}}. The same
  declarations for $\typ E_7$, $\typ E_8$ give the other two cases.~\qed\end{proof}
\end{proof}

\begin{namedresult}{Computation}{C20 at the apex}\label{comp:414-Nc}
 \lean{LineSystem.E₆.c_mem}\lean{LineSystem.E₇.c_mem}\lean{LineSystem.E₈.c_mem}
 \lean{LineSystem.E₆.N_c}\lean{LineSystem.E₇.N_c}\lean{LineSystem.E₈.N_c}
 \lean{LineSystem.E₆.card_CEdge}\lean{LineSystem.E₇.card_CEdge}\lean{LineSystem.E₈.card_CEdge}
 \uses{def:3.2,def:3.3,comp:C20}
 \assume{$t\in\{6,7,8\}$, and $N_\Zc=\{\Zx\in\Phi(\typ E_t):(\Zx,\Zc)=1\}$.}
 \prove{$\Zc\in\Phi(\typ E_t)$, $N_\Zc=\{\Ze_P:P\in\Ecal_t\}\cup\{\Zc-\Ze_P:P\in\Ecal_t\}$, and
 $|\Ecal_t|$ is 10, 16 and 28 for $t=6,7,8$.}
 This is LN6 of \path{cgss/prerequisites.md}, at the root $\Zc$.
\end{namedresult}
\begin{proof}
 \pfsketch\ Lean identifies the pairs of $N_\Zc$ that Computation C20 lists with the members of
 $\Ecal_t$, in \path{LineGraphs/Computations/C20Apex.lean}.
 \qedstep
 \begin{proof}\pf\ \leancite{LineSystem.E₆.c_mem}{\texttt{LineSystem.E\ensuremath{{}_6}.c\_mem}},
  \leancite{LineSystem.E₆.N_c}{\texttt{E\ensuremath{{}_6}.N\_c}} and
  \leancite{LineSystem.E₆.card_CEdge}{\texttt{E\ensuremath{{}_6}.card\_CEdge}}, and their analogues for
  $\typ E_7$ and $\typ E_8$; in Lean the coordinates are numbered from 0, and $\Ecal_t$ is the
  subtype \leancite{LineSystem.E₆.CEdge}{\texttt{E\ensuremath{{}_6}.CEdge}} of ordered pairs
  $i<j$.~\qed\end{proof}
\end{proof}

\begin{namedresult}{Lemma}{(transitivity on roots)}\label{lem:414-transitive}
 \lean{LineSystem.E₆.exists_isometry_apply_eq_c}\lean{LineSystem.E₇.exists_isometry_apply_eq_c}
 \lean{LineSystem.E₈.exists_isometry_apply_eq_c}
 \uses{facts:414,comp:414-Nc,comp:C4,lem:3.4,def:3.3,def:indecomposable,facts:inner}
 \assume{$t\in\{6,7,8\}$ and $\Zr\in\Phi(\typ E_t)$.}
 \prove{There is a linear isometry $\psi$ of $\re^8$ with $\psi(\Phi(\typ E_t))=\Phi(\typ E_t)$ and
 $\psi(\Zr)=\Zc$.}
 \new\ This replaces LN5 of \path{cgss/prerequisites.md}, a transitivity by reflections, by the
 transitivity on ordered stars that the Lean of Computation C4 proves.
\end{namedresult}
\begin{proof}
 \pfsketch\ Both $\Zr$ and $\Zc$ lie in ordered stars, and the stars form one orbit.
 \step{s414trr}{There is $\Zs\in\Phi(\typ E_t)$ with $(\Zr,\Zs)=-1$.}
 \begin{proof}\pf\ $\typ E_t$ is indecomposable (Computation~\ref{comp:C4}) and has more than one
  line (Lemma~\ref{lem:3.4}). If every root other than $\pm\Zr$ were orthogonal to $\Zr$, then
  $\{\pm\Zr\}$ and the other roots would split $\Phi(\typ E_t)$ orthogonally into nonempty parts
  (the definition of indecomposability). So $(\Zr,\Zs')\ne0$ for some $\Zs'\ne\pm\Zr$, and then
  $(\Zr,\Zs')=\pm1$ by \eqref{eq:endpoints}; take $\Zs=\mp\Zs'$.~\qed\end{proof}
 \step{s414trc}{$\Zc$ and $-\Ze_1-\Ze_2$ lie in $\Phi(\typ E_t)$, and
  $(\Zc,-\Ze_1-\Ze_2)=-1$.}
 \begin{proof}\pf\ $\Zc\in\Phi(\typ E_t)$ by Computation~\ref{comp:414-Nc}; $-\Ze_1-\Ze_2$ lies in
  $\Phi(\typ D_8)$ and has $x_6=x_7=x_8=0$, so it lies in $\Phi(\typ E_6)\sbs\Phi(\typ E_7)$ by
  Definition~\ref{def:3.3}; and $(\Zc,\Ze_1+\Ze_2)=\frac12+\frac12$.~\qed\end{proof}
 \qedstep
 \begin{proof}\pf\ The transitivity on ordered stars of Computation~\ref{comp:C4} (in
  Lean; see the facts at the start of this section), applied to the star
  $(\Zr,\Zs)$ of \stepref{s414trr} and to the star $(\Zc,-\Ze_1-\Ze_2)$ of \stepref{s414trc},
  gives isometries $g_1,g_2$ preserving $\Phi(\typ E_t)$ with $g_1(\Za_0)=\Zr$ and
  $g_2(\Za_0)=\Zc$; take $\psi=g_2g_1^{-1}$.~\qed\end{proof}
\end{proof}

\begin{namedresult}{Lemma}{(switching from the apex)}\label{lem:414-transversal}
 \lean{SimpleGraph.IsRepresentation.nonempty_iso_lineGraph_switch}
 \uses{def:represented,def:414-switch,def:2.2,thm:3.6,lem:4-lines,lem:4-transport,lem:414-lattice,lem:414-apex,lem:414-transitive,comp:414-Nc,facts:414}
 \assume{\begin{pfenum}
  \item $X$ is a connected $k$-regular graph on $n$ vertices with a representation $(\Zx_u)$ in a
   finite-dimensional real inner product space $V$.
  \item $\Lambda$, $W$ and $R$ are as in Lemma~\ref{lem:414-lattice}, and $\Zw$ is as in
   Lemma~\ref{lem:414-apex}.
  \item $t\in\{6,7,8\}$, and $f:\Phi(R)\to\Phi(\typ E_t)$ is a bijection that preserves inner
   products.
  \item $(\Zw,\Zw)=2$ and $n=|\Ecal_t|$.
 \end{pfenum}}
 \prove{There are a bijection $\pi:V(X)\to\Ecal_t$ and a set $H\sbs\Ecal_t$ such that $\pi$ is
 an isomorphism from $X$ to $L(\Delta_t)^H$.}
 \new\ These are steps $\langle1\rangle7$ to $\langle1\rangle9$ of
 \path{cgss/prerequisites.md}, Section 7.3; we read the switched graph off the inner products
 directly, without SW4.
\end{namedresult}
\begin{proof}
 \pfsketch\ The apex is a root, the $\Zx_u$ are a transversal of the pairs of $N_\Zw$, and an
 isometry takes $\Zw$ to $\Zc$, where the pairs are indexed by $\Ecal_t$.
 \step{s414trmax}{$R$ is maximal in $W$.}
 \begin{proof}\pf\ $R$ is nonempty and indecomposable and spans $W$ (Lemma~\ref{lem:414-lattice});
  it is isomorphic to $\typ E_t$ by \asm[3]{0}, which is in the list of
  Theorem~\ref{thm:3.6}.~\qed\end{proof}
 \step{s414trw}{$\Zw\in\Phi(R)$.}
 \begin{proof}\pf\ Lemma~\ref{lem:414-apex} applies, since $n=|\Ecal_t|\ge1$ by \asm[4]{0} and
  Computation~\ref{comp:414-Nc}; so $\Zw\in W$, and $(\Zw,\Zw)=2$ by \asm[4]{0}. For
  $\Zy\in\Phi(R)\sbs \Lambda$, $(\Zw,\Zy)$ is an integer (Lemma~\ref{lem:414-apex}) with
  $|(\Zw,\Zy)|\le|\Zw||\Zy|=2$ (\path{abs_real_inner_le_norm}); apply \stepref{s414trmax} and
  Definition~\ref{def:2.2}.~\qed\end{proof}
 \step{s414trg}{There is a bijection $g:\Phi(R)\to\Phi(\typ E_t)$ that preserves inner products,
  with $g(\Zw)=\Zc$.}
 \begin{proof}\pf\ $f(\Zw)\in\Phi(\typ E_t)$ by \stepref{s414trw}; Lemma~\ref{lem:414-transitive}
  gives $\psi$ with $\psi(f(\Zw))=\Zc$, and $g=\psi\circ f$ will do.~\qed\end{proof}
 \pflet{$\Zy_u=g(\Zx_u)$.}
 \step{s414trrep}{$(\Zy_u)$ represents $X$ in $\typ E_t$, and $(\Zy_u,\Zc)=1$ for every $u$.}
 \begin{proof}\pf\ $X$ is represented in $R$ by $(\Zx_u)$ (Lemma~\ref{lem:414-lattice}); apply
  Lemma~\ref{lem:4-transport} to $g$. Further
  $(\Zy_u,\Zc)=(g(\Zx_u),g(\Zw))=(\Zx_u,\Zw)=1$ by Lemma~\ref{lem:414-apex}.~\qed\end{proof}
 \step{s414trpick}{\pick\ $\pi(u)\in\Ecal_t$ for each $u$, and $D\sbs V(X)$, such that
  $\Zy_u=\Zc-\Ze_{\pi(u)}$ for $u\in D$ and $\Zy_u=\Ze_{\pi(u)}$ for $u\notin D$.}
 \begin{proof}\pf\ $\Zy_u\in N_\Zc$ by \stepref{s414trrep}; apply
  Computation~\ref{comp:414-Nc}.~\qed\end{proof}
 \step{s414trip}{For $P,P'\in\Ecal_t$, $(\Ze_P,\Ze_{P'})=(\Zc-\Ze_P,\Zc-\Ze_{P'})=|P\cap P'|$
  and $(\Ze_P,\Zc-\Ze_{P'})=1-|P\cap P'|$.}
 \begin{proof}\pf\ Expand, with $(\Zc,\Zc)=2$, $(\Zc,\Ze_i)=\frac12$ and
  $(\Ze_i,\Ze_j)=[i=j]$.~\qed\end{proof}
 \step{s414trinj}{$\pi$ is injective.}
 \begin{proof}\pf\ If $u\ne v$ and $\pi(u)=\pi(v)=P$, then $\Zy_u\ne\Zy_v$
  (Lemma~\ref{lem:4-lines}, assertion 1), so one is $\Ze_P$ and the other $\Zc-\Ze_P$, and their
  inner product is $1-2=-1$ by \stepref{s414trip}; this is against
  Definition~\ref{def:represented} and \stepref{s414trrep}.~\qed\end{proof}
 \step{s414trbij}{$\pi$ is a bijection.}
 \begin{proof}\pf\ \stepref{s414trinj}, and $|V(X)|=n=|\Ecal_t|$ by \asm[4]{0}.~\qed\end{proof}
 \pflet{$H=\pi(D)$.}
 \step{s414tradj}{Distinct $u,v$ are adjacent in $X$ if and only if $\pi(u)$ and $\pi(v)$ are
  adjacent in $L(\Delta_t)^H$.}
 \begin{proof}\pf\ By \stepref{s414trinj}, $|\pi(u)\cap\pi(v)|\le1$, and $u\in D$ exactly when
  $\pi(u)\in H$. By \stepref{s414trrep} and \stepref{s414trip}, $(\Zy_u,\Zy_v)=1$ exactly when
  $\pi(u),\pi(v)$ meet and both or neither lie in $H$, or are disjoint and exactly one lies in
  $H$; this is adjacency in $L(\Delta_t)^H$ (Definition~\ref{def:414-switch}).~\qed\end{proof}
 \qedstep
 \begin{proof}\pf\ \stepref{s414trbij} and \stepref{s414tradj}.~\qed\end{proof}
\end{proof}

\subsubsection*{Regular switchings}

\begin{namedresult}{Lemma}{(degrees in a switched line graph)}\label{lem:414-degree}
 \lean{SimpleGraph.card_adj_lineGraph_switch_of_complete}
 \uses{def:414-switch}
 \assume{\begin{pfenum}
  \item $Y$ is a graph, $M\sbs V(Y)$ has $m\ge2$ elements, and every 2-subset of $M$ is an edge
   of $Y$; $K$ is the set of these edges.
  \item $H\sbs E(Y)$, $h=|H\cap K|$, and $\deg_{H\cap K}$ is taken on $M$.
  \item $e=\{i,j\}\in K$ and $s=\deg_{H\cap K}(i)+\deg_{H\cap K}(j)$.
 \end{pfenum}}
 \prove{The number of neighbours of $e$ in $L(Y)^H$ that lie in $K$ is
 $2s-h-3+\binom{m-2}2$ if $e\in H$, and $2(m-2)+h-2s$ if $e\notin H$.}
 \new\ This replaces SW3 of \path{cgss/prerequisites.md} in step $\langle1\rangle10$ of its
 Section 7.3 by a direct count.
\end{namedresult}
\begin{proof}
 \pfsketch\ We will count the members of $K$ that meet $e$ and those that miss it, inside and
 outside $H$.
 \step{s414degmeet}{The members of $H\cap K$ other than $e$ that meet $e$ number $s-2$ if $e\in H$
  and $s$ if $e\notin H$.}
 \begin{proof}\pf\ No member of $K$ other than $e$ contains both $i$ and $j$, so the members of
  $H\cap K$ meeting $e$ number $\deg_{H\cap K}(i)+\deg_{H\cap K}(j)$, counting $e$ twice when
  $e\in H$.~\qed\end{proof}
 \step{s414degall}{The members of $K$ other than $e$ that meet $e$ number $2(m-2)$, and those
  disjoint from $e$ number $\binom{m-2}2$.}
 \begin{proof}\pf\ They are $\{i,l\}$ and $\{j,l\}$ with $l\in M\setminus e$, and the 2-subsets
  of $M\setminus e$.~\qed\end{proof}
 \step{s414degin}{\case{$e\in H$.}\prove{The count is $2s-h-3+\binom{m-2}2$.}}
 \begin{proof}\pf\ By Definition~\ref{def:414-switch}, the neighbours of $e$ in $K$ are the
  $f\in H\cap K$ that meet $e$, $s-2$ of them by \stepref{s414degmeet}, and the $f\in K\setminus H$
  disjoint from $e$. Of the $\binom{m-2}2$ members of $K$ disjoint from $e$
  (\stepref{s414degall}), $h-1-(s-2)$ lie in $H$. The total is
  $(s-2)+\binom{m-2}2-(h-s+1)$.~\qed\end{proof}
 \step{s414degout}{\case{$e\notin H$.}\prove{The count is $2(m-2)+h-2s$.}}
 \begin{proof}\pf\ By Definition~\ref{def:414-switch}, the neighbours are the $f\in K\setminus H$
  with $f\ne e$ that meet $e$, which number $2(m-2)-s$ by \stepref{s414degmeet} and
  \stepref{s414degall}, and the $f\in H\cap K$ disjoint from $e$, which number
  $h-s$.~\qed\end{proof}
 \qedstep
 \begin{proof}\pf\ The two cases are exhaustive.~\qed\end{proof}
\end{proof}

\begin{namedresult}{Lemma}{(constant pair sums)}\label{lem:414-constsum}
 \lean{Finset.two_mul_eq_of_forall_add_eq}
 \uses{}
 \assume{$M$ is a finite set with at least 3 elements, $d:M\to\ints$ and $\sigma\in\ints$, and
 $d(i)+d(j)=\sigma$ for all distinct $i,j\in M$.}
 \prove{$2d(i)=\sigma$ for every $i\in M$.}
\end{namedresult}
\begin{proof}
 \pfsketch\ Compare two pairs with a common member.
 \step{s414csconst}{$d(i)=d(j)$ for all $i,j\in M$.}
 \begin{proof}\pf\ For $i\ne j$ pick $l\in M\setminus\{i,j\}$; then
  $d(i)+d(l)=\sigma=d(j)+d(l)$ by \asm{0}.~\qed\end{proof}
 \qedstep
 \begin{proof}\pf\ For $i\in M$ pick $j\ne i$; then $\sigma=d(i)+d(j)=2d(i)$ by
  \stepref{s414csconst}.~\qed\end{proof}
\end{proof}

\begin{namedresult}{Theorem}{(SE4: regular graphs on 8 vertices)}\label{res:414-SE4}
 \lean{SimpleGraph.IsRegularOfDegree.nonempty_iso_cycle}
 \lean{SimpleGraph.IsRegularOfDegree.eq_bot_of_zero}
 \lean{SimpleGraph.IsRegularOfDegree.nonempty_iso_fourK2}
 \lean{SimpleGraph.IsRegularOfDegree.nonempty_iso_of_two_of_card_eq_eight}
 \lean{SimpleGraph.IsRegularOfDegree.nonempty_iso_of_three_of_card_eq_eight}
 \uses{}
 \assume{$Z$ is a finite graph.}
 \prove{\begin{pfenum}
  \item If $Z$ is connected and 2-regular, then $Z$ is a cycle.
  \item If $Z$ is $d$-regular on 8 vertices with $d\le3$, then $Z$ is isomorphic to one of the
   following eleven graphs: for $d=0$, $\comp{K_8}$; for $d=1$, $4K_2$; for $d=2$, $C_8$,
   $C_3+C_5$ or $2C_4$; for $d=3$, $2K_4$, the cube $Q_3$, the Wagner graph $M_8$, or the graphs
   $Y_1$, $Y_2$, $Y_4$ of \path{cgss/proofs/reg8/reg8.tex}, Definition~1.1.
 \end{pfenum}}
 This is an input, proved in \path{cgss/proofs/reg8/reg8.tex} (Lemma~2.1 and Theorem~1.2); in Lean its
 graphs are \texttt{SimpleGraph.fourK2}, \texttt{cycleThreeFive}, \texttt{twoC4}, \texttt{twoK4},
 \texttt{cubeGraph}, \texttt{wagnerGraph}, \texttt{truncatedK33Graph},
 \texttt{twiceTruncatedK4Graph} and \texttt{twoDiamondsGraph} on \texttt{Fin 8}. Its vertex $i$ is our vertex $i+1$.
\end{namedresult}
\begin{proof}
 \pfsketch\ We cite the structured case split of \path{cgss/proofs/reg8/reg8.tex}.
 \qedstep
 \begin{proof}\pf\ Lemma~2.1 of \path{cgss/proofs/reg8/reg8.tex} for assertion 1, and its Theorem~1.2
  for assertion 2.~\qed\end{proof}
\end{proof}

\begin{namedresult}{Lemma}{(2-regular graphs on 5 and 6 vertices)}\label{lem:414-tworeg}
 \lean{SimpleGraph.IsRegularOfDegree.nonempty_iso_cycle_of_two_of_card_eq_five}
 \lean{SimpleGraph.IsRegularOfDegree.nonempty_iso_of_two_of_card_eq_six}
 \uses{res:414-SE4}
 \assume{$Z$ is a 2-regular graph on $p\in\{5,6\}$ vertices.}
 \prove{If $p=5$ then $Z\cong C_5$; if $p=6$ then $Z\cong C_6$ or $Z\cong C_3+C_3$.}
\end{namedresult}
\begin{proof}
 \pfsketch\ The components are cycles of length at least 3.
 \step{s414twocomp}{Each component of $Z$ is a connected 2-regular graph with at least 3
  vertices.}
 \begin{proof}\pf\ The neighbours of a vertex lie in its component, so the component is
  2-regular; a vertex and its two neighbours are three vertices.~\qed\end{proof}
 \step{s414twocyc}{Each component of $Z$ is a cycle.}
 \begin{proof}\pf\ Theorem~\ref{res:414-SE4}, assertion 1, with
  \stepref{s414twocomp}.~\qed\end{proof}
 \qedstep
 \begin{proof}\pf\ By \stepref{s414twocomp} the orders of the components are at least 3 and add up
  to $p$: for $p=5$ this is $5$ alone, and for $p=6$ it is $6$ or $3+3$. Apply
  \stepref{s414twocyc}.~\qed\end{proof}
\end{proof}

\begin{namedresult}{Computation}{C22}\label{comp:C22}
 \lean{Computations.C22.e6_switch_cycle5}\lean{Computations.C22.e7_switch_cycle6}
 \lean{Computations.C22.e7_switch_twoC3}\lean{Computations.C22.e8_switch_fourK2}
 \lean{Computations.C22.e8_switch_twoC4}\lean{Computations.C22.e8_switch_cube}
 \lean{Computations.C22.e8_switch_twoDiamonds}\lean{Computations.C22.e8_switch_cycle8}
 \lean{Computations.C22.e8_switch_wagner}\lean{Computations.C22.e8_switch_twiceTruncatedK4}
 \lean{Computations.C22.e8_switch_cycleThreeFive}\lean{Computations.C22.e8_switch_truncatedK33}
 \lean{Computations.C22.e8_switch_twoK4}\lean{Computations.C22.e8_switch_empty}
 \lean{Computations.C22.cliqueProfile_t8}\lean{Computations.C22.cliqueProfile_changI}
 \lean{Computations.C22.cliqueProfile_changII}\lean{Computations.C22.cliqueProfile_changIII}
 \lean{Computations.C22.cliqueCount_shrikhande}
 \lean{Computations.C22.cliqueCount_lineGraph_completeBipartite}
 \uses{ex:4.5,ex:4.7,ex:4.9,def:414-switch,res:414-SE4,facts:414}
 \assume{$H_5$ is the edge set of the 5-cycle $1,2,3,4,5$; $H_6$ that of the 6-cycle
 $1,\dots,6$; $H_{3,3}=\{12,13,23,45,46,56\}$; the graphs of Theorem~\ref{res:414-SE4},
 assertion 2, are placed on $\{1,\dots,8\}$; and $\cprof(Z)$ is the multiset of the numbers of
 4-cliques of $Z$ through its vertices.}
 \prove{\begin{pfenum}
  \item $L(\Delta_6)^{H_5}$ is isomorphic to the Petersen graph of Example~\ref{ex:4.9}.
  \item $L(\Delta_7)^{H_6}$ is isomorphic to the Shrikhande graph of Example~\ref{ex:4.7}, and
   $L(\Delta_7)^{H_{3,3}}$ to $L(K_{4,4})$.
  \item For the graphs $Z$ of Theorem~\ref{res:414-SE4}, assertion 2, $L(\Delta_8)^{E(Z)}$ is
   isomorphic to $T(8)$ for $Z=\comp{K_8},2K_4$; to the Chang graph I of Example~\ref{ex:4.5}
   for $Z=4K_2,2C_4,Q_3,Y_4$; to the Chang graph II for $Z=C_8,M_8,Y_2$; and to the Chang graph
   III for $Z=C_3+C_5,Y_1$.
  \item $\cprof(T(8))=\ms{40^{28}}$, and $\cprof$ of the Chang graphs I, II and III is
   $\ms{32^4,36^{24}}$, $\ms{34^{24},36^4}$ and $\ms{33^{10},35^{18}}$.
  \item The Shrikhande graph has no 4-clique, and $L(K_{4,4})$ has 8.
 \end{pfenum}}
 These are SE5 and the part of SE6 that we need (\path{cgss/prerequisites.md}), with the
 4-clique profile in place of the maximal cliques (\path{cgss/certificates/README.md}, C22).
\end{namedresult}
\begin{proof}
 \pfsketch\ Each isomorphism is an explicit permutation of at most 28 vertices, found by
 SageMath (\path{cgss/certificates/c22_switching.sage}, certificates in
 \path{cgss/certificates/data/c22.json}); Lean checks it by \path{Computations.checkGraphIso} and
 \texttt{decide +kernel}. The profiles are evaluated by \path{Computations.cliqueProfile}, which
 enumerates only the 4-subsets that are cliques.
 \step{s414c22iso}{Assertions 1 to 3 hold.}
 \begin{proof}\raggedright\pf\ For each of the fourteen switched graphs, the adjacency table and
  the certificate permutation pass \path{Computations.checkGraphIso}; apply
  \path{Computations.nonempty_graphOfMatrix_iso}. For $Z=\comp{K_8}$ the permutation is the
  identity, since $L(\Delta_8)^{\varnothing}=T(8)$.~\qed\end{proof}
 \step{s414c22prof}{Assertions 4 and 5 hold.}
 \begin{proof}\raggedright\pf\ \path{Computations.cliqueProfile} and \path{Computations.cliqueCount} on the
  adjacency tables, with their soundness theorems; the profiles of assertion 4 are those of
  \path{Computations.Examples.t8_clique_profile} in
  \path{LineGraphs/Computations/CliqueBenchmarkT8.lean} and of its analogues for the Chang graphs in
  \path{LineGraphs/Computations/CliqueBenchmarkChangI.lean},
  \path{LineGraphs/Computations/CliqueBenchmarkChangII.lean} and
  \path{LineGraphs/Computations/CliqueBenchmarkChangIII.lean}.~\qed\end{proof}
 \qedstep
 \begin{proof}\pf\ \stepref{s414c22iso} and \stepref{s414c22prof}.~\qed\end{proof}
\end{proof}

\begin{namedresult}{Lemma}{(regular switchings of $T(5)$)}\label{lem:414-E6}
 \lean{SimpleGraph.nonempty_iso_petersen_of_switch_triangular5}
 \uses{def:414-switch,lem:414-switch,lem:414-degree,lem:414-constsum,lem:414-tworeg,comp:C22,ex:4.9,facts:414}
 \assume{$H$ is a subset of $\Ecal_6$, and the graph $L(\Delta_6)^H$ is 3-regular.}
 \prove{$L(\Delta_6)^H$ is isomorphic to the Petersen graph of Example~\ref{ex:4.9}.}
\end{namedresult}
\begin{proof}
 \pfsketch\ The degree count forces $H$ to be the edge set of a 2-regular graph on
 $\{1,\dots,5\}$, which is a 5-cycle.
 \pflet{$h=|H|$ and $M=\{1,\dots,5\}$.}
 \step{s414esixsum}{$2(\deg_H(i)+\deg_H(j))=h+3$ for all distinct $i,j\in M$.}
 \begin{proof}\pf\ Every vertex of $L(\Delta_6)$ lies in $K=\Ecal_6$. Lemma~\ref{lem:414-degree} with
  $m=5$, $\binom32=3$ and $2(m-2)=6$ gives degree $2s-h$ if $\{i,j\}\in H$ and $6+h-2s$ if not;
  both equal 3 by \asm{0}.~\qed\end{proof}
 \step{s414esixreg}{$4\deg_H(i)=h+3$ for every $i\in M$.}
 \begin{proof}\pf\ Lemma~\ref{lem:414-constsum}, applied to $d=2\deg_H$ on $M$ with $\sigma=h+3$,
  and \stepref{s414esixsum}: $2\cdot2\deg_H(i)=h+3$.~\qed\end{proof}
 \step{s414esixtwo}{$(M,H)$ is 2-regular.}
 \begin{proof}\pf\ Summing \stepref{s414esixreg} over $M$, with $\sum_i\deg_H(i)=2h$
  (\path{SimpleGraph.sum_degrees_eq_twice_card_edges}), gives $8h=5(h+3)$, so $h=5$ and
  $\deg_H(i)=2$.~\qed\end{proof}
 \step{s414esixrelabel}{$L(\Delta_6)^H\cong L(\Delta_6)^{H_5}$.}
 \begin{proof}\pf\ By \stepref{s414esixtwo} and Lemma~\ref{lem:414-tworeg}, some bijection $\sigma$ of
  $M$ carries $H$ onto $H_5$; extended by fixing 6, 7, 8, it is an automorphism of $\Delta_6$. Apply
  Lemma~\ref{lem:414-switch}, assertion 3.~\qed\end{proof}
 \qedstep
 \begin{proof}\pf\ \stepref{s414esixrelabel} and Computation~\ref{comp:C22}, assertion
  1.~\qed\end{proof}
\end{proof}

\begin{namedresult}{Lemma}{(regular switchings of $T(6)$ with an isolated vertex)}\label{lem:414-E7}
 \lean{SimpleGraph.nonempty_iso_shrikhande_or_lineGraph_completeBipartite_of_switch_triangular6AddOne}
 \uses{def:414-switch,lem:414-switch,lem:414-degree,lem:414-constsum,lem:414-tworeg,comp:C22,ex:4.7}
 \assume{$H$ is a subset of $\Ecal_7$, and the graph $L(\Delta_7)^H$ is 6-regular.}
 \prove{$L(\Delta_7)^H$ is isomorphic to the Shrikhande graph of Example~\ref{ex:4.7} or to
 $L(K_{4,4})$.}
\end{namedresult}
\begin{proof}
 \pfsketch\ After complementing $H$ we may assume that $\varepsilon\notin H$; then the degree
 of $\varepsilon$ gives $|H|=6$, and the degree count makes $H$ 2-regular on
 $\{1,\dots,6\}$.
 \pflet{$M=\{1,\dots,6\}$, $K=\Ecal_7\setminus\{\varepsilon\}$, and $H^*=H$ if
 $\varepsilon\notin H$, $H^*=\Ecal_7\setminus H$ if $\varepsilon\in H$.}
 \step{s414esevcompl}{$\varepsilon\notin H^*$, $H^*\sbs K$, and $L(\Delta_7)^H=L(\Delta_7)^{H^*}$.}
 \begin{proof}\pf\ The definition of $H^*$, and Lemma~\ref{lem:414-switch}, assertion 1, when
  $H^*=\Ecal_7\setminus H$.~\qed\end{proof}
 \step{s414esevsix}{$|H^*|=6$.}
 \begin{proof}\pf\ $\varepsilon$ is disjoint from every other member of $\Ecal_7$, so it is
  isolated in $L(\Delta_7)$; by Definition~\ref{def:414-switch} and \stepref{s414esevcompl} its
  neighbours in $L(\Delta_7)^{H^*}$ are the members of $H^*$. Its degree is 6 by \asm{0} and
  \stepref{s414esevcompl}.~\qed\end{proof}
 \step{s414esevsum}{$\deg_{H^*}(i)+\deg_{H^*}(j)=4$ for all distinct $i,j\in M$.}
 \begin{proof}\pf\ Let $e=\{i,j\}$ and $s=\deg_{H^*}(i)+\deg_{H^*}(j)$. By
  Lemma~\ref{lem:414-degree} with $m=6$, $\binom42=6$, $2(m-2)=8$ and $h=6$
  (\stepref{s414esevsix}), $e$ has $2s-3$ neighbours in $K$ if $e\in H^*$ and $14-2s$ if not. By
  Definition~\ref{def:414-switch}, $\varepsilon$ is a further neighbour exactly when $e\in H^*$.
  So the degree is $2s-2$ or $14-2s$, and it is 6.~\qed\end{proof}
 \step{s414esevtwo}{$(M,H^*)$ is 2-regular.}
 \begin{proof}\pf\ Lemma~\ref{lem:414-constsum} with $\sigma=4$ and
  \stepref{s414esevsum}.~\qed\end{proof}
 \step{s414esevrelabel}{$L(\Delta_7)^{H^*}\cong L(\Delta_7)^{H_6}$ or $L(\Delta_7)^{H^*}\cong
  L(\Delta_7)^{H_{3,3}}$.}
 \begin{proof}\pf\ By \stepref{s414esevtwo} and Lemma~\ref{lem:414-tworeg}, some bijection
  $\sigma$ of $M$ carries $H^*$ onto $H_6$ or onto $H_{3,3}$; extended by fixing 7 and 8, it is
  an automorphism of $\Delta_7$. Apply Lemma~\ref{lem:414-switch}, assertion 3.~\qed\end{proof}
 \qedstep
 \begin{proof}\pf\ \stepref{s414esevcompl}, \stepref{s414esevrelabel} and
  Computation~\ref{comp:C22}, assertion 2.~\qed\end{proof}
\end{proof}

\begin{namedresult}{Lemma}{(regular switchings of $T(8)$)}\label{lem:414-E8}
 \lean{SimpleGraph.nonempty_iso_triangular8_or_chang_of_switch_triangular8}
 \uses{def:414-switch,lem:414-switch,lem:414-degree,lem:414-constsum,res:414-SE4,comp:C22,ex:4.5,facts:414}
 \assume{$H$ is a subset of $\Ecal_8$, and the graph $L(\Delta_8)^H$ is 12-regular.}
 \prove{$L(\Delta_8)^H$ is isomorphic to $T(8)$ or to one of the Chang graphs I, II and III of
 Example~\ref{ex:4.5}.}
\end{namedresult}
\begin{proof}
 \pfsketch\ Here the degree count says only that $H$ is regular; after complementing, its valency
 is at most 3, and SE4 leaves eleven graphs.
 \pflet{$h=|H|$ and $M=\{1,\dots,8\}$.}
 \step{s414eeightsum}{$2(\deg_H(i)+\deg_H(j))=h$ for all distinct $i,j\in M$.}
 \begin{proof}\pf\ Every vertex of $L(\Delta_8)$ lies in $K=\Ecal_8$. Lemma~\ref{lem:414-degree} with
  $m=8$, $\binom62=15$ and $2(m-2)=12$ gives degree $2s-h+12$ if $\{i,j\}\in H$ and $12+h-2s$ if
  not; both equal 12 by \asm{0}.~\qed\end{proof}
 \step{s414eeightreg}{$(M,H)$ is regular of some valency $d\le7$.}
 \begin{proof}\pf\ Lemma~\ref{lem:414-constsum}, applied to $d'=2\deg_H$ on $M$ with $\sigma=h$, and
  \stepref{s414eeightsum} give $4\deg_H(i)=h$ for every $i\in M$; so $d=h/4$, and $d\le7$ since
  $|M|=8$.~\qed\end{proof}
 \pflet{$H^*=H$ if $d\le3$, and $H^*=\Ecal_8\setminus H$ if $d\ge4$.}
 \step{s414eeightcompl}{$(M,H^*)$ is regular of valency at most 3, and
  $L(\Delta_8)^H=L(\Delta_8)^{H^*}$.}
 \begin{proof}\pf\ $\Ecal_8\setminus H$ is the edge set of the complement of $(M,H)$, which is
  $(7-d)$-regular (\path{SimpleGraph.IsRegularOfDegree.compl}); apply
  Lemma~\ref{lem:414-switch}, assertion 1.~\qed\end{proof}
 \step{s414eeightrelabel}{$L(\Delta_8)^{H^*}\cong L(\Delta_8)^{E(Z)}$ for one of the eleven graphs $Z$ of
  Theorem~\ref{res:414-SE4}, assertion 2, placed on $M$.}
 \begin{proof}\pf\ By \stepref{s414eeightcompl} and Theorem~\ref{res:414-SE4}, some bijection
  $\sigma$ of $M$ is an isomorphism from $(M,H^*)$ to such a $Z$; it is an automorphism of
  $\Delta_8=K_8$. Apply Lemma~\ref{lem:414-switch}, assertion 3.~\qed\end{proof}
 \qedstep
 \begin{proof}\pf\ \stepref{s414eeightcompl}, \stepref{s414eeightrelabel} and
  Computation~\ref{comp:C22}, assertion 3.~\qed\end{proof}
\end{proof}

\subsubsection*{The Schl\"afli and Clebsch graphs}

\begin{namedresult}{Lemma}{(the triangle property from the parameters, GQ1)}\label{lem:414-GQ1}
 \lean{SimpleGraph.IsSRGWith.hasTriangleProperty}
 \uses{def:triangle-property}
 \assume{$Z$ is strongly regular with parameters $(n,k,1,c)$, and $n+3=3k$.}
 \prove{$Z$ has the triangle property.}
 \new\ This is the converse half of GQ1 of \path{cgss/prerequisites.md}, with the triangle
 property of Section 2 in place of a generalized quadrangle whose lines are the triangles;
 [BCN] Lemma 1.15.1 says only ``Straightforward''. The hypothesis is $n-3=3(k-2)$, written
 without subtraction as in Lean.
\end{namedresult}
\begin{proof}
 \pfsketch\ The edges leaving a triangle reach every other vertex exactly once.
 \step{s414gqone}{Every edge $xy$ has exactly one common neighbour $z$.}
 \begin{proof}\pf\ The parameter $a=1$.~\qed\end{proof}
 \step{s414gqmost}{A vertex $v\notin\{x,y,z\}$ is adjacent to at most one of $x,y,z$.}
 \begin{proof}\pf\ If $v$ were adjacent to two of them, say to $x$ and $y$, then $v$ and $z$ would
  be two common neighbours of the edge $xy$, against \stepref{s414gqone}; the other pairs are
  edges of the triangle as well.~\qed\end{proof}
 \qedstep
 \begin{proof}\pf\ Each of $x,y,z$ has $k-2$ neighbours outside $\{x,y,z\}$, so there are
  $3(k-2)=n-3$ edges from the triangle to the other $n-3$ vertices, by \asm{0}; by
  \stepref{s414gqmost} no vertex receives two, so every vertex receives exactly one. With
  \stepref{s414gqone}, this is Definition~\ref{def:triangle-property}.~\qed\end{proof}
\end{proof}

\begin{namedresult}{Definition}{(the folded 5-cube)}\label{def:414-folded}
 \lean{SimpleGraph.foldedFiveCube}\lean{SimpleGraph.compl_foldedFiveCube_adj_iff}
 \uses{}
 \assume{$\Omega_5=\{1,\dots,5\}$, and $V_{16}$ is the set of the subsets of $\Omega_5$ of even
 size.}
 \pflet{The \emph{folded 5-cube} $\square_5$ is the graph on $V_{16}$ in which $s\sim t$ when
 $|s\sd t|=4$.}
 This is Definition 1.1 of \path{cgss/proofs/clebsch/clebsch.tex}. $V_{16}$ consists of $\varnothing$,
 the ten 2-subsets and the five 4-subsets of $\Omega_5$. In Lean the vertex type is
 \texttt{$\{$s : Finset (Fin 5) // Even s.card$\}$}, whose \texttt{i} is our $i+1$. In the complement, distinct $s,t$ are adjacent
 exactly when $|s\sd t|=2$, since $|s\sd t|$ is even and at most 5.
\end{namedresult}

\begin{namedresult}{Definition}{(the double-six graph)}\label{def:414-gq24}
 \lean{SimpleGraph.gq24}\lean{SimpleGraph.compl_gq24_adj_iff}
 \uses{}
 \assume{$\Omega_6=\{1,\dots,6\}$, and $V_{27}$ is the set of the 27 symbols $a_i$, $b_i$
 ($i\in\Omega_6$) and $c_{ij}=c_{ji}$ ($i\ne j$ in $\Omega_6$).}
 \pflet{$Q_{27}$ be the graph on $V_{27}$ in which $a_i\sim b_j$ for $i\ne j$; $a_i\sim c_{jl}$ and
 $b_i\sim c_{jl}$ for $i\in\{j,l\}$; $c_{ij}\sim c_{lm}$ when $\{i,j\}\cap\{l,m\}=\varnothing$;
 and there are no other edges.}
 This is the graph $Q$ of \path{cgss/proofs/clebsch/clebsch.tex}, Definition 1.2, the collinearity
 graph of the generalized quadrangle of order $(2,4)$; we write $Q_{27}$, since $Q$ is $2I+A$ in
 this section. In Lean the vertex type is \texttt{Fin 2} $\times$ \texttt{Fin 6} $\oplus$
 \texttt{T} with \texttt{T = $\{$s : Finset (Fin 6) // s.card = 2$\}$}; $(0,i)$ is $a_{i+1}$, $(1,i)$
 is $b_{i+1}$, and $\{i,j\}$ is $c_{i+1,j+1}$. The second Lean name spells out adjacency in
 $\comp{Q_{27}}$ by the same cases.
\end{namedresult}

\begin{namedresult}{Computation}{(the models of SE2 and SE3)}\label{comp:414-models}
 \lean{SimpleGraph.foldedFiveCube_isSRGWith}\lean{SimpleGraph.gq24_isSRGWith}
 \uses{def:414-folded,def:414-gq24,comp:C9,ex:4.6,ex:4.8}
 \assume{$X_{\mathrm{Cl}}$ and $X_{\mathrm{Sch}}$ are the Clebsch and Schl\"afli graphs of
 Examples~\ref{ex:4.8} and~\ref{ex:4.6}, that is, the graphs of the vectors of their
 representations (Computation~\ref{comp:C9}).}
 \prove{\begin{pfenum}
  \item $\square_5$ is strongly regular with parameters $(16,5,0,2)$, and $Q_{27}$ with
   parameters $(27,10,1,5)$.
  \item $X_{\mathrm{Cl}}\cong\comp{\square_5}$ and $X_{\mathrm{Sch}}\cong\comp{Q_{27}}$.
 \end{pfenum}}
 The Clebsch and Schl\"afli graphs stay the graphs of the vectors; $\square_5$ and $Q_{27}$ are
 the combinatorial models of \path{cgss/proofs/clebsch/clebsch.tex}. The isomorphisms of assertion 2
 are the certificates of Computation~\ref{comp:C9}, which carry their Lean names there; we
 repeat them here beside the parameters.
\end{namedresult}
\begin{proof}
 \pfsketch\ These are finite checks on at most 27 vertices, as for Computation C9.
 \step{s414modpar}{Assertion 1 holds.}
 \begin{proof}\pf\ In Lean, \texttt{decide} counts the neighbours of every vertex and the common
  neighbours of every pair, at most $27^3$ triples. By hand, Lemmas 3.1 and 4.2 of
  \path{cgss/proofs/clebsch/clebsch.tex} count them at one vertex of each orbit of an automorphism
  group.~\qed\end{proof}
 \step{s414modiso}{Assertion 2 holds.}
 \begin{proof}\pf\ Computation~\ref{comp:C9}, its last sentence.~\qed\end{proof}
 \qedstep
 \begin{proof}\pf\ \stepref{s414modpar} and \stepref{s414modiso}.~\qed\end{proof}
\end{proof}

\begin{namedresult}{Lemma}{(the double-six graph from its parameters)}\label{lem:414-gq24iso}
 \lean{SimpleGraph.IsSRGWith.nonempty_iso_gq24}
 \uses{lem:414-GQ1,lem:triangle-iso,def:414-gq24,comp:414-models}
 \assume{$Z$ is a strongly regular graph with parameters $(27,10,1,5)$.}
 \prove{$Z\cong Q_{27}$.}
 \new\ SE3 of \path{cgss/prerequisites.md} identifies $Z$ by GQ2, the classification of the
 generalized quadrangles of order $(2,4)$; we use instead the isomorphism from the order of
 Section 2, as in Theorem 1.4 of \path{cgss/proofs/clebsch/clebsch.tex}.
\end{namedresult}
\begin{proof}
 \pfsketch\ Both graphs have the triangle property, and Section 2 shows that such a graph is
 determined by its order.
 \step{s414gqtri}{$Z$ and $Q_{27}$ have the triangle property.}
 \begin{proof}\pf\ Lemma~\ref{lem:414-GQ1}, with \asm{0} and Computation~\ref{comp:414-models},
  assertion 1, since $27+3=3\cdot10$.~\qed\end{proof}
 \step{s414gqedge}{$Z$ and $Q_{27}$ have 27 vertices each, an edge, and no vertex adjacent to
  every other vertex.}
 \begin{proof}\pf\ Both are 10-regular on 27 vertices, by \asm{0} and
  Computation~\ref{comp:414-models}, assertion 1, and $0<10<26$.~\qed\end{proof}
 \qedstep
 \begin{proof}\pf\ Lemma~\ref{lem:triangle-iso}, with \stepref{s414gqtri} and
  \stepref{s414gqedge}.~\qed\end{proof}
\end{proof}

\begin{namedresult}{Lemma}{(SE3: the Schl\"afli graph)}\label{lem:414-SE3}
 \lean{SimpleGraph.IsSRGWith.nonempty_iso_schlafli}
 \uses{lem:414-iso,lem:414-gq24iso,comp:414-models,ex:4.6,facts:414}
 \assume{$X$ is a strongly regular graph with parameters $(27,16,10,8)$.}
 \prove{$X$ is isomorphic to the Schl\"afli graph of Example~\ref{ex:4.6}.}
 \new\ SE3 of \path{cgss/prerequisites.md} goes through GQ1 and GQ2; we go through
 Lemma~\ref{lem:414-gq24iso} and the certificate of Computation~\ref{comp:414-models}.
\end{namedresult}
\begin{proof}
 \pfsketch\ We will pass to the complement, which is the double-six graph.
 \pflet{$X_{\mathrm{Sch}}$ be the Schl\"afli graph of Example~\ref{ex:4.6}.}
 \step{s414sthreecompl}{$\comp X$ is strongly regular with parameters $(27,10,1,5)$.}
 \begin{proof}\pf\ \path{SimpleGraph.IsSRGWith.compl}: $(27,27-16-1,27-(32-8)-2,27-(32-10))$, with no
  truncated subtraction.~\qed\end{proof}
 \step{s414sthreeiso}{$\comp X\cong Q_{27}$.}
 \begin{proof}\pf\ Lemma~\ref{lem:414-gq24iso} with \stepref{s414sthreecompl}.~\qed\end{proof}
 \qedstep
 \begin{proof}\pf\ Lemma~\ref{lem:414-iso}, assertion 1 (\texttt{SimpleGraph.Iso.compl}), applied
  to \stepref{s414sthreeiso}, with \path{compl_compl}, gives $X\cong\comp{Q_{27}}$; and
  $\comp{Q_{27}}\cong X_{\mathrm{Sch}}$ by Computation~\ref{comp:414-models},
  assertion 2.~\qed\end{proof}
\end{proof}

\begin{namedresult}{Lemma}{(the folded 5-cube from its parameters)}\label{lem:414-foldiso}
 \lean{SimpleGraph.IsSRGWith.nonempty_iso_foldedFiveCube}
 \uses{def:414-folded}
 \assume{$Z$ is a strongly regular graph with parameters $(16,5,0,2)$.}
 \prove{$Z\cong\square_5$.}
 This is the first half of SE2 of \path{cgss/prerequisites.md}, after [C17] Theorem 2.2, and
 Theorem 1.3 of \path{cgss/proofs/clebsch/clebsch.tex}. [BCN] 3.11 (iii) prints these parameters as
 $(16,5,0,1)$; they are $(16,5,0,2)$.
\end{namedresult}
\begin{proof}
 \pfsketch\ The non-neighbours of a vertex $u$ correspond to the 2-subsets of its
 neighbourhood, and adjacency among them is disjointness.
 \pflet{$u$ be a vertex of $Z$, $N$ its neighbourhood and $O=V(Z)\setminus(\{u\}\cup N)$.}
 \step{s414foldcocl}{$N$ is a coclique of 5 vertices, and $|O|=10$.}
 \begin{proof}\pf\ $|N|=5$ and $|O|=16-1-5$. Two adjacent vertices of $N$ would have the common
  neighbour $u$, against $a=0$.~\qed\end{proof}
 \step{s414foldphi}{For $w\in O$, $\phi(w)=N(w)\cap N$ has exactly two elements.}
 \begin{proof}\pf\ $w\ne u$ and $w\not\sim u$, so they have $c=2$ common neighbours, which are
  the members of $N(w)\cap N$.~\qed\end{proof}
 \step{s414foldbij}{$\phi$ is a bijection from $O$ onto $\binom N2$.}
 \begin{proof}\pf\ If $w\ne w'$ in $O$ had $\phi(w)=\phi(w')=\{x,y\}$, the distinct non-adjacent
  vertices $x,y$ (\stepref{s414foldcocl}) would have the three common neighbours $u,w,w'$,
  against $c=2$. So $\phi$ is injective, and $|O|=10=\binom52$.~\qed\end{proof}
 \step{s414foldmeet}{\assume{$w,w'\in O$ and $\phi(w)\cap\phi(w')\ne\varnothing$.}
  \prove{$w\not\sim w'$.}}
 \begin{proof}\pf\ Take $x\in\phi(w)\cap\phi(w')$ by \asm{1}; an edge $ww'$ would have the common
  neighbour $x$, against $a=0$.~\qed\end{proof}
 \step{s414foldO}{Each $w\in O$ has exactly 3 neighbours in $O$.}
 \begin{proof}\pf\ $w$ has 5 neighbours; two lie in $N$ (\stepref{s414foldphi}), and $u$ is not
  one.~\qed\end{proof}
 \step{s414folddisj}{Distinct $w,w'\in O$ are adjacent if and only if $\phi(w)$ and $\phi(w')$
  are disjoint.}
 \begin{proof}\pf\ Exactly $\binom32=3$ members of $\binom N2$ are disjoint from $\phi(w)$, so by
  \stepref{s414foldbij} exactly 3 vertices $w'\in O$ have $\phi(w')$ disjoint from $\phi(w)$. By
  \stepref{s414foldmeet} the 3 neighbours of $w$ in $O$ (\stepref{s414foldO}) are among them;
  so they are all of them.~\qed\end{proof}
 \pflet{$\beta:N\to\Omega_5$ be a bijection, and $f:V(Z)\to V_{16}$ the map with $f(u)=\varnothing$,
 $f(x)=\Omega_5\setminus\{\beta(x)\}$ for $x\in N$ and $f(w)=\beta(\phi(w))$ for $w\in O$.}
 \step{s414foldf}{$f$ is a bijection.}
 \begin{proof}\pf\ $f$ maps $u$ to $\varnothing$, $N$ onto the five 4-subsets, since $\beta$ is a
  bijection, and $O$ onto the ten 2-subsets, by \stepref{s414foldbij}; these are all of $V_{16}$
  (Definition~\ref{def:414-folded}).~\qed\end{proof}
 \step{s414foldsd}{For $i\in\Omega_5$ and 2-subsets $P\ne P'$ of $\Omega_5$,
  $|(\Omega_5\setminus\{i\})\sd P|$ is 4 if $i\in P$ and 2 if $i\notin P$, and
  $|P\sd P'|=4-2|P\cap P'|$.}
 \begin{proof}\pf\ $|s\sd s'|=|s|+|s'|-2|s\cap s'|$, and $(\Omega_5\setminus\{i\})\cap P$ has 1
  element if $i\in P$ and 2 if not.~\qed\end{proof}
 \step{s414foldadj}{Distinct vertices $p,q$ of $Z$ are adjacent if and only if
  $|f(p)\sd f(q)|=4$.}
 \begin{proof}
  \step{s414foldcaseone}{\case{$p=u$ and $q=x\in N$.}\prove{Both hold.}}
  \begin{proof}\pf\ $u\sim x$, and $|\varnothing\sd(\Omega_5\setminus\{\beta(x)\})|=4$.~\qed\end{proof}
  \step{s414foldcasetwo}{\case{$p=u$ and $q=w\in O$.}\prove{Neither holds.}}
  \begin{proof}\pf\ $w\notin N$, and $|\varnothing\sd f(w)|=|f(w)|=2$.~\qed\end{proof}
  \step{s414foldcasethree}{\case{$p=x$ and $q=y$ lie in $N$.}\prove{Neither holds.}}
  \begin{proof}\pf\ $x\not\sim y$ by \stepref{s414foldcocl}, and
   $|f(x)\sd f(y)|=|\{\beta(x),\beta(y)\}|=2$.~\qed\end{proof}
  \step{s414foldcasefour}{\case{$p=x\in N$ and $q=w\in O$.}\prove{The two conditions agree.}}
  \begin{proof}\pf\ $x\sim w$ exactly when $x\in\phi(w)$, that is, when $\beta(x)\in f(w)$, and
   by \stepref{s414foldsd} this is when $|f(x)\sd f(w)|=4$.~\qed\end{proof}
  \step{s414foldcasefive}{\case{$p=w$ and $q=w'$ lie in $O$.}\prove{The two conditions
   agree.}}
  \begin{proof}\pf\ $w\sim w'$ exactly when $f(w)\cap f(w')=\varnothing$, by
   \stepref{s414folddisj} and the injectivity of $\beta$, and by \stepref{s414foldsd} this is
   when $|f(w)\sd f(w')|=4$.~\qed\end{proof}
  \qedstep
  \begin{proof}\pf\ $V(Z)$ is the disjoint union of $\{u\}$, $N$ and $O$, so up to the order of
   $p$ and $q$ the five cases are exhaustive.~\qed\end{proof}
 \end{proof}
 \qedstep
 \begin{proof}\pf\ By \stepref{s414foldf} and \stepref{s414foldadj}, $f$ is an isomorphism from
  $Z$ to $\square_5$ (Definition~\ref{def:414-folded}).~\qed\end{proof}
\end{proof}

\begin{namedresult}{Lemma}{(SE2: the Clebsch graph)}\label{lem:414-SE2}
 \lean{SimpleGraph.IsSRGWith.nonempty_iso_clebsch}
 \uses{lem:414-iso,lem:414-foldiso,comp:414-models,ex:4.8,facts:414}
 \assume{$X$ is a strongly regular graph with parameters $(16,10,6,6)$.}
 \prove{$X$ is isomorphic to the Clebsch graph of Example~\ref{ex:4.8}.}
 This is SE2 of \path{cgss/prerequisites.md}.
\end{namedresult}
\begin{proof}
 \pfsketch\ We will pass to the complement, which is the folded 5-cube.
 \pflet{$X_{\mathrm{Cl}}$ be the Clebsch graph of Example~\ref{ex:4.8}.}
 \step{s414stwocompl}{$\comp X$ is strongly regular with parameters $(16,5,0,2)$.}
 \begin{proof}\pf\ \path{SimpleGraph.IsSRGWith.compl}: $(16,16-10-1,16-(20-6)-2,16-(20-6))$, with no
  truncated subtraction.~\qed\end{proof}
 \step{s414stwofold}{$\comp X\cong\square_5$.}
 \begin{proof}\pf\ Lemma~\ref{lem:414-foldiso} with \stepref{s414stwocompl}.~\qed\end{proof}
 \qedstep
 \begin{proof}\pf\ Lemma~\ref{lem:414-iso}, assertion 1 (\texttt{SimpleGraph.Iso.compl}), applied
  to \stepref{s414stwofold}, with \path{compl_compl}, gives $X\cong\comp{\square_5}$; and
  $\comp{\square_5}\cong X_{\mathrm{Cl}}$ by Computation~\ref{comp:414-models},
  assertion 2.~\qed\end{proof}
\end{proof}

\subsubsection*{The seven graphs}

\begin{namedresult}{Proposition}{(the seven graphs)}\label{prop:414-seven}
 \lean{SimpleGraph.seidel_seven_properties}\lean{SimpleGraph.seidel_seven_pairwise_not_iso}
 \uses{def:CP,def:represented,lem:4-gram,comp:C9,comp:C22,ex:4.5,ex:4.6,ex:4.7,ex:4.8,ex:4.9,facts:414}
 \assume{$\cS$ is the list of the seven graphs of Examples 4.5 to 4.9: the Chang graphs I, II
 and III, the Schl\"afli graph, the Shrikhande graph, the Clebsch graph and the Petersen
 graph.}
 \prove{\begin{pfenum}
  \item Each $Z\in\cS$ is connected and strongly regular, has $\lambda_{\min}(Z)=-2$, and is
   represented in $\typ E_8$.
  \item No $Z\in\cS$ is isomorphic to $L(K_m)$, to $L(K_{m,m})$ or to $CP(m)$, for any $m$.
  \item No two graphs of $\cS$ are isomorphic.
 \end{pfenum}}
 \repair\ ``The seven graphs'' presupposes that the three Chang graphs are pairwise
 non-isomorphic, which the paper does not show (errata, Gap 12). We will separate them by their
 4-clique profiles (Computation C22), which the kernel can evaluate by extending cliques one
 vertex at a time; their numbers of maximal cliques, used in SE6 of
 \path{cgss/prerequisites.md}, can be enumerated in the same way, but every clique found must
 then be checked for maximality, which costs the kernel much more.
\end{namedresult}
\begin{proof}
 \pfsketch\ The parameters separate all pairs except the ones sharing $(n,k)$, and C22 separates
 those.
 \step{s414sevparams}{Each $Z\in\cS$ is represented in $\typ E_8$, and is strongly regular with
  parameters $(28,12,6,4)$ (Chang), $(27,16,10,8)$, $(16,6,2,2)$, $(16,10,6,6)$ and $(10,3,0,1)$.}
 \begin{proof}\pf\ Examples~\ref{ex:4.5} to~\ref{ex:4.9}, and
  Computation~\ref{comp:C9}.~\qed\end{proof}
 \step{s414sevconn}{Each $Z\in\cS$ is connected.}
 \begin{proof}\pf\ In each case $c\ge1$ by \stepref{s414sevparams}, so any two distinct
  non-adjacent vertices have a common neighbour.~\qed\end{proof}
 \step{s414sevlmin}{$\lambda_{\min}(Z)=-2$ for each $Z\in\cS$.}
 \begin{proof}\pf\ $Z$ is represented in $\Phi(\typ E_8)\sbs\re^8$ (\stepref{s414sevparams}), a
  space of dimension $8<10\le|V(Z)|$; apply Lemma~\ref{lem:4-gram}, assertion
  2.~\qed\end{proof}
 \step{s414sevpar}{$L(K_m)$ has $\binom m2$ vertices and valency $2(m-2)$, $L(K_{m,m})$ has $m^2$
  vertices and valency $2(m-1)$, and $CP(m)$ has $2m$ vertices and valency $2m-2$.}
 \begin{proof}\pf\ An edge $\{i,j\}$ of $K_m$ meets the $2(m-2)$ edges $\{i,l\}$, $\{j,l\}$; an
  edge of $K_{m,m}$ meets $m-1$ others at each end; and Definition~\ref{def:CP}.~\qed\end{proof}
 \step{s414sevnum}{The Schl\"afli, Clebsch and Petersen graphs are isomorphic to none of $L(K_m)$,
  $L(K_{m,m})$, $CP(m)$, and the Chang and Shrikhande graphs to no $CP(m)$.}
 \begin{proof}\pf\ Compare \stepref{s414sevparams} with \stepref{s414sevpar}. For 27 vertices,
  $\binom m2=27$, $m^2=27$ and $2m=27$ have no solution. For 16 vertices, $\binom m2=16$ has none,
  and $m^2=16$, $2m=16$ give valencies 6 and 14, not 10. For 10 vertices, $\binom m2=10$ gives
  valency 6, $m^2=10$ has no solution and $2m=10$ gives 8, not 3. Finally $2m=28$ and $2m=16$
  give valencies 26 and 14, not 12 and 6.~\qed\end{proof}
 \step{s414sevchang}{The Chang graphs I, II and III and $T(8)$ are pairwise non-isomorphic.}
 \begin{proof}\pf\ Their 4-clique profiles differ (Computation~\ref{comp:C22}, assertion 4), and
  an isomorphism preserves the profile (\path{SimpleGraph.Iso.cliqueProfile_eq}).~\qed\end{proof}
 \step{s414sevshr}{The Shrikhande graph is not isomorphic to $L(K_{4,4})$.}
 \begin{proof}\pf\ By Computation~\ref{comp:C22}, assertion 5, $L(K_{4,4})$ has a 4-clique; an
  isomorphism from $L(K_{4,4})$ to the Shrikhande graph would map it to a 4-clique
  (\path{SimpleGraph.IsNClique.map}), which the Shrikhande graph does not have.~\qed\end{proof}
 \step{s414sevlk}{No Chang graph is isomorphic to $L(K_{m,m})$ or to $L(K_m)$ with $m\ne8$, and the
  Shrikhande graph is isomorphic to no $L(K_m)$ and to no $L(K_{m,m})$ with $m\ne4$.}
 \begin{proof}\pf\ By \stepref{s414sevparams} and \stepref{s414sevpar}: $m^2=28$ has no solution,
  $\binom m2=28$ only $m=8$, $\binom m2=16$ none, and $m^2=16$ only $m=4$.~\qed\end{proof}
 \qedstep
 \begin{proof}\pf\ Assertion 1 is \stepref{s414sevparams}, \stepref{s414sevconn} and
  \stepref{s414sevlmin}; assertion 2 is \stepref{s414sevnum}, \stepref{s414sevlk},
  \stepref{s414sevchang} and \stepref{s414sevshr}. For assertion 3, an isomorphism preserves
  $(n,k)$, and by \stepref{s414sevparams} only the Chang graphs share it; apply
  \stepref{s414sevchang}.~\qed\end{proof}
\end{proof}

\begin{namedresult}{Corollary}{(the seven graphs are not line graphs)}\label{cor:414-notline}
 \lean{SimpleGraph.seidelSeven_not_iso_lineGraph}
 \uses{prop:414-seven,lem:414-linegraph}
 \assume{$\cS$ is the list of the seven graphs of Proposition~\ref{prop:414-seven}, and $Y$ is a
 graph, possibly infinite or with isolated vertices.}
 \prove{No $Z\in\cS$ is isomorphic to $L(Y)$.}
 This is the claim of the paper, before Example 4.5, that these strongly regular graphs are not
 line graphs.
\end{namedresult}
\begin{proof}
 \pfsketch\ We will apply Lemma~\ref{lem:414-linegraph}; assertion 2 of
 Proposition~\ref{prop:414-seven} excludes both of its alternatives.
 \step{s414nlcase}{\assume{$Z\in\cS$ and $Z\cong L(Y)$.}\prove{A contradiction.}}
 \begin{proof}
  \step{s414nldet}{$Z$ is connected and strongly regular, and $\det(2I+A(Z))=0$.}
  \begin{proof}\pf\ Proposition~\ref{prop:414-seven}, assertion 1, with $Z\in\cS$ by \asm{1};
   since $\lambda_{\min}(Z)=-2$ is an eigenvalue of $A(Z)$, the matrix $2I+A(Z)$ is
   singular.~\qed\end{proof}
  \step{s414nllg}{$Z\cong L(K_m)$ or $Z\cong L(K_{m,m})$ for some $m\ge2$.}
  \begin{proof}\pf\ Lemma~\ref{lem:414-linegraph}, with \stepref{s414nldet} and
   \asm{1}.~\qed\end{proof}
  \qedstep
  \begin{proof}\pf\ Proposition~\ref{prop:414-seven}, assertion 2, excludes both alternatives
   of \stepref{s414nllg}.~\qed\end{proof}
 \end{proof}
 \qedstep
 \begin{proof}\pf\ Step \stepref{s414nlcase} applies to every $Z\in\cS$.~\qed\end{proof}
\end{proof}

\subsubsection*{The classification}

\begin{namedresult}{Remark}{4.14, restated}\label{rem:414-proof}
 \lean{SimpleGraph.seidel_seven_classification}
 \uses{rem:4.14,def:represented,def:CP,def:3.2,def:3.3,lem:4-gramfactor,lem:4-transport,lem:4-regglg,thm:4.1,thm:4.2,thm:4.4,thm:3.5,lem:sec3-transport,lem:sec3-spans,lem:414-CP,lem:414-srg,comp:C18,lem:414-linegraph,lem:414-lattice,lem:414-apex,lem:414-dual,comp:C21,comp:414-Nc,lem:414-transversal,lem:414-E6,lem:414-E7,lem:414-E8,lem:414-SE2,lem:414-SE3,comp:414-models,prop:414-seven,facts:414}
 \assume{\begin{pfenum}
  \item $X$ is a connected graph, strongly regular with parameters $(n,k,a,c)$.
  \item $\lambda_{\min}(X)=-2$.
  \item $X$ is represented in $\typ E_8$ by $(\Zx_u)$.
  \item For no $m$ is $X$ isomorphic to $L(K_m)$, to $L(K_{m,m})$ or to $CP(m)$.
 \end{pfenum}}
 \prove{\begin{pfenum}
  \item $X$ is isomorphic to exactly one of the Chang graphs I, II and III, the Schl\"afli graph,
   the Shrikhande graph, the Clebsch graph and the Petersen graph.
  \item These seven graphs satisfy hypotheses 1 to 4.
 \end{pfenum}}
 \new\ The paper cites Seidel [19] and gives no proof. We follow the route of
 \path{cgss/prerequisites.md}, Section 7.3 (decision 3 of its Section 13); the departures are
 tagged in the lemmas above. \repair\ The proof fills Gap 12 (errata): it proves the
 classification, and ``exactly one'' is the non-isomorphism of Proposition~\ref{prop:414-seven}.
 In Lean, the statement \texttt{seidel\_seven} of Remark~\ref{rem:4.14} is the
 conjunction of \texttt{seidel\_seven\_classification}, of
 \texttt{seidel\_seven\_pairwise\_not\_iso} and of \texttt{seidel\_seven\_properties}
 (Proposition~\ref{prop:414-seven}), all four in the namespace \texttt{SimpleGraph}; Lean
 names are unique, so the two results carry different ones.
\end{namedresult}
\begin{proof}
 \pfsketch\ We will pin down the parameters, find the type of the lattice roots $R$, exclude the
 types $\typ A$ and $\typ D$ by the line graph theorems, exclude two parameter sets by the
 integrality of the apex, and identify the remaining five cases by switching and by uniqueness.
 \pflet{$A=A(X)$, $Q=2I+A$, $\rho=\rk Q$, and let $\Lambda$, $W$, $R$ be as in
 Lemma~\ref{lem:414-lattice} and $\Zw$ as in Lemma~\ref{lem:414-apex} for $(\Zx_u)$.}
 \step{s414mdet}{$n\ge1$ and $\det Q=0$.}
 \begin{proof}\pf\ A connected graph has a vertex; \asm[2]{0} and
  \path{SimpleGraph.leastEigenvalue_eq_neg_two_iff}.~\qed\end{proof}
 \step{s414mbounds}{$n\le28$ and $k\le16$.}
 \begin{proof}\pf\ Theorem~\ref{thm:4.4}, with \asm[1]{0}, \asm[3]{0} and
  \stepref{s414mdet}.~\qed\end{proof}
 \step{s414mrank}{$\rho=\dim W\le8$.}
 \begin{proof}\pf\ $Q$ is the Gram matrix of $(\Zx_u)$, so $\rho=\dim W$ by
  Lemma~\ref{lem:4-gramfactor}, assertion 1; and $W\sbs\re^8$
  (\path{Submodule.finrank_le}).~\qed\end{proof}
 \step{s414msrg}{$1\le k\le n-2$, $a\le k-1$, $1\le c\le k$, and conditions 2 to 4 of
  Computation~\ref{comp:C18} hold.}
 \begin{proof}\pf\ Lemma~\ref{lem:414-srg} with \asm[1]{0} and \stepref{s414mdet}.~\qed\end{proof}
 \step{s414mparam}{$k(k-a-1)=c(n-k-1)$, condition 1 of Computation~\ref{comp:C18}.}
 \begin{proof}\pf\ \path{SimpleGraph.IsSRGWith.param_eq}, which needs $0<n$ (\stepref{s414mdet});
  it is stated in $\mathbb N$ with truncated subtraction, and by \stepref{s414msrg} neither
  $k-a-1$ nor $n-k-1$ is truncated.~\qed\end{proof}
 \step{s414mCeighteen}{$(n,k,a,c,\rho)$ is one of the 15 tuples of Computation C18.}
 \begin{proof}\pf\ Computation~\ref{comp:C18}: its hypotheses are \stepref{s414msrg},
  \stepref{s414mparam}, and $n\le28\le36$, $k\le16$ and $\rho\le8$ from \stepref{s414mbounds} and
  \stepref{s414mrank}.~\qed\end{proof}
 \step{s414mnotCP}{$(n,k,a,c)\ne(2m,2m-2,2m-4,2m-2)$ for every $m$.}
 \begin{proof}\pf\ Otherwise $X$ is $(2m-2)$-regular on $2m$ vertices, and $X\cong CP(m)$ by
  Lemma~\ref{lem:414-CP}, against \asm[4]{0}.~\qed\end{proof}
 \step{s414mR}{$R$ is a nonempty, indecomposable, star-closed line system spanning $W$, and $X$
  is represented in $R$ by $(\Zx_u)$.}
 \begin{proof}\pf\ Lemma~\ref{lem:414-lattice}, with \asm[1]{0} and \asm[3]{0}.~\qed\end{proof}
 \step{s414mtype}{\pick\ a type $\tau$ among $\typ A_\ell$ ($\ell\ge1$), $\typ D_\ell$
  ($\ell\ge4$), $\typ E_6$, $\typ E_7$, $\typ E_8$, and a bijection $f:\Phi(R)\to\Phi(\tau)$
  preserving inner products.}
 \begin{proof}\pf\ Theorem~\ref{thm:3.5} with \stepref{s414mR}.~\qed\end{proof}
 \step{s414mspan}{\pick\ a linear isometry $F$ from $W$ onto $\opk{span}\Phi(\tau)$ that extends
  $f$.}
 \begin{proof}\pf\ Lemma~\ref{lem:sec3-transport} with \stepref{s414mtype}, since $\Phi(R)$ spans
  $W$ (\stepref{s414mR}).~\qed\end{proof}
 \step{s414mrk}{The rank of $\tau$ is $\rho$.}
 \begin{proof}\pf\ $F$ of \stepref{s414mspan} is an isometry onto $\opk{span}\Phi(\tau)$, and
  $\dim W=\rho$ by \stepref{s414mrank}.~\qed\end{proof}
 \step{s414mAD}{\assume{$\tau$ is $\typ A_\ell$ or $\typ D_\ell$.}\prove{A contradiction.}}
 \begin{proof}
  \step{s414mA}{\case{$\tau=\typ A_\ell$.}\prove{$X\cong L(Y)$ for some graph $Y$.}}
  \begin{proof}\pf\ $X$ is represented in $\typ A_\ell$ by Lemma~\ref{lem:4-transport} with $f$
   and \stepref{s414mR}; apply Theorem~\ref{thm:4.1}.~\qed\end{proof}
  \step{s414mD}{\case{$\tau=\typ D_\ell$.}\prove{$X\cong L(Y)$ for some graph $Y$, or $X\cong
   CP(b)$ for some $b$.}}
  \begin{proof}\pf\ $X$ is represented in $\typ D_\ell$ by Lemma~\ref{lem:4-transport} with $f$
   and \stepref{s414mR}, so it is a generalized line graph by Theorem~\ref{thm:4.2}, assertion 1;
   it is connected and regular by \asm[1]{0}, so apply
   Lemma~\ref{lem:4-regglg}.~\qed\end{proof}
  \step{s414mnotL}{$X\not\cong L(Y)$ for every graph $Y$.}
  \begin{proof}\pf\ Otherwise Lemma~\ref{lem:414-linegraph}, with \asm[1]{0} and
   \stepref{s414mdet}, gives $X\cong L(K_m)$ or $L(K_{m,m})$, against \asm[4]{0}.~\qed\end{proof}
  \qedstep
  \begin{proof}\pf\ In both cases \stepref{s414mA}, \stepref{s414mD}, the conclusion contradicts
   \stepref{s414mnotL} or \asm[4]{0}.~\qed\end{proof}
 \end{proof}
 \step{s414mE}{$\tau=\typ E_t$ with $t=\rho\in\{6,7,8\}$, and $(n,k,a,c)$ is one of
  $(10,3,0,1)$, $(15,8,4,4)$, $(16,10,6,6)$ (for $t=6$), $(16,6,2,2)$, $(21,10,5,4)$,
  $(27,16,10,8)$ (for $t=7$) and $(28,12,6,4)$ (for $t=8$).}
 \begin{proof}\pf\ By \stepref{s414mAD}, $\tau$ is one of $\typ E_6$, $\typ E_7$, $\typ E_8$, of
  ranks 6, 7, 8 (Lemma~\ref{lem:sec3-spans}), so $t=\rho$ by \stepref{s414mrk}. Of the 15
  tuples of \stepref{s414mCeighteen}, \stepref{s414mnotCP} removes six, and $\rho\in\{6,7,8\}$
  removes $(9,4,1,2,5)$ and $(10,6,3,4,5)$.~\qed\end{proof}
 \step{s414mint}{$\det(B_t)\,n/(k+2)$ is an integer, where $\det B_t$ is 3, 2, 1 for $t=6,7,8$.}
 \begin{proof}
  \step{s414mbasis}{\pick\ $\Zb_i=F^{-1}(\Zs_i)$, $1\le i\le t$, with the roots $\Zs_i$ of
   Computation C21; they lie in $\Phi(R)$ and form a basis of $W$ with Gram matrix $B_t$.}
  \begin{proof}\pf\ Computation~\ref{comp:C21} gives the $\Zs_i$, a basis of
   $\opk{span}\Phi(\typ E_t)$; $F^{-1}$ is a surjective linear isometry onto $W$ that maps
   $\Phi(\typ E_t)$ onto $\Phi(R)$ (\stepref{s414mtype} and
   \stepref{s414mspan}).~\qed\end{proof}
  \step{s414mwb}{$\Zw\in W$, $(\Zw,\Zb_i)\in\ints$ for every $i$, and $(\Zw,\Zw)=n/(k+2)$.}
  \begin{proof}\pf\ Lemma~\ref{lem:414-apex}, with \asm[1]{0} and \stepref{s414mdet}; the $\Zb_i$
   lie in $\Phi(R)\sbs \Lambda$ by \stepref{s414mbasis}.~\qed\end{proof}
  \qedstep
  \begin{proof}\pf\ Lemma~\ref{lem:414-dual}, with \stepref{s414mbasis} and \stepref{s414mwb};
   the determinants are those of Computation~\ref{comp:C21}.~\qed\end{proof}
 \end{proof}
 \step{s414mexcl}{$(n,k,a,c)$ is neither $(15,8,4,4)$ nor $(21,10,5,4)$.}
 \begin{proof}\pf\ These have $t=6$ and $t=7$ by \stepref{s414mE}, and
  $3\cdot\frac{15}{10}=\frac92$ and $2\cdot\frac{21}{12}=\frac72$ are not integers, against
  \stepref{s414mint}.~\qed\end{proof}
 \step{s414mclebsch}{\case{$(n,k,a,c)=(16,10,6,6)$.}\prove{$X$ is isomorphic to the Clebsch
  graph.}}
 \begin{proof}\pf\ Lemma~\ref{lem:414-SE2} with \asm[1]{0}.~\qed\end{proof}
 \step{s414mschlafli}{\case{$(n,k,a,c)=(27,16,10,8)$.}\prove{$X$ is isomorphic to the Schl\"afli
  graph.}}
 \begin{proof}\pf\ Lemma~\ref{lem:414-SE3} with \asm[1]{0}.~\qed\end{proof}
 \step{s414mpet}{\case{$(n,k,a,c)=(10,3,0,1)$.}\prove{$X$ is isomorphic to the Petersen graph.}}
 \begin{proof}
  \step{s414mpetsw}{$X\cong L(\Delta_6)^H$ for some $H\sbs\Ecal_6$.}
  \begin{proof}\pf\ $t=6$ by \stepref{s414mE}, $(\Zw,\Zw)=10/5=2$ by Lemma~\ref{lem:414-apex}, and
   $n=10=|\Ecal_6|$ by Computation~\ref{comp:414-Nc}; apply Lemma~\ref{lem:414-transversal} with
   \asm[1]{0} and $f$ of \stepref{s414mtype}.~\qed\end{proof}
  \qedstep
  \begin{proof}\pf\ $L(\Delta_6)^H\cong X$ is 3-regular by \stepref{s414mpetsw}; apply
   Lemma~\ref{lem:414-E6}.~\qed\end{proof}
 \end{proof}
 \step{s414mshr}{\case{$(n,k,a,c)=(16,6,2,2)$.}\prove{$X$ is isomorphic to the Shrikhande
  graph.}}
 \begin{proof}
  \step{s414mshrsw}{$X\cong L(\Delta_7)^H$ for some $H\sbs\Ecal_7$.}
  \begin{proof}\pf\ $t=7$ by \stepref{s414mE}, $(\Zw,\Zw)=16/8=2$ by Lemma~\ref{lem:414-apex},
   and $n=16=|\Ecal_7|$ by Computation~\ref{comp:414-Nc}; apply Lemma~\ref{lem:414-transversal}
   with \asm[1]{0} and $f$ of \stepref{s414mtype}.~\qed\end{proof}
  \qedstep
  \begin{proof}\pf\ $L(\Delta_7)^H\cong X$ is 6-regular by \stepref{s414mshrsw}, so it is the
   Shrikhande graph or $L(K_{4,4})$ by Lemma~\ref{lem:414-E7}; the latter is excluded by
   \asm[4]{0}.~\qed\end{proof}
 \end{proof}
 \step{s414mchang}{\case{$(n,k,a,c)=(28,12,6,4)$.}\prove{$X$ is isomorphic to one of the Chang
  graphs.}}
 \begin{proof}
  \step{s414mchangsw}{$X\cong L(\Delta_8)^H$ for some $H\sbs\Ecal_8$.}
  \begin{proof}\pf\ $t=8$ by \stepref{s414mE}, $(\Zw,\Zw)=28/14=2$ by Lemma~\ref{lem:414-apex},
   and $n=28=|\Ecal_8|$ by Computation~\ref{comp:414-Nc}; apply Lemma~\ref{lem:414-transversal}
   with \asm[1]{0} and $f$ of \stepref{s414mtype}.~\qed\end{proof}
  \qedstep
  \begin{proof}\pf\ $L(\Delta_8)^H\cong X$ is 12-regular by \stepref{s414mchangsw}, so it is $T(8)$ or
   a Chang graph by Lemma~\ref{lem:414-E8}; $T(8)=L(K_8)$ is excluded by
   \asm[4]{0}.~\qed\end{proof}
 \end{proof}
 \qedstep
 \begin{proof}\pf\ By \stepref{s414mE} and \stepref{s414mexcl}, the cases \stepref{s414mclebsch},
  \stepref{s414mschlafli}, \stepref{s414mpet}, \stepref{s414mshr} and \stepref{s414mchang} are
  exhaustive, and each gives one of the seven graphs; it is only one of them by
  Proposition~\ref{prop:414-seven}, assertion 3. Assertion 2 is
  Proposition~\ref{prop:414-seven}, assertions 1 and 2.~\qed\end{proof}
\end{proof}

\medskip\noindent\textbf{What the proof relies on.}
Besides Sections 2 to 4, the proof uses the computations C4 (with its transitivity on ordered
stars),
C9 (with its two certificate isomorphisms), C18, C20 (at the apex), C21 and C22, and the classification SE4 of
\path{cgss/proofs/reg8/reg8.tex}. All of it is in the Lean formalization
(\path{LineGraphs/Seidel/}, \path{LineGraphs/Computations/C18.lean}, \path{LineGraphs/Computations/C22.lean} and its data files): C18 as a
checked search, C22's fourteen isomorphism certificates with the clique profiles and counts, SE4
by a case split with checked certificates, and the restated remark as
\path{SimpleGraph.seidel_seven_classification} and \path{SimpleGraph.seidel_seven}.

\endgroup

\section{Spectral characterization}\label{sec:5}
\begingroup
\emergencystretch=3em
\raggedbottom

We will follow Section 5 of [CGSS]. All graphs are finite and simple. For a graph $X$ we write
$V(X)$ and $E(X)$ for its vertices and edges, $A(X)$ for its adjacency matrix and $L(X)$ for
its line graph. The paper's graphs $G$, $G_1$, $G_2$ are our $X$, $X_1$, $X_2$, and its
auxiliary graph $H$ in Theorem 5.8 is our $Y$ (\path{cgss/notation.md}, Section 1). Valencies
are $k$, $k_1$, $k_2$ in Lemmas 5.1 to 5.6 and in Remark 5.7; from Theorem 5.8 on, $k$ is the
block size of a design and $(v,k,b,r,\lambda)$ are the design parameters, as in the paper. A
generic eigenvalue is $\mu$, since $\lambda$ is a design parameter.

The spectrum $\Spec(M)$ of a real symmetric matrix $M$ is the multiset of the roots of its
characteristic polynomial, and $\mult_M(\mu)$ is the number of copies of $\mu$ in it. In a
multiset $\ms{\mu^{m},\dots}$ the exponent $m$ counts copies, and a sum of multisets adds the
counts. The spectrum of a graph is that of its adjacency matrix, two graphs are
\emph{cospectral} when their adjacency matrices have the same characteristic polynomial, and
$\lmax$ is the largest eigenvalue. Where the paper writes that two graphs have the same
spectrum, we will say that they are cospectral. We cite the spectral toolkit \path{cgss/proofs/spectraltools/spectraltools.tex} as [ST],
by the numbers of its results; its structured version \path{cgss/proofs/spectraltools/spectraltools-lamport.tex}
has the same numbering. The toolkit is not part of this document, so we list the facts that we
take from it, with its proposed Lean names, before we use them.

The paper quotes Lemmas 5.1 to 5.4 as well known (errata, Gap 8). The toolkit proves them as
SP1 to SP4, and we will state them in the paper's form and cite the toolkit. Lemmas 5.5 and 5.6
and Remark 5.7 have proofs in the paper, and we will follow those proofs, with the toolkit's
facts at the leaves. For Lemmas 5.1, 5.2, 5.3 and 5.5 the toolkit proposes the same Lean
declarations (its Theorems 5.3, 5.4, 5.6 and 6.5), so each of these results and the
corresponding toolkit result are one declaration. Lemma 5.4, Lemma 5.6 and Remark 5.7 differ in
statement from the toolkit's Theorem 6.4, Theorem 6.7 and Corollary 6.8, in the forms that
Theorem 5.8 needs: Lemma 5.4 does not assume an edge, and asserts $k_1\ne k_2$; Lemma 5.6 (i) holds for
every biregular bipartition of $X_1$ with $|P_1|\le|Q_1|$; and Remark 5.7 assumes only that
$X_2$ has no isolated vertices and is not bipartite. We will derive them from the toolkit's
results, which we cite as inputs, and give them their own Lean names, derived from the
toolkit's by the rules of \path{cgss/notation.md}, Section 10:
\path{SimpleGraph.regular_or_isBiregularWith_of_lineGraph_regular},
\path{SimpleGraph.lineGraph_cospectral_bipartite_forall_isBiregularWith} and
\path{SimpleGraph.lineGraph_cospectral_bipartite_spectrum_of_not_isBipartite}.

\begingroup\raggedright
\medskip\noindent\textbf{Inputs from the spectral toolkit.}\label{input:st}
\repair\ These are the toolkit's proofs of the facts that the paper quotes (errata, Gap 8, and
\path{cgss/prerequisites.md}, SP1 to SP5). Matrices are real, square matrices are symmetric,
and index sets are finite.
\begin{itemize}
 \item \factlabel{[ST, Theorem 2.2]}{st:2.2}\ST{Theorem 2.2}
  (\path{Matrix.IsHermitian.eigenvalueMultiplicity_eq_finrank},
  \path{Matrix.IsHermitian.mem_spectrumMultiset_iff}). If $M$ is symmetric of order $n$, then
  $\Spec(M)$ has $n$ members counted with multiplicity, $\mult_M(\mu)=\dim\ker(M-\mu I)$, and
  $\mu\in\Spec(M)$ exactly when $M\Zx=\mu\Zx$ for some $\Zx\ne0$.
 \item \factlabel{[ST, Lemma 2.4]}{st:2.4}\ST{Lemma 2.4}
  (\path{Matrix.IsHermitian.largestEigenvalue_mem_spectrumMultiset},
  \path{Matrix.IsHermitian.le_largestEigenvalue_of_mem}). If $M$ is symmetric of positive
  order, then $\lmax(M)$ is the largest member of $\Spec(M)$.
 \item \factlabel{[ST, Proposition 2.6]}{st:2.6}\ST{Proposition 2.6}
  (\path{Matrix.IsHermitian.spectrumMultiset_eq_iff_charpoly_eq},
  \path{SimpleGraph.cospectral_iff_spectrumMultiset_eq}, \path{SimpleGraph.Iso.cospectral}).
  For symmetric $M$, $M'$, possibly on different index sets, $\chi_M=\chi_{M'}$ if and only if
  $\Spec(M)=\Spec(M')$, if and only if $\mult_M(\mu)=\mult_{M'}(\mu)$ for every real $\mu$.
  Cospectral graphs have the same number of vertices, and isomorphic graphs are cospectral.
 \item \factlabel{[ST, Theorem 3.1]}{st:3.1}\ST{Theorem 3.1}
  (\path{Matrix.IsHermitian.spectrumMultiset_aeval}). If $M$ is symmetric and $f$ is a real
  polynomial, then $\Spec(f(M))$ is the image of $\Spec(M)$ under $f$, with multiplicities. In
  particular $\Spec(cI+M)$ is $\Spec(M)$ shifted by $c$, and $\mult_{cI+M}(\mu+c)=\mult_M(\mu)$.
 \item \factlabel{[ST, Theorem 3.2]}{st:3.2}\ST{Theorem 3.2}
  (\path{Matrix.IsHermitian.rank_aeval_eq_sum_multiplicity}). If $M$ is symmetric of order $n$,
  then $\rk M=n-\mult_M(0)$, and for a real polynomial $f$, $\rk f(M)$ is the sum of
  $\mult_M(\theta)$ over the distinct eigenvalues $\theta$ of $M$ with $f(\theta)\ne0$.
 \item \factlabel{[ST, Corollary 3.4]}{st:3.4}\ST{Corollary 3.4}
  (\path{SimpleGraph.posSemidef_two_add_adjMatrix_iff}). For a graph $X$, the matrix
  $2I+A(X)$ is positive semidefinite if and only if every member of $\Spec(X)$ is at least $-2$.
 \item \factlabel{[ST, Theorem 3.5]}{st:3.5}\ST{Theorem 3.5}
  (\path{Matrix.rank_gram_eq_finrank_span}). For every matrix $B$,
  $\rk(B^TB)=\rk B=\rk(BB^T)$. The Gram matrix of a finite family of vectors in a real inner
  product space is positive semidefinite, and its rank is the dimension of the span of the family.
 \item \factlabel{[ST, Theorem 5.1]}{st:5.1}\ST{Theorem 5.1}
  (\path{SimpleGraph.adjMatrix_cube_apply_self_eq_twice_triangles}). For a graph $X$ with
  $A=A(X)$ and a vertex $u$, $A^3(u,u)$ is twice the number of triangles of $X$ that
  contain $u$.
 \item \factlabel{[ST, Theorem 5.3]}{st:5.3}\ST{Theorem 5.3} (SP1;
  \path{SimpleGraph.IsRegularOfDegree.largestEigenvalue_eq},
  \path{SimpleGraph.IsRegularOfDegree.multiplicity_eq_card_components}). If $X$ is
  $k$-regular on a nonempty vertex set, then $\lmax(X)=k$ and $\mult_X(k)$ is the number of
  connected components of $X$.
 \item \factlabel{[ST, Theorem 5.4]}{st:5.4}\ST{Theorem 5.4} (SP2;
  \path{SimpleGraph.isRegular_iff_sum_sq_eigenvalues}). A graph $X$ on $n\ge1$ vertices is
  regular if and only if $\sum_{\mu\in\Spec(X)}\mu^2=n\lmax(X)$.
 \item \factlabel{[ST, Corollary 5.5]}{st:5.5}\ST{Corollary 5.5}
  (\path{SimpleGraph.Cospectral.isRegularOfDegree},
  \path{SimpleGraph.Cospectral.connected_of_regular}). If $X$ is $k$-regular on a nonempty
  vertex set and $X'$ is cospectral with $X$, then $X'$ is $k$-regular, and $X'$ is connected
  if $X$ is.
 \item \factlabel{[ST, Theorem 5.6]}{st:5.6}\ST{Theorem 5.6} (SP3;
  \path{SimpleGraph.IsRegularOfDegree.isBipartite_iff_neg_eigenvalue}). If $X$ is connected and
  $k$-regular, then $X$ is bipartite if and only if $-k\in\Spec(X)$, and then $\mult_X(-k)=1$.
 \item \factlabel{[ST, Corollary 5.9]}{st:5.9}\ST{Corollary 5.9}
  (\path{Matrix.spectrumMultiset_smul_one_add_smul_allOnes}). For real $c$, $t$ and order
  $n>0$, $\Spec(cI+tJ)=\ms{(c+tn)^{1},c^{n-1}}$.
 \item \factlabel{[ST, Theorem 6.1]}{st:6.1}\ST{Theorem 6.1}
  (\path{Matrix.eigenvalueMultiplicity_mul_comm_of_ne_zero},
  \path{Matrix.spectrumMultiset_gram_products}). If $M$ is $P\times Q$ and $N$ is $Q\times P$,
  then $\mult_{MN}(\mu)=\mult_{NM}(\mu)$ for every $\mu\ne0$. If $B$ is $P\times Q$ of rank
  $\rho$, then $BB^T$ and $B^TB$ are positive semidefinite, have the same positive eigenvalues
  with multiplicities, and have $\mult_{BB^T}(0)=|P|-\rho$ and $\mult_{B^TB}(0)=|Q|-\rho$.
 \item \factlabel{[ST, Theorem 6.2]}{st:6.2}\ST{Theorem 6.2}
  (\path{SimpleGraph.edgeIncMatrix_mul_transpose}, \path{SimpleGraph.transpose_edgeIncMatrix_mul},
  \path{SimpleGraph.multiplicity_neg_two_lineGraph}). Let $M$ be the $V(X)\times E(X)$
  incidence matrix of a graph $X$ and $D$ its diagonal matrix of degrees. Then $MM^T=D+A(X)$,
  $M^TM=2I+A(L(X))$, and $\rk M=|V(X)|-c$, where $c$ is the number of bipartite connected
  components of $X$. Every eigenvalue of $L(X)$ is at least $-2$.
 \item \factlabel{[ST, Lemma 6.3]}{st:6.3}\ST{Lemma 6.3}
  (\path{SimpleGraph.lineGraph_connected_iff_of_no_isolated}, \path{SimpleGraph.degree_lineGraph}).
  If $X$ has at least one edge and no isolated vertices, then $X$ is connected if and only if
  $L(X)$ is connected, and the edge $\{u,w\}$ has degree $\deg(u)+\deg(w)-2$ in $L(X)$.
 \item \factlabel{[ST, Theorem 6.4]}{st:6.4}\ST{Theorem 6.4} (SP4;
  \path{SimpleGraph.regular_or_semiregular_of_lineGraph_regular}). If $X$ has at least one edge
  and no isolated vertices, and $L(X)$ is connected and $h$-regular, then $X$ is regular, or $X$
  is bipartite with constant degree on each part.
 \item \factlabel{[ST, Theorem 6.6]}{st:6.6}\ST{Theorem 6.6}
  (\path{Matrix.spectrumMultiset_fromBlocks_scalar}). Let $B$ be $P\times Q$ of rank $\rho$,
  let $\mu_1,\dots,\mu_\rho$ be the positive eigenvalues of $BB^T$ with repetition, and
  let $s,t\in\re$. Then the spectrum of $C=\pmat{sI&B\\B^T&tI}$ consists of
  $\gamma_j^{\pm}=\frac12\bigl(s+t\pm\sqrt{(s-t)^2+4\mu_j}\bigr)$ for $1\le j\le\rho$,
  of $s$ with multiplicity $|P|-\rho$ and of $t$ with multiplicity $|Q|-\rho$.
 \item \factlabel{[ST, Theorem 6.7]}{st:6.7}\ST{Theorem 6.7}
  (\path{SimpleGraph.lineGraph_cospectral_bipartite}). Let $X_1$ be a connected bipartite graph
  with an edge such that $L(X_1)$ is regular, and let $X_2$ be a graph without isolated vertices
  such that $L(X_2)$ is cospectral with $L(X_1)$. Then $X_2$ is connected, and precisely one of
  the following holds: $X_2$ is bipartite, and with the part sizes ordered $p\le q$ in each graph,
  the part sizes, the respective degrees and the spectra of $B_1B_1^T$ and $B_2B_2^T$ agree; or
  $X_1$ is not regular, $X_2$ is regular and not bipartite, and $|V(X_1)|=|V(X_2)|+1$.
 \item \factlabel{[ST, Corollary 6.8]}{st:6.8}\ST{Corollary 6.8}
  (\path{SimpleGraph.lineGraph_cospectral_bipartite_spectrum}). Let $X_1$ be a connected graph
  with a bipartition $(P,Q)$, biregular with degrees $(k_1,k_2)$, with $k_1>k_2>0$,
  $p=|P|\le q=|Q|$ and biadjacency matrix $B$ of rank $\rho$, and let $X_2$ be a connected
  regular graph that is not bipartite with $L(X_2)$ cospectral with $L(X_1)$. Then $k_1k_2$ is a
  positive eigenvalue of $BB^T$ of multiplicity 1, and if $\mu_2,\dots,\mu_\rho$ are the other
  positive eigenvalues of $BB^T$ with repetition, then $\Spec(X_2)$ consists of
  $\frac12(k_1+k_2)$, of $\pm\sqrt{\mu_j+\frac14(k_1-k_2)^2}$ for $2\le j\le\rho$, of
  $\frac12(k_1-k_2)$ with multiplicity $p-\rho$ and of $-\frac12(k_1-k_2)$ with multiplicity
  $q-\rho$.
 \item \factlabel{[ST, Theorem 7.7]}{st:7.7}\ST{Theorem 7.7}
  (\path{Matrix.gram_eq_of_two_eigenvalues}). Let $N$ be a $v\times b$ zero-one matrix with
  $v>1$, constant row sum $r>0$ and constant column sum $k>1$, and let $\lambda>0$ with
  $r(k-1)=\lambda(v-1)$. If $\Spec(NN^T)=\ms{(rk)^{1},(r-\lambda)^{v-1}}$, then
  $NN^T=(r-\lambda)I+\lambda J$. If moreover $r>\lambda$, then $b\ge v$ (Fisher's inequality,
  DS3).
 \item \factlabel{[ST, Theorem 7.10]}{st:7.10}\ST{Theorem 7.10}
  (\path{SimpleGraph.IsRegularOfDegree.aeval_eq_smul_allOnes_of_connected}). If $X$ is connected
  and $k$-regular on $n>0$ vertices, and the real polynomial $f$ vanishes at every eigenvalue
  of $X$ other than $k$, then $f(A(X))=\frac{f(k)}{n}J$.
\end{itemize}
\endgroup

\begingroup\raggedright
\medskip\noindent\textbf{Inputs from Mathlib and from the Hall--Connor section.}\label{facts:sec5}
\uses{hc:facts:finite}
We will use the following declarations of Mathlib at commit \path{082e2d37}
(4 September 2026), the version pinned in \path{lake-manifest.json}: \path{SimpleGraph.Iso.lineGraph} (an isomorphism of graphs induces one of
their line graphs), \path{SimpleGraph.lineGraph_adj_iff_exists} (two edges are adjacent in the
line graph exactly when they are distinct and share an endpoint),
\path{SimpleGraph.sum_degrees_eq_twice_card_edges}, \path{SimpleGraph.degree_compl},
\path{SimpleGraph.adjMatrix_mul_self_apply_self} (the diagonal of $A(X)^2$ is the degree),
\path{SimpleGraph.cycleGraph_adj'}, \path{SimpleGraph.Iso.completeGraph},
\path{SimpleGraph.completeBipartiteGraphCongr}, \path{Equiv.sumComm}, \path{Fintype.equivFin},
\path{Finset.card_le_univ}, \path{Finset.eq_of_subset_of_card_le},
\path{LinearEquiv.finrank_eq}, \path{Submodule.finrank_le}, \path{finrank_euclideanSpace_fin}
and the structure \path{Matrix.IsHadamard}. Double counting is the paragraph
\hyperref[hc:facts:finite]{(finite counting and order)} of the Hall--Connor section. The
design interface is Definition~\ref{hc:def:design} of that section, with its incidence matrix
$N$ of Definition~\ref{hc:def:matrices}. This is the one design interface of Section 5; the
toolkit's facts concern real matrices, and we apply them to the entrywise real cast of the
integer matrix $N$, which we also write $N$. The cast is taken in Lemma~\ref{lem:design},
assertion 3, and in Theorem~\ref{thm:5.8}, where \ST{Theorem 7.7} is applied to a real zero-one
matrix. The toolkit's Section 7 proposes the same names \path{Design} and
\path{Design.incMatrix} with a real incidence matrix; the shared declarations are those of the
Hall--Connor section, over $\ints$, and we use no toolkit result stated for the toolkit's
designs.\par
\endgroup

\begingroup\raggedright
\medskip\noindent\textbf{Inputs from Section 4.}\label{input:sec4}
Theorem 5.8 rests on Theorem 4.11. We recall what we will use from Section 4, where these
results are proved: Definition~\ref{def:CP}, Definition~\ref{def:represented} and
Theorem~\ref{thm:4.11}. Theorem~\ref{thm:5.8} and Lemma~\ref{lem:dense} depend on those labels
directly, and none of them depends on Section 5.
\begin{itemize}
 \item The \emph{cocktail party graph} $CP(m)$ (\path{SimpleGraph.cocktailParty}) has vertex set
  $\{1,\dots,m\}\times\{0,1\}$, and $(i,a)$ is adjacent to $(j,a')$ exactly when $i\ne j$. It is
  $(2m-2)$-regular on $2m$ vertices.
 \item A graph $X$ is \emph{represented in $\typ E_8$} (\path{SimpleGraph.IsRepresentedIn}) if
  there is a family $(\Zx_u)_{u\in V(X)}$ of vectors of $\Phi(\typ E_8)\sbs\re^8$ (model (ii) of
  Definition~\ref{def:3.2}) whose Gram matrix is $2I+A(X)$.
 \item Theorem 4.11 (Hoffman and Ray-Chaudhuri; \path{SimpleGraph.hoffman_rayChaudhuri}). If
  $X$ is connected and regular and $2I+A(X)$ is positive semidefinite, then $X$ is isomorphic to
  $L(Y)$ for some graph $Y$, or $X$ is isomorphic to $CP(m)$ for some $m$, or $X$ is represented
  in $\typ E_8$.
\end{itemize}
\endgroup

\begingroup\raggedright
\medskip\noindent\textbf{Inputs for Example 5.13.}\label{input:ex513}
The example needs one fact that this section does not prove, EX4 of
\path{cgss/prerequisites.md}, and the existence results EX1 to EX3 for its instances. We record
them as deferred inputs; none of them may depend on Section 5.
\begin{itemize}
 \item (EX4, deferred.) If $Y$ is a connected graph on more than three vertices that is not
  bipartite, and $Z$ is a bipartite graph without isolated vertices, then $L(Y)\not\cong L(Z)$.
  This is a case of Whitney's theorem on line graphs, which Mathlib does not have. The
  prerequisites prove it from four facts, which belong to Sections 3 and 4: the star closure of
  the standard representation of $L(Y)$ is $\typ D_n$ when $Y$ is connected, not bipartite and on
  $n$ vertices; that of $L(Z)$ is $\typ A_{m-1}$ when $Z$ is connected, bipartite and on $m$
  vertices; isomorphic graphs have isometric representations, so the star closures are
  isomorphic (LN3 and LA2 there), and $\typ D_n\not\cong\typ A_n$ for $n\ge4$; and $Z$ is
  connected, since $L(Z)\cong L(Y)$ is connected and $Z$ has no isolated vertices. \new\ The
  Lean proof of Example 5.13 avoids EX4, by a clique bound recorded at the step
  on non-isomorphism in the proof of Example~\ref{ex:5.13}; we keep EX4 for the route of the paper.
 \item (EX1 to EX3; computation C23.) Resolvable designs with parameters
  $(k(2k-1),k,4k^2-1,2k+1,1)$ exist for $k=2^d$ with $d\ge1$ and for $k=3$, and designs with
  parameters $(6,3,10,5,2)$ and $(28,21,36,27,20)$ exist. We will use these only in the remark after
  Example 5.13. Computation C23 gives in Lean the designs that the paper's instances $k=2$ and $k=4$
  need; their existence for every $k=2^d$ and for $k=3$ is a citation, outside the scope.
\end{itemize}
\endgroup

\subsection{Lemmas 5.1 to 5.7}

\medskip\noindent\textbf{Definitions and auxiliary lemmas.}

\begin{namedresult}{Definition}{(bipartitions and biadjacency matrices)}\label{def:bipartition}
 \lean{SimpleGraph.IsBiregularWith}\lean{SimpleGraph.biadjMatrix}\uses{}
 \assume{$X$ is a graph.}
 A \emph{bipartition} of $X$ is a pair $(P,Q)$ of disjoint sets with $P\cup Q=V(X)$ such that
 every edge has one end in $P$ and one in $Q$; in Mathlib, \path{X.IsBipartiteWith P Q}
 (\path{SimpleGraph.IsBipartiteWith}) with $Q=V(X)\diff P$. $X$ is bipartite if it has a bipartition. The bipartition is \emph{biregular
 with degrees} $(k_1,k_2)$ if every vertex of $P$ has degree $k_1$ and every vertex of $Q$ has
 degree $k_2$. The \emph{biadjacency matrix} of $X$ for $(P,Q)$ is the $P\times Q$ matrix $B$
 with $B(u,w)=A(X)(u,w)$.
\end{namedresult}

\begin{namedresult}{Definition}{(incidence graphs, flag graphs, characterization)}\label{def:flag}
 \lean{Design.incidenceGraph}\lean{Design.flagGraph}\lean{Design.FlagGraphCharacterized}
 \uses{hc:def:design}
 \leavevmode\par
 \assume{$D$ is a design with parameters $(v,k,b,r,\lambda)$.}
 Here $D$ is as in Definition~\ref{hc:def:design}, with points $P$ and block labels $\cB$. The
 \emph{incidence graph} $I(D)$ has vertex set $P\sqcup\cB$, and a point $p$ is adjacent to a
 block $B$ when $p\in B$; there are no other edges. The \emph{flag graph} is $F(D)=L(I(D))$. We
 say that $F(D)$ is \emph{characterized by its spectrum} if every graph cospectral with $F(D)$
 is isomorphic to $F(D_2)$ for some design $D_2$ with the parameters $(v,k,b,r,\lambda)$, on
 some finite types $P_2$ of points and $\cB_2$ of block labels with decidable equality. In Lean,
 $P_2$ and $\cB_2$ range over \texttt{Type}, as do the vertex types of all graphs of this section;
 the design $D_2$ of Theorem 5.8 lives on subtypes of the vertex type of a graph.
 \repair\ This is the sense of the paper's proofs and of its sources (errata, Implicit
 hypothesis 6); in the stronger sense ``determined up to isomorphism'' Corollaries 5.10 and 5.11
 and Remark 5.12 are false.
\end{namedresult}

\begin{namedresult}{Lemma}{(design identities)}\label{lem:design}
 \lean{Design.b_mul_k_eq_v_mul_r}\lean{Design.r_mul_k_sub_one_eq}\lean{Design.k_le_v}
 \lean{Design.one_le_lam}\lean{Design.lam_le_r}\lean{Design.lam_eq_r_iff}
 \lean{Design.incMatrix_mul_transpose}\lean{Design.spectrumMultiset_incMatrix_mul_transpose}
 \lean{Design.v_le_b_of_lam_lt_r}
 \uses{hc:def:design,hc:def:matrices,hc:facts:finite,facts:sec5,st:5.9,st:7.7}
 \assume{$D$ is a design with parameters $(v,k,b,r,\lambda)$ on $P$ and $\cB$, with incidence
  matrix $N$.}
 \prove{\begin{pfenum}
  \item $bk=vr$ and $r(k-1)=\lambda(v-1)$.
  \item $k\le v$, $\lambda\ge1$ and $r\ge\lambda$; moreover $r=\lambda$ if and only if $k=v$,
   if and only if every block contains every point.
  \item $NN^T=(r-\lambda)I+\lambda J$ and $\Spec(NN^T)=\ms{(rk)^{1},(r-\lambda)^{v-1}}$.
  \item If $r>\lambda$, then $b\ge v$.
 \end{pfenum}}
\end{namedresult}
\begin{proof}
 \pfsketch\ These are DS1 and DS3 of \path{cgss/prerequisites.md}: double counting, and the
 Fisher inequality of the toolkit's Theorem 7.7.
 \step{s5dsflags}{$bk=vr$.}
 \begin{proof}\pf\ Count the pairs $(p,B)$ with $p\in B$ by blocks and by points, with the
  double counting of \hyperref[hc:facts:finite]{(finite counting and order)}.~\qed\end{proof}
 \step{s5dspairs}{$r(k-1)=\lambda(v-1)$.}
 \begin{proof}\pf\ Fix a point $p$ and count the pairs $(q,B)$ with $q\ne p$ and $p,q\in B$, by
  blocks through $p$ and by points $q$, with the same double counting.~\qed\end{proof}
 \step{s5dskv}{$k\le v$.}
 \begin{proof}\pf\ A block is a subset of $P$ of size $k$; apply \path{Finset.card_le_univ}.~\qed\end{proof}
 \step{s5dslam}{$\lambda\ge1$.}
 \begin{proof}
  \step{s5dslamzero}{\assume{$\lambda=0$.}\prove{\textsc{False}.}}
  \begin{proof}\pf\ By \stepref{s5dspairs}, $r(k-1)=0$, and $k>1$ gives $r=0$. Then
   \stepref{s5dsflags} gives $bk=0$, against $b>0$ and $k>1$.~\qed\end{proof}
  \qedstep
  \begin{proof}\pf\ $\lambda$ is a natural number and \stepref{s5dslamzero}.~\qed\end{proof}
 \end{proof}
 \step{s5dsrlam}{$r\ge\lambda$, with equality if and only if $k=v$.}
 \begin{proof}\pf\ By \stepref{s5dspairs} and \stepref{s5dskv},
  $\lambda(v-1)=r(k-1)\le r(v-1)$, with equality exactly when $k=v$, since $r\ge1$ by
  \stepref{s5dspairs} and \stepref{s5dslam}; divide by $v-1>0$.~\qed\end{proof}
 \step{s5dsfull}{$k=v$ if and only if every block contains every point.}
 \begin{proof}\pf\ A block of size $k$ is a subset of $P$; apply
  \path{Finset.eq_of_subset_of_card_le} to it and $P$.~\qed\end{proof}
 \step{s5dsgram}{$NN^T=(r-\lambda)I+\lambda J$.}
 \begin{proof}\pf\ $(NN^T)(p,q)$ is the number of blocks containing $p$ and $q$, which is $r$
  for $p=q$ and $\lambda$ otherwise, by Definition~\ref{hc:def:design}.~\qed\end{proof}
 \step{s5dsspec}{$\Spec(NN^T)=\ms{(rk)^{1},(r-\lambda)^{v-1}}$.}
 \begin{proof}\pf\ \ST{Corollary 5.9} with $c=r-\lambda$, $t=\lambda$ and $n=v$ gives the
  eigenvalue $r-\lambda+\lambda v$ once, and $r-\lambda+\lambda v=rk$ by
  \stepref{s5dspairs}.~\qed\end{proof}
 \step{s5dsfisher}{\assume{$r>\lambda$.}\prove{$b\ge v$.}}
 \begin{proof}\pf\ \ST{Theorem 7.7} for the real zero-one $v\times b$ matrix $N$, whose row
  sums are $r>\lambda\ge1$ and column sums $k>1$ (Definition~\ref{hc:def:matrices}), with $v>1$,
  $\lambda>0$ by \stepref{s5dslam}, \stepref{s5dspairs}, \stepref{s5dsspec} and \asm{1}.~\qed\end{proof}
 \qedstep
 \begin{proof}\pf\ The preceding steps are the four assertions.~\qed\end{proof}
\end{proof}

\begin{namedresult}{Lemma}{(incidence and flag graphs)}\label{lem:incidence}
 \lean{Design.incidenceGraph_isBiregularWith}\lean{Design.incidenceGraph_connected}
 \lean{Design.incidenceGraph_isRegular_iff}\lean{Design.flagGraph_connected}
 \lean{Design.flagGraph_isRegularOfDegree}\lean{Design.card_flagGraph}
 \lean{Design.posSemidef_two_add_adjMatrix_flagGraph}\lean{Design.rank_two_add_adjMatrix_flagGraph}
 \uses{def:flag,def:bipartition,hc:def:matrices,lem:design,facts:sec5,st:3.5,st:6.2,st:6.3}
 \assume{$D$ is a design with parameters $(v,k,b,r,\lambda)$ on $P$ and $\cB$, with incidence
  matrix $N$.}
 \prove{\begin{pfenum}
  \item $(P,\cB)$ is a bipartition of $I(D)$, biregular with degrees $(r,k)$, whose
   biadjacency matrix is $N$; $I(D)$ has $vr$ edges and no isolated vertices.
  \item $I(D)$ is connected, and it is regular if and only if $r=k$.
  \item $F(D)$ is connected and $(r+k-2)$-regular on $vr$ vertices.
  \item $2I+A(F(D))$ is positive semidefinite of rank $v+b-1$.
 \end{pfenum}}
\end{namedresult}
\begin{proof}
 \pfsketch\ We will read the degrees off the definitions, and then apply the toolkit's facts on
 line graphs and incidence matrices.
 \step{s5igbip}{Assertion 1 holds.}
 \begin{proof}
  \step{s5igparts}{$(P,\cB)$ is a bipartition with biadjacency matrix $N$.}
  \begin{proof}\pf\ Every edge of $I(D)$ joins a point to a block, by Definition~\ref{def:flag},
   and $N(p,B)=1$ exactly when $p\in B$ (Definition~\ref{hc:def:matrices}).~\qed\end{proof}
  \step{s5igdeg}{Points have degree $r\ge1$ and blocks have degree $k\ge2$.}
  \begin{proof}\pf\ The neighbours of $p$ are the blocks through $p$, and those of $B$ are its
   $k$ points; $r\ge\lambda\ge1$ by Lemma~\ref{lem:design}, and $k>1$.~\qed\end{proof}
  \qedstep
  \begin{proof}\pf\ The number of edges is the number of pairs $(p,B)$ with $p\in B$, which is
   $vr$ by \stepref{s5igdeg} and double counting; there are no isolated vertices by
   \stepref{s5igdeg}.~\qed\end{proof}
 \end{proof}
 \step{s5igconn}{$I(D)$ is connected.}
 \begin{proof}
  \step{s5igpoints}{Any two points are joined by a path.}
  \begin{proof}\pf\ Two distinct points lie in $\lambda\ge1$ common blocks, by
   Lemma~\ref{lem:design}; such a block is a common neighbour.~\qed\end{proof}
  \qedstep
  \begin{proof}\pf\ $P$ is nonempty since $v>1$, and every block has a neighbouring point, since
   $k\ge2$; combine this with \stepref{s5igpoints}.~\qed\end{proof}
 \end{proof}
 \step{s5igreg}{$I(D)$ is regular if and only if $r=k$.}
 \begin{proof}\pf\ Both parts are nonempty ($v>1$, $b>0$), and the degrees are $r$ and $k$ by
  \stepref{s5igbip}.~\qed\end{proof}
 \step{s5igflag}{Assertion 3 holds.}
 \begin{proof}\pf\ Apply \ST{Lemma 6.3} to $I(D)$, which has an edge and no isolated vertices
  by \stepref{s5igbip} and is connected by \stepref{s5igconn}: $F(D)$ is connected, and an edge
  $\{p,B\}$ has degree $r+k-2$ in it. $F(D)$ has $|E(I(D))|=vr$ vertices.~\qed\end{proof}
 \step{s5igrank}{Assertion 4 holds.}
 \begin{proof}\pf\ Let $M$ be the incidence matrix of $I(D)$. By \ST{Theorem 6.2},
  $M^TM=2I+A(F(D))$ and $\rk M=(v+b)-1$, since $I(D)$ is one bipartite component by
  \stepref{s5igbip} and \stepref{s5igconn}. $M^TM$ is a Gram matrix, positive semidefinite
  with $\rk(M^TM)=\rk M$ by \ST{Theorem 3.5}.~\qed\end{proof}
 \qedstep
 \begin{proof}\pf\ Steps \stepref{s5igbip} to \stepref{s5igrank}.~\qed\end{proof}
\end{proof}

\begin{namedresult}{Lemma}{(regular graphs of large valency)}\label{lem:dense}
 \lean{SimpleGraph.eq_top_of_isRegularOfDegree_card_sub_one}
 \lean{SimpleGraph.nonempty_iso_cocktailParty}
 \uses{def:CP,facts:sec5}
 \assume{$X$ is a graph on $n$ vertices.}
 \prove{\begin{pfenum}
  \item If $n\ge1$ and $X$ is $(n-1)$-regular, then $X$ is complete, and $X\cong K_n$.
  \item If $n=2m$ with $m\ge1$ and $X$ is $(2m-2)$-regular, then $X\cong CP(m)$.
 \end{pfenum}}
\end{namedresult}
\begin{proof}
 \pfsketch\ \repair\ The paper does not state these, but its corollaries leave them to the
 reader (errata, Gaps 15 and 16). We will look at the non-neighbours of each vertex.
 \step{s5dncompl}{Every vertex has $n-1-\deg(u)$ non-neighbours other than itself.}
 \begin{proof}\pf\ \path{SimpleGraph.degree_compl}.~\qed\end{proof}
 \step{s5dncomplete}{Assertion 1 holds.}
 \begin{proof}\pf\ By \stepref{s5dncompl}, each vertex is adjacent to every other vertex, so
  $X$ is the complete graph on $V(X)$; apply \path{SimpleGraph.Iso.completeGraph} to
  \path{Fintype.equivFin}.~\qed\end{proof}
 \step{s5dncp}{Assertion 2 holds.}
 \begin{proof}
  \step{s5dnpartner}{Each vertex $u$ has a unique non-neighbour $u^*\ne u$, and $u^{**}=u$.}
  \begin{proof}\pf\ \stepref{s5dncompl} gives exactly one, and non-adjacency is
   symmetric.~\qed\end{proof}
  \step{s5dnenum}{\pick\ a bijection $e:V(X)\to\{1,\dots,m\}\times\{0,1\}$ such that $e(u)$ and
   $e(u^*)$ have the same first coordinate for every $u$.}
  \begin{proof}\pf\ The pairs $\{u,u^*\}$ partition $V(X)$ into $m$ pairs, by
   \stepref{s5dnpartner}; number the pairs with \path{Fintype.equivFin} and the two members of
   each pair by $0$ and $1$.~\qed\end{proof}
  \qedstep
  \begin{proof}\pf\ For $u\ne w$, $u$ and $w$ are adjacent if and only if $w\ne u^*$, by
   \stepref{s5dnpartner}, that is, when $e(u)$ and $e(w)$ have different first coordinates by
   \stepref{s5dnenum}; this is adjacency in $CP(m)$.~\qed\end{proof}
 \end{proof}
 \qedstep
 \begin{proof}\pf\ Steps \stepref{s5dncomplete} and \stepref{s5dncp}.~\qed\end{proof}
\end{proof}

\begin{namedresult}{Lemma}{(connected $2$-regular graphs)}\label{lem:cycle}
 \lean{SimpleGraph.nonempty_iso_cycleGraph}\uses{facts:sec5}
 \assume{$X$ is a connected $2$-regular graph on $n$ vertices.}
 \prove{$n\ge3$ and $X\cong C_n$, the cycle \path{SimpleGraph.cycleGraph n}.}
\end{namedresult}
\begin{proof}
 \pfsketch\ \repair\ The paper leaves this to the reader in Corollaries 5.10 and 5.11 (errata,
 Gap 16). We will walk around $X$ without turning back.
 \step{s5cyseq}{\pick\ adjacent vertices $u_0,u_1$, and for $i\ge1$ let $u_{i+1}$ be the
  neighbour of $u_i$ other than $u_{i-1}$.}
 \begin{proof}\pf\ $X$ is nonempty since it is connected, and a vertex has two neighbours; each
  $u_i$ with $i\ge1$ is adjacent to $u_{i-1}$, so it has exactly one other
  neighbour.~\qed\end{proof}
 \step{s5cyfirst}{\pick\ the least $j$ such that $u_j=u_i$ for some $i<j$. Then $i=0$ and
  $j\ge3$.}
 \begin{proof}
  \step{s5cyexists}{Such a $j$ exists, with $j\le n$.}
  \begin{proof}\pf\ The vertices $u_0,\dots,u_n$ are $n+1$ vertices of a set of size $n$.~\qed\end{proof}
  \step{s5cynotlater}{\assume{$u_j=u_i$ with $1\le i<j$.}\prove{\textsc{False}.}}
  \begin{proof}\pf\ By \stepref{s5cyseq}, $u_{i+1}\ne u_i$ and $u_{i+2}\ne u_i$, so $j\ge i+3$;
   by the minimality of $j$, $u_{i-1}$, $u_{i+1}$ and $u_{j-1}$ are then three distinct
   neighbours of $u_i=u_j$, against $2$-regularity.~\qed\end{proof}
  \qedstep
  \begin{proof}\pf\ Steps \stepref{s5cyexists} and \stepref{s5cynotlater} give $i=0$; and
   $u_1\ne u_0$, $u_2\ne u_0$ by \stepref{s5cyseq}, so $j\ge3$.~\qed\end{proof}
 \end{proof}
 \step{s5cyclosed}{The set $W=\{u_0,\dots,u_{j-1}\}$ contains all neighbours of its members.}
 \begin{proof}\pf\ For $1\le i\le j-1$ the two neighbours of $u_i$ are $u_{i-1}$ and $u_{i+1}$,
  with $u_j=u_0$; the two neighbours of $u_0$ are $u_1$ and $u_{j-1}$, distinct since $j\ge3$,
  by \stepref{s5cyfirst}.~\qed\end{proof}
 \qedstep
 \begin{proof}\pf\ By \stepref{s5cyclosed} and connectedness, $W=V(X)$, so $j=n\ge3$. By the
  same step, $u_i\mapsto i$ is a bijection $V(X)\to\{0,\dots,n-1\}$ under which $u_i$ and $u_{i'}$
  are adjacent exactly when $i-i'\equiv\pm1\pmod n$; this is adjacency in
  $C_n$ by \path{SimpleGraph.cycleGraph_adj'}.~\qed\end{proof}
\end{proof}

\begin{namedresult}{Lemma}{(isolated vertices)}\label{lem:isolated}
 \lean{SimpleGraph.exists_forall_degree_pos_nonempty_iso_lineGraph}\uses{facts:sec5}
 \assume{$Y$ is a graph.}
 \prove{There is a graph $Y'$ without isolated vertices such that $L(Y')\cong L(Y)$.}
\end{namedresult}
\begin{proof}
 \pfsketch\ \repair\ Isolated vertices are invisible to the line graph (errata, Implicit
 hypotheses 3 and 5). We will delete them.
 \pflet{$W$ be the set of vertices of $Y$ of positive degree, and $Y'$ the subgraph of $Y$
  induced on $W$.}
 \step{s5isnone}{$Y'$ has no isolated vertices.}
 \begin{proof}\pf\ A neighbour of $w\in W$ has positive degree, so it lies in $W$.~\qed\end{proof}
 \step{s5isedges}{Every edge of $Y$ has both ends in $W$, so $\{u,w\}\mapsto\{u,w\}$ is a
  bijection $E(Y')\to E(Y)$.}
 \begin{proof}\pf\ The ends of an edge have positive degree; $Y'$ is induced.~\qed\end{proof}
 \qedstep
 \begin{proof}\pf\ The bijection of \stepref{s5isedges} preserves the ends of edges, so it
  preserves adjacency in the line graphs by \path{SimpleGraph.lineGraph_adj_iff_exists}; with
  \stepref{s5isnone} this is the claim.~\qed\end{proof}
\end{proof}

\medskip\noindent\textbf{Lemmas 5.1 to 5.5.}

\begin{cgss}{Lemma}{5.1}\label{lem:5.1}
 \lean{SimpleGraph.IsRegularOfDegree.largestEigenvalue_eq}
 \lean{SimpleGraph.IsRegularOfDegree.multiplicity_eq_card_components}\uses{st:5.3}
 \assume{$X$ is a $k$-regular graph on a nonempty vertex set.}
 \prove{$\lmax(X)=k$, and $\mult_X(k)$ is the number of connected components of $X$.}
\end{cgss}
\begin{proof}
 \pfsketch\ \repair\ The paper quotes this as well known (errata, Gap 8); the toolkit proves it
 through the Laplacian $kI-A(X)$.
 \qedstep
 \begin{proof}\pf\ \ST{Theorem 5.3}.~\qed\end{proof}
\end{proof}

\begin{cgss}{Lemma}{5.2}\label{lem:5.2}
 \lean{SimpleGraph.isRegular_iff_sum_sq_eigenvalues}\uses{st:5.4}
 \assume{$X$ is a graph on $n\ge1$ vertices.}
 \prove{$X$ is regular if and only if $\sum_{\mu\in\Spec(X)}\mu^2=n\lmax(X)$, where the sum
  counts multiplicities.}
\end{cgss}
\begin{proof}
 \pfsketch\ \repair\ The paper quotes this as well known (errata, Gap 8); the toolkit proves it
 from the Rayleigh quotient of $\Zj$.
 \qedstep
 \begin{proof}\pf\ \ST{Theorem 5.4}.~\qed\end{proof}
\end{proof}

\begin{cgss}{Lemma}{5.3}\label{lem:5.3}
 \lean{SimpleGraph.IsRegularOfDegree.isBipartite_iff_neg_eigenvalue}\uses{st:5.6}
 \assume{$X$ is a connected $k$-regular graph.}
 \prove{$-k\in\Spec(X)$ if and only if $X$ is bipartite, and then $\mult_X(-k)=1$.}
\end{cgss}
\begin{proof}
 \pfsketch\ \repair\ The paper quotes this as well known (errata, Gap 8); the toolkit proves it
 through the signless Laplacian $kI+A(X)$.
 \qedstep
 \begin{proof}\pf\ \ST{Theorem 5.6}.~\qed\end{proof}
\end{proof}

\begin{cgss}{Lemma}{5.4}\label{lem:5.4}
 \lean{SimpleGraph.regular_or_isBiregularWith_of_lineGraph_regular}\uses{def:bipartition,st:6.4}
 \assume{\begin{pfenum}
  \item $X$ is a graph without isolated vertices.
  \item $L(X)$ is connected and regular.
 \end{pfenum}}
 \prove{$X$ is regular, or $X$ has a bipartition that is biregular with degrees $(k_1,k_2)$
  for some $k_1\ne k_2$.}
\end{cgss}
\begin{proof}
 \pfsketch\ \repair\ The paper quotes this as well known (errata, Gap 8), and it needs
 \asm[1]{0} (errata, Implicit hypothesis 3: $K_2$ with an isolated vertex is a counterexample).
 The toolkit proves it by counting, as its Theorem 6.4, which also assumes an edge and asserts
 $k_1+k_2=h+2$ and $pk_1=qk_2=|E(X)|$ for the part sizes $p$, $q$. We will derive the edge and
 $k_1\ne k_2$; the form without an edge is the one that Theorem 5.8 uses, so this lemma has its
 own Lean name.
 \step{s554edge}{$X$ has an edge.}
 \begin{proof}\pf\ A connected graph is nonempty, and the vertices of $L(X)$ are the edges of
  $X$.~\qed\end{proof}
 \step{s554reg}{\case{$X$ is regular.}}
 \begin{proof}\pf\ This is the first alternative.~\qed\end{proof}
 \step{s554notreg}{\case{$X$ is not regular.}}
 \begin{proof}\pf\ By \ST{Theorem 6.4}, with \stepref{s554edge}, \asm[1]{0} and \asm[2]{0}, $X$
  has a bipartition that is biregular with some degrees $(k_1,k_2)$, and $k_1\ne k_2$ since $X$
  is not regular.~\qed\end{proof}
 \qedstep
 \begin{proof}\pf\ The cases \stepref{s554reg} and \stepref{s554notreg} are
  exhaustive.~\qed\end{proof}
\end{proof}

\begin{cgss}{Lemma}{5.5}\label{lem:5.5}
 \lean{SimpleGraph.cospectral_of_lineGraph_cospectral}
 \uses{lem:5.1,facts:sec5,st:2.2,st:2.4,st:2.6,st:3.1,st:6.1,st:6.2,st:6.3}
 \assume{\begin{pfenum}
  \item $X$ and $X'$ are regular graphs, each with at least one edge.
  \item $L(X)$ and $L(X')$ are cospectral.
 \end{pfenum}}
 \prove{$X$ and $X'$ are cospectral.}
\end{cgss}
\begin{proof}
 \pfsketch\ \repair\ The paper states that the spectrum of a regular graph is determined by that
 of its line graph; this needs an edge, since all edgeless graphs have the same line graph
 (errata, Implicit hypothesis 3). We will follow the paper: the largest eigenvalue of the line
 graph gives the valency, its order gives the number of vertices, and $MM^T$ and $M^TM$ have the
 same nonzero eigenvalues.
 \pflet{$k$, $k'$ be the valencies, $n$, $n'$ the numbers of vertices and $m$, $m'$ the numbers
  of edges of $X$, $X'$, and $M$, $M'$ their incidence matrices.}
 \step{s555m}{$m=m'$.}
 \begin{proof}\pf\ The cospectral graphs $L(X)$, $L(X')$ have the same number of vertices, by
  \ST{Proposition 2.6}.~\qed\end{proof}
 \step{s555lk}{$L(X)$ is $(2k-2)$-regular and $L(X')$ is $(2k'-2)$-regular, with $k,k'\ge1$.}
 \begin{proof}\pf\ An edge gives $k,k'\ge1$, so there are no isolated vertices; apply the
  degree formula of \ST{Lemma 6.3}.~\qed\end{proof}
 \step{s555k}{$k=k'$.}
 \begin{proof}\pf\ By Lemma~\ref{lem:5.1} and \stepref{s555lk}, $\lmax(L(X))=2k-2$ and
  $\lmax(L(X'))=2k'-2$; cospectral graphs have the same largest eigenvalue by
  \ST{Proposition 2.6} and \ST{Lemma 2.4}.~\qed\end{proof}
 \step{s555n}{$n=n'$.}
 \begin{proof}\pf\ $nk=2m$ and $n'k'=2m'$ by \path{SimpleGraph.sum_degrees_eq_twice_card_edges};
  use \stepref{s555m}, \stepref{s555k} and $k\ge1$.~\qed\end{proof}
 \step{s555gram}{$MM^T=kI+A(X)$, $M^TM=2I+A(L(X))$, $M'M'^T=k'I+A(X')$ and
  $M'^TM'=2I+A(L(X'))$.}
 \begin{proof}\pf\ \ST{Theorem 6.2} for $X$ and for $X'$, with the degree matrices $kI$ and
  $k'I$.~\qed\end{proof}
 \step{s555nonzero}{$\mult_{kI+A(X)}(\mu)=\mult_{kI+A(X')}(\mu)$ for every $\mu\ne0$.}
 \begin{proof}
  \step{s555nzX}{$\mult_{kI+A(X)}(\mu)=\mult_{2I+A(L(X))}(\mu)$ and
   $\mult_{k'I+A(X')}(\mu)=\mult_{2I+A(L(X'))}(\mu)$.}
  \begin{proof}\pf\ \ST{Theorem 6.1} for $M$, $M^T$ and for $M'$, $M'^T$, with
   \stepref{s555gram}, since $\mu\ne0$.~\qed\end{proof}
  \step{s555nzL}{$\mult_{2I+A(L(X))}(\mu)=\mult_{2I+A(L(X'))}(\mu)$.}
  \begin{proof}\pf\ Both equal the multiplicity of $\mu-2$ in the common spectrum of $L(X)$
   and $L(X')$, by \ST{Theorem 3.1} and \ST{Proposition 2.6}.~\qed\end{proof}
  \qedstep
  \begin{proof}\pf\ Steps \stepref{s555nzX} and \stepref{s555nzL}, with $k=k'$ from
   \stepref{s555k}.~\qed\end{proof}
 \end{proof}
 \step{s555zero}{$\mult_{kI+A(X)}(0)=\mult_{kI+A(X')}(0)$.}
 \begin{proof}\pf\ By \ST{Theorem 2.2}, each multiplicity at $0$ is the order $n=n'$ minus the
  sum of the multiplicities at the nonzero eigenvalues; use \stepref{s555n} and
  \stepref{s555nonzero}.~\qed\end{proof}
 \qedstep
 \begin{proof}\pf\ By \stepref{s555nonzero}, \stepref{s555zero} and \ST{Proposition 2.6},
  $\Spec(kI+A(X))=\Spec(kI+A(X'))$. Shifting by $-k$ with \ST{Theorem 3.1} gives
  $\Spec(X)=\Spec(X')$, and \ST{Proposition 2.6} gives cospectrality.~\qed\end{proof}
\end{proof}

\medskip\noindent\textbf{Lemma 5.6 and Remark 5.7.}

We will need the paper's analysis of a connected bipartite graph with a regular line graph
twice, for $X_1$ and for $X_2$ in Lemma 5.6. We state it as two lemmas, so that no step is cited
from inside another proof.

\begin{namedresult}{Lemma}{(a bipartite root graph)}\label{lem:rootgraph}
 \lean{SimpleGraph.IsBiregularWith.lineGraph_isRegularOfDegree}
 \lean{SimpleGraph.IsBiregularWith.signlessLapMatrix_eq_fromBlocks}
 \lean{SimpleGraph.IsBiregularWith.multiplicity_zero_signlessLapMatrix}
 \lean{SimpleGraph.IsBiregularWith.card_eq_rank_add_one}
 \lean{SimpleGraph.IsBiregularWith.spectrumMultiset_map_quadratic_signlessLapMatrix}
 \uses{def:bipartition,lem:5.1,facts:sec5,hc:facts:finite,st:2.2,st:3.1,st:3.2,st:3.5,st:6.1,st:6.2,st:6.3,st:6.6}
 \assume{\begin{pfenum}
  \item $Z$ is a connected graph with at least one edge.
  \item $(P,Q)$ is a bipartition of $Z$, biregular with degrees $(k_1,k_2)$, with biadjacency
   matrix $B$.
  \item $M$ is the incidence matrix of $Z$ and $C=MM^T$.
 \end{pfenum}}
 \prove{\begin{pfenum}
  \item $Z$ has no isolated vertices, $k_1,k_2\ge1$, and $|P|k_1=|Q|k_2=|E(Z)|$.
  \item $L(Z)$ is connected and $(k_1+k_2-2)$-regular, and $\lmax(L(Z))=k_1+k_2-2$.
  \item In block form for $V(Z)=P\cup Q$, $C=\pmat{k_1I&B\\B^T&k_2I}$, and
   $BB^T\Zj=k_1k_2\Zj$.
  \item $\mult_C(0)=1$.
  \item $|V(Z)|=\rk(2I+A(L(Z)))+1$.
  \item Let $\rho=\rk B$, let $\Lambda$ be the multiset of positive eigenvalues of $BB^T$, and
   let $f(z)=(z-k_1)(z-k_2)$. Then $\Lambda$ has $\rho$ members, and the image of $\Spec(C)$
   under $f$ is $\Lambda+\Lambda+\ms{0^{|V(Z)|-2\rho}}$.
 \end{pfenum}}
\end{namedresult}
\begin{proof}
 \pfsketch\ This is the first part of the paper's proof of Lemma 5.6. The kernel of $C$ is the
 image of the eigenspace of $C$ for $k_1+k_2$ under an invertible map, and $k_1+k_2$ is a simple
 eigenvalue since $L(Z)$ is connected.
 \step{s5rgiso}{Assertion 1 holds.}
 \begin{proof}
  \step{s5rgnoiso}{$Z$ has no isolated vertices.}
  \begin{proof}\pf\ An edge gives two vertices, so every vertex $u$ has a vertex $w\ne u$ and a
   walk from $u$ to $w$, by \asm[1]{0}; its first edge is at $u$.~\qed\end{proof}
  \step{s5rgdeg}{$P$ and $Q$ are nonempty and $k_1,k_2\ge1$.}
  \begin{proof}\pf\ An edge has an end in $P$ and one in $Q$ by \asm[2]{0}; the degrees of
   these ends are positive by \stepref{s5rgnoiso}.~\qed\end{proof}
  \qedstep
  \begin{proof}\pf\ Each edge has exactly one end in $P$ and one in $Q$, so double counting the
   pairs $(u,e)$ with $u\in P$ an end of $e$, and the pairs $(w,e)$ with $w\in Q$ an end of $e$,
   gives $|P|k_1=|E(Z)|$ and $|Q|k_2=|E(Z)|$, by
   \hyperref[hc:facts:finite]{(finite counting and order)}.~\qed\end{proof}
 \end{proof}
 \step{s5rgline}{Assertion 2 holds.}
 \begin{proof}\pf\ \ST{Lemma 6.3} applies by \asm[1]{0} and \stepref{s5rgiso}; an edge has one
  end in $P$ and one in $Q$, so its degree in $L(Z)$ is $k_1+k_2-2$. Then
  Lemma~\ref{lem:5.1} gives $\lmax$.~\qed\end{proof}
 \step{s5rgblock}{Assertion 3 holds.}
 \begin{proof}
  \step{s5rgblockC}{$C=\pmat{k_1I&B\\B^T&k_2I}$.}
  \begin{proof}\pf\ By \ST{Theorem 6.2}, $C=D+A(Z)$ with $D$ the degree matrix, which is $k_1I$
   on $P$ and $k_2I$ on $Q$ by \asm[2]{0}; $A(Z)$ vanishes on $P\times P$ and $Q\times Q$ by
   \asm[2]{0}, and its $P\times Q$ block is $B$ by Definition~\ref{def:bipartition}.~\qed\end{proof}
  \step{s5rgblockj}{$B\Zj=k_1\Zj$ and $B^T\Zj=k_2\Zj$.}
  \begin{proof}\pf\ The row sum of $B$ at $u\in P$ is $\deg(u)=k_1$, and the column sum at $w\in
   Q$ is $\deg(w)=k_2$, by \asm[2]{0}.~\qed\end{proof}
  \qedstep
  \begin{proof}\pf\ \stepref{s5rgblockC}, and $BB^T\Zj=k_2B\Zj=k_1k_2\Zj$ by
   \stepref{s5rgblockj}.~\qed\end{proof}
 \end{proof}
 \step{s5rgsimple}{$\mult_C(k_1+k_2)=1$.}
 \begin{proof}
  \step{s5rgsimpleA}{$\mult_C(k_1+k_2)=\mult_{2I+A(L(Z))}(k_1+k_2)$.}
  \begin{proof}\pf\ \ST{Theorem 6.1} for $M$ and $M^T$, since $k_1+k_2\ne0$, with
   $M^TM=2I+A(L(Z))$ from \ST{Theorem 6.2}.~\qed\end{proof}
  \step{s5rgsimpleB}{$\mult_{2I+A(L(Z))}(k_1+k_2)=\mult_{L(Z)}(k_1+k_2-2)=1$.}
  \begin{proof}\pf\ The first equality is \ST{Theorem 3.1}; the second is Lemma~\ref{lem:5.1},
   since $L(Z)$ is connected and $(k_1+k_2-2)$-regular by \stepref{s5rgline}.~\qed\end{proof}
  \qedstep
  \begin{proof}\pf\ Steps \stepref{s5rgsimpleA} and \stepref{s5rgsimpleB}.~\qed\end{proof}
 \end{proof}
 \step{s5rgmap}{The map $\phi(\Zx,\Zy)=(-k_2\Zx,k_1\Zy)$ on $\re^P\oplus\re^Q$ is invertible and
  carries $\ker(C-(k_1+k_2)I)$ onto $\ker C$.}
 \begin{proof}
  \step{s5rgeigen}{$C(\Zx,\Zy)=(k_1+k_2)(\Zx,\Zy)$ if and only if $B\Zy=k_2\Zx$ and
   $B^T\Zx=k_1\Zy$.}
  \begin{proof}\pf\ Block multiplication with \stepref{s5rgblock}: $C(\Zx,\Zy)=(k_1\Zx+B\Zy,
   B^T\Zx+k_2\Zy)$.~\qed\end{proof}
  \step{s5rgkernel}{$C\phi(\Zx,\Zy)=0$ if and only if $B\Zy=k_2\Zx$ and $B^T\Zx=k_1\Zy$.}
  \begin{proof}\pf\ Block multiplication gives
   $C\phi(\Zx,\Zy)=\bigl(k_1(B\Zy-k_2\Zx),\,k_2(k_1\Zy-B^T\Zx)\bigr)$, and $k_1,k_2\ge1$ by
   \stepref{s5rgiso}.~\qed\end{proof}
  \qedstep
  \begin{proof}\pf\ $\phi$ is invertible since $k_1,k_2\ne0$; use \stepref{s5rgeigen} and
   \stepref{s5rgkernel}.~\qed\end{proof}
 \end{proof}
 \step{s5rgzero}{Assertion 4 holds.}
 \begin{proof}\pf\ By \ST{Theorem 2.2}, $\mult_C(0)=\dim\ker C$ and
  $\mult_C(k_1+k_2)=\dim\ker(C-(k_1+k_2)I)$; these are equal by \stepref{s5rgmap} and
  \path{LinearEquiv.finrank_eq}, and the second is 1 by \stepref{s5rgsimple}.~\qed\end{proof}
 \step{s5rgrank}{Assertion 5 holds.}
 \begin{proof}\pf\ $\rk C=|V(Z)|-1$ by \ST{Theorem 3.2} and \stepref{s5rgzero}, and
  $\rk C=\rk(MM^T)=\rk(M^TM)=\rk(2I+A(L(Z)))$ by \ST{Theorem 3.5} and
  \ST{Theorem 6.2}.~\qed\end{proof}
 \step{s5rgquad}{Assertion 6 holds.}
 \begin{proof}
  \pflet{$\mu_1,\dots,\mu_\rho$ be the members of $\Lambda$, which has $\rho$ members by
   \ST{Theorem 6.1}, and
   $\gamma_j^{\pm}=\frac12\bigl(k_1+k_2\pm\sqrt{(k_1-k_2)^2+4\mu_j}\bigr)$.}
  \step{s5rgroots}{$f(\gamma_j^+)=f(\gamma_j^-)=\mu_j$ for $1\le j\le\rho$.}
  \begin{proof}\pf\ $\gamma_j^{\pm}$ are the roots of $z^2-(k_1+k_2)z+k_1k_2-\mu_j$, by the
   quadratic formula, and $f(z)-\mu_j$ is this polynomial.~\qed\end{proof}
  \step{s5rgspecC}{$\Spec(C)$ consists of the $\gamma_j^{\pm}$, of $k_1$ with multiplicity
   $|P|-\rho$ and of $k_2$ with multiplicity $|Q|-\rho$.}
  \begin{proof}\pf\ \ST{Theorem 6.6} with $s=k_1$, $t=k_2$, and \stepref{s5rgblock}.~\qed\end{proof}
  \qedstep
  \begin{proof}\pf\ Apply $f$ to \stepref{s5rgspecC}: by \stepref{s5rgroots} and
   $f(k_1)=f(k_2)=0$, the image is $\Lambda+\Lambda+\ms{0^{|P|+|Q|-2\rho}}$, and
   $|P|+|Q|=|V(Z)|$.~\qed\end{proof}
 \end{proof}
 \qedstep
 \begin{proof}\pf\ Steps \stepref{s5rgiso}, \stepref{s5rgline}, \stepref{s5rgblock},
  \stepref{s5rgzero}, \stepref{s5rgrank} and \stepref{s5rgquad}.~\qed\end{proof}
\end{proof}

\begin{namedresult}{Lemma}{(two bipartite root graphs)}\label{lem:tworoot}
 \lean{SimpleGraph.IsBiregularWith.eq_of_lineGraph_cospectral}
 \uses{lem:rootgraph,facts:sec5,st:2.2,st:2.4,st:2.6,st:3.1,st:3.2,st:6.1,st:6.2}
 \assume{\begin{pfenum}
  \item $Z$ and $Z'$ are connected graphs, each with at least one edge.
  \item $(P,Q)$ is a bipartition of $Z$, biregular with degrees $(k_1,k_2)$, with biadjacency
   matrix $B$, and $|P|\le|Q|$.
  \item $(P',Q')$ is a bipartition of $Z'$, biregular with degrees $(k_1',k_2')$, with
   biadjacency matrix $B'$, and $|P'|\le|Q'|$.
  \item $L(Z)$ and $L(Z')$ are cospectral.
 \end{pfenum}}
 \prove{$|P|=|P'|$, $|Q|=|Q'|$, $k_1=k_1'$, $k_2=k_2'$, and $\Spec(BB^T)=\Spec(B'B'^T)$.}
\end{namedresult}
\begin{proof}
 \pfsketch\ This is the paper's claim that the spectrum of $L(X_1)$ determines $p$, $q$, $k_1$,
 $k_2$ and $\Spec(B_1B_1^T)$. We will recover the sizes and degrees from the number of
 edges, the valency and the rank, and then read $\Spec(BB^T)$ off $\Spec(C)$ through the
 quadratic $f(z)=(z-k_1)(z-k_2)$.
 \pflet{$m=|E(Z)|$, $N=|V(Z)|$, $C=MM^T$ for the incidence matrix $M$ of $Z$, and
  $m'=|E(Z')|$, $N'=|V(Z')|$, $C'=M'M'^T$ for the incidence matrix $M'$ of $Z'$.}
 \step{s5trm}{$m=m'$.}
 \begin{proof}\pf\ These are the numbers of vertices of the cospectral graphs $L(Z)$, $L(Z')$;
  apply \ST{Proposition 2.6}.~\qed\end{proof}
 \step{s5trs}{$k_1+k_2=k_1'+k_2'$.}
 \begin{proof}\pf\ $k_1+k_2-2$ and $k_1'+k_2'-2$ are the largest eigenvalues of $L(Z)$ and
  $L(Z')$, by Lemma~\ref{lem:rootgraph}; cospectral graphs have the same largest eigenvalue by
  \ST{Proposition 2.6} and \ST{Lemma 2.4}.~\qed\end{proof}
 \step{s5trN}{$N=N'$.}
 \begin{proof}
  \step{s5trrank}{$\rk(2I+A(L(Z)))=\rk(2I+A(L(Z')))$.}
  \begin{proof}\pf\ \ST{Theorem 3.2} with $f(z)=z+2$ expresses each rank through the common
   multiplicities of the two line graphs, which agree by \ST{Proposition 2.6}.~\qed\end{proof}
  \qedstep
  \begin{proof}\pf\ Lemma~\ref{lem:rootgraph}, assertion 5, for $Z$ and $Z'$, and
   \stepref{s5trrank}.~\qed\end{proof}
 \end{proof}
 \step{s5trprod}{$k_1k_2N=m(k_1+k_2)$ and $k_1'k_2'N'=m'(k_1'+k_2')$.}
 \begin{proof}\pf\ By Lemma~\ref{lem:rootgraph}, assertion 1, for $Z$ and for $Z'$,
  $N=|P|+|Q|=m/k_1+m/k_2$ and $N'=m'/k_1'+m'/k_2'$; multiply by $k_1k_2$ and by
  $k_1'k_2'$.~\qed\end{proof}
 \step{s5trorder}{$k_1\ge k_2$ and $k_1'\ge k_2'$.}
 \begin{proof}\pf\ By Lemma~\ref{lem:rootgraph}, assertion 1, for $Z$ and for $Z'$,
  $|P|k_1=|Q|k_2=m>0$ and $|P'|k_1'=|Q'|k_2'=m'>0$, with $|P|\le|Q|$ and
  $|P'|\le|Q'|$.~\qed\end{proof}
 \step{s5trdeg}{$k_1=k_1'$ and $k_2=k_2'$.}
 \begin{proof}\pf\ By \stepref{s5trm} to \stepref{s5trprod}, the pairs have the same sum and the
  same product, so $(k_1-k_2)^2=(k_1'-k_2')^2$; both differences are nonnegative by
  \stepref{s5trorder}.~\qed\end{proof}
 \step{s5trsize}{$|P|=|P'|$ and $|Q|=|Q'|$.}
 \begin{proof}\pf\ $|P|=m/k_1$, $|Q|=m/k_2$, $|P'|=m'/k_1'$ and $|Q'|=m'/k_2'$ by
  Lemma~\ref{lem:rootgraph}, assertion 1, for $Z$ and $Z'$; use \stepref{s5trm} and
  \stepref{s5trdeg}.~\qed\end{proof}
 \step{s5trC}{$\Spec(C)=\Spec(C')$.}
 \begin{proof}
  \step{s5trCnz}{$\mult_C(\gamma)=\mult_{C'}(\gamma)$ for every $\gamma\ne0$.}
  \begin{proof}\pf\ By \ST{Theorem 6.1} and \ST{Theorem 6.2}, $\mult_C(\gamma)=
   \mult_{2I+A(L(Z))}(\gamma)$ and $\mult_{C'}(\gamma)=\mult_{2I+A(L(Z'))}(\gamma)$; these agree
   by \ST{Theorem 3.1} and \ST{Proposition 2.6}.~\qed\end{proof}
  \qedstep
  \begin{proof}\pf\ $\mult_C(0)=\mult_{C'}(0)=1$ by Lemma~\ref{lem:rootgraph}, assertion 4; with
   \stepref{s5trCnz} and \ST{Proposition 2.6}.~\qed\end{proof}
 \end{proof}
 \pflet{$f(z)=(z-k_1)(z-k_2)$, and let $\Lambda$ and $\Lambda'$ be the multisets of positive
  eigenvalues of $BB^T$ and $B'B'^T$, and $\rho=\rk B$ and $\rho'=\rk B'$.}
 \step{s5trmap}{The images of $\Spec(C)$ and $\Spec(C')$ under $f$ are
  $\Lambda+\Lambda+\ms{0^{N-2\rho}}$ and $\Lambda'+\Lambda'+\ms{0^{N-2\rho'}}$.}
 \begin{proof}\pf\ Lemma~\ref{lem:rootgraph}, assertion 6, for $Z$, and for $Z'$, whose
  polynomial $(z-k_1')(z-k_2')$ is $f$ by \stepref{s5trdeg}, with $N'=N$ by
  \stepref{s5trN}.~\qed\end{proof}
 \step{s5trLambda}{$\Lambda=\Lambda'$.}
 \begin{proof}\pf\ By \stepref{s5trC} the two images agree. For $\mu>0$, the count of $\mu$ in
  them is $2\,\mathrm{count}_\Lambda(\mu)$ and $2\,\mathrm{count}_{\Lambda'}(\mu)$, by
  \stepref{s5trmap}; and $\Lambda$, $\Lambda'$ have only positive
  members.~\qed\end{proof}
 \qedstep
 \begin{proof}\pf\ By \ST{Theorem 6.1}, $\Spec(BB^T)=\Lambda+\ms{0^{|P|-\rho}}$ and
  $\Spec(B'B'^T)=\Lambda'+\ms{0^{|P'|-\rho'}}$; these agree by \stepref{s5trLambda}, which gives
  $\rho=\rho'$, and \stepref{s5trsize}. With \stepref{s5trdeg} and \stepref{s5trsize} this is
  the claim.~\qed\end{proof}
\end{proof}

\begin{namedresult}{Lemma}{(biregular bipartitions)}\label{lem:biregular}
 \lean{SimpleGraph.exists_isBiregularWith_of_lineGraph_regular}
 \uses{def:bipartition,lem:5.4}
 \assume{\begin{pfenum}
  \item $Z$ is a bipartite graph without isolated vertices.
  \item $L(Z)$ is connected and regular.
 \end{pfenum}}
 \prove{$Z$ has a bipartition $(P,Q)$ with $|P|\le|Q|$ that is biregular with some degrees
  $(k_1,k_2)$.}
\end{namedresult}
\begin{proof}
 \pfsketch\ Lemma 5.4 gives a biregular bipartition or regularity, and in a regular bipartite
 graph every bipartition is biregular.
 \step{s5brget}{$Z$ has a biregular bipartition.}
 \begin{proof}
  \step{s5brreg}{\case{$Z$ is $k$-regular.}}
  \begin{proof}\pf\ Any bipartition, which exists by \asm[1]{0}, is biregular with degrees
   $(k,k)$.~\qed\end{proof}
  \step{s5brsemi}{\case{$Z$ has a bipartition that is biregular with degrees $(k_1,k_2)$,
   $k_1\ne k_2$.}}
  \begin{proof}\pf\ This is the claim.~\qed\end{proof}
  \qedstep
  \begin{proof}\pf\ By Lemma~\ref{lem:5.4}, with \asm[1]{0} and \asm[2]{0}, the cases
   \stepref{s5brreg} and \stepref{s5brsemi} are exhaustive.~\qed\end{proof}
 \end{proof}
 \qedstep
 \begin{proof}\pf\ Exchange the two parts of the bipartition of \stepref{s5brget} if
  necessary; this exchanges the two degrees.~\qed\end{proof}
\end{proof}

\begin{cgss}{Lemma}{5.6}\label{lem:5.6}
 \lean{SimpleGraph.lineGraph_cospectral_bipartite_forall_isBiregularWith}
 \uses{def:bipartition,lem:rootgraph,lem:tworoot,lem:biregular,st:5.5,st:6.3,st:6.7}
 \assume{\begin{pfenum}
  \item $X_1$ is a connected bipartite graph with at least one edge, and $L(X_1)$ is
   $h$-regular.
  \item $X_2$ is a graph without isolated vertices.
  \item $L(X_1)$ and $L(X_2)$ are cospectral.
 \end{pfenum}}
 \prove{$X_2$ is connected, and one of the following holds.
  \begin{pfenum}
   \item $X_2$ is bipartite, and for every bipartition $(P_1,Q_1)$ of $X_1$ with
    $|P_1|\le|Q_1|$ that is biregular with degrees $(k_1,k_2)$ there is a bipartition
    $(P_2,Q_2)$ of $X_2$ with $|P_2|=|P_1|$ and $|Q_2|=|Q_1|$ that is biregular with degrees
    $(k_1,k_2)$, such that $\Spec(B_1B_1^T)=\Spec(B_2B_2^T)$ for the biadjacency matrices
    $B_1$, $B_2$ of $X_1$, $X_2$ for these bipartitions.
   \item $X_1$ is not regular; $X_2$ is $\kappa$-regular and not bipartite, for the natural
    number $\kappa$ with $2\kappa=h+2$; and $|V(X_1)|=|V(X_2)|+1$.
  \end{pfenum}}
\end{cgss}
\begin{proof}
 \pfsketch\ \repair\ We need \asm[2]{0} (errata, Implicit hypothesis 5: $X_1=K_{1,3}$ and $X_2$ a
 triangle with an isolated vertex have the same line graph). \new\ We state (i) for every
 biregular bipartition of $X_1$, so that Theorem 5.8 can choose the points of a design as $P_1$;
 the toolkit's Theorem 6.7 orders the parts by size instead, and so this lemma has its own Lean
 name. We will take connectedness and (ii) from the toolkit's Theorem 6.7, which follows the
 paper: in the nonbipartite case $M_2M_2^T$ is nonsingular by Lemma 5.3, so $X_2$ has $p+q-1$
 vertices. For (i) we will follow the first part of the paper's proof, which is
 Lemmas~\ref{lem:rootgraph} and~\ref{lem:tworoot}.
 \step{s556L1}{$X_1$ has no isolated vertices, and $L(X_1)$ is connected and $h$-regular on a
  nonempty vertex set.}
 \begin{proof}\pf\ By \asm[1]{0} and \ST{Lemma 6.3}: a connected graph with an edge has no
  isolated vertices, since each vertex has a walk to an end of the edge that is different from
  it.~\qed\end{proof}
 \step{s556L2}{$L(X_2)$ is connected and $h$-regular.}
 \begin{proof}\pf\ \ST{Corollary 5.5}, with \stepref{s556L1} and \asm[3]{0}.~\qed\end{proof}
 \step{s556st}{$X_2$ has an edge and is connected, and if $X_2$ is not bipartite, then $X_1$ is
  not regular, $X_2$ is regular, and $|V(X_1)|=|V(X_2)|+1$.}
 \begin{proof}\pf\ \ST{Theorem 6.7}, with \asm[1]{0}, \asm[2]{0} and \asm[3]{0}; $L(X_2)$ is
  nonempty by \stepref{s556L2}, so $X_2$ has an edge.~\qed\end{proof}
 \step{s556bip}{\case{$X_2$ is bipartite.}\prove{(i) holds.}}
 \begin{proof}
  \step{s556given}{\assume{$(P_1,Q_1)$ is a bipartition of $X_1$ with $|P_1|\le|Q_1|$ that is
   biregular with degrees $(k_1,k_2)$.}
   \prove{There is a bipartition $(P_2,Q_2)$ of $X_2$ as in (i).}}
  \begin{proof}
   \step{s556pick}{\pick\ a bipartition $(P_2,Q_2)$ of $X_2$ with $|P_2|\le|Q_2|$ that is
    biregular with some degrees $(k_1',k_2')$.}
   \begin{proof}\pf\ Lemma~\ref{lem:biregular} for $X_2$, with \asm{1}, \asm[2]{0}
    and \stepref{s556L2}.~\qed\end{proof}
   \qedstep
   \begin{proof}\pf\ Lemma~\ref{lem:tworoot} with $Z=X_1$, $Z'=X_2$ applies by \asm[1]{0},
    \stepref{s556st}, \asm{2}, \stepref{s556pick} and \asm[3]{0}; it gives the sizes, the
    degrees $k_1'=k_1$, $k_2'=k_2$ and the spectra.~\qed\end{proof}
  \end{proof}
  \qedstep
  \begin{proof}\pf\ \asm{1} and \stepref{s556given} are (i).~\qed\end{proof}
 \end{proof}
 \step{s556nonbip}{\case{$X_2$ is not bipartite.}\prove{(ii) holds.}}
 \begin{proof}
  \step{s556reg}{$X_2$ is $\kappa$-regular for a natural number $\kappa\ge1$ with $2\kappa=h+2$.}
  \begin{proof}\pf\ $X_2$ is regular by \stepref{s556st} and \asm{1}, of some valency $\kappa$;
   $\kappa\ge1$ since $X_2$ has an edge, by \stepref{s556st}. By the degree formula of
   \ST{Lemma 6.3}, with \asm[2]{0}, an edge of $X_2$ has degree $2\kappa-2$ in $L(X_2)$, which is
   $h$ by \stepref{s556L2}.~\qed\end{proof}
  \qedstep
  \begin{proof}\pf\ Steps \stepref{s556st} and \stepref{s556reg}, with \asm{1}, are
   (ii).~\qed\end{proof}
 \end{proof}
 \qedstep
 \begin{proof}\pf\ Step \stepref{s556st}, and the cases \stepref{s556bip} and
  \stepref{s556nonbip}, which are exhaustive.~\qed\end{proof}
\end{proof}

\begin{cgss}{Remark}{5.7}\label{rem:5.7}
 \lean{SimpleGraph.lineGraph_cospectral_bipartite_spectrum_of_not_isBipartite}
 \uses{def:bipartition,lem:5.6,lem:rootgraph,st:6.8}
 \assume{\begin{pfenum}
  \item $X_1$ is a connected graph with at least one edge, and $(P,Q)$ is a bipartition of
   $X_1$, biregular with degrees $(k_1,k_2)$ with $k_1\ne k_2$, with $p=|P|\le q=|Q|$ and
   biadjacency matrix $B$ of rank $\rho$.
  \item $X_2$ is a graph without isolated vertices that is not bipartite, and $L(X_1)$ and
   $L(X_2)$ are cospectral.
 \end{pfenum}}
 \prove{$k_1>k_2\ge1$; $k_1k_2$ is an eigenvalue of $BB^T$ of multiplicity 1; and if
  $\Lambda'$ is the multiset of positive eigenvalues of $BB^T$ with one copy of $k_1k_2$ removed,
  then
  \[
   \Spec(X_2)=\ms{\tfrac12(k_1+k_2)}
    +\sum_{\mu\in\Lambda'}\ms{\delta_\mu,-\delta_\mu}
    +\ms{\bigl(\tfrac12(k_1-k_2)\bigr)^{p-\rho},\bigl(-\tfrac12(k_1-k_2)\bigr)^{q-\rho}},
  \]
  where $\delta_\mu=\sqrt{\mu+\frac14(k_1-k_2)^2}$.}
 We state only the conclusion of the remark, the spectrum of $A(X_2)$. Its intermediate
 assertions, that each $\mu$ gives the eigenvalues $\pm\sqrt\mu$ of $A(X_1)$, $\gamma$ and
 $\gamma'$ of $C$, and $\gamma-2$ and $\gamma'-2$ of $A(L(X_1))$, have no Lean statements of
 their own, and neither has the paper's comment that the multiplicities of $\pm\frac12(k_1-k_2)$
 need not be equal (the statement above gives them, as $p-\rho$ and $q-\rho$). The spectrum of a
 matrix of the form of $C$ is \path{Matrix.spectrumMultiset_fromBlocks_scalar_of_ne}.
\end{cgss}
\begin{proof}
 \pfsketch\ \repair\ The paper says that each eigenvalue $\mu\ne0$ of $B_1B_1^T$ gives two
 eigenvalues $\gamma$, $\gamma'$ of $C$, the roots of $(\gamma-k_1)(\gamma-k_2)=\mu$, and two
 eigenvalues $\delta=\gamma-\frac12(k_1+k_2)$ and $\delta'=\gamma'-\frac12(k_1+k_2)$ of $A(X_2)$
 (errata, Error 13, item 8, for $\delta'$), with $\delta^2=\mu+\frac14(k_1-k_2)^2$. This fails for
 $\mu=k_1k_2$, whose roots are $k_1+k_2$ and $0$, and $0$ is not an eigenvalue of the nonsingular
 $M_2M_2^T$; and the multiplicities of $\pm\frac12(k_1-k_2)$ are $p-\rho$ and $q-\rho$, whose
 difference, not sum, is $q-p$ (errata, Error 14). The toolkit's Corollary 6.8 is the paper's
 computation with these corrections, for $X_2$ connected, regular and not bipartite. We will
 derive those hypotheses from Lemma 5.6; the form that assumes only that $X_2$ has no isolated
 vertices is the one that Theorem 5.8 uses, so this remark has its own Lean name.
 \step{s557deg}{$k_1>k_2\ge1$.}
 \begin{proof}\pf\ By Lemma~\ref{lem:rootgraph}, assertion 1, $k_2\ge1$ and $pk_1=qk_2$, and
  $p\le q$, $k_1\ne k_2$.~\qed\end{proof}
 \step{s557ii}{$X_2$ is connected and regular, and not bipartite.}
 \begin{proof}\pf\ Lemma~\ref{lem:5.6} applies, since $X_1$ is bipartite by \asm[1]{0} and
  $L(X_1)$ is $(k_1+k_2-2)$-regular by Lemma~\ref{lem:rootgraph}, assertion 2; its alternative
  (i) fails since $X_2$ is not bipartite, so (ii) holds.~\qed\end{proof}
 \qedstep
 \begin{proof}\pf\ \ST{Corollary 6.8}, with \asm[1]{0}, \stepref{s557deg}, \stepref{s557ii} and
  \asm[2]{0}; its $\mu_2,\dots,\mu_\rho$ are the members of $\Lambda'$, and its
  $\pm\sqrt{\mu_j+\frac14(k_1-k_2)^2}$ are $\pm\delta_{\mu_j}$.~\qed\end{proof}
\end{proof}

\subsection{Theorem 5.8 and its consequences}

\begin{cgss}{Theorem}{5.8}\label{thm:5.8}
 \lean{Design.flagGraph_cospectral}
 \uses{def:flag,def:bipartition,hc:def:design,lem:design,lem:incidence,lem:dense,lem:isolated,lem:5.1,lem:5.2,lem:5.6,rem:5.7,def:CP,def:represented,thm:4.11,def:3.2,facts:sec5,st:2.6,st:3.2,st:3.4,st:3.5,st:5.1,st:5.5,st:7.7,st:7.10}
 \assume{\begin{pfenum}
  \item $D_1$ is a design with parameters $(v,k,b,r,\lambda)$ on $P$ and $\cB$, with incidence
   matrix $N$, and $b\ge v$.
  \item $X_1=F(D_1)$, and $X_2$ is a graph cospectral with $X_1$.
 \end{pfenum}}
 \prove{One of the following holds.
  \begin{pfenum}
   \item $v+b\le9$, and $X_2$ is represented in $\typ E_8$.
   \item $X_2\cong F(D_2)$ for a design $D_2$ with the parameters $(v,k,b,r,\lambda)$.
   \item There are integers $s\ge2$ and $t\ge1$ with $t(s+1)$ even, $s\ge2t-1$ and
    $(s-1)\mid2t(t-1)$, such that
    \[
     v=\tfrac12s(s+1),\quad k=st,\quad b=\tfrac12(s+1)(s+2),\quad r=(s+2)t,\quad
     \lambda(s-1)=2t(st-1),
    \]
    and $X_2\cong L(Y)$ for a connected $(s+1)t$-regular graph $Y$ on $s(s+2)$ vertices with
    \[
     \Spec(Y)=\ms{\bigl((s+1)t\bigr)^{1},(-t)^{s+1},(\sqrt c)^{\ell},(-\sqrt c)^{\ell}},\qquad
     c=\frac{t(s+1)(s-t)}{s-1},\quad \ell=\tfrac12(s-1)(s+2).
    \]
  \end{pfenum}}
\end{cgss}
\begin{proof}
 \pfsketch\ \repair\ The hypothesis $b\ge v$ is the paper's remark after Theorem 5.8: by
 Lemma~\ref{lem:design} it holds when $\lambda<r$, and otherwise every block contains every point
 and one passes to the dual design, as Corollary 5.9 does. Without it the theorem is false
 (errata, Error 16): the design $(10,10,6,6,6)$ has flag graph $L(K_{6,10})$, cospectral with
 $L(T(6))$, and none of (i), (ii), (iii) holds. Here $v+b=16$; (ii) fails since the only design
 with these parameters has flag graph $L(K_{6,10})$, and $T(6)$ is not bipartite, so Whitney's
 theorem excludes $L(T(6))\cong L(K_{6,10})$; and $v=10$ gives $s=4$ and $t=k/s=\frac52$. The
 details are in Error 16. We will follow the paper. The spectrum makes $X_2$
 regular and connected with $2I+A(X_2)\succeq0$ of rank $v+b-1$, so Theorem 4.11 applies. A
 representation in $\typ E_8$ gives (i), a cocktail party graph gives (ii), and a line graph
 $L(Y)$ gives (ii) or (iii) according to the two cases of Lemma 5.6.
 \step{s558X1}{$X_1$ is connected and $(r+k-2)$-regular on $vr$ vertices, and $2I+A(X_1)$ is
  positive semidefinite of rank $v+b-1$.}
 \begin{proof}\pf\ Lemma~\ref{lem:incidence}, assertions 3 and 4.~\qed\end{proof}
 \step{s558X2}{$X_2$ is connected and $(r+k-2)$-regular on $vr$ vertices, and $2I+A(X_2)$ is
  positive semidefinite of rank $v+b-1$.}
 \begin{proof}
  \step{s558X2reg}{$X_2$ is connected and $(r+k-2)$-regular on $vr$ vertices.}
  \begin{proof}\pf\ \ST{Corollary 5.5} and \ST{Proposition 2.6}, with \stepref{s558X1}; the
   corollary is Lemma~\ref{lem:5.2} and Lemma~\ref{lem:5.1} for both graphs, as in the
   paper.~\qed\end{proof}
  \step{s558X2psd}{$2I+A(X_2)$ is positive semidefinite.}
  \begin{proof}\pf\ By \ST{Corollary 3.4} and \stepref{s558X1}, every eigenvalue of $X_1$ is at
   least $-2$; $X_2$ has the same eigenvalues by \ST{Proposition 2.6}; apply \ST{Corollary 3.4}
   again.~\qed\end{proof}
  \step{s558X2rank}{$\rk(2I+A(X_2))=v+b-1$.}
  \begin{proof}\pf\ \ST{Theorem 3.2} with $f(z)=z+2$ gives the rank from the multiplicities,
   which agree by \ST{Proposition 2.6}; use \stepref{s558X1}.~\qed\end{proof}
  \qedstep
  \begin{proof}\pf\ Steps \stepref{s558X2reg}, \stepref{s558X2psd} and
   \stepref{s558X2rank}.~\qed\end{proof}
 \end{proof}
 \step{s558hrc}{$X_2\cong L(Y)$ for some graph $Y$, or $X_2\cong CP(m)$ for some $m$, or $X_2$ is
  represented in $\typ E_8$.}
 \begin{proof}\pf\ Theorem~\ref{thm:4.11}, with
  \stepref{s558X2}.~\qed\end{proof}
 \step{s558E8}{\case{$X_2$ is represented in $\typ E_8$.}\prove{(i) holds.}}
 \begin{proof}
  \step{s558E8rank}{$\rk(2I+A(X_2))\le8$.}
  \begin{proof}\pf\ $2I+A(X_2)$ is the Gram matrix of a family in $\re^8$
   (Definition~\ref{def:3.2}), so by \ST{Theorem 3.5} its rank is the dimension of a subspace of
   $\re^8$; apply \path{Submodule.finrank_le} and \path{finrank_euclideanSpace_fin}.~\qed\end{proof}
  \qedstep
  \begin{proof}\pf\ $v+b-1\le8$ by \stepref{s558X2} and \stepref{s558E8rank}, and $X_2$ is
   represented in $\typ E_8$ by \asm{1}.~\qed\end{proof}
 \end{proof}
 \step{s558CP}{\case{$X_2\cong CP(m)$.}\prove{(ii) holds, with $D_2=D_1$.}}
 \begin{proof}
  \pfsketch\ \repair\ The paper does not exclude this case of Theorem 4.11 (errata, Gap 15). The
  errata's fill forces $v=k=b=r=2$ and hence case (i), which would also need a representation in
  $\typ E_8$; we reach (ii) instead, since both graphs are then cocktail party graphs.
  \step{s558CPX2}{$X_2$ is $(2m-2)$-regular on $2m$ vertices, and $m\ge1$.}
  \begin{proof}\pf\ $CP(m)$ is, by Definition~\ref{def:CP}, and
   $X_2$ has $vr\ge1$ vertices by \stepref{s558X2}.~\qed\end{proof}
  \step{s558CPX1}{$X_1\cong CP(m)$.}
  \begin{proof}\pf\ By \ST{Corollary 5.5} and \ST{Proposition 2.6}, applied to $X_2$ and its
   cospectral mate $X_1$, $X_1$ is $(2m-2)$-regular on $2m$ vertices by \stepref{s558CPX2}; apply
   Lemma~\ref{lem:dense}, assertion 2.~\qed\end{proof}
  \qedstep
  \begin{proof}\pf\ $X_2\cong CP(m)\cong X_1=F(D_1)$ by \asm{1} and
   \stepref{s558CPX1}.~\qed\end{proof}
 \end{proof}
 \step{s558line}{\case{$X_2\cong L(Y)$.}\prove{(ii) or (iii) holds.}}
 \begin{proof}
  \step{s558Yp}{\pick\ a graph $Y'$ without isolated vertices with $X_2\cong L(Y')$.}
  \begin{proof}\pf\ Lemma~\ref{lem:isolated} and \asm{1}.~\qed\end{proof}
  \step{s558cosp}{$L(I(D_1))$ and $L(Y')$ are cospectral.}
  \begin{proof}\pf\ $L(I(D_1))=X_1$ is cospectral with $X_2$ by \asm[2]{0}, and $X_2$ with
   $L(Y')$ by \stepref{s558Yp} and \ST{Proposition 2.6}; cospectrality is equality of
   characteristic polynomials.~\qed\end{proof}
  \step{s558apply}{Lemma~\ref{lem:5.6} applies to $X_1'=I(D_1)$, with $h=r+k-2$, and
   $X_2'=Y'$; in particular $Y'$ is connected.}
  \begin{proof}\pf\ $I(D_1)$ is connected and bipartite with an edge by
   Lemma~\ref{lem:incidence}, assertions 1 and 2, and $L(I(D_1))=X_1$ is $(r+k-2)$-regular by
   \stepref{s558X1}; with \stepref{s558Yp} and \stepref{s558cosp}.~\qed\end{proof}
  \step{s558bip}{\case{Alternative (i) of Lemma~\ref{lem:5.6} holds.}\prove{(ii) holds.}}
  \begin{proof}
   \step{s558bippick}{\pick\ a bipartition $(P_2,Q_2)$ of $Y'$ with $|P_2|=v$, $|Q_2|=b$,
    biregular with degrees $(r,k)$, whose biadjacency matrix $B_2$ has
    $\Spec(B_2B_2^T)=\Spec(NN^T)$.}
   \begin{proof}\pf\ Apply (i) to the bipartition $(P,\cB)$ of $I(D_1)$, which is biregular with
    degrees $(r,k)$ and has biadjacency matrix $N$ by Lemma~\ref{lem:incidence}, and has
    $|P|=v\le b=|\cB|$ by \asm[1]{0}.~\qed\end{proof}
   \step{s558bipgram}{$B_2B_2^T=(r-\lambda)I+\lambda J$.}
   \begin{proof}\pf\ \ST{Theorem 7.7} for the zero-one $v\times b$ matrix $B_2$, whose row sums
    are $r$ and column sums $k$ by \stepref{s558bippick}; its hypotheses $v>1$, $k>1$, $r>0$,
    $\lambda>0$, $r(k-1)=\lambda(v-1)$ and $\Spec(B_2B_2^T)=\ms{(rk)^{1},(r-\lambda)^{v-1}}$ hold
    by Definition~\ref{hc:def:design}, Lemma~\ref{lem:design} and
    \stepref{s558bippick}.~\qed\end{proof}
   \step{s558bipD2}{The rule $D_2(w)=\{u\in P_2:u\sim w\}$ for $w\in Q_2$ is a design with
    parameters $(v,k,b,r,\lambda)$ on $P_2$ and $Q_2$.}
   \begin{proof}\pf\ By Definition~\ref{hc:def:design}: $|P_2|=v$, $|Q_2|=b$; a block has $k$
    points and a point lies in $r$ blocks by \stepref{s558bippick}; two points lie in
    $(B_2B_2^T)(u,u')=\lambda$ common blocks by \stepref{s558bipgram}.~\qed\end{proof}
   \step{s558bipiso}{$I(D_2)\cong Y'$.}
   \begin{proof}\pf\ The map $P_2\sqcup Q_2\to V(Y')$ that is the inclusion on each part is a
    bijection, and it carries the incidences of $D_2$ onto the edges of $Y'$, since every edge of
    $Y'$ joins $P_2$ to $Q_2$ (Definition~\ref{def:bipartition}).~\qed\end{proof}
   \qedstep
   \begin{proof}\pf\ $F(D_2)=L(I(D_2))\cong L(Y')\cong X_2$ by \stepref{s558bipiso},
    \path{SimpleGraph.Iso.lineGraph} and \stepref{s558Yp}.~\qed\end{proof}
  \end{proof}
  \step{s558nb}{\case{Alternative (ii) of Lemma~\ref{lem:5.6} holds.}\prove{(iii) holds.}}
  \begin{proof}
   \step{s558nbfacts}{$r>k$, and $Y'$ is connected, not bipartite and $\kappa$-regular on
    $v+b-1$ vertices, with $vr$ edges, for the natural number $\kappa$ with $2\kappa=r+k$.}
   \begin{proof}
    \step{s558nbrk}{$r>k$.}
    \begin{proof}\pf\ $I(D_1)$ is not regular by (ii), so $r\ne k$ by
     Lemma~\ref{lem:incidence}; and $r=bk/v\ge k$ by Lemma~\ref{lem:design} and
     \asm[1]{0}.~\qed\end{proof}
    \qedstep
    \begin{proof}\pf\ By (ii), $Y'$ is not bipartite and $\kappa$-regular with
     $2\kappa=h+2=r+k$, and $|V(Y')|=(v+b)-1$; $Y'$ is connected by \stepref{s558apply}; and
     $|E(Y')|=|V(L(Y'))|=|V(X_2)|=vr$ by \stepref{s558Yp} and \stepref{s558X2}.~\qed\end{proof}
   \end{proof}
   \step{s558nbhand}{$(r+k)(v+b-1)=4vr$.}
   \begin{proof}\pf\ \path{SimpleGraph.sum_degrees_eq_twice_card_edges} for $Y'$ gives
    $\kappa(v+b-1)=2vr$ by \stepref{s558nbfacts}; multiply by 2.~\qed\end{proof}
   \step{s558nbv}{$v(r-k)^2=k(r+k)$.}
   \lean{Computations.C11.parameter_equation}
   \begin{proof}\pf\ Multiply \stepref{s558nbhand} by $k$ and substitute $bk=vr$ from
    Lemma~\ref{lem:design}: $(r+k)(vk+vr-k)=4vrk$, which rearranges to the claim (computation
    C11).~\qed\end{proof}
   \step{s558nbt}{$t=\frac12(r-k)$ is a positive integer.}
   \lean{Computations.C11.exists_positive_half_difference}
   \begin{proof}\pf\ $r+k=2\kappa$ is even by \stepref{s558nbfacts}, so $r-k=(r+k)-2k$ is even;
    and $r>k$.~\qed\end{proof}
   \step{s558nbtk}{$2t^2v=k(k+t)$ and $t(b-v)=k+t$; thus $t\mid k$.}
   \lean{Computations.C11.two_mul_sq_mul_points}\lean{Computations.C11.mul_blocks_sub_points}
   \lean{Computations.C11.dvd_block_size}
   \begin{proof}\pf\ Substitute $r=k+2t$ in \stepref{s558nbv} for the first; then
    $b-v=v(r-k)/k=2tv/k=(k+t)/t$ by $bk=vr$ (computation C11). As $t\mid k+t$, also
    $t\mid k$.~\qed\end{proof}
   \pflet{$s=k/t$, a positive integer by \stepref{s558nbtk}.}
   \step{s558nbparams}{$v=\frac12s(s+1)$, $k=st$, $b=\frac12(s+1)(s+2)$ and $r=(s+2)t$.}
   \lean{Computations.C11.parameters}
   \begin{proof}\pf\ Substitute $k=st$ in \stepref{s558nbtk} and in $r=k+2t$ and $bk=vr$
    (computation C11).~\qed\end{proof}
   \step{s558nbs}{$s\ge2$.}
   \lean{Computations.C11.two_le_parameter}
   \begin{proof}\pf\ $\frac12s(s+1)=v>1$ by \stepref{s558nbparams}.~\qed\end{proof}
   \step{s558nblam}{$\lambda(s-1)=2t(st-1)$ and $(s-1)\mid2t(t-1)$.}
   \lean{Computations.C11.lambda_equation}\lean{Computations.C11.lambda_remainder}
   \lean{Computations.C11.parameter_dvd}
   \begin{proof}\pf\ Substituting \stepref{s558nbparams} in $r(k-1)=\lambda(v-1)$ of
    Lemma~\ref{lem:design} gives $(s+2)t(st-1)=\frac12\lambda(s+2)(s-1)$; divide by $s+2$. Then
    $2t(t-1)=(\lambda-2t^2)(s-1)$ (computation C11).~\qed\end{proof}
   \step{s558nbineq}{$s\ge2t-1$.}
   \lean{Computations.C11.two_mul_sub_one_le}
   \begin{proof}\pf\ \repair\ The paper prints $s\le2t-1$ (errata, Error 4). By
    Lemma~\ref{lem:design}, $k\le v$, that is $st\le\frac12s(s+1)$ by \stepref{s558nbparams};
    divide by $s>0$.~\qed\end{proof}
   \step{s558nbY}{$Y'$ is connected and $(s+1)t$-regular on $s(s+2)$ vertices.}
   \begin{proof}\pf\ By \stepref{s558nbfacts} and \stepref{s558nbparams},
    $\kappa=\frac12(r+k)=(s+1)t$ and $v+b-1=(s+1)^2-1=s(s+2)$.~\qed\end{proof}
   \step{s558nbspec}{$\Spec(Y')=\ms{\bigl((s+1)t\bigr)^{1},(-t)^{s+1},(\sqrt c)^{\ell},
    (-\sqrt c)^{\ell}}$.}
   \lean{Computations.C11.spectrum_of_parameters}
   \begin{proof}
    \step{s558nbrem}{$\Spec(Y')=\ms{(s+1)t}+\sum_{\mu\in\Lambda'}\ms{\delta_\mu,-\delta_\mu}
     +\ms{t^{v-\rho},(-t)^{b-\rho}}$, where $\rho=\rk N$, $\delta_\mu=\sqrt{\mu+t^2}$, and
     $\Lambda'$ is the multiset of positive eigenvalues of $NN^T$ with one copy of $rk$ removed.}
    \begin{proof}\pf\ Remark~\ref{rem:5.7} with $X_1=I(D_1)$ and the bipartition $(P,\cB)$,
     biregular with degrees $(r,k)$, $r\ne k$, $v\le b$ by \asm[1]{0}, and biadjacency matrix $N$
     (Lemma~\ref{lem:incidence}), and $X_2=Y'$ by \stepref{s558Yp}, \stepref{s558cosp} and
     \stepref{s558nbfacts}; here $\frac12(r+k)=(s+1)t$ and $\frac12(r-k)=t$.~\qed\end{proof}
    \step{s558nbc}{$r-\lambda+t^2=c$, $v-1=\ell$ and $b-v=s+1$.}
    \lean{Computations.C11.nonprincipal_square}\lean{Computations.C11.spectral_parameters}
    \begin{proof}\pf\ Substitute \stepref{s558nbparams} and \stepref{s558nblam}
     (computation C11).~\qed\end{proof}
    \step{s558nbfish}{\case{$r>\lambda$.}}
    \begin{proof}
     \step{s558nbfishrho}{$\rho=v$ and $\Lambda'=\ms{(r-\lambda)^{v-1}}$.}
     \begin{proof}\pf\ By Lemma~\ref{lem:design}, $\Spec(NN^T)=\ms{(rk)^{1},(r-\lambda)^{v-1}}$
      with both values positive, so $\rk(NN^T)=v$ by \ST{Theorem 3.2}, and $\rho=\rk(NN^T)$ by
      \ST{Theorem 3.5}.~\qed\end{proof}
     \qedstep
     \begin{proof}\pf\ By \stepref{s558nbrem}, \stepref{s558nbfishrho} and \stepref{s558nbc},
      the terms are $(s+1)t$, $(\pm\sqrt c)^{v-1}$ with $v-1=\ell$, $t^{0}$ and
      $(-t)^{b-v}$ with $b-v=s+1$. (Computation C11,
      \path{Computations.C11.spectrum_of_parameters}, is this substitution; it takes the spectrum
      of \stepref{s558nbrem} in this form as a hypothesis, so the rank cases and
      Remark~\ref{rem:5.7} stay with this proof.)~\qed\end{proof}
    \end{proof}
    \step{s558nbtriv}{\case{$r=\lambda$.}}
    \begin{proof}
     \step{s558nbtrivrho}{$\rho=1$ and $\Lambda'$ is empty.}
     \begin{proof}\pf\ By Lemma~\ref{lem:design}, $\Spec(NN^T)=\ms{(rk)^{1},0^{v-1}}$ with
      $rk>0$, so $\rk(NN^T)=1$ by \ST{Theorem 3.2}, and $\rho=\rk(NN^T)$ by
      \ST{Theorem 3.5}.~\qed\end{proof}
     \step{s558nbtrivc}{$c=t^2$, so $\sqrt c=t$.}
     \begin{proof}\pf\ $c=r-\lambda+t^2$ by \stepref{s558nbc}, and $r=\lambda$.~\qed\end{proof}
     \qedstep
     \begin{proof}\pf\ By \stepref{s558nbrem} and \stepref{s558nbtrivrho},
      $\Spec(Y')=\ms{(s+1)t,t^{v-1},(-t)^{b-1}}$, and $v-1=\ell$, $b-1=(s+1)+\ell$ by
      \stepref{s558nbc}; with \stepref{s558nbtrivc} this is the claim.~\qed\end{proof}
    \end{proof}
    \qedstep
    \begin{proof}\pf\ $r\ge\lambda$ by Lemma~\ref{lem:design}, so the cases
     \stepref{s558nbfish} and \stepref{s558nbtriv} are exhaustive.~\qed\end{proof}
   \end{proof}
   \step{s558nbeven}{$t(s+1)$ is even.}
   \lean{Computations.C11.even_cube_diagonal}\lean{Computations.C11.even_parameter_cube}
   \lean{Computations.C11.even_mul_of_even_cube}
   \begin{proof}
    \pflet{$A=A(Y')$, $f(z)=(z+t)(z^2-c)$ and $e=t^2(s+1)(st-1)/(s-1)$.}
    \step{s558nbhp}{$f(A)=eJ$.}
    \lean{Computations.C11.matrix_identity}\lean{Computations.C11.hoffman_value}
    \begin{proof}\pf\ $f$ vanishes at $-t$ and $\pm\sqrt c$, which by \stepref{s558nbspec} are
     the eigenvalues other than $(s+1)t$; \ST{Theorem 7.10} with \stepref{s558nbY} gives
     $f(A)=\frac{f((s+1)t)}{s(s+2)}J$, and $f((s+1)t)=s(s+2)e$ (computation C11).~\qed\end{proof}
    \step{s558nbdiag}{$A^3(u,u)=t^3(s+1)$ for every vertex $u$.}
    \lean{Computations.C11.cube_diagonal_eq}\lean{Computations.C11.triangle_diagonal}
    \begin{proof}\pf\ The $(u,u)$ entry of $f(A)=A^3+tA^2-cA-ctI$ is
     $A^3(u,u)+t^2(s+1)-ct$, since $A(u,u)=0$ and $A^2(u,u)=\deg(u)=(s+1)t$ by
     \path{SimpleGraph.adjMatrix_mul_self_apply_self};
     by \stepref{s558nbhp} it equals $e$, and $e+ct-t^2(s+1)=t^3(s+1)$ (computation
     C11).~\qed\end{proof}
    \qedstep
    \begin{proof}\pf\ $Y'$ has a vertex, and $A^3(u,u)$ is twice a number of triangles by
     \ST{Theorem 5.1}; so $t^3(s+1)$ is even by \stepref{s558nbdiag}. The difference
     $t^3(s+1)-t(s+1)=(t-1)t(t+1)(s+1)$ is even, as a multiple of two consecutive
     integers. (In computation C11, \path{Computations.C11.even_parameter_cube} takes the
     integer form of \stepref{s558nbdiag} as a hypothesis; the passage from the real identity of
     \stepref{s558nbhp} to it stays with this proof.)~\qed\end{proof}
   \end{proof}
   \qedstep
   \begin{proof}\pf\ Take $Y=Y'$: $X_2\cong L(Y')$ by \stepref{s558Yp}; the conditions are
    \stepref{s558nbparams} to \stepref{s558nbeven}, with $t\ge1$ by \stepref{s558nbt}.~\qed\end{proof}
  \end{proof}
  \qedstep
  \begin{proof}\pf\ By \stepref{s558apply}, the cases \stepref{s558bip} and \stepref{s558nb} are
   the two alternatives of Lemma~\ref{lem:5.6}.~\qed\end{proof}
 \end{proof}
 \qedstep
 \begin{proof}\pf\ By \stepref{s558hrc}, the cases \stepref{s558E8}, \stepref{s558CP} and
  \stepref{s558line} are exhaustive.~\qed\end{proof}
\end{proof}

\medskip\noindent\textbf{Corollaries 5.9 to 5.11 and Remark 5.12.}

The next two lemmas settle the small cases that the corollaries leave to the reader.

\begin{namedresult}{Lemma}{(a single block)}\label{lem:oneblock}
 \lean{Design.flagGraphCharacterized_of_b_eq_one}
 \lean{SimpleGraph.Cospectral.nonempty_iso_of_isRegularOfDegree_card_sub_one}
 \uses{def:flag,lem:incidence,lem:dense,st:2.6,st:5.5}
 \prove{\begin{pfenum}
  \item If $X$ is an $(n-1)$-regular graph on $n\ge1$ vertices, then every graph cospectral
   with $X$ is isomorphic to $X$.
  \item If $D$ is a design with parameters $(v,v,1,1,1)$, then $F(D)$ is characterized by its
   spectrum.
 \end{pfenum}}
\end{namedresult}
\begin{proof}
 \pfsketch\ \repair\ The paper's Corollary 5.11 omits the designs with a single block (errata,
 Gap 16). A cospectral mate is regular of the same valency on the same number of vertices, so
 both graphs are complete.
 \step{s5obmate}{\assume{$X$ is an $(n-1)$-regular graph on $n\ge1$ vertices, and $X'$ is
  cospectral with $X$.}\prove{$X'\cong X$.}}
 \begin{proof}\pf\ $X'$ is $(n-1)$-regular on $n$ vertices by \ST{Corollary 5.5} and
  \ST{Proposition 2.6}; both graphs are isomorphic to $K_n$ by Lemma~\ref{lem:dense},
  assertion 1.~\qed\end{proof}
 \step{s5obflag}{\assume{$D$ is a design with parameters $(v,v,1,1,1)$.}
  \prove{$F(D)$ is characterized by its spectrum.}}
 \begin{proof}
  \step{s5obflagreg}{$F(D)$ is $(v-1)$-regular on $v$ vertices.}
  \begin{proof}\pf\ Lemma~\ref{lem:incidence}, assertion 3, with $r+k-2=v-1$ and
   $vr=v$.~\qed\end{proof}
  \qedstep
  \begin{proof}\pf\ Apply \stepref{s5obmate} to $F(D)$, by \stepref{s5obflagreg}, and take
   $D_2=D$ in Definition~\ref{def:flag}.~\qed\end{proof}
 \end{proof}
 \qedstep
 \begin{proof}\pf\ Steps \stepref{s5obmate} and \stepref{s5obflag}.~\qed\end{proof}
\end{proof}

\begin{namedresult}{Lemma}{(the hexagon)}\label{lem:hexagon}
 \lean{Design.flagGraphCharacterized_of_parameters_3_2_1}
 \uses{def:flag,lem:incidence,lem:cycle,st:2.6,st:5.5}
 \assume{$D$ is a design with parameters $(3,2,3,2,1)$.}
 \prove{$F(D)$ is characterized by its spectrum.}
\end{namedresult}
\begin{proof}
 \pfsketch\ \repair\ The paper leaves this to the reader in Corollaries 5.10 and 5.11 (errata,
 Gap 16): a cospectral mate of the hexagon is connected and $2$-regular on six vertices.
 \step{s5hxflag}{$F(D)\cong C_6$.}
 \begin{proof}\pf\ $F(D)$ is connected and $2$-regular on $6$ vertices by
  Lemma~\ref{lem:incidence}, assertion 3; apply Lemma~\ref{lem:cycle}.~\qed\end{proof}
 \step{s5hxmate}{\assume{$X_2$ is cospectral with $F(D)$.}\prove{$X_2\cong F(D)$.}}
 \begin{proof}\pf\ $X_2$ is connected and $2$-regular on $6$ vertices by \ST{Corollary 5.5} and
  \ST{Proposition 2.6}, so $X_2\cong C_6$ by Lemma~\ref{lem:cycle}; use
  \stepref{s5hxflag}.~\qed\end{proof}
 \qedstep
 \begin{proof}\pf\ Step \stepref{s5hxmate} with $D_2=D$ in Definition~\ref{def:flag}.~\qed\end{proof}
\end{proof}

\begin{namedresult}{Lemma}{(designs with full blocks)}\label{lem:fullblock}
 \lean{Design.trivial}\lean{Design.nonempty_iso_flagGraph_lineGraph_completeBipartite}
 \uses{def:flag,hc:def:design,lem:design,facts:sec5}
 \assume{$m\ge1$ and $n\ge2$ are integers.}
 \prove{\begin{pfenum}
  \item The \emph{trivial design} $T$ on the points $\{1,\dots,n\}$ with $m$ block labels, each
   block being the whole point set, is a design with parameters $(n,n,m,m,m)$.
  \item If $D$ is a design with parameters $(n,n,m,m,m)$, then $F(D)\cong L(K_{m,n})$.
 \end{pfenum}}
\end{namedresult}
\begin{proof}
 \pfsketch\ Every point lies in every block, so the incidence graph is complete bipartite.
 \step{s5fbT}{Assertion 1 holds.}
 \begin{proof}\pf\ By Definition~\ref{hc:def:design}: there are $n>1$ points and $m>0$ block
  labels, every block has $n$ points, and every point and every pair of points lies in all $m$
  blocks.~\qed\end{proof}
 \step{s5fbiso}{\assume{$D$ is a design with parameters $(n,n,m,m,m)$ on $P$ and
  $\cB$.}\prove{$F(D)\cong L(K_{m,n})$.}}
 \begin{proof}
  \step{s5fbfull}{Every block of $D$ contains every point.}
  \begin{proof}\pf\ Lemma~\ref{lem:design}, assertion 2, since $k=v=n$.~\qed\end{proof}
  \step{s5fbI}{$I(D)\cong K_{m,n}$.}
  \begin{proof}\pf\ By \stepref{s5fbfull}, every point is adjacent to every block in $I(D)$,
   so $I(D)$ is the complete bipartite graph on $P\sqcup\cB$ (Definition~\ref{def:flag}), with
   $|P|=n$ and $|\cB|=m$. We obtain an isomorphism with $K_{m,n}$ from \path{Equiv.sumComm},
   \path{Fintype.equivFin} and \path{SimpleGraph.completeBipartiteGraphCongr}.~\qed\end{proof}
  \qedstep
  \begin{proof}\pf\ $F(D)=L(I(D))$, and apply \path{SimpleGraph.Iso.lineGraph} to
   \stepref{s5fbI}.~\qed\end{proof}
 \end{proof}
 \qedstep
 \begin{proof}\pf\ Steps \stepref{s5fbT} and \stepref{s5fbiso}.~\qed\end{proof}
\end{proof}

\begin{cgss}{Corollary}{5.9}\label{cor:5.9}
 \lean{SimpleGraph.lineGraph_completeBipartite_characterized}
 \uses{thm:5.8,def:flag,hc:def:design,lem:design,lem:oneblock,lem:fullblock,facts:sec5,st:2.6,st:6.3,st:7.10}
 \assume{\begin{pfenum}
  \item $m,n\ge1$ are integers with $m+n>9$.
  \item There is no integer $t\ge1$ with $\{m,n\}=\{2t^2+t,2t^2-t\}$ for which a symmetric
   Hadamard matrix of order $4t^2$ with constant diagonal exists.
  \item $X_2$ is a graph cospectral with $L(K_{m,n})$.
 \end{pfenum}}
 \prove{$X_2\cong L(K_{m,n})$.}
\end{cgss}
\begin{proof}
 \pfsketch\ \repair\ The two readings of ``characterized'' agree here (errata, Implicit
 hypothesis 6), since a design with $k=v$ is unique. \repair\ We write the exception as a
 condition on $\{m,n\}$ (errata, Error 17): as printed, with $m=6$, $n=10$ the exception does not
 apply, but $L(K_{6,10})$ is cospectral with $L(T(6))$ (errata, Error 14, evidence). We will follow the paper: $L(K_{m,n})$
 is the flag graph of the trivial design, and in case (iii) of Theorem 5.8 bordering a
 $(\pm1)$ form of $A(Y)$ gives the Hadamard matrix.
 \step{s559ordered}{\assume{$m\ge n$, and \asm[1]{0} to \asm[3]{0}.}\prove{$X_2\cong
  L(K_{m,n})$.}}
 \begin{proof}
  \step{s559one}{\case{$n=1$.}}
  \begin{proof}\pf\ $L(K_{m,1})$ is $(m-1)$-regular on $m$ vertices by the degree formula of
   \ST{Lemma 6.3}; apply Lemma~\ref{lem:oneblock}.~\qed\end{proof}
  \step{s559two}{\case{$n\ge2$.}}
  \begin{proof}
   \pflet{$T$ be the trivial design of Lemma~\ref{lem:fullblock} on $n$ points with $m$ block
    labels.}
   \step{s559T}{$T$ is a design with parameters $(n,n,m,m,m)$, $b=m\ge n=v$, and
    $F(T)\cong L(K_{m,n})$.}
   \begin{proof}\pf\ Lemma~\ref{lem:fullblock}, both assertions, with $n\ge2$ by \asm{2} and
    $m\ge n$ by \asm{1}.~\qed\end{proof}
   \step{s559cases}{One of (i), (ii) and (iii) of Theorem~\ref{thm:5.8} holds for $D_1=T$ and
    $X_2$.}
   \begin{proof}\pf\ $X_2$ is cospectral with $F(T)$ by \asm[3]{0}, \stepref{s559T} and
    \ST{Proposition 2.6}; apply Theorem~\ref{thm:5.8}.~\qed\end{proof}
   \step{s559i}{\case{(i) holds.}}
   \begin{proof}\pf\ Then $n+m=v+b\le9$, against \asm[1]{0}.~\qed\end{proof}
   \step{s559ii}{\case{(ii) holds.}}
   \begin{proof}\pf\ Then $X_2\cong F(D_2)$ for a design $D_2$ with parameters $(n,n,m,m,m)$,
    and $F(D_2)\cong L(K_{m,n})$ by Lemma~\ref{lem:fullblock}, assertion 2.~\qed\end{proof}
   \step{s559iii}{\case{(iii) holds.}\prove{\textsc{False}.}}
   \begin{proof}
    \pfsketch\ We will show that the exception of \asm[2]{0} occurs.
    \step{s559st}{$s=2t-1$, $n=2t^2-t$, $m=2t^2+t$, $t\ge2$ and $c=t^2$.}
    \lean{Computations.C12.full_block_parameters}
    \begin{proof}\pf\ From $v=k$ in (iii), $\frac12s(s+1)=st$, so $s+1=2t$; substitute in
     $v$, $b$ and $c$, and $2t-1=s\ge2$ (computation C12). (Computation C12 proves the whole
     step, $t\ge2$ and $c=t^2$ included, for the integers $s\ge2$ and $t\ge1$ of (iii), with the
     fractions evaluated over $\rats$.)~\qed\end{proof}
    \step{s559Y}{$A=A(Y)$ satisfies $A^2=t^2(I+J)$ and $AJ=2t^2J$, on $4t^2-1$ vertices.}
    \begin{proof}\pf\ By (iii) and \stepref{s559st}, $Y$ is connected and $2t^2$-regular on
     $s(s+2)=4t^2-1$ vertices with eigenvalues $2t^2$ and $\pm t$. \ST{Theorem 7.10} with
     $f(z)=z^2-t^2$ gives $A^2-t^2I=\frac{4t^4-t^2}{4t^2-1}J=t^2J$.~\qed\end{proof}
    \pflet{$S=2A-J$ and $H=\pmat{-1&-\Zj^T\\-\Zj&S}$, of order $4t^2$.}
    \lean{Computations.C12.core}\lean{Computations.C12.border}
    \step{s559S}{$S$ is symmetric with entries $\pm1$ and diagonal $-1$, $S\Zj=\Zj$ and
     $S^2=4t^2I-J$.}
    \lean{Computations.C12.core_row_sum}\lean{Computations.C12.core_mul}
    \begin{proof}\pf\ $A$ is a symmetric zero-one matrix with zero diagonal; by \stepref{s559Y},
     $S\Zj=(4t^2-(4t^2-1))\Zj$ and $S^2=4A^2-2AJ-2JA+J^2=4t^2I-J$, since $J^2=(4t^2-1)J$
     (computation C12).~\qed\end{proof}
    \step{s559H}{$H$ is a symmetric Hadamard matrix with constant diagonal $-1$.}
    \lean{Computations.C12.border_isHadamard}\lean{Computations.C12.hadamard_border}
    \begin{proof}\pf\ By \stepref{s559S}, $H$ is symmetric with entries $\pm1$ and diagonal
     $-1$, and block multiplication gives $H^2=4t^2I$: the corner is $1+(4t^2-1)$, the border
     is $\Zj-S\Zj=0$, and the rest is $J+S^2$. This is \path{Matrix.IsHadamard} over
     $\re$.~\qed\end{proof}
    \qedstep
    \begin{proof}\pf\ \stepref{s559st} and \stepref{s559H} contradict \asm[2]{0}.~\qed\end{proof}
   \end{proof}
   \qedstep
   \begin{proof}\pf\ By \stepref{s559cases}, the cases \stepref{s559i}, \stepref{s559ii} and
    \stepref{s559iii} are exhaustive, and \stepref{s559i} and \stepref{s559iii} are
    impossible.~\qed\end{proof}
  \end{proof}
  \qedstep
  \begin{proof}\pf\ The cases \stepref{s559one} and \stepref{s559two} are exhaustive, since
   $n\ge1$.~\qed\end{proof}
 \end{proof}
 \step{s559ge}{\case{$m\ge n$.}}
 \begin{proof}\pf\ Step \stepref{s559ordered}.~\qed\end{proof}
 \step{s559lt}{\case{$m<n$.}}
 \begin{proof}
  \step{s559swap}{$L(K_{m,n})\cong L(K_{n,m})$.}
  \begin{proof}\pf\ \path{Equiv.sumComm} is an isomorphism $K_{m,n}\cong K_{n,m}$ by the
   definition of the complete bipartite graph; apply \path{SimpleGraph.Iso.lineGraph}.~\qed\end{proof}
  \qedstep
  \begin{proof}\pf\ \asm[1]{0} and \asm[2]{0} are symmetric in $m$ and $n$, and $X_2$ is
   cospectral with $L(K_{n,m})$ by \stepref{s559swap} and \ST{Proposition 2.6}; apply
   \stepref{s559ordered} to $(n,m)$, and then \stepref{s559swap}.~\qed\end{proof}
 \end{proof}
 \qedstep
 \begin{proof}\pf\ The cases \stepref{s559ge} and \stepref{s559lt} are exhaustive.~\qed\end{proof}
\end{proof}

\begin{cgss}{Corollary}{5.10}\label{cor:5.10}
 \lean{Design.IsSymmetric.flagGraphCharacterized}
 \uses{thm:5.8,def:flag,hc:def:design,lem:design,lem:hexagon}
 \assume{\begin{pfenum}
  \item $D_1$ is a design with parameters $(v,k,v,r,\lambda)$, that is, a symmetric design.
  \item $k<v$ and $(v,k,\lambda)\ne(4,3,2)$.
 \end{pfenum}}
 \prove{$F(D_1)$ is characterized by its spectrum.}
\end{cgss}
\begin{proof}
 \pfsketch\ We will follow the paper: case (iii) of Theorem 5.8 cannot occur, and case (i)
 leaves only $(3,2,1)$, the hexagon, since $(4,3,2)$ is excluded.
 \step{s5510rk}{$r=k$ and $b\ge v$.}
 \begin{proof}\pf\ $vk=vr$ by Lemma~\ref{lem:design} with $b=v$.~\qed\end{proof}
 \step{s5510mate}{\assume{$X_2$ is cospectral with $F(D_1)$.}
  \prove{$X_2\cong F(D_2)$ for a design $D_2$ with the parameters of $D_1$.}}
 \begin{proof}
  \step{s5510i}{\case{(i) of Theorem~\ref{thm:5.8} holds.}}
  \begin{proof}
   \step{s5510small}{$(v,k,\lambda)=(3,2,1)$.}
   \begin{proof}
    \step{s5510v}{$v\in\{3,4\}$.}
    \begin{proof}\pf\ $2v=v+b\le9$, and $v>k>1$.~\qed\end{proof}
    \step{s5510three}{\case{$v=3$.}}
    \begin{proof}\pf\ $k=2$, and $2\lambda=\lambda(v-1)=r(k-1)=2$ by Lemma~\ref{lem:design}
     and \stepref{s5510rk}, so $\lambda=1$.~\qed\end{proof}
    \step{s5510four}{\case{$v=4$.}\prove{\textsc{False}.}}
    \begin{proof}
     \step{s5510fourlam}{$3\lambda=k(k-1)$ and $k\in\{2,3\}$.}
     \begin{proof}\pf\ $\lambda(v-1)=r(k-1)$ by Lemma~\ref{lem:design}, with $r=k$ by
      \stepref{s5510rk}; and $1<k<v=4$.~\qed\end{proof}
     \step{s5510fourtwo}{\case{$k=2$.}}
     \begin{proof}\pf\ Then $3\lambda=2$ by \stepref{s5510fourlam}, which has no integer
      solution.~\qed\end{proof}
     \step{s5510fourthree}{\case{$k=3$.}}
     \begin{proof}\pf\ Then $\lambda=2$ by \stepref{s5510fourlam}, and $(v,k,\lambda)=(4,3,2)$
      is excluded by \asm[2]{0}.~\qed\end{proof}
     \qedstep
     \begin{proof}\pf\ The cases \stepref{s5510fourtwo} and \stepref{s5510fourthree} are
      exhaustive by \stepref{s5510fourlam}.~\qed\end{proof}
    \end{proof}
    \qedstep
    \begin{proof}\pf\ The cases \stepref{s5510three} and \stepref{s5510four} are exhaustive by
     \stepref{s5510v}.~\qed\end{proof}
   \end{proof}
   \qedstep
   \begin{proof}\pf\ Lemma~\ref{lem:hexagon} with \stepref{s5510small} and $b=v=3$,
    $r=k=2$.~\qed\end{proof}
  \end{proof}
  \step{s5510ii}{\case{(ii) of Theorem~\ref{thm:5.8} holds.}}
  \begin{proof}\pf\ This is the claim.~\qed\end{proof}
  \step{s5510iii}{\case{(iii) of Theorem~\ref{thm:5.8} holds.}\prove{\textsc{False}.}}
  \begin{proof}\pf\ $r=(s+2)t>st=k$ since $t\ge1$, against \stepref{s5510rk}.~\qed\end{proof}
  \qedstep
  \begin{proof}\pf\ Theorem~\ref{thm:5.8} applies by \asm[1]{0}, \stepref{s5510rk} and
   \asm{1}, so the cases are exhaustive.~\qed\end{proof}
 \end{proof}
 \qedstep
 \begin{proof}\pf\ Step \stepref{s5510mate} is Definition~\ref{def:flag}.~\qed\end{proof}
\end{proof}

\begin{remark}
 \lean{Computations.C13.graph1}\lean{Computations.C13.graph2}\lean{Computations.C13.spectrum_graph1}
 \lean{Computations.C13.spectrum_graph2}
 \lean{Computations.C13.cospectral}\lean{Computations.C13.not_nonempty_iso}
 \lean{Computations.C13.design}\lean{Computations.C13.nonempty_iso_flagGraph}
 \lean{Computations.C13.representation_graph1}\lean{Computations.C13.representation_graph2}
 \lean{Computations.C13.isRepresentedIn_E₇_graph1}\lean{Computations.C13.isRepresentedIn_E₇_graph2}
 \lean{Computations.C13.cube}\lean{Computations.C13.cubeLineGraphIso}
 \lean{Computations.C13.nonempty_iso_cube_lineGraph}
 \lean{Computations.C13.not_flagGraphCharacterized}
 The exception is genuine. The paper's two sets of 12 vectors in $\typ E_7$ give two graphs with
 spectrum $4^1\,2^3\,0^3\,(-2)^5$: the flag graph of the $(4,3,2)$ design, which is
 $L(Q_3)$, and a second graph, which computation C13 shows is not isomorphic to the first
 (\path{cgss/prerequisites.md}, HR1). A design with parameters $(4,3,4,3,2)$ is unique: each of
 its four blocks omits one point, each point is omitted by $b-r=1$ block, and so the blocks are
 the four $3$-subsets of the points. Every flag graph of such a design is therefore isomorphic to
 the first graph, and the second graph shows, in the sense of Definition~\ref{def:flag}, that the
 flag graph is not characterized by its spectrum. The uniqueness of the exception (HR2) is a side
 claim, and we do not formalize it here. Computation C13 checks the spectra, that the two graphs are
 represented by the paper's vectors in $\typ E_7$, and that they are not isomorphic (the first
 has no induced claw and the second has one); it identifies the first graph with the line graph
 of the incidence graph of the design on the four $3$-subsets, built in
 \path{LineGraphs/Computations/C13Design.lean} rather than through Definition~\ref{def:flag}.
 It also identifies the first graph with $L(Q_3)$, by an explicit isomorphism, where $Q_3$ has
 the eight binary triples as vertices, two of them adjacent when they differ in exactly one
 coordinate; we do not use this identification. The non-characterization itself is
 \path{Computations.C13.not_flagGraphCharacterized}: the second graph is cospectral with the
 flag graph of the design and has an induced claw, while the flag graph of every design, being a
 line graph, has none.
\end{remark}

\begin{cgss}{Corollary}{5.11}\label{cor:5.11}
 \lean{Design.flagGraphCharacterized_of_lam_eq_one}
 \uses{thm:5.8,def:flag,hc:def:design,lem:design,lem:oneblock,lem:hexagon}
 \assume{$D_1$ is a design with parameters $(v,k,b,r,1)$.}
 \prove{$F(D_1)$ is characterized by its spectrum.}
\end{cgss}
\begin{proof}
 \pfsketch\ \repair\ The paper says that case (i) of Theorem 5.8 occurs only for $(3,2,1)$; it
 also occurs for a single block containing every point, $(v,v,1,1,1)$ with $2\le v\le8$, whose
 flag graph is complete (errata, Gap 16). We will treat that case first, and then follow the
 paper.
 \step{s5511full}{\case{$k=v$.}}
 \begin{proof}\pf\ Then $r=\lambda=1$ by Lemma~\ref{lem:design}, and $b=vr/k=1$, so $D_1$ has
  parameters $(v,v,1,1,1)$; apply Lemma~\ref{lem:oneblock}.~\qed\end{proof}
 \step{s5511proper}{\case{$k<v$.}}
 \begin{proof}
  \step{s5511fisher}{$r>1$ and $b\ge v$.}
  \begin{proof}\pf\ Lemma~\ref{lem:design}: $r\ge\lambda=1$ with equality only if $k=v$, and
   then $b\ge v$ since $r>\lambda$.~\qed\end{proof}
  \step{s5511mate}{\assume{$X_2$ is cospectral with $F(D_1)$.}
   \prove{$X_2\cong F(D_2)$ for a design $D_2$ with the parameters of $D_1$.}}
  \begin{proof}
   \step{s5511i}{\case{(i) of Theorem~\ref{thm:5.8} holds.}}
   \begin{proof}
    \step{s5511small}{$(v,k,b,r)=(3,2,3,2)$.}
    \begin{proof}
     \step{s5511v}{$v\in\{3,4\}$.}
     \begin{proof}\pf\ $2v\le v+b\le9$ by \stepref{s5511fisher}, and $v>k>1$.~\qed\end{proof}
     \step{s5511three}{\case{$v=3$.}}
     \begin{proof}\pf\ $k=2$, $r=\lambda(v-1)/(k-1)=2$ and $b=vr/k=3$ by
      Lemma~\ref{lem:design}.~\qed\end{proof}
     \step{s5511four}{\case{$v=4$.}\prove{\textsc{False}.}}
     \begin{proof}
      \step{s5511fourr}{$r(k-1)=3$ and $k\in\{2,3\}$.}
      \begin{proof}\pf\ $r(k-1)=\lambda(v-1)$ by Lemma~\ref{lem:design}, with $\lambda=1$ and
       $1<k<v=4$.~\qed\end{proof}
      \step{s5511fourtwo}{\case{$k=2$.}}
      \begin{proof}\pf\ Then $r=3$ by \stepref{s5511fourr}, and $b=vr/k=6$ by
       Lemma~\ref{lem:design}, so $v+b=10>9$, against (i).~\qed\end{proof}
      \step{s5511fourthree}{\case{$k=3$.}}
      \begin{proof}\pf\ Then $2r=3$ by \stepref{s5511fourr}, which has no integer
       solution.~\qed\end{proof}
      \qedstep
      \begin{proof}\pf\ The cases \stepref{s5511fourtwo} and \stepref{s5511fourthree} are
       exhaustive by \stepref{s5511fourr}.~\qed\end{proof}
     \end{proof}
     \qedstep
     \begin{proof}\pf\ The cases \stepref{s5511three} and \stepref{s5511four} are exhaustive by
      \stepref{s5511v}.~\qed\end{proof}
    \end{proof}
    \qedstep
    \begin{proof}\pf\ Lemma~\ref{lem:hexagon} with \stepref{s5511small}.~\qed\end{proof}
   \end{proof}
   \step{s5511ii}{\case{(ii) of Theorem~\ref{thm:5.8} holds.}}
   \begin{proof}\pf\ This is the claim.~\qed\end{proof}
   \step{s5511iii}{\case{(iii) of Theorem~\ref{thm:5.8} holds.}\prove{\textsc{False}.}}
   \begin{proof}\pf\ Then $\lambda(s-1)=2t(st-1)\ge2(s-1)$, since $t\ge1$ and $st-1\ge s-1$, so
    $\lambda\ge2$, as $s\ge2$; this contradicts $\lambda=1$.~\qed\end{proof}
   \qedstep
   \begin{proof}\pf\ Theorem~\ref{thm:5.8} applies by \stepref{s5511fisher} and \asm{1}, so the
    cases are exhaustive.~\qed\end{proof}
  \end{proof}
  \qedstep
  \begin{proof}\pf\ Step \stepref{s5511mate} is Definition~\ref{def:flag}.~\qed\end{proof}
 \end{proof}
 \qedstep
 \begin{proof}\pf\ $k\le v$ by Lemma~\ref{lem:design}, so the cases \stepref{s5511full} and
  \stepref{s5511proper} are exhaustive.~\qed\end{proof}
\end{proof}

\begin{cgss}{Remark}{5.12}\label{rem:5.12}
 \lean{Design.flagGraphCharacterized_of_lam_lt_nineteen}
 \uses{thm:5.8,def:flag,hc:def:design,lem:design,hc:cor:complement}
 \assume{$D_1$ is a design with parameters $(v,k,b,r,\lambda)$, with $k<v$ and $2<\lambda<19$.}
 \prove{$F(D_1)$ is characterized by its spectrum.}
\end{cgss}
\begin{proof}
 \pfsketch\ We will follow the paper. In case (iii) of Theorem 5.8,
 $\lambda=2t^2+2t(t-1)/(s-1)$, so $t=1$ gives $\lambda=2$, $t\ge3$ gives $\lambda\ge19$, and
 $t=2$ gives $s\in\{3,5\}$; $s=3$ is a design with $k=v$, and $s=5$ gives the parameters
 $(15,10,21,14,9)$, whose complement $(15,5,21,7,2)$ does not exist. The paper cites Hall's book
 for this; we use the Hall--Connor section. \new\ The paper does not mention case (i); it cannot
 occur for $\lambda\ge3$.
 \step{s5512fisher}{$r>\lambda$ and $b\ge v$.}
 \begin{proof}\pf\ Lemma~\ref{lem:design}, since $k<v$.~\qed\end{proof}
 \step{s5512mate}{\assume{$X_2$ is cospectral with $F(D_1)$.}
  \prove{$X_2\cong F(D_2)$ for a design $D_2$ with the parameters of $D_1$.}}
 \begin{proof}
  \step{s5512i}{\case{(i) of Theorem~\ref{thm:5.8} holds.}\prove{\textsc{False}.}}
  \begin{proof}
   \step{s5512b}{$v\in\{3,4\}$ and $b=\lambda v(v-1)/(k(k-1))$.}
   \begin{proof}\pf\ $2v\le v+b\le9$ by \stepref{s5512fisher}, and $v>k>1$; the formula is
    Lemma~\ref{lem:design}.~\qed\end{proof}
   \step{s5512three}{\case{$v=3$.}}
   \begin{proof}\pf\ $k=2$, so $b=3\lambda\ge9$ and $v+b\ge12$.~\qed\end{proof}
   \step{s5512four}{\case{$v=4$.}}
   \begin{proof}\pf\ $k\in\{2,3\}$, so $b\ge12\lambda/6=2\lambda\ge6$ and
    $v+b\ge10$.~\qed\end{proof}
   \qedstep
   \begin{proof}\pf\ The cases \stepref{s5512three} and \stepref{s5512four} are exhaustive by
    \stepref{s5512b}, and both contradict $v+b\le9$.~\qed\end{proof}
  \end{proof}
  \step{s5512ii}{\case{(ii) of Theorem~\ref{thm:5.8} holds.}}
  \begin{proof}\pf\ This is the claim.~\qed\end{proof}
  \step{s5512iii}{\case{(iii) of Theorem~\ref{thm:5.8} holds.}\prove{\textsc{False}.}}
  \begin{proof}
   \step{s5512lam}{$\lambda(s-1)=2t^2(s-1)+2t(t-1)$.}
   \lean{Computations.C14.lambda_remainder}
   \begin{proof}\pf\ $2t(st-1)=2t^2(s-1)+2t(t-1)$, and (iii) (computation C14).~\qed\end{proof}
   \step{s5512t1}{\case{$t=1$.}}
   \lean{Computations.C14.lambda_eq_two}
   \begin{proof}\pf\ Then $\lambda=2$ by \stepref{s5512lam}, against $\lambda>2$.~\qed\end{proof}
   \step{s5512t3}{\case{$t\ge3$.}}
   \lean{Computations.C14.lambda_ge_nineteen}
   \begin{proof}\pf\ Then $2t(t-1)>0$, so $\lambda>2t^2\ge18$ by \stepref{s5512lam}, against
    $\lambda<19$.~\qed\end{proof}
   \step{s5512t2}{\case{$t=2$.}}
   \lean{Computations.C14.t2_parameters}
   \begin{proof}
    \step{s5512s}{$s\in\{3,5\}$.}
    \lean{Computations.C14.eq_three_or_five}
    \begin{proof}\pf\ By (iii), $(s-1)\mid4$ and $s\ge2t-1=3$ (computation C14).~\qed\end{proof}
    \step{s5512s3}{\case{$s=3$.}}
    \begin{proof}\pf\ Then $v=6=k$ by (iii), against $k<v$.~\qed\end{proof}
    \step{s5512s5}{\case{$s=5$.}}
    \begin{proof}\pf\ Then $(v,k,b,r,\lambda)=(15,10,21,14,9)$ by (iii), and
     Corollary~\hyperref[hc:cor:complement]{(the complementary parameters)} of the Hall--Connor
     section gives \textsc{False}.~\qed\end{proof}
    \qedstep
    \begin{proof}\pf\ The cases \stepref{s5512s3} and \stepref{s5512s5} are exhaustive by
     \stepref{s5512s}.~\qed\end{proof}
   \end{proof}
   \qedstep
   \begin{proof}\pf\ $t\ge1$ by (iii), so the cases \stepref{s5512t1}, \stepref{s5512t3} and
    \stepref{s5512t2} are exhaustive.~\qed\end{proof}
  \end{proof}
  \qedstep
  \begin{proof}\pf\ Theorem~\ref{thm:5.8} applies by \stepref{s5512fisher} and \asm{1}, so the
   cases are exhaustive.~\qed\end{proof}
 \end{proof}
 \qedstep
 \begin{proof}\pf\ Step \stepref{s5512mate} is Definition~\ref{def:flag}.~\qed\end{proof}
\end{proof}

\begin{remark}
 \lean{Computations.TrivialDesignException.design}
 \lean{Computations.TrivialDesignException.incidenceSignless}
 \lean{Computations.TrivialDesignException.completeAdjacency}
 \lean{Computations.TrivialDesignException.incidence_spectrum_checked}
 \lean{Computations.TrivialDesignException.complete_spectrum_checked}
 \lean{Computations.TrivialDesignException.incidenceIndex}
 \lean{Computations.TrivialDesignException.incidence_signless_eq}
 \lean{Computations.TrivialDesignException.spectrum_incidence_signless}
 \lean{Computations.TrivialDesignException.spectrum_complete}
 \lean{Computations.TrivialDesignException.spectrum_triangular}
 \lean{Computations.TrivialDesignException.spectrum_flagGraph}
 \lean{Computations.TrivialDesignException.spectrum_double_lineGraph}
 \lean{Computations.TrivialDesignException.flagGraph_cospectral_double_lineGraph}
 \lean{Computations.TrivialDesignException.lineGraph_completeBipartite_cospectral_double_lineGraph}
 \lean{Computations.TrivialDesignException.flagGraph_not_nonempty_iso_double_lineGraph}
 \lean{Computations.TrivialDesignException.not_flagGraphCharacterized}
 \lean{SimpleGraph.spectrumMultiset_lineGraph_of_card_le}
 The paper adds that for $t=2$ and $s=3$ the parameters are those of the trivial design with
 $(v,k,b,r)=(6,6,10,10)$, and that $L(L(K_6))$ provides an exception. The exception is genuine,
 and we prove it in \path{LineGraphs/Computations/TrivialDesignException.lean}. Let $D_1$ be the
 design $(6,6,10,10,10)$, whose ten blocks all equal its set of six points
 (\path{Computations.TrivialDesignException.design}); its incidence graph is $K_{10,6}$, so
 $F(D_1)\cong L(K_{10,6})$. For a graph $Y$ with $n$ vertices and $m\ge n$ edges, the spectrum of
 $L(Y)$ is the spectrum of the signless Laplacian of $Y$ shifted by $-2$, together with the
 eigenvalue $-2$ of multiplicity $m-n$
 (\path{SimpleGraph.spectrumMultiset_lineGraph_of_card_le}). Certificates for the signless
 Laplacian of $K_{10,6}$, of order 16, and for $A(K_6)$, of order 6, give through this identity
 the spectrum $14^1\,8^5\,4^9\,(-2)^{45}$ for both $F(D_1)$ and $L(L(K_6))=L(T(6))$
 (\path{Computations.TrivialDesignException.flagGraph_cospectral_double_lineGraph}). Thus
 $L(K_{10,6})$ and $L(T(6))$ are cospectral, the pair behind Errors 14, 16 and 17 (errata). They
 are not isomorphic, and $F(D_1)$ is not characterized by its spectrum in the sense of
 Definition~\ref{def:flag}
 (\path{Computations.TrivialDesignException.not_flagGraphCharacterized}): the flag graph of a
 design with replication number 10 has a clique of 10 vertices, the flags through one point,
 while a clique of $L(T(6))$ consists of edges of $T(6)$ through one vertex or of the three edges
 of a triangle, and so has at most 8 vertices, since $T(6)$ is 8-regular.
\end{remark}

\medskip\noindent\textbf{Example 5.13.}

\begin{cgss}{Example}{5.13}\label{ex:5.13}
 \lean{Design.not_flagGraphCharacterized_of_resolvable}
 \uses{def:flag,hc:def:design,hc:def:matrices,lem:design,lem:incidence,lem:rootgraph,lem:5.1,lem:5.3,thm:5.8,input:ex513,facts:sec5,st:2.2,st:2.6,st:3.1,st:3.2,st:3.5,st:6.1,st:6.2,st:6.6}
 \assume{\begin{pfenum}
  \item $k\ge2$ is an integer, and $D'$ is a design with parameters
   $(k(2k-1),k,4k^2-1,2k+1,1)$ on $P'$ and $\cB'$, with incidence matrix $N'$.
  \item $\pi:\cB'\to\{1,\dots,2k+1\}$ is a \emph{resolution} of $D'$: for each $i$, every point
   lies in exactly one block $B$ with $\pi(B)=i$.
  \item $H$ is the graph on $\cB'$ in which $B\ne B'$ are adjacent when $B$ and $B'$ are disjoint
   and $\pi(B)\ne\pi(B')$.
  \item $D_1$ is a design with parameters $(v_1,k_1,b_1,r_1,\lambda_1)=
   (k(2k-1),(k-1)(2k-1),k(2k+1),(k-1)(2k+1),k(2k-3))$.
 \end{pfenum}}
 \prove{\begin{pfenum}
  \item $H$ is connected and $\kappa$-regular on $4k^2-1$ vertices, where $\kappa=2k(k-1)$, and
   $\Spec(H)=\ms{\kappa^{1},(-(k-1))^{2k},k^{\ell},(-k)^{\ell}}$, where $\ell=(k-1)(2k+1)$.
   This is the spectrum of case (iii) of Theorem~\ref{thm:5.8} with $s=2k-1$ and $t=k-1$.
  \item $L(H)$ is cospectral with $F(D_1)$.
  \item $F(D_1)$ is not characterized by its spectrum.
 \end{pfenum}}
\end{cgss}
\begin{proof}
 \pfsketch\ \repair\ The paper computes the spectrum of $H$ and leaves the rest to the reader;
 ``not characterized'' is meant in the sense of Definition~\ref{def:flag} (errata, Implicit
 hypothesis 6), and then it is EX4. We will split $\re^{\cB'}$ into four eigenspaces of $A(H)$,
 compare the signless Laplacians of $H$ and of $I(D_1)$ as in Lemma 5.5, and then apply EX4. With
 $s=2k-1$, $t=k-1$ we have $(s+1)t=\kappa$, $s(s+2)=4k^2-1$, $c=k^2$ and
 $\frac12(s-1)(s+2)=\ell$ in Theorem 5.8 (iii).
 \pflet{$c_i\in\re^{\cB'}$ be the indicator vector of $\pi^{-1}(i)$ for $1\le i\le2k+1$, $R$ the
  $\cB'\times\cB'$ matrix with $R(B,B')=1$ when $B\ne B'$ and $\pi(B)=\pi(B')$ and $0$ otherwise,
  and $J$ the all-ones $\cB'\times\cB'$ matrix.}
 \step{s5exclass}{$N'c_i=\Zj$, each $\pi^{-1}(i)$ has $2k-1$ blocks, $\sum_ic_i=\Zj$ and
  $R=\sum_ic_ic_i^T-I$.}
 \begin{proof}\pf\ By \asm[2]{0}, the blocks of $\pi^{-1}(i)$ partition $P'$; each has $k$
  points and $|P'|=k(2k-1)$. Every block lies in exactly one class.~\qed\end{proof}
 \step{s5exadj}{$A(H)=J-I-R-(N'^TN'-kI)$.}
 \begin{proof}
  \step{s5exadjd}{$(N'^TN')(B,B')=|D'(B)\cap D'(B')|$, and $(N'^TN')(B,B)=k$.}
  \begin{proof}\pf\ The entry is the number of points on both blocks, and every block has $k$
   points by \asm[1]{0}.~\qed\end{proof}
  \step{s5exadj01}{$(N'^TN')(B,B')\in\{0,1\}$ for $B\ne B'$.}
  \begin{proof}\pf\ Two distinct points lie in exactly $\lambda=1$ common block, by \asm[1]{0},
   so distinct blocks share at most one point; apply \stepref{s5exadjd}.~\qed\end{proof}
  \step{s5exadjc}{$(N'^TN')(B,B')=0$ for $B\ne B'$ with $\pi(B)=\pi(B')$.}
  \begin{proof}\pf\ The blocks of $\pi^{-1}(\pi(B))$ are disjoint by \stepref{s5exclass}; apply
   \stepref{s5exadjd}.~\qed\end{proof}
  \qedstep
  \begin{proof}\pf\ By Steps \stepref{s5exadjd} to \stepref{s5exadjc}, the matrix
   $J-I-R-(N'^TN'-kI)$ has $(B,B')$ entry $0$ when $B=B'$ or $\pi(B)=\pi(B')$, and
   $1-|D'(B)\cap D'(B')|$ otherwise, which is $1$ exactly when $B$ and $B'$ are disjoint. Compare
   the entries with \asm[3]{0}.~\qed\end{proof}
 \end{proof}
 \step{s5exgram}{$N'N'^T=2kI+J$, and $\rk N'=k(2k-1)$.}
 \begin{proof}\pf\ Lemma~\ref{lem:design}, assertion 3, with $r-\lambda=2k$; its spectrum is
  positive, so $\rk(N'N'^T)=k(2k-1)$ by \ST{Theorem 3.2}, and $\rk N'=\rk(N'N'^T)$ by
  \ST{Theorem 3.5}.~\qed\end{proof}
 \pflet{$W_1=\{\sum_ia_ic_i:\sum_ia_i=0\}$, $W_2=\{N'^T\Zy:\Zy^T\Zj=0\}$ and
  $W_3=\{\Zx:N'\Zx=0,\ c_i^T\Zx=0\text{ for all }i\}$.}
 \step{s5exvalues}{$A(H)\Zj=\kappa\Zj$, and $A(H)$ acts on $W_1$, $W_2$, $W_3$ as $-(k-1)$, $-k$
  and $k$.}
 \begin{proof}
  \step{s5exvj}{$A(H)\Zj=\kappa\Zj$.}
  \begin{proof}\pf\ $J\Zj=(4k^2-1)\Zj$, $R\Zj=(2k-2)\Zj$ by \stepref{s5exclass}, and
   $N'^TN'\Zj=(2k+1)k\Zj$ from the row and column sums of $N'$; by \stepref{s5exadj},
   $A(H)\Zj=(4k^2-2-(2k-2)-(2k+1)k+k)\Zj=\kappa\Zj$.~\qed\end{proof}
  \step{s5exv1}{$A(H)\Zx=-(k-1)\Zx$ for $\Zx\in W_1$.}
  \begin{proof}\pf\ By \stepref{s5exclass}, $J\Zx=(2k-1)(\sum_ia_i)\Zj=0$,
   $R\Zx=(2k-2)\Zx$ and $N'\Zx=(\sum_ia_i)\Zj=0$; apply \stepref{s5exadj}.~\qed\end{proof}
  \step{s5exv2}{$A(H)\Zx=-k\Zx$ for $\Zx\in W_2$.}
  \begin{proof}\pf\ For $\Zx=N'^T\Zy$ with $\Zy^T\Zj=0$, $c_i^T\Zx=\Zj^T\Zy=0$ by
   \stepref{s5exclass}, so $J\Zx=0$ and $R\Zx=-\Zx$; and $N'^TN'\Zx=N'^T(2kI+J)\Zy=2k\Zx$ by
   \stepref{s5exgram}. Apply \stepref{s5exadj}.~\qed\end{proof}
  \step{s5exv3}{$A(H)\Zx=k\Zx$ for $\Zx\in W_3$.}
  \begin{proof}\pf\ $J\Zx=0$ and $R\Zx=-\Zx$ by \stepref{s5exclass}, and $N'^TN'\Zx=0$; apply
   \stepref{s5exadj}.~\qed\end{proof}
  \qedstep
  \begin{proof}\pf\ Steps \stepref{s5exvj} to \stepref{s5exv3}.~\qed\end{proof}
 \end{proof}
 \step{s5exdims}{$\dim W_1=2k$ and $\dim W_2=\dim W_3=\ell$.}
 \begin{proof}
  \step{s5exd12}{$\dim W_1=2k$ and $\dim W_2=\ell$.}
  \begin{proof}\pf\ By \stepref{s5exclass}, the $2k+1$ vectors $c_i$ have disjoint nonempty
   supports, so they are independent, and $W_1$ is a hyperplane of their span. $N'^T$ is
   injective by \stepref{s5exgram}, so $\dim W_2=k(2k-1)-1=\ell$.~\qed\end{proof}
  \step{s5exd3}{$\dim W_3=\ell$.}
  \begin{proof}
   \step{s5exd3k}{$\dim\ker N'=(4k^2-1)-k(2k-1)$.}
   \begin{proof}\pf\ $N'$ has $4k^2-1$ columns and rank $k(2k-1)$ by
    \stepref{s5exgram}.~\qed\end{proof}
   \step{s5exd3w}{$W_1\sbs\ker N'$.}
   \begin{proof}\pf\ $N'\sum_ia_ic_i=(\sum_ia_i)\Zj$ by \stepref{s5exclass}.~\qed\end{proof}
   \step{s5exd3p}{$W_3$ is the orthogonal complement of $W_1$ in $\ker N'$.}
   \begin{proof}\pf\ A vector $\Zx$ is orthogonal to $W_1$ exactly when all $c_i^T\Zx$ are
    equal, so $W_3$ lies in the complement. Conversely, if $\Zx\in\ker N'$ and $\Zx\perp W_1$,
    then all $c_i^T\Zx$ are equal, and their sum is $\Zj^T\Zx=\frac1k\Zj^TN'\Zx=0$ by
    \stepref{s5exclass}, since the column sums of $N'$ are $k$; so $\Zx\in
    W_3$.~\qed\end{proof}
   \qedstep
   \begin{proof}\pf\ By Steps \stepref{s5exd3k} to \stepref{s5exd3p} and \stepref{s5exd12},
    $\dim W_3=(4k^2-1)-k(2k-1)-2k=\ell$.~\qed\end{proof}
  \end{proof}
  \qedstep
  \begin{proof}\pf\ Steps \stepref{s5exd12} and \stepref{s5exd3}.~\qed\end{proof}
 \end{proof}
 \step{s5exspec}{Assertion 1 holds.}
 \begin{proof}
  \step{s5exspecH}{$\Spec(H)=\ms{\kappa^{1},(-(k-1))^{2k},k^{\ell},(-k)^{\ell}}$.}
  \begin{proof}\pf\ The values $\kappa$, $-(k-1)$, $-k$, $k$ are distinct since $k\ge2$, so by
   \ST{Theorem 2.2}, \stepref{s5exvalues} and \stepref{s5exdims}, their multiplicities are at
   least $1$, $2k$, $\ell$, $\ell$. These add up to $4k^2-1$, the order of $A(H)$, so by
   \ST{Theorem 2.2} they are the multiplicities and there is no other eigenvalue.~\qed\end{proof}
  \step{s5exspec58}{$\Spec(H)$ is the spectrum of case (iii) of Theorem~\ref{thm:5.8} with
   $s=2k-1$ and $t=k-1$.}
  \begin{proof}\pf\ With these values, $(s+1)t=2k(k-1)=\kappa$, $s+1=2k$, $-t=-(k-1)$,
   $c=t(s+1)(s-t)/(s-1)=(k-1)\cdot2k\cdot k/(2k-2)=k^2$, so $\sqrt c=k$, and
   $\frac12(s-1)(s+2)=(k-1)(2k+1)=\ell$. Compare with \stepref{s5exspecH}.~\qed\end{proof}
  \qedstep
  \begin{proof}\pf\ $H$ is $\kappa$-regular on $|\cB'|=4k^2-1$ vertices by \stepref{s5exvalues}, and
   connected by Lemma~\ref{lem:5.1}, since $\mult_H(\kappa)=1$ by
   \stepref{s5exspecH}. The spectrum is given by Steps \stepref{s5exspecH} and
   \stepref{s5exspec58}.~\qed\end{proof}
 \end{proof}
 \step{s5exline}{Assertion 2 holds.}
 \begin{proof}
  \pflet{$M_H$ and $M_1$ be the incidence matrices of $H$ and of $I(D_1)$.}
  \step{s5exCH}{$\Spec(M_HM_H^T)=\ms{(2\kappa)^{1},(\kappa-k+1)^{2k},(\kappa+k)^{\ell},
   (\kappa-k)^{\ell}}$, with all members positive.}
  \begin{proof}\pf\ $M_HM_H^T=\kappa I+A(H)$ by \ST{Theorem 6.2} and \stepref{s5exspec}; shift
   $\Spec(H)$ by $\kappa$ with \ST{Theorem 3.1}. $\kappa-k=k(2k-3)>0$ as $k\ge2$.~\qed\end{proof}
  \step{s5exC1}{$\Spec(M_1M_1^T)$ is $\Spec(M_HM_H^T)$ with one copy of $0$ added.}
  \begin{proof}
   \step{s5exC1a}{$\Spec(N_1N_1^T)=\ms{(r_1k_1)^{1},(2k-1)^{v_1-1}}$ and $\rk N_1=v_1$, for the
    incidence matrix $N_1$ of $D_1$.}
   \begin{proof}\pf\ Lemma~\ref{lem:design}, assertion 3, with $r_1-\lambda_1=2k-1>0$; then
    \ST{Theorem 3.2} and \ST{Theorem 3.5}.~\qed\end{proof}
   \step{s5exC1b}{$\Spec(M_1M_1^T)$ consists of $0$ and $r_1+k_1$, of
    $\frac12\bigl(r_1+k_1\pm\sqrt{(r_1-k_1)^2+4(2k-1)}\bigr)$ with multiplicity $v_1-1$ each, and
    of $k_1$ with multiplicity $b_1-v_1$.}
   \begin{proof}\pf\ Lemma~\ref{lem:rootgraph}, assertion 3, for $Z=I(D_1)$, which applies by
    Lemma~\ref{lem:incidence}; then \ST{Theorem 6.6} with $s=r_1$, $t=k_1$, $\rho=v_1$ and
    \stepref{s5exC1a}. The roots for $r_1k_1$ are $0$ and $r_1+k_1$.~\qed\end{proof}
   \qedstep
   \begin{proof}\pf\ Substitute \asm[4]{0}: $r_1+k_1=2\kappa$, $r_1-k_1=2(k-1)$, so the square
    root is $2k$ and the pair is $\kappa\pm k$; $v_1-1=\ell$, $k_1=\kappa-k+1$ and
    $b_1-v_1=2k$. Compare with \stepref{s5exCH}.~\qed\end{proof}
  \end{proof}
  \step{s5exedges}{$|E(H)|=|E(I(D_1))|=v_1r_1$.}
  \begin{proof}\pf\ $|E(H)|=\frac12\kappa(4k^2-1)=k(k-1)(2k-1)(2k+1)$ by
   \path{SimpleGraph.sum_degrees_eq_twice_card_edges} and \stepref{s5exspec}, and this is
   $v_1r_1$, which is $|E(I(D_1))|$ by Lemma~\ref{lem:incidence}.~\qed\end{proof}
  \step{s5exnz}{$\mult_{2I+A(L(H))}(\gamma)=\mult_{2I+A(F(D_1))}(\gamma)$ for every
   $\gamma\ne0$.}
  \begin{proof}\pf\ By \ST{Theorem 6.2}, $M_H^TM_H=2I+A(L(H))$ and $M_1^TM_1=2I+A(F(D_1))$; by
   \ST{Theorem 6.1}, their multiplicities at $\gamma\ne0$ are those of $M_HM_H^T$ and
   $M_1M_1^T$, which agree by \stepref{s5exC1}.~\qed\end{proof}
  \qedstep
  \begin{proof}\pf\ Both matrices have order $v_1r_1$ by \stepref{s5exedges}, so their
   multiplicities at $0$ agree by \ST{Theorem 2.2} and \stepref{s5exnz}. Thus
   $\Spec(2I+A(L(H)))=\Spec(2I+A(F(D_1)))$ by \ST{Proposition 2.6}; shift by $-2$ with
   \ST{Theorem 3.1}, and apply \ST{Proposition 2.6}.~\qed\end{proof}
 \end{proof}
 \step{s5exnot}{Assertion 3 holds.}
 \begin{proof}
  \step{s5exnonbip}{$H$ is not bipartite.}
  \begin{proof}\pf\ $H$ is connected and $\kappa$-regular, and $-\kappa\notin\Spec(H)$ by
   \stepref{s5exspec}, since $\kappa>k$ for $k\ge2$; apply Lemma~\ref{lem:5.3}.~\qed\end{proof}
  \step{s5exnoiso}{\assume{$D_2$ is a design with the parameters of $D_1$, and
   $L(H)\cong F(D_2)$.}\prove{\textsc{False}.}}
  \begin{proof}\pf\ $F(D_2)=L(I(D_2))$, and $I(D_2)$ is bipartite without isolated vertices by
   Lemma~\ref{lem:incidence}. $H$ is connected on $4k^2-1\ge15$ vertices by \stepref{s5exspec},
   and not bipartite by \stepref{s5exnonbip}. This contradicts EX4 of the
   \hyperref[input:ex513]{inputs for Example 5.13}. \new\ The Lean proof,
   \path{Design.not_nonempty_iso_flagGraph_lineGraph_of_degree}, uses a clique bound in place of
   EX4, so the example does not need EX4: the $r_1$ flags on a point of $D_2$ form an
   $r_1$-clique of $F(D_2)$ (\path{Design.exists_clique_flagGraph_card_replication}), while
   every clique of $L(H)$ has at most $\kappa$ vertices, since $H$ is $\kappa$-regular with
   $\kappa\ge3$ (\path{SimpleGraph.card_clique_lineGraph_le}); and
   $r_1=(k-1)(2k+1)=\kappa+k-1>\kappa=2k(k-1)\ge4$ for $k\ge2$.~\qed\end{proof}
  \qedstep
  \begin{proof}\pf\ $L(H)$ is cospectral with $F(D_1)$ by \stepref{s5exline}, and is isomorphic
   to no $F(D_2)$ by \stepref{s5exnoiso}; this is the negation of Definition~\ref{def:flag}
   for $D_1$.~\qed\end{proof}
 \end{proof}
 \qedstep
 \begin{proof}\pf\ Steps \stepref{s5exspec}, \stepref{s5exline} and \stepref{s5exnot}.~\qed\end{proof}
\end{proof}

\begin{remark}
 \lean{Computations.C15.flagGraph_cospectral_double_lineGraph}
 \lean{Computations.C15.flagGraph_not_nonempty_iso_double_lineGraph}
 \lean{Computations.C15.design}\lean{Computations.C15.petersen}
 \lean{Computations.C15.petersenIso}\lean{Computations.C15.nonempty_iso_petersen}
 \lean{Computations.C15.four_cliques_graph1}\lean{Computations.C15.four_cliques_graph2}
 By EX1 and EX2, the resolvable designs of the example exist for $k=2^d$ with $d\ge1$ and for
 $k=3$, and by EX3 the designs $D_1$ exist for $k=2$ and $k=4$, with parameters $(6,3,10,5,2)$
 and $(28,21,36,27,20)$; for these two, the flag graph is not characterized by its spectrum. For
 $k=2$ the paper describes $H$ as the line graph of the Petersen graph; we do not check this.
 Computation C15 checks that the flag graph of the design $(6,3,10,5,2)$ is cospectral with
 $L(L(\text{Petersen}))$ and not isomorphic to it; the two graphs have 30 and 15 cliques of
 size 4. In Lean the flag graph is the line graph of the incidence graph of
 \path{LineGraphs/Computations/C13Design.lean}, and the Petersen graph is the settled
 \path{SimpleGraph.petersen}. The certificate is computed on \path{Computations.C15.petersen},
 the pentagon, the pentagram and five spokes, and \path{Computations.C15.petersenIso} is an
 explicit isomorphism from it to \path{SimpleGraph.petersen}, through which the two statements
 are transported. For $k=3$ the parameters of $D_1$ are $(15,10,21,14,9)$, and no such design exists, by the Hall--Connor section, as in the proof of
 Remark 5.12. Computation C23 (\path{LineGraphs/Computations/C23.lean}, with the data of
 \path{cgss/certificates/c23_lean_data.py}) builds the designs of the instances $k=2$ and $k=4$ and proves
 both instances, and the case $k=3$ above, in Lean; the existence of the designs of EX1 for every
 $k=2^d$ is a citation, outside the scope.
\end{remark}

\endgroup

\subsection{The Hall--Connor nonexistence theorem}\label{hc:sec:main}
\begin{mergedproofsection}
\raggedbottom

We will follow the ordinary proof of
\path{cgss/proofs/hallconnor/hallconnor.tex}. The argument proves the
nonexistence assertion used in Remark 5.12 and Example 5.13 of
\cite{CGSS}, together with its complementary parameters. We use
Hall and Connor's argument \cite{HC}, including proofs of the instances
they take from Connor. We have not consulted Connor's paper directly.
We omit the two comparison remarks of the ordinary proof, which the
argument does not use.

We follow the blueprint conventions and Lamport \cite{L95}:
\textsc{Assume} introduces hypotheses, \textsc{Prove} the conclusion,
\textsc{Let} a definition, and \textsc{Pick} a witness.
Each numbered step has a proof. A reference $\langle0\rangle{:}i$
means hypothesis $i$ of the result; $\langle k\rangle$ means the
assumption of the enclosing level-$k$ step. Closing \textsc{Qed}
steps state why the preceding steps suffice. The tags \new\ and
\repair\ mark a change of argument and a filled gap, respectively.
The no-op commands \verb|\lean| and \verb|\uses| record the Lean
declarations and the direct dependencies.

We use the notation of \path{cgss/notation.md}, Sections 7 and 10:
\texttt{Design P B v k b r lam}, \path{D.incMatrix}, and \path{D.compl}.
The Hall--Connor names not listed individually there retain the names
of the ordinary proof, following the namespace and
naming rules of Section 10. Points and block labels are separate finite
types with decidable equality. All sets used in counts are finite sets;
their cardinalities are natural numbers. Moment identities and
subtractions in weighted sums are evaluated in $\ints$ or $\rats$,
never by truncated natural subtraction. Matrix entries are $M(u,w)$.
We write $I_U$ for the identity, $J_{U,V}$ for the all-ones matrix,
and $\Zj_U$ for the all-ones vector. Omitted index types are determined
by the products. All ranks are over $\rats$.

\subsubsection{Definitions and Mathlib inputs}

\begin{namedresult}{Definition}{(labelled designs)}\label{hc:def:design}
 \lean{Design}\lean{Design.IsSymmetric}\uses{}
 \leavevmode\par
 \assume{\begin{pfenum}
  \item $P,\cB$ are finite types with decidable equality.
  \item $v,k,b,r,\lambda\in\nn$, with $v>1$, $k>1$, and $b>0$.
 \end{pfenum}}
 A design $D$ with these parameters consists of a function
 $B\mapsto D(B)\subseteq P$ such that
 \[
 \begin{gathered}
 |P|=v,\quad |\cB|=b,\quad |D(B)|=k\quad(B\in\cB),\\
 |\{B:p\in D(B)\}|=r\quad(p\in P),\\
 |\{B:p,q\in D(B)\}|=\lambda\quad(p,q\in P,\ p\ne q).
 \end{gathered}
 \]
 We write $p\in B$ for $p\in D(B)$; the map $D$ need not be injective.
 The predicate $D.\mathrm{IsSymmetric}$ means $v=b$.
 In Lean the parameter names are \texttt{v k b r lam}; any positivity
 fields of the shared interface are satisfied by the parameters here.
\end{namedresult}

\begin{namedresult}{Definition}{(incidence and block matrices)}\label{hc:def:matrices}
 \lean{Design.incMatrix}\lean{Design.G}\lean{Design.C}\uses{hc:def:design}
 \leavevmode\par
 \assume{\begin{pfenum}
  \item $P,\cB$ are finite types with decidable equality.
  \item $D$ is a design with parameters $(v,k,b,r,\lambda)$.
 \end{pfenum}}
 \pflet{$N\triangleq D.\mathrm{incMatrix}\in\mat{P}{\cB}{\ints}$, where}
 \[
 N(p,B)\triangleq\begin{cases}1,&p\in B,\\0,&p\notin B.\end{cases}
 \]
 Write $N_{\rats}$ for its entrywise rational cast. Put
 $G\triangleq N_{\rats}^TN_{\rats}$ and
 $C\triangleq35I_{\cB}-7G+10J_{\cB}$, both over $\rats$.
 A product containing $G$ or $C$ uses $N_{\rats}$; we abbreviate it
 by $N$ in rational calculations. Products asserted integral use the
 original $N$.
\end{namedresult}

\begin{namedresult}{Definition}{(the block graph and all six-cliques)}\label{hc:def:graph}
 \lean{Design.Z}\lean{Design.sixCliques}\uses{hc:def:design}
 \leavevmode\par
 \assume{\begin{pfenum}
  \item $P,\cB$ are finite types with decidable equality.
  \item $D$ is a design with parameters $(v,k,b,r,\lambda)$.
 \end{pfenum}}
 \pflet{$Z$ be the simple graph on $\cB$ with}
 \[
 B\sim B'\quad\Longleftrightarrow\quad B\ne B'\ \text{and}\ |B\cap B'|=1.
 \]
 Here $B\cap B'$ means $D(B)\cap D(B')$.
 Put $\Gamma(B)\triangleq\{B':B\sim B'\}$ and
 \[
 \cS\triangleq D.\mathrm{sixCliques}
   \triangleq\{S\subseteq\cB:|S|=6\ \text{and}\
      \forall X,Y\in S,\ X\ne Y\Rightarrow X\sim Y\}.
 \]
 Thus $\cS$ is the finite set of \emph{all} six-cliques, not a choice
 of local cliques. Its members form a finite subtype when used as indices.
\end{namedresult}

\begingroup\raggedright
\paragraph{Finite counting and order.}\factlabel{F1}{hc:facts:finite}
\uses{}
All declarations below were checked in Mathlib at commit
\path{082e2d37} (4 September 2026), the version pinned in
\path{lake-manifest.json}; paths are relative to \path{Mathlib/}. Finite double counting is
interchange of indicator sums, followed by grouping fibres:
\begin{itemize}
 \item \path{Finset.sum_comm}, in
 \path{Algebra/BigOperators/Group/Finset/Sigma.lean};
 \item \path{Finset.card_eq_sum_card_fiberwise}, in
 \path{Algebra/BigOperators/Group/Finset/Basic.lean};
 \item \path{Finset.card_sdiff_of_subset},
 \path{Finset.card_union_add_card_inter},
 \path{Finset.eq_of_subset_of_card_le}, \path{Finset.card_pos}, and
 \path{Finset.card_le_one}, \path{Finset.card_eq_one},
 \path{Finset.card_eq_two}, \path{Finset.card_eq_three}, and
 \path{Finset.card_image_of_injective}, in \path{Data/Finset/Card.lean};
 \item \path{Fintype.bijective_iff_injective_and_card} and
 \path{Fintype.equivFin}, in \path{Data/Fintype/EquivFin.lean};
 \item \path{Fintype.card_sum}, in \path{Data/Fintype/Sum.lean},
 and \path{Fintype.card_option}, in \path{Data/Fintype/Option.lean};
 \item \path{Fintype.sum_sum_type} and \path{Fintype.sum_option}, in
 \path{Data/Fintype/BigOperators.lean};
 \item \path{Finset.sum_eq_zero_iff_of_nonneg}, in
 \path{Algebra/Order/BigOperators/Group/Finset.lean}.
\end{itemize}
The last declaration says that a sum of nonnegative terms is zero
exactly when every term is zero. The equal-cardinality theorem turns
an injection between finite types of equal size into a bijection.
An enumeration of a finite subset is an equivalence to its subtype;
its inclusion in the ambient type is injective. Integer and rational
identities below use ring arithmetic, nonnegative squares, and finite
sum distributivity. Small numerical tables are explicit computations.

\paragraph{Matrix operations and reindexing.}\factlabel{F2}{hc:facts:matrix}
\uses{}
We use \path{Matrix.mul_apply}, \path{Matrix.mul_assoc},
\path{Matrix.transpose_mul}, \path{Matrix.dotProduct_mulVec},
\path{Matrix.mulVec_mulVec}, \path{Matrix.map_mul}, and
\path{Matrix.submatrix_mul_equiv}, all in \path{Data/Matrix/Mul.lean}.
The last identity says that, for
an equivalence $e$ of the summed index type,
\[
 (A\mathbin{\mathrm{submatrix}}f\,e)
 (B\mathbin{\mathrm{submatrix}}e\,g)
       =(AB)\mathbin{\mathrm{submatrix}}f\,g.
\]
We also use \path{Matrix.transpose_map} and
\path{Matrix.transpose_submatrix}, in
\path{LinearAlgebra/Matrix/Defs.lean}, and \path{Matrix.map_one}, in
\path{Data/Matrix/Diagonal.lean}. Applied to the ring homomorphism
$\ints\to\rats$, these transport products, transposes, and $I$.
The cast of $J$ is $J$ entrywise because $1$ maps to $1$; addition
and scalar multiplication commute with the cast entrywise.

\paragraph{Ranks and determinants.}\factlabel{F3}{hc:facts:rank}
\uses{}
The following declarations are in
\path{LinearAlgebra/Matrix/Rank.lean}:
\path{Matrix.rank_self_mul_transpose}, \path{Matrix.rank_transpose},
\path{Matrix.rank_of_isUnit}, \path{Matrix.rank_of_det_ne_zero},
\path{Matrix.rank_add_rank_le_card_of_mul_eq_zero}, and
\path{Matrix.rank_submatrix_le}.
The identity $\operatorname{rk}(NN^T)=\operatorname{rk}N$ applies over
the linearly ordered field $\rats$; rank under transpose applies over
a field. A nonzero determinant gives full rank. If $AB=0$, the sum
of their ranks is at most the size of the shared index type.
\path{Matrix.isUnit_iff_isUnit_det}, in
\path{LinearAlgebra/Matrix/NonsingularInverse.lean}, turns a nonzero
rational determinant into an invertible matrix.

\paragraph{Determinant arithmetic.}\factlabel{F4}{hc:facts:det}
\uses{}
We use \path{Matrix.det_mul}, \path{Matrix.det_transpose},
\path{Matrix.det_smul}, \path{Matrix.det_one}, and \path{Int.cast_det}, in
\path{LinearAlgebra/Matrix/Determinant/Basic.lean}.
We use the matrix determinant lemma
\path{Matrix.det_one_add_replicateCol_mul_replicateRow}, in
\path{LinearAlgebra/Matrix/SchurComplement.lean}, in the form
$\det(I+\Zu\Zv^T)=1+\Zv^T\Zu$ (the shared singleton index is
\texttt{Unit}). These give multiplicativity, scalar homogeneity,
and compatibility with the cast $\ints\to\rats$. In particular, the
determinant over $\rats$ of an integral matrix is the cast of its
integer determinant. We do not use the analytic
positive-semidefinite determinant theorem, whose scalar hypotheses
do not include $\rats$.

\paragraph{Integer divisibility.}\factlabel{F5}{hc:facts:prime}
\uses{}
We use \path{Prime.dvd_of_dvd_pow}, in \path{Algebra/Prime/Defs.lean},
\path{Int.prime_iff_natAbs_prime}, in \path{RingTheory/Int/Basic.lean},
and \path{Nat.prime_five}, in \path{Data/Nat/Prime/Defs.lean}.
The last two prove that $5$ is prime in $\ints$; the first gives
$5\mid d^2\Rightarrow5\mid d$ for $d\in\ints$.

\par\endgroup

\subsubsection{Counts and intersections}

\begin{namedresult}{Lemma}{(the forced parameters)}\label{hc:lem:parameters}
 \lean{Design.exists_of_pair_counts_15_5_2}
 \uses{hc:def:design,hc:facts:finite}
 \leavevmode\par
 \assume{\begin{pfenum}
  \item $P,\cB$ are finite types with decidable equality and $|P|=15$.
  \item $F:\cB\to\operatorname{Finset}(P)$ has $|F(B)|=5$ for every $B$.
  \item Each pair of distinct points lies in exactly two labelled blocks of $F$.
 \end{pfenum}}
 \prove{Every point is in seven labels, $|\cB|=21$, and
 \[
 \exists D:\texttt{Design P B 15 5 21 7 2},\
       \forall B\in\cB,\ D.\mathrm{block}(B)=F(B).
 \]
 Here the Lean type variable \texttt{B} denotes $\cB$.}
\end{namedresult}
\begin{proof}
 \pfsketch\ We will count first through one point, then over all points.
 \step{hcParPoint}{\assume{$p\in P$.}
  \prove{$r_p\triangleq|\{B:p\in F(B)\}|=7$.}}
 \begin{proof}
  \step{hcParBlocks}{The set of pairs $(q,B)$ with $q\ne p$ and
   $p,q\in F(B)$ has size $4r_p$.}
  \begin{proof}\pf\ For each of the $r_p$ labels, deleting $p$ from
   its five points leaves four choices, by \asm[2]{0} and
   the fibre count in \ref{hc:facts:finite}.~\qed\end{proof}
  \step{hcParPairs}{The same set has size $28$.}
  \begin{proof}\pf\ There are $15-1=14$ choices of $q$ and two labels
   for each, by \asm[1]{0}, \asm[3]{0} and
   \ref{hc:facts:finite}.~\qed\end{proof}
  \qedstep
  \begin{proof}\pf\ The equality $4r_p=28$ from
   \stepref{hcParBlocks} and \stepref{hcParPairs} gives $r_p=7$.~\qed\end{proof}
 \end{proof}
 \step{hcParTotal}{$5|\cB|=105$ and $|\cB|=21$.}
 \begin{proof}\pf\ Count $(p,B)$ with $p\in F(B)$ by blocks and by
  points, using \asm[1]{0}, \asm[2]{0}, \stepref{hcParPoint}, and
  \ref{hc:facts:finite}: $5|\cB|=15\cdot7$. Division by five gives
  the second assertion.~\qed\end{proof}
 \qedstep
 \begin{proof}\pf\ Steps \stepref{hcParPoint} and \stepref{hcParTotal},
  with \asm[1]{0}, \asm[2]{0}, and \asm[3]{0}, give every field of
  Definition~\ref{hc:def:design}. Choose the structure with block
  function $F$; its block equality is reflexivity, and
  $15>1$, $5>1$, $21>0$ are numerical.~\qed\end{proof}
\end{proof}

\begin{namedresult}{Lemma}{(incidence identities)}\label{hc:lem:incidence}
 \lean{Design.incMatrix_mul_transpose_15_5_2}
 \lean{Design.incMatrix_row_sum_15_5_2}\lean{Design.incMatrix_col_sum_15_5_2}
 \lean{Design.G_apply}\lean{Design.G_mul_G}\lean{Design.G_mul_J}
 \uses{hc:def:design,hc:def:matrices,hc:facts:finite,hc:facts:matrix}
 \leavevmode\par
 \assume{\begin{pfenum}
  \item $P,\cB$ are finite types with decidable equality.
  \item $D$ is a design with parameters $(15,5,21,7,2)$.
 \end{pfenum}}
 \prove{The integral incidence matrix has
 \[
 NN^T=5I_P+2J_P,\quad N\Zj_{\cB}=7\Zj_P,\quad
 N^T\Zj_P=5\Zj_{\cB}.
 \]
 Over $\rats$, $G(B,B')=|B\cap B'|$, $G^T=G$, and
 \[
 G^2=5G+50J_{\cB},\quad GJ_{\cB}=J_{\cB}G=35J_{\cB},\quad
 J_{\cB}^2=21J_{\cB}.
 \]}
\end{namedresult}
\begin{proof}
 \pfsketch\ The matrix entries count points or labels; associativity
 then gives the identities for $G$.
 \step{hcIncDiag}{$(NN^T)(p,p)=7$ for every $p\in P$.}
 \begin{proof}\pf\ By Definitions~\ref{hc:def:matrices} and
  \ref{hc:def:design}, the entry is the sum of the squares of the
  seven incidence indicators equal to one; use \ref{hc:facts:matrix}.~\qed\end{proof}
 \step{hcIncOff}{\assume{$p,q\in P$ and $p\ne q$.}
  \prove{$(NN^T)(p,q)=2$.}}
 \begin{proof}\pf\ The product of the two indicators is one exactly
  on labels containing $p,q$, of which there are two by
  Definition~\ref{hc:def:design}; expand by \ref{hc:facts:matrix}.~\qed\end{proof}
 \step{hcIncNN}{$NN^T=5I_P+2J_P$.}
 \begin{proof}\pf\ Steps \stepref{hcIncDiag} and \stepref{hcIncOff}
  treat the exhaustive alternatives $p=q$ and $p\ne q$.~\qed\end{proof}
 \step{hcIncRow}{$N\Zj_{\cB}=7\Zj_P$.}
 \begin{proof}\pf\ The row sum counts the seven labels on $p$,
  by Definition~\ref{hc:def:design} and \ref{hc:facts:matrix}.~\qed\end{proof}
 \step{hcIncCol}{$N^T\Zj_P=5\Zj_{\cB}$.}
 \begin{proof}\pf\ The column sum counts the five points of a block,
  by Definition~\ref{hc:def:design} and \ref{hc:facts:matrix}.~\qed\end{proof}
 \step{hcIncG}{$G(B,B')=|B\cap B'|$ and $G^T=G$.}
 \begin{proof}\pf\ In Definition~\ref{hc:def:matrices}, the summand
  $N(p,B)N(p,B')$ is the intersection indicator. Its symmetry in
  $B,B'$ gives symmetry of $G$, by \ref{hc:facts:matrix}.~\qed\end{proof}
 \step{hcIncSquare}{$G^2=5G+50J_{\cB}$.}
 \begin{proof}
  \pf\ Cast \stepref{hcIncNN} and \stepref{hcIncCol} entrywise
  using the product, transpose, and constant-matrix rules of
  \ref{hc:facts:matrix}. Then
  \begin{align*}
   G^2&=N^T(NN^T)N &&\text{[\ref{hc:facts:matrix}]}\\
      &=5G+2N^TJ_PN &&\text{[\stepref{hcIncNN}]}\\
      &=5G+50J_{\cB} &&\text{[\stepref{hcIncCol}, $5\cdot5=25$].}
  \end{align*}\qed
 \end{proof}
 \step{hcIncGJ}{$GJ_{\cB}=35J_{\cB}$.}
 \begin{proof}\pf\ Steps \stepref{hcIncRow} and \stepref{hcIncCol}
  give $G\Zj_{\cB}=N^T(7\Zj_P)=35\Zj_{\cB}$;
  each column of $J$ is $\Zj$, by \ref{hc:facts:matrix}.~\qed\end{proof}
 \step{hcIncJG}{$J_{\cB}G=35J_{\cB}$.}
 \begin{proof}\pf\ Transpose \stepref{hcIncGJ}, using the symmetry
  in \stepref{hcIncG} and \ref{hc:facts:matrix}.~\qed\end{proof}
 \step{hcIncJJ}{$J_{\cB}^2=21J_{\cB}$.}
 \begin{proof}\pf\ Each entry sums $|\cB|=21$ ones, by
  Definition~\ref{hc:def:design} and \ref{hc:facts:matrix}.~\qed\end{proof}
 \qedstep
 \begin{proof}\pf\ Steps \stepref{hcIncNN}--\stepref{hcIncJJ}
  give every asserted identity.~\qed\end{proof}
\end{proof}

\begin{namedresult}{Lemma}{(intersection moments)}\label{hc:lem:moments}
 \lean{Design.intersection_moments_15_5_2}
 \uses{hc:def:design,hc:facts:finite}
 \leavevmode\par
 \assume{\begin{pfenum}
  \item $P,\cB$ are finite types with decidable equality and
  $D$ is a design with parameters $(15,5,21,7,2)$.
  \item $B\in\cB$ and $n_i\triangleq|\{B'\ne B:|B\cap B'|=i\}|$ for $0\le i\le5$.
 \end{pfenum}}
 \prove{Over $\ints$,
 $\sum_{i=0}^5n_i=20$, $\sum_{i=0}^5in_i=30$, and
 $\sum_{i=0}^5i(i-1)n_i=20$.
 In particular $\sum_{B'\ne B}|B\cap B'|=30$.}
\end{namedresult}
\begin{proof}
 \pfsketch\ We will count labels, incidences, and ordered pairs of
 distinct points. The moment statements will be cited outside this proof.
 \step{hcMomZero}{$\sum_{i=0}^5n_i=20$.}
 \begin{proof}\pf\ Intersection size is between zero and five;
  the fibres partition the $21-1$ other labels, by
  Definition~\ref{hc:def:design} and \ref{hc:facts:finite}.~\qed\end{proof}
 \step{hcMomOne}{$\sum_{i=0}^5in_i=30$ and
  $\sum_{B'\ne B}|B\cap B'|=30$.}
 \begin{proof}\pf\ Count $(p,B')$ with $p\in B\cap B'$ and $B'\ne B$.
  For each of five points, six of its seven labels remain.
  Grouping by $B'$ and then by $i$ gives both sums, using
  Definition~\ref{hc:def:design} and \ref{hc:facts:finite}.~\qed\end{proof}
 \step{hcMomTwo}{$\sum_{i=0}^5i(i-1)n_i=20$.}
 \begin{proof}\pf\ Count $(p,q,B')$ with $p,q\in B\cap B'$, $p\ne q$,
  $B'\ne B$. Each of $5\cdot4$ ordered pairs in $B$ has exactly
  one further label; a label with intersection size $i$ contributes
  $i(i-1)$, using Definition~\ref{hc:def:design} and
  \ref{hc:facts:finite}.~\qed\end{proof}
 \qedstep
 \begin{proof}\pf\ Steps \stepref{hcMomZero}, \stepref{hcMomOne},
  and \stepref{hcMomTwo} prove the four assertions.~\qed\end{proof}
\end{proof}

\begin{namedresult}{Lemma}{(block intersections)}\label{hc:lem:intersections}
 \textup{[\cite{HC}, Lemma 4.1 at $k=5$.]}
 \lean{Design.inter_card_eq_one_or_two}\lean{Design.card_neighbors}
 \lean{Design.card_filter_inter_card_eq_two}
 \uses{hc:def:design,hc:def:graph,hc:lem:moments,hc:facts:finite}
 \leavevmode\par
 \assume{\begin{pfenum}
  \item $P,\cB$ are finite types with decidable equality.
  \item $D$ is a design with parameters $(15,5,21,7,2)$.
 \end{pfenum}}
 \prove{Distinct labels meet in one or two points. Every block has
 exactly ten other labels of each kind; in particular $|\Gamma(B)|=10$.}
\end{namedresult}
\begin{proof}
 \pfsketch\ A nonnegative linear combination of the moments excludes
 intersection sizes other than one and two.
 \step{hcInterBlock}{\assume{$B\in\cB$, with $n_i$ as in
  Lemma~\ref{hc:lem:moments}.}
  \prove{Every other label meets $B$ in one or two points,
  and exactly ten labels have each intersection size.}}
 \begin{proof}
  \step{hcInterPolynomial}{$\sum_{i=0}^5(i-1)(i-2)n_i=0$.}
  \begin{proof}\pf\ The polynomial equals $i(i-1)-2i+2$;
   Lemma~\ref{hc:lem:moments} gives $20-2\cdot30+2\cdot20=0$.~\qed\end{proof}
  \step{hcInterSupport}{$n_0=n_3=n_4=n_5=0$.}
  \begin{proof}
   \step{hcInterCoeffs}{The coefficients $(i-1)(i-2)$, in order
    $i=0,1,2,3,4,5$, are $2,0,0,2,6,12$.}
   \begin{proof}\pf\ Direct integer substitution in the polynomial
    of \stepref{hcInterPolynomial}.~\qed\end{proof}
   \qedstep
   \begin{proof}\pf\ The $n_i$ are nonnegative cardinalities.
    Apply the zero-sum criterion of \ref{hc:facts:finite} to
    \stepref{hcInterPolynomial} and \stepref{hcInterCoeffs}; the
    four indicated coefficients are positive.~\qed\end{proof}
  \end{proof}
  \step{hcInterCounts}{$n_1=n_2=10$.}
  \begin{proof}\pf\ Lemma~\ref{hc:lem:moments} and
   \stepref{hcInterSupport} give $n_1+n_2=20$ and $n_1+2n_2=30$.
   Subtract the equations.~\qed\end{proof}
  \qedstep
  \begin{proof}\pf\ The zero fibres of \stepref{hcInterSupport}
   exclude the other intersection sizes; \stepref{hcInterCounts}
   and Definition~\ref{hc:def:graph} give the multiplicities and degree.~\qed\end{proof}
 \end{proof}
 \qedstep
 \begin{proof}\pf\ Step \stepref{hcInterBlock} applies to every $B$;
  Definition~\ref{hc:def:graph} identifies its ten one-point
  intersections with $\Gamma(B)$.~\qed\end{proof}
\end{proof}

\begin{namedresult}{Lemma}{(common-neighbour counts)}\label{hc:lem:common}
 \textup{[\cite{HC}, Lemma 4.2 at $k=5$.]}
 \lean{Design.intersection_pair_counts}\lean{Design.card_commonNeighbors}
 \uses{hc:def:design,hc:def:matrices,hc:def:graph,hc:lem:incidence,
 hc:lem:moments,hc:lem:intersections,hc:facts:finite,hc:facts:matrix}
 \leavevmode\par
 \assume{\begin{pfenum}
  \item $P,\cB$ are finite types with decidable equality and
  $D$ is a design with parameters $(15,5,21,7,2)$.
  \item $B_1\ne B_2$ are block labels, $s\triangleq|B_1\cap B_2|$, and
  $x_{ij}\triangleq|\{B\notin\{B_1,B_2\}:|B_1\cap B|=i,\ |B_2\cap B|=j\}|$.
 \end{pfenum}}
 \prove{The four counts are
 \[
 \begin{array}{c|rrrr}
 s&x_{11}&x_{12}&x_{21}&x_{22}\\\hline
 1&5&4&4&6\\2&4&6&6&3
 \end{array}.
 \]
 Thus adjacent vertices have five common neighbours and distinct
 nonadjacent vertices have four.}
\end{namedresult}
\begin{proof}
 \pfsketch\ We will expose the four filtered-sum equations, including
 both terms removed from the entry of $G^2$, before solving them.
 \step{hcCommonTotal}{The four fibres partition
  $\cB\diff\{B_1,B_2\}$, and $x_{11}+x_{12}+x_{21}+x_{22}=19$.}
 \begin{proof}\pf\ Lemma~\ref{hc:lem:intersections} partitions the
  $21-2$ remaining labels into the four disjoint fibres; use
  \asm[2]{0}, \ref{hc:facts:finite}, and
  Definition~\ref{hc:def:design}.~\qed\end{proof}
 \step{hcCommonFirst}{$x_{11}+x_{12}+2x_{21}+2x_{22}=30-s$.}
 \begin{proof}\pf\ From the last conclusion of Lemma~\ref{hc:lem:moments}
  for $B_1$, remove the term $B_2$ of size $s$. The four fibres have
  first weights $1,1,2,2$, using \stepref{hcCommonTotal},
  Lemma~\ref{hc:lem:intersections}, \asm[2]{0}, and \ref{hc:facts:finite}.~\qed\end{proof}
 \step{hcCommonSecond}{$x_{11}+2x_{12}+x_{21}+2x_{22}=30-s$.}
 \begin{proof}\pf\ Apply Lemma~\ref{hc:lem:moments} to $B_2$ and remove
  the term $B_1$ of size $s$. The second weights are $1,2,1,2$,
  using \stepref{hcCommonTotal}, Lemma~\ref{hc:lem:intersections},
  \asm[2]{0}, and \ref{hc:facts:finite}.~\qed\end{proof}
 \step{hcCommonProduct}{$x_{11}+2x_{12}+2x_{21}+4x_{22}=50-5s$.}
 \begin{proof}
  \step{hcCommonEntry}{$(G^2)(B_1,B_2)=5s+50$.}
  \begin{proof}\pf\ Substitute $G(B_1,B_2)=s$ into the identity
   $G^2=5G+50J$ of Lemma~\ref{hc:lem:incidence}.~\qed\end{proof}
  \step{hcCommonRemoved}{The terms indexed by $B_1$ and $B_2$ in
   $(G^2)(B_1,B_2)$ are each $5s$.}
  \begin{proof}\pf\ The products are $G(B_1,B_1)G(B_1,B_2)$ and
   $G(B_1,B_2)G(B_2,B_2)$; both diagonals equal five by
   Definition~\ref{hc:def:design} and Lemma~\ref{hc:lem:incidence}.~\qed\end{proof}
  \qedstep
  \begin{proof}\pf\ Expand the matrix entry using \ref{hc:facts:matrix},
   remove the two distinct terms of \stepref{hcCommonRemoved} from
   \stepref{hcCommonEntry}, then group the remaining terms by their
   weights $ij$ using \stepref{hcCommonTotal},
   Lemma~\ref{hc:lem:intersections}, \asm[2]{0}, and
   \ref{hc:facts:finite}. This gives
   $(5s+50)-10s=50-5s$.~\qed\end{proof}
 \end{proof}
 \step{hcCommonSolve}{$x_{22}=9-3s$, $x_{12}=x_{21}=2+2s$, and $x_{11}=6-s$.}
 \begin{proof}\pf\ Subtract \stepref{hcCommonFirst} and
  \stepref{hcCommonSecond} from \stepref{hcCommonProduct}, then add
  \stepref{hcCommonTotal}, to get $x_{22}=9-3s$.
  Subtract \stepref{hcCommonTotal} and this value from each weighted
  marginal, then use the total for $x_{11}$.~\qed\end{proof}
 \step{hcCommonOne}{\case{$s=1$.} The count vector is $(5,4,4,6)$.}
 \begin{proof}\pf\ Substitute \asm{1} in \stepref{hcCommonSolve}.~\qed\end{proof}
 \step{hcCommonTwo}{\case{$s=2$.} The count vector is $(4,6,6,3)$.}
 \begin{proof}\pf\ Substitute \asm{1} in \stepref{hcCommonSolve}.~\qed\end{proof}
 \step{hcCommonNeighbours}{$\Gamma(B_1)\cap\Gamma(B_2)$ is exactly
  the $x_{11}$ fibre; $B_1\sim B_2$ if and only if $s=1$.}
 \begin{proof}\pf\ By Definition~\ref{hc:def:graph}, membership
  in both neighbourhoods is intersection size one with each block;
  irreflexivity excludes $B_1,B_2$. Conversely the fibre excludes
  these two labels and gives both adjacencies. For the pair itself,
  \asm[2]{0} and Lemma~\ref{hc:lem:intersections} leave precisely
  $s=1$ for adjacency and $s=2$ for nonadjacency.~\qed\end{proof}
 \qedstep
 \begin{proof}\pf\ Lemma~\ref{hc:lem:intersections} makes the cases
  \stepref{hcCommonOne}, \stepref{hcCommonTwo} exhaustive.
  Step \stepref{hcCommonNeighbours} identifies their $x_{11}$
  counts with the asserted common-neighbour counts.~\qed\end{proof}
\end{proof}

\subsubsection{Connor's matrix and its rank}

\begin{namedresult}{Lemma}{(a scalar matrix plus a constant matrix)}\label{hc:lem:det}
 \lean{Matrix.det_smul_one_add_smul_J}
 \uses{hc:facts:det}
 \leavevmode\par
 \assume{\begin{pfenum}
  \item $U$ is a finite type with decidable equality, $n=|U|\ge1$.
  \item $\alpha,\beta\in\rats$ and $\alpha\ne0$.
 \end{pfenum}}
 \prove{$\det(\alpha I_U+\beta J_U)=\alpha^{n-1}(\alpha+n\beta)$.}
\end{namedresult}
\begin{proof}
 \pfsketch\ Factor out the nonzero scalar and apply the matrix
 determinant lemma from \ref{hc:facts:det} to a rank-one product.
 \step{hcDetFactor}{$\alpha I+\beta J=
  \alpha\bigl(I+(\tfrac\beta\alpha\Zj_U)\Zj_U^T\bigr)$.}
 \begin{proof}\pf\ Every entry of $\Zj_U\Zj_U^T$ is one.
  Distribute $\alpha$ and cancel $\alpha\ne0$ from \asm[2]{0}.~\qed\end{proof}
 \step{hcDetRankOne}{\new\ $\det\bigl(I+(\tfrac\beta\alpha\Zj_U)\Zj_U^T\bigr)
  =1+n\tfrac\beta\alpha$.}
 \begin{proof}\pf\ Apply
  \path{Matrix.det_one_add_replicateCol_mul_replicateRow}
  from \ref{hc:facts:det}, with singleton index \texttt{Unit},
  $\Zu=\tfrac\beta\alpha\Zj_U$, and $\Zv=\Zj_U$.
  The dot product sums $n$ copies of $\beta/\alpha$.~\qed\end{proof}
 \step{hcDetScale}{$\det(\alpha I+\beta J)=\alpha^{n-1}(\alpha+n\beta)$.}
 \begin{proof}\pf\ By \path{Matrix.det_smul} in \ref{hc:facts:det},
  steps \stepref{hcDetFactor} and \stepref{hcDetRankOne} give
  $\alpha^n(1+n\beta/\alpha)$.
  Use $n=(n-1)+1$ from \asm[1]{0} and cancel $\alpha\ne0$.~\qed\end{proof}
 \qedstep
 \begin{proof}\pf\ Step \stepref{hcDetScale}.~\qed\end{proof}
\end{proof}

\begin{namedresult}{Lemma}{(Connor identities and entries)}\label{hc:lem:Connor}
 \lean{Design.C_transpose}\lean{Design.C_apply}
 \lean{Design.C_mul_C}\lean{Design.C_mul_incMatrix_transpose}
 \uses{hc:def:design,hc:def:matrices,hc:def:graph,hc:lem:incidence,
 hc:lem:intersections,hc:facts:matrix}
 \leavevmode\par
 \assume{\begin{pfenum}
  \item $P,\cB$ are finite types with decidable equality.
  \item $D$ is a design with parameters $(15,5,21,7,2)$.
 \end{pfenum}}
 \prove{Over $\rats$, $C^T=C$, $C^2=35C$, $CN^T=0$, and
 \[
 C(B,B')=\begin{cases}
 10,&B=B',\\3,&B\sim B',\\-4,&B\ne B'\text{ and }B\not\sim B'.
 \end{cases}
 \]}
\end{namedresult}
\begin{proof}
 \pfsketch\ We will substitute the identities for $G$ into
 $C=35I-7G+10J$ and check its three entry types separately.
 \step{hcCSym}{$C^T=C$.}
 \begin{proof}\pf\ $I,J,G$ are symmetric by
  Lemma~\ref{hc:lem:incidence}; transpose the defining linear
  combination in Definition~\ref{hc:def:matrices}.~\qed\end{proof}
 \step{hcCDiag}{\assume{$B=B'$.}\prove{$C(B,B')=10$.}}
 \begin{proof}\pf\ Lemma~\ref{hc:lem:incidence} and
  Definition~\ref{hc:def:design} give $G(B,B)=5$, so
  Definition~\ref{hc:def:matrices} gives $35-7\cdot5+10=10$.~\qed\end{proof}
 \step{hcCAdj}{\assume{$B\sim B'$.}\prove{$C(B,B')=3$.}}
 \begin{proof}\pf\ Definition~\ref{hc:def:graph} gives distinct
  labels and intersection size one. Definition~\ref{hc:def:matrices}
  and Lemma~\ref{hc:lem:incidence} give $0-7\cdot1+10=3$.~\qed\end{proof}
 \step{hcCNonadj}{\assume{$B\ne B'$ and $B\not\sim B'$.}
  \prove{$C(B,B')=-4$.}}
 \begin{proof}\pf\ Lemma~\ref{hc:lem:intersections} and
  Definition~\ref{hc:def:graph} force intersection size two.
  Definition~\ref{hc:def:matrices} gives $0-7\cdot2+10=-4$.~\qed\end{proof}
 \step{hcCSquare}{$C^2=35C$.}
 \begin{proof}
  \pf\ In the following calculation $I,G,J$ are indexed by $\cB$:
  \begin{align*}
   C^2&=1225I+49G^2+100J^2-490G+700J-140GJ
     &&\text{[\ref{hc:lem:incidence}, $GJ=JG$]}\\
   &=1225I+49(5G+50J)-490G-2100J
     &&\text{[\ref{hc:lem:incidence}, arithmetic]}\\
   &=1225I-245G+350J &&\text{[ring arithmetic]}\\
   &=35C &&\text{[\ref{hc:def:matrices}].}
  \end{align*}\qed
 \end{proof}
 \step{hcCZero}{$CN^T=0$.}
 \begin{proof}
  \step{hcCOnes}{$N^TJ_P=5J_{\cB,P}$ and $J_{\cB}N^T=7J_{\cB,P}$.}
  \begin{proof}\pf\ Each entry in the first product is a column
   sum of $N$, and each in the second a row sum; apply
   Lemma~\ref{hc:lem:incidence} and \ref{hc:facts:matrix}.~\qed\end{proof}
  \qedstep
  \begin{proof}
   \pf\ By associativity from \ref{hc:facts:matrix},
   \begin{align*}
    CN^T&=35N^T-7N^T(NN^T)+10J_{\cB}N^T
      &&\text{[\ref{hc:def:matrices}]}\\
    &=35N^T-7N^T(5I_P+2J_P)+70J_{\cB,P}
      &&\text{[\ref{hc:lem:incidence}, \stepref{hcCOnes}]}\\
    &=0 &&\text{[\stepref{hcCOnes}, ring arithmetic].}
   \end{align*}\qed
  \end{proof}
 \end{proof}
 \qedstep
 \begin{proof}\pf\ Steps \stepref{hcCDiag}, \stepref{hcCAdj},
  \stepref{hcCNonadj} cover every pair of indices.
  Steps \stepref{hcCSym}, \stepref{hcCSquare}, \stepref{hcCZero}
  give the other conclusions.~\qed\end{proof}
\end{proof}

\begin{namedresult}{Lemma}{(quadratic inequality and rank bound)}\label{hc:lem:rank}
 \lean{Design.C_quadratic_nonneg}\lean{Design.rank_C_le_six}
 \uses{hc:def:design,hc:def:matrices,hc:lem:incidence,hc:lem:Connor,
 hc:lem:det,hc:facts:matrix,hc:facts:rank}
 \leavevmode\par
 \assume{\begin{pfenum}
  \item $P,\cB$ are finite types with decidable equality.
  \item $D$ is a design with parameters $(15,5,21,7,2)$.
 \end{pfenum}}
 \prove{For every $\Zx\in\rats^{\cB}$,
 $35\Zx^TC\Zx=(C\Zx)^T(C\Zx)\ge0$, so $\Zx^TC\Zx\ge0$;
 also $\operatorname{rk}C\le6$.}
\end{namedresult}
\begin{proof}
 \pfsketch\ The equation $C^2=35C$ gives a sum of squares. The
 equation $CN^T=0$ gives the rank bound after proving $N$ has rank fifteen.
 \step{hcRankQuad}{For every $\Zx$, $35\Zx^TC\Zx=(C\Zx)^T(C\Zx)\ge0$.}
 \begin{proof}\pf\ Lemma~\ref{hc:lem:Connor} gives $35C=C^TC$.
  Multiply by $\Zx^T,\Zx$, using \ref{hc:facts:matrix}; the resulting
  value is a sum of rational squares.~\qed\end{proof}
 \step{hcRankNonneg}{For every $\Zx$, $\Zx^TC\Zx\ge0$.}
 \begin{proof}\pf\ Divide \stepref{hcRankQuad} by $35>0$.~\qed\end{proof}
 \step{hcRankDet}{$\det(NN^T)=5^{14}\cdot35\ne0$.}
 \begin{proof}\pf\ First cast Lemma~\ref{hc:lem:incidence} to $\rats$
  by \ref{hc:facts:matrix}. Apply Lemma~\ref{hc:lem:det}.
  Definition~\ref{hc:def:design} gives
  $|P|=15$; take $\alpha=5\ne0$ and $\beta=2$.~\qed\end{proof}
 \step{hcRankN}{$\operatorname{rk}N=\operatorname{rk}N^T=15$.}
 \begin{proof}
  \step{hcRankFull}{$\operatorname{rk}(NN^T)=15$.}
  \begin{proof}\pf\ The determinant in \stepref{hcRankDet} is a
   unit in $\rats$. The unit determinant and full-rank criteria of
   \ref{hc:facts:rank} give rank $|P|=15$.~\qed\end{proof}
  \qedstep
  \begin{proof}\pf\ Apply the two transpose-rank identities of
   \ref{hc:facts:rank} over $\rats$ to \stepref{hcRankFull}.~\qed\end{proof}
 \end{proof}
 \step{hcRankBound}{$\operatorname{rk}C\le6$.}
 \begin{proof}\pf\ By $CN^T=0$ in Lemma~\ref{hc:lem:Connor} and
  the zero-product inequality of \ref{hc:facts:rank},
  $\operatorname{rk}C+\operatorname{rk}N^T\le|\cB|=21$.
  Substitute \stepref{hcRankN} and subtract fifteen.~\qed\end{proof}
 \qedstep
 \begin{proof}\pf\ Steps \stepref{hcRankQuad}, \stepref{hcRankNonneg},
  and \stepref{hcRankBound}.~\qed\end{proof}
\end{proof}

\begin{namedresult}{Lemma}{(principal matrices on labelled subsets)}\label{hc:lem:principal}
 \lean{Design.C_submatrix_quadratic_nonneg}
 \lean{Design.rank_C_submatrix_le_six}\lean{Design.det_C_submatrix_eq_zero}
 \uses{hc:def:design,hc:def:matrices,hc:lem:rank,hc:facts:finite,
 hc:facts:matrix,hc:facts:rank}
 \leavevmode\par
 \assume{\begin{pfenum}
  \item $P,\cB,U$ are finite types with decidable equality and
  $D$ is a design with parameters $(15,5,21,7,2)$.
  \item $e:U\to\cB$ is injective, and $C_e(u,w)\triangleq C(e(u),e(w))$.
 \end{pfenum}}
 \prove{For all $\Zx\in\rats^U$, $\Zx^TC_e\Zx\ge0$;
 $\operatorname{rk}C_e\le6$; if $|U|\ge7$, then $\det C_e=0$.}
\end{namedresult}
\begin{proof}
 \pfsketch\ We will extend vectors by zero along the injection.
 Rank monotonicity also applies to the specified submatrix, so later
 enumerations need no unstated transport of singularity.
 \step{hcPrinExtension}{For $\Zx\in\rats^U$, define
  $\widehat\Zx(e(u))=\Zx(u)$ and $\widehat\Zx(B)=0$ off the image.
  This defines $\widehat\Zx\in\rats^{\cB}$.}
 \begin{proof}\pf\ Injectivity in \asm[2]{0} gives at most one
  preimage; the two domains partition $\cB$.~\qed\end{proof}
 \step{hcPrinValue}{$\widehat\Zx^TC\widehat\Zx=\Zx^TC_e\Zx\ge0$.}
 \begin{proof}\pf\ Expand the quadratic sum using \ref{hc:facts:matrix}.
  Terms off the image vanish by \stepref{hcPrinExtension}; change both
  finite indices along $e$ using \ref{hc:facts:finite}. Apply
  Lemma~\ref{hc:lem:rank} to $\widehat\Zx$.~\qed\end{proof}
 \step{hcPrinRank}{$\operatorname{rk}C_e\le6$.}
 \begin{proof}\pf\ Rank monotonicity for submatrices in
  \ref{hc:facts:rank} bounds this rank by $\operatorname{rk}C$;
  use Lemma~\ref{hc:lem:rank}.~\qed\end{proof}
 \step{hcPrinSingular}{\assume{$|U|\ge7$.}\prove{$\det C_e=0$.}}
 \begin{proof}\pf\ If the determinant were nonzero, the full-rank
  criterion in \ref{hc:facts:rank} would give
  $\operatorname{rk}C_e=|U|\ge7$, contradicting \stepref{hcPrinRank}.~\qed\end{proof}
 \qedstep
 \begin{proof}\pf\ Steps \stepref{hcPrinValue}, \stepref{hcPrinRank},
  and \stepref{hcPrinSingular}.~\qed\end{proof}
\end{proof}

\subsubsection{Finite certificates and local restrictions}

\begin{namedresult}{Definition}{(the C19 matrices)}\label{hc:def:certificates}
 \lean{Computations.C19.clawMatrix}\lean{Computations.C19.fiveBlockMatrix}
 \lean{Computations.C19.Q1}\lean{Computations.C19.Q2}
 \lean{Computations.C19.K1}\lean{Computations.C19.K2}\uses{}
 The following matrices have integer entries; their rational casts
 are denoted by the same letters. Indices in the displays run from
 $1$ to the indicated order, corresponding to Lean's
 $\mathrm{Fin}\,n$ in increasing order.

\[
 Q_{\rm claw}=
% BEGIN MATRIX claw
 \begin{pmatrix}
 10&3&3&3\\3&10&-4&-4\\3&-4&10&-4\\3&-4&-4&10
 \end{pmatrix}
% END MATRIX claw
\]

\[
 Q_{\rm five}=
% BEGIN MATRIX five
 \begin{pmatrix}
 10&3&3&3&-4\\3&10&3&-4&3\\3&3&10&-4&3\\
 3&-4&-4&10&3\\-4&3&3&3&10
 \end{pmatrix}
% END MATRIX five
\]

\[
 Q_1=
% BEGIN MATRIX Q1
 \begin{pmatrix}
 10&3&3&3&-4&-4&-4\\
 3&10&3&-4&3&3&3\\
 3&3&10&-4&3&3&3\\
 3&-4&-4&10&-4&-4&-4\\
 -4&3&3&-4&10&3&3\\
 -4&3&3&-4&3&10&3\\
 -4&3&3&-4&3&3&10
 \end{pmatrix}
% END MATRIX Q1
\]

\[
 Q_2=
% BEGIN MATRIX Q2
 \begin{pmatrix}
 10&3&3&3&3&-4&-4\\
 3&10&3&-4&-4&3&3\\
 3&3&10&-4&-4&3&3\\
 3&-4&-4&10&3&-4&-4\\
 3&-4&-4&3&10&-4&-4\\
 -4&3&3&-4&-4&10&3\\
 -4&3&3&-4&-4&3&10
 \end{pmatrix}
% END MATRIX Q2
\]

\[
 K_1=
% BEGIN MATRIX K1
 \begin{pmatrix}
 20&-12&-12&-6&8&8&8\\
 -12&9&7&4&-5&-5&-5\\
 -12&7&9&4&-5&-5&-5\\
 -6&4&4&4&-2&-2&-2\\
 8&-5&-5&-2&5&3&3\\
 8&-5&-5&-2&3&5&3\\
 8&-5&-5&-2&3&3&5
 \end{pmatrix}
% END MATRIX K1
\]

\[
 K_2=
% BEGIN MATRIX K2
 \begin{pmatrix}
 17&-11&-11&-7&-7&6&6\\
 -11&9&7&5&5&-4&-4\\
 -11&7&9&5&5&-4&-4\\
 -7&5&5&5&3&-2&-2\\
 -7&5&5&3&5&-2&-2\\
 6&-4&-4&-2&-2&4&2\\
 6&-4&-4&-2&-2&2&4
 \end{pmatrix}
% END MATRIX K2
\]
\end{namedresult}

\begin{namedresult}{Computation}{(C19: quadratic values and inverse certificates)}\label{hc:comp:C19}
 \lean{Computations.C19.claw_quadratic}
 \lean{Computations.C19.five_block_quadratic}
 \lean{Computations.C19.left_inverse_Q1}\lean{Computations.C19.left_inverse_Q2}
 \uses{hc:def:certificates,hc:facts:matrix}
 \leavevmode\par
 \pflet{$\Zx\triangleq(-1,1,1,1)^T$ and
 $\Zy\triangleq(-1,1,1,1,-1)^T$.}
 \prove{$\Zx^TQ_{\rm claw}\Zx=-2$, $\Zy^TQ_{\rm five}\Zy=-4$, and
 $K_jQ_j=14I_7$ for each $j\in\{1,2\}$.}
\end{namedresult}
\begin{proof}
 \pfsketch\ These are finite integer certificates, suitable for
 kernel-checked arithmetic. The script \path{blueprint/checks/hallconnor/check_proof.py} reads every matrix
 from Definition~\ref{hc:def:certificates} and recomputes every entry.
 \step{hcCertClaw}{$\Zx^TQ_{\rm claw}\Zx=-2$.}
 \begin{proof}\pf\ From Definition~\ref{hc:def:certificates}, the
  diagonal sum is $40$ and the upper-triangle sum is $-9-12=-21$.
  The quadratic sum is $40+2(-21)=-2$, by \ref{hc:facts:matrix}.~\qed\end{proof}
 \step{hcCertFive}{$\Zy^TQ_{\rm five}\Zy=-4$.}
 \begin{proof}\pf\ From Definition~\ref{hc:def:certificates}, the
  diagonal sum is $50$. The upper-triangle row contributions are
  $-13,-4,-7,-3$, totalling $-27$. Thus the value is
  $50+2(-27)=-4$, by \ref{hc:facts:matrix}.~\qed\end{proof}
 \step{hcCertOne}{$K_1Q_1=14I_7$.}
 \begin{proof}\pf\ For all $u,w\in\{1,\ldots,7\}$, substitute the
  integer entries in Definition~\ref{hc:def:certificates} into
  $\sum_{i=1}^7K_1(u,i)Q_1(i,w)$, using \ref{hc:facts:matrix}.
  The seven diagonal sums are $14$ and all $42$ other sums are zero.
  For example the first row starts with
  $200-36-36-18-32-32-32=14$ and
  $60-120-36+24+24+24+24=0$.~\qed\end{proof}
 \step{hcCertTwo}{$K_2Q_2=14I_7$.}
 \begin{proof}\pf\ For all $u,w\in\{1,\ldots,7\}$, the sums
  $\sum_{i=1}^7K_2(u,i)Q_2(i,w)$ from
  Definition~\ref{hc:def:certificates} are $14$ on the diagonal
  and zero off it, by integer multiplication and \ref{hc:facts:matrix}.
  The first two sums are
  $170-33-33-21-21-24-24=14$ and
  $51-110-33+28+28+18+18=0$.~\qed\end{proof}
 \qedstep
 \begin{proof}\pf\ Steps \stepref{hcCertClaw}, \stepref{hcCertFive},
  \stepref{hcCertOne}, and \stepref{hcCertTwo}.~\qed\end{proof}
\end{proof}

\begin{namedresult}{Lemma}{(the two nonsingular matrices)}\label{hc:lem:inverse}
 \lean{Computations.C19.det_Q1_ne_zero}\lean{Computations.C19.det_Q2_ne_zero}
 \uses{hc:def:certificates,hc:comp:C19,hc:facts:matrix,hc:facts:det}
 \leavevmode\par
 \assume{$j\in\{1,2\}$, with $Q_j,K_j$ as in Definition~\ref{hc:def:certificates}.}
 \prove{$\det Q_j\ne0$ over $\rats$.}
\end{namedresult}
\begin{proof}
 \pfsketch\ The certified product has nonzero determinant.
 \step{hcInvProduct}{\new\ $(\det K_j)(\det Q_j)=14^7$.}
 \begin{proof}\pf\ Cast the integral identity of
  Computation~\ref{hc:comp:C19} to $\rats$ by \ref{hc:facts:matrix},
  then take determinants;
  multiplicativity and scalar multiplication in \ref{hc:facts:det}
  give $\det(14I_7)=14^7$. This replaces the ordinary proof's
  zero-kernel argument with the determinant of the same certificate.~\qed\end{proof}
 \qedstep
 \begin{proof}\pf\ If $\det Q_j=0$, the left side of
  \stepref{hcInvProduct} vanishes, whereas $14^7\ne0$ in $\rats$.~\qed\end{proof}
\end{proof}

\begin{namedresult}{Lemma}{(no induced claw)}\label{hc:lem:claw}
 \lean{Design.Z_clawFree}
 \uses{hc:def:design,hc:def:matrices,hc:def:graph,hc:def:certificates,hc:lem:Connor,
 hc:lem:principal,hc:comp:C19}
 \leavevmode\par
 \assume{\begin{pfenum}
  \item $P,\cB$ are finite types with decidable equality.
  \item $D$ is a design with parameters $(15,5,21,7,2)$.
 \end{pfenum}}
 \prove{There are no four distinct blocks $B,X,Y,W$ for which
 $B\sim X$, $B\sim Y$, $B\sim W$, and $X,Y,W$ are pairwise nonadjacent.}
\end{namedresult}
\begin{proof}
 \pfsketch\ The four specified labels would give a negative
 quadratic value of a principal matrix of $C$.
 \step{hcClawBad}{\assume{Four blocks $B,X,Y,W$ satisfy the forbidden pattern.}
  \prove{\textsc{False}.}}
 \begin{proof}
  \step{hcClawIndex}{The ordered list $(B,X,Y,W)$ defines an
   injection $e:\mathrm{Fin}\,4\to\cB$.}
  \begin{proof}\pf\ The blocks are distinct by \asm{1}.~\qed\end{proof}
  \step{hcClawDiag}{$C_e(i,i)=10$ for each $i$.}
  \begin{proof}\pf\ The diagonal clause of Lemma~\ref{hc:lem:Connor}.~\qed\end{proof}
  \step{hcClawCentre}{All entries from the centre to a leaf are $3$.}
  \begin{proof}\pf\ The three adjacencies in \asm{1} and the
   adjacency clause of Lemma~\ref{hc:lem:Connor}.~\qed\end{proof}
  \step{hcClawLeaves}{All entries between distinct leaves are $-4$.}
  \begin{proof}\pf\ Their distinctness and nonadjacency in \asm{1}
   give the last clause of Lemma~\ref{hc:lem:Connor}.~\qed\end{proof}
  \step{hcClawMatrix}{$C_e=Q_{\rm claw}$.}
  \begin{proof}\pf\ The diagonal, centre-leaf, and leaf-leaf entries
   are exactly those of Definition~\ref{hc:def:certificates}, by
   \stepref{hcClawDiag}, \stepref{hcClawCentre}, \stepref{hcClawLeaves}.
   Symmetry in Lemma~\ref{hc:lem:Connor} supplies the reversed entries.~\qed\end{proof}
  \qedstep
  \begin{proof}\pf\ Apply Lemma~\ref{hc:lem:principal} with
   \stepref{hcClawIndex} to the vector of Computation~\ref{hc:comp:C19}.
   Step \stepref{hcClawMatrix} makes its value $-2\ge0$, a contradiction.~\qed\end{proof}
 \end{proof}
 \qedstep
 \begin{proof}\pf\ Step \stepref{hcClawBad} negates the forbidden existence.~\qed\end{proof}
\end{proof}

\begin{namedresult}{Lemma}{(the five-block restriction)}\label{hc:lem:five}
 \textup{[\cite{HC}, Lemma 4.5 at $k=5$.]}
 \lean{Design.not_adj_of_five_block_pattern}
 \uses{hc:def:design,hc:def:matrices,hc:def:graph,hc:def:certificates,hc:lem:Connor,
 hc:lem:principal,hc:comp:C19}
 \leavevmode\par
 \assume{\begin{pfenum}
  \item $P,\cB$ are finite types with decidable equality and
  $D$ is a design with parameters $(15,5,21,7,2)$.
  \item $B_1,B_2,B_3,Y,W$ are five distinct blocks; $B_1,B_2,B_3$ form a triangle.
  \item $Y\sim B_1$, $Y\not\sim B_2,B_3$, and
  $W\sim B_2,B_3$, $W\not\sim B_1$.
 \end{pfenum}}
 \prove{$Y\not\sim W$.}
\end{namedresult}
\begin{proof}
 \pfsketch\ An edge $YW$ gives the second negative quadratic certificate.
 \step{hcFiveBad}{\assume{$Y\sim W$.}\prove{\textsc{False}.}}
 \begin{proof}
  \step{hcFiveIndex}{The ordered list $(B_1,B_2,B_3,Y,W)$ defines an
   injection $e:\mathrm{Fin}\,5\to\cB$.}
  \begin{proof}\pf\ Distinctness is \asm[2]{0}.~\qed\end{proof}
  \step{hcFiveDiag}{Every diagonal entry of $C_e$ is $10$.}
  \begin{proof}\pf\ Lemma~\ref{hc:lem:Connor}.~\qed\end{proof}
  \step{hcFiveTriangle}{The three entries between $B_1,B_2,B_3$ are $3$.}
  \begin{proof}\pf\ The triangle in \asm[2]{0} and
   Lemma~\ref{hc:lem:Connor}.~\qed\end{proof}
  \step{hcFiveY}{The entries from $Y$ to $(B_1,B_2,B_3)$ are $(3,-4,-4)$.}
  \begin{proof}\pf\ Apply the adjacency and nonadjacency clauses of
   Lemma~\ref{hc:lem:Connor} to the three stipulated relations in
   \asm[3]{0}; their labels are distinct by \asm[2]{0}.~\qed\end{proof}
  \step{hcFiveW}{The entries from $W$ to $(B_1,B_2,B_3)$ are $(-4,3,3)$.}
  \begin{proof}\pf\ Apply Lemma~\ref{hc:lem:Connor} to the three
   stipulated relations of $W$ in \asm[3]{0}, using \asm[2]{0}.~\qed\end{proof}
  \step{hcFiveYW}{$C(Y,W)=3$.}
  \begin{proof}\pf\ By \asm{1} and Lemma~\ref{hc:lem:Connor}.~\qed\end{proof}
  \step{hcFiveMatrix}{$C_e=Q_{\rm five}$.}
  \begin{proof}\pf\ Steps \stepref{hcFiveDiag}--\stepref{hcFiveYW}
   give every diagonal and upper-triangle entry in
   Definition~\ref{hc:def:certificates}; symmetry from
   Lemma~\ref{hc:lem:Connor} supplies the lower triangle.~\qed\end{proof}
  \qedstep
  \begin{proof}\pf\ Lemma~\ref{hc:lem:principal},
   \stepref{hcFiveIndex}, \stepref{hcFiveMatrix}, and
   Computation~\ref{hc:comp:C19} give $-4\ge0$, a contradiction.~\qed\end{proof}
 \end{proof}
 \qedstep
 \begin{proof}\pf\ Step \stepref{hcFiveBad} proves the negation.~\qed\end{proof}
\end{proof}

\subsubsection{The local six-cliques}

\begin{namedresult}{Lemma}{(splitting a neighbourhood)}\label{hc:lem:split}
 \lean{Design.exists_sixClique_neighborhood_split}
 \uses{hc:def:design,hc:def:matrices,hc:def:graph,hc:def:certificates,hc:lem:intersections,
 hc:lem:common,hc:lem:claw,hc:lem:five,hc:lem:Connor,hc:lem:principal,
 hc:lem:inverse,hc:facts:finite}
 \leavevmode\par
 \assume{\begin{pfenum}
  \item $P,\cB$ are finite types with decidable equality and
  $D$ is a design with parameters $(15,5,21,7,2)$.
  \item $B_1,B_2\in\cB$ and $B_1\sim B_2$.
 \end{pfenum}}
 \prove{There are $S_1,S_2\in\cS$ such that
 \[
 B_1,B_2\in S_1,\quad S_1\cap S_2=\{B_1\},\quad
 S_1\cup S_2=\{B_1\}\cup\Gamma(B_1).
 \]
 Moreover every $K\in\cS$ containing $B_1,B_2$ equals $S_1$.}
\end{namedresult}
\begin{proof}
 \pfsketch\ We will partition the other blocks by adjacency to $B_1,B_2$.
 The claw restriction gives two cliques and an exact edge count.
 A suitable third vertex then has a defect $j\le2$. The two positive
 defects give the certified nonsingular seven-block matrices, in
 conflict with the rank bound. Defect zero gives the split and uniqueness.
 \pflet{$\cM\triangleq\Gamma(B_1)\cap\Gamma(B_2)$,
 $\cU\triangleq\Gamma(B_1)\diff(\Gamma(B_2)\cup\{B_2\})$,
 $\cV\triangleq\Gamma(B_2)\diff(\Gamma(B_1)\cup\{B_1\})$, and
 $\cR\triangleq\cB\diff(\{B_1,B_2\}\cup\cM\cup\cU\cup\cV)$.}
 Write $\cF_{ij}\triangleq\{B\notin\{B_1,B_2\}:
 |B_1\cap B|=i,\ |B_2\cap B|=j\}$ for the fibres whose
 cardinalities are $x_{ij}$ in Lemma~\ref{hc:lem:common}.
 \step{hcLocalFibreM}{$\cM=\cF_{11}$.}
 \begin{proof}\pf\ Definition~\ref{hc:def:graph} identifies each
  adjacency with a one-point intersection and excludes the two
  endpoints by irreflexivity. Conversely the fibre gives both
  adjacencies. Lemma~\ref{hc:lem:intersections} places these
  intersections among the two possible sizes for distinct labels.~\qed\end{proof}
 \step{hcLocalFibreU}{$\cU=\cF_{12}$.}
 \begin{proof}\pf\ Membership in $\cU$ gives adjacency to $B_1$
  and excludes $B_2$ and adjacency to $B_2$.
  Definition~\ref{hc:def:graph} gives the first intersection size
  one; Lemma~\ref{hc:lem:intersections} forces the second to be two.
  Conversely these two sizes and the exclusions in $\cF_{12}$ give
  precisely the defining membership conditions of $\cU$.~\qed\end{proof}
 \step{hcLocalFibreV}{$\cV=\cF_{21}$.}
 \begin{proof}\pf\ Membership in $\cV$ gives adjacency to $B_2$
  and excludes $B_1$ and adjacency to $B_1$.
  Definition~\ref{hc:def:graph} gives the second size one;
  Lemma~\ref{hc:lem:intersections} forces the first to be two.
  Conversely the two sizes and endpoint exclusions in $\cF_{21}$
  give precisely the defining conditions of $\cV$.~\qed\end{proof}
 \step{hcLocalFibreR}{$\cR=\cF_{22}$.}
 \begin{proof}\pf\ Outside $\{B_1,B_2\}$,
  Lemma~\ref{hc:lem:intersections} leaves four possible pairs of
  sizes. Steps \stepref{hcLocalFibreM}--\stepref{hcLocalFibreV}
  remove the first three, leaving $(2,2)$. Conversely this pair
  gives no adjacency to either endpoint by Definition~\ref{hc:def:graph},
  so its label lies in none of $\cM,\cU,\cV$ and belongs to $\cR$.~\qed\end{proof}
 \step{hcLocalSizes}{$|\cM|=5$, $|\cU|=4$, $|\cV|=4$, and $|\cR|=6$.}
 \begin{proof}\pf\ Hypothesis \asm[2]{0} and
  Definition~\ref{hc:def:graph} give $s=1$.
  Apply the $s=1$ row of Lemma~\ref{hc:lem:common} to
  steps \stepref{hcLocalFibreM}--\stepref{hcLocalFibreR}.~\qed\end{proof}
 \step{hcLocalDisjoint}{$\cM,\cU,\cV,\cR$ are pairwise disjoint.
  Each is disjoint from $\{B_1,B_2\}$.}
 \begin{proof}\pf\ Their four different pairs of intersection
  sizes in steps \stepref{hcLocalFibreM}--\stepref{hcLocalFibreR}
  are mutually exclusive, and each fibre excludes the endpoints.~\qed\end{proof}
 \step{hcLocalCoverParts}{$\cM\cup\cU\cup\cV\cup\cR=
  \cB\diff\{B_1,B_2\}$.}
 \begin{proof}\pf\ By Lemma~\ref{hc:lem:intersections}, every label
  outside the endpoints has one of the four pairs of sizes.
  Use steps \stepref{hcLocalFibreM}--\stepref{hcLocalFibreR}.~\qed\end{proof}
 \step{hcLocalNeighboursOne}{$\Gamma(B_1)=\{B_2\}\cup\cM\cup\cU$,
  with pairwise disjoint parts.}
 \begin{proof}\pf\ The endpoint $B_2$ is a neighbour by \asm[2]{0}.
  Every other neighbour has first size one, so lies in $\cF_{11}$
  or $\cF_{12}$ by Lemma~\ref{hc:lem:intersections} and
  Definition~\ref{hc:def:graph}; both fibres consist of neighbours.
  Apply \stepref{hcLocalFibreM}, \stepref{hcLocalFibreU}, and
  \stepref{hcLocalDisjoint}.~\qed\end{proof}
 \step{hcLocalNeighboursTwo}{$\Gamma(B_2)=\{B_1\}\cup\cM\cup\cV$,
  with pairwise disjoint parts.}
 \begin{proof}\pf\ The endpoint $B_1$ is a neighbour by \asm[2]{0}.
  Every other neighbour has second size one, so lies in $\cF_{11}$
  or $\cF_{21}$ by Lemma~\ref{hc:lem:intersections} and
  Definition~\ref{hc:def:graph}; both fibres consist of neighbours.
  Apply \stepref{hcLocalFibreM}, \stepref{hcLocalFibreV}, and
  \stepref{hcLocalDisjoint}.~\qed\end{proof}
 \step{hcLocalU}{$\cU$ is a clique.}
 \begin{proof}
  \step{hcLocalUPair}{\assume{$Y,Y'\in\cU$, $Y\ne Y'$, and $Y\not\sim Y'$.}
   \prove{\textsc{False}.}}
  \begin{proof}\pf\ In \stepref{hcLocalFibreU}, $B_2,Y,Y'$ are
   distinct neighbours of $B_1$. Neither $Y$ nor $Y'$ is adjacent
   to $B_2$; \asm{2} excludes their mutual edge. This is the claw
   forbidden by Lemma~\ref{hc:lem:claw}.~\qed\end{proof}
  \qedstep
  \begin{proof}\pf\ Step \stepref{hcLocalUPair} excludes every
   missing edge between distinct members, as required by
   Definition~\ref{hc:def:graph}.~\qed\end{proof}
 \end{proof}
 \step{hcLocalV}{$\cV$ is a clique.}
 \begin{proof}
  \step{hcLocalVPair}{\assume{$W,W'\in\cV$, $W\ne W'$, and $W\not\sim W'$.}
   \prove{\textsc{False}.}}
  \begin{proof}\pf\ Step \stepref{hcLocalFibreV} makes
   $B_1,W,W'$ distinct neighbours of $B_2$; $B_1$ is adjacent to
   neither $W$ nor $W'$. With \asm{2} these are the three leaves
   of a claw, contradicting Lemma~\ref{hc:lem:claw}.~\qed\end{proof}
  \qedstep
  \begin{proof}\pf\ Step \stepref{hcLocalVPair} proves the clique
   predicate of Definition~\ref{hc:def:graph}.~\qed\end{proof}
 \end{proof}
 \step{hcLocalTwo}{For every $Y\in\cU$, $|\Gamma(Y)\cap\cM|=2$.}
 \begin{proof}
  \step{hcLocalYCommon}{For each $Y\in\cU$, the common neighbours of
   $B_1,Y$ are the disjoint union of $\cU\diff\{Y\}$ and
   $\Gamma(Y)\cap\cM$.}
  \begin{proof}\pf\ In the first neighbourhood decomposition of
   \stepref{hcLocalNeighboursOne}, $B_2\not\sim Y$ by membership in
   $\cU$; step \stepref{hcLocalU} gives every member of
   $\cU\diff\{Y\}$. These leave precisely the stated part of $\cM$.~\qed\end{proof}
  \qedstep
  \begin{proof}\pf\ Lemma~\ref{hc:lem:common} gives five common
   neighbours of the adjacent pair $B_1,Y$. The first part in
   \stepref{hcLocalYCommon} has size $4-1=3$ by
   \stepref{hcLocalSizes} and \ref{hc:facts:finite}, leaving two.~\qed\end{proof}
 \end{proof}
 \step{hcLocalEdgeSwap}{$\displaystyle
  \sum_{X\in\cM}|\Gamma(X)\cap\cU|
  =\sum_{Y\in\cU}|\Gamma(Y)\cap\cM|$.}
 \begin{proof}\pf\ Write each cardinality as a sum of adjacency
  indicators. Interchange the finite sums by \path{Finset.sum_comm}
  in \ref{hc:facts:finite}; symmetry of adjacency from
  Definition~\ref{hc:def:graph} gives the other sum.~\qed\end{proof}
 \step{hcLocalEight}{$\displaystyle
  \sum_{X\in\cM}|\Gamma(X)\cap\cU|=8$.}
 \begin{proof}\pf\ By \stepref{hcLocalEdgeSwap} the sum equals
  $\sum_{Y\in\cU}2$, using \stepref{hcLocalTwo}.
  Step \stepref{hcLocalSizes} gives $4\cdot2=8$.~\qed\end{proof}
 \step{hcLocalPick}{\pick\ $B_3\in\cM$ such that $|\Gamma(B_3)\cap\cU|\ge2$.}
 \begin{proof}\pf\ Otherwise each summand of \stepref{hcLocalEight}
  is at most one. Termwise summation gives
  $8\le\sum_{X\in\cM}1=|\cM|=5$, using
  \stepref{hcLocalSizes} and \ref{hc:facts:finite}, a contradiction.~\qed\end{proof}
 \step{hcLocalDefect}{\pflet{$j\triangleq|\cU\diff\Gamma(B_3)|$, and}
  \[
  h\triangleq|\Gamma(B_3)\cap(\cM\diff\{B_3\})|,\quad
  e\triangleq|\Gamma(B_3)\cap\cV|,\quad
  \ell\triangleq|\Gamma(B_3)\cap\cR|.
  \]
  Then $0\le j\le2$ and $|\Gamma(B_3)\cap\cU|=4-j$.}
 \begin{proof}\pf\ The neighbours and nonneighbours of $B_3$ partition
  the four-element set $\cU$, by \stepref{hcLocalSizes} and
  \ref{hc:facts:finite}; \stepref{hcLocalPick} supplies at least two neighbours.~\qed\end{proof}
 \step{hcLocalCounts}{$h=j$, $e=4-j$, and $\ell=j$.}
 \begin{proof}
  \step{hcLocalFirst}{$1+h+(4-j)=5$.}
  \begin{proof}\pf\ Count common neighbours of $B_1,B_3$ using
   Lemma~\ref{hc:lem:common} and \stepref{hcLocalNeighboursOne}:
   $B_2$ contributes one, $\cM\diff\{B_3\}$ contributes $h$, and
   $\cU$ contributes $4-j$ by \stepref{hcLocalDefect}.~\qed\end{proof}
  \step{hcLocalSecond}{$1+h+e=5$.}
  \begin{proof}\pf\ Count common neighbours of $B_2,B_3$ using
   Lemma~\ref{hc:lem:common}: $B_1$ contributes one, the part
   $\cM\diff\{B_3\}$ contributes $h$, and $\cV$ contributes $e$,
   by \stepref{hcLocalNeighboursTwo} and \stepref{hcLocalDefect}.~\qed\end{proof}
  \step{hcLocalDegree}{$2+h+(4-j)+e+\ell=10$.}
  \begin{proof}\pf\ Count all neighbours of $B_3$ in the partition
   \stepref{hcLocalCoverParts}, whose parts are disjoint by
   \stepref{hcLocalDisjoint}. Both $B_1,B_2$ are neighbours by
   \stepref{hcLocalPick}; the other contributions are the four
   counts in \stepref{hcLocalDefect}. The degree is ten by
   Lemma~\ref{hc:lem:intersections}.~\qed\end{proof}
  \qedstep
  \begin{proof}\pf\ Solve \stepref{hcLocalFirst} for $h$, then
   \stepref{hcLocalSecond} for $e$, then \stepref{hcLocalDegree}
   for $\ell$, over $\ints$.~\qed\end{proof}
 \end{proof}
 \step{hcLocalParts}{\pflet{$\cU_0\triangleq\cU\diff\Gamma(B_3)$ and
  $\cV_1\triangleq\cV\cap\Gamma(B_3)$.}
  They are disjoint cliques of sizes $j,4-j$, disjoint also from
  $\{B_1,B_2,B_3\}$.}
 \begin{proof}\pf\ Disjointness is \stepref{hcLocalDisjoint} and
  \stepref{hcLocalPick}; the sizes are \stepref{hcLocalDefect} and
  \stepref{hcLocalCounts}. They are subsets of the cliques in
  \stepref{hcLocalU} and \stepref{hcLocalV}.~\qed\end{proof}
 \step{hcLocalNoCross}{No member of $\cU_0$ is adjacent to a member of $\cV_1$.}
 \begin{proof}
  \step{hcLocalCrossData}{\assume{$Y\in\cU_0$ and $W\in\cV_1$.}
   \prove{$Y\not\sim W$.}}
  \begin{proof}
   \step{hcLocalFiveDistinct}{$B_1,B_2,B_3,Y,W$ are distinct.}
   \begin{proof}\pf\ Step \stepref{hcLocalParts} gives disjointness,
    and \stepref{hcLocalDisjoint}, \stepref{hcLocalPick} separate
    $B_3$ from the adjacent pair $B_1,B_2$.~\qed\end{proof}
   \step{hcLocalFiveTriangle}{$B_1,B_2,B_3$ form a triangle.}
   \begin{proof}\pf\ Use \asm[2]{0} and $B_3\in\cM$ in
    \stepref{hcLocalPick}.~\qed\end{proof}
   \step{hcLocalFiveY}{$Y\sim B_1$, $Y\not\sim B_2$, and $Y\not\sim B_3$.}
   \begin{proof}\pf\ Membership in $\cU_0$ in \asm{2}, and the
    definition in \stepref{hcLocalParts} and the fibre description
    in \stepref{hcLocalFibreU}.~\qed\end{proof}
   \step{hcLocalFiveW}{$W\sim B_2$, $W\sim B_3$, and $W\not\sim B_1$.}
   \begin{proof}\pf\ Membership in $\cV_1$ in \asm{2}, and the
    definition in \stepref{hcLocalParts} and the fibre description
    in \stepref{hcLocalFibreV}.~\qed\end{proof}
   \qedstep
   \begin{proof}\pf\ Apply Lemma~\ref{hc:lem:five} to
    \stepref{hcLocalFiveDistinct}--\stepref{hcLocalFiveW}.~\qed\end{proof}
  \end{proof}
  \qedstep
  \begin{proof}\pf\ Step \stepref{hcLocalCrossData} is quantified
   over every pair of the indicated parts.~\qed\end{proof}
 \end{proof}
 \step{hcLocalPattern}{On the named blocks $B_1,B_2,B_3$ and the
  members of $\cU_0,\cV_1$, $C$ is symmetric with diagonal $10$.
  Its off-diagonal entries are
  \[
  \begin{array}{c|ccc|cc}
       &B_1&B_2&B_3&\cU_0&\cV_1\\\hline
  B_1&*&3&3&3&-4\\
  B_2&3&*&3&-4&3\\
  B_3&3&3&*&-4&3\\\hline
  \cU_0&3&-4&-4&3&-4\\
  \cV_1&-4&3&3&-4&3
  \end{array}
  \]
  A set in this table denotes any of its members, with distinct
  blocks required for within-set entries.}
 \begin{proof}
  \step{hcPatternDiagonal}{For every block $X$, $C(X,X)=10$.}
  \begin{proof}\pf\ Lemma~\ref{hc:lem:Connor}.~\qed\end{proof}
  \step{hcPatternTriangle}{The three entries among $B_1,B_2,B_3$ are $3$.}
  \begin{proof}\pf\ The triangle follows from \asm[2]{0},
   \stepref{hcLocalFibreM}, and \stepref{hcLocalPick};
   use Lemma~\ref{hc:lem:Connor}.~\qed\end{proof}
  \step{hcPatternOneU}{For $Y\in\cU_0$, $C(B_1,Y)=3$.}
  \begin{proof}\pf\ Membership in $\cU$ gives adjacency to $B_1$
   by \stepref{hcLocalFibreU} and \stepref{hcLocalParts};
   apply Lemma~\ref{hc:lem:Connor}.~\qed\end{proof}
  \step{hcPatternTwoU}{For $Y\in\cU_0$, $C(B_2,Y)=-4$.}
  \begin{proof}\pf\ Step \stepref{hcLocalFibreU} excludes adjacency
   to $B_2$; distinctness in \stepref{hcLocalParts} permits
   Lemma~\ref{hc:lem:Connor}.~\qed\end{proof}
  \step{hcPatternThreeU}{For $Y\in\cU_0$, $C(B_3,Y)=-4$.}
  \begin{proof}\pf\ The definition and disjointness in
   \stepref{hcLocalParts} give nonadjacency and distinctness;
   apply Lemma~\ref{hc:lem:Connor}.~\qed\end{proof}
  \step{hcPatternOneV}{For $W\in\cV_1$, $C(B_1,W)=-4$.}
  \begin{proof}\pf\ Step \stepref{hcLocalFibreV} excludes adjacency
   to $B_1$; use the distinctness in \stepref{hcLocalParts} and
   Lemma~\ref{hc:lem:Connor}.~\qed\end{proof}
  \step{hcPatternTwoV}{For $W\in\cV_1$, $C(B_2,W)=3$.}
  \begin{proof}\pf\ Step \stepref{hcLocalFibreV} gives adjacency to
   $B_2$ for every member of the subset in \stepref{hcLocalParts};
   use Lemma~\ref{hc:lem:Connor}.~\qed\end{proof}
  \step{hcPatternThreeV}{For $W\in\cV_1$, $C(B_3,W)=3$.}
  \begin{proof}\pf\ The definition of $\cV_1$ in
   \stepref{hcLocalParts} gives adjacency; use Lemma~\ref{hc:lem:Connor}.~\qed\end{proof}
  \step{hcPatternUU}{For distinct $Y,Y'\in\cU_0$, $C(Y,Y')=3$.}
  \begin{proof}\pf\ The clique in \stepref{hcLocalParts} gives
   adjacency, so Lemma~\ref{hc:lem:Connor} applies.~\qed\end{proof}
  \step{hcPatternVV}{For distinct $W,W'\in\cV_1$, $C(W,W')=3$.}
  \begin{proof}\pf\ The clique in \stepref{hcLocalParts} gives
   adjacency; use Lemma~\ref{hc:lem:Connor}.~\qed\end{proof}
  \step{hcPatternUV}{For $Y\in\cU_0$ and $W\in\cV_1$, $C(Y,W)=-4$.}
  \begin{proof}\pf\ Step \stepref{hcLocalNoCross} excludes the edge;
   \stepref{hcLocalParts} gives distinctness, so
   Lemma~\ref{hc:lem:Connor} applies.~\qed\end{proof}
  \qedstep
  \begin{proof}\pf\ Steps \stepref{hcPatternDiagonal}--\stepref{hcPatternUV}
   cover the diagonal and all ten unordered pair types.
   Symmetry from Lemma~\ref{hc:lem:Connor} gives the reversed entries.~\qed\end{proof}
 \end{proof}
 \step{hcLocalNotOne}{\assume{$j=1$.}\prove{\textsc{False}.}}
 \begin{proof}
  \step{hcLocalEnumerateOne}{\pick\ $Y,W_1,W_2,W_3$ with
   $\cU_0=\{Y\}$, $\cV_1=\{W_1,W_2,W_3\}$, and the $W_i$ pairwise distinct.}
  \begin{proof}\pf\ By \asm{1} and \stepref{hcLocalParts}, the sizes
   are one and three. Apply \path{Finset.card_eq_one} and
   \path{Finset.card_eq_three} from \ref{hc:facts:finite}.~\qed\end{proof}
  \step{hcLocalInjectionOne}{Put
   $a\triangleq(B_1,B_2,B_3,Y,W_1,W_2,W_3):\mathrm{Fin}\,7\to\cB$.
   Then $a$ is injective.}
  \begin{proof}\pf\ The first three blocks are distinct by
   \asm[2]{0}, \stepref{hcLocalDisjoint}, and \stepref{hcLocalPick}.
   Step \stepref{hcLocalParts} separates them from the two disjoint
   parts, and \stepref{hcLocalEnumerateOne} separates the $W_i$.
   Thus the literal Lean vector \texttt{![B1, B2, B3, Y, W1, W2, W3]}
   has seven distinct entries.~\qed\end{proof}
  \step{hcLocalQOne}{$C_a=Q_1$.}
  \begin{proof}\pf\ Use the fixed list in \stepref{hcLocalInjectionOne}.
   In zero-based indices, the named blocks occupy $0,1,2$, $Y$ occupies
   $3$, and $W_1,W_2,W_3$ occupy $4,5,6$ by \stepref{hcLocalEnumerateOne}.
   For each of the $49$ pairs, the corresponding entry of
   \stepref{hcLocalPattern} equals the displayed entry of $Q_1$ in
   Definition~\ref{hc:def:certificates}: the diagonal is $10$,
   triangle and within-$W$ entries are $3$, the three entries from
   $Y$ to $B_1,B_2,B_3$ are $(3,-4,-4)$, each $W_i$ has
   $(-4,3,3)$ to those blocks, and all $Y,W_i$ entries are $-4$.
   These cases and their reverses exhaust the entries.~\qed\end{proof}
  \step{hcLocalSingularOne}{$\det C_a=0$.}
  \begin{proof}\pf\ Apply Lemma~\ref{hc:lem:principal} to
   \stepref{hcLocalInjectionOne}, whose domain has size seven.~\qed\end{proof}
  \qedstep
  \begin{proof}\pf\ Steps \stepref{hcLocalQOne} and
   \stepref{hcLocalSingularOne} contradict Lemma~\ref{hc:lem:inverse} at $j=1$.~\qed\end{proof}
 \end{proof}
 \step{hcLocalNotTwo}{\assume{$j=2$.}\prove{\textsc{False}.}}
 \begin{proof}
  \step{hcLocalEnumerateTwo}{\pick\ $Y_1,Y_2,W_1,W_2$ with
   $\cU_0=\{Y_1,Y_2\}$, $\cV_1=\{W_1,W_2\}$,
   $Y_1\ne Y_2$, and $W_1\ne W_2$.}
  \begin{proof}\pf\ The two sizes in \stepref{hcLocalParts} are
   both two by \asm{1}. Apply \path{Finset.card_eq_two} from
   \ref{hc:facts:finite} to each set.~\qed\end{proof}
  \step{hcLocalInjectionTwo}{Put
   $a\triangleq(B_1,B_2,B_3,Y_1,Y_2,W_1,W_2):\mathrm{Fin}\,7\to\cB$.
   Then $a$ is injective.}
  \begin{proof}\pf\ The initial blocks are distinct by
   \asm[2]{0}, \stepref{hcLocalDisjoint}, and \stepref{hcLocalPick}.
   The disjointness in \stepref{hcLocalParts} and the two inequalities
   in \stepref{hcLocalEnumerateTwo} separate all remaining pairs.
   This is the literal Lean vector \texttt{![B1, B2, B3, Y1, Y2, W1, W2]}.~\qed\end{proof}
  \step{hcLocalQTwo}{$C_a=Q_2$.}
  \begin{proof}\pf\ The list in \stepref{hcLocalInjectionTwo} puts
   $B_1,B_2,B_3$ at indices $0,1,2$, $Y_1,Y_2$ at $3,4$, and
   $W_1,W_2$ at $5,6$, using \stepref{hcLocalEnumerateTwo}.
   For each of the $49$ pairs, \stepref{hcLocalPattern} matches
   $Q_2$ in Definition~\ref{hc:def:certificates}: the diagonal is
   $10$, the triangle, $Y_1,Y_2$, and $W_1,W_2$ entries are $3$,
   each $Y_i$ has $(3,-4,-4)$ to the triangle, each $W_i$ has
   $(-4,3,3)$, and the four $Y_i,W_h$ entries are $-4$.
   The reverse entries agree by symmetry.~\qed\end{proof}
  \step{hcLocalSingularTwo}{$\det C_a=0$.}
  \begin{proof}\pf\ Lemma~\ref{hc:lem:principal} applies to
   \stepref{hcLocalInjectionTwo} with domain size seven.~\qed\end{proof}
  \qedstep
  \begin{proof}\pf\ Steps \stepref{hcLocalQTwo} and
   \stepref{hcLocalSingularTwo} contradict Lemma~\ref{hc:lem:inverse} at $j=2$.~\qed\end{proof}
 \end{proof}
 \step{hcLocalZero}{$j=0$; $B_3$ is adjacent to every member of
  $\cU$ and to no member of $\cM\diff\{B_3\}$.}
 \begin{proof}\pf\ The integer bound $0\le j\le2$ of
  \stepref{hcLocalDefect} and the exclusions
  \stepref{hcLocalNotOne}, \stepref{hcLocalNotTwo} force zero.
  The counts in \stepref{hcLocalDefect}, \stepref{hcLocalCounts}
  are then four and zero; use \ref{hc:facts:finite}.~\qed\end{proof}
 \step{hcLocalM}{$\cM\diff\{B_3\}$ is a clique of size four.}
 \begin{proof}
  \step{hcLocalMSize}{$|\cM\diff\{B_3\}|=4$.}
  \begin{proof}\pf\ Use $|\cM|=5$, $B_3\in\cM$ from
   \stepref{hcLocalSizes}, \stepref{hcLocalPick}, and
   deletion of a singleton in \ref{hc:facts:finite}.~\qed\end{proof}
  \step{hcLocalMClaw}{\assume{$X,Y\in\cM\diff\{B_3\}$ are distinct and nonadjacent.}
   \prove{\textsc{False}.}}
  \begin{proof}\pf\ The three distinct vertices $B_3,X,Y$ are
   neighbours of $B_1$ by \stepref{hcLocalFibreM} and
   \stepref{hcLocalPick}; $B_3$ has no edges to $X,Y$ by
   \stepref{hcLocalZero}. With \asm{2} they give a claw,
   contradicting Lemma~\ref{hc:lem:claw}.~\qed\end{proof}
  \qedstep
  \begin{proof}\pf\ Step \stepref{hcLocalMSize} gives the size and
   \stepref{hcLocalMClaw} gives every required edge.~\qed\end{proof}
 \end{proof}
 \step{hcLocalSOne}{Put $S_1\triangleq\{B_1,B_2\}\cup(\cM\diff\{B_3\})$.
  Then $S_1\in\cS$ and $B_1,B_2\in S_1$.}
 \begin{proof}
  \step{hcLocalSOneSize}{$|S_1|=6$.}
  \begin{proof}\pf\ The union is disjoint by
   \stepref{hcLocalDisjoint}, and \stepref{hcLocalM} gives $2+4=6$.~\qed\end{proof}
  \step{hcLocalSOneCore}{$B_1,B_2$ are adjacent.}
  \begin{proof}\pf\ Hypothesis \asm[2]{0}.~\qed\end{proof}
  \step{hcLocalSOneMixed}{Every member of $\cM\diff\{B_3\}$ is
   adjacent to each of $B_1,B_2$.}
  \begin{proof}\pf\ The definition of $\cM$ in \stepref{hcLocalFibreM}.~\qed\end{proof}
  \step{hcLocalSOneInternal}{Distinct members of $\cM\diff\{B_3\}$ are adjacent.}
  \begin{proof}\pf\ Step \stepref{hcLocalM}.~\qed\end{proof}
  \qedstep
  \begin{proof}\pf\ The three pair types in
   \stepref{hcLocalSOneCore}--\stepref{hcLocalSOneInternal} exhaust
   pairs in $S_1$. With \stepref{hcLocalSOneSize} this is membership
   in $\cS$, by Definition~\ref{hc:def:graph}; the two containments
   follow from the displayed union.~\qed\end{proof}
 \end{proof}
 \step{hcLocalSTwo}{Put $S_2\triangleq\{B_1,B_3\}\cup\cU$.
  Then $S_2\in\cS$.}
 \begin{proof}
  \step{hcLocalSTwoSize}{$|S_2|=6$.}
  \begin{proof}\pf\ The two-element set is disjoint from $\cU$
   by \stepref{hcLocalDisjoint} and \stepref{hcLocalPick}; its
   union has size $2+4=6$ by \stepref{hcLocalSizes}.~\qed\end{proof}
  \step{hcLocalSTwoCore}{$B_1\sim B_3$.}
  \begin{proof}\pf\ Step \stepref{hcLocalPick} puts $B_3\in\cM$.~\qed\end{proof}
  \step{hcLocalSTwoFirst}{Every member of $\cU$ is adjacent to $B_1$.}
  \begin{proof}\pf\ Step \stepref{hcLocalFibreU}.~\qed\end{proof}
  \step{hcLocalSTwoThird}{Every member of $\cU$ is adjacent to $B_3$.}
  \begin{proof}\pf\ Step \stepref{hcLocalZero}.~\qed\end{proof}
  \step{hcLocalSTwoInternal}{Every pair of distinct members of $\cU$ is adjacent.}
  \begin{proof}\pf\ Step \stepref{hcLocalU}.~\qed\end{proof}
  \qedstep
  \begin{proof}\pf\ Steps \stepref{hcLocalSTwoCore}--\stepref{hcLocalSTwoInternal}
   cover all pair types. Together with \stepref{hcLocalSTwoSize},
   they are the predicate of Definition~\ref{hc:def:graph}.~\qed\end{proof}
 \end{proof}
 \step{hcLocalMeet}{$S_1\cap S_2=\{B_1\}$.}
 \begin{proof}\pf\ In \stepref{hcLocalSOne} and \stepref{hcLocalSTwo},
  the parts $\{B_2\}$, $\cM\diff\{B_3\}$, $\{B_3\}$, and
  $\cU$ are pairwise disjoint by \stepref{hcLocalDisjoint} and
  \stepref{hcLocalPick}; only $B_1$ is shared.~\qed\end{proof}
 \step{hcLocalCover}{$S_1\cup S_2=\{B_1\}\cup\Gamma(B_1)$.}
 \begin{proof}\pf\ The two unions in \stepref{hcLocalSOne},
  \stepref{hcLocalSTwo} combine to $\{B_1,B_2\}\cup\cM\cup\cU$,
  since $B_3\in\cM$. Use the first neighbourhood decomposition
  of \stepref{hcLocalNeighboursOne}.~\qed\end{proof}
 \step{hcLocalUnique}{\assume{$K\in\cS$ and $B_1,B_2\in K$.}
  \prove{$K=S_1$.}}
 \begin{proof}
  \step{hcLocalKFour}{$K\diff\{B_1,B_2\}\subseteq\cM$ and
   $|K\diff\{B_1,B_2\}|=4$.}
  \begin{proof}\pf\ Every remaining vertex is adjacent to both
   initial blocks by Definition~\ref{hc:def:graph}. The two
   distinct blocks are in the six-element set $K$, so
   \ref{hc:facts:finite} gives size four.~\qed\end{proof}
  \step{hcLocalKExclude}{$B_3\notin K$.}
  \begin{proof}
   \step{hcLocalKBad}{\assume{$B_3\in K$.}\prove{\textsc{False}.}}
   \begin{proof}\pf\ By \stepref{hcLocalKFour}, removing $B_3$
    leaves three vertices in $\cM\diff\{B_3\}$, thus at least
    one, by \ref{hc:facts:finite}. The clique $K$ makes it adjacent
    to $B_3$, contradicting \stepref{hcLocalZero}.~\qed\end{proof}
   \qedstep
   \begin{proof}\pf\ Step \stepref{hcLocalKBad} proves the negation.~\qed\end{proof}
  \end{proof}
  \step{hcLocalKEqual}{$K\diff\{B_1,B_2\}=\cM\diff\{B_3\}$.}
  \begin{proof}\pf\ Steps \stepref{hcLocalKFour}, \stepref{hcLocalKExclude}
   give inclusion; both sets have size four by \stepref{hcLocalM}.
   Apply the subset and cardinality criterion in \ref{hc:facts:finite}.~\qed\end{proof}
  \qedstep
  \begin{proof}\pf\ Reinsert the two members $B_1,B_2$ into
   \stepref{hcLocalKEqual}; the resulting union is $S_1$ by
   \stepref{hcLocalSOne}.~\qed\end{proof}
 \end{proof}
 \qedstep
 \begin{proof}\pf\ Steps \stepref{hcLocalSOne}, \stepref{hcLocalSTwo},
  \stepref{hcLocalMeet}, \stepref{hcLocalCover}, and
  \stepref{hcLocalUnique} give the stated witnesses and uniqueness.~\qed\end{proof}
\end{proof}

\begin{namedresult}{Corollary}{(the unique six-clique through an edge)}\label{hc:cor:unique}
 \lean{Design.exists_unique_sixClique_of_adj}\lean{Design.sixClique_eq_of_adj}
 \uses{hc:def:design,hc:def:graph,hc:lem:split}
 \leavevmode\par
 \assume{\begin{pfenum}
  \item $P,\cB$ are finite types with decidable equality and
  $D$ is a design with parameters $(15,5,21,7,2)$.
  \item $B_1\sim B_2$.
 \end{pfenum}}
 \prove{There is a unique $S\in\cS$ containing $B_1,B_2$.
 In particular, if $K,L\in\cS$ both contain $B_1,B_2$, then $K=L$.}
\end{namedresult}
\begin{proof}
 \pfsketch\ Choose one local witness and compare every other
 six-clique through the edge with that same witness.
 \step{hcUniqueWitness}{\pick\ $S_1,S_2$ with the conclusions of
  Lemma~\ref{hc:lem:split} for $B_1,B_2$.}
 \begin{proof}\pf\ Apply that lemma to \asm[2]{0}.~\qed\end{proof}
 \step{hcUniqueExists}{There is a unique $S\in\cS$ containing $B_1,B_2$.}
 \begin{proof}\pf\ Step \stepref{hcUniqueWitness} gives the witness
  $S_1$ and states that every such $S$ equals $S_1$.~\qed\end{proof}
 \step{hcUniqueEqual}{\assume{$K,L\in\cS$ and $B_1,B_2\in K\cap L$.}
  \prove{$K=L$.}}
 \begin{proof}\pf\ The uniqueness clause of \stepref{hcUniqueWitness}
  gives $K=S_1$ and $L=S_1$. Transitivity gives $K=L$.~\qed\end{proof}
 \qedstep
 \begin{proof}\pf\ Steps \stepref{hcUniqueExists} and \stepref{hcUniqueEqual}.~\qed\end{proof}
\end{proof}

\begin{namedresult}{Corollary}{(exactly two six-cliques through a block)}\label{hc:cor:two}
 \lean{Design.card_sixCliques_through_block}
 \uses{hc:def:design,hc:def:graph,hc:lem:intersections,hc:lem:split,hc:cor:unique,hc:facts:finite}
 \leavevmode\par
 \assume{\begin{pfenum}
  \item $P,\cB$ are finite types with decidable equality and
  $D$ is a design with parameters $(15,5,21,7,2)$.
  \item $B_1\in\cB$.
 \end{pfenum}}
 \prove{There are distinct $S_1,S_2\in\cS$ with
 $\{S\in\cS:B_1\in S\}=\{S_1,S_2\}$ and
 $S_1\cap S_2=\{B_1\}$. In particular the filtered family has size two.}
\end{namedresult}
\begin{proof}
 \pfsketch\ We will use one local split, then apply the stated
 adjacent-pair uniqueness lemma to each possible six-clique.
 \step{hcTwoPick}{\pick\ $B_2\in\Gamma(B_1)$.}
 \begin{proof}\pf\ Lemma~\ref{hc:lem:intersections} gives ten
  neighbours; positive cardinality in \ref{hc:facts:finite} gives a witness.~\qed\end{proof}
 \step{hcTwoSplit}{\pick\ $S_1,S_2$ with all the conclusions of
  Lemma~\ref{hc:lem:split} for $B_1,B_2$.}
 \begin{proof}\pf\ Apply that lemma using \stepref{hcTwoPick}.~\qed\end{proof}
 \step{hcTwoDistinct}{$S_1\ne S_2$, and both contain $B_1$.}
 \begin{proof}\pf\ Their intersection is $\{B_1\}$ by
  \stepref{hcTwoSplit}. Equality would give a singleton of size
  six, contrary to Definition~\ref{hc:def:graph}.~\qed\end{proof}
 \step{hcTwoExhaust}{\assume{$K\in\cS$ and $B_1\in K$.}
  \prove{$K=S_1$ or $K=S_2$.}}
 \begin{proof}
  \step{hcTwoOther}{\pick\ $X\in K\diff\{B_1\}$.}
  \begin{proof}\pf\ This set has size five by \asm{1},
   Definition~\ref{hc:def:graph}, and \ref{hc:facts:finite}.~\qed\end{proof}
  \step{hcTwoEdge}{$B_1\sim X$ and
   $X\in(S_1\diff\{B_1\})\cup(S_2\diff\{B_1\})$.}
  \begin{proof}\pf\ The clique $K$ in \asm{1} and
   \stepref{hcTwoOther} give adjacency. The covering conclusion
   of \stepref{hcTwoSplit} gives the indicated membership.~\qed\end{proof}
  \step{hcTwoFirst}{\case{$X\in S_1$.} $K=S_1$.}
  \begin{proof}\pf\ Both six-cliques contain the adjacent pair
   $B_1,X$ by \asm{1}, \stepref{hcTwoDistinct}, and
   \stepref{hcTwoEdge}. Apply the uniqueness conclusion of
   Corollary~\ref{hc:cor:unique} to this pair.~\qed\end{proof}
  \step{hcTwoSecond}{\case{$X\in S_2$.} $K=S_2$.}
  \begin{proof}\pf\ The cliques $K,S_2$ contain $B_1,X$ by
   \asm{1}, \stepref{hcTwoDistinct}, and \stepref{hcTwoEdge}.
   Corollary~\ref{hc:cor:unique} applied to $B_1,X$ identifies them.~\qed\end{proof}
  \qedstep
  \begin{proof}\pf\ Step \stepref{hcTwoEdge} makes
   \stepref{hcTwoFirst}, \stepref{hcTwoSecond} exhaustive.~\qed\end{proof}
 \end{proof}
 \qedstep
 \begin{proof}\pf\ Step \stepref{hcTwoExhaust} gives containment
  of the filtered family in $\{S_1,S_2\}$; \stepref{hcTwoDistinct}
  gives the reverse containment and size two. The intersection
  is part of \stepref{hcTwoSplit}.~\qed\end{proof}
\end{proof}

\subsubsection{The global family}

\begin{namedresult}{Lemma}{(seven six-cliques)}\label{hc:lem:seven}
 \lean{Design.card_sixCliques}
 \uses{hc:def:design,hc:def:graph,hc:cor:two,hc:facts:finite}
 \leavevmode\par
 \assume{\begin{pfenum}
  \item $P,\cB$ are finite types with decidable equality.
  \item $D$ is a design with parameters $(15,5,21,7,2)$.
 \end{pfenum}}
 \prove{$|\cS|=7$.}
\end{namedresult}
\begin{proof}
 \pfsketch\ Count incidences between block labels and six-cliques.
 \step{hcSevenByClique}{There are $6|\cS|$ pairs $(B,S)$ with $B\in S\in\cS$.}
 \begin{proof}\pf\ Every member of $\cS$ has size six by
  Definition~\ref{hc:def:graph}; apply \ref{hc:facts:finite}.~\qed\end{proof}
 \step{hcSevenByBlock}{There are $42$ such pairs.}
 \begin{proof}\pf\ Each of the $21$ block labels belongs to two
  members of $\cS$ by Corollary~\ref{hc:cor:two} and
  Definition~\ref{hc:def:design}; use \ref{hc:facts:finite}.~\qed\end{proof}
 \qedstep
 \begin{proof}\pf\ Equate \stepref{hcSevenByClique} and
  \stepref{hcSevenByBlock}, then divide $6|\cS|=42$ by six.~\qed\end{proof}
\end{proof}

\begin{namedresult}{Lemma}{(intersections of six-cliques)}\label{hc:lem:clique-meet}
 \lean{Design.sixClique_inter_card}
 \uses{hc:def:design,hc:def:graph,hc:cor:unique,hc:cor:two,hc:lem:seven,hc:facts:finite}
 \leavevmode\par
 \assume{\begin{pfenum}
  \item $P,\cB$ are finite types with decidable equality and
  $D$ is a design with parameters $(15,5,21,7,2)$.
  \item $S_0,S_1\in\cS$ and $S_0\ne S_1$.
 \end{pfenum}}
 \prove{$|S_0\cap S_1|=1$.}
\end{namedresult}
\begin{proof}
 \pfsketch\ Any two distinct cliques share at most one block.
 The other clique through each of the six blocks of $S_0$ then
 enumerates all six cliques different from $S_0$.
 \step{hcMeetAtMost}{For all distinct $S,T\in\cS$, $|S\cap T|\le1$.}
 \begin{proof}
  \step{hcMeetTwo}{\assume{$S,T\in\cS$, $S\ne T$, and distinct $B,B'\in S\cap T$.}
   \prove{\textsc{False}.}}
  \begin{proof}\pf\ $B\sim B'$ by the clique predicate in
   Definition~\ref{hc:def:graph}. Corollary~\ref{hc:cor:unique}
   identifies $S,T$ through this edge, contradicting $S\ne T$.~\qed\end{proof}
  \qedstep
  \begin{proof}\pf\ The cardinality-at-most-one criterion of
   \ref{hc:facts:finite} applies to \stepref{hcMeetTwo}.~\qed\end{proof}
 \end{proof}
 \step{hcMeetMap}{For every $B\in S_0$ there is a unique
  $f(B)\in\cS\diff\{S_0\}$ containing $B$.}
 \begin{proof}\pf\ Corollary~\ref{hc:cor:two} supplies exactly
  two cliques through $B$, one of which is $S_0$ by \asm[2]{0}.
  Deleting that one leaves a singleton by \ref{hc:facts:finite}.~\qed\end{proof}
 \step{hcMeetInjective}{$f:S_0\to\cS\diff\{S_0\}$ is injective.}
 \begin{proof}\pf\ If $f(B)=f(B')$ with $B\ne B'$, then
  $S_0$ and $f(B)\ne S_0$ share two blocks by \stepref{hcMeetMap},
  contradicting \stepref{hcMeetAtMost}.~\qed\end{proof}
 \step{hcMeetSurjective}{$f$ is surjective.}
 \begin{proof}\pf\ Its domain has size six by
  Definition~\ref{hc:def:graph}; its codomain has size $7-1=6$
  by Lemma~\ref{hc:lem:seven} and \asm[2]{0}.
  The finite bijectivity criterion of \ref{hc:facts:finite}
  applies to \stepref{hcMeetInjective}.~\qed\end{proof}
 \step{hcMeetNonempty}{$S_0\cap S_1$ is nonempty.}
 \begin{proof}\pf\ By \asm[2]{0} and \stepref{hcMeetSurjective},
  choose $B\in S_0$ with $f(B)=S_1$. Step \stepref{hcMeetMap}
  gives $B\in S_1$.~\qed\end{proof}
 \qedstep
 \begin{proof}\pf\ Combine the lower bound from
  \stepref{hcMeetNonempty} with \stepref{hcMeetAtMost}, using
  \ref{hc:facts:finite}.~\qed\end{proof}
\end{proof}

\begin{namedresult}{Lemma}{(points in a six-clique)}\label{hc:lem:point-clique}
 \lean{Design.card_blocksThrough_in_sixClique}
 \uses{hc:def:design,hc:def:graph,hc:facts:finite}
 \leavevmode\par
 \assume{\begin{pfenum}
  \item $P,\cB$ are finite types with decidable equality and
  $D$ is a design with parameters $(15,5,21,7,2)$.
  \item $S_0\in\cS$.
 \end{pfenum}}
 \prove{For every $p\in P$, $|\{B\in S_0:p\in B\}|=2$.}
\end{namedresult}
\begin{proof}
 \pfsketch\ Two moments of the point multiplicities give a sum
 of squares equal to zero.
 \step{hcPointDefine}{\pflet{$m(p)\triangleq|\{B\in S_0:p\in B\}|$.}
  Then $\sum_{p\in P}m(p)=30$.}
 \begin{proof}\pf\ Count $(p,B)$ with $B\in S_0$ and $p\in B$.
  There are six labels, each with five points, by
  Definitions~\ref{hc:def:graph}, \ref{hc:def:design}, and
  \ref{hc:facts:finite}.~\qed\end{proof}
 \step{hcPointSecond}{$\sum_{p\in P}m(p)(m(p)-1)=30$ in $\ints$.}
 \begin{proof}\pf\ Count $(p,B,B')$ with distinct $B,B'\in S_0$
  and $p\in B\cap B'$. Each of the $6\cdot5$ ordered pairs
  meets in exactly one point by Definition~\ref{hc:def:graph}.
  Counting by $p$ gives the sum, using \ref{hc:facts:finite}.~\qed\end{proof}
 \step{hcPointSquares}{$\sum_{p\in P}(m(p)-2)^2=0$ in $\ints$.}
 \begin{proof}
  \pf\ Integer expansion gives
  \begin{align*}
   \sum_p(m(p)-2)^2
    &=\sum_p m(p)(m(p)-1)-3\sum_p m(p)+4|P|
        &&\text{[ring arithmetic]}\\
    &=30-90+60 &&\text{[\stepref{hcPointDefine}, \stepref{hcPointSecond},
                           \ref{hc:def:design}]}\\
    &=0 &&\text{[integer arithmetic].}
  \end{align*}\qed
 \end{proof}
 \step{hcPointZero}{For every $p\in P$, $(m(p)-2)^2=0$.}
 \begin{proof}\pf\ Each integer square is nonnegative; apply the
  zero-sum criterion of \ref{hc:facts:finite} to \stepref{hcPointSquares}.~\qed\end{proof}
 \qedstep
 \begin{proof}\pf\ An integer square vanishes only at zero, so
  \stepref{hcPointZero} gives $m(p)-2=0$ and $m(p)=2$ for every $p$.~\qed\end{proof}
\end{proof}

\subsubsection{The labelled extension and the matrix it supplies}

\begin{namedresult}{Definition}{(clique incidence and extended blocks)}\label{hc:def:extension}
 \lean{Design.sixCliqueIncMatrix}\lean{Design.hallConnorBlocks}
 \lean{Design.hallConnorIncMatrix}
 \uses{hc:def:design,hc:def:matrices,hc:def:graph}
 \leavevmode\par
 \assume{\begin{pfenum}
  \item $P,\cB$ are finite types with decidable equality.
  \item $D$ is a design with parameters $(15,5,21,7,2)$.
 \end{pfenum}}
 \pflet{$T\in\mat{\cB}{\cS}{\ints}$ have $T(B,S)=1$ for $B\in S$
 and $0$ otherwise. Put $P^*\triangleq P\sqcup\cS$ and
 $\cB^*\triangleq\operatorname{Option}(\cB)$.}
 Denote their injections by $\iota:P\to P^*$ and $\eta:\cS\to P^*$,
 the old labels by $\operatorname{some}(B)$, and the new label by $\operatorname{none}$.
 Define the block function $F:\cB^*\to\operatorname{Finset}(P^*)$ by
 \[
 \begin{split}
 F(\operatorname{some}(B))&\triangleq
    \iota(B)\cup\{\eta(S):S\in\cS,\ B\in S\},\\
 F(\operatorname{none})&\triangleq\{\eta(S):S\in\cS\}.
 \end{split}
 \]
 Here $\iota(B)$ is the image of the point set of $B$.
 Define $N^*\triangleq D.\mathrm{hallConnorIncMatrix}
 \in\mat{P^*}{\cB^*}{\ints}$ by the four entry formulas
 \[
 N^*(\iota(p),\operatorname{some}(B))\triangleq N(p,B),\quad N^*(\iota(p),\operatorname{none})\triangleq0,
 \]
 \[
 N^*(\eta(S),\operatorname{some}(B))\triangleq T(B,S),\quad N^*(\eta(S),\operatorname{none})\triangleq1.
 \]
\end{namedresult}

\begin{namedresult}{Lemma}{(parameters of the extended blocks)}\label{hc:lem:extension-parameters}
 \lean{Design.hallConnorBlocks_parameters}
 \uses{hc:def:design,hc:def:graph,hc:def:extension,hc:cor:two,hc:lem:seven,
 hc:lem:clique-meet,hc:lem:point-clique,hc:facts:finite}
 \leavevmode\par
 \assume{\begin{pfenum}
  \item $P,\cB$ are finite types with decidable equality.
  \item $D$ is a design with parameters $(15,5,21,7,2)$.
 \end{pfenum}}
 \prove{For the types and block function of Definition~\ref{hc:def:extension},
 \[
 \begin{gathered}
 |P^*|=22,\qquad|\cB^*|=22,\qquad
 \forall A\in\cB^*,\ |F(A)|=7,\\
 \forall x\in P^*,\ |\{A:x\in F(A)\}|=7,\\
 \forall x,y\in P^*,\ x\ne y\Rightarrow
                  |\{A:x,y\in F(A)\}|=2.
 \end{gathered}
 \]}
\end{namedresult}
\begin{proof}
 \pfsketch\ We will verify every cardinality and each of the three
 unordered types of point pair. The disjoint sum distinguishes new
 points from old points even if their underlying representations agree.
 \step{hcExtPoints}{$|P^*|=22$ and $|\cB^*|=22$.}
 \begin{proof}\pf\ Definition~\ref{hc:def:extension} gives the sums
  $|P|+|\cS|=15+7$ and $|\cB|+1=21+1$, using
  Definition~\ref{hc:def:design}, Lemma~\ref{hc:lem:seven}, and
  \path{Fintype.card_sum} and \path{Fintype.card_option} in
  \ref{hc:facts:finite}.~\qed\end{proof}
 \step{hcExtOldBlock}{For every $B\in\cB$, $|F(\operatorname{some}(B))|=7$.}
 \begin{proof}\pf\ There are five old points by
  Definition~\ref{hc:def:design} and two new points by
  Corollary~\ref{hc:cor:two}. Both $\iota$ and $\eta$ are injective,
  and their images are disjoint by Definition~\ref{hc:def:extension}.
  Apply \path{Finset.card_image_of_injective} and the disjoint-union
  count in \ref{hc:facts:finite} to get $5+2=7$.~\qed\end{proof}
 \step{hcExtNewBlock}{$|F(\operatorname{none})|=7$.}
 \begin{proof}\pf\ The map $\eta$ is injective, and
  Lemma~\ref{hc:lem:seven} gives seven members of $\cS$;
  use Definition~\ref{hc:def:extension} and
  \path{Finset.card_image_of_injective} in \ref{hc:facts:finite}.~\qed\end{proof}
 \step{hcExtOldReplication}{Each old point $\iota(p)$ lies in seven extended blocks.}
 \begin{proof}\pf\ Definition~\ref{hc:def:extension} preserves all
  incidences of $p$ with old labels and puts no old point in $F(\operatorname{none})$.
  Its old replication number is seven by Definition~\ref{hc:def:design}.~\qed\end{proof}
 \step{hcExtNewReplication}{Each new point $\eta(S)$ lies in seven extended blocks.}
 \begin{proof}\pf\ Definition~\ref{hc:def:extension} puts it in
  exactly the old labels $B\in S$, of which there are six by
  Definition~\ref{hc:def:graph}, and also in the distinct new
  label $\operatorname{none}$. The count is $6+1=7$.~\qed\end{proof}
 \step{hcExtOldOld}{\assume{$p,q\in P$ and $p\ne q$.}
  \prove{$\iota(p),\iota(q)$ lie together in exactly two extended blocks.}}
 \begin{proof}
  \step{hcExtOldOldOld}{The common old labels are precisely
   $\{\operatorname{some}(B):p,q\in B\}$.}
  \begin{proof}\pf\ The old-incidence clause of
   Definition~\ref{hc:def:extension}.~\qed\end{proof}
  \step{hcExtOldOldNew}{The new block contains neither old point.}
  \begin{proof}\pf\ The two injections have disjoint images in
   Definition~\ref{hc:def:extension}.~\qed\end{proof}
  \qedstep
  \begin{proof}\pf\ Step \stepref{hcExtOldOldOld} has cardinality two
   by Definition~\ref{hc:def:design} and \asm{1};
   \stepref{hcExtOldOldNew} excludes any further label.~\qed\end{proof}
 \end{proof}
 \step{hcExtOldNew}{\assume{$p\in P$ and $S\in\cS$.}
  \prove{$\iota(p),\eta(S)$ lie together in exactly two extended blocks.}}
 \begin{proof}
  \step{hcExtMixedOld}{The common old labels are precisely
   $\{\operatorname{some}(B):B\in S,\ p\in B\}$.}
  \begin{proof}\pf\ Use both old-label incidence clauses in
   Definition~\ref{hc:def:extension}.~\qed\end{proof}
  \step{hcExtMixedNew}{The new block does not contain $\iota(p)$.}
  \begin{proof}\pf\ Definition~\ref{hc:def:extension}.~\qed\end{proof}
  \qedstep
  \begin{proof}\pf\ Lemma~\ref{hc:lem:point-clique} gives two labels
   in \stepref{hcExtMixedOld}; \stepref{hcExtMixedNew} excludes
   a contribution from the new label.~\qed\end{proof}
 \end{proof}
 \step{hcExtNewOld}{\assume{$S\in\cS$ and $p\in P$.}
  \prove{$\eta(S),\iota(p)$ lie together in exactly two extended blocks.}}
 \begin{proof}\pf\ The set of common labels is unchanged when the
  two membership conditions are interchanged. Apply
  \stepref{hcExtOldNew} to $p,S$.~\qed\end{proof}
 \step{hcExtNewNew}{\assume{$S,T\in\cS$ and $S\ne T$.}
  \prove{$\eta(S),\eta(T)$ lie together in exactly two extended blocks.}}
 \begin{proof}
  \step{hcExtNewNewOld}{The common old labels are
   $\{\operatorname{some}(B):B\in S\cap T\}$, of size one.}
  \begin{proof}\pf\ Definition~\ref{hc:def:extension} gives this
   set, and Lemma~\ref{hc:lem:clique-meet} with \asm{1} gives its size.~\qed\end{proof}
  \step{hcExtNewNewNew}{Both points lie in the new block.}
  \begin{proof}\pf\ The clause for $F(\operatorname{none})$ in
   Definition~\ref{hc:def:extension}.~\qed\end{proof}
  \qedstep
  \begin{proof}\pf\ The old label of \stepref{hcExtNewNewOld}
   and the new label of \stepref{hcExtNewNewNew} are distinct,
   and together exhaust the common labels by
   Definition~\ref{hc:def:extension}.~\qed\end{proof}
 \end{proof}
 \qedstep
 \begin{proof}\pf\ Steps \stepref{hcExtPoints}--\stepref{hcExtNewReplication}
  verify the cardinalities, block sizes, and replication numbers.
  Steps \stepref{hcExtOldOld}, \stepref{hcExtOldNew},
  \stepref{hcExtNewOld}, \stepref{hcExtNewNew} exhaust ordered
  distinct point pairs: each point is in exactly one summand.
  Counting labels uses \path{Fintype.sum_option}, and splitting
  point sums uses \path{Fintype.sum_sum_type}, from
  \ref{hc:facts:finite}. These give every asserted equality.~\qed\end{proof}
\end{proof}

\begin{namedresult}{Definition}{(the Hall--Connor extension)}\label{hc:def:extension-design}
 \lean{Design.hallConnorExtension}
 \uses{hc:def:design,hc:def:extension,hc:lem:extension-parameters}
 \leavevmode\par
 \assume{\begin{pfenum}
  \item $P,\cB$ are finite types with decidable equality.
  \item $D$ is a design with parameters $(15,5,21,7,2)$.
 \end{pfenum}}
 Define $D^*\triangleq D.\mathrm{hallConnorExtension}$ of type
 \[
 \texttt{Design (Sum P D.sixCliques) (Option B) 22 7 22 7 2}
 \]
 with $D^*.\mathrm{block}\triangleq D.\mathrm{hallConnorBlocks}=F$
 from Definition~\ref{hc:def:extension}. The subtype of members
 is intended by \texttt{D.sixCliques} in the point type.
 The cardinality and incidence fields are exactly the equalities in
 Lemma~\ref{hc:lem:extension-parameters}; the positivity fields of
 Definition~\ref{hc:def:design} are $22>1$, $7>1$, and $22>0$.
\end{namedresult}

\begin{namedresult}{Theorem}{(the Hall--Connor residual identity)}\label{hc:thm:extension}
 \textup{[\cite{HC}, Theorems 2.1 and 3.2 at $k=5$.]}
 \lean{Design.hallConnorExtension_isSymmetric}
 \lean{Design.hallConnorExtension_residual}
 \uses{hc:def:design,hc:def:extension,hc:def:extension-design}
 \leavevmode\par
 \assume{\begin{pfenum}
  \item $P,\cB$ are finite types with decidable equality.
  \item $D$ is a design with parameters $(15,5,21,7,2)$.
 \end{pfenum}}
 \prove{The design $D^*=D.\mathrm{hallConnorExtension}$ is symmetric,
 and for every $B\in\cB$,
 \[
 D^*.\mathrm{block}(\operatorname{some}(B))\diff
 D^*.\mathrm{block}(\operatorname{none})=\iota(B).
 \]
 Thus deleting the new block and all its points recovers $D$ label by label.}
\end{namedresult}
\begin{proof}
 \pfsketch\ The parameter type gives symmetry; the two disjoint
 point images give the residual identity.
 \step{hcExtSymmetric}{$D^*$ is symmetric.}
 \begin{proof}\pf\ Its parameters in Definition~\ref{hc:def:extension-design}
  have $v=b=22$, which is Definition~\ref{hc:def:design}.~\qed\end{proof}
 \step{hcExtResidual}{For every $B$, the asserted block difference is $\iota(B)$.}
 \begin{proof}\pf\ Unfold the block function in
  Definition~\ref{hc:def:extension-design}. In
  Definition~\ref{hc:def:extension}, every added point belongs to
  $F(\operatorname{none})$ and no old point does. Removing that
  block from $F(\operatorname{some}(B))$ leaves exactly $\iota(B)$.~\qed\end{proof}
 \qedstep
 \begin{proof}\pf\ Steps \stepref{hcExtSymmetric} and \stepref{hcExtResidual}.~\qed\end{proof}
\end{proof}

\begin{namedresult}{Lemma}{(integral products for the extension)}\label{hc:lem:extension-products}
 \lean{Design.incMatrix_mul_sixCliqueIncMatrix}
 \lean{Design.sixCliqueIncMatrix_transpose_mul}
 \lean{Design.hallConnorIncMatrix_mul_transpose}
 \uses{hc:def:design,hc:def:matrices,hc:def:graph,hc:def:extension,hc:lem:incidence,
 hc:lem:clique-meet,hc:lem:point-clique,hc:facts:finite,hc:facts:matrix}
 \leavevmode\par
 \assume{\begin{pfenum}
  \item $P,\cB$ are finite types with decidable equality.
  \item $D$ is a design with parameters $(15,5,21,7,2)$.
 \end{pfenum}}
 \prove{All of the following are identities over $\ints$:
 \[
 NT=2J_{P,\cS},\qquad T^TT=5I_{\cS}+J_{\cS},\qquad
 N^*N^{*T}=5I_{P^*}+2J_{P^*}.
 \]}
\end{namedresult}
\begin{proof}
 \pfsketch\ We will expose the four matrix blocks so that the main
 contradiction can consume the integral identity directly.
 \step{hcProductsNT}{$NT=2J_{P,\cS}$.}
 \begin{proof}\pf\ By Definition~\ref{hc:def:extension} and
  \ref{hc:facts:matrix}, the $(p,S)$ entry counts $B\in S$ with
  $p\in B$. Lemma~\ref{hc:lem:point-clique} gives two.~\qed\end{proof}
 \step{hcProductsTTDiag}{$(T^TT)(S,S)=6$ for every $S\in\cS$.}
 \begin{proof}\pf\ This entry counts the six labels in $S$, by
  Definitions~\ref{hc:def:graph}, \ref{hc:def:extension}, and
  \ref{hc:facts:matrix}.~\qed\end{proof}
 \step{hcProductsTTOff}{\assume{$S,T_0\in\cS$ and $S\ne T_0$.}
  \prove{$(T^TT)(S,T_0)=1$.}}
 \begin{proof}\pf\ The entry counts $S\cap T_0$, by
  Definition~\ref{hc:def:extension} and \ref{hc:facts:matrix};
  Lemma~\ref{hc:lem:clique-meet} gives one.~\qed\end{proof}
 \step{hcProductsTT}{$T^TT=5I_{\cS}+J_{\cS}$.}
 \begin{proof}\pf\ The diagonal and off-diagonal cases
  \stepref{hcProductsTTDiag} and \stepref{hcProductsTTOff} are
  exhaustive and have the asserted entries $6$ and $1$.~\qed\end{proof}
 \step{hcProductsOO}{The old-old block of $N^*N^{*T}$ is $5I_P+2J_P$.}
 \begin{proof}\pf\ Sum over old labels and the new label using
  Definition~\ref{hc:def:extension}, \ref{hc:facts:matrix}, and
  \path{Fintype.sum_option} in \ref{hc:facts:finite}.
  The result is $NN^T+0$, equal to $5I_P+2J_P$ by
  Lemma~\ref{hc:lem:incidence}.~\qed\end{proof}
 \step{hcProductsON}{The old-new block of $N^*N^{*T}$ is $2J_{P,\cS}$.}
 \begin{proof}\pf\ The two label contributions are $NT$ and zero,
  by Definition~\ref{hc:def:extension}, \ref{hc:facts:matrix}, and
  \path{Fintype.sum_option} in \ref{hc:facts:finite};
  use \stepref{hcProductsNT}.~\qed\end{proof}
 \step{hcProductsNO}{The new-old block of $N^*N^{*T}$ is $2J_{\cS,P}$.}
 \begin{proof}\pf\ It is the transpose of the old-new block,
  since a matrix times its transpose is symmetric by
  \ref{hc:facts:matrix}; apply \stepref{hcProductsON}.~\qed\end{proof}
 \step{hcProductsNN}{The new-new block of $N^*N^{*T}$ is $5I_{\cS}+2J_{\cS}$.}
 \begin{proof}\pf\ Old labels contribute $T^TT$ and the new label
  contributes $J_{\cS}$, by Definition~\ref{hc:def:extension} and
  \ref{hc:facts:matrix}, using \path{Fintype.sum_option} in
  \ref{hc:facts:finite}. Step \stepref{hcProductsTT} gives their sum.~\qed\end{proof}
 \qedstep
 \begin{proof}\pf\ Steps \stepref{hcProductsOO}--\stepref{hcProductsNN}
  give the four blocks of $5I_{P^*}+2J_{P^*}$, because distinct
  summands never give a diagonal entry. Together with
  \stepref{hcProductsNT} and \stepref{hcProductsTT}, they prove all assertions.~\qed\end{proof}
\end{proof}

\begin{namedresult}{Theorem}{(Hall--Connor, matrix form)}\label{hc:thm:matrix}
 \lean{Design.hall_connor}
 \uses{hc:def:design,hc:def:extension,hc:lem:seven,hc:lem:extension-products,
 hc:facts:finite,hc:facts:matrix}
 \leavevmode\par
 \assume{\begin{pfenum}
  \item $P,\cB$ are finite types with decidable equality.
  \item $D$ is a design with parameters $(15,5,21,7,2)$.
 \end{pfenum}}
 \prove{There is $M\in\mat{\mathrm{Fin}\,22}{\mathrm{Fin}\,22}{\ints}$
 with $MM^T=5I_{22}+2J_{22}$.}
\end{namedresult}
\begin{proof}
 \pfsketch\ We will reindex both finite types in the integral
 identity, keeping the summed block index the same in both factors.
 \step{hcMatrixSizes}{$|P^*|=22$ and $|\cB^*|=22$.}
 \begin{proof}\pf\ Definition~\ref{hc:def:extension} gives disjoint
  sums of sizes $15+7$ and $21+1$, by
  Definition~\ref{hc:def:design}, Lemma~\ref{hc:lem:seven}, and
  \ref{hc:facts:finite}.~\qed\end{proof}
 \step{hcMatrixEquiv}{\pick\ equivalences
  $a:\mathrm{Fin}\,22\to P^*$ and $b:\mathrm{Fin}\,22\to\cB^*$.}
 \begin{proof}\pf\ Apply the enumeration fact in
  \ref{hc:facts:finite} to \stepref{hcMatrixSizes}.~\qed\end{proof}
 \step{hcMatrixDefine}{\pflet{$M(i,j)\triangleq N^*(a(i),b(j))$.}
  This is an integral matrix.}
 \begin{proof}\pf\ Every entry of $N^*$ is an integer by
  Definition~\ref{hc:def:extension}; composition preserves its scalar type.~\qed\end{proof}
 \step{hcMatrixProduct}{$(MM^T)(i,j)=(N^*N^{*T})(a(i),a(j))$ for all $i,j$.}
 \begin{proof}\pf\ Both factors use the bijection $b$ on the
  summed label index. Apply \path{Matrix.transpose_submatrix}
  and \path{Matrix.submatrix_mul_equiv} from \ref{hc:facts:matrix}
  to \stepref{hcMatrixDefine}, using \stepref{hcMatrixEquiv}.~\qed\end{proof}
 \step{hcMatrixTransport}{$MM^T=5I_{22}+2J_{22}$.}
 \begin{proof}\pf\ Lemma~\ref{hc:lem:extension-products} and
  \stepref{hcMatrixProduct} give entries
  $5\,1_{a(i)=a(j)}+2$. The injection $a$ in
  \stepref{hcMatrixEquiv} makes $a(i)=a(j)$ equivalent to $i=j$,
  so these are exactly the asserted entries.~\qed\end{proof}
 \qedstep
 \begin{proof}\pf\ The witness \stepref{hcMatrixDefine} and its
  equation \stepref{hcMatrixTransport} prove the existential conclusion.~\qed\end{proof}
\end{proof}

\subsubsection{The determinant obstruction and the conclusion}

\begin{namedresult}{Lemma}{(an odd power of five)}\label{hc:lem:odd-power}
 \lean{Computations.C19.not_sq_eq_49_mul_5_pow_odd}
 \lean{Computations.C19.not_sq_eq_49_mul_5_pow_21}
 \uses{hc:facts:prime}
 \leavevmode\par
 \assume{$m\in\nn$ and $d\in\ints$.}
 \prove{$d^2\ne49\cdot5^{2m+1}$. In particular $d^2\ne49\cdot5^{21}$.}
\end{namedresult}
\begin{proof}
 \pfsketch\ We will induct on $m$, universally quantifying the
 integer at each stage so that the induction hypothesis applies to $d/5$.
 \step{hcOddAll}{For every $m\in\nn$ and every $d\in\ints$,
  $d^2\ne49\cdot5^{2m+1}$.}
 \begin{proof}
  \step{hcOddBase}{\assume{$d\in\ints$ and $d^2=49\cdot5$.}
   \prove{\textsc{False}.}}
  \begin{proof}
   \step{hcOddBaseDivide}{\pick\ $e\in\ints$ with $d=5e$.}
   \begin{proof}\pf\ The equality in \asm{2} gives $5\mid d^2$;
    \ref{hc:facts:prime} gives $5\mid d$.~\qed\end{proof}
   \step{hcOddBaseCancel}{$5e^2=49$.}
   \begin{proof}\pf\ Substitute \stepref{hcOddBaseDivide} into
    \asm{2} and cancel the nonzero integer factor five.~\qed\end{proof}
   \qedstep
   \begin{proof}\pf\ Step \stepref{hcOddBaseCancel} would make
    $49$ divisible by five, whereas $49=5\cdot9+4$.~\qed\end{proof}
  \end{proof}
  \step{hcOddStep}{\assume{\begin{pfenum}
    \item $m\in\nn$ and $e^2\ne49\cdot5^{2m+1}$ for every $e\in\ints$.
    \item $d\in\ints$ and $d^2=49\cdot5^{2m+3}$.
   \end{pfenum}}
   \prove{\textsc{False}.}}
  \begin{proof}
   \step{hcOddStepDivide}{\pick\ $e\in\ints$ with $d=5e$.}
   \begin{proof}\pf\ The exponent $2m+3$ is positive, so
    \asm[2]{2} gives $5\mid d^2$. Apply
    \ref{hc:facts:prime}.~\qed\end{proof}
   \step{hcOddStepCancel}{$e^2=49\cdot5^{2m+1}$.}
   \begin{proof}\pf\ Substitute \stepref{hcOddStepDivide} into
    \asm[2]{2}, use $5^{2m+3}=25\cdot5^{2m+1}$, and cancel $25\ne0$.~\qed\end{proof}
   \qedstep
   \begin{proof}\pf\ Apply \asm[1]{2} to the integer $e$ in
    \stepref{hcOddStepCancel}.~\qed\end{proof}
  \end{proof}
  \qedstep
  \begin{proof}\pf\ Natural induction using \stepref{hcOddBase} and
   \stepref{hcOddStep}, since $2(m+1)+1=2m+3$.~\qed\end{proof}
 \end{proof}
 \qedstep
 \begin{proof}\pf\ Step \stepref{hcOddAll} proves the general assertion;
  its instance $m=10$ gives exponent $21$.~\qed\end{proof}
\end{proof}

\begin{namedresult}{Theorem}{(the Bruck--Ryser--Chowla instance)}\label{hc:thm:BRC}
 \lean{Matrix.bruck_ryser_chowla_twenty_two}
 \uses{hc:lem:det,hc:lem:odd-power,hc:facts:matrix,hc:facts:det}
 \leavevmode\par
 \assume{$M\in\mat{\mathrm{Fin}\,22}{\mathrm{Fin}\,22}{\ints}$.}
 \prove{$MM^T\ne5I_{22}+2J_{22}$.}
\end{namedresult}
\begin{proof}
 \pfsketch\ We will take determinants over $\rats$, retaining
 the integer determinant to apply the descent at five.
 \step{hcBRCBad}{\assume{$MM^T=5I_{22}+2J_{22}$.}\prove{\textsc{False}.}}
 \begin{proof}
  \step{hcBRCInteger}{Put $d\triangleq\det_{\ints}M$.
   Then $d\in\ints$ and $\det_{\rats}(M_{\rats})=(d:\rats)$.}
  \begin{proof}\pf\ The integer determinant is an integer by its
   scalar type; the cast identity is \path{Int.cast_det} in
   \ref{hc:facts:det}.~\qed\end{proof}
  \step{hcBRCSquare}{$(d:\rats)^2=49\cdot5^{21}$.}
  \begin{proof}
   \pf\ Cast \asm{1} to $\rats$ using \ref{hc:facts:matrix}.
   Since $5\ne0$ in $\rats$, the determinant lemma applies, giving
   \begin{align*}
    (d:\rats)^2&=\det(M_{\rats}M_{\rats}^T)
       &&\text{[\stepref{hcBRCInteger}, \ref{hc:facts:det}]}\\
     &=\det(5I_{22}+2J_{22}) &&\text{[\asm{1}]}\\
     &=5^{21}(5+22\cdot2) &&\text{[\ref{hc:lem:det}]}\\
     &=49\cdot5^{21} &&\text{[rational arithmetic].}
   \end{align*}\qed
  \end{proof}
  \step{hcBRCBack}{$d^2=49\cdot5^{21}$ in $\ints$.}
  \begin{proof}\pf\ The cast $\ints\to\rats$ is injective and
   preserves sums and products; apply it to \stepref{hcBRCSquare}.~\qed\end{proof}
  \qedstep
  \begin{proof}\pf\ Lemma~\ref{hc:lem:odd-power} contradicts
   \stepref{hcBRCBack}.~\qed\end{proof}
 \end{proof}
 \qedstep
 \begin{proof}\pf\ Step \stepref{hcBRCBad} proves the required negation.~\qed\end{proof}
\end{proof}

\begin{namedresult}{Theorem}{(the parameter contradiction)}\label{hc:thm:false}
 \lean{Design.false_of_parameters_15_5_2}
 \uses{hc:def:design,hc:thm:matrix,hc:thm:BRC}
 \leavevmode\par
 \assume{\begin{pfenum}
  \item $P,\cB$ are finite types with decidable equality.
  \item $D$ is a design with parameters $(15,5,21,7,2)$.
 \end{pfenum}}
 \prove{\textsc{False}.}
\end{namedresult}
\begin{proof}
 \pfsketch\ The matrix supplied by the embedding is forbidden by the obstruction.
 \step{hcFalseMatrix}{\pick\ an integer matrix $M$ of order $22$
  with $MM^T=5I_{22}+2J_{22}$.}
 \begin{proof}\pf\ Apply Theorem~\ref{hc:thm:matrix} to $D$.~\qed\end{proof}
 \qedstep
 \begin{proof}\pf\ Apply Theorem~\ref{hc:thm:BRC} to the witness
  and equation in \stepref{hcFalseMatrix}.~\qed\end{proof}
\end{proof}

\begin{namedresult}{Corollary}{(nonexistence of a $2$-$(15,5,2)$ design)}\label{hc:cor:main}
 \lean{Design.not_nonempty_15_5_2}\lean{Design.not_pair_counts_15_5_2}
 \uses{hc:def:design,hc:lem:parameters,hc:thm:false}
 \leavevmode\par
 \prove{The type
 \[\texttt{Design (Fin 15) (Fin 21) 15 5 21 7 2}\]
 is empty.
 Moreover, for every finite point type $P$ of size $15$ and every
 finite label type $\cB$, both with decidable equality, there is no
 function $F:\cB\to\operatorname{Finset}(P)$ whose blocks have size
 five and whose distinct point pairs lie in exactly two labels.}
\end{namedresult}
\begin{proof}
 \pfsketch\ We will state both the bundled and unbundled forms.
 \step{hcMainBundled}{The specified design type is empty.}
 \begin{proof}\pf\ Every inhabitant would give \textsc{False}
  by Theorem~\ref{hc:thm:false} with $P=\mathrm{Fin}\,15$,
  $\cB=\mathrm{Fin}\,21$.~\qed\end{proof}
 \step{hcMainUnbundled}{\assume{$P,\cB,F$ have the properties in the second assertion.}
  \prove{\textsc{False}.}}
 \begin{proof}
  \step{hcMainParameters}{\pick\ $D:\texttt{Design P B 15 5 21 7 2}$
   with $D.\mathrm{block}(B)=F(B)$ for every $B\in\cB$.}
  \begin{proof}\pf\ Apply Lemma~\ref{hc:lem:parameters} to \asm{1}.~\qed\end{proof}
  \qedstep
  \begin{proof}\pf\ Apply Theorem~\ref{hc:thm:false} to
   \stepref{hcMainParameters}.~\qed\end{proof}
 \end{proof}
 \qedstep
 \begin{proof}\pf\ Steps \stepref{hcMainBundled} and \stepref{hcMainUnbundled}.~\qed\end{proof}
\end{proof}

\begin{namedresult}{Definition}{(complementary labelled blocks)}\label{hc:def:compl}
 \lean{Design.compl}\uses{hc:def:design,hc:facts:finite}
 \leavevmode\par
 \assume{\begin{pfenum}
  \item $P,\cB$ are finite types with decidable equality.
  \item $D$ is a design with parameters $(v,k,b,r,\lambda)$.
  \item $2\le v-k$.
 \end{pfenum}}
 Define $D.\mathrm{compl}$ with the same point and label types by
 \[
 \begin{gathered}
 D.\mathrm{compl}:\texttt{Design P B v (v-k) b (b-r) (b+lam-2*r)},\\
 D.\mathrm{compl}.\mathrm{block}(B)\triangleq P\diff D(B).
 \end{gathered}
 \]
 The differences are natural differences. No injectivity of the
 block function is assumed. The proof below constructs its fields.
\end{namedresult}
\begin{proof}
 \pfsketch\ We will give the complement constructor its full
 parameter type. The bound on the new block size is exactly the
 condition needed by Definition~\ref{hc:def:design}.
 \step{hcComplGeneralSize}{Every complementary block has size $v-k$.}
 \begin{proof}\pf\ Its old block is a subset of $P$ with size $k$;
  subtract from $|P|=v$ using Definition~\ref{hc:def:design}
  and the finite-difference formula of \ref{hc:facts:finite}.~\qed\end{proof}
 \step{hcComplGeneralPoint}{Every point belongs to $b-r$ complementary labels.}
 \begin{proof}\pf\ The $r$ old labels through it form a subset of
  the $b$ labels. The other labels are exactly its complementary
  incidences, by Definition~\ref{hc:def:design} and \ref{hc:facts:finite}.~\qed\end{proof}
 \step{hcComplGeneralPair}{\new\ \assume{$p,q\in P$ and $p\ne q$.}
  \prove{Exactly $b+\lambda-2r$ complementary labels contain $p,q$.}}
 \begin{proof}
  \step{hcComplGeneralUnion}{Let $A_p,A_q$ be the sets of old
   incident labels of $p,q$. Then
   $|A_p\cup A_q|+\lambda=2r$.}
  \begin{proof}\pf\ Both sets have size $r$ and their intersection
   has size $\lambda$, by \asm{1} and Definition~\ref{hc:def:design}; use
   the union-intersection identity in \ref{hc:facts:finite}.~\qed\end{proof}
  \step{hcComplGeneralCount}{If $c$ counts labels containing
   neither point, then $c+2r=b+\lambda$.}
  \begin{proof}\pf\ The complement of $A_p\cup A_q$ among labels
   has size $c$, so $c+|A_p\cup A_q|=b$ by
   \ref{hc:facts:finite}. Add \stepref{hcComplGeneralUnion}.~\qed\end{proof}
  \qedstep
  \begin{proof}\pf\ Step \stepref{hcComplGeneralCount} gives
   $2r\le b+\lambda$ and $c=b+\lambda-2r$ in $\nn$.
   These are the common complementary labels by the block definition.
   This is the ordinary proof's inclusion-exclusion argument stated
   for general parameters, so the constructor has a precise type.~\qed\end{proof}
 \end{proof}
 \qedstep
 \begin{proof}\pf\ The types keep cardinalities $v,b$;
  steps \stepref{hcComplGeneralSize}, \stepref{hcComplGeneralPoint},
  and \stepref{hcComplGeneralPair} give the three incidence fields.
  The new block size is greater than one by \asm[3]{0}; $v>1,b>0$
  come from the old design in Definition~\ref{hc:def:design}.~\qed\end{proof}
\end{proof}

\begin{namedresult}{Corollary}{(the complementary parameters)}\label{hc:cor:complement}
 \lean{Design.false_of_parameters_15_10_9}\lean{Design.not_nonempty_15_10_9}
 \uses{hc:def:design,hc:def:compl,hc:thm:false}
 \leavevmode\par
 \assume{\begin{pfenum}
  \item $P,\cB$ are finite types with decidable equality.
  \item $D$ is a design with parameters $(15,10,21,14,9)$.
 \end{pfenum}}
 \prove{\textsc{False}. Thus
 \texttt{Design (Fin 15) (Fin 21) 15 10 21 14 9} is empty.}
\end{namedresult}
\begin{proof}
 \pfsketch\ Complementing every labelled block produces the forbidden parameters.
 \step{hcComplDesign}{Put $D'\triangleq D.\mathrm{compl}$.
  Then $D':\texttt{Design P B 15 5 21 7 2}$.}
 \begin{proof}\pf\ Apply the constructor of Definition~\ref{hc:def:compl}
  to \asm[2]{0}, with $2\le15-10=5$.
  Its parameters reduce to $15-10=5$, $21-14=7$, and
  $21+9-2\cdot14=2$. These are definitionally equal numerals in Lean.~\qed\end{proof}
 \qedstep
 \begin{proof}\pf\ Apply Theorem~\ref{hc:thm:false} to
  \stepref{hcComplDesign}. Since the argument applies to arbitrary
  finite $P,\cB$, it also excludes an inhabitant on
  $\mathrm{Fin}\,15$ and $\mathrm{Fin}\,21$.~\qed\end{proof}
\end{proof}

\subsubsection{Interfaces, provenance, and checks}

The statements above make the design-to-matrix bridge explicit:
\path{Design.hall_connor} returns the integer matrix equation that
\path{Matrix.bruck_ryser_chowla_twenty_two} negates. The constructor
\path{Design.hallConnorExtension} has the separately named theorem
\path{Design.hallConnorExtension_residual}, with point type \texttt{Sum P D.sixCliques} and label type
\texttt{Option B}. Here \path{D.sixCliques} denotes the subtype of
members when used as a type. The existential matrix theorem does
not need a general residual-design API.

No mathematical repair of the ordinary proof is needed.
We have made its moments, entry identities, principal-matrix bounds,
and clique incidence counts separately citable, put the determinant
lemma before its first use, and expanded the local and extension
cases. These are changes of proof structure. The ordinary
proof already proves the clique property before counting eight
edges, already defines the global family as all six-cliques, and
already uses rank to exclude $Q_1,Q_2$. Those repairs of the source
are retained here. Neither the unused rational determinant-sign
remark nor the unused source-local patterns are dependencies.
Three steps are tagged \new: using the rank-one determinant lemma,
taking the determinant of each certified inverse identity, and
stating the complement count for general parameters to give
\path{Design.compl} an exact interface. They retain the same
nonexistence argument.

All direct dependencies have local labels prefixed by
\path{hc:}; there are no deferred blueprint labels. The names in
\verb|\lean| are declarations of the Lean formalization, except for the
Mathlib API cited explicitly with paths. Mathlib
was inspected at commit
\par\noindent\path{082e2d37e8b0463410cdb532e111cd43d5a66174}\par\noindent
(4 September 2026), the version pinned in \path{lake-manifest.json}.
Its rank identities apply over $\rats$; the analytic theorem about
nonnegative determinants over an \path{RCLike} field is unnecessary.
The determinant obstruction uses only the even-order instance,
proved here without a quadratic-form local-global theorem.

The six displayed matrices are those of the ordinary proof. The exact
checks \path{blueprint/checks/hallconnor/check_proof.py}, with the $T(7)$
control \path{check_t7.py}, read them from this file and check all six
matrices, the two quadratic values, and every entry of both inverse
certificates; they are the checks of the ordinary proof in
\path{cgss/proofs/hallconnor/checks/}. The control
exercises all $210$ ordered adjacent pairs; no existing design can
exercise the hypothetical incidence structure. These checks do not
establish the quantified counting and clique arguments, which are
provided by the proofs above.
\end{mergedproofsection}

% Section 6 of the blueprint: Hadamard matrices.
% Local macros; macros.tex is shared and is not edited here.
\providecommand{\scl}{\opk{sc}}
\providecommand{\jp}{\Zj^\perp}
\providecommand{\agr}{\opk{ag}}

\section{Hadamard matrices}\label{sec:6}
\begingroup
\emergencystretch=3em
\raggedbottom

We will follow Section 6 of [CGSS]. Throughout, $n\geq1$ is an integer, the rows and the columns
of an $n\times n$ matrix are indexed by $\{1,\dots,n\}$, and $H$ is a Hadamard matrix of order
$n$: an integer matrix with entries $\pm1$ and $HH^T=H^TH=nI$. This is Mathlib's
\texttt{Matrix.IsHadamard} over $\ints$ (\path{Mathlib/LinearAlgebra/Matrix/HadamardMatrix.lean}),
whose fields are \texttt{apply\_mem}, \texttt{mul\_conjTranspose} and
\texttt{conjTranspose\_mul}; over $\ints$ the first says that every entry is $1$ or $-1$
(\path{Unitary.mem_iff_eq_one_or_eq_neg_one}). The paper writes $H_n$; we write $H$, and
$H_i\in\re^n$ for its row $i$, so that $H_i(c)=H(i,c)$. As in the rest of the blueprint,
$(\Zx,\Zy)$ is the standard inner product of $\re^n$
(\texttt{EuclideanSpace \ensuremath{\mathbb R} (Fin n)} in
Lean), $\Zj$ is the all-ones vector and $J$ the all-ones matrix. The paper's matrix $A$ with
$PP^T=nI+xA$ is our $C_H$ (\path{cgss/notation.md}, Section 7), since $A$ is an adjacency matrix
elsewhere and $C_H$ is the matrix $C$ of Theorem~\ref{thm:3.7} when $d=2$.

The paper's background paragraph asserts five things: the formula for $P^TP$, the quadratic
equation for $C_H$, the divisibility $x\mid n$, that the lines lie in $\jp$, and that the lines
through a fixed row are mutually orthogonal. We will prove each of them. The equations
$xAJ=-nJ$ and $AJ=-dJ$ are false in general (errata, Error 9), and we will not use them; the
divisibility comes instead from the count of fourth rows, which the paper also states. The
paragraph also gives the eigenvalues of $P^TP$ and of $PP^T$, and two claims about the Sylvester
matrices: that they satisfy Hypothesis 6.1 for every order $2^s$, and what their lines are. These
are side claims, which Theorem 6.2 does not use; we prove them in \path{LineGraphs/SideClaims/},
and we will record them in the remark after Lemma~\ref{lem:6-P} and in the last remark of this
section. Norman's theorem and the uniqueness of $H_{12}$ are side claims that the author has set
aside (\path{cgss/prerequisites.md}, Section 10). The examples of orders 2, 4, 8, 12 and 16
belong to computation C16, which Theorem 6.2 does not use either; we will record what it proves
in the first remark after the proof of Theorem 6.2.

The proof of Theorem 6.2 in the paper has two gaps. It assumes $n>9$ (errata, Gap 10), and it
asserts without proof that the lines form an indecomposable set (errata, Gap 17). We will fill
both. Only the orders with $16\mid n$ need an argument (Lemma~\ref{lem:6-small}). For these we
will split the lines into the components of their graph of non-orthogonality, and show that
each component of rank $m$ has exactly $\frac12nm$ lines (Lemma~\ref{lem:6-components}). The
classification of Theorem~\ref{thm:3.5} and the line counts of Lemma~\ref{lem:3.4} then leave
two possibilities (Lemma~\ref{lem:6-decomposition}): a single component, of type
$\typ A_{n-1}$ or $\typ D_{n-1}$, which the paper's own argument excludes
(Lemma~\ref{lem:6-orthogonal}); or $n=16$ and a component of type $\typ E_7$, which a count of
stars excludes (Lemma~\ref{lem:6-E7}). The last lemma is new.

The Lean names of Hypothesis 6.1 and Theorem 6.2, and of the matrices $P$ and $C_H$, are the
settled names of \path{cgss/notation.md}, Section 10, which also lists the other names of this
section.

\medskip\noindent\textbf{Inputs.}\label{input:sec6}
\begingroup\raggedright
We will use the following facts from other sections and from Mathlib.
\begin{itemize}
 \item \factlabel{LN3}{in6:LN3}(LN3 of \path{cgss/prerequisites.md}, deferred to the opening of
  Section 4.) Let $\Sigma$ be a finite set of vectors of squared norm 2 in a finite-dimensional
  real inner product space, with $\Sigma=-\Sigma$ and pairwise inner products in
  $\{0,\pm1,\pm2\}$. There is a smallest star-closed set $\scl(\Sigma)\supseteq\Sigma$ of vectors
  of squared norm 2; it is the vector set of a line system, it spans the same subspace as
  $\Sigma$, and it is indecomposable if the graph on $\Sigma$ in which two vectors are adjacent
  when their inner product is nonzero is connected. Lean: \path{LineSystem.starClosure}.
 \item \factlabel{Section 3}{in6:sec3}(Section 3.) Theorem~\ref{thm:3.5}, with the isomorphism
  induced by a surjective linear isometry of spans; the ranks of the models, $\ell$ for
  $\typ A_\ell$ and $\typ D_\ell$ and $6,7,8$ for $\typ E_6,\typ E_7,\typ E_8$
  (Lemma~\ref{lem:sec3-spans}); the line counts of Lemma~\ref{lem:3.4}; the coordinates of
  Definition~\ref{def:3.1}; and Theorem~\ref{thm:3.7} for $\typ E_7$, with
  Definition~\ref{def:sec3-signing}. An isomorphism of line systems is a bijection of vector sets
  that preserves inner products (Section~\ref{sec:2}), so it preserves orthogonality and the sets
  $N_\Zr=\{\Zu\in\Phi:(\Zu,\Zr)=1\}$. Lean: \path{LineSystem.classification},
  \path{LineSystem.RootType.rank},
  \leancite{LineSystem.numLines_E₇}{\texttt{LineSystem.numLines\_E\ensuremath{{}_7}}}
  and the other counts,
  \leancite{LineSystem.signedAdj_quadratic_E₇}%
   {\texttt{LineSystem.signedAdj\_quadratic\_E\ensuremath{{}_7}}};
  and, for $|N_\Zr|=32$ in the model,
  \leancite{LineSystem.E₇.card_N}{\texttt{LineSystem.E\ensuremath{{}_7}.card\_N}}
  (C20, \path{LineGraphs/Computations/C20E7.lean}).
 \item \factlabel{Mathlib}{in6:mathlib}(Mathlib.) \path{Matrix.IsHadamard.four_dvd_card}: an
  integer Hadamard matrix of order greater than 2 has order divisible by 4.
  \path{linearIndependent_of_ne_zero_of_inner_eq_zero}
  (\path{Mathlib/Analysis/InnerProductSpace/Basic.lean}): nonzero pairwise orthogonal vectors
  are linearly independent. \path{LinearIndependent.fintype_card_le_finrank}
  (\path{Mathlib/LinearAlgebra/Dimension/Finite.lean}),
  \path{LinearIndependent.span_eq_top_of_card_eq_finrank}
  (\path{Mathlib/LinearAlgebra/FiniteDimensional/Lemmas.lean}), \path{finrank_span_eq_card}
  (\path{Mathlib/LinearAlgebra/Dimension/Constructions.lean}) and
  \path{finrank_euclideanSpace_fin}. \path{ext_inner_left} and \path{inner_self_eq_zero}.
  \path{Nat.pow_dvd_pow_iff}: for $k\neq0$, $a^k\mid b^k$ if and only if $a\mid b$.
  \path{Nat.Coprime.dvd_of_dvd_mul_left}. The identities of Section~\ref{facts:inner}.
\end{itemize}
\endgroup

\subsection{The Hadamard products of pairs of rows}

\begin{namedresult}{Definition}{(Hadamard products and agreements)}\label{def:6-pairs}
 \lean{Matrix.hadamardPairs}\lean{Matrix.fourRowAgreeCount}\uses{}
 \assume{$H$ is an $n\times n$ Hadamard matrix.}
 \pflet{$\cP$ be the set of 2-subsets of $\{1,\dots,n\}$, so $|\cP|=\frac12n(n-1)$. For
 $\{i,j\}\in\cP$ let $\Zp_{ij}=H_i\circ H_j\in\re^n$, the \emph{Hadamard product} of the rows,
 with $\Zp_{ij}(c)=H(i,c)H(j,c)$. Let $P$ be the $\cP\times n$ matrix with rows $\Zp_{ij}$. For
 distinct rows $i,j,k,l$ let $\agr(i,j,k,l)$ be the number of columns $c$ with
 $H(i,c)=H(j,c)=H(k,c)=H(l,c)$; we say that the four rows \emph{agree} in these places.}
 Thus $\Zp_{ij}=\Zp_{ji}$, and $\agr$ is symmetric in its four arguments. In Lean the index
 type of $\cP$ is a proposal for the declaration \texttt{Matrix.hadamardPairs}; the pairs
 $(i,j)$ with $i<j$ will do.
\end{namedresult}

\begin{namedresult}{Lemma}{(the inner products of the Hadamard products)}\label{lem:6-gram}
 \lean{Matrix.IsHadamard.inner_hadamardPairs_self}
 \lean{Matrix.IsHadamard.inner_hadamardPairs_of_card_inter_eq_one}
 \lean{Matrix.IsHadamard.inner_hadamardPairs_of_disjoint}
 \lean{Matrix.IsHadamard.inner_hadamardPairs_one}
 \uses{def:6-pairs,in6:mathlib}
 \assume{\begin{pfenum}
  \item $H$ is a Hadamard matrix of order $n$.
  \item $\Zp_{ij}$, $\cP$ and $\agr$ are as in Definition~\ref{def:6-pairs}.
 \end{pfenum}}
 \prove{\begin{pfenum}
  \item Every entry of every $\Zp_{ij}$ is $\pm1$, and $(\Zp_{ij},\Zp_{ij})=n$.
  \item If $j\ne k$ and $i\notin\{j,k\}$, then $(\Zp_{ij},\Zp_{ik})=0$.
  \item If $i,j,k,l$ are distinct, then $(\Zp_{ij},\Zp_{kl})=8\agr(i,j,k,l)-n$.
  \item $(\Zp_{ij},\Zj)=0$ for every $\{i,j\}\in\cP$.
 \end{pfenum}}
\end{namedresult}
\begin{proof}
 \pfsketch\ Each part is a sum over the columns. Part 3 is the one computation; it rests on an
 identity for four signs.
 \step{s6gorth}{For distinct rows $u,w$, $\sum_cH(u,c)H(w,c)=0$.}
 \begin{proof}\pf\ This is the $(u,w)$ entry of $HH^T=nI$.~\qed\end{proof}
 \step{s6gsq}{For all $u$ and $c$, $H(u,c)^2=1$.}
 \begin{proof}\pf\ The entries are $\pm1$.~\qed\end{proof}
 \step{s6gone}{Part 1 holds.}
 \begin{proof}\pf\ $\Zp_{ij}(c)$ is a product of two signs, and
  $(\Zp_{ij},\Zp_{ij})=\sum_cH(i,c)^2H(j,c)^2=n$ by \stepref{s6gsq}.~\qed\end{proof}
 \step{s6gtwo}{Part 2 holds.}
 \begin{proof}\pf\ $(\Zp_{ij},\Zp_{ik})=\sum_cH(i,c)^2H(j,c)H(k,c)=\sum_cH(j,c)H(k,c)$ by
  \stepref{s6gsq}, and this is 0 by \stepref{s6gorth}, since $j\neq k$.~\qed\end{proof}
 \step{s6gthree}{Part 3 holds.}
 \begin{proof}
  \step{s6gsigns}{For $s_1,s_2,s_3,s_4\in\{1,-1\}$,
   \[
    8\,[s_1=s_2=s_3=s_4]=1+\sum_{1\leq a<b\leq4}s_as_b+s_1s_2s_3s_4,
   \]
   where $[\,\cdot\,]$ is 1 if the condition holds and 0 otherwise.}
  \begin{proof}\pf\ A check of the 16 sign vectors (in Lean, \texttt{decide}). If all four
   are equal, the right side is $1+6+1$; if exactly one differs, it is $1+(3-3)-1$; if two
   differ from the other two, it is $1+(2-4)+1$.~\qed\end{proof}
  \step{s6gsum}{$8\agr(i,j,k,l)=n+\sum_{\{u,w\}}\sum_cH(u,c)H(w,c)+(\Zp_{ij},\Zp_{kl})$,
   where $\{u,w\}$ runs over the six 2-subsets of $\{i,j,k,l\}$.}
  \begin{proof}\pf\ Sum \stepref{s6gsigns} over the columns $c$, with
   $s=(H(i,c),H(j,c),H(k,c),H(l,c))$; the left side sums to $8\agr(i,j,k,l)$ by
   Definition~\ref{def:6-pairs}, and $\sum_cH(i,c)H(j,c)H(k,c)H(l,c)=(\Zp_{ij},\Zp_{kl})$.~\qed
  \end{proof}
  \qedstep
  \begin{proof}\pf\ The six inner sums in \stepref{s6gsum} vanish by \stepref{s6gorth}, since
   the four rows are distinct.~\qed\end{proof}
 \end{proof}
 \step{s6gfour}{Part 4 holds.}
 \begin{proof}\pf\ $(\Zp_{ij},\Zj)=\sum_cH(i,c)H(j,c)=0$ by \stepref{s6gorth}.~\qed\end{proof}
 \qedstep
 \begin{proof}\pf\ Steps \stepref{s6gone}, \stepref{s6gtwo}, \stepref{s6gthree} and
  \stepref{s6gfour} are the four parts.~\qed\end{proof}
\end{proof}

\begin{namedresult}{Lemma}{(the matrix $P$)}\label{lem:6-P}
 \lean{Matrix.IsHadamard.transpose_hadamardPairs_mul}
 \lean{Matrix.IsHadamard.hadamardPairs_mulVec_one}
 \lean{Matrix.IsHadamard.hadamardPairs_gram_sq}
 \uses{def:6-pairs,lem:6-gram}
 \assume{$H$ is an $n\times n$ Hadamard matrix, $P$ is as in Definition~\ref{def:6-pairs}, and
 $G=PP^T$.}
 \prove{\begin{pfenum}
  \item $P^TP=\frac12n^2I-\frac12nJ$.
  \item $P\Zj=0$.
  \item $G^2=\frac12n^2G$.
 \end{pfenum}}
\end{namedresult}
\begin{proof}
 \pfsketch\ Part 1 is the paper's display; we compute an entry as a sum over pairs of rows,
 through the square of a sum. Part 3 follows from parts 1 and 2 by associativity.
 \step{s6Pone}{Part 1 holds.}
 \begin{proof}
  \pflet{$c,c'$ be columns, and $u_i=H(i,c)H(i,c')$ for each row $i$.}
  \step{s6Pentry}{$(P^TP)(c,c')=\sum_{\{i,j\}\in\cP}u_iu_j$.}
  \begin{proof}\pf\ By Definition~\ref{def:6-pairs}, we have that
   \[\Zp_{ij}(c)\Zp_{ij}(c')=H(i,c)H(i,c')H(j,c)H(j,c').~\qed\]
  \end{proof}
  \step{s6Psquare}{$2\sum_{\{i,j\}\in\cP}u_iu_j=\bigl(\sum_iu_i\bigr)^2-\sum_iu_i^2$.}
  \begin{proof}\pf\ Expand the square of the sum.~\qed\end{proof}
  \step{s6Psums}{$\sum_iu_i=(H^TH)(c,c')$, which is $n$ if $c=c'$ and 0 otherwise, and
   $\sum_iu_i^2=n$.}
  \begin{proof}\pf\ $H^TH=nI$, and each $u_i$ is $\pm1$.~\qed\end{proof}
  \qedstep
  \begin{proof}\pf\ By \stepref{s6Pentry}, \stepref{s6Psquare} and \stepref{s6Psums},
   $(P^TP)(c,c')$ is $\frac12(n^2-n)$ if $c=c'$ and $-\frac12n$ otherwise; these are the entries
   of $\frac12n^2I-\frac12nJ$.~\qed\end{proof}
 \end{proof}
 \step{s6Ptwo}{Part 2 holds.}
 \begin{proof}\pf\ The entry of $P\Zj$ at $\{i,j\}$ is $(\Zp_{ij},\Zj)$, which is 0 by
  Lemma~\ref{lem:6-gram}, part 4.~\qed\end{proof}
 \step{s6Pthree}{Part 3 holds.}
 \begin{proof}\pf\ We have that
  \[
   G^2=P(P^TP)P^T=\tfrac12n^2PP^T-\tfrac12n(P\Zj)(P\Zj)^T=\tfrac12n^2G,
  \]
  by associativity, by \stepref{s6Pone} with $J=\Zj\Zj^T$, and by \stepref{s6Ptwo}.~\qed
 \end{proof}
 \qedstep
 \begin{proof}\pf\ Steps \stepref{s6Pone}, \stepref{s6Ptwo} and \stepref{s6Pthree}.~\qed\end{proof}
\end{proof}

\begin{remark}
 \lean{Matrix.IsHadamard.charpoly_transpose_hadamardPairs_mul}
 \lean{Matrix.IsHadamard.charpoly_hadamardPairs_mul_transpose}
 \lean{Matrix.IsHadamard.spectrumMultiset_transpose_hadamardPairs_mul}
 \lean{Matrix.IsHadamard.spectrumMultiset_hadamardPairs_mul_transpose}
 \lean{Matrix.IsHadamard.eigenvalueMultiplicity_transpose_hadamardPairs_mul_zero}
 \lean{Matrix.IsHadamard.eigenvalueMultiplicity_transpose_hadamardPairs_mul_half_sq}
 \lean{Matrix.IsHadamard.eigenvalueMultiplicity_hadamardPairs_mul_transpose_zero}
 \lean{Matrix.IsHadamard.eigenvalueMultiplicity_hadamardPairs_mul_transpose_half_sq}
 \lean{Matrix.IsHadamard.finrank_eigenspace_transpose_hadamardPairs_mul_zero}
 \lean{Matrix.IsHadamard.finrank_eigenspace_transpose_hadamardPairs_mul_half_sq}
 \lean{Matrix.IsHadamard.finrank_eigenspace_hadamardPairs_mul_transpose_zero}
 \lean{Matrix.IsHadamard.finrank_eigenspace_hadamardPairs_mul_transpose_half_sq}
The paper deduces from part 1 that $P^TP$ has the eigenvalue 0, which is simple, and the
eigenvalue $\frac12n^2$, of multiplicity $n-1$; and that $PP^T$ has the eigenvalues 0 and
$\frac12n^2$, with multiplicities $\frac12(n-1)(n-2)$ and $n-1$. The second also follows from
part 3, since $\frac2{n^2}G$ is then an idempotent symmetric matrix whose trace, by
Lemma~\ref{lem:6-gram}, part 1, is $\frac2{n^2}\cdot n\cdot\frac12n(n-1)=n-1$. These are side
claims, and we prove them in \path{LineGraphs/SideClaims/HadamardSpectra.lean} for every $n\geq1$,
in the three forms of the spectral toolkit: the characteristic polynomials, the multisets of
eigenvalues, and the multiplicities, which are the dimensions of the eigenspaces. For instance,
\path{Matrix.IsHadamard.spectrumMultiset_hadamardPairs_mul_transpose} gives the spectrum of
$PP^T$. We will not use the spectrum in the proof of Theorem 6.2.
\end{remark}

\begin{namedresult}{Lemma}{(the lines through a fixed row)}\label{lem:6-fixedrow}
 \lean{Matrix.IsHadamard.hadamardPairs_orthogonal_of_mem}
 \lean{Matrix.IsHadamard.eq_sum_inner_hadamardPairs}
 \uses{def:6-pairs,lem:6-gram,in6:mathlib,facts:inner}
 \assume{\begin{pfenum}
  \item $H$ is an $n\times n$ Hadamard matrix, with $n\geq2$, and $\Zp_{ij}$ is as in
   Definition~\ref{def:6-pairs}.
  \item $i$ is a row.
 \end{pfenum}}
 \prove{\begin{pfenum}
  \item The $n-1$ vectors $\Zp_{ij}$ with $j\neq i$ are nonzero and pairwise orthogonal, and
   every $\Zp_{kl}$ lies in $\jp$.
  \item For every $\Zv\in\jp$,
   \[
    \Zv=\frac1n\sum_{j\neq i}(\Zv,\Zp_{ij})\,\Zp_{ij}.
   \]
  \item The vectors $\Zp_{ij}$ with $j\neq i$ span $\jp$, which has dimension $n-1$.
 \end{pfenum}}
\end{namedresult}
\begin{proof}
 \pfsketch\ The $n$ vectors $\Zj$ and $\Zp_{ij}$ ($j\neq i$) are nonzero and pairwise orthogonal
 in $\re^n$, and thus a basis. We will expand in this basis.
 \step{s6frone}{Part 1 holds.}
 \begin{proof}\pf\ Lemma~\ref{lem:6-gram}: part 1 gives $(\Zp_{ij},\Zp_{ij})=n\neq0$, part 2
  orthogonality, and part 4 that $\Zp_{kl}\in\jp$.~\qed\end{proof}
 \step{s6frbasis}{The $n$ vectors $\Zj$ and $\Zp_{ij}$ ($j\neq i$) are linearly independent and span
  $\re^n$.}
 \begin{proof}
  \step{s6frind}{They are linearly independent.}
  \begin{proof}\pf\ They are nonzero and pairwise orthogonal by \stepref{s6frone} and
   $(\Zj,\Zj)=n$; apply \path{linearIndependent_of_ne_zero_of_inner_eq_zero}.~\qed\end{proof}
  \qedstep
  \begin{proof}\pf\ Independence is \stepref{s6frind}. There are $n=\dim\re^n$ of them, so
   \path{LinearIndependent.span_eq_top_of_card_eq_finrank} applied to \stepref{s6frind} gives
   the span.~\qed\end{proof}
 \end{proof}
 \step{s6frtwo}{\assume{$\Zv\in\jp$.}\prove{Part 2 holds for $\Zv$.}}
 \begin{proof}
  \pflet{$\Zw=\Zv-\frac1n\sum_{j\neq i}(\Zv,\Zp_{ij})\Zp_{ij}$.}
  \step{s6frwp}{$(\Zw,\Zp_{ik})=0$ for each $k\neq i$.}
  \begin{proof}\pf\ By bilinearity and \stepref{s6frone}, only the term $j=k$ of the sum
   survives, and it gives $(\Zw,\Zp_{ik})=(\Zv,\Zp_{ik})-\frac1n(\Zv,\Zp_{ik})\,n$.~\qed
  \end{proof}
  \step{s6frwj}{$(\Zw,\Zj)=0$.}
  \begin{proof}\pf\ $\Zv\in\jp$ by \asm{1}, and each $\Zp_{ij}\in\jp$ by \stepref{s6frone}.~\qed
  \end{proof}
  \qedstep
  \begin{proof}\pf\ By \stepref{s6frwp}, \stepref{s6frwj} and \stepref{s6frbasis}, $\Zw$ is
   orthogonal to a spanning set, so $(\Zw,\Zw)=0$ and $\Zw=0$ by \path{inner_self_eq_zero}.~\qed
  \end{proof}
 \end{proof}
 \step{s6frthree}{Part 3 holds.}
 \begin{proof}\pf\ Part 2 writes each $\Zv\in\jp$ in the span, and the vectors lie in $\jp$ by
  part 1. By \stepref{s6frbasis} they are $n-1$ linearly independent vectors, and thus
  the dimension is $n-1$
  by \path{finrank_span_eq_card}.~\qed\end{proof}
 \qedstep
 \begin{proof}\pf\ Steps \stepref{s6frone}, \stepref{s6frtwo} and \stepref{s6frthree}.~\qed
 \end{proof}
\end{proof}

\subsection{Hypothesis 6.1}

\begin{cgss}{Hypothesis}{6.1}\label{hyp:6.1}
 \lean{Matrix.IsHadamard.FourRowAgreement}\uses{def:6-pairs}
 \assume{$H$ is an $n\times n$ Hadamard matrix and $x\geq0$ is an integer.}
 We say that $H$ satisfies Hypothesis 6.1 with parameter $x$ if, for all distinct rows
 $i,j,k,l$,
 \[
  8\agr(i,j,k,l)\in\{n,\ n-x,\ n+x\}.
 \]
 The paper allows $x\in\ints$; replacing $x$ by $-x$ does not change the condition, so we take
 $x\geq0$, as in the Lean statement \texttt{H.FourRowAgreement x} with
 \texttt{x : \ensuremath{\mathbb N}}.
\end{cgss}

\begin{namedresult}{Definition}{(the matrix $C_H$)}\label{def:6-CH}
 \lean{Matrix.hadamardSignedAdj}\uses{def:6-pairs}
 \assume{$H$ is an $n\times n$ Hadamard matrix, $x>0$ is an integer, and $P$ is as in
 Definition~\ref{def:6-pairs}.}
 \pflet{$C_H$ be the $\cP\times\cP$ matrix with $C_H=\frac1x(PP^T-nI)$.}
 Thus $PP^T=nI+xC_H$, which is the paper's display with $C_H$ for $A$. The diagonal of $C_H$ is
 0 by Lemma~\ref{lem:6-gram}, part 1. The Lean declaration takes $x$ as an argument,
 \texttt{Matrix.hadamardSignedAdj H x}, the settled name of \path{cgss/notation.md}, which
 now takes $x$.
\end{namedresult}

\begin{namedresult}{Lemma}{(the entries of $C_H$)}\label{lem:6-CHentries}
 \lean{Matrix.IsHadamard.FourRowAgreement.hadamardSignedAdj_mem}
 \uses{def:6-CH,hyp:6.1,lem:6-gram}
 \assume{\begin{pfenum}
  \item $H$ is an $n\times n$ Hadamard matrix and $x>0$ is an integer.
  \item $H$ satisfies Hypothesis 6.1 with parameter $x$.
 \end{pfenum}}
 \prove{Every entry of $C_H$ lies in $\{0,1,-1\}$. For $\{i,j\}\neq\{k,l\}$ we have
 $C_H(\{i,j\},\{k,l\})=0$ when the two pairs meet, and $C_H(\{i,j\},\{k,l\})=0$ exactly when
 $8\agr(i,j,k,l)=n$ when they are disjoint.}
\end{namedresult}
\begin{proof}
 \step{s6Cdiag}{The diagonal entries are 0.}
 \begin{proof}\pf\ Definition~\ref{def:6-CH}.~\qed\end{proof}
 \step{s6Cmeet}{\case{The pairs are distinct and meet.}}
 \begin{proof}\pf\ They share exactly one row, and $(\Zp_{ij},\Zp_{kl})=0$ by
  Lemma~\ref{lem:6-gram}, part 2.~\qed\end{proof}
 \step{s6Cdisj}{\case{The pairs are disjoint.}}
 \begin{proof}\pf\ $xC_H(\{i,j\},\{k,l\})=(\Zp_{ij},\Zp_{kl})=8\agr(i,j,k,l)-n$ by
  Lemma~\ref{lem:6-gram}, part 3, and this lies in $\{0,-x,x\}$ by \asm[2]{0}, with 0 exactly
  when $8\agr(i,j,k,l)=n$; divide by $x>0$.~\qed\end{proof}
 \qedstep
 \begin{proof}\pf\ Two entries are on the diagonal or are indexed by two distinct pairs, which
  meet or are disjoint.~\qed\end{proof}
\end{proof}

\begin{namedresult}{Lemma}{(the quadratic equation)}\label{lem:6-quadratic}
 \lean{Matrix.IsHadamard.hadamardSignedAdj_quadratic}
 \uses{def:6-CH,lem:6-P}
 \assume{$H$ is an $n\times n$ Hadamard matrix, $x>0$ is an integer, and $d=n/x\in\rats$.}
 \prove{\[
  (xC_H+nI)\bigl(xC_H-\tfrac12n(n-2)I\bigr)=0,\qquad
  (C_H+dI)\bigl(C_H-\tfrac12d(n-2)I\bigr)=0.
 \]}
\end{namedresult}
\begin{proof}
 \pfsketch\ This is Lemma~\ref{lem:6-P}, part 3, rewritten. It does not need Hypothesis 6.1,
 which only makes $C_H$ a $(0,\pm1)$ matrix.
 \pflet{$G=PP^T$.}
 \step{s6qG}{$G(G-\frac12n^2I)=0$.}
 \begin{proof}\pf\ Lemma~\ref{lem:6-P}, part 3.~\qed\end{proof}
 \step{s6qsub}{$G=xC_H+nI$ and $G-\frac12n^2I=xC_H-\frac12n(n-2)I$.}
 \begin{proof}\pf\ Definition~\ref{def:6-CH}, and $n-\frac12n^2=-\frac12n(n-2)$.~\qed\end{proof}
 \qedstep
 \begin{proof}\pf\ Substituting \stepref{s6qsub} in \stepref{s6qG} gives the first equation;
  dividing it by $x^2>0$ gives the second.~\qed\end{proof}
\end{proof}

\begin{namedresult}{Lemma}{(the fourth rows)}\label{lem:6-fourth}
 \lean{Matrix.IsHadamard.FourRowAgreement.card_fourthRows}
 \lean{Matrix.IsHadamard.FourRowAgreement.dvd_card}
 \uses{hyp:6.1,lem:6-gram,lem:6-fixedrow,in6:mathlib}
 \assume{\begin{pfenum}
  \item $H$ is an $n\times n$ Hadamard matrix, with $n\geq3$, and $x>0$ is an integer.
  \item $H$ satisfies Hypothesis 6.1 with parameter $x$.
  \item $i,k,l$ are distinct rows, and $m$ is the number of rows $j\notin\{i,k,l\}$ with
   $8\agr(i,j,k,l)\neq n$.
 \end{pfenum}}
 \prove{$mx^2=n^2$. Thus $x\mid n$, and $d=n/x$ is an integer with $d^2=m$.}
\end{namedresult}
\begin{proof}
 \pfsketch\ \repair\ The paper derives $x\mid n$ from $xAJ=-nJ$, which is false (errata,
 Error 9). We will expand $\Zp_{kl}$ in the orthogonal basis of $\jp$ attached to the row $i$,
 and compare squared norms. This is the paper's remark that $d^2$ is the number of fourth rows
 agreeing with three given rows in $\frac18(n\pm x)$ places.
 \step{s6fparseval}{$n^2=\sum_{j\neq i}(\Zp_{kl},\Zp_{ij})^2$.}
 \begin{proof}\pf\ Take the inner product of the expansion of Lemma~\ref{lem:6-fixedrow}, part 2,
  for $\Zv=\Zp_{kl}\in\jp$, with $\Zp_{kl}$; the left side is $(\Zp_{kl},\Zp_{kl})=n$ by
  Lemma~\ref{lem:6-gram}, part 1. Multiply by $n$.~\qed\end{proof}
 \step{s6fkl}{The terms $j=k$ and $j=l$ vanish.}
 \begin{proof}\pf\ The pairs $\{i,k\}$ and $\{k,l\}$ share only $k$, and $\{i,l\}$ and
  $\{k,l\}$ share only $l$; apply Lemma~\ref{lem:6-gram}, part 2.~\qed\end{proof}
 \step{s6fother}{For $j\notin\{i,k,l\}$, $(\Zp_{kl},\Zp_{ij})^2$ is $x^2$ if
  $8\agr(i,j,k,l)\neq n$ and 0 otherwise.}
 \begin{proof}\pf\ By Lemma~\ref{lem:6-gram}, part 3, $(\Zp_{ij},\Zp_{kl})=8\agr(i,j,k,l)-n$,
  which is $0$ or $\pm x$ by \asm[2]{0}.~\qed\end{proof}
 \step{s6fcount}{$mx^2=n^2$.}
 \begin{proof}\pf\ Substitute \stepref{s6fkl} and \stepref{s6fother} in \stepref{s6fparseval},
  and use \asm[3]{0}.~\qed\end{proof}
 \qedstep
 \begin{proof}\pf\ By \stepref{s6fcount}, $x^2\mid n^2$, so $x\mid n$ by
  \path{Nat.pow_dvd_pow_iff}; then $d^2=n^2/x^2=m$.~\qed\end{proof}
\end{proof}

\begin{remark}
\repair\ (Errata, Error 9.) The paper's equation $xAJ=-nJ$ says that $P^T\Zj=0$. The entry of
$P^T\Zj$ at a column $c$ is $\frac12(\sigma(c)^2-n)$, where $\sigma(c)$ is the sum of column $c$
of $H$, by the computation of Lemma~\ref{lem:6-P}, part 1, with $u_i=H(i,c)$. It vanishes only
if every column sum is $\pm\sqrt n$, which fails unless $n$ is a square. None of our results uses
it. A second route to $x\mid n$ for $n>2$: the eigenvalue $-n/x$ of the integer matrix $C_H$ is
rational and is a root of its monic characteristic polynomial, so it is an integer by the
rational root theorem ([ST, Theorem 9.6]; [ST, Theorem 10.4] proves $P^TP$ and $x\mid n$ by this
route).
\end{remark}

\subsection{Theorem 6.2}

Theorem 6.2 is the case $x=\frac12n$, $d=2$: any four rows agree in $\frac n{16}$, $\frac n8$ or
$\frac{3n}{16}$ places. We will first dispose of the orders not divisible by 16, then build the
line system and study its components.

\begin{namedresult}{Lemma}{(orders not divisible by 16)}\label{lem:6-small}
 \lean{Matrix.IsHadamard.sixteen_dvd_of_agree_sixteenths}
 \uses{def:6-pairs,lem:6-gram,in6:mathlib}
 \assume{\begin{pfenum}
  \item $H$ is an $n\times n$ Hadamard matrix, with $n>2$.
  \item For all distinct rows $i,j,k,l$, $16\agr(i,j,k,l)\in\{n,2n,3n\}$.
 \end{pfenum}}
 \prove{$16\mid n$, and $n\geq16$.}
\end{namedresult}
\begin{proof}
 \pfsketch\ \repair\ (Errata, Gap 10; \path{cgss/prerequisites.md}, HD1.) If $16\nmid n$, every
 four rows agree in exactly $\frac n8$ places, and then the $\frac12n(n-1)$ vectors $\Zp_{ij}$
 are pairwise orthogonal in $\re^n$, which is too many.
 \step{s6sfour}{$4\mid n$, so $n\geq4$.}
 \begin{proof}\pf\ \path{Matrix.IsHadamard.four_dvd_card} with \asm[1]{0}.~\qed\end{proof}
 \step{s6scontra}{\assume{$16\nmid n$.}\prove{FALSE.}}
 \begin{proof}
  \step{s6stwo}{For all distinct rows $i,j,k,l$, $16\agr(i,j,k,l)=2n$.}
  \begin{proof}
   \step{s6sone}{$16\agr(i,j,k,l)\neq n$.}
   \begin{proof}\pf\ Otherwise $16\mid n$, against \asm{1}.~\qed\end{proof}
   \step{s6sthree}{$16\agr(i,j,k,l)\neq3n$.}
   \begin{proof}\pf\ Otherwise $16\mid3n$, so $16\mid n$ by
    \path{Nat.Coprime.dvd_of_dvd_mul_left}, since 3 and 16 are coprime, against
    \asm{1}.~\qed\end{proof}
   \qedstep
   \begin{proof}\pf\ \asm[2]{0}, \stepref{s6sone} and \stepref{s6sthree}.~\qed\end{proof}
  \end{proof}
  \step{s6sorth}{The vectors $\Zp_\pi$, $\pi\in\cP$, are nonzero and pairwise orthogonal.}
  \begin{proof}\pf\ Nonzero by Lemma~\ref{lem:6-gram}, part 1. Two distinct pairs meet in one row
   or are disjoint; Lemma~\ref{lem:6-gram}, parts 2 and 3, gives inner product 0, or
   $8\agr-n=0$ by \stepref{s6stwo}.~\qed\end{proof}
  \step{s6sbound}{$\frac12n(n-1)\leq n$.}
  \begin{proof}\pf\ \path{linearIndependent_of_ne_zero_of_inner_eq_zero} makes the family of
   \stepref{s6sorth}, indexed by $\cP$, linearly independent, and
   \path{LinearIndependent.fintype_card_le_finrank} with \path{finrank_euclideanSpace_fin}
   bounds $|\cP|$ by $n$.~\qed\end{proof}
  \qedstep
  \begin{proof}\pf\ \stepref{s6sbound} gives $n\leq3$, against \stepref{s6sfour}.~\qed\end{proof}
 \end{proof}
 \qedstep
 \begin{proof}\pf\ By \stepref{s6scontra}, $16\mid n$; and $n>0$, so $n\geq16$.~\qed\end{proof}
\end{proof}

\begin{namedresult}{Lemma}{(the line system of $H$)}\label{lem:6-lines}
 \lean{Matrix.IsHadamard.lineSystem}\lean{Matrix.IsHadamard.lineSystem_signedAdj}
 \lean{Matrix.IsHadamard.card_adj_lineSystem}
 \uses{def:6-pairs,lem:6-gram,hyp:6.1,def:6-CH,lem:6-CHentries,lem:6-quadratic,lem:6-fixedrow,
 def:line-system,def:sec3-signing,facts:inner}
 \assume{\begin{pfenum}
  \item $H$ is an $n\times n$ Hadamard matrix, with $16\mid n$ and $n\geq16$.
  \item For all distinct rows $i,j,k,l$, $16\agr(i,j,k,l)\in\{n,2n,3n\}$.
 \end{pfenum}}
 \pflet{$\Zr_\pi=\sqrt{2/n}\,\Zp_\pi$ for $\pi\in\cP$, and $\Phi_H=\{\pm\Zr_\pi:\pi\in\cP\}$.}
 \prove{\begin{pfenum}
  \item $H$ satisfies Hypothesis 6.1 with parameter $x=\frac12n$, and then $d=2$.
  \item $(\Zr_\pi,\Zr_\pi)=2$, and for $\pi\neq\pi'$,
   $(\Zr_\pi,\Zr_{\pi'})=C_H(\pi,\pi')\in\{0,1,-1\}$.
  \item $\Phi_H$ is the vector set of a line system $S_H$ in $\re^n$ with $\frac12n(n-1)$ lines,
   and $\pi\mapsto\Zr_\pi$ is a signing of $S_H$ whose matrix $C$ of
   Definition~\ref{def:sec3-signing} is $C_H$.
  \item $(C_H+2I)(C_H-(n-2)I)=0$.
  \item For each $\pi\in\cP$ there are exactly $2n-4$ pairs $\pi'\neq\pi$ with
   $(\Zr_\pi,\Zr_{\pi'})\neq0$.
  \item Every entry of every $\Zr_\pi$ is $\pm\sqrt{2/n}$; every vector of $\Phi_H$ lies in
   $\jp$; and for each row $i$, the vectors $\Zr_{ij}$, $j\neq i$, are pairwise orthogonal.
 \end{pfenum}}
\end{namedresult}
\begin{proof}
 \pfsketch\ With $x=\frac12n$ we rescale the Hadamard products to squared norm 2, and read off
 each part from the lemmas above.
 \step{s6lone}{Part 1 holds.}
 \begin{proof}\pf\ $x=\frac12n$ is an integer, and $\{n,2n,3n\}/16=\{n-x,n,n+x\}/8$; so
  \asm[2]{0} is Hypothesis 6.1. Then $d=n/x=2$.~\qed\end{proof}
 \step{s6ltwo}{Part 2 holds.}
 \begin{proof}\pf\ $(\Zr_\pi,\Zr_{\pi'})=\frac2n(\Zp_\pi,\Zp_{\pi'})$. For $\pi=\pi'$ this is
  $\frac2n\cdot n$ by Lemma~\ref{lem:6-gram}, part 1. For $\pi\neq\pi'$ it is
  $\frac2n\cdot xC_H(\pi,\pi')=C_H(\pi,\pi')$ by Definition~\ref{def:6-CH}, and
  Lemma~\ref{lem:6-CHentries} with \stepref{s6lone} gives the values.~\qed\end{proof}
 \step{s6lthree}{Part 3 holds.}
 \begin{proof}
  \step{s6lls}{$\Phi_H$ is a line system.}
  \begin{proof}\pf\ It is finite and $\Phi_H=-\Phi_H$ by definition. By \stepref{s6ltwo} and
   bilinearity, the vectors have squared norm 2 and pairwise inner products in
   $\{0,\pm1,\pm2\}$; this is the definition in Section~\ref{sec:2}.~\qed\end{proof}
  \step{s6ldistinct}{For $\pi\neq\pi'$, $\Zr_{\pi'}\neq\pm\Zr_\pi$.}
  \begin{proof}\pf\ By \stepref{s6ltwo}, $(\Zr_\pi,\Zr_{\pi'})\neq\pm2$; apply
   \eqref{eq:endpoints}.~\qed\end{proof}
  \step{s6lsigning}{$\pi\mapsto\Zr_\pi$ is a signing of $S_H$, and $S_H$ has $|\cP|$ lines.}
  \begin{proof}\pf\ By \stepref{s6ldistinct}, $(\pi,\epsilon)\mapsto\epsilon\Zr_\pi$ is a
   bijection $\cP\times\{1,-1\}\to\Phi_H$.~\qed\end{proof}
  \qedstep
  \begin{proof}\pf\ By \stepref{s6lsigning}, the matrix $C$ of the signing has off-diagonal
   entries $(\Zr_\pi,\Zr_{\pi'})=C_H(\pi,\pi')$ by \stepref{s6ltwo}, and zero diagonal, as $C_H$
   has.~\qed\end{proof}
 \end{proof}
 \step{s6lfour}{Part 4 holds.}
 \begin{proof}\pf\ Lemma~\ref{lem:6-quadratic} with $d=2$ by \stepref{s6lone}:
  $\frac12d(n-2)=n-2$.~\qed\end{proof}
 \step{s6lfive}{Part 5 holds.}
 \begin{proof}
  \step{s6lsq}{$C_H^2=(n-4)C_H+2(n-2)I$.}
  \begin{proof}\pf\ Expand \stepref{s6lfour}.~\qed\end{proof}
  \step{s6ldiag}{$C_H^2(\pi,\pi)=\sum_{\pi'}C_H(\pi,\pi')^2$ is the number of $\pi'$ with
   $C_H(\pi,\pi')\neq0$.}
  \begin{proof}\pf\ $C_H$ is symmetric with entries in $\{0,\pm1\}$ by
   Lemma~\ref{lem:6-CHentries}.~\qed\end{proof}
  \qedstep
  \begin{proof}\pf\ The diagonal of the right side of \stepref{s6lsq} is $2(n-2)$, since
   $C_H(\pi,\pi)=0$; use \stepref{s6ldiag}, \stepref{s6ltwo}, and that such $\pi'$ differ from
   $\pi$.~\qed\end{proof}
 \end{proof}
 \step{s6lsix}{Part 6 holds.}
 \begin{proof}\pf\ Scale Lemma~\ref{lem:6-gram}, part 1, and Lemma~\ref{lem:6-fixedrow},
  part 1, by $\pm\sqrt{2/n}$.~\qed\end{proof}
 \qedstep
 \begin{proof}\pf\ Steps \stepref{s6lone} to \stepref{s6lsix}.~\qed\end{proof}
\end{proof}

We keep the hypotheses of Lemma~\ref{lem:6-lines} and its notation for the rest of the section.
Let $\Lambda_H$ be the graph on $\cP$ in which $\pi\neq\pi'$ are adjacent when
$(\Zr_\pi,\Zr_{\pi'})\neq0$. A \emph{component} is the vertex set $\cK\sbs\cP$ of a connected
component of $\Lambda_H$ (\texttt{SimpleGraph.ConnectedComponent}); it is nonempty. For a
component $\cK$ we write $\Phi(\cK)=\{\pm\Zr_\pi:\pi\in\cK\}$, $U_\cK$ for its span,
$m_\cK=\dim U_\cK$ for its \emph{rank}, and $R_\cK=\scl(\Phi(\cK))$ (LN3).

\begin{namedresult}{Lemma}{(the components)}\label{lem:6-components}
 \lean{Matrix.IsHadamard.card_component_row}\lean{Matrix.IsHadamard.two_mul_card_component}
 \lean{Matrix.IsHadamard.component_rootType}
 \uses{lem:6-lines,lem:6-fixedrow,in6:LN3,in6:sec3,thm:3.5,lem:3.4,lem:sec3-spans,in6:mathlib,
 facts:inner}
 \assume{\begin{pfenum}
  \item $H$ is an $n\times n$ Hadamard matrix, with $16\mid n$ and $n\geq16$.
  \item For all distinct rows $i,j,k,l$, $16\agr(i,j,k,l)\in\{n,2n,3n\}$.
  \item $\cK$ is a component of $\Lambda_H$, with rank $m=m_\cK$.
 \end{pfenum}}
 \prove{\begin{pfenum}
  \item For every row $i$, the number of rows $j$ with $\{i,j\}\in\cK$ is $m$.
  \item $2|\cK|=nm$.
  \item $1\leq m\leq n-1$.
  \item $R_\cK$ is a nonempty indecomposable star-closed line system in $\re^n$, with
   $\Phi(\cK)\sbs\Phi(R_\cK)$, $|\Phi(\cK)|=2|\cK|$, and span $U_\cK$.
  \item There is a type $t$ of Theorem~\ref{thm:3.5}, namely $\typ A_\ell$ ($\ell\geq1$),
   $\typ D_\ell$ ($\ell\geq4$), $\typ E_6$, $\typ E_7$ or $\typ E_8$, with an isomorphism
   $R_\cK\cong t$ induced by a surjective linear isometry of spans; $t$ is unique, by
   Theorem~\ref{thm:3.5}, and we call it the \emph{type} of $\cK$. The rank of $t$ is $m$,
   $R_\cK$ and $t$ have the same number $L_t$ of lines, and $\frac12nm\leq L_t$.
 \end{pfenum}}
\end{namedresult}
\begin{proof}
 \pfsketch\ \repair\ (Errata, Gap 17.) For a fixed row $i$, the vectors $\Zr_{ij}$ with
 $\{i,j\}\in\cK$ are an orthogonal basis of $U_\cK$, by the expansion of
 Lemma~\ref{lem:6-fixedrow}; so each row lies in exactly $m$ pairs of $\cK$. The rest is double
 counting and Theorem~\ref{thm:3.5}.
 \step{s6cone}{Part 1 holds.}
 \begin{proof}
  \pflet{$i$ be a row, $B=\{\{i,j\}\in\cK\}$, and $U_B$ the span of $\{\Zr_\beta:\beta\in B\}$.}
  \step{s6cexp}{\assume{$\pi\in\cK$.}\prove{$\Zr_\pi\in U_B$.}}
  \begin{proof}
   \step{s6cterm}{If $(\Zr_\pi,\Zr_{ij})\neq0$ for a row $j\neq i$, then $\{i,j\}\in B$.}
   \begin{proof}\pf\ Either $\{i,j\}=\pi$, or $\{i,j\}\neq\pi$ is adjacent to $\pi$ in
    $\Lambda_H$; in both cases $\{i,j\}$ lies in the component $\cK$ of $\pi$.~\qed\end{proof}
   \step{s6csum}{$\Zr_\pi=\frac12\sum_{j\neq i}(\Zr_\pi,\Zr_{ij})\Zr_{ij}$.}
   \begin{proof}\pf\ Scale the expansion of Lemma~\ref{lem:6-fixedrow}, part 2, for
    $\Zv=\Zp_\pi\in\jp$ by $\sqrt{2/n}$.~\qed\end{proof}
   \qedstep
   \begin{proof}\pf\ By \stepref{s6cterm} each term of \stepref{s6csum} with a nonzero
    coefficient lies in $U_B$.~\qed\end{proof}
  \end{proof}
  \step{s6cspan}{$U_B=U_\cK$.}
  \begin{proof}\pf\ $U_B\sbs U_\cK$ since $B\sbs\cK$; $U_\cK\sbs U_B$ by \stepref{s6cexp},
   since $U_B$ is closed under negation.~\qed\end{proof}
  \step{s6cind}{The vectors $\Zr_\beta$, $\beta\in B$, are linearly independent.}
  \begin{proof}\pf\ They are nonzero and pairwise orthogonal by Lemma~\ref{lem:6-lines}, part 6;
   apply \path{linearIndependent_of_ne_zero_of_inner_eq_zero}.~\qed\end{proof}
  \qedstep
  \begin{proof}\pf\ $|B|=\dim U_B=m$ by \stepref{s6cind}, \path{finrank_span_eq_card} and
   \stepref{s6cspan}.~\qed\end{proof}
 \end{proof}
 \step{s6ctwo}{Part 2 holds.}
 \begin{proof}\pf\ Count the pairs (row $i$, pair $\pi\in\cK$ with $i\in\pi$) in two ways: each
  $\pi$ has two rows, and each row lies in $m$ pairs of $\cK$ by \stepref{s6cone}.~\qed
 \end{proof}
 \step{s6cthree}{Part 3 holds.}
 \begin{proof}\pf\ $\cK$ is nonempty, so $2|\cK|=nm>0$ by \stepref{s6ctwo}; and a row $i$ lies in
  at most $n-1$ pairs, so $m\leq n-1$ by \stepref{s6cone}.~\qed\end{proof}
 \step{s6cfour}{Part 4 holds.}
 \begin{proof}
  \step{s6cfsigma}{$\Phi(\cK)$ satisfies the hypotheses of LN3, and its graph of nonzero inner
   products is connected.}
  \begin{proof}\pf\ Lemma~\ref{lem:6-lines}, parts 2 and 3, give the norms and inner products
   and $\Phi(\cK)=-\Phi(\cK)$. The vectors $\Zr_\pi$ and $-\Zr_\pi$ have inner product $-2$, and
   $\pm\Zr_\pi$, $\pm\Zr_{\pi'}$ are joined when $\pi,\pi'$ are adjacent in $\Lambda_H$; $\cK$ is
   connected in $\Lambda_H$.~\qed\end{proof}
  \step{s6cfcard}{$|\Phi(\cK)|=2|\cK|$.}
  \begin{proof}\pf\ The vectors $\epsilon\Zr_\pi$ are distinct by Lemma~\ref{lem:6-lines},
   part 3.~\qed\end{proof}
  \qedstep
  \begin{proof}\pf\ Apply LN3 to \stepref{s6cfsigma}; its span is that of $\Phi(\cK)$, which is
   $U_\cK$; use \stepref{s6cfcard}.~\qed\end{proof}
 \end{proof}
 \step{s6cfive}{Part 5 holds.}
 \begin{proof}
  \step{s6ctype}{There is exactly one type $t$ with $R_\cK\cong t$, and the isomorphism is
   induced by a surjective linear isometry from $U_\cK$ onto the span of the model of $t$.}
  \begin{proof}\pf\ Theorem~\ref{thm:3.5}, applied to \stepref{s6cfour} in $V=\re^n$.~\qed
  \end{proof}
  \step{s6crank}{The rank of $t$ is $m$.}
  \begin{proof}\pf\ A surjective linear isometry preserves dimension; the rank of the model is
   the dimension of its span (Lemma~\ref{lem:sec3-spans}), and $\dim U_\cK=m$.~\qed\end{proof}
  \step{s6clines}{$R_\cK$ has $L_t$ lines, and $|\cK|\leq L_t$.}
  \begin{proof}\pf\ An isomorphism is a bijection of vector sets, so the numbers of lines agree
   (Lemma~\ref{lem:3.4} gives $L_t$); and $2|\cK|=|\Phi(\cK)|\leq|\Phi(R_\cK)|=2L_t$ by
   \stepref{s6cfour}.~\qed\end{proof}
  \qedstep
  \begin{proof}\pf\ Steps \stepref{s6ctype}, \stepref{s6crank} and \stepref{s6clines}, with
   $|\cK|=\frac12nm$ by \stepref{s6ctwo}.~\qed\end{proof}
 \end{proof}
 \qedstep
 \begin{proof}\pf\ Steps \stepref{s6cone} to \stepref{s6cfive}.~\qed\end{proof}
\end{proof}

\begin{namedresult}{Lemma}{(the possible decompositions)}\label{lem:6-decomposition}
 \lean{Matrix.IsHadamard.component_cases}
 \uses{lem:6-components,lem:3.4,thm:3.5,lem:sec3-spans}
 \assume{\begin{pfenum}
  \item $H$ is an $n\times n$ Hadamard matrix, with $16\mid n$ and $n\geq16$.
  \item For all distinct rows $i,j,k,l$, $16\agr(i,j,k,l)\in\{n,2n,3n\}$.
 \end{pfenum}}
 \prove{One of the following holds.
 \begin{pfenum}
  \item $\cP$ is a single component, and its type in Lemma~\ref{lem:6-components} is
   $\typ A_{n-1}$ or $\typ D_{n-1}$.
  \item $n=16$, and some component has type $\typ E_7$.
 \end{pfenum}}
\end{namedresult}
\begin{proof}
 \pfsketch\ This is the analysis in Gap 17 of the errata. Each type forces a bound on the rank
 through $\frac12nm\leq L_t$, and the ranks of the components add up to $n-1$.
 \step{s6dsum}{The ranks of the components add up to $n-1$.}
 \begin{proof}\pf\ The $n-1$ pairs $\{1,j\}$ lie in the components, each in exactly one, and a
  component $\cK$ contains $m_\cK$ of them by Lemma~\ref{lem:6-components}, part 1.~\qed
 \end{proof}
 \step{s6dtypes}{\assume{$\cK$ is a component, with rank $m$ and type $t$.}
  \prove{\begin{pfenum}
   \item If $t=\typ A_m$, then $m=n-1$.
   \item If $t=\typ D_m$, then $m\geq\frac12n+1$.
   \item $t\neq\typ E_6$.
   \item If $t$ is $\typ E_7$ or $\typ E_8$, then $n=16$.
  \end{pfenum}}}
 \begin{proof}
  \step{s6dA}{Part 1 holds.}
  \begin{proof}\pf\ $\frac12nm\leq\frac12m(m+1)$ by Lemma~\ref{lem:6-components}, part 5, and
   Lemma~\ref{lem:3.4}; divide by $\frac12m>0$ to get $n-1\leq m$, and $m\leq n-1$ by part 3 of
   that lemma.~\qed\end{proof}
  \step{s6dD}{Part 2 holds.}
  \begin{proof}\pf\ $\frac12nm\leq m(m-1)$ in the same way; divide by $m>0$.~\qed\end{proof}
  \step{s6dEsix}{Part 3 holds.}
  \begin{proof}\pf\ Otherwise $m=6$ and $3n\leq36$, so $n\leq12$, against \asm[1]{0}.~\qed
  \end{proof}
  \step{s6dEseveneight}{Part 4 holds.}
  \begin{proof}\pf\ For $\typ E_7$, $m=7$ and $\frac72n\leq63$, so $n\leq18$; for $\typ E_8$,
   $m=8$ and $4n\leq120$, so $n\leq30$. The only multiple of 16 in $[16,30]$ is 16.~\qed
  \end{proof}
  \qedstep
  \begin{proof}\pf\ Steps \stepref{s6dA} to \stepref{s6dEseveneight}.~\qed\end{proof}
 \end{proof}
 \step{s6dcaseA}{\case{Some component $\cK$ has type $\typ A_m$.}}
 \begin{proof}\pf\ Then $m=n-1$ by \stepref{s6dtypes}. By \stepref{s6dsum} every other component
  has rank 0, which Lemma~\ref{lem:6-components}, part 3, forbids; so $\cK=\cP$, and
  alternative 1 holds.~\qed\end{proof}
 \step{s6dcaseD}{\case{Some component $\cK$ has type $\typ D_m$.}}
 \begin{proof}
  \step{s6dalone}{$\cK$ is the only component.}
  \begin{proof}
   \step{s6dother}{\assume{$\cK'\neq\cK$ is a component, of rank $m'$ and type $t'$.}
    \prove{FALSE.}}
   \begin{proof}
    \step{s6dsmall}{$m'\leq\frac12n-2$.}
    \begin{proof}\pf\ $m+m'\leq n-1$ by \stepref{s6dsum}, and $m\geq\frac12n+1$ by
     \stepref{s6dtypes}.~\qed\end{proof}
    \step{s6dnotAD}{$t'$ is neither $\typ A_{m'}$ nor $\typ D_{m'}$.}
    \begin{proof}\pf\ By \stepref{s6dtypes}, these would give $m'=n-1$ or
     $m'\geq\frac12n+1$, against \stepref{s6dsmall}.~\qed\end{proof}
    \step{s6dnotE}{$t'$ is not exceptional.}
    \begin{proof}\pf\ By \stepref{s6dtypes}, $t'\neq\typ E_6$, and for $\typ E_7$ or $\typ E_8$
     we get $n=16$ and $m'\geq7>6=\frac12n-2$, against \stepref{s6dsmall}.~\qed\end{proof}
    \qedstep
    \begin{proof}\pf\ Theorem~\ref{thm:3.5} has no other types.~\qed\end{proof}
   \end{proof}
   \qedstep
   \begin{proof}\pf\ Step \stepref{s6dother}.~\qed\end{proof}
  \end{proof}
  \qedstep
  \begin{proof}\pf\ By \stepref{s6dalone} and \stepref{s6dsum}, $\cK=\cP$ and $m=n-1$;
   alternative 1 holds.~\qed\end{proof}
 \end{proof}
 \step{s6dcaseE}{\case{No component has type $\typ A_m$ or $\typ D_m$.}}
 \begin{proof}
  \step{s6dEonly}{$n=16$, and every component has rank 7 or 8.}
  \begin{proof}\pf\ $\cP$ is nonempty, so there is a component. Each component $\cK$ has a type
   $t$ by Lemma~\ref{lem:6-components}, part 5, of rank $m_\cK$, and the model $\typ A_\ell$
   or $\typ D_\ell$ has rank $\ell$ (Lemma~\ref{lem:sec3-spans}); so $t$ is neither
   $\typ A_{m_\cK}$ nor $\typ D_{m_\cK}$ by the case hypothesis, and not $\typ E_6$ by
   \stepref{s6dtypes}, part 3. In the list of Theorem~\ref{thm:3.5} this leaves $\typ E_7$ and
   $\typ E_8$, of ranks 7 and 8 (Lemma~\ref{lem:sec3-spans}); then $n=16$ by
   \stepref{s6dtypes}, part 4.~\qed\end{proof}
  \step{s6dseven}{Some component has rank 7.}
  \begin{proof}\pf\ Otherwise the ranks are all 8 and add up to $15$ by \stepref{s6dsum}, but
   $8\nmid15$.~\qed\end{proof}
  \qedstep
  \begin{proof}\pf\ The component of \stepref{s6dseven} has type $\typ E_7$, the only type of
   rank 7 left; alternative 2 holds.~\qed\end{proof}
 \end{proof}
 \qedstep
 \begin{proof}\pf\ The three cases are exhaustive.~\qed\end{proof}
\end{proof}

\begin{namedresult}{Lemma}{(orthogonal lines in $\typ A_\ell$ and $\typ D_\ell$)}
 \label{lem:6-orthogonal}
 \lean{LineSystem.An.not_card_pairwise_orthogonal}
 \lean{LineSystem.Dn.even_of_card_pairwise_orthogonal}
 \uses{def:3.1,facts:inner}
 \assume{\begin{pfenum}
  \item $\ell\geq2$ is an integer, and $M$ is the model $\typ A_\ell$ in $\re^{\ell+1}$ or
   $\typ D_\ell$ in $\re^\ell$ of Definition~\ref{def:3.1}.
  \item $F\sbs\Phi(M)$ has $\ell$ elements, and any two distinct elements of $F$ are orthogonal.
 \end{pfenum}}
 \prove{$M$ is not $\typ A_\ell$, and if $M$ is $\typ D_\ell$, then $\ell$ is even.}
\end{namedresult}
\begin{proof}
 \pfsketch\ \new\ The paper states this without proof, and no result of Section 3 contains it;
 it needs only the coordinates of Definition~\ref{def:3.1}. Each root has a support of two
 coordinates, and two orthogonal roots have disjoint supports in $\typ A_\ell$, and disjoint or
 equal supports in $\typ D_\ell$. Theorem 3.6 is not needed.
 \pflet{the support of a root be the set of its two nonzero coordinates.}
 \step{s6oA}{\assume{$M=\typ A_\ell$.}\prove{FALSE.}}
 \begin{proof}
  \step{s6oAdisj}{Distinct $\Zr,\Zs\in F$ have disjoint supports.}
  \begin{proof}\pf\ For $\Zr=\Ze_a-\Ze_b$ and $\Zs=\Ze_c-\Ze_e$,
   $(\Zr,\Zs)=[a=c]-[a=e]-[b=c]+[b=e]$. If the supports share exactly one coordinate, this is
   $\pm1$; if they are equal, it is $\pm2$; neither is 0, against \asm[2]{0}.~\qed\end{proof}
  \qedstep
  \begin{proof}\pf\ By \stepref{s6oAdisj}, the $\ell$ supports are disjoint 2-subsets of
   $\{1,\dots,\ell+1\}$, so $2\ell\leq\ell+1$ and $\ell\leq1$, against \asm[1]{0}.~\qed\end{proof}
 \end{proof}
 \step{s6oD}{\assume{$M=\typ D_\ell$.}\prove{$\ell$ is even.}}
 \begin{proof}
  \step{s6oDsupp}{Distinct $\Zr,\Zs\in F$ have disjoint or equal supports.}
  \begin{proof}\pf\ If the supports share exactly one coordinate, $(\Zr,\Zs)$ is the product of
   the two entries there, which is $\pm1$, against \asm[2]{0}.~\qed\end{proof}
  \step{s6oDtwo}{At most two elements of $F$ have a given support $\{a,b\}$.}
  \begin{proof}
   \step{s6oDperp}{If $\Zr,\Zs\in F$ are distinct with support $\{a,b\}$, then
    $\Zs=\pm(\Zr(a)\Ze_a-\Zr(b)\Ze_b)$.}
   \begin{proof}\pf\ $0=(\Zr,\Zs)=\Zr(a)\Zs(a)+\Zr(b)\Zs(b)$ with all four entries $\pm1$, so
    $\Zs(a)\Zs(b)=-\Zr(a)\Zr(b)$.~\qed\end{proof}
   \qedstep
   \begin{proof}\pf\ If $\Zr,\Zs,\Zt\in F$ were distinct with support $\{a,b\}$, then
    $\Zs=\pm\Zt$ by \stepref{s6oDperp}, so $(\Zs,\Zt)=\pm2\neq0$, against \asm[2]{0}.~\qed
   \end{proof}
  \end{proof}
  \pflet{$k$ be the number of distinct supports of the elements of $F$.}
  \step{s6oDk}{$2k\leq\ell$ and $\ell\leq2k$.}
  \begin{proof}\pf\ The $k$ supports are pairwise disjoint 2-subsets of $\{1,\dots,\ell\}$ by
   \stepref{s6oDsupp}; and $|F|=\ell$ with at most two elements per support by
   \stepref{s6oDtwo}.~\qed\end{proof}
  \qedstep
  \begin{proof}\pf\ By \stepref{s6oDk}, $\ell=2k$.~\qed\end{proof}
 \end{proof}
 \qedstep
 \begin{proof}\pf\ Steps \stepref{s6oA} and \stepref{s6oD}.~\qed\end{proof}
\end{proof}

\begin{namedresult}{Lemma}{(no component of type $\typ E_7$)}\label{lem:6-E7}
 \lean{Matrix.IsHadamard.not_component_E₇}
 \uses{lem:6-lines,lem:6-components,in6:LN3,in6:sec3,thm:3.7,lem:3.4,def:2.1,def:sec3-signing,
 facts:inner}
 \assume{\begin{pfenum}
  \item $H$ is a $16\times16$ Hadamard matrix.
  \item For all distinct rows $i,j,k,l$, $\agr(i,j,k,l)\in\{1,2,3\}$.
  \item $\cK$ is a component of $\Lambda_H$ whose type in Lemma~\ref{lem:6-components} is
   $\typ E_7$, and $R=R_\cK$.
 \end{pfenum}}
 \prove{FALSE.}
\end{namedresult}
\begin{proof}
 \pfsketch\ \new\ \repair\ (Errata, Gap 17, the repair.) Two vectors $\Zr,\Zr'$ of
 $\Phi(\cK)$ with $(\Zr,\Zr')=-1$ span a star of $R$ whose third vector $\Zr+\Zr'$ has zero
 entries, so it lies in $R$ but not in $\Phi(\cK)$; and $\Zr$ lies in
 $N_{\Zr+\Zr'}$. Thus the 3136 such ordered pairs inject into the 448 pairs $(\Zt,\Zu)$ with
 $\Zt\in\Phi(R)\diff\Phi(\cK)$ and $\Zu\in N_\Zt$. This is the errata's count $784\leq7\cdot16$,
 of pairs of lines against the 16 stars through each of the 7 lines of $R$ outside $\cK$,
 multiplied by 4; we count signed vectors so that no choice of orientation, and no definition of
 a star as a set of lines, is needed.
 \step{s6ehyp}{The hypotheses of Lemmas~\ref{lem:6-lines} and~\ref{lem:6-components} hold,
  with $n=16$.}
 \begin{proof}\pf\ $16\mid16$, and $\agr\in\{1,2,3\}$ is $16\agr\in\{n,2n,3n\}$ for $n=16$;
  use \asm[1]{0}, \asm[2]{0} and \asm[3]{0}.~\qed\end{proof}
 \pflet{$\Phi=\Phi(\cK)$ and $Q=\{(\Zr,\Zr')\in\Phi\times\Phi:(\Zr,\Zr')=-1\}$.}
 \step{s6ecounts}{$|\Phi|=112$ and $|\Phi(R)|=126$, so $|\Phi(R)\diff\Phi|=14$.}
 \begin{proof}\pf\ The rank of $\cK$ is 7, the rank of $\typ E_7$, so $|\Phi|=2|\cK|=16\cdot7=112$
  by Lemma~\ref{lem:6-components}, parts 2, 4 and 5; $R$ has 63 lines by that lemma and
  Lemma~\ref{lem:3.4}; and $\Phi\sbs\Phi(R)$.~\qed\end{proof}
 \step{s6eQ}{$|Q|=112\cdot28=3136$.}
 \begin{proof}
  \step{s6eQr}{\assume{$\Zr\in\Phi$.}\prove{Exactly 28 vectors $\Zr'\in\Phi$ satisfy
   $(\Zr,\Zr')=-1$.}}
  \begin{proof}
   \pflet{$\Zr=\epsilon\Zr_\pi$, with $\pi\in\cK$ and $\epsilon=\pm1$.}
   \step{s6eQlines}{Exactly $2n-4=28$ pairs $\pi'\neq\pi$ have $(\Zr_\pi,\Zr_{\pi'})\neq0$, and
    they lie in $\cK$.}
   \begin{proof}\pf\ Lemma~\ref{lem:6-lines}, part 5; such $\pi'$ are adjacent to $\pi$ in
    $\Lambda_H$.~\qed\end{proof}
   \step{s6eQone}{For each such $\pi'$, exactly one of $\pm\Zr_{\pi'}$ has inner product $-1$
    with $\Zr$.}
   \begin{proof}\pf\ $(\Zr,\Zr_{\pi'})=\pm1$ by Lemma~\ref{lem:6-lines}, part 2, and
    $(\Zr,-\Zr_{\pi'})=-(\Zr,\Zr_{\pi'})$.~\qed\end{proof}
   \step{s6eQnone}{No other vector of $\Phi$ has inner product $-1$ with $\Zr$.}
   \begin{proof}\pf\ $(\Zr,\pm\Zr_\pi)=\pm2$, and $(\Zr,\pm\Zr_{\pi'})=0$ for the remaining
    $\pi'$.~\qed\end{proof}
   \qedstep
   \begin{proof}\pf\ Steps \stepref{s6eQlines}, \stepref{s6eQone} and \stepref{s6eQnone}, with
    the distinctness of Lemma~\ref{lem:6-lines}, part 3.~\qed\end{proof}
  \end{proof}
  \qedstep
  \begin{proof}\pf\ Sum \stepref{s6eQr} over the 112 vectors of \stepref{s6ecounts}.~\qed
  \end{proof}
 \end{proof}
 \step{s6ethird}{\assume{$(\Zr,\Zr')\in Q$, and $\Zw=\Zr+\Zr'$.}
  \prove{$\Zw\in\Phi(R)\diff\Phi$, and $\Zr\in N_\Zw(R)$.}}
 \begin{proof}
  \step{s6einR}{$\Zw\in\Phi(R)$.}
  \begin{proof}\pf\ $\Zr,\Zr'\in\Phi\sbs\Phi(R)$ have inner product $-1$, and $R$ is star-closed
   (Lemma~\ref{lem:6-components}, part 4, and Definition~\ref{def:2.1}).~\qed\end{proof}
  \step{s6eentries}{Every entry of $\Zw$ is $0$ or $\pm2\sqrt{2/16}$, and at most four are
   nonzero.}
  \begin{proof}\pf\ Each entry of $\Zr$ and of $\Zr'$ is $\pm\sqrt{2/16}$ by
   Lemma~\ref{lem:6-lines}, part 6, so each entry of the sum is $0$ or $\pm2\sqrt{2/16}$. A
   nonzero entry contributes $4\cdot\frac2{16}=\frac12$ to $(\Zw,\Zw)$, which is $2+2-2=2$ by
   bilinearity; so there are exactly four.~\qed\end{proof}
  \step{s6enotK}{$\Zw\notin\Phi$.}
  \begin{proof}\pf\ Every vector of $\Phi$ has 16 nonzero entries by Lemma~\ref{lem:6-lines},
   part 6, and $\Zw$ has four by \stepref{s6eentries}.~\qed\end{proof}
  \step{s6eN}{$(\Zr,\Zw)=1$.}
  \begin{proof}\pf\ $(\Zr,\Zr)+(\Zr,\Zr')=2-1$.~\qed\end{proof}
  \qedstep
  \begin{proof}\pf\ Steps \stepref{s6einR}, \stepref{s6enotK} and \stepref{s6eN}, and
   $\Zr\in\Phi(R)$.~\qed\end{proof}
 \end{proof}
 \step{s6eNt}{For every $\Zt\in\Phi(R)$, $|N_\Zt(R)|=32$.}
 \begin{proof}
  \step{s6eNsign}{$R$ has a signing $\sigma:I\to\Phi(R)$, $|I|=63$, whose matrix $C$
   satisfies $(C+2I_{63})(C-16I_{63})=0$.}
  \begin{proof}\pf\ Choose one vector on each line. $R\cong\typ E_7$ through a surjective
   linear isometry of spans by Lemma~\ref{lem:6-components}, part 5, so $R$ is an isometric
   realization of $\typ E_7$; apply Theorem~\ref{thm:3.7}.~\qed\end{proof}
  \step{s6eNrow}{Each row of $C$ has exactly 32 nonzero entries.}
  \begin{proof}\pf\ From \stepref{s6eNsign}, $C^2=14C+32I_{63}$; the diagonal entry
   $C^2(a,a)=\sum_bC(a,b)^2$ counts the nonzero entries of row $a$, since the entries of $C$ lie
   in $\{0,\pm1\}$ and $C(a,a)=0$ (Definition~\ref{def:sec3-signing}).~\qed\end{proof}
  \pflet{$\Zt=\delta\sigma(a)$ with $a\in I$ and $\delta=\pm1$.}
  \step{s6eNone}{For $b\in I$, exactly one of $\pm\sigma(b)$ lies in $N_\Zt(R)$ if
   $C(a,b)\neq0$, and neither does otherwise.}
  \begin{proof}\pf\ $(\pm\sigma(b),\Zt)=\pm\delta C(a,b)$ for $b\neq a$, and
   $(\pm\sigma(a),\Zt)=\pm2$.~\qed\end{proof}
  \qedstep
  \begin{proof}\pf\ Every vector of $\Phi(R)$ is $\pm\sigma(b)$ for some $b\in I$; so
   $|N_\Zt(R)|=32$ by \stepref{s6eNone} and \stepref{s6eNrow}. (In Lean the leaf may instead
   transport \leancite{LineSystem.E₇.card_N}{\texttt{LineSystem.E\ensuremath{{}_7}.card\_N}}
   of C20 through the isomorphism, which preserves the sets $N_\Zt$.)~\qed\end{proof}
 \end{proof}
 \step{s6einj}{The map $(\Zr,\Zr')\mapsto(\Zr+\Zr',\Zr)$ is an injection of $Q$ into
  $T=\{(\Zt,\Zu):\Zt\in\Phi(R)\diff\Phi,\ \Zu\in N_\Zt(R)\}$.}
 \begin{proof}\pf\ It maps into $T$ by \stepref{s6ethird}, and $(\Zr,\Zr')$ is recovered from its
  image $(\Zt,\Zu)$ as $(\Zu,\Zt-\Zu)$.~\qed\end{proof}
 \step{s6eT}{$|T|=14\cdot32=448$.}
 \begin{proof}\pf\ Steps \stepref{s6ecounts} and \stepref{s6eNt}.~\qed\end{proof}
 \qedstep
 \begin{proof}\pf\ By \stepref{s6eQ}, \stepref{s6einj} and \stepref{s6eT}, $3136\leq448$, which
  is false.~\qed\end{proof}
\end{proof}

\begin{remark}
The same count excludes the type $\typ A_{n-1}$ without Lemma~\ref{lem:6-orthogonal}: then
$\cK=\cP$ has $\frac12n(n-1)$ lines, which is all of $\typ A_{n-1}$, so $\Phi(R)\diff\Phi$ is
empty while $Q$ is not. It does not exclude $\typ D_{n-1}$, nor the component of type
$\typ E_8$ when $n=16$ (there $|Q|=128\cdot28$ and $|T|=112\cdot56$). We checked the counts of
$\typ E_7$ (126 vectors, $|N_\Zt|=32$ for each) numerically in model (ii) of
Definition~\ref{def:3.3}, in a scratch script outside the trust base.
\end{remark}

\begin{cgss}{Theorem}{6.2}\label{thm:6.2}
 \lean{Matrix.IsHadamard.not_fourRow_agree_sixteenths}
 \uses{def:6-pairs,lem:6-small,lem:6-lines,lem:6-components,lem:6-decomposition,lem:6-orthogonal,
 lem:6-E7,thm:3.5,in6:sec3}
 \assume{\begin{pfenum}
  \item $H$ is an $n\times n$ Hadamard matrix, with $n>2$.
  \item For all distinct rows $i,j,k,l$, $16\agr(i,j,k,l)\in\{n,2n,3n\}$; that is, any four rows
   agree in $\frac n{16}$, $\frac n8$ or $\frac{3n}{16}$ places.
 \end{pfenum}}
 \prove{FALSE. Thus there is no Hadamard matrix of order $n>2$ for which any four rows agree in
 $\frac n{16}$, $\frac n8$ or $\frac{3n}{16}$ places.}
\end{cgss}
\begin{proof}
 \pfsketch\ We follow the paper: the $\frac12n(n-1)$ lines lie at $60^\circ$ and $90^\circ$ in
 $\jp$, and the $n-1$ lines through a fixed row are mutually orthogonal, which the systems
 $\typ A_{n-1}$ and $\typ D_{n-1}$ do not allow for odd $n-1$. \repair\ The small orders are
 Lemma~\ref{lem:6-small} (errata, Gap 10), and the indecomposability that the paper asserts is
 replaced by Lemmas~\ref{lem:6-components}, \ref{lem:6-decomposition} and~\ref{lem:6-E7}
 (errata, Gap 17). With a single component, the star closure is $\typ A_{n-1}$ or
 $\typ D_{n-1}$ by Theorem~\ref{thm:3.5}, as the paper says; the paper's appeal to $n>9$ is
 replaced by the count of lines.
 \step{s62sixteen}{$16\mid n$ and $n\geq16$.}
 \begin{proof}\pf\ Lemma~\ref{lem:6-small}, with \asm[1]{0} and \asm[2]{0}.~\qed\end{proof}
 \step{s62cases}{Alternative 1 or alternative 2 of Lemma~\ref{lem:6-decomposition} holds.}
 \begin{proof}\pf\ Lemma~\ref{lem:6-decomposition}, with \stepref{s62sixteen} and
  \asm[2]{0}.~\qed\end{proof}
 \step{s62single}{\case{$\cP$ is a single component, of type $\typ A_{n-1}$ or
  $\typ D_{n-1}$.}}
 \begin{proof}
  \pflet{$R=R_\cP$, $f:\Phi(R)\to\Phi(M)$ the isomorphism of Lemma~\ref{lem:6-components},
   part 5, onto the model $M$, and $F=\{f(\Zr_{1j}):2\leq j\leq n\}$.}
  \step{s62F}{$F$ has $n-1$ elements, and any two distinct elements of $F$ are orthogonal.}
  \begin{proof}\pf\ The $\Zr_{1j}$ are distinct and pairwise orthogonal vectors of
   $\Phi(\cP)\sbs\Phi(R)$ by Lemma~\ref{lem:6-lines}, parts 3 and 6, and
   Lemma~\ref{lem:6-components}, part 4; $f$ is injective and preserves inner products.~\qed
  \end{proof}
  \step{s62odd}{$n-1\geq2$ is odd.}
  \begin{proof}\pf\ $n\geq16$ is even by \stepref{s62sixteen}.~\qed\end{proof}
  \qedstep
  \begin{proof}\pf\ Lemma~\ref{lem:6-orthogonal}, with $\ell=n-1$, \stepref{s62F} and
   \stepref{s62odd}, excludes $M=\typ A_{n-1}$, and gives $n-1$ even if $M=\typ D_{n-1}$,
   against \stepref{s62odd}.~\qed\end{proof}
 \end{proof}
 \step{s62E}{\case{$n=16$, and some component has type $\typ E_7$.}}
 \begin{proof}\pf\ With $n=16$, \asm[2]{0} says $\agr(i,j,k,l)\in\{1,2,3\}$; apply
  Lemma~\ref{lem:6-E7}.~\qed\end{proof}
 \qedstep
 \begin{proof}\pf\ The two cases are those of \stepref{s62cases}.~\qed\end{proof}
\end{proof}

\begin{remark}
 \lean{Computations.C16.h2}\lean{Computations.C16.h4}\lean{Computations.C16.h8}
 \lean{Computations.C16.h12}\lean{Computations.C16.h16}\lean{Computations.C16.h2_isHadamard}
 \lean{Computations.C16.h4_isHadamard}\lean{Computations.C16.h8_isHadamard}
 \lean{Computations.C16.h12_isHadamard}\lean{Computations.C16.h16_isHadamard}
 \lean{Computations.C16.checkFourRowAgreement_sound}\lean{Computations.C16.fourRow_agreements}
 \lean{Computations.C16.h2_fourRowAgreement}\lean{Computations.C16.h4_fourRowAgreement}
 \lean{Computations.C16.h8_fourRowAgreement}\lean{Computations.C16.h12_fourRowAgreement}
 \lean{Computations.C16.h16_fourRowAgreement}
 \lean{Computations.C16.h12_transpose_hadamardPairs_mul}
 \lean{Computations.C16.h12_root_mem_hyperplane}\lean{Computations.C16.h12_hyperplane_finrank}
 \lean{Computations.C16.h12HyperplaneIsometry}\lean{Computations.C16.h12_lines_injective}
 \lean{Computations.C16.h12_numLines}\lean{Computations.C16.h12_line_angles}
 \lean{Computations.C16.h12_inner_values_attained}
 \uses{def:6-pairs,hyp:6.1,lem:6-gram,lem:6-P,lem:6-fixedrow}
Computation C16 (\path{LineGraphs/Computations/C16.lean}) concerns the Sylvester matrices of
orders 2, 4, 8 and 16 and the Paley matrix $H_{12}$ of order 12, given as integer matrices in
\path{LineGraphs/Computations/C16Data.lean}. We see that each of the five is a Hadamard matrix,
by \texttt{decide +kernel}. A Boolean checker counts the agreements of every four distinct rows,
and \path{Computations.C16.checkFourRowAgreement_sound} proves that its answer gives
Hypothesis~\ref{hyp:6.1}; thus $H_{12}$ satisfies Hypothesis 6.1 with $x=4$, and the Sylvester
matrix of order $n$ satisfies it with $x=n$ (for $n=2$ the condition is vacuous, since there
are no four distinct rows). For $H_{12}$, the formula of Lemma~\ref{lem:6-P}, part 1, is
$P^TP=72I-6J$, a case of \path{Matrix.IsHadamard.transpose_hadamardPairs_mul}. The vectors
$\sqrt{2/12}\,\Zp_\pi$, for the 66 pairs $\pi\in\cP$, span 66 distinct lines in $\jp$; this
hyperplane has dimension 11 by Lemma~\ref{lem:6-fixedrow} and is isometric to $\re^{11}$, and
any two of the lines are at $90^\circ$ or at $\arccos(1/3)$, both angles occurring. We note
that computation C16 does not prove the uniqueness of $H_{12}$ or Norman's theorem.
\end{remark}

\begin{remark}
 \lean{Matrix.sylvester}\lean{Matrix.sylvesterKronecker}
 \lean{Matrix.isHadamard_sylvester}\lean{Matrix.sylvester_succ}
 \lean{Matrix.sylvesterKronecker_eq_sylvester}\lean{Matrix.sylvester_fourRowAgreement}
 \lean{Matrix.hadamardRowProduct_sylvester_sylvesterEquiv}
 \lean{Matrix.hadamardRowProduct_sylvester}
 \lean{Matrix.exists_hadamardRowProduct_sylvester_eq}
 \lean{Matrix.image_hadamardRowProduct_sylvester}
 \lean{Matrix.card_image_hadamardRowProduct_sylvester}
 \lean{Matrix.two_mul_card_filter_hadamardRowProduct_sylvester}
 \lean{Matrix.inner_hadamardRowProduct_sylvester_of_ne}
 \lean{Matrix.sylvester_lines}
 \lean{Computations.C16.h2_eq_sylvester}\lean{Computations.C16.h4_eq_sylvester}
 \lean{Computations.C16.h8_eq_sylvester}\lean{Computations.C16.h16_eq_sylvester}
The paper says that the only known examples are the Hadamard matrices of Sylvester type, with
$n=2^s$ and $x=n$, and the unique $H_{12}$, with $n=12$ and $x=4$; and that the Sylvester case
corresponds to $d=1$ and yields $n-1$ mutually orthogonal lines, each line occurring $\frac12n$
times. The claims about the Sylvester matrices are side claims, and we prove them for every $s$
in \path{LineGraphs/SideClaims/Sylvester.lean}; the words ``only known'' are history, and we do
not formalize them. The Sylvester matrix $H$ of order $n=2^s$ has its rows and columns indexed by
$\{0,1\}^s$, with $H(a,b)=(-1)^{a\cdot b}$; it is the Kronecker product of $s$ copies of the
matrix $H_2$ with rows $(1,1)$ and $(1,-1)$ (\path{Matrix.sylvesterKronecker_eq_sylvester}), and
it is a Hadamard matrix (\path{Matrix.isHadamard_sylvester}). The Hadamard product of rows $a$
and $b$ is row $a+b$, with addition modulo 2
(\path{Matrix.hadamardRowProduct_sylvester_sylvesterEquiv}); in Lean the rows are also indexed
by $\{0,\dots,n-1\}$ through binary digits, and then the product of rows $i$ and $j$ is the row
of their bitwise exclusive or (\path{Matrix.hadamardRowProduct_sylvester}). Thus
the rows of $P$ are the $n-1$ rows of $H$ other than the all-ones row, each of them is the
product of exactly $\frac12n$ pairs of rows, and distinct ones are orthogonal; for the lines
that the rows of $P$ span, this is \path{Matrix.sylvester_lines}. By Lemma~\ref{lem:6-gram},
part 3, four distinct rows $i,j,k,l$ then have $8\agr(i,j,k,l)=n+(\Zp_{ij},\Zp_{kl})$, which is
$n$ or $2n$, and $H$ satisfies Hypothesis~\ref{hyp:6.1} with $x=n$
(\path{Matrix.sylvester_fourRowAgreement}); so $d=n/x=1$. The four Sylvester matrices of
computation C16 are these matrices for $s=1,2,3,4$ (\path{Computations.C16.h16_eq_sylvester} and
the three others).
\end{remark}
\endgroup


\begin{thebibliography}{CGSS}
\bibitem[CGSS]{CGSS} P. J. Cameron, J. M. Goethals, J. J. Seidel and E. E. Shult, Line graphs,
  root systems, and elliptic geometry, \emph{J. Algebra} \textbf{43} (1976), 305--327.
\bibitem[BCN]{BCN} A. E. Brouwer, A. M. Cohen and A. Neumaier, \emph{Distance-Regular Graphs},
  Springer, 1989.
\bibitem[HC]{HC} M. Hall, Jr. and W. S. Connor, An embedding theorem for balanced incomplete
  block designs, \emph{Canad. J. Math.} \textbf{6} (1954), 35--41.
\bibitem[L95]{L95} L. Lamport, How to write a proof, \emph{Amer. Math. Monthly} \textbf{102}
  (1995), 600--608.
\end{thebibliography}

\end{document}
